% Options for packages loaded elsewhere
\PassOptionsToPackage{unicode}{hyperref}
\PassOptionsToPackage{hyphens}{url}
\documentclass[
  british,
  a4paper,
  openany]{book}
\usepackage{xcolor}
\usepackage[margin=1in]{geometry}
\usepackage{amsmath,amssymb}
\setcounter{secnumdepth}{5}
\usepackage{iftex}
\ifPDFTeX
  \usepackage[T1]{fontenc}
  \usepackage[utf8]{inputenc}
  \usepackage{textcomp} % provide euro and other symbols
\else % if luatex or xetex
  \usepackage{unicode-math} % this also loads fontspec
  \defaultfontfeatures{Scale=MatchLowercase}
  \defaultfontfeatures[\rmfamily]{Ligatures=TeX,Scale=1}
\fi
\usepackage{lmodern}
\ifPDFTeX\else
  % xetex/luatex font selection
  \setmainfont[]{Noto Serif}
  \setmonofont[]{Noto Sans Mono}
\fi
% Use upquote if available, for straight quotes in verbatim environments
\IfFileExists{upquote.sty}{\usepackage{upquote}}{}
\IfFileExists{microtype.sty}{% use microtype if available
  \usepackage[]{microtype}
  \UseMicrotypeSet[protrusion]{basicmath} % disable protrusion for tt fonts
}{}
\makeatletter
\@ifundefined{KOMAClassName}{% if non-KOMA class
  \IfFileExists{parskip.sty}{%
    \usepackage{parskip}
  }{% else
    \setlength{\parindent}{0pt}
    \setlength{\parskip}{6pt plus 2pt minus 1pt}}
}{% if KOMA class
  \KOMAoptions{parskip=half}}
\makeatother
\usepackage{longtable,booktabs,array}
\newcounter{none} % for unnumbered tables
\usepackage{calc} % for calculating minipage widths
% Correct order of tables after \paragraph or \subparagraph
\usepackage{etoolbox}
\makeatletter
\patchcmd\longtable{\par}{\if@noskipsec\mbox{}\fi\par}{}{}
\makeatother
% Allow footnotes in longtable head/foot
\IfFileExists{footnotehyper.sty}{\usepackage{footnotehyper}}{\usepackage{footnote}}
\makesavenoteenv{longtable}
% definitions for citeproc citations
\NewDocumentCommand\citeproctext{}{}
\NewDocumentCommand\citeproc{mm}{%
  \begingroup\def\citeproctext{#2}\cite{#1}\endgroup}
\makeatletter
 % allow citations to break across lines
 \let\@cite@ofmt\@firstofone
 % avoid brackets around text for \cite:
 \def\@biblabel#1{}
 \def\@cite#1#2{{#1\if@tempswa , #2\fi}}
\makeatother
\newlength{\cslhangindent}
\setlength{\cslhangindent}{1.5em}
\newlength{\csllabelwidth}
\setlength{\csllabelwidth}{3em}
\newenvironment{CSLReferences}[2] % #1 hanging-indent, #2 entry-spacing
 {\begin{list}{}{%
  \setlength{\itemindent}{0pt}
  \setlength{\leftmargin}{0pt}
  \setlength{\parsep}{0pt}
  % turn on hanging indent if param 1 is 1
  \ifodd #1
   \setlength{\leftmargin}{\cslhangindent}
   \setlength{\itemindent}{-1\cslhangindent}
  \fi
  % set entry spacing
  \setlength{\itemsep}{#2\baselineskip}}}
 {\end{list}}
\usepackage{calc}
\newcommand{\CSLBlock}[1]{\hfill\break\parbox[t]{\linewidth}{\strut\ignorespaces#1\strut}}
\newcommand{\CSLLeftMargin}[1]{\parbox[t]{\csllabelwidth}{\strut#1\strut}}
\newcommand{\CSLRightInline}[1]{\parbox[t]{\linewidth - \csllabelwidth}{\strut#1\strut}}
\newcommand{\CSLIndent}[1]{\hspace{\cslhangindent}#1}
\ifLuaTeX
\usepackage[bidi=basic,shorthands=off]{babel}
\else
\usepackage[bidi=default,shorthands=off]{babel}
\fi
\ifPDFTeX
\else
\babelfont{rm}[]{Noto Serif}
\fi
\ifLuaTeX
  \usepackage{selnolig} % disable illegal ligatures
\fi
\setlength{\emergencystretch}{3em} % prevent overfull lines
\providecommand{\tightlist}{%
  \setlength{\itemsep}{0pt}\setlength{\parskip}{0pt}}
\usepackage{array}
\usepackage{tabularx}
\usepackage{morewrites}
\usepackage{imakeidx}
\indexsetup{level=\relax,noclearpage}
\setcounter{tocdepth}{1}
\setcounter{secnumdepth}{1}
\renewcommand{\sectionmark}[1]{\markright{\thesection.\ #1}}
% Predicted (counterfactual) form marker: superscript dagger immediately before the form
\providecommand{\Pred}[1]{\textsuperscript{\dag}\emph{#1}}
% Reconstructed lexical form: asterisk + italic (asterisk supplied by macro)
\providecommand{\Recon}[1]{\emph{*#1}}
% U+2192 RIGHTWARDS ARROW — Noto Serif lacks this glyph; remap to math-mode arrow
{\catcode`→=\active \global\def→{\ensuremath{\rightarrow}}}
\catcode`→=\active
\makeindex[name=oe,title={Old English},columns=3]
\makeindex[name=pgmc,title={Proto-Germanic},columns=3]
\makeindex[name=pnwgmc,title={Proto-Northwest Germanic},columns=3]
\makeindex[name=pwgmc,title={Proto-West Germanic},columns=3]
\makeindex[name=preoe,title={Pre-Old English / prehistoric English},columns=3]
\makeindex[name=on,title={Old Norse},columns=3]
\makeindex[name=ohg,title={Old High German},columns=3]
\makeindex[name=ofris,title={Old Frisian},columns=3]
\makeindex[name=goth,title={Gothic},columns=3]
\makeindex[name=os,title={Old Saxon},columns=3]
\makeindex[name=dutch,title={Dutch},columns=3]
\makeindex[name=german,title={German},columns=3]
\makeindex[name=me,title={Middle English},columns=3]
\makeindex[name=modeng,title={Modern English},columns=3]
\makeindex[name=mlg,title={Middle Low German},columns=3]
\usepackage{fvextra}
\providecommand{\texorpdfstring}[2]{#1}
\DeclareRobustCommand{\CAPRRuleHeading}[2]{%
  #1\texorpdfstring{\\(#2)}{ (#2)}%
}
\DeclareRobustCommand{\CAPRHeadingBreak}{%
  \texorpdfstring{\\}{ }%
}
\DefineVerbatimEnvironment{ReaderFacingFoma}{Verbatim}{
  breaklines=true,
  breakanywhere=true,
  fontsize=\small,
  frame=single,
  framesep=2mm
}
% Predicted (counterfactual) form marker: superscript dagger immediately before the form
\providecommand{\Pred}[1]{\textsuperscript{\dag}\emph{#1}}
% Reconstructed lexical form: asterisk + italic (asterisk supplied by macro)
\providecommand{\Recon}[1]{\emph{*#1}}
% U+2192 RIGHTWARDS ARROW — Noto Serif lacks this glyph; remap to math-mode arrow
{\catcode`→=\active \global\def→{\ensuremath{\rightarrow}}}
\catcode`→=\active
\usepackage{bookmark}
\IfFileExists{xurl.sty}{\usepackage{xurl}}{} % add URL line breaks if available
\urlstyle{same}
% fallback for those not using the hyperref driver hyperxmp:
\makeatletter
\@ifundefined{xmpquote}{\newcommand{\xmpquote}[1]{#1}}{}
\makeatother
\hypersetup{
  pdflang={en-GB},
  hidelinks,
  pdfcreator={LaTeX via pandoc}}

\author{}
\date{}

\begin{document}
\frontmatter

{
\setcounter{tocdepth}{0}
\tableofcontents
}
\mainmatter
\mainmatter

\chapter{Introduction}\label{introduction}

\section{From sound law to
derivation}\label{from-sound-law-to-derivation}

Historical linguists ordinarily test an etymology by carrying a
reconstructed form through the sound changes that separate it from its
alleged reflex. Most such derivations remain implicit. A scholar knows
that Proto-Germanic \emph{*p} yields Old English \emph{f}, that West
Germanic \emph{*z} became \emph{r} under the appropriate conditions, and
that one change must have preceded another because the reverse order
produces the wrong form. For a single word this mental arithmetic
presents little difficulty. Across hundreds of words and scores of
interacting changes it becomes treacherous. Each step may look familiar
while the derivation as a whole is false.

In computer-assisted phonological reconstruction (CAPR), I represent
each sound law as a finite-state transducer and compose the transducers
in the order argued for below. The resulting cascade applies every
change without discretion: the same rule, with the same environment,
reaches every eligible form. A proposed reconstruction either yields the
Old English comparator or it does not. Moving a rule may repair one
derivation and destroy another. Both results deserve attention, for a
failure locates the point at which phonology ceases to suffice and
philology must decide among analogy, borrowing, dialect mixture,
uncertain attestation, or a faulty reconstruction.

No new principle is involved. The Neogrammarians already demanded
exceptionless sound laws and an account of apparent exceptions
(\citeproc{ref-OsthoffBrugmann1881}{Osthoff and Brugmann 1878-\/-1910}).
Their formula was the \emph{Ausnahmslosigkeit der Lautgesetze}.
Formalization merely enforces that demand mechanically: every eligible
form bears the stated consequences. The computer cannot decide which
reconstruction is historically defensible, but it prevents the
investigator from forgetting what that reconstruction entails.

I apply this method to the development of Proto-Germanic and early West
Germanic forms into Old English. Germanic makes a severe test. Its
historical grammar rests on two centuries of philological labor, while
Old English offers abundant but orthographically and dialectally varied
testimony (\citeproc{ref-Campbell1959}{Campbell 1959};
\citeproc{ref-Hogg1992}{Hogg 1992}; \citeproc{ref-RingeTaylor2014}{Ringe
and Taylor 2014}; \citeproc{ref-Fulk2018}{Fulk 2018}). A computational
account cannot plead scarcity of evidence. It must reproduce familiar
developments, identify the evidence for their order, and explain why its
input sometimes differs from the headword printed in an etymological
dictionary.

\section{The formal claim}\label{the-formal-claim}

A sound change defines a relation between strings. For each etymology I
supply a reconstructed input to an ordered series of such relations and
compare the result with an Old English form. Kaplan and Kay demonstrated
that the familiar rewrite rules of phonology admit a finite-state
interpretation (\citeproc{ref-KaplanKay1994}{Kaplan and Kay 1994}); Foma
gives that interpretation executable form
(\citeproc{ref-Hulden2009}{Hulden 2009}). A toy rule deleting final
\emph{*z} may be written:

\begin{ReaderFacingFoma}
define ToyFinalZLoss [{*z} -> 0 || _ .#.];
\end{ReaderFacingFoma}

The notation states a target, an output, and an environment. Actual
rules also refer to segment classes, stress, syllable weight, word
boundaries, and the output of earlier changes. A loose prose formulation
can conceal choices that code must resolve. Does the rule affect every
vowel or only short vowels? Does a following consonant block it? Does it
apply before or after the loss of a final syllable? A transducer forces
each question into the open.

Order gives the cascade its historical content. An early change may
create the environment for a later one and thus feed it; alternatively,
it may remove that environment and bleed it. A word affected by both
changes can therefore establish their relative chronology. Other changes
commute across the available lexicon: reversing them alters no checked
output. Such a result does not prove simultaneity or historical
indifference. It shows only that the present evidence fails to order
them. Throughout this book I distinguish chronology compelled by the
tested forms from placement adopted on the authority of the historical
grammars.

Backward reconstruction requires a further restriction. An unrestricted
inverse transducer will propose formally possible strings that no
Germanic language could have inherited. I therefore restrict backward
reconstruction with an inventory and a statement of permissible
ancestral forms. Reconstruction always combines correspondences with a
theory of what could count as a word in the \emph{Grundsprache}; this
restriction states that theory rather than leaving it tacit.

\section{Inputs, targets, and success}\label{inputs-targets-and-success}

I compare a selected earlier Germanic form with a selected Old English
target. Neither selection is innocent. Dictionaries cite lexemes, but
sound change operates on word-forms. The ancestor of an Old English
plural, preterite, or oblique case may differ from the reconstructed
lemma in precisely the material on which a later sound law acts. Kroonen
and Orel provide indispensable lexical reconstructions, while the
grammars often supply the paradigm history needed to choose the actual
input (\citeproc{ref-Orel2003}{Orel 2003};
\citeproc{ref-Kroonen2013}{Kroonen 2013};
\citeproc{ref-RingeTaylor2014}{Ringe and Taylor 2014}).

I therefore distinguish the citation reconstruction from the transducer
input. The former identifies the etymon; the latter represents the
paradigm cell or remodeled stem whose history is at issue. Where the two
differ, the lexical entry states the difference and argues for it. This
prevents a convenient input from masquerading as a received
reconstruction. An unmotivated alteration made solely to secure the
desired output would empty the exercise of historical meaning.

The distinction is concrete in Old English
{\emph{sċuldrum}}\index[iv]{02oe@\textbf{Old English}!sculdrum@\emph{sċuldrum}}
`shoulders, dative plural'. Its inflectional history requires an input
of the type
\Recon{skúldramiz}\index[iv]{03pgmc@\textbf{Proto-Germanic}!skuldramiz@*skúldramiz}
`shoulders, reconstructed dative-plural input'; the singular headword
represents a different paradigm cell.

The target also requires judgment. Old English spelling varies by date,
dialect, manuscript, and editorial practice. The cascade may produce an
internal phonological symbol that an orthographic transducer then maps
to a normalized written form. A string match at this final stage cannot
by itself establish an etymology; conversely, a superficial spelling
mismatch need not disprove one. I treat phonological development and
orthographic normalization as separate operations so that exact
computation does not confuse notation with history.

Within these limits a successful derivation has three senses. It
succeeds formally when the output string matches the target. It succeeds
philologically when the chosen input and comparator are the proper forms
to compare. It succeeds historically when the proposed path agrees with
the wider Germanic evidence. The first kind of success is cheap. The
argument of the book concerns the conjunction of all three.

\section{The evidence of failure}\label{the-evidence-of-failure}

Regular sound change makes irregularity legible. If an inherited form
refuses to pass through an otherwise successful cascade, the mismatch
demands a name. Analogy may have replaced the expected reflex with a
form drawn from another paradigm cell. Borrowing may have introduced the
word after the relevant changes. A dialectal form may lie outside the
modeled West Saxon path. The target may be late, corrupt, or normalized
beyond what the manuscript evidence warrants. Finally, either the
reconstruction or the rule may be wrong.

These possibilities should not be suppressed by narrow, lexeme-specific
``sound laws.'' A system that derives every target by multiplying
exceptions has only encoded its answers. I instead distinguish several
classes of non-regular result: attested variants, early and late
analogy, reconstructed Old English comparators, known but unmodelled
phonological and morphological developments, and unexplained exceptions.
The categories are claims, not housekeeping labels. They identify where
regular phonology ends and what additional history the evidence
requires.

This treatment follows the original Burmish CAPR work, in which
resistant forms often disclosed loans or mistaken cognate assignments.
Old English shifts the balance toward morphology and analogy, but the
methodological advantage remains the same. Failure concentrates inquiry.
It tells us which assumption---input, target, environment, order, or
lexical history---must bear the explanation.

\section{Evidence for relative
chronology}\label{evidence-for-relative-chronology}

The chronology chapters combine three kinds of evidence. First come the
statements of the standard historical grammars. These establish the
received description and often the broad order of developments
(\citeproc{ref-Campbell1959}{Campbell 1959};
\citeproc{ref-Hogg1992}{Hogg 1992}; \citeproc{ref-RingeTaylor2014}{Ringe
and Taylor 2014}; \citeproc{ref-Fulk2018}{Fulk 2018}). Second come
individual witness words. A derivation that succeeds under one order and
fails under the reverse order supplies direct lexical evidence for that
relation. Third come exhaustive order tests across the active dataset.
These reveal whether an apparently decisive relation is local, whether
other words contradict it, and how far a rule can move without
disturbing any output.

Suppose that a consonant change creates the environment for a later
vowel change. Under the received order both apply and the Old English
target emerges. Reverse them and the vowel change misses its
environment. The word then supports the priority of the consonant
change. The converse case is equally informative: if an early rule
creates a segment that a later rule would wrongly alter, the creating
rule must follow the other. When every checked form remains unchanged
under reversal, the lexical evidence leaves the order open even if
philological considerations still favor one placement.

A note on notation: a superscript dagger placed immediately before an
italicized form marks a concrete predicted output generated by a
counterfactual manipulation of the rule order (for example, moving a
change earlier or later than argued). This dagger indicates that the
marked form is a model prediction, not an independently attested
historical form; reconstructions remain starred in the usual way.

Sims-Williams argues for mechanizing precisely this kind of reasoning
(\citeproc{ref-SimsWilliams2018}{Sims-Williams 2018}). Computation does
not replace the historical argument; it makes the extent of that
argument measurable. A traditional chronology may prove correct but less
tightly constrained than its customary presentation suggests. A relation
described as local may in fact rest only on a broad terminus. Such
negative results are salutary. They separate what the data demonstrate
from what a convenient exposition merely presupposes.

\section{Rules and words}\label{rules-and-words}

The book accordingly moves twice through the same history. Part I begins
with the ordered rules. For each development it states the historical
problem, gives the rule in formal notation, and examines the evidence
for placement. The code appears because it is part of the claim, but the
code-name is never an explanation. \texttt{OEIUmlaut}, for example, is
only a rule name; its linguistic content lies in the stated environment,
the historical discussion, and the words whose derivations depend upon
it.

Part II begins with the words. Each entry identifies the reconstruction,
the selected input, the Old English comparator, and the derivational
class. A trace shows which rules altered the form; the accompanying
prose explains the philological decisions that the trace cannot make.
Short entries record ordinary derivations. Longer ones treat doubtful
reconstructions, variant attestations, paradigm-cell selection, or
analogy. The asymmetry is intentional: equal database rows do not
warrant equal historical discussion.

The reader can thus move in either direction. A chronology chapter names
the lexical witnesses that constrain a rule; their entries display the
complete derivations. A lexical entry invokes a change; the
corresponding chapter explains its formulation and place in the cascade.
The index verborum provides a third route through the material,
gathering reconstructed and attested forms by language.

\section{Reproducibility and
disagreement}\label{reproducibility-and-disagreement}

An executable derivation identifies the exact point of disagreement. One
reader may accept a sound change but reject its environment; another may
accept the rule and dispute its order; a third may object to the
selected paradigm cell or to the normalization of the Old English
target. Each objection addresses a recorded decision. Given the same
inputs, rules, and order, the stated outputs follow. Reproducibility
here concerns those consequences, not the surrender of philological
judgment to an algorithm.

Prose can sound settled while concealing several incompatible
derivations. Code is less accommodating. A rule that is too broad
damages words it was never designed to explain; one that is too narrow
leaves legitimate reflexes untouched. A reconstruction plausible in
isolation may fail after three later changes. The resulting
embarrassment is a virtue of the method. Applying every rule to the
whole lexicon subjects local analyses to a global test.

The converse danger is false precision. A deterministic cascade may
tempt the reader to mistake exact strings for exact history. Every
result still depends on choices about segmentation, symbol inventory,
reconstruction, morphology, dialect, chronology, and orthography. I
therefore cite the philological sources, display the selected forms,
preserve unresolved exceptions, and distinguish tested order from
source-based order. Those choices remain accountable to philology.

\section{The argument}\label{the-argument}

A book must advance an argument. Mine is that traditional rule-based
reconstruction becomes clearer when every proposed derivation can be
executed, every rule must face the whole lexicon, and every failure is
reported. The Germanic-to-Old-English case shows both the power and the
boundary of that claim. Much of the history admits a coherent ordered
account. The residue does not disappear: it resolves into morphology,
analogy, variation, borrowing, imperfect attestation, and a small number
of genuine problems.

The sound-change chapters ask how tightly the lexical evidence
constrains the cascade. The lexical chapters ask whether individual
etymologies survive that cascade without philological sleight of hand.
Neither inquiry suffices alone. A formal system that reaches the targets
by hiding doubtful inputs is bad historical linguistics; elegant prose
that never reproduces its derivations is an incomplete formal account.
Executable derivations make those two standards answer to one another.

\part{Sound changes, formalization, and relative chronology}

\chapter*{Sound-change overview}\label{sound-change-overview}
\addcontentsline{toc}{chapter}{Sound-change overview}

\section*{Scope and orientation}\label{scope-and-orientation}
\addcontentsline{toc}{section}{Scope and orientation}

The sequence begins with early West Germanic consonant and vowel changes
and ends with Old English r-metathesis.

Rhotacism, brightening, breaking, umlaut, and apocope alternate with
narrowly conditioned changes whose relative order rests on particular
witness words.

The evidence ranges from broadly attested sound laws to lexical
constraints that establish only one chronological boundary.

\section*{Numbering note}\label{numbering-note}
\addcontentsline{toc}{section}{Numbering note}

SC numbers remain the established legacy identifiers. The Version 1 book
presents the changes in historical chapter order, which differs from the
computational cascade order for several rules.

SC038, SC062, and SC084 mark technical or prosodic stages rather than
sound changes; SC077 is unused.

\newpage

\chapter{From Proto-Germanic to Proto-Northwest
Germanic}\label{from-proto-germanic-to-proto-northwest-germanic}

\section{Historical interval}\label{historical-interval}

This chapter covers developments that took place within the
Proto-Germanic period, from the inherited consonant system to the first
changes that separate the Northwest Germanic line from the rest of the
Germanic family.

The reconstruction labelled Proto-Germanic here is the common ancestor
of Gothic, North Germanic, and West Germanic, reconstructed through the
classical comparative method from attested descendant languages. The
label Proto-Northwest Germanic designates the hypothetical node linking
the ancestors of North Germanic (Old Norse and its relatives) and West
Germanic (Old English, Old High German, Old Saxon, and Old Frisian among
others), to the exclusion of Gothic and the other East Germanic
varieties.

\section{What this chapter contains}\label{what-this-chapter-contains}

Chapter 1 contains one reader-facing sound-change section: the
positional allophony of Proto-Germanic \emph{*b}, implemented as
\texttt{SC049\ PGmcBAllophony}. The rule governs the distribution of
\emph{*b} as a stop versus a voiced bilabial fricative \emph{*β}
depending on syllabic environment. Hogg, Ringe and Taylor, and Luick all
identify this distribution as a Proto-Germanic feature
(\citeproc{ref-Hogg1992}{Hogg 1992, 101--2};
\citeproc{ref-RingeTaylor2014}{Ringe and Taylor 2014, 121};
\citeproc{ref-Luick1914}{Luick 1914-\/-1940, 107}).

CAPR implements this rule late in the computational cascade because the
alternation interacts with environments shaped by intermediate rule
applications. The cascade placement therefore diverges from the
historical stage; the reader-facing section notes this divergence
explicitly.

One other historically Proto-Germanic change, Gm-simplification
(\texttt{SC002\ PGmcGmSimplification}), is documented in the book-entry
plan and its literature dossier confirms the source base is narrow (two
lexical families:
\Recon{draugma-}\index[iv]{03pgmc@\textbf{Proto-Germanic}!draugma@draugma-}
`dream' and
\Recon{taugma-}\index[iv]{03pgmc@\textbf{Proto-Germanic}!taugma@taugma-}
`team'; (\citeproc{ref-Kroonen2013}{Kroonen 2013, 101, 511})). A
reader-facing section for SC002 awaits a stronger explanatory source
base and is not yet assembled in the reader-facing sequence.

\section{Scope and genealogical
context}\label{scope-and-genealogical-context}

Changes in this chapter are pan-Germanic in scope: they apply to
ancestral forms that feed both the North Germanic and West Germanic
descendants, or they represent internal Proto-Germanic processes visible
across the Germanic family.

The boundary between Proto-Germanic and Proto-Northwest Germanic is not
sharp in the textbook literature. Ringe and Taylor treat many of the
traditionally ``Proto-Germanic'' changes as part of a shared innovation
package that is diagnostically older than North--West Germanic
divergence but not necessarily earlier than the separation of the East
Germanic line (\citeproc{ref-RingeTaylor2014}{Ringe and Taylor 2014,
1--30}). For book purposes, the distinction matters primarily because it
separates the features inherited uniformly from all Germanic from those
shared selectively by North and West Germanic to the exclusion of
Gothic.

\section{A note on the rule names}\label{a-note-on-the-rule-names}

The CAPR rules implemented in this chapter carry names beginning with
\texttt{PGmc}. Those names are intended as stable internal identifiers,
not as claims about the precise historical stage of every rule so
labelled. A rule named \texttt{PGmcX} may in some cases be a later
development that affects only the West Germanic or Northwest Germanic
branch; the chapter assignment in this staging map takes priority over
the rule-name prefix for historical organization purposes.

\section{B allophony}\label{b-allophony}

\subsection{Historical discussion}\label{historical-discussion}

The positional alternation of Germanic \emph{*b} is a Proto-Germanic
distributional feature. Hogg states the Old English distribution
clearly: /b/ is a stop initially, after nasals, and in gemination, while
the same segment is otherwise realized as a voiced bilabial fricative
(\citeproc{ref-Hogg1992}{Hogg 1992, 101--2}). Ringe and Taylor support
the broader West Germanic background by treating Proto-West-Germanic
\emph{*b} as a segment whose stop and fricative values depend on
position (\citeproc{ref-RingeTaylor2014}{Ringe and Taylor 2014, 121}),
and Luick's spelling evidence shows the same labial fricative pattern in
Old English (\citeproc{ref-Luick1914}{Luick 1914-\/-1940, 107}).

The distribution is narrow, but later changes presuppose the
stop-fricative alternation. CAPR implements the rule at a late cascade
position for computational reasons: the alternation must interact with
consonant environments shaped by intermediate rule applications. Its
historical stage is Proto-Germanic.

\subsection{\texorpdfstring{\CAPRRuleHeading{SC049. Distribution of \emph{*b} after vowels and liquids}{PGmcBAllophony}}{}}\label{rule-PGmcBAllophony}

\begin{ReaderFacingFoma}
define PGmcBAllophony [
    {*b} -> {*β} || PGmcStarVocalic _,
    {*b} -> {*β} || [{*l} | {*r}] _
] .o. [
    {*β} -> {*b} || _ {*b}
];
\end{ReaderFacingFoma}

The handbooks describe \emph{*b}/\emph{*bb} as a positional alternation
within the consonant system, and one compound supplies its chronological
consequence. Before \hyperref[rule-OECompoundLinkingSyncope]{SC037
OECompoundLinkingSyncope}, \emph{reġnboga} `rainbow' develops as
\Pred{reġnfoga} rather than expected OE \emph{reġnboga}; later placement
creates no comparable failure. The witness places b-allophony after
compound-linking syncope without turning the alternation into an
independent sound law.

\newpage

\chapter{From Proto-Northwest Germanic to Proto-West
Germanic}\label{from-proto-northwest-germanic-to-proto-west-germanic}

\section{Historical interval}\label{historical-interval-1}

This chapter covers the sound changes that took place in the
proto-language shared by the West Germanic languages --- Old English,
Old Frisian, Old Saxon, Old High German, and Old Dutch --- before the
individual languages diverged. The starting reconstruction is
Proto-Northwest Germanic (PNWGmc), the hypothetical common ancestor of
North Germanic and West Germanic together; the ending reconstruction is
Proto-West Germanic (PWGmc), the immediate common ancestor of the West
Germanic languages specifically.

\section{Scope and internal
diversity}\label{scope-and-internal-diversity}

Changes in this chapter are not all equally pan-West-Germanic in scope.
They may be grouped broadly as follows:

\textbf{Northwest Germanic innovations} (shared by both North and West
Germanic): innovations in the unstressed vowel system, certain
final-syllable vowel changes, and selected consonant cluster
simplifications. Changes labelled \texttt{NWGmc} in the CAPR rule names
fall here, though rule prefixes are not always reliable guides to
historical scope.

\textbf{Proto-West-Germanic innovations} (shared within West Germanic
but not in North Germanic): the cluster of morphological and
phonological changes that distinguish Old English, Old High German, Old
Saxon, and Old Frisian from Old Norse. Changes labelled \texttt{PWGmc}
in the CAPR rule names generally fall here. They include early apocope
rules, certain consonant assimilations, and the West Germanic gemination
of consonants before \texttt{*j}.

\section{Major changes}\label{major-changes}

The \texttt{*ai} monophthongization (SC004) represents one of the most
pervasive shared NW--West Germanic vowel shifts, turning unstressed
endings such as the dative singular and strong-adjective plural to
longer vowels. Ringe and Taylor treat this as one of the clearest
post-PNWGmc shared developments (\citeproc{ref-RingeTaylor2014}{Ringe
and Taylor 2014, 40--41}); Fulk groups it among the North/West-Germanic
shared innovations that distinguish the period from Gothic
(\citeproc{ref-Fulk2018}{Fulk 2018, sect. 5.2}).

The West Germanic consonant changes of this chapter --- j-gemination
(SC010), early i-apocope (SC006), coronal-w assimilation (SC008), and
related rules --- represent the most productive phonological territory
for the CAPR derivations. They feed a large proportion of the
distinctive consonant clusters of Old English. Handbooks vary in exactly
how they group and name these changes
(\citeproc{ref-Campbell1959}{Campbell 1959, sects 404, 406};
\citeproc{ref-Hogg1992}{Hogg 1992, sect. 7.1}).

The nasal spirant corridor (SC026--SC027) illustrates a type of change
common in historical grammars of the ``Ingvaeonic'' or ``North Sea
Germanic'' area: nasals disappear before voiceless fricatives, with
compensatory vowel lengthening (\citeproc{ref-Campbell1959}{Campbell
1959, sects 462--463}; \citeproc{ref-Hogg1992}{Hogg 1992, sect. 7.77}).
The CAPR model splits this into two ordered steps to make the vowel
effect computationally tractable; the book prose explains that split
against the handbook tradition, which typically presents the change as a
single process.

Several changes in this chapter carry \texttt{PWGmc} labels in the CAPR
implementation but appear later in the computational cascade than their
historical stage would suggest: final bare-\texttt{*a} loss (SC041),
surviving bimoric \texttt{*ō} unrounding (SC042), and Sievers-law
syncope (SC050) are placed late in the transducer for computational
reasons. Their chapter assignment here reflects their historical stage,
not their cascade position; the individual sound-change sections note
the divergence.

\section{A note on source terminology and
subgrouping}\label{a-note-on-source-terminology-and-subgrouping}

The literature uses several partly overlapping stage labels for this
period:

\begin{itemize}
\tightlist
\item
  \textbf{Northwest Germanic}: the node uniting North and West Germanic.
\item
  \textbf{Proto-West Germanic}: the node uniting only the West Germanic
  languages.
\item
  \textbf{North Sea Germanic} or \textbf{Ingvaeonic}: a proposed
  subgroup within West Germanic covering Old English, Old Frisian, and
  Old Saxon (and sometimes Old Low Franconian), sharing certain
  innovations over a broader area.
\item
  \textbf{Anglo-Frisian}: a narrower proposed subgroup linking only Old
  English and Old Frisian.
\end{itemize}

These labels are not always used consistently across sources. Ringe and
Taylor are cautious about reconstructing a discrete Proto-West-Germanic
node (\citeproc{ref-RingeTaylor2014}{Ringe and Taylor 2014, 50--55}).
Campbell notes that many of the ``West Germanic'' shared features could
alternatively be treated as parallel developments rather than common
inheritance (\citeproc{ref-Campbell1959}{Campbell 1959, sects 1--5}).

CAPR uses \texttt{PWGmc} and \texttt{NWGmc} as organizing labels for
this chapter without claiming to have settled all questions about West
Germanic subgrouping. Changes that appear in the literature under
``Ingvaeonic'' labels but affect the Old English--to-Proto-Germanic
derivation chain are treated here as late expressions of the same West
Germanic developmental period unless existing CAPR dossier research
specifically argues for Anglo-Frisian or English-specific placement.

\section{Rule names}\label{rule-names}

The CAPR rules in this chapter carry names beginning with \texttt{NWGmc}
or \texttt{PWGmc}. These names are stable internal identifiers. A name
beginning with \texttt{NWGmc} does not guarantee that the change is
exclusive to Northwest Germanic, and a name beginning with
\texttt{PWGmc} does not guarantee that it is absent from North Germanic.
The historical analysis in each sound-change section takes priority over
the name prefix.

\section{Proto-West-Germanic
ai-monophthongization}\label{proto-west-germanic-ai-monophthongization}

\subsection{Historical discussion}\label{historical-discussion-1}

Ringe and Taylor treat the reduction of unstressed \emph{*ai} as one of
the major early vowel shifts shared across the Northwest Germanic area
(\citeproc{ref-RingeTaylor2014}{Ringe and Taylor 2014, 40--41}).

The historical support is strongest for unstressed \emph{*ai},
especially word-finally. The rule extends the change to nonfinal
\emph{*ai > *ā}, a generalization stated more sharply than in the
current handbook discussion. Both developments have inherited \emph{*ai}
as their input.

\subsection{\texorpdfstring{\CAPRRuleHeading{SC004. Proto-West-Germanic ai-monophthongization}{PWGmcAiMonophthongization}}{}}\label{rule-PWGmcAiMonophthongization}

\begin{ReaderFacingFoma}
define PWGmcAiMonophthongization [
    [{*ai} -> {*ē} || _ .#.]
    .o.
    [{*ai} -> {*ā}]
    .o.
    [{*ái} -> {*ā}]
];
\end{ReaderFacingFoma}

The soul form fixes the relation to interstress raising. If
monophthongization is delayed until after that change, PGmc
\Recon{sáiwalō} `soul' yields \Pred{sāwel} rather than expected OE
\emph{sāwol} `soul'. No earlier placement changes a checked output.

This witness proves that monophthongization preceded interstress
raising; it says nothing about the date of the wider nonfinal
\emph{*ai > *ā} generalization. The word-final merger with long mid
\emph{*ē} belongs among the early Northwest Germanic vowel shifts; the
broader chronology remains less certain.

\newpage

\section{\texorpdfstring{Unstressed \emph{*a}-raising before final
\emph{*m}}{Unstressed -raising before final }}\label{unstressed--raising-before-final}

\subsection{Historical discussion}\label{historical-discussion-2}

Campbell notes that unstressed \emph{u} is especially well preserved
before \emph{m}, with dat.pl. \emph{-um} and related endings as the
clearest evidence (\citeproc{ref-Campbell1959}{Campbell 1959, 156},
§373). Fulk likewise includes the development of early unstressed
\emph{*o} to \emph{u} before \emph{m} among the similarities shared by
North and West Germanic (\citeproc{ref-Fulk2018}{Fulk 2018, 16}, §5.2).

I restrict the change to unstressed vowels in inflectional material
because the strongest evidence concerns noninitial unstressed material
before final \emph{*m}. Final \emph{*m} conditions the raising.

\subsection{\texorpdfstring{SC005. Unstressed \emph{*a}-raising before
final \emph{*m}
(\texttt{NWGmcAToUBeforeM})}{SC005. Unstressed -raising before final  (NWGmcAToUBeforeM)}}\label{rule-NWGmcAToUBeforeM}

\begin{ReaderFacingFoma}
define NWGmcAToUBeforeM [
    {*a} -> {*u} || EnglishStarVocalic EnglishStarConsonant+ _ {*m} ({*i})? ({*z})? .#.
];
\end{ReaderFacingFoma}

Here the witness word and the comparative evidence serve different
purposes. If raising is delayed until after
\hyperref[rule-NWGmcULowering]{SC017 NWGmcULowering}, PGmc
\Recon{skúldramiz} `shoulders' yields \Pred{sċoldrum} rather than
expected OE \emph{sċuldrum} `shoulders'; earlier placements converge on
the expected output. The \texttt{shoulder} family therefore tests the
chronology, while the inflectional endings justify restricting the
change to noninitial unstressed material before \emph{*m}.

The evidence is confined to inflectional material.

\newpage

\section{Early i-apocope}\label{early-i-apocope}

\subsection{Historical discussion}\label{historical-discussion-3}

Sievers/Brunner treats the early loss of final \emph{*i} after
unstressed syllables as established by the fact that these endings no
longer trigger later i-umlaut in Old English, and Ringe and Taylor make
the same point through the pathway to \emph{geoguþ} `youth'
(\citeproc{ref-SieversBrunner1965}{Brunner 1965, sects 145--146};
\citeproc{ref-RingeTaylor2014}{Ringe and Taylor 2014, 141}). Campbell's
\emph{dugup} `troop' and \emph{geogup} `youth' examples belong to the
same pattern (\citeproc{ref-Campbell1959}{Campbell 1959, sect. 332}).

The ending vowel disappears in a weak suffixal environment early enough
to block later umlaut. This anti-umlaut chronology distinguishes the
change from later final-vowel losses.

\subsection{\texorpdfstring{SC006. Early i-apocope
(\texttt{PWGmcEarlyIApocope})}{SC006. Early i-apocope (PWGmcEarlyIApocope)}}\label{rule-PWGmcEarlyIApocope}

\begin{ReaderFacingFoma}
define PWGmcEarlyIApocope [
    {*i} -> 0 || PGmcStarStressedVowel PGmcStarConsonant+ PGmcStarVocalic PGmcStarConsonant+ _ .#.,
    {*i} -> 0 || PGmcStarStressedVowel PGmcStarConsonant+ PGmcStarVocalic PGmcStarConsonant+ _ {*z} .#.
];
\end{ReaderFacingFoma}

The absence of umlaut in \emph{geoguþ} `youth' provides the historical
argument for early deletion. The ordered derivation supplies a different
test: if apocope is delayed until after
\hyperref[rule-OEAwLongDiphthong]{SC034 OEAwLongDiphthong}, PGmc
\Recon{skáwōθi} `shows' yields \Pred{sċēaweþ} rather than expected OE
\emph{sċēawaþ} `shows'.

Early i-apocope must therefore precede the long-diphthong development.
Moving it earlier within the tested range leaves every checked output
unchanged; its early date rests on the anti-umlaut evidence, not on a
lower boundary supplied by the witness words.

\newpage

\section{\texorpdfstring{Final \emph{*ō}-lowering before
\emph{*r}}{Final -lowering before }}\label{final--lowering-before}

\subsection{Historical discussion}\label{historical-discussion-4}

Ringe and Taylor treat the West Germanic lowering of final bimoric
\emph{*ō} before word-final \emph{*r} as a specific inherited
development and illustrate it above all with the families behind
\emph{fēower} `four' and \emph{wæter} `water'
(\citeproc{ref-RingeTaylor2014}{Ringe and Taylor 2014, 58--59}).

The rule is historically secure but narrow: final or pre-final \emph{*ō}
before word-final \emph{*r}. The clearest evidence remains concentrated
in the \texttt{four} and \texttt{water} material. No broader environment
for \emph{*ō} is attested.

\subsection{\texorpdfstring{\CAPRRuleHeading{SC007. Lowering of final bimoric \emph{*ō} before \emph{*r}}{PWGmcFinalOrLowering}}{}}\label{rule-PWGmcFinalOrLowering}

\begin{ReaderFacingFoma}
define PWGmcFinalOrLowering [
    {*ō} -> {*a} || _ {*r} .#.
];
\end{ReaderFacingFoma}

OE \emph{wæter} `water' reveals why lowering must precede
\hyperref[rule-AngloFrisianBrightening]{SC043 AngloFrisianBrightening}.
If \hyperref[rule-PWGmcFinalOrLowering]{SC007 PWGmcFinalOrLowering} is
delayed until afterwards, PGmc \Recon{wátōr} `water' yields \Pred{water}
rather than expected OE \emph{wæter} `water': brightening can affect the
vowel only after lowering has created its input. Moving the change
earlier within the tested range alters no checked output.

The witness thus supplies a terminus ante quem at brightening but no
earlier boundary. The \emph{fēower} `four' and \emph{wæter} `water'
families support the narrow environment before word-final \emph{*r}; no
broader lowering of \emph{*ō} is attested.

\newpage

\section{Coronal-w assimilation}\label{coronal-w-assimilation}

\subsection{Historical discussion}\label{historical-discussion-5}

Ringe and Taylor treat the assimilation of \texttt{*dw} and \texttt{*zw}
to \texttt{*ww} as a shared Proto-West-Germanic innovation and support
it through the \texttt{four} family and plural-pronominal forms such as
\texttt{you} and \texttt{your} (\citeproc{ref-RingeTaylor2014}{Ringe and
Taylor 2014, 56--57}).

The historical support rests on a small witness set. Both coronal inputs
assimilate before \emph{*w}, but only the numeral and the pronominal
forms directly support the generalization.

\subsection{\texorpdfstring{\CAPRRuleHeading{SC008. Assimilation of coronal consonants before \emph{*w}}{PWGmcCoronalWAssimilation}}{}}\label{rule-PWGmcCoronalWAssimilation}

\begin{ReaderFacingFoma}
define PWGmcCoronalWAssimilation [
    {*d} -> {*w} || _ {*w},
    {*z} -> {*w} || _ {*w}
];
\end{ReaderFacingFoma}

OE \emph{fēower} `four' exposes a feeding relation: coronal assimilation
must create \emph{*ww} while simplification can still reduce it. If
\hyperref[rule-PWGmcCoronalWAssimilation]{SC008
PWGmcCoronalWAssimilation} is delayed until after
\hyperref[rule-OEWWSimplification]{SC031 OEWWSimplification}, PGmc
\Recon{fédwōr} `four' yields \Pred{fēowwer} rather than expected OE
\emph{fēower} `four'. Earlier placements alter no checked output.

The numeral fixes that relative order. The pronouns extend both input
clusters beyond \texttt{four}; the earlier boundary remains
undetermined.

\newpage

\section{\texorpdfstring{\emph{ij}-contraction in
\emph{friend}}{-contraction in }}\label{contraction-in}

\subsection{Historical discussion}\label{historical-discussion-6}

Ringe and Taylor describe a change of \texttt{*ijo} to \texttt{*iu} in
the \texttt{friend} family, with the pathway PGmc \emph{*frijond-}
\textgreater{} PWGmc \Recon{friund} `friend' \textgreater{} OE
\emph{frēond} `friend' (\citeproc{ref-RingeTaylor2014}{Ringe and Taylor
2014, 62}). The same source immediately warns that the \texttt{*ijo}
sequence is unique enough that wider generalization is inadvisable
(\citeproc{ref-RingeTaylor2014}{Ringe and Taylor 2014, 62}).

The change concerns a rare sequence confined to the \texttt{friend}
family and cannot safely be generalized into a broadly productive rule.

\subsection{\texorpdfstring{SC009. \emph{ij}-contraction in
\emph{friend}
(\texttt{PWGmcIjContraction})}{SC009. -contraction in  (PWGmcIjContraction)}}\label{rule-PWGmcIjContraction}

\begin{ReaderFacingFoma}
define PWGmcIjContraction [
    {*i} {*j} {*ō} -> {*iu} || _ EnglishStarConsonant,
    {*í} {*j} {*ō} -> {*íu} || _ EnglishStarConsonant
];
\end{ReaderFacingFoma}

Only the \texttt{friend} family tests this contraction. If the rare
\emph{*ijō} sequence survives until after
\hyperref[rule-OEDiphthongLeveling]{SC032 OEDiphthongLeveling}, PGmc
\Recon{fríjōndz} `friend' yields \Pred{friund} rather than expected OE
\emph{frēond} `friend'; moving contraction earlier within the tested
range changes no checked output.

That single contrast places \hyperref[rule-PWGmcIjContraction]{SC009
PWGmcIjContraction} before diphthong leveling but gives no lower
boundary. It cannot establish a productive sound law beyond this family,
precisely the reservation made by Ringe and Taylor.

\newpage

\section{West Germanic j-gemination}\label{west-germanic-j-gemination}

\subsection{Historical discussion}\label{historical-discussion-7}

Fulk treats West Germanic consonant gemination before \texttt{*j} after
a short vowel as a regular development and illustrates it with forms
such as OE \emph{settan} `set' and \emph{lecgan} `lay'
(\citeproc{ref-Fulk2018}{Fulk 2018, 127}, §6.15).

The change applies specifically after a short vowel before \emph{*j},
not to geminate consonants generally.

\subsection{\texorpdfstring{SC010. West Germanic j-gemination
(\texttt{PWGmcJGemination})}{SC010. West Germanic j-gemination (PWGmcJGemination)}}\label{rule-PWGmcJGemination}

\begin{ReaderFacingFoma}
define PWGmcJGemination [
    {*p} -> {*p} {*p} || EnglishStarShortVowel _ {*j},
    {*b} -> {*b} {*b} || EnglishStarShortVowel _ {*j},
    {*t} -> {*t} {*t} || EnglishStarShortVowel _ {*j},
    {*d} -> {*d} {*d} || EnglishStarShortVowel _ {*j},
    {*k} -> {*k} {*k} || EnglishStarShortVowel _ {*j},
    {*g} -> {*g} {*g} || EnglishStarShortVowel _ {*j},
    {*f} -> {*f} {*f} || EnglishStarShortVowel _ {*j},
    {*s} -> {*s} {*s} || EnglishStarShortVowel _ {*j},
    {*m} -> {*m} {*m} || EnglishStarShortVowel _ {*j},
    {*n} -> {*n} {*n} || EnglishStarShortVowel _ {*j},
    {*l} -> {*l} {*l} || EnglishStarShortVowel _ {*j},
    {*ŋ} -> {*ŋ} {*ŋ} || EnglishStarShortVowel _ {*j},
    {*x} -> {*x} {*x} || EnglishStarShortVowel _ {*j}
];
\end{ReaderFacingFoma}

OE \emph{nett} `net' fixes the order because the syllabic-\emph{j}
development would remove the glide that conditions gemination. If
\hyperref[rule-PWGmcSyllabicJ]{SC011 PWGmcSyllabicJ} precedes
\hyperref[rule-PWGmcJGemination]{SC010 PWGmcJGemination}, PGmc
\Recon{nátją} `net' yields \Pred{nete} rather than expected OE
\emph{nett} `net'. Earlier movement of gemination changes no checked
output.

The chronology is phonologically transparent: the consonant must
geminate before \emph{*j} ceases to be consonantal. The witness
establishes no earlier boundary.

\newpage

\section{Syllabic j after final-vowel
loss}\label{syllabic-j-after-final-vowel-loss}

\subsection{Historical discussion}\label{historical-discussion-8}

Ringe and Taylor state directly that after final unstressed \texttt{*a}
and \texttt{*ą} were lost, postconsonantal \texttt{*j} became syllabic
\texttt{*i}, with outcomes behind OE \emph{here} `army' and \emph{rice}
`kingdom' (\citeproc{ref-RingeTaylor2014}{Ringe and Taylor 2014, 46}).

The sources establish the development, although the checked lexicon
supplies little independent evidence for its position. Its scope is
postconsonantal j, not high-vowel vocalization generally.

\subsection{\texorpdfstring{SC011. Syllabic \emph{*j} after final-vowel
loss
(\texttt{PWGmcSyllabicJ})}{SC011. Syllabic  after final-vowel loss (PWGmcSyllabicJ)}}\label{rule-PWGmcSyllabicJ}

\begin{ReaderFacingFoma}
define PWGmcSyllabicJ [
    {*j} {*a} -> {*i} || EnglishStarShortVowel EnglishStarConsonant _ .#.,
    {*j} {*ą} -> {*i} || EnglishStarShortVowel EnglishStarConsonant _ .#.
];
\end{ReaderFacingFoma}

The same PGmc \Recon{nátją} `net' witness supplies the only firm
boundary. Placing \hyperref[rule-PWGmcSyllabicJ]{SC011 PWGmcSyllabicJ}
before \hyperref[rule-PWGmcJGemination]{SC010 PWGmcJGemination} yields
\Pred{nete} rather than expected OE \emph{nett} `net'; moving it later
changes no checked output.

Comparative evidence establishes postconsonantal \emph{*j} to syllabic
\emph{*i} after final unstressed \emph{*a} or \emph{*ą} loss, with
\emph{here} `army' and \emph{rice} `kingdom' as outcomes. The lexicon
adds only that vocalization followed gemination, not where it falls
among subsequent changes.

\newpage

\section{\texorpdfstring{\emph{lþ}-voicing}{-voicing}}\label{voicing}

\subsection{Historical discussion}\label{historical-discussion-9}

Ringe and Taylor treat word-internal \emph{*lþ} \textgreater{}
\emph{*ld} as a regular sound change in northern West Germanic and
illustrate it with forms such as \emph{fealdan} `fold', \emph{beald}
`bold', \emph{wuldor} `glory', and \emph{gylden} `golden'
(\citeproc{ref-RingeTaylor2014}{Ringe and Taylor 2014, 170--71}).
Campbell gives a similar West-Germanic-facing formulation with examples
such as \emph{fealdan}, \emph{wuldor}, \emph{beald}, \emph{gold} `gold',
and \emph{feld} `field' (\citeproc{ref-Campbell1959}{Campbell 1959,
169}, §414).

The comparative evidence supports \emph{lþ > ld} most clearly in
northern West Germanic, not as an unqualified pan-PWGmc development.

\subsection{\texorpdfstring{SC012. \emph{lþ}-voicing
(\texttt{PWGmcLThVoicing})}{SC012. -voicing (PWGmcLThVoicing)}}\label{rule-PWGmcLThVoicing}

\begin{ReaderFacingFoma}
define PWGmcLThVoicing [
    {*θ} -> {*d} || {*l} _
];
\end{ReaderFacingFoma}

The \texttt{field}, \texttt{fold}, \texttt{gold}, and \texttt{wold}
families preserve \emph{*lþ} to \emph{*ld}, but none dates the change
against a neighboring rule. Every checked output remains unchanged when
the voicing is moved in either direction.

Comparative reconstruction therefore establishes northern West Germanic
\emph{lþ > ld}, but the witness forms fix no date. Neither a pan-PWGmc
attribution nor an exact local placement follows from the evidence
presented here.

\newpage

\section{Dental hardening}\label{dental-hardening}

\subsection{Historical discussion}\label{historical-discussion-10}

Ringe and Taylor state directly that in PWGmc voiced dental fricative
\texttt{*ð} became stop \texttt{*d} in all positions
(\citeproc{ref-RingeTaylor2014}{Ringe and Taylor 2014, 43}).

The change is systemic across early West Germanic and extends beyond any
one lexical family.

\subsection{\texorpdfstring{SC013. Dental hardening
(\texttt{PWGmcDentalHardening})}{SC013. Dental hardening (PWGmcDentalHardening)}}\label{rule-PWGmcDentalHardening}

\begin{ReaderFacingFoma}
define PWGmcDentalHardening [
    {*ð} -> {*d}
];
\end{ReaderFacingFoma}

Dental hardening has systemic scope: voiced fricative \emph{*ð} became
stop \emph{*d} throughout early West Germanic. Moving
\hyperref[rule-PWGmcDentalHardening]{SC013 PWGmcDentalHardening} earlier
or later changes no checked output.

Comparative evidence establishes the sound law; the present lexicon
leaves its exact position approximate.

\newpage

\section{Early unstressed vowel
changes}\label{early-unstressed-vowel-changes}

\subsection{Historical discussion of the earliest unstressed vowel
changes}\label{historical-discussion-of-the-earliest-unstressed-vowel-changes}

The first change removes the remaining diphthongal quality of unstressed
\emph{*ai}; the second carries early unstressed front-vowel leveling
farther in forms such as \emph{weorold} `world'. Their chronological
evidence differs: monophthongization is historically clear but not
closely dated by the witness forms, whereas \emph{*i}-lowering has a
diagnostic later boundary.

\subsection{\texorpdfstring{Historical discussion of unstressed
\emph{*ai}
monophthongization}{Historical discussion of unstressed  monophthongization}}\label{historical-discussion-of-unstressed-monophthongization}

Ringe and Taylor describe the broad Northwest Germanic reduction of
unstressed \emph{*ai} to a long mid vowel that merges with unstressed
\emph{*e} (\citeproc{ref-RingeTaylor2014}{Ringe and Taylor 2014,
37--41}). The historical change is thus established, although the order
test determines no closer relative position.

\subsection{\texorpdfstring{\CAPRRuleHeading{SC014. Monophthongization of unstressed \emph{*ai}}{NWGmcUnstressedAiMonophthongization}}{}}\label{rule-NWGmcUnstressedAiMonophthongization}

\begin{ReaderFacingFoma}
define NWGmcUnstressedAiMonophthongization [
    {*ăi} -> {*ē}
];
\end{ReaderFacingFoma}

Moving \hyperref[rule-NWGmcUnstressedAiMonophthongization]{SC014
NWGmcUnstressedAiMonophthongization} earlier or later changes no checked
form. The lexicon therefore cannot refine its source-based placement
among the earliest Northwest Germanic simplifications of unstressed
vowels.

Ringe and Taylor's merger of unstressed \emph{*ai} with \emph{*e}
establishes the historical development; the current witnesses do not
distinguish its position relative to neighboring changes.

\subsection{Historical discussion of early unstressed front-vowel
leveling}\label{historical-discussion-of-early-unstressed-front-vowel-leveling}

Campbell treats the merger of unstressed front vowels directly and also
records the variation of \emph{weorold} `world' and \emph{weoruld}
`world' (\citeproc{ref-Campbell1959}{Campbell 1959, 141--42, 154--55}).
These forms supply \hyperref[rule-NWGmcILowering]{SC015 NWGmcILowering}
with a firmer lexical basis than the preceding change.

\subsection{\texorpdfstring{SC015. Leveling of early unstressed front
vowels
(\texttt{NWGmcILowering})}{SC015. Leveling of early unstressed front vowels (NWGmcILowering)}}\label{rule-NWGmcILowering}

\begin{ReaderFacingFoma}
define NWGmcILowering [
    {*i} -> {*e}
        || .#. EnglishStarNonVelarConsonant* _
           EnglishStarCoronal+ EnglishStarNonHighVowel,
    {*í} -> {*é}
        || .#. EnglishStarNonVelarConsonant* _
           EnglishStarCoronal+ EnglishStarNonHighVowel
];
\end{ReaderFacingFoma}

The \emph{weorold} `world' and \emph{weoruld} `world' variants turn the
general source claim into an ordering test. If
\hyperref[rule-NWGmcILowering]{SC015 NWGmcILowering} is delayed until
after \hyperref[rule-OEInterStressRaising]{SC036 OEInterStressRaising},
PGmc \Recon{wír-àldu} `world' yields \Pred{wuruld} rather than expected
OE \emph{weorold} `world'; earlier movement changes no checked output.

The derivation thus fixes front-vowel leveling before interstress
raising while leaving its earlier boundary open.

\hyperref[rule-OEWsPalatalGlide]{SC016 OEWsPalatalGlide} and
\hyperref[rule-NWGmcULowering]{SC017 NWGmcULowering} follow with a more
tightly constrained local chronology.

\newpage

\section{Northwest Germanic
u-lowering}\label{northwest-germanic-u-lowering}

\subsection{Historical discussion}\label{historical-discussion-11}

The derivation of \emph{ġeoc} `yoke' passes through both this change and
the preceding West Saxon palatal glide
(\hyperref[rule-OEWsPalatalGlide]{SC016 OEWsPalatalGlide}). Campbell
treats the West Saxon rising-diphthong spellings before back vowels,
while the same handbook tradition describes the lowering of \emph{u}
before a following non-high vowel separately
(\citeproc{ref-Campbell1959}{Campbell 1959, 17}, §44;
\citeproc{ref-Campbell1959}{Campbell 1959, 42--43}, §115;
\citeproc{ref-Fulk2018}{Fulk 2018, 56}, §4.3). The first change creates
the West Saxon \emph{ġeoc} type; u-lowering then carries the same
material into the subsequent vowel history.

After the glide-conditioned West Saxon spellings are in place, the
broader Northwest Germanic lowering of \emph{u} to \emph{o} before a
following non-high vowel provides the clearest standard sound change in
this small region. Campbell and Fulk both describe that change directly
(\citeproc{ref-Campbell1959}{Campbell 1959, 42--43}, §115;
\citeproc{ref-Fulk2018}{Fulk 2018, 56}, §4.3).

\hyperref[rule-NWGmcULowering]{SC017 NWGmcULowering} thus rests on a
broader source base than the preceding West Saxon rule.

\subsection{\texorpdfstring{\CAPRRuleHeading{SC017. Lowering of \emph{*u} before following non-high vowels}{NWGmcULowering}}{}}\label{rule-NWGmcULowering}

\begin{ReaderFacingFoma}
define NWGmcULowering [
    {*u} -> {*o}
        || .#. EnglishStarConsonant* _
           [EnglishStarConsonantNoJ - EnglishStarNasal]
           EnglishStarConsonantNoJ* EnglishStarNonHighVowel,
    {*ú} -> {*ó}
        || .#. EnglishStarConsonant* _
           [EnglishStarConsonantNoJ - EnglishStarNasal]
           EnglishStarConsonantNoJ* EnglishStarNonHighVowel
];
\end{ReaderFacingFoma}

Lowering of \emph{u} to \emph{o} is fixed on both sides by \emph{ġeoc}
`yoke', \emph{nosu} `nose', \emph{sċofl} `shovel', and \emph{sorg}
`sorrow'.

Before \hyperref[rule-OEWsPalatalGlide]{SC016 OEWsPalatalGlide}, PGmc
\Recon{júką} `yoke' yields \Pred{ġoc} rather than expected OE
\emph{ġeoc} `yoke'. After \hyperref[rule-NWGmcFinalLongORaising]{SC019
NWGmcFinalLongORaising}, PGmc \Recon{núsō} `nose' yields \Pred{nusu}
rather than expected \emph{nosu} `nose', PGmc \Recon{skúflō} `shovel'
yields \Pred{sċufl} rather than expected \emph{sċofl} `shovel', and PGmc
\Recon{súrgō} `sorrow' yields \Pred{surg} rather than expected
\emph{sorg} `sorrow'. The two witness sets place
\hyperref[rule-NWGmcULowering]{SC017 NWGmcULowering} after glide
formation and before final long-\emph{o} raising.

\newpage

\section{\texorpdfstring{Stressed monosyllable
\emph{*ō}-raising}{Stressed monosyllable -raising}}\label{stressed-monosyllable--raising}

\subsection{Historical discussion}\label{historical-discussion-12}

Campbell treats the development of final accented \emph{ō} to \emph{ū}
in stressed monosyllables directly, with the familiar outcomes behind
\emph{cū} `cow', \emph{hū} `how', \emph{tū} `two', and \emph{bū} `both'
(\citeproc{ref-Campbell1959}{Campbell 1959, 47}, §122).

The change is historically secure, but the tested forms determine no
close relative position for it. Its input is final \emph{*ō} in a
stressed monosyllable.

\subsection{\texorpdfstring{\CAPRRuleHeading{SC018. Raising of final stressed monosyllabic \emph{*ō}}{NWGmcStressedMonosyllableORaising}}{}}\label{rule-NWGmcStressedMonosyllableORaising}

\begin{ReaderFacingFoma}
define NWGmcStressedMonosyllableORaising [
    {*ō} -> {*ū} || .#. [EnglishStarConsonant | EnglishPalatalConsonant]* _ .#.
];
\end{ReaderFacingFoma}

Campbell's \emph{cū} `cow', \emph{hū} `how', and \emph{tū} `two'
establish final stressed monosyllabic \emph{*ō} \textgreater{}
\emph{*ū}.

Reversing \hyperref[rule-NWGmcStressedMonosyllableORaising]{SC018
NWGmcStressedMonosyllableORaising} with neighboring changes leaves every
checked output unchanged. The sound change is secure, but its exact
position in the early history of long vowels rests on the handbooks.

\newpage

\section{\texorpdfstring{Raising of final unstressed long
\emph{*ō}}{Raising of final unstressed long }}\label{raising-of-final-unstressed-long}

\subsection{Historical discussion}\label{historical-discussion-13}

Ringe and Taylor describe the change of unstressed final non-nasalized
long \emph{*ō} to short \emph{*u} as a Northwest Germanic development
(\citeproc{ref-RingeTaylor2014}{Ringe and Taylor 2014, 30}). It applies
in the same final-syllable environment as the subsequent loss of
word-final \emph{*z}. The derivation of \emph{ræste} `rest' fixes their
local order: \hyperref[rule-NWGmcFinalLongORaising]{SC019
NWGmcFinalLongORaising} must still see final \emph{*ō}, and word-final
\emph{*z}-deletion removes the following \emph{*z} only afterward.

The change supplies the final vowel of forms such as \emph{nosu} `nose',
\emph{sċofl} `shovel', and \emph{sorg} `sorrow'.

\subsection{\texorpdfstring{\CAPRRuleHeading{SC019. Raising of final unstressed long \emph{*ō}}{NWGmcFinalLongORaising}}{}}\label{rule-NWGmcFinalLongORaising}

\begin{ReaderFacingFoma}
define NWGmcFinalLongORaising [
    {*ō} -> {*u}
        || EnglishStarVocalic
           [EnglishStarConsonant | EnglishPalatalConsonant]+ _ .#.
];
\end{ReaderFacingFoma}

Two groups of witnesses confine final unstressed long \emph{*ō}
\textgreater{} \emph{*u}. The forms \emph{nosu} `nose', \emph{sċofl}
`shovel', and \emph{sorg} `sorrow' fix its lower boundary.

Before \hyperref[rule-NWGmcULowering]{SC017 NWGmcULowering}, PGmc
\Recon{núsō} `nose' yields \Pred{nusu} rather than expected OE
\emph{nosu} `nose', PGmc \Recon{skúflō} `shovel' yields \Pred{sċufl}
rather than expected \emph{sċofl} `shovel', and PGmc \Recon{súrgō}
`sorrow' yields \Pred{surg} rather than expected \emph{sorg} `sorrow'.
After word-final \emph{*z}-deletion
(\hyperref[rule-PGmcFinalZDeletion]{SC020 PGmcFinalZDeletion}), PGmc
\Recon{rástōz} `rest' yields \Pred{rast} rather than expected
\emph{ræste} `rest'. These failures place
\hyperref[rule-NWGmcFinalLongORaising]{SC019 NWGmcFinalLongORaising}
after u-lowering and before final \emph{z}-loss.

\newpage

\section{\texorpdfstring{Unstressed
\emph{*o}-raising}{Unstressed -raising}}\label{unstressed--raising}

\subsection{Historical discussion}\label{historical-discussion-14}

The older history of \emph{heofon} `heaven' requires an unstressed-vowel
adjustment before the later reshaping of medial vowels in Old English.
Campbell derives the \emph{-o-} from an earlier unstressed environment,
and Ringe and Taylor place the same family within the wider West
Germanic record (\citeproc{ref-Campbell1959}{Campbell 1959, 155--56},
§373; \citeproc{ref-RingeTaylor2014}{Ringe and Taylor 2014, 287}).

The change is historically recognizable, but the checked forms provide
only a later boundary.

\subsection{\texorpdfstring{\CAPRRuleHeading{SC021. Raising of unstressed \emph{*o} before later \emph{*u}}{NWGmcUnstressedORaising}}{}}\label{rule-NWGmcUnstressedORaising}

\begin{ReaderFacingFoma}
define NWGmcUnstressedORaising [
    {*o} -> {*u} || EnglishStarVocalic EnglishStarConsonant+ _ EnglishStarConsonant* {*ų}
];
\end{ReaderFacingFoma}

After \hyperref[rule-OEMedUnstressedULowering]{SC040
OEMedUnstressedULowering}, PGmc \Recon{xémonų} `heaven' yields
\Pred{heofun} rather than expected OE \emph{heofon} `heaven'; earlier
placement changes no checked output. The witness therefore places
\hyperref[rule-NWGmcUnstressedORaising]{SC021 NWGmcUnstressedORaising}
before medial unstressed-\emph{u} lowering.

Nothing in the present lexicon supplies the corresponding earlier
boundary.

\newpage

\section{\texorpdfstring{\emph{mn}-dissimilation}{-dissimilation}}\label{dissimilation}

\subsection{Historical discussion}\label{historical-discussion-15}

The handbooks describe the history of \emph{mn} sequences as a limited
descriptive pattern. Campbell discusses both loss of unstressed material
and later assimilation in forms of this type, including the special
status of \emph{mōnaþ} `month'-type evidence
(\citeproc{ref-Campbell1959}{Campbell 1959, 189, 195}, §§470, 484).

The pattern is historically established, but the checked forms do not
constrain its position.

\subsection{\texorpdfstring{SC022. Dissimilation of \emph{mn} sequences
(\texttt{NWGmcMnDissimilation})}{SC022. Dissimilation of  sequences (NWGmcMnDissimilation)}}\label{rule-NWGmcMnDissimilation}

\begin{ReaderFacingFoma}
define NWGmcMnDissimilation [
    {*m} -> {*β}
        || EnglishStarVocalic _
           EnglishStarVocalic EnglishStarConsonant* EnglishStarNasal
];
\end{ReaderFacingFoma}

Campbell's \emph{heofon} `heaven' and \emph{mōnaþ} `month' material
supports early \emph{m} \textgreater{} \emph{β} before a later nasal,
but supplies no ordering witness.

Moving \hyperref[rule-NWGmcMnDissimilation]{SC022 NWGmcMnDissimilation}
earlier or later leaves every checked output unchanged. Its place among
the early consonantal changes rests on the handbook account of
\emph{mn}-dissimilation.

\newpage

\section{\texorpdfstring{N-stem
\emph{n}-loss}{N-stem -loss}}\label{n-stem--loss}

\subsection{Historical discussion}\label{historical-discussion-16}

The broader history is the reduction and leveling of older n-stem
endings in West Germanic. Ringe and Taylor describe the resulting
syncretism in the n-stems, which is the wider morphological setting for
the narrower step isolated here (\citeproc{ref-RingeTaylor2014}{Ringe
and Taylor 2014, 72}).

The path to \emph{dōn} `do' provides the clearest witness, but the
change remains narrow in scope.

\subsection{\texorpdfstring{SC023. Loss of n-stem \emph{*n} in final
position
(\texttt{NWGmcNStemNLoss})}{SC023. Loss of n-stem  in final position (NWGmcNStemNLoss)}}\label{rule-NWGmcNStemNLoss}

\begin{ReaderFacingFoma}
define NWGmcNStemNLoss [
    {*ō} {*n} -> {*ǭ} || _ .#.
];
\end{ReaderFacingFoma}

After \hyperref[rule-OEHeavySyllableNasalApocope]{SC047
OEHeavySyllableNasalApocope}, PGmc \Recon{dōną} `do' fails entirely
(\emph{+?}) instead of yielding expected OE \emph{dōn} `do'; earlier
placement changes no checked output. Thus
\hyperref[rule-NWGmcNStemNLoss]{SC023 NWGmcNStemNLoss} must feed the
later apocope.

This failed derivation supplies a terminus ante quem, while the lower
boundary remains unattested.

\newpage

\section{\texorpdfstring{Long
\emph{ē}-lowering}{Long -lowering}}\label{long--lowering}

\subsection{Historical discussion}\label{historical-discussion-17}

The later West Saxon forms \emph{sċēap} `sheep' and \emph{ġēar} `year'
imply an earlier lowering of long \emph{ē} before the palatal
diphthongal outcomes described more fully later in the sequence.
Campbell and Ringe and Taylor discuss those later West Saxon outputs
directly (\citeproc{ref-Campbell1959}{Campbell 1959, 69--70}, §185;
\citeproc{ref-RingeTaylor2014}{Ringe and Taylor 2014, 215--16}, §6.5.1).

The change is historically recognizable, but the checked forms provide
only a later boundary.

\subsection{\texorpdfstring{\CAPRRuleHeading{SC024. Lowering of long \emph{ē} before non-nasal consonants}{NWGmcLongELowering}}{}}\label{rule-NWGmcLongELowering}

\begin{ReaderFacingFoma}
define NWGmcLongELowering [
    {*ē} -> {*ǣ} || _ [EnglishStarConsonant - EnglishStarNasal],
    {*ḗ} -> {*ǣ} || _ [EnglishStarConsonant - EnglishStarNasal]
];
\end{ReaderFacingFoma}

After \hyperref[rule-OEWsPalatalDiphthongization]{SC056
OEWsPalatalDiphthongization}, long \emph{ē} \textgreater{} \emph{ǣ} can
no longer produce the expected West Saxon forms: PGmc \Recon{skḗpą}
`sheep' yields \Pred{sċīep} rather than OE \emph{sċēap} `sheep', and
PGmc \Recon{jḗrą} `year' yields \Pred{ġīer} rather than \emph{ġēar}
`year'. Earlier placement changes no checked output, so
\hyperref[rule-NWGmcLongELowering]{SC024 NWGmcLongELowering} has a
secure upper boundary.

Its lower boundary remains a matter of handbook chronology.

\newpage

\section{\texorpdfstring{Long \emph{ē}
nasal-rounding}{Long  nasal-rounding}}\label{long-nasal-rounding}

\subsection{Historical discussion}\label{historical-discussion-18}

Before nasals, older long \emph{ē} can round toward the
\emph{ō}-vocalism seen later in \emph{mōnaþ} `month' and \emph{mōna}
`moon' / \emph{mōn} `moon'-type material. Campbell treats this split
directly in his discussion of Germanic long \emph{ē} before nasal
consonants (\citeproc{ref-Campbell1959}{Campbell 1959, 53}, §129).

The change is historically recognizable, but the tested forms supply no
close relative chronology.

\subsection{\texorpdfstring{SC025. Rounding of long \emph{ē} before
nasals
(\texttt{NWGmcLongENasalRounding})}{SC025. Rounding of long  before nasals (NWGmcLongENasalRounding)}}\label{rule-NWGmcLongENasalRounding}

\begin{ReaderFacingFoma}
define NWGmcLongENasalRounding [
    {*ē} -> {*ō} || _ EnglishStarNasal,
    {*ḗ} -> {*ō} || _ EnglishStarNasal
];
\end{ReaderFacingFoma}

Reversing \hyperref[rule-NWGmcLongENasalRounding]{SC025
NWGmcLongENasalRounding} with neighboring changes leaves every checked
output unchanged. Its position beside the other \emph{ē}-developments
therefore follows the handbooks.

\newpage

\section{Nasal spirant changes}\label{nasal-spirant-changes}

\subsection{Historical discussion of nasal loss before spirants and
compensatory
lengthening}\label{historical-discussion-of-nasal-loss-before-spirants-and-compensatory-lengthening}

The two rules state successive phases of a single development. Campbell
describes nasal loss before voiceless spirants with compensatory
lengthening and nasalization of the preceding vowel. Ringe and Taylor
assign the same outcomes to inherited northern West Germanic, before
late Old English (\citeproc{ref-Campbell1959}{Campbell 1959, 47}, §121;
\citeproc{ref-RingeTaylor2014}{Ringe and Taylor 2014, 140--41}).

\hyperref[rule-NWGmcNasalSpirantLengthening]{SC026
NWGmcNasalSpirantLengthening} adjusts the vowel while the
nasal-plus-spirant sequence remains present;
\hyperref[rule-NWGmcNasalSpirantLoss]{SC027 NWGmcNasalSpirantLoss} then
removes the nasal. The first rule must therefore precede the second.

\subsection{\texorpdfstring{\CAPRRuleHeading{SC026. Lengthening before nasal plus spirant}{NWGmcNasalSpirantLengthening}}{}}\label{rule-NWGmcNasalSpirantLengthening}

\begin{ReaderFacingFoma}
define NWGmcNasalSpirantLengthening [
    {*a} -> {*ō} || _ EnglishStarNasal EnglishStarVoicelessFricative,
    {*e} -> {*ē} || _ EnglishStarNasal EnglishStarVoicelessFricative,
    {*i} -> {*ī} || _ EnglishStarNasal EnglishStarVoicelessFricative,
    {*o} -> {*ō} || _ EnglishStarNasal EnglishStarVoicelessFricative,
    {*u} -> {*ū} || _ EnglishStarNasal EnglishStarVoicelessFricative,
    {*æ} -> {*ē} || _ EnglishStarNasal EnglishStarVoicelessFricative,
    {*á} -> {*ō} || _ EnglishStarNasal EnglishStarVoicelessFricative,
    {*é} -> {*ḗ} || _ EnglishStarNasal EnglishStarVoicelessFricative,
    {*í} -> {*ī} || _ EnglishStarNasal EnglishStarVoicelessFricative,
    {*ó} -> {*ō} || _ EnglishStarNasal EnglishStarVoicelessFricative,
    {*ú} -> {*ū} || _ EnglishStarNasal EnglishStarVoicelessFricative
];
\end{ReaderFacingFoma}

All three witnesses require the vowel adjustment while the nasal is
still present. If \hyperref[rule-NWGmcNasalSpirantLengthening]{SC026
NWGmcNasalSpirantLengthening} follows
\hyperref[rule-NWGmcNasalSpirantLoss]{SC027 NWGmcNasalSpirantLoss}, PGmc
\Recon{fúnxstiz} `fist' yields \Pred{fyst} rather than expected OE
\emph{fȳst} `fist', PGmc \Recon{gánsz} `goose' yields \Pred{ġeas} rather
than expected \emph{gōs} `goose', and PGmc \Recon{júgunθ} `youth' yields
\Pred{ġeogoþ} rather than expected \emph{ġeoguþ} `youth'. Earlier
placement changes no checked output. The evidence requires lengthening
to precede nasal loss without supplying a lower boundary, in agreement
with the handbook treatment of the two as successive phases.

\subsection{\texorpdfstring{SC027. Loss of the nasal before spirants
(\texttt{NWGmcNasalSpirantLoss})}{SC027. Loss of the nasal before spirants (NWGmcNasalSpirantLoss)}}\label{rule-NWGmcNasalSpirantLoss}

\begin{ReaderFacingFoma}
define NWGmcNasalSpirantLoss [
    EnglishStarNasal -> 0 || _ EnglishStarVoicelessFricative
];
\end{ReaderFacingFoma}

The converse test fixes the same boundary: placing
\hyperref[rule-NWGmcNasalSpirantLoss]{SC027 NWGmcNasalSpirantLoss}
before \hyperref[rule-NWGmcNasalSpirantLengthening]{SC026
NWGmcNasalSpirantLengthening} produces the same errors in \emph{fȳst}
`fist', \emph{gōs} `goose', and \emph{ġeoguþ} `youth'. Later placement
changes no checked output. These forms prove that the vowel was adjusted
before the nasal disappeared; they provide no upper boundary for the
loss.

\newpage

\section{\texorpdfstring{Preconsonantal
\emph{*x}-loss}{Preconsonantal -loss}}\label{preconsonantal--loss}

\subsection{Historical discussion}\label{historical-discussion-19}

Campbell explicitly treats loss of \emph{x} and gives forms such as
\emph{fléam} `flight' and \emph{hēla} `heel' as examples of the same
broad development (\citeproc{ref-Campbell1959}{Campbell 1959, 186},
§461).

The historical evidence is firmer than the chronology: the checked forms
do not constrain the rule's position.

\subsection{\texorpdfstring{SC028. Loss of preconsonantal \emph{*x}
(\texttt{NWGmcPreconsonantalXLoss})}{SC028. Loss of preconsonantal  (NWGmcPreconsonantalXLoss)}}\label{rule-NWGmcPreconsonantalXLoss}

\begin{ReaderFacingFoma}
define NWGmcPreconsonantalXLoss [
    {*x} -> 0 || _ {*s} EnglishStarConsonant
];
\end{ReaderFacingFoma}

No witness word dates preconsonantal \emph{*x}-loss before \emph{*s}
plus another consonant: moving
\hyperref[rule-NWGmcPreconsonantalXLoss]{SC028 NWGmcPreconsonantalXLoss}
in either direction leaves every checked output unchanged. Its position
within this stretch therefore rests on the handbook chronology for
\emph{x}-loss.

\newpage

\section{\texorpdfstring{Final bare-\emph{a}
loss}{Final bare- loss}}\label{final-bare--loss}

\subsection{Historical discussion}\label{historical-discussion-20}

I isolate the loss of final short low vowels within the broader erosion
of final syllables described by the handbooks
(\citeproc{ref-Campbell1959}{Campbell 1959, 143}, §341;
\citeproc{ref-RingeTaylor2014}{Ringe and Taylor 2014, 60--61}).

Final bare-a loss follows the medial unstressed vowel changes and
precedes restoration, which depends on the environment left by the loss.

\subsection{\texorpdfstring{SC041. Loss of final bare \emph{*a}
(\texttt{PWGmcFinalBareALoss})}{SC041. Loss of final bare  (PWGmcFinalBareALoss)}}\label{rule-PWGmcFinalBareALoss}

\begin{ReaderFacingFoma}
define PWGmcFinalBareALoss [
    {*a} -> 0 || _ .#.
];
\end{ReaderFacingFoma}

The two sides of final bare-\emph{a} loss rest on different evidence.
Applied before final \emph{z}-deletion, the change gives the wrong
outputs: PGmc \Recon{bárdaz} `beard' yields \Pred{bearda} rather than
expected OE \emph{beard} `beard', and PGmc \Recon{kámbaz} `comb' yields
\Pred{camba} rather than expected \emph{camb} `comb'. Applied after
restoration, PGmc \Recon{kráftaz} `craft' yields \Pred{craft} rather
than expected OE \emph{cræft} `craft', and PGmc \Recon{dágaz} `day'
yields \Pred{dag} rather than expected \emph{dæġ} `day'. The distant
lower limit follows final \emph{z}-loss; the local feeding relation
precedes restoration, which requires the environment created by the
vowel loss.

\newpage

\section{\texorpdfstring{Surviving bimoric \emph{*ō}
unrounding}{Surviving bimoric  unrounding}}\label{surviving-bimoric-unrounding}

\subsection{Historical discussion}\label{historical-discussion-21}

The handbooks do not isolate a large independent sound change under this
label. The surviving bimoric \emph{*ō} in the pathway to \emph{ræste}
`rest' nevertheless undergoes unrounding before
\hyperref[rule-AngloFrisianBrightening]{SC043 AngloFrisianBrightening}.
Campbell, Hogg, and Ringe and Taylor describe the surrounding fronting
and restoration history without naming this feeder separately
(\citeproc{ref-Campbell1959}{Campbell 1959, 52, 60}, §§131, 157--158;
\citeproc{ref-Hogg1992}{Hogg 1992, 101, 119};
\citeproc{ref-RingeTaylor2014}{Ringe and Taylor 2014, 157--58,
189--90}).

The sole witness establishes a local relation to brightening but
supports no broader generalization.

\subsection{\texorpdfstring{\CAPRRuleHeading{SC042. Unrounding of the surviving bimoric \emph{*ō}}{PWGmcSurvivingBimoricOUnrounding}}{}}\label{rule-PWGmcSurvivingBimoricOUnrounding}

\begin{ReaderFacingFoma}
define PWGmcSurvivingBimoricOUnrounding [
    {*ō} -> {*ā} || EnglishStarVocalic [EnglishStarConsonant | EnglishPalatalConsonant]+ _ .#.
];
\end{ReaderFacingFoma}

The single \emph{ræste} `rest' derivation carries the chronology of
bimoric \emph{*ō} \textgreater{} \emph{*ā}. Before
\hyperref[rule-PGmcFinalZDeletion]{SC020 PGmcFinalZDeletion} or after
\hyperref[rule-AngloFrisianBrightening]{SC043 AngloFrisianBrightening},
PGmc \Recon{rástōz} `rest' yields \Pred{rasta} rather than expected OE
\emph{ræste}. Unrounding must therefore follow final \emph{z}-loss and
precede brightening, although only the relation to brightening is local.

\newpage

\section{Sievers-law syncope}\label{sievers-law-syncope}

\subsection{Historical discussion}\label{historical-discussion-22}

Sievers' Law concerns a prosodic and morphological adjustment in heavy
stems. It is a distributional rule distinct from b-allophony
(\hyperref[rule-PGmcBAllophony]{SC049 PGmcBAllophony}). Adamczyk treats
the Old English reflexes of the law as historical evidence from weak
verbs and related formations (\citeproc{ref-Adamczyk2001}{Adamczyk 2001,
61--72}). Fulk gives the compact comparative summary through familiar
forms such as \emph{biddan} `ask', \emph{sellan} `give', and
\emph{nerian} `save' (\citeproc{ref-Fulk2018}{Fulk 2018, 127}, §6.15).

Sievers-law syncope is narrow in scope, but its relation to the
following palatalization is lexically secure. Its earlier limit is less
sharply defined than that of the preceding allophony rule.

\subsection{\texorpdfstring{SC050. Sievers-law syncope
(\texttt{SieversLawSyncope})}{SC050. Sievers-law syncope (SieversLawSyncope)}}\label{rule-SieversLawSyncope}

\begin{ReaderFacingFoma}
define SieversLawSyncope [
    {*i} -> 0 || [EnglishStarConsonant | EnglishPalatalConsonant] _ {*j}
];
\end{ReaderFacingFoma}

The Sievers-law reduction \emph{*-CijV-*} \textgreater{} \emph{*-CjV-*},
including loss of \emph{*i} before \emph{*j}, must precede
palatalization. If \hyperref[rule-SieversLawSyncope]{SC050
SieversLawSyncope} follows \hyperref[rule-OEVelarPalatalization]{SC052
OEVelarPalatalization}, PGmc \Recon{strákkijaną} `stretch' yields
\Pred{strecċan} rather than expected OE \emph{streċċan} `stretch';
earlier placement creates no comparably precise error. The single
cluster witness therefore places syncope before velar palatalization.

\newpage

\chapter{From Proto-West Germanic to
Anglo-Frisian}\label{from-proto-west-germanic-to-anglo-frisian}

\section{Historical interval}\label{historical-interval-2}

This chapter covers the sound changes that occurred after the Proto-West
Germanic period and before, or during the emergence of, the specifically
English line. The starting reconstruction is Proto-West Germanic; the
end point is the Proto-Anglo-Frisian stage --- or more precisely, the
cluster of innovations that define the English and Frisian branch within
West Germanic.

\section{A necessary terminological
caution}\label{a-necessary-terminological-caution}

The title of this chapter uses ``Anglo-Frisian'' as an organizing
historical concept. That choice requires an explicit qualification.

The scholarly literature uses several overlapping terms for this
developmental period:

\begin{itemize}
\tightlist
\item
  \textbf{North Sea Germanic} and \textbf{Ingvaeonic}: labels used by
  some scholars for a proposed subgroup comprising Old English, Old
  Frisian, and Old Saxon (or more narrowly, Old English and Old Frisian
  only). The innovations associated with this label --- especially the
  nasal spirant changes and certain vowel developments --- are sometimes
  described as diffusion rather than shared inheritance
  (\citeproc{ref-Campbell1959}{Campbell 1959, sects 1--3}).
\item
  \textbf{Anglo-Frisian}: a label used specifically for the Old English
  / Old Frisian branch, or for innovations shared between the two
  languages. Its use presupposes a tighter relationship between English
  and Frisian than between either and Old Saxon.
\item
  \textbf{Proto-Anglo-Frisian} (PAF): the strongest interpretation,
  positing a discrete reconstructed common ancestor for Old English and
  Old Frisian specifically. This is the position of Ringe and Taylor,
  who reconstruct a PAF stage between Proto-West Germanic and
  Proto-Old-English (\citeproc{ref-RingeTaylor2014}{Ringe and Taylor
  2014, 54--68}).
\end{itemize}

\textbf{CAPR does not commit to a universally accepted discrete PAF
node.} The chapter title uses ``Anglo-Frisian'' because the key changes
of this period --- especially West Germanic rhotacism (SC003),
word-final \texttt{*z} deletion (SC020), and Anglo-Frisian brightening
(SC043) --- are most prominently associated with that label in the
handbook literature. But the analysis does not require that every change
passed through a single genealogical PAF stage. Some changes may be West
Germanic broadly; others may reflect areal diffusion. The existing CAPR
dossiers record the source-by-source picture where these distinctions
matter.

\section{Major changes and their historical
basis}\label{major-changes-and-their-historical-basis}

\subsection{West Germanic rhotacism
(SC003)}\label{west-germanic-rhotacism-sc003}

The medial change of \texttt{*z} to \texttt{*r} in environments such as
\Recon{déuzaz}\index[iv]{03pgmc@\textbf{Proto-Germanic}!deuzaz@déuzaz}
`deer',
\Recon{xúrdaz}\index[iv]{03pgmc@\textbf{Proto-Germanic}!xurdaz@xúrdaz}
`hoard', and
\Recon{líznōjaną}\index[iv]{03pgmc@\textbf{Proto-Germanic}!liznojana@líznōjaną}
`learn' is historically a post-Proto-West-Germanic development. Ringe
and Taylor argue that rhotacism was not inherited from Proto-Northwest
Germanic and was not uniform within West Germanic
(\citeproc{ref-RingeTaylor2014}{Ringe and Taylor 2014, 52, 98, 102}).
Crist separates this change explicitly from the deletion of word-final
\texttt{*z} and argues that rhotacism must follow the deletion rules
(\citeproc{ref-Crist2001}{Crist 2001, 104--6};
\citeproc{ref-Crist2002}{2002, 1, 4}). Hogg gives the standard Old
English--facing summary: \texttt{*z} yielded \texttt{*r} in intervocalic
position but was generally lost in final position
(\citeproc{ref-Hogg1992}{Hogg 1992, 37}).

The CAPR rule is named \texttt{PGmcRhotacism}, which is historically
misleading: the change is not Proto-Germanic. The internal rule name is
a stable identifier that will not be changed during this phase; the
reader-facing chapter label corrects the historical description.

\subsection{\texorpdfstring{Word-final \texttt{*z} deletion
(SC020)}{Word-final *z deletion (SC020)}}\label{word-final-z-deletion-sc020}

The deletion of word-final \texttt{*z} in forms such as
\Recon{rástōz}\index[iv]{03pgmc@\textbf{Proto-Germanic}!rastoz@rástōz}
`rest (nom.sg.)' is placed here on the basis of Crist's analysis, which
distinguishes a pan-West-Germanic loss of \texttt{*z} after unstressed
vowels from the earlier NWGmc changes and from the later narrower
Ingvaeonic deletion rules (\citeproc{ref-Crist2002}{Crist 2002, 1, 4}).
The standard handbooks confirm a West Germanic deletion: Campbell notes
that \texttt{*z} is ``later lost or changed to \texttt{r}''
(\citeproc{ref-Campbell1959}{Campbell 1959}); Hogg gives a clean
statement that Germanic \texttt{*z} is generally lost in final position
(\citeproc{ref-Hogg1992}{Hogg 1992, 37}).

The CAPR rule is named \texttt{PGmcFinalZDeletion}, which is
historically misleading. The existing reader-facing prose already notes
this: the current assembled section correctly describes SC003 as
presupposing an earlier loss of final \texttt{*z}, creating a
presentational tension with SC020's place at cascade position 20. This
tension is a targeted audit item for the next phase of the chronology
work.

\subsection{Anglo-Frisian brightening
(SC043)}\label{anglo-frisian-brightening-sc043}

The fronting of low \texttt{*a} to \texttt{*æ} outside nasal
environments is the defining Anglo-Frisian change and the central event
of this chapter. Campbell gives the classical statement: ``By a very
early change Prim. Gmc. \texttt{a\ \textgreater{}\ æ} in OE and OFris.
when not followed by a nasal consonant''
(\citeproc{ref-Campbell1959}{Campbell 1959, sects 163--165}). Hogg gives
the most familiar modern label pair: ``This vowel normally fronted to
/ae/ by the sound change of Anglo-Frisian Brightening (or First
Fronting)'' (\citeproc{ref-Hogg1992}{Hogg 1992, sect. 5.8}).

The change is notable for what follows it: OE Breaking presupposes the
fronted input; OE a-Restoration partially undoes it in back-vowel
environments. The three-change sequence (brightening, breaking,
restoration) is one of the clearest relative-chronology chains in the
Old English historical grammar.

Campbell notes that English and Frisian may not simply reflect one
undifferentiated shared prehistoric event, and Ringe and Taylor leave
open whether the wider spread of fronted outcomes happened mainly on the
continent or in Britain (\citeproc{ref-RingeTaylor2014}{Ringe and Taylor
2014, 60--62}). CAPR's implementation treats the change as a single
rule; the book prose acknowledges the uncertainty about its exact
geographical scope.

The current CAPR inventory has this change labeled ``Old English'' in
the pipeline taxonomy (no separate Anglo-Frisian bucket previously
existed). The historical staging map now places it in Chapter 3,
correcting that provisional label.

\section{Cascade vs.~historical order in this
chapter}\label{cascade-vs.-historical-order-in-this-chapter}

The three changes in this chapter currently occur at cascade positions
3, 19--20, and 43 respectively. These cascade positions reflect
computational dependencies, not historical sequence. The reader-facing
book order in this chapter places rhotacism first (as the post-PWGmc
WGmc change), then word-final \texttt{*z} deletion (SC020, closely
related to rhotacism and discussed in the same reader-facing section as
SC019), then Anglo-Frisian brightening as the culminating change of the
pre-OE period.

The divergence between cascade order and book order in this chapter is
one of the clearest illustrations of the principle that FST dependency
does not automatically equal historical sequence.

\section{West Germanic rhotacism}\label{west-germanic-rhotacism}

\subsection{Historical discussion}\label{historical-discussion-23}

Hogg states that Germanic \emph{*z} yielded \emph{*r} in intervocalic
position in Old English, while final \emph{*z} was generally lost
(\citeproc{ref-Hogg1992}{Hogg 1992, 37}). Ringe and Taylor argue that
this merger of \emph{*z} with \emph{*r} was independent in Norse and
West Germanic and belongs after the Proto-West-Germanic stage
(\citeproc{ref-RingeTaylor2014}{Ringe and Taylor 2014, 52, 98, 102}).
Crist likewise places rhotacism after earlier West Germanic
\emph{*z}-deletion rules and rejects treating it as an inherited
Proto-Northwest-Germanic innovation (\citeproc{ref-Crist2001}{Crist
2001, 104--6}; \citeproc{ref-Crist2002}{2002, 1, 4}).

The label \hyperref[rule-PGmcRhotacism]{SC003 PGmcRhotacism} is
historically misleading: the change is a later West Germanic rhotacism,
not a Proto-Germanic one. It is also distinct from
\hyperref[rule-PGmcFinalZDeletion]{SC020 PGmcFinalZDeletion}, which
removes final \emph{*z} before the surviving medial consonant becomes
\emph{*r}.

\subsection{\texorpdfstring{SC003. West Germanic rhotacism
(\texttt{PGmcRhotacism})}{SC003. West Germanic rhotacism (PGmcRhotacism)}}\label{rule-PGmcRhotacism}

\begin{ReaderFacingFoma}
define PGmcRhotacism [
    {*z} -> {*r} || EnglishStarVocalic _ ?
];
\end{ReaderFacingFoma}

Breaking supplies the decisive upper boundary. If rhotacism is delayed
until after \hyperref[rule-OEBreaking]{SC044 OEBreaking}, PGmc
\Recon{líznōjaną} `learn' yields \Pred{lirnian} rather than expected OE
\emph{liornian} `learn', PGmc \Recon{líznōθi} `learns' yields
\Pred{lirnaþ} rather than expected \emph{liornaþ} `learns', PGmc
\Recon{líznô} `learn' yields \Pred{lirna} rather than expected
\emph{liorna} `learn', and PGmc \Recon{mízdai} `meed' yields
\Pred{merde} rather than expected OE \emph{meorde} `meed'. Moving
rhotacism earlier within the tested range changes none of the checked
forms.

The lexical evidence thus supplies a terminus ante quem but no terminus
post quem. Its placement after the earlier loss of final \emph{*z} rests
on the historical analyses cited above.

\newpage

\section{\texorpdfstring{Deletion of word-final
\emph{*z}}{Deletion of word-final }}\label{deletion-of-word-final}

\subsection{Historical discussion}\label{historical-discussion-24}

The loss of word-final \emph{*z} is a West Germanic development.
Standard handbook tradition and Crist's West Germanic discussion
establish the development within broader accounts of inflectional
morphology (\citeproc{ref-Hogg1992}{Hogg 1992, 37};
\citeproc{ref-Crist2002}{Crist 2002, 1}). The derivation of \emph{ræste}
`rest' demonstrates the local order: final \emph{*ō}-raising
(\hyperref[rule-NWGmcFinalLongORaising]{SC019 NWGmcFinalLongORaising},
Chapter 2) must precede final \emph{*z}-loss.

The internal CAPR rule is labelled
\hyperref[rule-PGmcFinalZDeletion]{SC020 PGmcFinalZDeletion}, but this
label reflects a legacy identifier and not the historical stage of the
change. The existing CAPR prose (and literature dossier) identifies the
development as a West Germanic pan-WGmc loss, not a Proto-Germanic one.
The exact scope (whether all West Germanic or specifically Ingvaeonic)
and the precise relationship between this rule and West Germanic
rhotacism (\hyperref[rule-PGmcRhotacism]{SC003 PGmcRhotacism}) remain
flagged for targeted chronology audit.

Final z-loss follows long-o raising and precedes the later changes in
weak syllables.

\subsection{\texorpdfstring{SC020. Deletion of word-final \emph{*z}
(\texttt{PGmcFinalZDeletion})}{SC020. Deletion of word-final  (PGmcFinalZDeletion)}}\label{rule-PGmcFinalZDeletion}

\begin{ReaderFacingFoma}
define PGmcFinalZDeletion [{*z} -> 0 || _ .#.];
\end{ReaderFacingFoma}

The chronology of word-final \emph{*z}-loss is unusually well delimited:
\emph{ræste} `rest' supplies its early boundary, while later weak
syllables supply its late boundary.

Before \hyperref[rule-NWGmcFinalLongORaising]{SC019
NWGmcFinalLongORaising}, PGmc \Recon{rástōz} `rest' yields \Pred{rast}
rather than expected OE \emph{ræste} `rest'. After
\hyperref[rule-OEMedUnstressedULowering]{SC040
OEMedUnstressedULowering}, PGmc \Recon{bébruz} `beaver' yields
\Pred{befro} rather than expected \emph{befer} `beaver', PGmc
\Recon{kwéðuz} `cud' yields \Pred{cwedo} rather than expected
\emph{cwedu} `cud', and PGmc \Recon{félθuz} `field' yields \Pred{feldo}
rather than expected \emph{feld} `field', alongside eight other newly
failing rows. Final \emph{z}-loss therefore follows
\hyperref[rule-NWGmcFinalLongORaising]{SC019 NWGmcFinalLongORaising} and
precedes \hyperref[rule-OEMedUnstressedULowering]{SC040
OEMedUnstressedULowering}.

The \Recon{rástōz} `rest' derivation fixes the local relation to
\hyperref[rule-NWGmcFinalLongORaising]{SC019 NWGmcFinalLongORaising}.
The distant boundary at \hyperref[rule-OEMedUnstressedULowering]{SC040
OEMedUnstressedULowering} shows only that word-final \emph{*z}-loss
precedes the later weak-syllable sequence; its placement within that
wider interval follows the handbook chronology after final
\emph{*ō}-raising.

\newpage

\section{Anglo-Frisian brightening}\label{anglo-frisian-brightening}

\subsection{Historical discussion}\label{historical-discussion-25}

Anglo-Frisian Brightening or First Fronting turns low \emph{*a} into
fronted \emph{*æ}-type outcomes outside nasal environments. Later Old
English developments presuppose this fronted stage even where they
partly conceal it. Campbell gives the classical statement of the change,
Hogg supplies the standard modern labels, and Ringe and Taylor establish
its local chronology with breaking and restoration
(\citeproc{ref-Campbell1959}{Campbell 1959, 52}, §131;
\citeproc{ref-Hogg1992}{Hogg 1992, 101, 119};
\citeproc{ref-RingeTaylor2014}{Ringe and Taylor 2014, 157--58, 189--90};
\citeproc{ref-Fulk2018}{Fulk 2018, 73--74}, §§4.12--4.13).

Brightening creates the input to \hyperref[rule-OEBreaking]{SC044
OEBreaking}, while \hyperref[rule-OEARestoration]{SC046 OEARestoration}
later partly reverses its outcome before back vowels.

\subsection{\texorpdfstring{\CAPRRuleHeading{SC043. Fronting of low \emph{*a} outside nasal environments}{AngloFrisianBrightening}}{}}\label{rule-AngloFrisianBrightening}

\begin{ReaderFacingFoma}
define AngloFrisianBrightening [
    AngloFrisianBrighteningUnstressed .o.
    AngloFrisianBrighteningStressed .o.
    AngloFrisianBrighteningLongFinal
];
\end{ReaderFacingFoma}

Two derivations place low \emph{*a} \textgreater{} \emph{*æ} between
unrounding and breaking. Before
\hyperref[rule-PWGmcSurvivingBimoricOUnrounding]{SC042
PWGmcSurvivingBimoricOUnrounding}, PGmc \Recon{rástōz} `rest' yields
\Pred{rasta} rather than expected OE \emph{ræste} `rest'. After
\hyperref[rule-OEBreaking]{SC044 OEBreaking}, PGmc \Recon{sláxaną}
`slay' yields \emph{sleaan | slēaan} rather than expected OE
\emph{slēan} `slay'. The first witness requires brightening to receive
the outcome of the surviving-bimoric \emph{*ō} development; the second
requires breaking to receive the fronted vowel.

\newpage

\chapter{From Anglo-Frisian to Old
English}\label{from-anglo-frisian-to-old-english}

\section{Historical interval}\label{historical-interval-3}

This chapter covers the sound changes that occurred within the Old
English period: the changes that produced attested Old English from the
prehistoric English forms that emerged from the Anglo-Frisian stage. The
starting point is the end of the Anglo-Frisian changes of Chapter 3; the
ending point is attested West Saxon Old English, the primary dialect of
the CAPR corpus.

\section{Scope and dialect variation}\label{scope-and-dialect-variation}

Not every change in this chapter has pan-Old-English scope. Some changes
--- most notably the West Saxon palatal-glide effects (SC016) and West
Saxon palatal umlaut (SC060), the back-mutation rules (SC059), and the
West Saxon diphthong chain (SC031--SC034) --- are specifically West
Saxon or more broadly southern Old English phenomena.

The CAPR derivations target West Saxon Old English citation forms as the
default comparator. Changes that belong to other dialects, or that are
absent from West Saxon, may appear in lexical entries as comparanda
rather than as derivational steps.

The existing reader-facing sound-change sections record which changes
have pan-Old-English scope and which are specifically West Saxon or
Anglian (\citeproc{ref-Campbell1959}{Campbell 1959, sects 1--10};
\citeproc{ref-Hogg1992}{Hogg 1992, sects 1.1--1.15}).

\section{Chapter structure}\label{chapter-structure}

The changes in this chapter fall into several natural historical
subgroups, though the boundaries between them are not always sharp:

\textbf{Early Old English changes linked to the Anglo-Frisian
inheritance}: Changes that feed directly on, or are closely related to,
the brightening of Chapter 3. The \texttt{*awj} glide formation (SC029),
\texttt{*au} fronting (SC030), and the West Saxon diphthong chain
(SC031--SC034) all operate on the vowel inventory shaped by
Anglo-Frisian brightening. OE Breaking (SC044) and a-Restoration (SC046)
similarly presuppose the fronted \texttt{*æ} input.

\textbf{Old English consonantal changes}: Velar palatalization (SC052),
palatalization of \texttt{*sk} (SC051), j-cluster coalescence (SC057),
and related changes produce the characteristically Old English consonant
phonemes. Hogg discusses these as OE consonant changes that are not
broadly West Germanic (\citeproc{ref-Hogg1992}{Hogg 1992, sects
7.18--7.23}).

\textbf{Old English i-umlaut and its context}: The i-umlaut (SC055) is
one of the most productive changes in the Old English nominal and verbal
morphology. Its relative chronology in relation to breaking,
palatalization, and back-mutation is carefully documented in the
existing CAPR chronology evidence audit and individual dossiers
(\citeproc{ref-Campbell1959}{Campbell 1959, sects 193--204};
\citeproc{ref-Hogg1992}{Hogg 1992, sects 5.62--5.68}).

\textbf{Late Old English syllabic reduction and apocope}: High-vowel
apocope (SC063), medial syncope (SC065), and the cluster of late
unstressed-vowel changes (SC069--SC078) represent the later stage of Old
English phonological history, when the syllabic structure of the
language began to shift toward the more reduced profile of Middle
English.

\section{Cascade positions and historical
order}\label{cascade-positions-and-historical-order}

The cascade positions of changes in this chapter broadly follow
historical sequence within the Old English period. Some divergences
exist and are documented in the individual sound-change sections:

\begin{itemize}
\tightlist
\item
  SC041 (PWGmc Final Bare-\texttt{*a} Loss) and SC042 (Surviving Bimoric
  \texttt{*ō} Unrounding) appear late in the cascade but are assigned to
  Chapter 2 (Proto-West Germanic) for historical reasons; their sections
  appear there.
\item
  SC064 (NWGmc \texttt{*-n} Stem \texttt{*n} Loss) carries a Northwest
  Germanic label and is assigned to Chapter 2 historically, but appears
  in the cascade after OE High-Vowel Apocope (SC063). Its section
  cross-references this divergence.
\item
  SC049 (PGmc B Allophony) is assigned to Chapter 1 historically but
  appears in the cascade here; its section cross-references Chapter 1.
\end{itemize}

\section{Sources}\label{sources}

Campbell's \emph{Old English Grammar} is the primary source for the
dating and scope of individual changes in this chapter
(\citeproc{ref-Campbell1959}{Campbell 1959}). Hogg's \emph{Grammar of
Old English} provides modern reassessments and additional
relative-chronology evidence (\citeproc{ref-Hogg1992}{Hogg 1992}). Ringe
and Taylor supply the most detailed relative-chronology analysis for the
earlier portion of the chapter, through back-mutation
(\citeproc{ref-RingeTaylor2014}{Ringe and Taylor 2014, 70--160}). Fulk's
\emph{Comparative Grammar} provides additional coverage for
morphological conditioning (\citeproc{ref-Fulk2018}{Fulk 2018}). For
individual changes, source-specific citations appear in the relevant
sound-change sections.

\section{West Saxon palatal glide before back
vowels}\label{west-saxon-palatal-glide-before-back-vowels}

\subsection{Historical discussion}\label{historical-discussion-26}

West Saxon spellings such as \emph{ġeoc} `yoke', \emph{ġeong} `young',
and \emph{ġeoguþ} `youth' reflect an early Old English development
before back vowels. Campbell gives the most direct handbook statement of
the phenomenon (\citeproc{ref-Campbell1959}{Campbell 1959, 17}, §44).

The sources establish \hyperref[rule-OEWsPalatalGlide]{SC016
OEWsPalatalGlide}, although the checked forms provide only a later
boundary. The rule is computationally positioned before
\hyperref[rule-NWGmcULowering]{SC017 NWGmcULowering} because the
derivation of \emph{ġeoc} `yoke' requires glide insertion before
u-lowering applies. That computational dependency places
\hyperref[rule-OEWsPalatalGlide]{SC016 OEWsPalatalGlide} in the Old
English section of the cascade even though the cascade position precedes
many Northwest Germanic changes.

\subsection{\texorpdfstring{\CAPRRuleHeading{SC016. West Saxon palatal glide before back vowels}{OEWsPalatalGlide}}{}}\label{rule-OEWsPalatalGlide}

\begin{ReaderFacingFoma}
define OEWsPalatalGlide [
    {*j} {*u} -> {*j} {*e} {*u} || .#. _,
    {*j} {*ú} -> {*j} {*é} {*u} || .#. _
] .o. [
    {*ʤ} {*u} -> {*ʤ} {*e} {*u} || .#. _,
    {*ʤ} {*ú} -> {*ʤ} {*é} {*u} || .#. _
] .o. [
    {*ʧ} {*u} -> {*ʧ} {*e} {*u} || .#. _,
    {*ʧ} {*ú} -> {*ʧ} {*é} {*u} || .#. _
] .o. [
    {*ʃ} {*u} -> {*ʃ} {*e} {*u} || .#. _,
    {*ʃ} {*ú} -> {*ʃ} {*é} {*u} || .#. _
];
\end{ReaderFacingFoma}

OE \emph{ġeoc} `yoke' fixes the close relation between glide insertion
before back-vocalic \emph{u} and the following change.

If glide insertion follows \hyperref[rule-NWGmcULowering]{SC017
NWGmcULowering}, PGmc \Recon{júką} `yoke' yields \Pred{ġoc} rather than
expected OE \emph{ġeoc} `yoke'; earlier placement changes no checked
output. The witness therefore dates
\hyperref[rule-OEWsPalatalGlide]{SC016 OEWsPalatalGlide} before
u-lowering without supplying an earlier boundary. The \emph{ġeoc}
`yoke', \emph{ġeong} `young', and \emph{ġeoguþ} `youth' material
establishes the lexical scope of the West Saxon development.

\newpage

\section{Awj glide formation and
au-fronting}\label{awj-glide-formation-and-au-fronting}

\subsection{Historical discussion of awj glide formation and
au-fronting}\label{historical-discussion-of-awj-glide-formation-and-au-fronting}

The \emph{hīeġ} `hay' and \emph{strīeġan} `strew' material undergoes
both changes. Glide formation reshapes the older \emph{awj} sequence,
and fronting then affects the resulting \emph{au}. Campbell's discussion
of these outcomes and Ringe and Taylor's derivations of \emph{hīeġ} and
\emph{strīeġan} describe the same sequence
(\citeproc{ref-Campbell1959}{Campbell 1959, 46}, §120;
\citeproc{ref-RingeTaylor2014}{Ringe and Taylor 2014, 188}).

Glide formation creates the input to fronting; diphthong leveling
follows both.

\subsection{Historical discussion of awj glide
formation}\label{historical-discussion-of-awj-glide-formation}

Older \emph{awj} sequences are the source of forms such as \emph{hīeġ}
`hay' and \emph{strīeġan} `strew'. Campbell treats the relevant
developments directly, and Ringe and Taylor likewise trace the same
material through intermediate \emph{auj}-type stages
(\citeproc{ref-Campbell1959}{Campbell 1959, 46}, §120;
\citeproc{ref-RingeTaylor2014}{Ringe and Taylor 2014, 188}).

The sources establish glide formation, while the witness forms supply
only a later boundary.

\subsection{\texorpdfstring{SC029. Glide formation in \emph{*awj}
(\texttt{OEAwjGlideFormation})}{SC029. Glide formation in  (OEAwjGlideFormation)}}\label{rule-OEAwjGlideFormation}

\begin{ReaderFacingFoma}
define OEAwjGlideFormation [
    {*á} {*w} {*w} {*j} -> {*áu} {*j},
    {*a} {*w} {*w} {*j} -> {*au} {*j},
    {*á} {*w}      {*j} -> {*áu} {*j},
    {*a} {*w}      {*j} -> {*au} {*j}
];
\end{ReaderFacingFoma}

The \emph{hīeġ} `hay' and \emph{strīeġan} `strew' derivations show that
\emph{awj} reshaping prepared the input to fronting. If fronting is
applied first, PGmc \Recon{xáwwją} `hay' yields \Pred{hauġ} rather than
expected OE \emph{hīeġ} `hay', and PGmc \Recon{stráwjaną} `strew' yields
\Pred{strauian} rather than expected \emph{strīeġan} `strew'. Earlier
placement of glide formation changes no checked output, so these forms
supply an upper boundary without a corresponding lower one.

\subsection{Historical discussion of
au-fronting}\label{historical-discussion-of-au-fronting}

Once the glide sequence is in place, \emph{au}-fronting produces the
fronted diphthongal outcomes of the broader West Saxon vowel history.
Campbell describes \emph{au} \textgreater{} \emph{ēa}
(\citeproc{ref-Campbell1959}{Campbell 1959, 53--54}, §135).

Fronting must follow glide formation and precede diphthong leveling,
which applies to a wider set of derivations.

\subsection{\texorpdfstring{SC030. Fronting of \emph{*au}
(\texttt{OEAuFronting})}{SC030. Fronting of  (OEAuFronting)}}\label{rule-OEAuFronting}

\begin{ReaderFacingFoma}
define OEAuFronting [
    {*au} -> {*aeu},
    {*áu} -> {*áeu}
];
\end{ReaderFacingFoma}

Two distinct failure sets confine fronting. Placed before glide
formation, it produces the wrong forms: PGmc \Recon{xáwwją} `hay' yields
\Pred{hauġ} rather than expected OE \emph{hīeġ} `hay', and PGmc
\Recon{stráwjaną} `strew' yields \Pred{strauian} rather than expected
\emph{strīeġan} `strew'. Placed after diphthong leveling, PGmc
\Recon{galáubijaną} `believe', \Recon{bráudą} `bread', and
\Recon{dráugmaz} `dream', together with sixteen other derivations, fail
to produce output at all (\emph{+?}) instead of yielding expected OE
\emph{ġelīefan} `believe', \emph{brēad} `bread', and \emph{drēam}
`dream'. The lexical errors require fronting to follow glide formation,
while the failed derivations require it to precede diphthong leveling.

The later failure set consists of failed derivations, not competing Old
English surface forms.

\newpage

\section{West Saxon diphthong
sequence}\label{west-saxon-diphthong-sequence}

\subsection{Historical discussion of the West Saxon diphthong
sequence}\label{historical-discussion-of-the-west-saxon-diphthong-sequence}

Four distinct developments shape the West Saxon diphthongal field.
Campbell discusses inherited \emph{aw}/\emph{ew} outcomes,
palatal-triggered diphthongization, and later Anglian smoothing in
connected but separate parts of the vowel history; Hogg likewise
distinguishes the palatal-diphthongal developments
(\citeproc{ref-Campbell1959}{Campbell 1959, 46, 53--54, 65--70, 95--96},
§§120, 135--136, 170--176, 185, 223--227; \citeproc{ref-Hogg1992}{Hogg
1992, 106--7, 111--12}).

The closest interaction joins \emph{ww}-simplification and
long-\emph{aw} diphthongization, which together shape \emph{dēaw} `dew'
and \emph{hēawan} `hew'. Diphthong leveling regularizes a wider field,
while long-\emph{ew} diphthongization carries \emph{ēow} into the later
environment of breaking.

\subsection{Historical discussion of WW
simplification}\label{historical-discussion-of-ww-simplification}

West Germanic \emph{ww} sequences lie behind forms such as \emph{dēaw}
`dew' and \emph{hēawan} `hew', and Campbell treats them as part of the
early West Germanic diphthong history
(\citeproc{ref-Campbell1959}{Campbell 1959, 46}, §120).

\hyperref[rule-OEWWSimplification]{SC031 OEWWSimplification} precedes
the later long-diphthong outcomes.

\subsection{\texorpdfstring{SC031. Simplification of \emph{*ww}
sequences
(\texttt{OEWWSimplification})}{SC031. Simplification of  sequences (OEWWSimplification)}}\label{rule-OEWWSimplification}

\begin{ReaderFacingFoma}
define OEWWSimplification [
    {*w} {*w} -> {*w}
];
\end{ReaderFacingFoma}

The \emph{dēaw} `dew' and \emph{hēawan} `hew' derivations establish that
doubled \emph{w} was simplified before the long \emph{ēaw} development.
If \hyperref[rule-OEWWSimplification]{SC031 OEWWSimplification} follows
\hyperref[rule-OEAwLongDiphthong]{SC034 OEAwLongDiphthong}, PGmc
\Recon{dáwwō} `dew' yields \Pred{dawu} rather than expected OE
\emph{dēaw} `dew', and PGmc \Recon{xáwwaną} `hew' yields \Pred{hawan}
rather than expected \emph{hēawan} `hew'. Earlier placement changes no
checked output. The witnesses require simplification before the
long-diphthong change and leave the lower boundary to the broader West
Saxon chronology.

\subsection{Historical discussion of diphthong
leveling}\label{historical-discussion-of-diphthong-leveling}

Forms such as \emph{hēafod} `head' reflect the redistribution of
diphthongal outcomes across a wider set of words. Campbell describes
smoothing and related later monophthongization, although the rule below
is more narrowly conditioned than any single textbook label
(\citeproc{ref-Campbell1959}{Campbell 1959, 95--96}, §§223--227).

The evidence for \hyperref[rule-OEDiphthongLeveling]{SC032
OEDiphthongLeveling} is less self-contained than that for the
\emph{dēaw} `dew' / \emph{hēawan} `hew' developments.

\subsection{\texorpdfstring{SC032. Leveling of diphthongal outputs
(\texttt{OEDiphthongLeveling})}{SC032. Leveling of diphthongal outputs (OEDiphthongLeveling)}}\label{rule-OEDiphthongLeveling}

\begin{ReaderFacingFoma}
define OEDiphthongLeveling [
    {*aeu} -> {*ēa},
    {*áeu} -> {*ēa},
    {*eu} -> {*ēo},
    {*éu} -> {*ēo},
    {*iu} -> {*ēo},
    {*íu} -> {*ēo},
    {*e} {*u} -> {*eo},
    {*é} {*u} -> {*éo},
    {*i} {*u} -> {*eo}
];
\end{ReaderFacingFoma}

The two edges of this interval fail differently. Before
\hyperref[rule-OEAuFronting]{SC030 OEAuFronting}, PGmc
\Recon{galáubijaną} `believe', \Recon{báug} `bow', and \Recon{bráudą}
`bread' produce no output (\emph{+?}) instead of expected OE
\emph{ġelīefan} `believe', \emph{bēag} `bow', and \emph{brēad} `bread',
alongside fifteen other failed derivations. After
\hyperref[rule-OEMedUnstressedULowering]{SC040
OEMedUnstressedULowering}, PGmc \Recon{xáubudą} `head' yields
\Pred{hēafud} rather than expected \emph{hēafod} `head'. Absence at the
lower edge places diphthong leveling after fronting; the wrong surface
form at the upper edge places it before medial unstressed-\emph{u}
lowering.

\subsection{\texorpdfstring{Historical discussion of long
\emph{ēow}}{Historical discussion of long }}\label{historical-discussion-of-long}

The long \emph{ēow} forms of \emph{ċēowan} `chew', \emph{fēower} `four',
and \emph{cnēow} `knee' form part of the West Saxon vowel history,
although their clearest ordering relation points forward. Campbell
describes early \emph{eu} in Old English, and Ringe and Taylor give the
corresponding examples from chew, four, and knee
(\citeproc{ref-Campbell1959}{Campbell 1959, 53--54}, §136;
\citeproc{ref-RingeTaylor2014}{Ringe and Taylor 2014, 188, 202}).

The only checked boundary for \hyperref[rule-OEEwLongDiphthong]{SC033
OEEwLongDiphthong} lies ahead at \hyperref[rule-OEBreaking]{SC044
OEBreaking}.

\subsection{\texorpdfstring{\CAPRRuleHeading{SC033. Long \emph{ēow} before following vowels and weak endings}{OEEwLongDiphthong}}{}}\label{rule-OEEwLongDiphthong}

\begin{ReaderFacingFoma}
define OEEwLongDiphthong [
    {*e} {*w} -> {*ēo} {*w} || _ OEEwLongContext,
    {*i} {*w} -> {*ēo} {*w} || _ OEEwLongContext,
    {*é} {*w} -> {*ēo} {*w} || _ OEEwLongContext,
    {*í} {*w} -> {*ēo} {*w} || _ OEEwLongContext
];
\end{ReaderFacingFoma}

The long \emph{ēow} of \emph{ċēowan} `chew', \emph{fēower} `four', and
\emph{cnēow} `knee' supplies only a terminus ante quem. If
\hyperref[rule-OEEwLongDiphthong]{SC033 OEEwLongDiphthong} follows
\hyperref[rule-OEBreaking]{SC044 OEBreaking}, PGmc \Recon{kéwwaną}
`chew' yields \Pred{ċeowan} rather than expected OE \emph{ċēowan}
`chew', PGmc \Recon{fédwōr} `four' yields \Pred{feower} rather than
expected \emph{fēower} `four', and PGmc \Recon{knéwą} `knee' yields
\Pred{cneow} rather than expected \emph{cnēow} `knee'. Earlier placement
changes no checked output. The sources associate \emph{ew} and \emph{iw}
with the same diphthongal history but furnish no lower boundary.

\subsection{\texorpdfstring{Historical discussion of long
\emph{ēaw}}{Historical discussion of long }}\label{historical-discussion-of-long-1}

After \hyperref[rule-OEWWSimplification]{SC031 OEWWSimplification} has
reduced \emph{ww} to single \emph{w}, the remaining \emph{aw} sequence
can develop into the long \emph{ēaw} seen in \emph{dēaw} `dew' and
\emph{hēawan} `hew'. Campbell treats these outputs in the early
diphthong history of West Germanic and Old English
(\citeproc{ref-Campbell1959}{Campbell 1959, 46, 53--54}, §§120,
135--136). The resulting long diphthong is \emph{ēaw}.

\hyperref[rule-OEAwLongDiphthong]{SC034 OEAwLongDiphthong} follows
\hyperref[rule-OEWWSimplification]{SC031 OEWWSimplification} locally and
must also precede \hyperref[rule-AngloFrisianBrightening]{SC043
AngloFrisianBrightening}.

\subsection{\texorpdfstring{SC034. Long \emph{ēaw} before following
vowels
(\texttt{OEAwLongDiphthong})}{SC034. Long  before following vowels (OEAwLongDiphthong)}}\label{rule-OEAwLongDiphthong}

\begin{ReaderFacingFoma}
define OEAwLongDiphthong [
    {*a} {*w} -> {*ēa} {*w} || _ [EnglishStarVocalic | {*ô}],
    {*á} {*w} -> {*ḗa} {*w} || _ [EnglishStarVocalic | {*ô}]
];
\end{ReaderFacingFoma}

A local feeding relation and a later vowel change confine \emph{aw}
\textgreater{} \emph{ēaw}. Before
\hyperref[rule-OEWWSimplification]{SC031 OEWWSimplification}, PGmc
\Recon{dáwwō} `dew' yields \Pred{dawu} rather than expected OE
\emph{dēaw} `dew', and PGmc \Recon{xáwwaną} `hew' yields \Pred{hawan}
rather than expected \emph{hēawan} `hew'. After
\hyperref[rule-AngloFrisianBrightening]{SC043 AngloFrisianBrightening},
PGmc \Recon{skáwōjaną} `show' yields \Pred{sċawian} rather than expected
OE \emph{sċēawian} `show', PGmc \Recon{skáwōθi} `shows' yields
\Pred{sċawaþ} rather than expected \emph{sċēawaþ} `shows', and PGmc
\Recon{stráwą} `straw' yields \Pred{stræw} rather than expected
\emph{strēaw} `straw'. The \emph{dēaw} and \emph{hēawan} forms require
long-diphthong formation after simplification, while \emph{sċēawian}
requires it before brightening; the handbooks assign the same interval
to the West Saxon development.

\newpage

\section{Prefix and compound
adjustments}\label{prefix-and-compound-adjustments}

\subsection{\texorpdfstring{Historical discussion of prefixal
\emph{*a}-reduction}{Historical discussion of prefixal -reduction}}\label{historical-discussion-of-prefixal--reduction}

Weakly stressed prefixes can lose their older low vowel early in Old
English, and that is the historical setting for
\hyperref[rule-OEPrefixAReduction]{SC035 OEPrefixAReduction}. Campbell
treats the small class of pretonic losses directly, while Ringe and
Taylor's derivation of \Recon{galaubijana} `believe' supplies the
comparative witness for the same development
(\citeproc{ref-Campbell1959}{Campbell 1959, 147}, §354;
\citeproc{ref-RingeTaylor2014}{Ringe and Taylor 2014, 245};
\citeproc{ref-RingeTaylor2014}{2014, 267}).

The rule has a narrow historical range and gives prefixed forms the weak
vowel inherited by later vocalic changes.

\subsection{\texorpdfstring{SC035. Reduction of prefixal \emph{*a}
(\texttt{OEPrefixAReduction})}{SC035. Reduction of prefixal  (OEPrefixAReduction)}}\label{rule-OEPrefixAReduction}

\begin{ReaderFacingFoma}
define OEPrefixAReduction [
    {*a} -> {*ĕ}
        || .#. {*g} _
           [EnglishStarConsonant | EnglishPalatalConsonant]
           EnglishStarVocalic
];
\end{ReaderFacingFoma}

The prefix of \emph{ġelīefan} `believe' supplies the upper boundary for
\emph{*ga-} \textgreater{} \emph{*ge-}. If
\hyperref[rule-OEPrefixAReduction]{SC035 OEPrefixAReduction} follows
\hyperref[rule-AngloFrisianBrightening]{SC043 AngloFrisianBrightening},
PGmc \Recon{galáubijaną} `believe' yields \Pred{ġealīefan} rather than
expected OE \emph{ġelīefan} `believe'. Earlier placement changes no
checked output, so the witness dates prefix reduction before brightening
without locating its beginning.

\subsection{Historical discussion of inter-stress
raising}\label{historical-discussion-of-inter-stress-raising}

\hyperref[rule-OEInterStressRaising]{SC036 OEInterStressRaising} has the
strongest evidence of the three. Campbell's discussion of \emph{weorold}
`world' / \emph{weoruld} `world' and Ringe and Taylor's derivation of
\Recon{weraldu} `world' \textgreater{} \Recon{weruldu} `world'
\textgreater{} OE \emph{weorold} place the rule squarely in the history
of low-stress medial vowels (\citeproc{ref-Campbell1959}{Campbell 1959,
141--42}, §§338--339; \citeproc{ref-RingeTaylor2014}{Ringe and Taylor
2014, 322}, §6.3.3).

The rule changes the vowel between stronger stress peaks, and its
witnesses consequently constrain the relative chronology.

\subsection{\texorpdfstring{\CAPRRuleHeading{SC036. Raising of medial \emph{*a} between stress peaks}{OEInterStressRaising}}{}}\label{rule-OEInterStressRaising}

\begin{ReaderFacingFoma}
define OEInterStressRaising [
    {*a} -> {*u}
        || PGmcStarVowel EnglishStarConsonant* _
           [EnglishStarConsonant - {*j}]+ [{*u}|{*ū}],
    {*à} -> {*u}
];
\end{ReaderFacingFoma}

The two boundaries have unequal force. Before
\hyperref[rule-NWGmcFinalLongORaising]{SC019 NWGmcFinalLongORaising},
PGmc \Recon{sáiwalō} `soul' yields \Pred{sāwel} rather than expected OE
\emph{sāwol} `soul'; after
\hyperref[rule-OEMedUnstressedULowering]{SC040
OEMedUnstressedULowering}, it yields \Pred{sāwul} rather than
\emph{sāwol}, while PGmc \Recon{wír-àldu} `world' yields \Pred{weoruld}
rather than \emph{weorold} `world'. The distant lower boundary places
inter-stress raising after final long-\emph{o} raising, and the local
upper boundary places it before medial unstressed-\emph{u} lowering. In
handbook terms, medial \emph{*a} \textgreater{} \emph{*u} belongs to the
\emph{world}- and \emph{soul}-type low-stress vocalism that followed the
earlier final-vowel changes.

\subsection{Historical discussion of compound linking
syncope}\label{historical-discussion-of-compound-linking-syncope}

Compound members with weakened force often lose or reshape their linking
vowels, and Campbell treats that broad pattern through reduced second
elements, connecting vowels, and obscured compounds
(\citeproc{ref-Campbell1959}{Campbell 1959, 148--49}, §§356--357;
\citeproc{ref-Campbell1959}{Campbell 1959, 153}, §367;
\citeproc{ref-Campbell1959}{Campbell 1959, 159}, §§386--387).

\hyperref[rule-OECompoundLinkingSyncope]{SC037 OECompoundLinkingSyncope}
captures this pattern in compounds such as \emph{reġnboga} `rainbow'.
Its only checked boundary is the immediately following technical
stress-stripping stage, which is not a sound change.

\subsection{\texorpdfstring{\CAPRRuleHeading{SC037. Syncope of compound linking vowels}{OECompoundLinkingSyncope}}{}}\label{rule-OECompoundLinkingSyncope}

\begin{ReaderFacingFoma}
define OECompoundLinkingSyncope [
    [{*a}|{*i}|{*u}] -> 0
        || PGmcStarAcuteVowel OEAnyConsonant+ _
           OEAnyConsonant+ PGmcStarGraveVowel
];
\end{ReaderFacingFoma}

The \emph{reġnboga} `rainbow' test exposes a bookkeeping dependency
rather than a historical sound-change boundary. After SC038
OEStripSecondaryStress, PGmc \Recon{régna-bùgô} `rainbow' yields
\Pred{reġnefoga} rather than expected OE \emph{reġnboga} `rainbow',
because the technical stage has erased the stress information that
licenses syncope. The handbooks instead place weakened compound
junctures with the behavior described under
\hyperref[rule-OEPrefixAReduction]{SC035 OEPrefixAReduction} and
\hyperref[rule-OEInterStressRaising]{SC036 OEInterStressRaising}.

\newpage

\section{Medial unstressed vowel
changes}\label{medial-unstressed-vowel-changes}

\subsection{Historical discussion of medial unstressed vowel
changes}\label{historical-discussion-of-medial-unstressed-vowel-changes}

The history of \emph{wuduwe} `widow' orders these two changes within the
same low-stress vocalic development. Campbell discusses both the
\emph{w}-conditioned \emph{u} forms and the later \emph{weorold} `world'
/ \emph{weoruld} `world' alternation, while Ringe and Taylor give the
same connection comparatively in \emph{*widuwon-}, \Recon{weraldu}
`world', and \Recon{jugunþi} `youth'
(\citeproc{ref-Campbell1959}{Campbell 1959, 92}, §218;
\citeproc{ref-Campbell1959}{Campbell 1959, 140}, §332;
\citeproc{ref-Campbell1959}{Campbell 1959, 141--42}, §§338--339;
\citeproc{ref-RingeTaylor2014}{Ringe and Taylor 2014, 267};
\citeproc{ref-RingeTaylor2014}{Ringe and Taylor 2014, 322}, §6.3.3).

\hyperref[rule-OEWICombinativeUUmlaut]{SC039 OEWICombinativeUUmlaut}
feeds the vowel sequence subsequently reshaped by
\hyperref[rule-OEMedUnstressedULowering]{SC040
OEMedUnstressedULowering}. Initial \emph{w} conditions the first change.

\subsection{\texorpdfstring{SC039. Combinative \emph{*u}-umlaut in
\emph{wi}-forms
(\texttt{OEWICombinativeUUmlaut})}{SC039. Combinative -umlaut in -forms (OEWICombinativeUUmlaut)}}\label{rule-OEWICombinativeUUmlaut}

\begin{ReaderFacingFoma}
define OEWICombinativeUUmlaut [
    {*í} -> {*ú}
        || .#. {*w} _ EnglishStarConsonant [{*u} | {*o}]
];
\end{ReaderFacingFoma}

The \emph{wuduwe} `widow' derivation answers one narrow question about
\emph{wi}-forms. If \hyperref[rule-OEWICombinativeUUmlaut]{SC039
OEWICombinativeUUmlaut} follows
\hyperref[rule-OEMedUnstressedULowering]{SC040
OEMedUnstressedULowering}, PGmc \Recon{wíduwōn} `widow' yields
\Pred{wudowe} rather than expected OE \emph{wuduwe}; earlier placement
changes no checked output. The witness requires combinative u-umlaut to
precede medial lowering and supplies no lower boundary.

\subsection{\texorpdfstring{\CAPRRuleHeading{SC040. Lowering of medial unstressed \emph{*u}}{OEMedUnstressedULowering}}{}}\label{rule-OEMedUnstressedULowering}

\begin{ReaderFacingFoma}
define OEMedUnstressedULowering [
    {*u} -> {*o}
        || [EnglishStarVocalic - [{*u}|{*ū}|{*ú}]]
           [EnglishStarConsonant | EnglishPalatalConsonant]+ _
           [[EnglishStarConsonant | EnglishPalatalConsonant] - {*m}]
];
\end{ReaderFacingFoma}

The two witnesses date medial unstressed \emph{*u} \textgreater{}
\emph{*o} at very different scales. Before
\hyperref[rule-OEWICombinativeUUmlaut]{SC039 OEWICombinativeUUmlaut},
PGmc \Recon{wíduwōn} `widow' yields \Pred{wudowe} rather than expected
OE \emph{wuduwe} `widow'; after
\hyperref[rule-OEUnstressedLongVowelShortening]{SC072
OEUnstressedLongVowelShortening}, PGmc \Recon{júgunθ} `youth' yields
\Pred{ġeogoþ} rather than expected \emph{ġeoguþ} `youth'. The local
\emph{weorold} `world' and widow evidence places lowering after
combinative u-umlaut, while the youth form supplies only the distant
requirement that lowering precede unstressed long-vowel shortening.

\newpage

\section{Breaking and velar-fricative
palatalization}\label{breaking-and-velar-fricative-palatalization}

\subsection{Historical discussion of breaking and velar-fricative
palatalization}\label{historical-discussion-of-breaking-and-velar-fricative-palatalization}

Breaking creates \emph{eo}-type outputs before \emph{h}, \emph{rC}, and
\emph{lC}; velar-fricative palatalization then operates in that reshaped
environment. Campbell, Ringe and Taylor, and Fulk place breaking after
brightening. The following fricative palatalization is more narrowly
conditioned (\citeproc{ref-Campbell1959}{Campbell 1959, 54, 166}, §§139,
405--406; \citeproc{ref-RingeTaylor2014}{Ringe and Taylor 2014, 168--69,
213--14}, §§6.2.1--6.2.3, 6.4.1--6.4.2; \citeproc{ref-Fulk2018}{Fulk
2018, 73--74}, §4.13).

Breaking has the fuller handbook treatment, while velar-fricative
palatalization follows it locally in the \emph{feoh} `cattle' and
\emph{feohtan} `fight' type derivations.

\subsection{\texorpdfstring{SC044. Breaking before \emph{h}, \emph{rC},
and \emph{lC}
(\texttt{OEBreaking})}{SC044. Breaking before , , and  (OEBreaking)}}\label{rule-OEBreaking}

\begin{ReaderFacingFoma}
define OEBreaking OEBreakingA
    .o. OEBreakingE
    .o. OEBreakingI;
\end{ReaderFacingFoma}

Breaking must encounter the vowel created by brightening and must
precede the fricative change seen in \emph{feoh} `fee' and
\emph{feohtan} `fight'. Before
\hyperref[rule-AngloFrisianBrightening]{SC043 AngloFrisianBrightening},
PGmc \Recon{sláxaną} `slay' yields \emph{sleaan | slēaan} rather than
expected OE \emph{slēan} `slay'. After
\hyperref[rule-OEVelarFricativePalatalization]{SC045
OEVelarFricativePalatalization}, PGmc \Recon{féxu} `cattle' yields
\Pred{fehu} rather than expected OE \emph{feoh}, and PGmc
\Recon{féxtaną} `fight' yields \Pred{fehtan} rather than expected
\emph{feohtan}. The two feeding relations place breaking between
brightening and velar-fricative palatalization.

\subsection{\texorpdfstring{\CAPRRuleHeading{SC045. Palatalization of velar fricatives beside front vowels}{OEVelarFricativePalatalization}}{}}\label{rule-OEVelarFricativePalatalization}

\begin{ReaderFacingFoma}
define OEVelarFricativePalatalization [
    {*x} -> {*ç} || _ EnglishStarFrontVowel,
    {*ɣ} -> {*j} || _ EnglishStarFrontVowel,
    {*x} -> {*ç} || EnglishStarFrontVowel _,
    {*ɣ} -> {*j} || EnglishStarFrontVowel _,
    {*x} -> {*ç} || _ {*j},
    {*ɣ} -> {*j} || _ {*j}
]
    .o. EnglishStarAlphabet*;
\end{ReaderFacingFoma}

The local chronology comes from \emph{feoh} `cattle' and \emph{feohtan}
`fight'. Before \hyperref[rule-OEBreaking]{SC044 OEBreaking},
palatalization of \emph{*x} and \emph{*ɣ} beside front vowels or
\emph{*j} makes PGmc \Recon{féxu} `cattle' yield \Pred{fehu} rather than
expected OE \emph{feoh}, and PGmc \Recon{féxtaną} `fight' yield
\Pred{fehtan} rather than expected \emph{feohtan}. The distant upper
limit comes from \emph{six} `six': after
\hyperref[rule-OEWsPalatalUmlaut]{SC060 OEWsPalatalUmlaut}, PGmc
\Recon{séxs} `six' yields \Pred{sihs} rather than expected OE
\emph{six}. Breaking therefore feeds velar-fricative palatalization
directly, while palatal umlaut supplies only the broader upper limit.

\newpage

\section{A-restoration and nasal
changes}\label{a-restoration-and-nasal-changes}

\subsection{Historical discussion of
A-restoration}\label{historical-discussion-of-a-restoration}

Campbell's restoration of \emph{a} before following back vowels and
Ringe and Taylor's later retraction describe the same post-brightening
development (\citeproc{ref-Campbell1959}{Campbell 1959, 60--61},
§§157--159; \citeproc{ref-RingeTaylor2014}{Ringe and Taylor 2014,
189--90}, §6.3.1; \citeproc{ref-Fulk2018}{Fulk 2018, 74}, §4.13). Some
outcomes of Anglo-Frisian fronting survive only in environments where
restoration does not return them to back \emph{a}.

\hyperref[rule-OEARestoration]{SC046 OEARestoration} has firmer handbook
support than the two following nasal rules.

\subsection{\texorpdfstring{\CAPRRuleHeading{SC046. Restoration of \emph{*a} before following back vowels}{OEARestoration}}{}}\label{rule-OEARestoration}

\begin{ReaderFacingFoma}
define OEARestoration (
    {*æ} -> {*a} || _
        OEARestorationIntervening OEARestorationTriggerVowel
        - OEARestorationIntervening OEARestorationWeakTailVowel
);
\end{ReaderFacingFoma}

Restoration must receive fronted \emph{*æ} and return \emph{*a} before
the nasal-tail changes. Before
\hyperref[rule-AngloFrisianBrightening]{SC043 AngloFrisianBrightening},
PGmc \Recon{bákaną} `bake' yields \Pred{bæcan} rather than expected OE
\emph{bacan} `bake', and PGmc \Recon{fáraną} `fare' yields \Pred{færan}
rather than expected \emph{faran} `fare'. After
\hyperref[rule-OESecondaryNasalization]{SC048 OESecondaryNasalization},
\Recon{bákaną} again yields \Pred{bæcan} instead of \emph{bacan}, while
PGmc \Recon{wádaną} `wade' yields \Pred{wædan} instead of \emph{wadan}
`wade'. These independent witness pairs place restoration after
brightening and before secondary nasalization.

\subsection{Historical discussion of heavy-syllable nasal loss and
secondary
nasalization}\label{historical-discussion-of-heavy-syllable-nasal-loss-and-secondary-nasalization}

Heavy-syllable nasal apocope removes the final nasalized vowel;
secondary nasalization then marks the preceding \emph{a} before final
\emph{n}. The handbooks do not isolate both developments under equally
prominent labels. Campbell describes later nasal loss and the
back-mutation environment; Ringe and Taylor provide the later relation
to back mutation (\citeproc{ref-Campbell1959}{Campbell 1959, 86, 166},
§§205--206, 403; \citeproc{ref-RingeTaylor2014}{Ringe and Taylor 2014,
319}, §6.9.4).

The reciprocal failure set fixes the order: apocope removes the ending
before secondary nasalization acts on the remaining structure.
Restoration receives the fuller historical treatment in the handbooks.

\subsection{\texorpdfstring{\CAPRRuleHeading{SC047. Heavy-syllable nasal apocope of final \emph{*ą}}{OEHeavySyllableNasalApocope}}{}}\label{rule-OEHeavySyllableNasalApocope}

\begin{ReaderFacingFoma}
define OEHeavySyllableNasalApocope [
    {*ą} -> 0 || OEAnyConsonant _ .#.
];
\end{ReaderFacingFoma}

The evidence for final nasalized \emph{*ą} loss is sharply asymmetric.
Before \hyperref[rule-OEAwLongDiphthong]{SC034 OEAwLongDiphthong}, the
single PGmc witness \Recon{stráwą} `straw' yields \Pred{stræw} rather
than expected OE \emph{strēaw} `straw'. After
\hyperref[rule-OESecondaryNasalization]{SC048 OESecondaryNasalization},
PGmc \Recon{bákaną} `bake' yields \Pred{bacen} rather than expected OE
\emph{bacan} `bake', and PGmc \Recon{bíndaną} `bind' yields
\Pred{binden} rather than expected \emph{bindan} `bind', alongside a
broad \emph{-en} failure set. One lower witness places apocope after
long-diphthong formation; many reciprocal upper failures place it before
secondary nasalization.

\subsection{\texorpdfstring{\CAPRRuleHeading{SC048. Secondary nasalization before final \emph{*n}}{OESecondaryNasalization}}{}}\label{rule-OESecondaryNasalization}

\begin{ReaderFacingFoma}
define OESecondaryNasalization [
    {*a} -> {*ą} || _ {*n} .#.
];
\end{ReaderFacingFoma}

The broad \emph{-an}/\emph{-en} split fixes the lower boundary of final
\emph{*a} nasalization before \emph{n}. Before
\hyperref[rule-OEHeavySyllableNasalApocope]{SC047
OEHeavySyllableNasalApocope}, PGmc \Recon{bákaną} `bake' yields
\Pred{bacen} rather than expected OE \emph{bacan} `bake', and PGmc
\Recon{bíndaną} `bind' yields \Pred{binden} rather than expected
\emph{bindan} `bind'. The upper boundary comes from back mutation. After
\hyperref[rule-OEBackMutation]{SC059 OEBackMutation}, PGmc
\Recon{stélaną} `steal' yields \Pred{steolan} rather than expected OE
\emph{stelan} `steal', and PGmc \Recon{wébaną} `weave' yields
\Pred{weofan} rather than expected \emph{wefan} `weave'. Reciprocal
nasal-tail failures place secondary nasalization after apocope, and the
later mutation witnesses place it before back mutation;
\hyperref[rule-OEARestoration]{SC046 OEARestoration} retains the
clearest independent historical support.

\newpage

\section{\texorpdfstring{Palatalization of \emph{*sk} to
\emph{*sc}}{Palatalization of  to }}\label{palatalization-of-to}

\subsection{Historical discussion}\label{historical-discussion-27}

The palatalization of \emph{*sk} to Old English \emph{*sc} is one of the
recognizable early cluster changes in the larger palatalization zone.
Campbell distinguishes the cluster from plain velars when he remarks
that \emph{*sk} is especially prone to palatalization and assibilation
(\citeproc{ref-Campbell1959}{Campbell 1959, 278}, §440). Hogg gives the
same change a clearer structural place by treating \emph{*sk} beside the
palatalization of plain velars and before the later West Saxon
diphthongal developments (\citeproc{ref-Hogg1992}{Hogg 1992, 106--7,
111--12}). Ringe and Taylor make the same sequence explicit when they
distinguish the earlier palatalization of velars and \emph{*sk} from the
later diphthongization after already palatal consonants
(\citeproc{ref-RingeTaylor2014}{Ringe and Taylor 2014, 213--16},
§§6.4.1, 6.5.1).

Luick places the cluster change within a broader early movement toward
palatal articulation, while still allowing later vowel consequences to
form a different chapter of the history (\citeproc{ref-Luick1914}{Luick
1914-\/-1940, 157}, §168). Fulk's summary is the most concise warning
against overextension: Old English \emph{*sc} is palatal except in the
well-known back-vowel environments that preserve harder outcomes
(\citeproc{ref-Fulk2018}{Fulk 2018, 28}). The result is a historically
clear rule, but not an identity between the cluster change and the later
umlautal developments.

\subsection{\texorpdfstring{SC051. Palatalization of \emph{*sk} to
\emph{*sc}
(\texttt{OESkPalatalization})}{SC051. Palatalization of  to  (OESkPalatalization)}}\label{rule-OESkPalatalization}

\begin{ReaderFacingFoma}
define OESkPalatalization [
    {*s} {*k} -> {*ʃ} || .#. _
] .o. [
    {*s} {*k} -> {*ʃ} || EnglishStarFrontVowel _ (EnglishStarConsonant | .#.)
] .o. [
    {*s} {*k} -> {*ʃ} || (EnglishStarConsonant | .#.) _ EnglishStarFrontVowel
] .o. [
    {*s} {*k} -> {*ʃ} || _ {*j}
] .o. [
    {*s} {*k} -> {*ʃ} || {*j} _
];
\end{ReaderFacingFoma}

The non-fronted vowels of \emph{flasce} `flask' and \emph{wascan} `wash'
fix the lower boundary of \emph{*sk} \textgreater{} \emph{*sc}. Before
\hyperref[rule-OEARestoration]{SC046 OEARestoration}, the forms are
fronted too soon, yielding \emph{flæsce} `flask' and \emph{wæscan}
`wash' rather than expected OE \emph{flasce} and \emph{wascan}. This
places \hyperref[rule-OESkPalatalization]{SC051 OESkPalatalization}
after restoration.

Five witnesses establish the upper boundary collectively. The palatal
cluster must already underlie \emph{sċeaft} `shaft', \emph{sċēar}
`shear', \emph{sċēaþ} `sheath', \emph{sċēap} `sheep', and \emph{sċield}
`shield' before \hyperref[rule-OEWsPalatalDiphthongization]{SC056
OEWsPalatalDiphthongization}. The \emph{*sċea-* 'sea'}/\emph{*sċie-*}
set therefore places cluster palatalization before the West Saxon vowel
change. The cluster change occupies the same palatalization zone as
\hyperref[rule-OEVelarPalatalization]{SC052 OEVelarPalatalization} while
remaining distinct from plain-velar palatalization and the later vowel
changes.

\newpage

\section{Velar palatalization before front
vowels}\label{velar-palatalization-before-front-vowels}

\subsection{Historical discussion}\label{historical-discussion-28}

Luick places the change inside a broad early palatalizing movement.
Under the heading ``Frühe Verschiebungen in palataler Richtung,'' he
treats English \texttt{k} and \texttt{g} before bright vowels together
with the larger field of palatal effects (\citeproc{ref-Luick1914}{Luick
1914-\/-1940, 157}, §168). His emphasis falls on the environment first:
velars before bright vowels and in the vicinity of the palatal glide
belong to one early phonological sequence. The examples associated with
that sequence, such as \emph{ceaster} `town', \emph{geaf} `gave',
\emph{giefan} `give', and \emph{giest} `guest', already show that
consonantal palatalization and later vowel effects stand close together
historically, even when they must be distinguished analytically
(\citeproc{ref-Luick1914}{Luick 1914-\/-1940, 157--67}, §§168--182).

Campbell narrows the picture by distinguishing plain velars from the
especially palatal-prone \texttt{sk} cluster. His remark that ``{[}sk{]}
is more prone to palatalization and assibilation than {[}k{]}'' is
brief, but it makes clear that different members of the larger palatal
field need not behave identically (\citeproc{ref-Campbell1959}{Campbell
1959, 278}, §440). Elsewhere in the same part of the grammar he uses
forms such as \emph{cild} `child', \emph{dæg} `day', \emph{giefan}
`give', and \emph{giest} `guest', which show how palatalized velars,
palatal influence, and later umlautal outcomes meet in the same region
of the lexicon without collapsing into one process
(\citeproc{ref-Campbell1959}{Campbell 1959, 69--72, 89}, §§170,
190--191).

Hogg makes the conditioning sharper still. He states that the change
takes place when the velar consonant is adjacent to and in the same
syllable as a front vowel or the palatal consonant \texttt{j}
(\citeproc{ref-Hogg1992}{Hogg 1992, 103--4}). This formulation replaces
a broad list of palatal outcomes with a phonological environment defined
by adjacency and syllable structure.

Ringe and Taylor make the chronological relation still clearer. When
they write that ``after initial velars and \emph{*sk} had been
palatalized'' West-Saxon diphthongization follows, plain velar
palatalization becomes an earlier consonantal stage presupposed by later
vowel developments (\citeproc{ref-RingeTaylor2014}{Ringe and Taylor
2014, 215}, §6.5.1). Their own examples of the plain-velar rule, such as
\emph{weccan} `wake', \emph{licgan} `lie', \emph{lecgan} `lay',
\emph{secg} `retainer', \emph{ecg} `edge', \emph{wicg} `horse', and
\emph{brycg} `bridge', illustrate the same point in lexical detail:
front vowels and \texttt{j} create the palatal environment in which
plain \texttt{k} and \texttt{g} cease to behave as plain velars
(\citeproc{ref-RingeTaylor2014}{Ringe and Taylor 2014, 213--14},
§6.4.1).

Luick describes a broad early movement; Campbell distinguishes plain
velars from the \texttt{sk} complex; Hogg specifies the adjacency and
syllable conditions; and Ringe and Taylor order the plain-velar change
before West-Saxon diphthongization. Plain-velar palatalization thus
forms part of a wider palatalizing environment without being identical
to its neighboring changes.

\subsection{\texorpdfstring{\CAPRRuleHeading{SC052. Palatalization of \emph{*k} before front vowels and \emph{*j}}{OEVelarPalatalizationKFront}}{}}\label{rule-OEVelarPalatalizationKFront}

\begin{ReaderFacingFoma}
define OEVelarPalatalizationKFront [
    {*k} -> {*ʧ} || .#. _ EnglishStarFrontVowel,
    {*k} -> {*ʧ} || _ [{*i} | {*ī}],
    {*k} -> {*ʧ} || _ {*ḯ},
    {*k} -> {*ʧ} || [{*i} | {*ī}] _ EnglishStarFrontVowel,
    {*k} -> {*ʧ} || {*ḯ} _ EnglishStarFrontVowel,
    {*k} -> {*ʧ} || [{*i} | {*ī}] _ .#.,
    {*k} -> {*ʧ} || {*ḯ} _ .#.
] .o. [
    {*k} {*k} -> {*ʧ} {*ʧ} || _ {*j}
] .o. [
    {*k} -> {*ʧ} || _ {*j}
] ;
\end{ReaderFacingFoma}

The \emph{weccan} `wake', \emph{licgan} `lie', and \emph{lecgan} `lay'
set identifies front vowels and \texttt{j} as the environment for
palatalization of \texttt{k} (\citeproc{ref-RingeTaylor2014}{Ringe and
Taylor 2014, 213--14}, §6.4.1). These forms establish the conditioning;
different witnesses establish the chronology.

Applied before Sievers-law syncope, PGmc \Recon{strákkijaną} `stretch'
yields \Pred{strecċan} rather than expected OE \emph{streċċan}
`stretch'. Applied after i-umlaut fronting, PGmc \Recon{kūi} `cow' and
\Recon{lúnganjō} `lungs' yield \emph{ċȳ} `cows' and \emph{lunġen}
`lungs' rather than expected OE \emph{cȳ} `cows' and \emph{lungen}
`lungs'. The front-vowel \texttt{k} change therefore follows Sievers-law
syncope and precedes i-umlaut fronting.

\subsection{\texorpdfstring{\CAPRRuleHeading{SC052. Velar palatalization before front vowels}{OEVelarPalatalization}}{}}\label{rule-OEVelarPalatalization}

\begin{ReaderFacingFoma}
define OEVelarPalatalization [
    OEVelarPalatalizationKFront
] .o. [
    {*g} -> {*ʤ} || _ EnglishStarFrontVowel,
    {*g} -> {*ʤ} || EnglishStarFrontVowel _ .#.,
    {*g} -> {*ʤ} || EnglishStarFrontVowel _ EnglishStarFrontVowel,
    {*g} -> {*ʤ} || EnglishStarFrontVowel _ [EnglishStarConsonant - {*j}],
    {*g} {*g} -> {*ʤ} {*ʤ} || _ {*j}
] .o. [
    {*g} -> {*ʤ} || _ {*j}
];
\end{ReaderFacingFoma}

Plain \texttt{k} and \texttt{g} palatalization in front-vocalic and
\texttt{j}-adjacent environments follows \texttt{sk}-palatalization and
occupies a sharply defined pre-umlaut interval. Applied before
Sievers-law syncope, PGmc \Recon{strákkijaną} `stretch' yields
\Pred{strecċan} rather than expected OE \emph{streċċan} `stretch'.
Applied after general i-umlaut, PGmc \Recon{kūi} `cow' yields \Pred{ċȳ}
rather than expected \emph{cȳ} `cows', and PGmc \Recon{lúnganjō} `lungs'
yields \Pred{lunġen} rather than expected \emph{lungen} `lungs'. These
witnesses place velar palatalization after Sievers-law syncope and
before umlaut.

Luick, Campbell, and Ringe and Taylor place \emph{cild} `child' and
\emph{dæg} `day' in a consonantal palatalization that precedes later
vowel fronting (\citeproc{ref-Luick1914}{Luick 1914-\/-1940, 157}, §168;
\citeproc{ref-Campbell1959}{Campbell 1959, 278}, §440;
\citeproc{ref-RingeTaylor2014}{Ringe and Taylor 2014, 203--15}, §§6.4.1,
6.5.1). The umlautal developments therefore receive plain \texttt{k} and
\texttt{g} already reshaped beside front vowels and \texttt{j}.

The \texttt{sk} change belongs to the same palatalizing region with a
separate scope. The \emph{streċċan} `stretch' evidence establishes a
specific dependency on earlier syncope; it does not merge the two
changes into one process.

\newpage

\section{\texorpdfstring{Post-velar \emph{*w}-loss and loss of \emph{*w}
before final
\emph{*i}}{Post-velar -loss and loss of  before final }}\label{post-velar--loss-and-loss-of-before-final}

\subsection{\texorpdfstring{Historical discussion of early
\emph{*w}-loss before
umlaut}{Historical discussion of early -loss before umlaut}}\label{historical-discussion-of-early--loss-before-umlaut}

The first rule is a narrow loss of \emph{*w} after velars in the
\emph{*ngw} sequence. Ringe and Taylor derive PGmc \Recon{singwan}
`sing' to Old English \emph{singan} `sing'
(\citeproc{ref-RingeTaylor2014}{Ringe and Taylor 2014, 214}, §6.4.2).
This comparative evidence establishes the change, although no checked
form fixes its order relative to a neighboring rule.

The second rule is historically more legible. Campbell notes the
recurring loss of \emph{*w} before \emph{*i} in unstressed position
(\citeproc{ref-Campbell1959}{Campbell 1959, 167}, §406). Ringe and
Taylor trace the development of \emph{sǣ} `sea' from earlier
\emph{*saiwi-} / \emph{*sawi-} (\citeproc{ref-RingeTaylor2014}{Ringe and
Taylor 2014, 257}, §6.7.1), and Luick gives the same trajectory in his
own historical grammar (\citeproc{ref-Luick1914}{Luick 1914-\/-1940,
173}, §187). The first rule is restricted to the \emph{*ngw} sequence;
the second has a specific lexical witness and defined earlier and later
limits.

\subsection{\texorpdfstring{SC053. Loss of \emph{*w} after velars
(\texttt{OEPostVelarWLoss})}{SC053. Loss of  after velars (OEPostVelarWLoss)}}\label{rule-OEPostVelarWLoss}

\begin{ReaderFacingFoma}
define OEPostVelarWLoss [
    {*w} -> 0 || {*n} {*g} _
];
\end{ReaderFacingFoma}

The comparative development \texttt{*singwan\ \textgreater{}\ singan}
establishes narrow post-velar \emph{*w}-loss in the \emph{*ngw}
sequence, yielding \emph{singan} `sing'. Moving
\hyperref[rule-OEPostVelarWLoss]{SC053 OEPostVelarWLoss} earlier or
later leaves every checked output unchanged. Its pre-umlaut position
therefore rests on comparative evidence, while the present lexicon
supplies no neighboring boundary.

\subsection{\texorpdfstring{SC054. Loss of \emph{*w} before final
\emph{*i}
(\texttt{OEWLossBeforeI})}{SC054. Loss of  before final  (OEWLossBeforeI)}}\label{rule-OEWLossBeforeI}

\begin{ReaderFacingFoma}
define OEWLossBeforeI [
    {*w} -> 0 || EnglishStarVocalic _ {*i} .#.
];
\end{ReaderFacingFoma}

The history of \emph{sǣ} `sea' explains why non-initial \emph{*w}
disappeared before final unstressed \emph{*i}. Campbell describes the
loss, Ringe and Taylor derive the form from
\emph{*saiwi-}/\emph{*sawi-}, and Luick gives the parallel trajectory
(\citeproc{ref-Campbell1959}{Campbell 1959, 167}, §406;
\citeproc{ref-RingeTaylor2014}{Ringe and Taylor 2014, 257}, §6.7.1;
\citeproc{ref-Luick1914}{Luick 1914-\/-1940, 173}, §187). Loss of the
glide allowed the preceding vowel to undergo the later fronting and
lengthening.

The same witness supplies two distant limits. Before
\hyperref[rule-PGmcFinalZDeletion]{SC020 PGmcFinalZDeletion} or after
\hyperref[rule-OEHighVowelApocope]{SC063 OEHighVowelApocope},
\hyperref[rule-OEWLossBeforeI]{SC054 OEWLossBeforeI} yields \Pred{sǣw}
rather than expected OE \emph{sǣ} `sea'. The loss must therefore follow
final \emph{z}-deletion and precede high-vowel apocope, while its exact
position within that broad interval remains source-based.

\newpage

\section{The Old English i-umlaut and West Saxon palatal
diphthongization}\label{the-old-english-i-umlaut-and-west-saxon-palatal-diphthongization}

\subsection{\texorpdfstring{Historical discussion of i-umlaut
\CAPRHeadingBreak and West Saxon palatal
diphthongization}{Historical discussion of i-umlaut and West Saxon palatal diphthongization}}\label{historical-discussion-of-i-umlaut-and-west-saxon-palatal-diphthongization}

Luick gives the change its traditional scale:

\begin{quote}
Der wichtigste Fall von palataler Beeinflussung \ldots{} war die
Veränderung der urenglischen Vokale durch i oder j der Folgesilbe.

(\citeproc{ref-Luick1914}{Luick 1914-\/-1940, 166--67}, §182)
\end{quote}

Campbell gives the most compact classical formulation in English when he
writes that ``the process known as i-umlaut or i-mutation operates on
practically all the sounds which it could theoretically affect in OE''
(\citeproc{ref-Campbell1959}{Campbell 1959, 69}, §190). He immediately
defines the core conditioning environment as a following \texttt{i} or
\texttt{j}, and he goes on to trace the consequences across much of the
vowel system, including forms such as \emph{giest} `guest',
\emph{giefan} `give', \emph{hierde} `shepherd', and \emph{ieldra}
`older' (\citeproc{ref-Campbell1959}{Campbell 1959, 69--72},
§§190--197).

Hogg continues in the same vein: ``we come now to a change which is
almost as uncontroversial as it is important''
(\citeproc{ref-Hogg1992}{Hogg 1992, 112}). His examples, such as
\emph{bryd} `bride', \emph{trymman} `strengthen', \emph{bedd} `bed',
\emph{ciest} `chest', and \emph{wiersa} `worse', likewise emphasize that
the change is a broad redistribution of vowel quality across the Old
English vowel system (\citeproc{ref-Hogg1992}{Hogg 1992, 112--14}).

The narrower palatal-diphthongal material is described differently.
Ringe and Taylor treat West-Saxon diphthongization after initial
palatals as a distinct process (\citeproc{ref-RingeTaylor2014}{Ringe and
Taylor 2014, 215}, §6.5.1), and Fulk is even more explicit about its
chronological delicacy when he calls it ``diphthongization by initial
palatal consonants (which precedes front umlaut but not breaking)''
(\citeproc{ref-Fulk2018}{Fulk 2018, 74}, §4.13). Ringe and Taylor's
examples such as \emph{gieldan} `pay', \emph{scield} `shield', and
\emph{scieppan} `create' show that this narrower process is triggered by
already palatal consonants and leads to specifically West-Saxon
diphthongal outputs (\citeproc{ref-RingeTaylor2014}{Ringe and Taylor
2014, 215--16}, §6.5.1).

Luick, Campbell, and Hogg treat i-umlaut as a system-wide change. Ringe
and Taylor and Fulk distinguish from it a narrower West-Saxon process
affecting words after initial palatals. The two changes act in different
environments and produce different lexical consequences.

\subsection{\texorpdfstring{SC055. Fronting under i-umlaut
(\texttt{OEIUmlautFronting})}{SC055. Fronting under i-umlaut (OEIUmlautFronting)}}\label{rule-OEIUmlautFronting}

\begin{ReaderFacingFoma}
define OEIUmlautFronting [
    {*a} -> {*æ} || _ EnglishIUmlautIntervening EnglishIUmlautTrigger,
    {*ā} -> {*ǣ} || _ EnglishIUmlautIntervening EnglishIUmlautTrigger,
    {*e} -> {*i} || _ EnglishIUmlautIntervening EnglishIUmlautTrigger,
    {*o} -> {*e} || _ EnglishIUmlautIntervening EnglishIUmlautTrigger,
    {*ō} -> {*ē} || _ EnglishIUmlautIntervening EnglishIUmlautTrigger,
    {*u} -> {*y} || _ EnglishIUmlautIntervening EnglishIUmlautTrigger,
    {*ū} -> {*ȳ} || _ EnglishIUmlautIntervening EnglishIUmlautTrigger,
    {*á} -> {*æ} || _ EnglishIUmlautIntervening EnglishIUmlautTrigger,
    {*é} -> {*i} || _ EnglishIUmlautIntervening EnglishIUmlautTrigger,
    {*ó} -> {*e} || _ EnglishIUmlautIntervening EnglishIUmlautTrigger,
    {*ú} -> {*y} || _ EnglishIUmlautIntervening EnglishIUmlautTrigger
];
\end{ReaderFacingFoma}

The breadth of i-umlaut appears in lexical classes that share only a
following high front vocoid. The forms \emph{fylgan} `follow',
\emph{gylden} `golden', \emph{wyrm} `worm', and \emph{giest} `guest'
exemplify the same \texttt{i}- or \texttt{j}-conditioned fronting across
different vowels (\citeproc{ref-RingeTaylor2014}{Ringe and Taylor 2014,
222}, §6.6.1; \citeproc{ref-Campbell1959}{Campbell 1959, 69--72},
§§190--191).

The cow and lung forms establish the lower boundary. If fronting
precedes velar palatalization, PGmc \Recon{kūi} `cow' yields \Pred{ċȳ}
rather than expected OE \emph{cȳ} `cows', and \Recon{lúnganjō} `lungs'
yields \Pred{lunġen} rather than expected OE \emph{lungen} `lungs'. The
consonantal change must therefore precede fronting.

The gift and sheath forms establish the upper boundary. If West Saxon
palatal diphthongization precedes fronting, PGmc \Recon{géftiz} `gift'
yields \Pred{ġieft} rather than expected OE \emph{ġift} `gift', and
\Recon{skáiθiz} `sheath' yields \Pred{sċǣþ} rather than expected
\emph{sċēaþ} `sheath'. Fronting consequently follows velar
palatalization and precedes the West Saxon change; the other components
of i-umlaut share those bounds.

\subsection{\texorpdfstring{SC055. Raising under i-umlaut
(\texttt{OEIUmlautRaising})}{SC055. Raising under i-umlaut (OEIUmlautRaising)}}\label{rule-OEIUmlautRaising}

\begin{ReaderFacingFoma}
define OEIUmlautRaising [
    {*æ} -> {*e} || _ EnglishIUmlautIntervening EnglishIUmlautTrigger
];
\end{ReaderFacingFoma}

Raising of umlauted \texttt{æ} to \texttt{e} continues the same
assimilatory event as fronting and therefore shares the chronology of
general i-umlaut.

The same four forms fix both boundaries. If raising precedes velar
palatalization, \Recon{kūi} `cow' yields \Pred{ċȳ} instead of expected
\emph{cȳ} `cows' and \Recon{lúnganjō} `lungs' yields \Pred{lunġen}
instead of expected \emph{lungen} `lungs'. If West Saxon palatal
diphthongization precedes raising, \Recon{géftiz} `gift' yields
\Pred{ġieft} rather than expected \emph{ġift} `gift', and
\Recon{skáiθiz} `sheath' yields \Pred{sċǣþ} rather than expected
\emph{sċēaþ} `sheath'. These forms place raising after velar
palatalization and before West Saxon palatal diphthongization.

The sources do not describe umlaut as simple fronting alone. Campbell
notes that the low front vowel changes again before \texttt{m} and
\texttt{n} in most dialects (\citeproc{ref-Campbell1959}{Campbell 1959,
69}, §190), and Hogg likewise treats short front vowels as part of the
same assimilatory system (\citeproc{ref-Hogg1992}{Hogg 1992, 112}).

\subsection{\texorpdfstring{SC055. Diphthongal outcomes under i-umlaut
(\texttt{OEIUmlautDiphthong})}{SC055. Diphthongal outcomes under i-umlaut (OEIUmlautDiphthong)}}\label{rule-OEIUmlautDiphthong}

\begin{ReaderFacingFoma}
define OEIUmlautDiphthong [
    {*ea} -> {*ie} || _ EnglishIUmlautIntervening EnglishIUmlautTrigger,
    {*ēa} -> {*īe} || _ EnglishIUmlautIntervening EnglishIUmlautTrigger,
    {*io} -> {*ie} || _ EnglishIUmlautIntervening EnglishIUmlautTrigger,
    {*īo} -> {*īe} || _ EnglishIUmlautIntervening EnglishIUmlautTrigger,
    {*eo} -> {*ie} || _ EnglishIUmlautIntervening EnglishIUmlautTrigger,
    {*ēo} -> {*īe} || _ EnglishIUmlautIntervening EnglishIUmlautTrigger,
    {*éa} -> {*íe} || _ EnglishIUmlautIntervening EnglishIUmlautTrigger,
    {*éo} -> {*íe} || _ EnglishIUmlautIntervening EnglishIUmlautTrigger,
    {*ío} -> {*íe} || _ EnglishIUmlautIntervening EnglishIUmlautTrigger
];
\end{ReaderFacingFoma}

Diphthongal outcomes belong to the same system-wide assimilation as
simple-vowel fronting and raising. All three therefore belong to a
single historical event.

The relevant examples are the recurring West-Saxon \texttt{ie} forms
cited in the handbooks, including \emph{giest} `guest', \emph{giefan}
`give', and \emph{hierde} `shepherd' in Campbell and \emph{ciest}
`chest' in Hogg (\citeproc{ref-Campbell1959}{Campbell 1959, 69--72,
78--80}, §§190--191, 248--251; \citeproc{ref-Hogg1992}{Hogg 1992,
112--14}). These diphthongal outcomes form a distinct part of the
general umlautal development alongside simple fronting.

The chronology comes from the cow/lung and gift/sheath contrasts. Placed
before velar palatalization, diphthongal mutation over-palatalizes
\Recon{kūi} `cow' and \Recon{lúnganjō} `lungs'; placed after West Saxon
palatal diphthongization, it yields \Pred{ġieft} and \Pred{sċǣþ} instead
of expected \emph{ġift} `gift' and \emph{sċēaþ} `sheath'. These failures
place diphthongal mutation after velar palatalization and before West
Saxon palatal diphthongization.

\subsection{\texorpdfstring{SC055. The composite i-umlaut rule
(\texttt{OEIUmlaut})}{SC055. The composite i-umlaut rule (OEIUmlaut)}}\label{rule-OEIUmlaut}

\begin{ReaderFacingFoma}
define OEIUmlaut OEIUmlautFronting
    .o. OEIUmlautRaising
    .o. OEIUmlautDiphthong;
\end{ReaderFacingFoma}

The literature presents fronting, raising, and diphthongal mutation as
effects of one historical development. They consequently occupy a single
place in the Old English chronology.

The lower boundary is consonantal. If general umlaut precedes velar
palatalization, PGmc \Recon{kūi} `cow' yields \Pred{ċȳ} rather than
expected \emph{cȳ} `cows', and PGmc \Recon{lúnganjō} `lungs' yields
\Pred{lunġen} rather than expected \emph{lungen} `lungs'. These
over-palatalized forms place general umlaut after velar palatalization.

The upper boundary separates general umlaut from the narrower West Saxon
process. If West Saxon palatal diphthongization precedes umlaut, PGmc
\Recon{géftiz} `gift' yields \Pred{ġieft} rather than expected OE
\emph{ġift} `gift', and \Recon{skáiθiz} `sheath' yields \Pred{sċǣþ}
rather than expected \emph{sċēaþ} `sheath'. Together the two witness
pairs place general umlaut after velar palatalization and before the
West Saxon process.

\subsection{\texorpdfstring{\CAPRRuleHeading{SC056. West Saxon palatal diphthongization}{OEWsPalatalDiphthongization}}{}}\label{rule-OEWsPalatalDiphthongization}

\begin{ReaderFacingFoma}
define OEWsPalatalDiphthongization [
    {*æ} -> {*ea} || .#. [{*ʧ} | {*ʤ} | {*ʃ} | {*j}] _ [EnglishStarConsonant | EnglishPalatalConsonant | .#.],
    {*ǣ} -> {*ēa} || .#. [{*ʧ} | {*ʤ} | {*ʃ} | {*j}] _ [EnglishStarConsonant | EnglishPalatalConsonant | .#.],
    {*e} -> {*ie} || .#. [{*ʧ} | {*ʤ} | {*ʃ} | {*j}] _ [EnglishStarConsonant | EnglishPalatalConsonant | .#.],
    {*ē} -> {*īe} || .#. [{*ʧ} | {*ʤ} | {*ʃ} | {*j}] _ [EnglishStarConsonant | EnglishPalatalConsonant | .#.],
    {*é} -> {*íe} || .#. [{*ʧ} | {*ʤ} | {*ʃ} | {*j}] _ [EnglishStarConsonant | EnglishPalatalConsonant | .#.],
    {*ḗ} -> {*īe} || .#. [{*ʧ} | {*ʤ} | {*ʃ} | {*j}] _ [EnglishStarConsonant | EnglishPalatalConsonant | .#.]
];
\end{ReaderFacingFoma}

West Saxon \emph{gieldan} `pay', \emph{scield} `shield', and
\emph{scieppan} `create' show diphthongization after an already palatal
consonant (\citeproc{ref-RingeTaylor2014}{Ringe and Taylor 2014,
215--16}, §6.5.1). Their dialectal and phonological restriction
separates this development from system-wide i-umlaut.

Hogg's \emph{giefan} `give' and \emph{sceap} `sheep' belong to the same
palatal-consonant environment (\citeproc{ref-Hogg1992}{Hogg 1992,
108--9}). Fulk likewise assigns this diphthongization a place before
front mutation and distinguishes the two processes
(\citeproc{ref-Fulk2018}{Fulk 2018, 74}, §4.13).

The forms \emph{ġift} `gift' and \emph{sċēaþ} `sheath' fix the lower
boundary. If West Saxon palatal diphthongization precedes general
i-umlaut, PGmc \Recon{géftiz} `gift' yields \Pred{ġieft} rather than
expected \emph{ġift}, and PGmc \Recon{skáiθiz} `sheath' yields
\Pred{sċǣþ} rather than expected \emph{sċēaþ}. These witnesses place
West Saxon palatal diphthongization after general umlaut; no tested
lexical item supplies a later terminus ante quem.

The one-sided chronology reflects the difference in scale. General
umlaut reorganizes the vowel system, whereas West Saxon palatal
diphthongization affects a narrower dialectal class after palatal
consonants. Its exact later placement remains undemonstrated by the
present lexicon.

\newpage

\section{J-cluster coalescence}\label{j-cluster-coalescence}

\subsection{Historical discussion}\label{historical-discussion-29}

Only a small lexical group reveals the coalescence of velars with
\emph{*j}. Plain-velar and \emph{*sk} palatalization must already have
run before \emph{*gj} and \emph{*kj} acquire their later outcomes.
Campbell, Ringe and Taylor, and Fulk discuss the palatalized and fronted
outcomes in \emph{bīeġan} `bend' and \emph{sēċan} `seek' without
assigning this later cluster adjustment the status of a major sound law
(\citeproc{ref-Campbell1959}{Campbell 1959, 89, 107--8}, §§170,
248--251; \citeproc{ref-RingeTaylor2014}{Ringe and Taylor 2014,
213--51}, §§6.4.1, 6.5.1, 6.6.1--6.6.4; \citeproc{ref-Fulk2018}{Fulk
2018, 65, 75}, §§4.7, 4.13).

\subsection{\texorpdfstring{SC057. Coalescence of velar + \emph{*j}
clusters
(\texttt{OEJClusterCoalescence})}{SC057. Coalescence of velar +  clusters (OEJClusterCoalescence)}}\label{rule-OEJClusterCoalescence}

\begin{ReaderFacingFoma}
define OEJClusterCoalescence (
    [{*g} {*j} -> {*ʤ}]
    .o. [{*k} {*j} -> {*ʧ}]
);
\end{ReaderFacingFoma}

The forms \emph{bīeġan} `bend' and \emph{sēċan} `seek' determine the
earlier boundary. If coalescence precedes
\hyperref[rule-OEVelarPalatalization]{SC052 OEVelarPalatalization}, the
developments behind \emph{bīeġan} `bend' and \emph{sēċan} `seek' are
lost. Related forms such as \emph{fylġan} `follow', \emph{heċġ} `hedge',
and \emph{sengan} `singe' fail in the same broader palatalization zone.
PGmc \Recon{báugijaną} `bow' yields \Pred{bēaġan} rather than expected
OE \emph{bīeġan}, and PGmc \Recon{sōkijaną} `seek' yields \Pred{sōċan}
rather than expected \emph{sēċan}. This demonstrates that velar
palatalization preceded coalescence. Nothing in the present lexicon
supplies a terminus ante quem.

\newpage

\section{Nasal dissimilation}\label{nasal-dissimilation}

\subsection{Historical discussion}\label{historical-discussion-30}

Most accounts introduce nasal dissimilation to explain individual forms
rather than as a regular sound law. Luick records \emph{enetre}
`yearling' (spelled \emph{enitre} `yearling' in his text)
(\citeproc{ref-Luick1914}{Luick 1914-\/-1940, 166}); Campbell discusses
\emph{heofon} `heaven' with suffixal variation
(\citeproc{ref-Campbell1959}{Campbell 1959, 155}); and Hogg encounters
the same form while treating back mutation (\citeproc{ref-Hogg1992}{Hogg
1992, 112}).

Fulk supplies the clearest general formulation: ``In the cluster mn, the
first consonant tends to lose its nasality by dissimilation, though the
results are hardly regular'' (\citeproc{ref-Fulk2018}{Fulk 2018, 121},
§6.11). Ringe and Taylor stay close to the lexical evidence and note
that \emph{enetre} `yearling' reflects ``loss of the second \emph{*n} by
dissimilation'' (\citeproc{ref-RingeTaylor2014}{Ringe and Taylor 2014,
282}).

The disagreement concerns scope. Fulk's formulation recognizes a
recurrent but irregular development in \texttt{mn}; the remaining
discussions stay with particular lexical outcomes. None warrants a sound
law comparable in scope to the major Old English vowel changes.

\subsection{\texorpdfstring{\CAPRRuleHeading{SC058. Nasal dissimilation in short-vowel environments}{OENasalDissimilation}}{}}\label{rule-OENasalDissimilation}

\begin{ReaderFacingFoma}
define OENasalDissimilation [
    {*m} -> {*f} || EnglishStarShortVowel _ EnglishStarShortVowel {*n} [EnglishStarShortVowel | .#.]
];
\end{ReaderFacingFoma}

I adopt a narrower environment than the handbook observations might
suggest. Fulk formulates the tendency at the level of \texttt{mn}
clusters and illustrates it with \emph{heofon} `heaven' and
\emph{fæstenn} `fasting' (\citeproc{ref-Fulk2018}{Fulk 2018, 121},
§6.11). Ringe and Taylor show the same kind of development in
\emph{enetre} `yearling' (\citeproc{ref-RingeTaylor2014}{Ringe and
Taylor 2014, 282}). Campbell's ``\emph{heofon} is for older
\emph{hefzen}'' and Hogg's sequence \emph{*hefon > heofon} preserve
outcomes of the same kind (\citeproc{ref-Campbell1959}{Campbell 1959,
155}; \citeproc{ref-Hogg1992}{Hogg 1992, 112}). The short-vowel
environment adopted here covers a recurrent subset of these outcomes,
not every dissimilatory development involving nasals.

No witness word fixes the position of nasal dissimilation within the Old
English sequence. Reversing its order with any tested neighbor leaves
every checked output unchanged. A more precise relative chronology would
therefore require lexical evidence not represented here.

\newpage

\section{Back mutation}\label{back-mutation}

\subsection{Historical discussion}\label{historical-discussion-31}

West Saxon \emph{giefan} `give' and \emph{wefan} `weave' stand against
non-West-Saxon \emph{geofad} `gave' and \emph{weofan} `weave'. Ringe and
Taylor use this contrast to define the dialectal profile of back
mutation (\citeproc{ref-RingeTaylor2014}{Ringe and Taylor 2014, 319},
§6.9.4). Campbell's treatment of diphthongization before following back
vowels includes \emph{heofon} `heaven'
(\citeproc{ref-Campbell1959}{Campbell 1959, 86}, §207), while Hogg draws
the instructive comparison with breaking (\citeproc{ref-Hogg1992}{Hogg
1992, 112}). Fulk accordingly separates back mutation from the earlier
umlautal changes (\citeproc{ref-Fulk2018}{Fulk 2018, 69}, §4.8).

\subsection{\texorpdfstring{SC059. Back mutation before labials and
liquids
(\texttt{OEBackMutation})}{SC059. Back mutation before labials and liquids (OEBackMutation)}}\label{rule-OEBackMutation}

\begin{ReaderFacingFoma}
define OEBackMutation [
    {*e} -> {*eo} || _ [EnglishStarLabial | EnglishStarLiquid] {*u},
    {*æ} -> {*ea} || _ [EnglishStarLabial | EnglishStarLiquid] EnglishBackMutationTrigger,
    {*é} -> {*éo} || _ [EnglishStarLabial | EnglishStarLiquid] {*u}
];
\end{ReaderFacingFoma}

Three witness forms bracket the chronology. If back mutation precedes
\hyperref[rule-OESecondaryNasalization]{SC048 OESecondaryNasalization},
forms such as \Recon{gébaną} `give' produce \emph{ġeofan} `give'; the
expected form is \emph{ġiefan} `give'. \Recon{stélaną} `steal' likewise
produces \emph{steolan} `steal'; the expected form is \emph{stelan}
`steal'. At the other edge, delaying back mutation until after
\hyperref[rule-OEWeakTailReduction]{SC078 OEWeakTailReduction} makes
\Recon{wébaną} `weave' yield \emph{weofan} `weave'; the expected form is
\emph{wefan} `weave'. Thus back mutation follows secondary nasalization
but precedes the weak-tail reductions.

\newpage

\section{West Saxon palatal umlaut}\label{west-saxon-palatal-umlaut}

\subsection{Historical discussion}\label{historical-discussion-32}

The reflexes \emph{miht} `might' and \emph{niht} `night' place West
Saxon palatal umlaut after the principal umlautal developments. Campbell
and Ringe and Taylor describe the forms themselves; Fulk supplies the
broader chronology of palatal-vowel change
(\citeproc{ref-Campbell1959}{Campbell 1959, 107--8}, §§248--251;
\citeproc{ref-RingeTaylor2014}{Ringe and Taylor 2014, 215--51}, §§6.5.1,
6.6.1--6.6.4; \citeproc{ref-Fulk2018}{Fulk 2018, 65, 75}, §§4.7, 4.13).

\subsection{\texorpdfstring{\CAPRRuleHeading{SC060. West Saxon palatal umlaut before \emph{*h}-clusters}{OEWsPalatalUmlaut}}{}}\label{rule-OEWsPalatalUmlaut}

\begin{ReaderFacingFoma}
define OEWsPalatalUmlaut [
    {*eo} -> {*i} || _ OEHCluster .#.,
    {*io} -> {*i} || _ OEHCluster .#.,
    {*ie} -> {*i} || _ OEHCluster .#.,
    {*eo} -> {*i} || _ OEHCluster EnglishStarFrontVowel,
    {*io} -> {*i} || _ OEHCluster EnglishStarFrontVowel,
    {*ie} -> {*i} || _ OEHCluster EnglishStarFrontVowel,
    {*éo} -> {*i} || _ OEHCluster .#.,
    {*ío} -> {*i} || _ OEHCluster .#.,
    {*íe} -> {*i} || _ OEHCluster .#.,
    {*éo} -> {*i} || _ OEHCluster EnglishStarFrontVowel,
    {*ío} -> {*i} || _ OEHCluster EnglishStarFrontVowel,
    {*íe} -> {*i} || _ OEHCluster EnglishStarFrontVowel
];
\end{ReaderFacingFoma}

The change to \emph{*i} before \emph{*h}-clusters can be ordered only on
its earlier side. If palatal umlaut precedes
\hyperref[rule-OEIUmlaut]{SC055 OEIUmlaut}, the forms behind \emph{miht}
`might' and \emph{niht} `night' remain at the overdeveloped stage
\Pred{mieht} and \Pred{nieht} rather than expected OE \emph{miht} and
\emph{niht}. Consequently, i-umlaut precedes palatal umlaut. Reordering
the latter against any tested later change leaves both witness forms
unchanged.

\newpage

\section{Weak-tail nasal loss}\label{weak-tail-nasal-loss}

\subsection{Historical discussion}\label{historical-discussion-33}

The pathway from \Recon{dōną} `do' to \emph{dōn} `do' supplies the sole
lexical thread through this reduction. Campbell, Hogg, and Fulk place
such weak-tail losses among apocope and related late reductions
(\citeproc{ref-Campbell1959}{Campbell 1959, 144--45}, §§345--349;
\citeproc{ref-Hogg1992}{Hogg 1992, 120--21};
\citeproc{ref-Fulk2018}{Fulk 2018, 91}, §5.6). The witness, however,
ties the change to a much older development. Its immediate neighbors
remain untested.

\subsection{\texorpdfstring{\CAPRRuleHeading{SC061. Reduction of final nasal weak-tail endings}{OEWeakTailNasalLoss}}{}}\label{rule-OEWeakTailNasalLoss}

\begin{ReaderFacingFoma}
define OEWeakTailNasalLoss [
    {*n} {*ą} -> {*n} || _ .#.,
    {*m} {*ą} -> {*m} || _ .#.
];
\end{ReaderFacingFoma}

Final weak-tail \emph{*-ną} and \emph{*-mą} accordingly yield plain
\emph{*-n} and \emph{*-m}.

Only \emph{dōn} `do' constrains the relative order. Placing this loss
before the older n-stem loss makes the derivation record no output
instead of expected OE \emph{dōn} `do'. The older loss must therefore
precede weak-tail nasal loss. Nothing in the current lexicon
distinguishes among its possible later positions, and one witness cannot
establish a wider historical development.

\newpage

\section{High-vowel apocope}\label{high-vowel-apocope}

\subsection{Historical discussion}\label{historical-discussion-34}

Final high vowels must survive long enough to condition umlaut before
apocope removes them after heavy syllables and in the relevant
trisyllabic patterns. Campbell, Hogg, Ringe and Taylor, and Fulk agree
on this Old English development, though they differ over the extent of
the surrounding syncope (\citeproc{ref-Campbell1959}{Campbell 1959,
144--45}, §§345--349; \citeproc{ref-Hogg1992}{Hogg 1992, 120};
\citeproc{ref-RingeTaylor2014}{Ringe and Taylor 2014, 284--303},
§§6.8.1, 6.8.4; \citeproc{ref-Fulk2018}{Fulk 2018, 91}, §5.6).

\subsection{\texorpdfstring{\CAPRRuleHeading{SC063. High-vowel apocope after heavy syllables and in trisyllables}{OEHighVowelApocope}}{}}\label{rule-OEHighVowelApocope}

\begin{ReaderFacingFoma}
define OEHighVowelApocope [
    {*i} -> 0 || EnglishStarLongVowel OEAnyConsonant+ _ .#.,
    {*u} -> 0 || EnglishStarLongVowel OEAnyConsonant+ _ .#.,
    {*ų} -> 0 || EnglishStarLongVowel OEAnyConsonant+ _ .#.,
    {*i} -> 0 || EnglishStarLongDiphthong OEAnyConsonant+ _ .#.,
    {*u} -> 0 || EnglishStarLongDiphthong OEAnyConsonant+ _ .#.,
    {*ų} -> 0 || EnglishStarLongDiphthong OEAnyConsonant+ _ .#.,
    {*i} -> 0 || EnglishStarShortDiphthong OEAnyConsonant OEAnyConsonant+ _ .#.,
    {*u} -> 0 || EnglishStarShortDiphthong OEAnyConsonant OEAnyConsonant+ _ .#.,
    {*ų} -> 0 || EnglishStarShortDiphthong OEAnyConsonant OEAnyConsonant+ _ .#.,
    {*i} -> 0 || EnglishStarShortVowel OEAnyConsonant OEAnyConsonant+ _ .#.,
    {*u} -> 0 || EnglishStarShortVowel OEAnyConsonant OEAnyConsonant+ _ .#.,
    {*ų} -> 0 || EnglishStarShortVowel OEAnyConsonant OEAnyConsonant+ _ .#.,
    {*i} -> 0 || EnglishStarLongVowel OEAnyConsonant+ EnglishStarShortVowel OEAnyConsonant+ _ .#.,
    {*u} -> 0 || EnglishStarLongVowel OEAnyConsonant+ EnglishStarShortVowel OEAnyConsonant+ _ .#.,
    {*ų} -> 0 || EnglishStarLongVowel OEAnyConsonant+ EnglishStarShortVowel OEAnyConsonant+ _ .#.,
    {*i} -> 0 || EnglishStarLongDiphthong OEAnyConsonant+ EnglishStarShortVowel OEAnyConsonant+ _ .#.,
    {*u} -> 0 || EnglishStarLongDiphthong OEAnyConsonant+ EnglishStarShortVowel OEAnyConsonant+ _ .#.,
    {*ų} -> 0 || EnglishStarLongDiphthong OEAnyConsonant+ EnglishStarShortVowel OEAnyConsonant+ _ .#.,
    {*i} -> 0 || EnglishStarShortDiphthong OEAnyConsonant OEAnyConsonant+ EnglishStarShortVowel OEAnyConsonant+ _ .#.,
    {*u} -> 0 || EnglishStarShortDiphthong OEAnyConsonant OEAnyConsonant+ EnglishStarShortVowel OEAnyConsonant+ _ .#.,
    {*ų} -> 0 || EnglishStarShortDiphthong OEAnyConsonant OEAnyConsonant+ EnglishStarShortVowel OEAnyConsonant+ _ .#.,
    {*i} -> 0 || EnglishStarShortDiphthong OEAnyConsonant EnglishStarShortVowel OEAnyConsonant+ _ .#.,
    {*u} -> 0 || EnglishStarShortDiphthong OEAnyConsonant EnglishStarShortVowel OEAnyConsonant+ _ .#.,
    {*ų} -> 0 || EnglishStarShortDiphthong OEAnyConsonant EnglishStarShortVowel OEAnyConsonant+ _ .#.,
    {*i} -> 0 || EnglishStarShortVowel OEAnyConsonant EnglishStarShortVowel OEAnyConsonant+ _ .#.,
    {*u} -> 0 || EnglishStarShortVowel OEAnyConsonant EnglishStarShortVowel OEAnyConsonant+ _ .#.,
    {*ų} -> 0 || EnglishStarShortVowel OEAnyConsonant EnglishStarShortVowel OEAnyConsonant+ _ .#.,
    {*i} -> 0 || EnglishStarShortVowel OEAnyConsonant OEAnyConsonant+ EnglishStarShortVowel OEAnyConsonant+ _ .#.,
    {*u} -> 0 || EnglishStarShortVowel OEAnyConsonant OEAnyConsonant+ EnglishStarShortVowel OEAnyConsonant+ _ .#.,
    {*ų} -> 0 || EnglishStarShortVowel OEAnyConsonant OEAnyConsonant+ EnglishStarShortVowel OEAnyConsonant+ _ .#.,
    {*i} -> 0 || EnglishStarLongVowel _ .#.,
    {*u} -> 0 || {*x} _ .#.,
    {*ų} -> 0 || {*x} _ .#.,
    {*i} -> 0 || {*x} _ .#.
];
\end{ReaderFacingFoma}

Final \emph{*i}, \emph{*u}, and \emph{*ų} cannot disappear before
completing their umlautal work. Applied before
\hyperref[rule-OEIUmlaut]{SC055 OEIUmlaut}, PGmc \Recon{kūi} `cow'
yields \Pred{cū} rather than expected OE \emph{cȳ} `cow', and PGmc
\Recon{brūdiz} `bride' yields \Pred{brūd} rather than expected OE
\emph{brȳd} `bride'. Conversely, if apocope waits until after
\hyperref[rule-OEUnstressedLongVowelShortening]{SC072
OEUnstressedLongVowelShortening}, PGmc \Recon{fúrxtīnaz} `fright' yields
\Pred{fyrht} rather than expected OE \emph{fyrhte} `fright'. The three
witnesses establish the sequence i-umlaut, high-vowel apocope,
unstressed long-vowel shortening.

\newpage

\section{\texorpdfstring{Post-apocope \emph{*n}-loss and medial
syncope}{Post-apocope -loss and medial syncope}}\label{post-apocope--loss-and-medial-syncope}

\subsection{\texorpdfstring{Historical discussion of post-apocope
\emph{*n}-loss and medial
syncope}{Historical discussion of post-apocope -loss and medial syncope}}\label{historical-discussion-of-post-apocope--loss-and-medial-syncope}

Evidence for post-apocope reduction is strikingly uneven. The inherited
\emph{*furht-} family makes the survival of one nasal diagnostic and
fixes both sides of stem-final n-loss
(\citeproc{ref-Kroonen2013}{Kroonen 2013, 201}). No comparable witness
orders the medial syncope that follows. Hogg, Ringe and Taylor, and Fulk
describe both processes within the late history of weak syllables
(\citeproc{ref-Hogg1992}{Hogg 1992, 120--21};
\citeproc{ref-RingeTaylor2014}{Ringe and Taylor 2014, 264--303},
§§6.7.3--6.8.4; \citeproc{ref-Fulk2018}{Fulk 2018, 91}, §5.6).

\subsection{\texorpdfstring{SC064. Loss of stem-final \emph{*n} after
long \emph{*ī}
(\texttt{NWGmcInStemNLoss})}{SC064. Loss of stem-final  after long  (NWGmcInStemNLoss)}}\label{rule-NWGmcInStemNLoss}

\begin{ReaderFacingFoma}
define NWGmcInStemNLoss [{*n} -> 0 || {*ī} _ .#.];
\end{ReaderFacingFoma}

Only final \emph{*n} after long \emph{*ī} is at issue, as in the
inherited family behind \emph{fyrhte} `fright'.

The same proto-form fixes both edges. Before
\hyperref[rule-PWGmcFinalBareALoss]{SC041 PWGmcFinalBareALoss}, PGmc
\Recon{fúrxtīnaz} `fright' yields \Pred{fyrhten} rather than expected OE
\emph{fyrhte} `fright'. After
\hyperref[rule-OEUnstressedLongVowelShortening]{SC072
OEUnstressedLongVowelShortening}, PGmc \Recon{fúrxtīnaz} again yields
\Pred{fyrhten} rather than expected \emph{fyrhte} `fright'. I therefore
order final bare-a loss, stem-final n-loss, and unstressed long-vowel
shortening in that sequence. Both boundaries are firm within the
derivation, but depend upon one lexical family.

\subsection{\texorpdfstring{\CAPRRuleHeading{SC065. Medial syncope before dentals after heavy syllables}{OEMedialSyncope}}{}}\label{rule-OEMedialSyncope}

Loss of medial \emph{*i} before dentals belongs to the late weak-tail
history described by Hogg, Ringe and Taylor, and Fulk
(\citeproc{ref-Hogg1992}{Hogg 1992, 120--21};
\citeproc{ref-RingeTaylor2014}{Ringe and Taylor 2014, 264--303},
§§6.7.3--6.8.4; \citeproc{ref-Fulk2018}{Fulk 2018, 91}, §5.6).

\begin{ReaderFacingFoma}
define OEMedialSyncope [
    {*i} -> 0 || EnglishStarLongVowel OEAnyConsonant+ _ [{*θ}|{*ð}|{*d}|{*t}],
    {*i} -> 0 || EnglishStarLongDiphthong OEAnyConsonant+ _ [{*θ}|{*ð}|{*d}|{*t}],
    {*i} -> 0 || EnglishStarShortVowel OEAnyConsonant OEAnyConsonant+ _ [{*θ}|{*ð}|{*d}|{*t}]
];
\end{ReaderFacingFoma}

No diagnostic word establishes a local chronology. Moving medial syncope
to either end of the tested range leaves every checked output unchanged.
Its handbook placement after apocope and before later cluster
simplification therefore remains preferable, but the present lexicon
cannot demonstrate it.

\newpage

\section{Late syncope and
degemination}\label{late-syncope-and-degemination}

\subsection{Historical discussion of late syncope and
degemination}\label{historical-discussion-of-late-syncope-and-degemination}

Vowel loss creates the clusters upon which later assimilation and
degemination operate. Hogg and Ringe and Taylor describe this
dependence, while Brunner's \emph{netle} `nettle' beside later
\emph{netele} `nettle' supplies a concrete lexical type
(\citeproc{ref-Hogg1992}{Hogg 1992, 120--21};
\citeproc{ref-RingeTaylor2014}{Ringe and Taylor 2014, 264--96},
§§6.7.3--6.8.2; \citeproc{ref-SieversBrunner1965}{Brunner 1965,
144--45}, §§158--159). Fulk places this syncope after i-umlaut
(\citeproc{ref-Fulk2018}{Fulk 2018, 91}, §5.6).

The three relations are not equally secure. Lexical evidence orders
syncope and degemination; the intervening dental assimilation has no
independent ordering witness.

\subsection{\texorpdfstring{\CAPRRuleHeading{SC066. L-adjacent syncope in medial syllables}{OELAdjacentSyncope}}{}}\label{rule-OELAdjacentSyncope}

\begin{ReaderFacingFoma}
define OELAdjacentSyncope [
    {*i} -> 0 || EnglishStarShortVowel OEAnyConsonant+ _ {*l},
    {*i} -> 0 || EnglishStarLongVowel OEAnyConsonant+ _ {*l},
    {*i} -> 0 || EnglishStarDiphthong OEAnyConsonant+ _ {*l}
];
\end{ReaderFacingFoma}

The loss of medial \emph{*i} before \emph{*l} is late enough to preserve
earlier umlaut, as \emph{netle} `nettle' and \emph{spinl} `spindle'
demonstrate.

Placed before i-umlaut, PGmc \Recon{nátilōn} `nettle' yields
\Pred{nætle} rather than expected OE \emph{netle} `nettle', and PGmc
\Recon{spénnilō} `spindle' yields \Pred{spenl} rather than expected
\emph{spinl} `spindle'. Placed after preconsonantal degemination, PGmc
\Recon{spénnilō} yields \Pred{spinnl} rather than expected \emph{spinl}.
The witnesses therefore establish the sequence i-umlaut, l-adjacent
syncope, preconsonantal degemination. The first relation separates two
historical phases; the second is a direct feeding relation, since
syncope creates the cluster that degemination simplifies.

\subsection{\texorpdfstring{\CAPRRuleHeading{SC067. Dental assimilation in newly formed clusters}{OEDentalAssimilation}}{}}\label{rule-OEDentalAssimilation}

\begin{ReaderFacingFoma}
define OEDentalAssimilation [
    {*θ} -> 0 || {*t} _
];
\end{ReaderFacingFoma}

Loss of \emph{*θ} after \emph{*t} resolves a dental cluster produced by
syncope (\citeproc{ref-Hogg1992}{Hogg 1992, 120--21};
\citeproc{ref-RingeTaylor2014}{Ringe and Taylor 2014, 279--96}, §§6.7.5,
6.8.2). No witness distinguishes its position: moving dental
assimilation across every tested neighbor leaves the outputs unchanged.
I nevertheless place it after syncope, which supplies its input, and
before the more general cluster simplification described in the
handbooks. This order is phonologically motivated, not established by a
lexical contrast.

\subsection{\texorpdfstring{\CAPRRuleHeading{SC068. Preconsonantal degemination before sonorants}{OEPreconsonantalDegemination}}{}}\label{rule-OEPreconsonantalDegemination}

\begin{ReaderFacingFoma}
define OEPreconsonantalDegemination OEPreconsonantalDegemTT .o. OEPreconsonantalDegemNN;
\end{ReaderFacingFoma}

Preconsonantal \emph{*tt} and \emph{*nn} simplify only after syncope has
created a following sonorant cluster, as in \emph{spinl} `spindle'
(\citeproc{ref-RingeTaylor2014}{Ringe and Taylor 2014, 279--96},
§§6.7.5, 6.8.2).

Placed before l-adjacent syncope, PGmc \Recon{spénnilō} `spindle' yields
\Pred{spinnl} rather than expected OE \emph{spinl} `spindle'. Syncope
must therefore create the cluster before degemination simplifies it.
Reordering degemination against any tested later change leaves the
witness unchanged, so no terminus ante quem is known.

\newpage

\section{Early o-shortening}\label{early-o-shortening}

\subsection{Historical discussion}\label{historical-discussion-35}

After the principal palatal and umlautal changes, unstressed vowels
undergo shortening, fronting, merger, and sometimes complete loss.
Campbell describes the early shortening of unaccented long vowels, while
Hogg, Ringe and Taylor, and Fulk relate it to apocope, syncope, and the
later reductions (\citeproc{ref-Campbell1959}{Campbell 1959, 148}, §355;
\citeproc{ref-Hogg1992}{Hogg 1992, 120--21};
\citeproc{ref-RingeTaylor2014}{Ringe and Taylor 2014, 298--314},
§§6.8.3--6.9.3; \citeproc{ref-Fulk2018}{Fulk 2018, 90--96}, §§5.6--5.7).

Early o-shortening has only a distant earlier boundary. The rules that
follow, especially \hyperref[rule-OEUnstressedFrontingEarly]{SC070
OEUnstressedFrontingEarly} and
\hyperref[rule-OEUnstressedLongVowelShortening]{SC072
OEUnstressedLongVowelShortening}, have more closely defined relations.

\subsection{\texorpdfstring{\CAPRRuleHeading{SC069. Early shortening of unstressed \emph{*ō} before nasals}{OEEarlyOShortening}}{}}\label{rule-OEEarlyOShortening}

\begin{ReaderFacingFoma}
define OEEarlyOShortening [
    {*ō} -> {*a} || EnglishStarVocalic [EnglishStarConsonant | EnglishPalatalConsonant]+ _ EnglishStarNasal
];
\end{ReaderFacingFoma}

The rule shortens unstressed long \emph{*ō} before a following nasal.
Because this shortening happens early, the resulting \emph{*a} can still
participate in the later fronting and merger that shape many weak final
syllables.

Moving the rule before \hyperref[rule-NWGmcNStemNLoss]{SC023
NWGmcNStemNLoss}, PGmc \Recon{nḗdrōn} `adder' yields \Pred{nǣdran}
rather than expected OE \emph{nǣdre} `adder', PGmc \Recon{érθōn} `earth'
yields \Pred{eorþan} rather than expected \emph{eorþe} `earth', and PGmc
\Recon{fláskōn} `flask' yields \Pred{flascan} rather than expected
\emph{flasce} `flask'. The same earlier shift also disrupts forms such
as \emph{heorte} `heart' and \emph{līne} `line'. This broad set of
failures requires \hyperref[rule-OEEarlyOShortening]{SC069
OEEarlyOShortening} to follow \hyperref[rule-NWGmcNStemNLoss]{SC023
NWGmcNStemNLoss}.

If the rule is moved later within the tested sequence, no checked form
yields a form different from the expected one. The checked forms
therefore do not identify a corresponding later constraint. The sources
place early \emph{*ō}-shortening before the later weak-tail changes
without fixing a closer local order.

\newpage

\section{Early unstressed fronting and later
o-shortening}\label{early-unstressed-fronting-and-later-o-shortening}

\subsection{Historical discussion of early unstressed fronting and later
o-shortening}\label{historical-discussion-of-early-unstressed-fronting-and-later-o-shortening}

Campbell distinguishes the shortening of unaccented long vowels, while
Hogg, Ringe and Taylor, and Fulk place fronting and shortening within a
later history of syncope and final-vowel adjustment
(\citeproc{ref-Campbell1959}{Campbell 1959, 148}, §355;
\citeproc{ref-Hogg1992}{Hogg 1992, 120--21};
\citeproc{ref-RingeTaylor2014}{Ringe and Taylor 2014, 298--314},
§§6.8.3--6.9.3; \citeproc{ref-Fulk2018}{Fulk 2018, 90--96}, §§5.6--5.7).
Earlier unstressed fronting precedes later o-shortening.

\hyperref[rule-OEUnstressedFrontingEarly]{SC070
OEUnstressedFrontingEarly} has both an earlier and a later lexical
breakpoint. \hyperref[rule-OELateOShortening]{SC071 OELateOShortening}
confirms their reciprocal order, but no checked form fixes its later
boundary.

\subsection{\texorpdfstring{\CAPRRuleHeading{SC070. Early fronting of unstressed \emph{*a}}{OEUnstressedFrontingEarly}}{}}\label{rule-OEUnstressedFrontingEarly}

\begin{ReaderFacingFoma}
define OEUnstressedFrontingEarly OEUnstressedAFronting;
\end{ReaderFacingFoma}

The rule fronts unstressed \emph{*a} to \emph{*æ} after the earlier
shortening has created a frontable vowel but before the later shortening
of unstressed \emph{*ō}. It produces endings such as OE \emph{-en} in
\emph{lungen} `lungs'.

If the rule is moved before \hyperref[rule-OEVelarPalatalization]{SC052
OEVelarPalatalization}, PGmc \Recon{lúnganjō} `lungs' yields
\Pred{lunġen} rather than expected OE \emph{lungen} `lungs'. If the rule
is delayed until after \hyperref[rule-OELateOShortening]{SC071
OELateOShortening}, PGmc \Recon{búrōθi} `bears' yields \Pred{boreþ}
rather than expected OE \emph{boraþ} `bears', and PGmc \Recon{mḗnōθz}
`month' yields \Pred{mōneþ} rather than expected \emph{mōnaþ} `month'.
The witness forms require
\hyperref[rule-OEUnstressedFrontingEarly]{SC070
OEUnstressedFrontingEarly} to follow
\hyperref[rule-OEVelarPalatalization]{SC052 OEVelarPalatalization} and
precede \hyperref[rule-OELateOShortening]{SC071 OELateOShortening}.

The relation to \hyperref[rule-OELateOShortening]{SC071
OELateOShortening} is local. The earlier boundary at
\hyperref[rule-OEVelarPalatalization]{SC052 OEVelarPalatalization}
places fronting after the older palatal developments.

\subsection{\texorpdfstring{SC071. Later shortening of unstressed
\emph{*ō}
(\texttt{OELateOShortening})}{SC071. Later shortening of unstressed  (OELateOShortening)}}\label{rule-OELateOShortening}

The following rule handles the later shortening stage.

\begin{ReaderFacingFoma}
define OELateOShortening [
    {*ō} -> {*a} || EnglishStarVocalic [EnglishStarConsonant | EnglishPalatalConsonant]+ _ [EnglishStarConsonant | EnglishPalatalConsonant]*
];
\end{ReaderFacingFoma}

The rule shortens the remaining unstressed long \emph{*ō} after
fronting, producing the later ``stable a'' endings in OE \emph{boraþ}
`bears' and \emph{liornaþ} `learns'.

Moving the rule before \hyperref[rule-OEUnstressedFrontingEarly]{SC070
OEUnstressedFrontingEarly} makes PGmc \Recon{búrōθi} `bears' yield
\Pred{boreþ} rather than expected OE \emph{boraþ} `bears', and PGmc
\Recon{líznōθi} `learns' yield \Pred{liorneþ} rather than expected
\emph{liornaþ} `learns'. The contrast requires
\hyperref[rule-OELateOShortening]{SC071 OELateOShortening} to follow
\hyperref[rule-OEUnstressedFrontingEarly]{SC070
OEUnstressedFrontingEarly}. Moving it later within the tested range
creates no equally sharp failure.

\newpage

\section{Unstressed long-vowel shortening and
ae-merger}\label{unstressed-long-vowel-shortening-and-ae-merger}

\subsection{Historical discussion of unstressed long-vowel shortening
and
ae-merger}\label{historical-discussion-of-unstressed-long-vowel-shortening-and-ae-merger}

Campbell describes the shortening of unaccented long vowels, and Ringe
and Taylor place it among the last prehistoric Old English changes
before the merger of unstressed \emph{*æ} with \emph{*e}
(\citeproc{ref-Campbell1959}{Campbell 1959, 148}, §355;
\citeproc{ref-Hogg1992}{Hogg 1992, 120--21};
\citeproc{ref-RingeTaylor2014}{Ringe and Taylor 2014, 298--314},
§§6.8.3--6.9.3; \citeproc{ref-Fulk2018}{Fulk 2018, 90--96}, §§5.6--5.7).

\hyperref[rule-OEUnstressedLongVowelShortening]{SC072
OEUnstressedLongVowelShortening} and
\hyperref[rule-OEUnstressedAEMerger]{SC073 OEUnstressedAEMerger} have a
reciprocal ordering relation. \hyperref[rule-NWGmcInStemNLoss]{SC064
NWGmcInStemNLoss} supplies the earlier boundary of shortening, and
\hyperref[rule-OEHLoss]{SC085 OEHLoss} the later boundary of the merger.

\subsection{\texorpdfstring{\CAPRRuleHeading{SC072. Shortening of unstressed long vowels}{OEUnstressedLongVowelShortening}}{}}\label{rule-OEUnstressedLongVowelShortening}

\begin{ReaderFacingFoma}
define OEUnstressedLongVowelShortening OEUnstressedLongVowelShortening1
    .o. OEUnstressedLongVowelShortening2
    .o. OEUnstressedLongVowelShortening3
    .o. OEUnstressedLongVowelShortening5
    .o. OEUnstressedLongVowelShortening6
    .o. OEUnstressedLongVowelShortening7
    .o. OEUnstressedLongVowelShortening8;
\end{ReaderFacingFoma}

The rule shortens the remaining unstressed long vowels before weak final
syllables reach their later forms. A small group of lexical witnesses
fixes its chronology.

If the rule is moved before \hyperref[rule-NWGmcInStemNLoss]{SC064
NWGmcInStemNLoss}, PGmc \Recon{fúrxtīnaz} `fright' yields \Pred{fyrhten}
rather than expected OE \emph{fyrhte} `fright'. If the rule is delayed
until after \hyperref[rule-OEUnstressedAEMerger]{SC073
OEUnstressedAEMerger}, PGmc \Recon{nḗdrōn} `adder' yields \Pred{nǣdræ}
rather than expected OE \emph{nǣdre} `adder', and PGmc \Recon{fádēr}
`father' yields \Pred{fædær} rather than expected \emph{fæder} `father'.
These outputs require
\hyperref[rule-OEUnstressedLongVowelShortening]{SC072
OEUnstressedLongVowelShortening} to follow
\hyperref[rule-NWGmcInStemNLoss]{SC064 NWGmcInStemNLoss} and precede
\hyperref[rule-OEUnstressedAEMerger]{SC073 OEUnstressedAEMerger}.

Shortening therefore follows the earlier weak-tail preparation and
immediately precedes the merger.

\subsection{\texorpdfstring{SC073. Merger of unstressed \emph{*æ} with
\emph{*e}
(\texttt{OEUnstressedAEMerger})}{SC073. Merger of unstressed  with  (OEUnstressedAEMerger)}}\label{rule-OEUnstressedAEMerger}

The following rule handles the merger stage.

\begin{ReaderFacingFoma}
define OEUnstressedAEMerger OEWeakTailReduction3;
\end{ReaderFacingFoma}

The rule merges unstressed \emph{*æ} with \emph{*e} after shortening has
produced the weak final vowels, yielding the ordinary OE \emph{-e}
spellings.

Its earlier and later relations are both concrete. If the rule is moved
before \hyperref[rule-OEUnstressedLongVowelShortening]{SC072
OEUnstressedLongVowelShortening}, PGmc \Recon{nḗdrōn} `adder' yields
\Pred{nǣdræ} rather than expected OE \emph{nǣdre} `adder', and PGmc
\Recon{fádēr} `father' yields \Pred{fædær} rather than expected
\emph{fæder} `father'. If the rule is delayed until after
\hyperref[rule-OEHLoss]{SC085 OEHLoss}, PGmc \Recon{táixōn} `toe' yields
\Pred{tāæ} rather than expected OE \emph{tā} `toe'. These failures show
that \hyperref[rule-OEUnstressedLongVowelShortening]{SC072
OEUnstressedLongVowelShortening} must come before
\hyperref[rule-OEUnstressedAEMerger]{SC073 OEUnstressedAEMerger}, and
that \hyperref[rule-OEUnstressedAEMerger]{SC073 OEUnstressedAEMerger}
must come before \hyperref[rule-OEHLoss]{SC085 OEHLoss}.

The checked forms fix the local order after
\hyperref[rule-OEUnstressedLongVowelShortening]{SC072
OEUnstressedLongVowelShortening} and place the merger before the later
h-loss and contraction.

\newpage

\section{Medial unstressed-i
lowering}\label{medial-unstressed-i-lowering}

\subsection{\texorpdfstring{Historical discussion of medial unstressed-i
lowering and \emph{*ng}
retention}{Historical discussion of medial unstressed-i lowering and  retention}}\label{historical-discussion-of-medial-unstressed-i-lowering-and-retention}

Hogg and Ringe and Taylor treat the late weakening and merger of
unstressed vowels as a continuing history (\citeproc{ref-Hogg1992}{Hogg
1992, 120--21}; \citeproc{ref-RingeTaylor2014}{Ringe and Taylor 2014,
327--32}, §§6.9.5--6.9.6).
\hyperref[rule-OEMedUnstressedILowering1]{SC074
OEMedUnstressedILowering1} lowers medial unstressed \emph{i};
\hyperref[rule-OEMedUnstressedILowering]{SC075 OEMedUnstressedILowering}
preserves \emph{i} before \emph{*ng} in words of the \emph{sċilling}
`shilling' type.

General lowering precedes the restricted restoration before \emph{*ng}.
The evidence is narrower than that for
\hyperref[rule-OEUnstressedLongVowelShortening]{SC072
OEUnstressedLongVowelShortening} and
\hyperref[rule-OEUnstressedAEMerger]{SC073 OEUnstressedAEMerger}.

\subsection{\texorpdfstring{\CAPRRuleHeading{SC074. First medial unstressed-\emph{i} lowering}{OEMedUnstressedILowering1}}{}}\label{rule-OEMedUnstressedILowering1}

\begin{ReaderFacingFoma}
define OEMedUnstressedILowering1 [
    {*i} -> {*e} || EnglishStarVocalic [EnglishStarConsonant | EnglishPalatalConsonant]+ _
];
\end{ReaderFacingFoma}

The rule lowers medial unstressed \emph{*i} to \emph{*e} after a
preceding vocalic syllable. The resulting \emph{e}-outcome is reversed
before \emph{*ng}.

If the rule is moved before
\hyperref[rule-OEUnstressedLongVowelShortening]{SC072
OEUnstressedLongVowelShortening}, PGmc \Recon{fúrxtīnaz} `fright' yields
\Pred{fyrhti} rather than expected OE \emph{fyrhte} `fright'. If it is
delayed until after \hyperref[rule-OEMedUnstressedILowering]{SC075
OEMedUnstressedILowering}, PGmc \Recon{skíllingaz} `shilling' yields
\Pred{sċilleng} rather than expected \emph{sċilling} `shilling'. The
derivations require \hyperref[rule-OEMedUnstressedILowering1]{SC074
OEMedUnstressedILowering1} to follow
\hyperref[rule-OEUnstressedLongVowelShortening]{SC072
OEUnstressedLongVowelShortening} and precede
\hyperref[rule-OEMedUnstressedILowering]{SC075
OEMedUnstressedILowering}.

The evidence is narrow on each side. The rule follows unstressed
long-vowel shortening and precedes the more specific \emph{*ng}
preservation.

\subsection{\texorpdfstring{\CAPRRuleHeading{SC075. Preservation of medial unstressed \emph{*i} before \emph{*ng}}{OEMedUnstressedILowering}}{}}\label{rule-OEMedUnstressedILowering}

The following rule reverses the lowering before \emph{*ng}.

\begin{ReaderFacingFoma}
define OEMedUnstressedILowering [
    {*e} -> {*i} || _ {*n} {*g}
];
\end{ReaderFacingFoma}

The rule restores \emph{*i} before \emph{*ng}, preventing the broader
lowering from producing the wrong medial vowel in forms such as
\emph{sċilling} `shilling'.

Moving the rule before \hyperref[rule-OEMedUnstressedILowering1]{SC074
OEMedUnstressedILowering1} makes PGmc \Recon{skíllingaz} `shilling'
yield \Pred{sċilleng} rather than expected OE \emph{sċilling}
`shilling'. On this evidence, I take
\hyperref[rule-OEMedUnstressedILowering]{SC075 OEMedUnstressedILowering}
to follow \hyperref[rule-OEMedUnstressedILowering1]{SC074
OEMedUnstressedILowering1}. Moving it later within the tested range
creates no equally sharp failure.

\newpage

\section{Prefix i-reduction}\label{prefix-i-reduction}

\subsection{Historical discussion}\label{historical-discussion-36}

Late weak-tail reduction affects unstressed prefixes as well as
inflectional endings and medial vowels. Fulk's discussion of prefix
vowels accounts for OE \emph{*be-} and \emph{*ne-}
(\citeproc{ref-Fulk2018}{Fulk 2018, 97}, §5.7). Hogg and Ringe and
Taylor place such weakening within the broader late history of
unstressed vowels (\citeproc{ref-Hogg1992}{Hogg 1992, 120--21};
\citeproc{ref-RingeTaylor2014}{Ringe and Taylor 2014, 298--332},
§§6.8.3--6.9.6).

The tested forms do not determine the rule's position relative to a
neighboring change.

\subsection{\texorpdfstring{\CAPRRuleHeading{SC076. Reduction of prefixal \emph{*i} in unstressed position}{OEPrefixIReduction}}{}}\label{rule-OEPrefixIReduction}

\begin{ReaderFacingFoma}
define OEPrefixIReduction [
    {*i} -> {*ĕ} || .#. [{*b} | {*n}] _ [EnglishStarConsonant | EnglishPalatalConsonant] EnglishStarVocalic
];
\end{ReaderFacingFoma}

The rule reduces unstressed prefixal \emph{*i} to a weaker vowel in the
\emph{bi-} and \emph{ni-} type prefixes before a consonant plus a
following vowel. The development accounts for later prefix spellings
such as OE \emph{*be-} and \emph{*ne-}.

If the rule is moved earlier or later within the tested sequence, no
checked form yields a form different from the expected one. The tested
forms therefore do not place \hyperref[rule-OEPrefixIReduction]{SC076
OEPrefixIReduction} before or after any specific neighboring change.

The handbooks attest late prefix-vowel weakening, but the precise
placement remains approximate. No lexical failure fixes it.

\newpage

\section{Weak-tail reduction}\label{weak-tail-reduction}

\subsection{Historical discussion}\label{historical-discussion-37}

Campbell, Hogg, Ringe and Taylor, and Fulk describe a late history in
which apocope, shortening, contraction, and further weak-tail reductions
reshape final syllables (\citeproc{ref-Campbell1959}{Campbell 1959,
148}, §355; \citeproc{ref-Hogg1992}{Hogg 1992, 120--21};
\citeproc{ref-RingeTaylor2014}{Ringe and Taylor 2014, 298--314},
§§6.8.3--6.9.3; \citeproc{ref-Fulk2018}{Fulk 2018, 90--91}, §5.6).
Lexical failures place the remaining weak-tail reduction after
unstressed fronting and before contraction.

\subsection{\texorpdfstring{\CAPRRuleHeading{SC078. Reduction of remaining weak-tail vowels}{OEWeakTailReduction}}{}}\label{rule-OEWeakTailReduction}

\begin{ReaderFacingFoma}
define OEWeakTailReduction OEWeakTailReduction1;
\end{ReaderFacingFoma}

The rule reduces the remaining weak-tail vowels, preventing a broad
class of \emph{-en} and extra-vowel outcomes.

I place the change after \hyperref[rule-OEUnstressedFrontingEarly]{SC070
OEUnstressedFrontingEarly} and before
\hyperref[rule-OEContraction]{SC086 OEContraction}. Moving it before
\hyperref[rule-OEUnstressedFrontingEarly]{SC070
OEUnstressedFrontingEarly}, PGmc \Recon{bákaną} `bake' yields
\Pred{bacen} rather than expected OE \emph{bacan} `bake', and PGmc
\Recon{bíndaną} `bind' yields \Pred{binden} rather than expected
\emph{bindan} `bind', alongside a much wider set of comparable
\emph{-en} failures. If the rule is delayed until after
\hyperref[rule-OEContraction]{SC086 OEContraction}, PGmc
\Recon{fléuxaną} `flee' yields \Pred{flēoan} rather than expected OE
\emph{flēon} `flee', and PGmc \Recon{sláxaną} `slay' yields
\Pred{sleaan} rather than expected \emph{slēan} `slay'.

The earlier boundary spans a wide interval and does not establish a
close neighboring relation. The later boundary is narrower:
\hyperref[rule-OEWeakTailReduction]{SC078 OEWeakTailReduction} precedes
\hyperref[rule-OEContraction]{SC086 OEContraction}.

\newpage

\section{Final-j loss and final geminate
simplification}\label{final-j-loss-and-final-geminate-simplification}

\subsection{Historical discussion of final-j loss and final geminate
simplification}\label{historical-discussion-of-final-j-loss-and-final-geminate-simplification}

After \hyperref[rule-OEJLossAfterHeavy]{SC079 OEJLossAfterHeavy} removes
\emph{*j} in heavy environments, forms such as \emph{lungen} `lungs'
acquire a final geminate.
\hyperref[rule-OEFinalGeminateSimplification]{SC080
OEFinalGeminateSimplification} then removes the second nasal.

\hyperref[rule-OEJLossAfterHeavy]{SC079 OEJLossAfterHeavy} has a broad
earlier boundary at \hyperref[rule-OEIUmlaut]{SC055 OEIUmlaut}.
\hyperref[rule-OEFinalGeminateSimplification]{SC080
OEFinalGeminateSimplification} is fixed only by the final \emph{nn}
outcome in the following derivation.

\subsection{\texorpdfstring{SC079. Loss of \emph{*j} after heavy
syllables
(\texttt{OEJLossAfterHeavy})}{SC079. Loss of  after heavy syllables (OEJLossAfterHeavy)}}\label{rule-OEJLossAfterHeavy}

\begin{ReaderFacingFoma}
define OEJLossAfterHeavy [
    {*j} -> 0 || (EnglishStarLongVowel | EnglishStarDiphthong) [EnglishStarConsonantNoR | EnglishPalatalConsonant] _,
    {*j} -> 0 || EnglishStarShortVowel [EnglishStarConsonant | EnglishPalatalConsonant] [EnglishStarConsonantNoR | EnglishPalatalConsonant] _
];
\end{ReaderFacingFoma}

The rule removes \emph{*j} after the relevant heavy-syllable
configurations, after the earlier umlaut-sensitive vocalism has
developed. The affected glide is \emph{*j}.

If the rule is moved before \hyperref[rule-OEIUmlaut]{SC055 OEIUmlaut},
PGmc \Recon{galáubijaną} `believe' yields \Pred{ġelēafan} rather than
expected OE \emph{ġelīefan} `believe', PGmc \Recon{báugijaną} `bow'
yields \Pred{bēaġan} rather than expected \emph{bīeġan} `bow', and PGmc
\Recon{fúlgijaną} `follow' yields \Pred{fulġan} rather than expected
\emph{fylġan} `follow'. If it is delayed until after
\hyperref[rule-OEFinalGeminateSimplification]{SC080
OEFinalGeminateSimplification}, PGmc \Recon{lúnganjō} `lungs' yields
\Pred{lungenn} rather than expected OE \emph{lungen} `lungs'. I
accordingly take \hyperref[rule-OEJLossAfterHeavy]{SC079
OEJLossAfterHeavy} to follow \hyperref[rule-OEIUmlaut]{SC055 OEIUmlaut}
and precede \hyperref[rule-OEFinalGeminateSimplification]{SC080
OEFinalGeminateSimplification}.

The earlier boundary is broad, but the relation to final geminate
simplification is local.

\subsection{\texorpdfstring{\CAPRRuleHeading{SC080. Simplification of final geminates}{OEFinalGeminateSimplification}}{}}\label{rule-OEFinalGeminateSimplification}

The following rule handles the final simplification directly.

\begin{ReaderFacingFoma}
define OEFinalGeminateSimplification [
    {*n} -> 0 || {*n} _ .#.
];
\end{ReaderFacingFoma}

The rule removes the extra final nasal in forms where the preceding
derivation has already created a final geminate.

Moving the rule before \hyperref[rule-OEJLossAfterHeavy]{SC079
OEJLossAfterHeavy} makes PGmc \Recon{lúnganjō} `lungs' yield
\Pred{lungenn} rather than expected OE \emph{lungen} `lungs'. These
failures require \hyperref[rule-OEFinalGeminateSimplification]{SC080
OEFinalGeminateSimplification} to follow
\hyperref[rule-OEJLossAfterHeavy]{SC079 OEJLossAfterHeavy}. Moving it
later within the tested range before \hyperref[rule-OERMetathesis]{SC087
OERMetathesis} creates no new failure.

\newpage

\section{J-strengthening, vocalization, and
ei-contraction}\label{j-strengthening-vocalization-and-ei-contraction}

\subsection{Historical discussion of j-strengthening, vocalization, and
ei-contraction}\label{historical-discussion-of-j-strengthening-vocalization-and-ei-contraction}

\hyperref[rule-OEJStrengtheningAfterFrontDiphthong]{SC081
OEJStrengtheningAfterFrontDiphthong} preserves a consonantal outcome
after front diphthongs.
\hyperref[rule-OEIntervocalicJVocalization]{SC082
OEIntervocalicJVocalization} then vocalizes the remaining intervocalic
\emph{*j}, and \hyperref[rule-OEUnstressedEIContraction]{SC083
OEUnstressedEIContraction} removes the resulting \emph{ei}-like sequence
in weak verbal endings.

The output of each rule conditions the next.
\hyperref[rule-OEIntervocalicJVocalization]{SC082
OEIntervocalicJVocalization} has local lexical evidence on both sides;
\hyperref[rule-OEJStrengtheningAfterFrontDiphthong]{SC081
OEJStrengtheningAfterFrontDiphthong} has a distant earlier boundary, and
\hyperref[rule-OEUnstressedEIContraction]{SC083
OEUnstressedEIContraction} has no tested later boundary.

\subsection{\texorpdfstring{\CAPRRuleHeading{SC081. Strengthening of \emph{*j} after front diphthongs}{OEJStrengtheningAfterFrontDiphthong}}{}}\label{rule-OEJStrengtheningAfterFrontDiphthong}

\begin{ReaderFacingFoma}
define OEJStrengtheningAfterFrontDiphthong [
    {*j} -> {*ʒ} || [{*ēa}|{*ḗa}|{*íe}|{*īe}|{*éa}] _ EnglishStarVocalic
];
\end{ReaderFacingFoma}

After the relevant front diphthongs, \emph{*j} first strengthened to a
consonantal outcome; otherwise it would have vocalized too early.

If the rule is moved before \hyperref[rule-OEIUmlaut]{SC055 OEIUmlaut},
PGmc \Recon{stráwjaną} `strew' yields \Pred{strēaġan} rather than
expected OE \emph{strīeġan} `strew'. If it is delayed until after
\hyperref[rule-OEIntervocalicJVocalization]{SC082
OEIntervocalicJVocalization}, the same PGmc form yields \Pred{strīeian}
rather than \emph{strīeġan}. The order test requires
\hyperref[rule-OEJStrengtheningAfterFrontDiphthong]{SC081
OEJStrengtheningAfterFrontDiphthong} to follow
\hyperref[rule-OEIUmlaut]{SC055 OEIUmlaut} and precede
\hyperref[rule-OEIntervocalicJVocalization]{SC082
OEIntervocalicJVocalization}.

The earlier constraint reaches back to \hyperref[rule-OEIUmlaut]{SC055
OEIUmlaut} and therefore defines a wide interval. The \emph{strīeġan}
`strew' derivation fixes the local relation to
\hyperref[rule-OEIntervocalicJVocalization]{SC082
OEIntervocalicJVocalization}.

\subsection{\texorpdfstring{\CAPRRuleHeading{SC082. Intervocalic vocalization of \emph{*j}}{OEIntervocalicJVocalization}}{}}\label{rule-OEIntervocalicJVocalization}

\begin{ReaderFacingFoma}
define OEIntervocalicJVocalization [
    {*j} -> {*i} || EnglishStarVocalic _ EnglishStarVocalic
];
\end{ReaderFacingFoma}

The rule vocalizes intervocalic \emph{*j} to \emph{*i}, creating the
\emph{ei}-like sequence later removed by
\hyperref[rule-OEUnstressedEIContraction]{SC083
OEUnstressedEIContraction} in many weak verb forms.

Moving the rule before
\hyperref[rule-OEJStrengtheningAfterFrontDiphthong]{SC081
OEJStrengtheningAfterFrontDiphthong} makes PGmc \Recon{stráwjaną}
`strew' yield \Pred{strīeian} rather than expected OE \emph{strīeġan}
`strew'. Delaying it until after
\hyperref[rule-OEUnstressedEIContraction]{SC083
OEUnstressedEIContraction} makes PGmc \Recon{búrōjaną} `bore' yield
\Pred{boreian} rather than expected OE \emph{borian} `bore', PGmc
\Recon{xándlōjaną} `handle' yield \Pred{handleian} rather than expected
\emph{handlian} `handle', and PGmc \Recon{mákōjaną} `make' yield
\Pred{maceian} rather than expected \emph{macian} `make'. The witness
forms require \hyperref[rule-OEIntervocalicJVocalization]{SC082
OEIntervocalicJVocalization} to follow
\hyperref[rule-OEJStrengtheningAfterFrontDiphthong]{SC081
OEJStrengtheningAfterFrontDiphthong} and precede
\hyperref[rule-OEUnstressedEIContraction]{SC083
OEUnstressedEIContraction}.

\hyperref[rule-OEIntervocalicJVocalization]{SC082
OEIntervocalicJVocalization} is therefore ordered between strengthening
and contraction.

\subsection{\texorpdfstring{SC083. Contraction of unstressed \emph{ei}
(\texttt{OEUnstressedEIContraction})}{SC083. Contraction of unstressed  (OEUnstressedEIContraction)}}\label{rule-OEUnstressedEIContraction}

The final rule removes the extra unstressed \emph{e} before \emph{i}.

\begin{ReaderFacingFoma}
define OEUnstressedEIContraction [
    {*e} -> 0 || EnglishStarVocalic [EnglishStarConsonant | EnglishPalatalConsonant]+ _ {*i}
];
\end{ReaderFacingFoma}

The rule contracts the unstressed \emph{ei}-like sequence that the
preceding vocalization would otherwise leave behind in forms such as
\emph{borian} `bore' and \emph{liccian} `lick'.

Moving the rule before \hyperref[rule-OEIntervocalicJVocalization]{SC082
OEIntervocalicJVocalization} makes PGmc \Recon{búrōjaną} `bore' yield
\Pred{boreian} rather than expected OE \emph{borian} `bore', PGmc
\Recon{líznōjaną} `learn' yield \Pred{liorneian} rather than expected
\emph{liornian} `learn', and PGmc \Recon{líkkōjaną} `lick' yield
\Pred{licceian} rather than expected \emph{liccian} `lick'. The contrast
requires \hyperref[rule-OEUnstressedEIContraction]{SC083
OEUnstressedEIContraction} to follow
\hyperref[rule-OEIntervocalicJVocalization]{SC082
OEIntervocalicJVocalization}. Moving it later within the tested range
before \hyperref[rule-OERMetathesis]{SC087 OERMetathesis} creates no new
failure.

\newpage

\section{H-loss and contraction}\label{h-loss-and-contraction}

\subsection{Historical discussion of h-loss and
contraction}\label{historical-discussion-of-h-loss-and-contraction}

When \hyperref[rule-OEHLoss]{SC085 OEHLoss} removes intervocalic
\emph{*h}, it creates hiatus. \hyperref[rule-OEContraction]{SC086
OEContraction} immediately resolves the resulting vowel sequence.

Ringe and Taylor describe this late sequence of \emph{h}-loss and
contraction (\citeproc{ref-RingeTaylor2014}{Ringe and Taylor 2014,
305--14}, §§6.9.1--6.9.3). Fulk places the contracted verbs in a broader
Germanic context (\citeproc{ref-Fulk2018}{Fulk 2018, 270}, §12.21), and
Luick describes the corresponding West Germanic contractions
(\citeproc{ref-Luick1914}{Luick 1914-\/-1940, 165}).

\subsection{\texorpdfstring{SC085. Loss of intervocalic \emph{*h}
(\texttt{OEHLoss})}{SC085. Loss of intervocalic  (OEHLoss)}}\label{rule-OEHLoss}

\begin{ReaderFacingFoma}
define OEHLoss [
    {*x} -> 0 || EnglishStarVocalic _ EnglishStarVocalic
];
\end{ReaderFacingFoma}

The rule removes intervocalic \emph{*h}, creating the hiatus that later
contraction must resolve.

If the rule is moved before \hyperref[rule-OEUnstressedAEMerger]{SC073
OEUnstressedAEMerger}, PGmc \Recon{táixōn} `toe' yields \Pred{tāæ}
rather than expected OE \emph{tā} `toe'. If it is delayed until after
\hyperref[rule-OEContraction]{SC086 OEContraction}, PGmc
\Recon{fléuxaną} `flee' yields \Pred{flēoan} rather than expected OE
\emph{flēon} `flee', PGmc \Recon{sláxaną} `slay' yields \Pred{sleaan}
rather than expected \emph{slēan} `slay', PGmc \Recon{téxun} `draw'
yields \Pred{teoon} rather than expected \emph{tēon} `draw', and PGmc
\Recon{táixōn} yields \Pred{tāe} rather than expected \emph{tā}. These
outputs require \hyperref[rule-OEHLoss]{SC085 OEHLoss} to follow
\hyperref[rule-OEUnstressedAEMerger]{SC073 OEUnstressedAEMerger} and
precede \hyperref[rule-OEContraction]{SC086 OEContraction}.

The earlier boundary rests on one witness; the four later witnesses
establish the immediate relation to contraction.

\subsection{\texorpdfstring{SC086. Contraction of the resulting hiatus
(\texttt{OEContraction})}{SC086. Contraction of the resulting hiatus (OEContraction)}}\label{rule-OEContraction}

The following rule contracts the hiatus left by
\hyperref[rule-OEHLoss]{SC085 OEHLoss}.

\begin{ReaderFacingFoma}
define OEContraction [
    {*a} {*a} -> {*ā},
    {*e} {*e} -> {*ē},
    {*i} {*i} -> {*ī},
    {*o} {*o} -> {*ō},
    {*u} {*u} -> {*ū},
    {*ea} {*a} -> {*ēa},
    {*ēa} {*a} -> {*ēa},
    {*eo} {*a} -> {*ēo},
    {*ēo} {*a} -> {*ēo},
    {*eo} {*o} -> {*ēo},
    {*ēo} {*o} -> {*ēo},
    {*éo} {*o} -> {*ḗo},
    {*ḗo} {*o} -> {*ḗo},
    {*ā} {*a} -> {*ā},
    {*ā} {*e} -> {*ā},
    {*ē} {*a} -> {*ē},
    {*ē} {*e} -> {*ē},
    {*ḗ} {*a} -> {*ḗ},
    {*ḗ} {*e} -> {*ḗ},
    {*ī} {*a} -> {*ī},
    {*ī} {*e} -> {*ī},
    {*ḯ} {*a} -> {*ḯ},
    {*ḯ} {*e} -> {*ḯ},
    {*ō} {*a} -> {*ō},
    {*ō} {*e} -> {*ō},
    {*ū} {*a} -> {*ū},
    {*ū} {*e} -> {*ū}
];
\end{ReaderFacingFoma}

The rule contracts the vowel sequences created after \emph{h}-loss,
producing \emph{flēon} `flee', \emph{slēan} `slay', and \emph{tēon}
`draw'.

Moving contraction before \hyperref[rule-OEHLoss]{SC085 OEHLoss} makes
PGmc \Recon{fléuxaną} `flee' yield \Pred{flēoan} rather than expected OE
\emph{flēon} `flee', PGmc \Recon{sláxaną} `slay' yield \Pred{sleaan}
rather than expected \emph{slēan} `slay', PGmc \Recon{téxun} `draw'
yield \Pred{teoon} rather than expected \emph{tēon} `draw', and PGmc
\Recon{táixōn} `toe' yield \Pred{tāe} rather than expected \emph{tā}
`toe'. The derivations require \hyperref[rule-OEContraction]{SC086
OEContraction} to follow \hyperref[rule-OEHLoss]{SC085 OEHLoss}. Moving
it later within the tested range before
\hyperref[rule-OERMetathesis]{SC087 OERMetathesis} creates no new
failure. The more distant \hyperref[rule-OEWeakTailReduction]{SC078
OEWeakTailReduction} relation establishes only that weak-tail reduction
precedes contraction.

\newpage

\section{R-metathesis}\label{r-metathesis}

\subsection{Historical discussion}\label{historical-discussion-38}

Sievers-Brunner describes r-metathesis in forms such as \emph{berstan}
`burst', \emph{forst} `frost', and \emph{cærse} `cress'
(\citeproc{ref-SieversBrunner1965}{Brunner 1965, 159}, §179). Luick
likewise treats it as a later rearrangement whose interaction with
breaking remains variable (\citeproc{ref-Luick1914}{Luick 1914-\/-1940,
201}).

The evidence establishes that breaking precedes metathesis. It does not
establish an ordering relation between
\hyperref[rule-OEContraction]{SC086 OEContraction} and
\hyperref[rule-OERMetathesis]{SC087 OERMetathesis}.

\subsection{\texorpdfstring{\CAPRRuleHeading{SC087. Metathesis of \emph{*r} with a following short vowel}{OERMetathesis}}{}}\label{rule-OERMetathesis}

\begin{ReaderFacingFoma}
define OERMetathesis [
    {*r} {*e} -> {*e} {*r} || EnglishStarConsonant _ {*s} {*t},
    {*r} {*u} -> {*u} {*r} || EnglishStarConsonant _ {*s} {*t},
    {*r} {*i} -> {*i} {*r} || EnglishStarConsonant _ {*s} {*t},
    {*r} {*o} -> {*o} {*r} || EnglishStarConsonant _ {*s} {*t},
    {*r} {*a} -> {*a} {*r} || EnglishStarConsonant _ {*s} {*t},
    {*r} {*é} -> {*é} {*r} || EnglishStarConsonant _ {*s} {*t},
    {*r} {*ó} -> {*ó} {*r} || EnglishStarConsonant _ {*s} {*t},
    {*r} {*á} -> {*á} {*r} || EnglishStarConsonant _ {*s} {*t}
];
\end{ReaderFacingFoma}

The rule moves \emph{*r} across a following short vowel in the relevant
late clusters, producing forms such as \emph{berstan} `burst' where an
earlier order would still show a broken vowel sequence.

Moving the rule before \hyperref[rule-OEBreaking]{SC044 OEBreaking}
makes PGmc \Recon{bréstaną} `burst' yield \Pred{beorstan} rather than
expected OE \emph{berstan} `burst'. On this evidence, I take
\hyperref[rule-OERMetathesis]{SC087 OERMetathesis} to follow
\hyperref[rule-OEBreaking]{SC044 OEBreaking}. Moving it later within the
tested sequence alters none of the checked outputs.

The checked forms fix the earlier relation but do not identify a
corresponding later constraint. The sources treat r-metathesis as a late
rearrangement after breaking without placing it immediately beside
contraction.

\newpage

\part{Lexical derivations}

\chapter*{Word-by-word derivations}\label{word-by-word-derivations}
\addcontentsline{toc}{chapter}{Word-by-word derivations}

\section*{Introduction}\label{introduction-1}
\addcontentsline{toc}{section}{Introduction}

The catalogue groups words by the kind of historical explanation they
require. Regular derivations establish the reach of the sound laws. The
remaining classes mark, without concealment, the places where
attestation, morphology, analogy, or an unresolved mismatch intervenes.

The catalogue proceeds from words; the preceding part proceeds from
rules. Together they allow each proposed sound law to be tested against
its lexical witnesses and each etymology against the complete ordered
history.

\section*{Data and sources}\label{data-and-sources}
\addcontentsline{toc}{section}{Data and sources}

This volume assembles the lexical corpus from the aligned Germanic
dataset and the compact derivation traces that accompany each entry.
Comparative dictionaries, Old English dictionaries, and historical
grammars are cited in the prose where they bear on particular lexical
arguments.

The result is a lexical catalogue rather than a separate report on
citation method or trace machinery.

\section*{Transducer and derivation
method}\label{transducer-and-derivation-method}
\addcontentsline{toc}{section}{Transducer and derivation method}

Four objects must be distinguished in every derivation: the citation
reconstruction, the selected input, the transducer outcome, and the Old
English target. The summary identifies them where they differ; the boxed
trace then divides the changes into Earlier Germanic and Old English
stages.

\section*{Derivation classes}\label{derivation-classes}
\addcontentsline{toc}{section}{Derivation classes}

The lexical catalogue is ordered by seven derivation classes in the
current manifest. Counts in this alpha are:

\begin{itemize}
\tightlist
\item
  Regular derivations: \textbf{70}
\item
  Attested variants: \textbf{4}
\item
  Early analogy: \textbf{34}
\item
  Late analogy: \textbf{28}
\item
  Reconstructed Old English comparators: \textbf{3}
\item
  Known but unmodelled developments: \textbf{3}
\item
  Unexplained or deliberately unmodelled exceptions: \textbf{5}
\end{itemize}

\clearpage

\chapter{Regular derivations}\label{regular-derivations}

In these entries the selected Germanic input yields the Old English
reflex by regular sound change. They establish the standard against
which the following analogical and exceptional histories must be judged.

\section{\texorpdfstring{adder --- OE
\emph{nǣdre}}{adder --- OE nǣdre}}\label{adder-oe-nux1e3dre}

\index[iv]{02oe@\textbf{Old English}!naedre@\emph{nǣdre}}
\index[iv]{03pgmc@\textbf{Proto-Germanic}!nedron@*nḗdrōn}

Derivation: \emph{*nḗdrōn} \textgreater{} \emph{nǣdre} (regular).

\subsection{Derivation trace}\label{derivation-trace}

Proto input: \emph{*nḗdrōn}

\begingroup
\setlength{\fboxsep}{6pt}

\noindent\fbox{%
\begin{minipage}{0.97\linewidth}
\small
\begin{tabularx}{\linewidth}{@{}>{\raggedright\arraybackslash}p{0.460\linewidth}>{\centering\arraybackslash}X>{\raggedright\arraybackslash}p{0.460\linewidth}@{}}
\begin{minipage}[t]{\linewidth}
\centering\textbf{Earlier Germanic changes}\par
\vspace{0.35em}
\raggedright
\centering\textbf{}\par
\raggedright
\vspace{0.2em}
\begin{tabular}{@{}>{\raggedright\arraybackslash}p{0.74\linewidth}@{\hspace{0.55em}}>{\raggedright\arraybackslash}p{0.16\linewidth}@{\hspace{0.25em}}}
\multicolumn{2}{@{}l@{}}{*Proto-West Germanic*} \\
\multicolumn{2}{@{}l@{}}{[no change]} \\
\multicolumn{2}{@{}l@{}}{*Northwest Germanic*} \\
NWGmc N Stem N Loss & \emph{*nḗdrǭ} \\
\mbox{NWGmc Long E Lowering} & \emph{*nǣdrǭ} \\
\end{tabular}
\end{minipage}
&
&
\begin{minipage}[t]{\linewidth}
\centering\textbf{Old English changes}\par
\vspace{0.35em}
\raggedright
\centering\textbf{}\par
\raggedright
\vspace{0.2em}
\begin{tabular}{@{}>{\raggedright\arraybackslash}p{0.74\linewidth}@{\hspace{0.55em}}>{\raggedright\arraybackslash}p{0.16\linewidth}@{\hspace{0.25em}}}
\multicolumn{2}{@{}l@{}}{*Old English*} \\
OE Unstressed Long Vowel Shortening & \emph{*nǣdræ} \\
OE Unstressed AE Merger & \emph{*nǣdre} \\
\end{tabular}
\end{minipage}
\\
\end{tabularx}
\end{minipage}%
} \endgroup

Old English form: \emph{nǣdre}

\subsection{Reconstruction and comparative
evidence}\label{reconstruction-and-comparative-evidence}

Kroonen distinguishes the masculine snake word {\emph{*nadra-}} from a
feminine ablauting formation {\emph{*nēdrōn-}}, and gives Old English
{\emph{nǣdre}} `adder', {\emph{næddre}} `adder' under the latter
(\citeproc{ref-Kroonen2013}{Kroonen 2013, 426}). Orel likewise points
from the masculine entry to a feminine {\Recon{nēdrōn} `adder'}
\textasciitilde{} {\Recon{nadrōn} `adder'} type
(\citeproc{ref-Orel2003}{Orel 2003, 325}).

The derivational input therefore is not a reshaped convenience form. It
is the comparative reconstruction that specifically underlies the Old
English noun.

\subsection{Old English evidence}\label{old-english-evidence}

The Old English word is securely represented by {\emph{nǣdre}} `adder',
with {\emph{næddre}} `adder' as a secondary variant. Clark Hall
cross-references {\emph{næddre}} to {\emph{nædre}} `adder', and Fulk
treats {\emph{næddre}} as the later geminated form beside the older base
(\citeproc{ref-ClarkHall1960}{Clark Hall 1960, 225};
\citeproc{ref-Fulk2018}{Fulk 2018, 149}).

\subsection{Development to Old
English}\label{development-to-old-english}

From {\Recon{nḗdrōn} `adder'}, the stressed long mid vowel develops to
Old English {\emph{nǣdre}} `adder', and the weak feminine ending remains
as final \emph{-e}, giving {\emph{nǣdre}}. The doubled consonant of
{\emph{næddre}} `adder' is secondary and does not alter the inherited
base form.

\section{\texorpdfstring{bake --- OE
\emph{bacan}}{bake --- OE bacan}}\label{bake-oe-bacan}

\index[iv]{02oe@\textbf{Old English}!bacan@\emph{bacan}}
\index[iv]{03pgmc@\textbf{Proto-Germanic}!bakana@*bákaną}
\index[iv]{09ohg@\textbf{Old High German}!backan@backan}
\index[iv]{09ohg@\textbf{Old High German}!bahhan@bahhan}

Derivation: \emph{*bákaną} \textgreater{} \emph{bacan} (regular).

\subsection{Derivation trace}\label{derivation-trace-1}

Proto input: \emph{*bákaną}

\begingroup
\setlength{\fboxsep}{6pt}

\noindent\fbox{%
\begin{minipage}{0.97\linewidth}
\small
\begin{tabularx}{\linewidth}{@{}>{\raggedright\arraybackslash}p{0.320\linewidth}>{\centering\arraybackslash}X>{\raggedright\arraybackslash}p{0.520\linewidth}@{}}
\begin{minipage}[t]{\linewidth}
\centering\textbf{Earlier Germanic changes}\par
\vspace{0.35em}
\raggedright
\centering\textbf{}\par
\raggedright
\vspace{0.2em}
\begin{tabular}{@{}>{\raggedright\arraybackslash}p{0.68\linewidth}@{\hspace{0.55em}}>{\raggedright\arraybackslash}p{0.16\linewidth}@{\hspace{0.25em}}}
\multicolumn{2}{@{}l@{}}{*Proto-West Germanic*} \\
\multicolumn{2}{@{}l@{}}{[no change]} \\
\multicolumn{2}{@{}l@{}}{*Northwest Germanic*} \\
\multicolumn{2}{@{}l@{}}{[no change]} \\
\end{tabular}
\end{minipage}
&
&
\begin{minipage}[t]{\linewidth}
\centering\textbf{Old English changes}\par
\vspace{0.35em}
\raggedright
\centering\textbf{}\par
\raggedright
\vspace{0.2em}
\begin{tabular}{@{}>{\raggedright\arraybackslash}p{0.74\linewidth}@{\hspace{0.55em}}>{\raggedright\arraybackslash}p{0.16\linewidth}@{\hspace{0.25em}}}
\multicolumn{2}{@{}l@{}}{*Old English*} \\
Anglo Frisian Brightening & \emph{*bækaną} \\
OE A Restoration & \emph{*bakaną} \\
OE Heavy Syllable Nasal Apocope & \emph{*bakan} \\
OE Secondary Nasalization & \emph{*bakąn} \\
OE Weak Tail Reduction & \emph{*bakan} \\
\end{tabular}
\end{minipage}
\\
\end{tabularx}
\end{minipage}%
} \endgroup

Old English form: \emph{bacan}

\subsection{Reconstruction and comparative
evidence}\label{reconstruction-and-comparative-evidence-1}

Orel reconstructs the verb as {\Recon{bakanan} `bake'} and cites Old
English {\emph{bacan}} `bake' beside Old High German \emph{backan,
bahhan} (\citeproc{ref-Orel2003}{Orel 2003}). Campbell gives
{\emph{bacan}} as one of the standard examples of Old English
A-restoration before a single consonant, and Ringe and Taylor state the
same development from {\Recon{bakan} `bake'} to Old English
{\emph{bacan}} (\citeproc{ref-Campbell1959}{Campbell 1959, 61};
\citeproc{ref-RingeTaylor2014}{Ringe and Taylor 2014}).

\subsection{Old English evidence}\label{old-english-evidence-1}

Bosworth-Toller and Clark Hall both record {\emph{bacan}} `bake' as the
ordinary Old English verb (\citeproc{ref-BosworthToller1898}{Bosworth
and Toller 1898, 72}; \citeproc{ref-ClarkHall1960}{Clark Hall 1960}).
The target in this entry is therefore the attested infinitive headword
itself, not a selected oblique or finite paradigm cell.

\subsection{Development to Old
English}\label{development-to-old-english-1}

From {\Recon{bákaną} `bake'}, Anglo-Frisian brightening first gives
\Recon{bækaną} `bake'. A-restoration then returns the stem vowel to
\emph{a} before single \emph{k} plus the back-vocalic infinitive suffix,
and later apocope and weak-tail reduction yield {\emph{bacan}} `bake'
(\citeproc{ref-Campbell1959}{Campbell 1959, 61};
\citeproc{ref-RingeTaylor2014}{Ringe and Taylor 2014}). The development
is therefore straightforward: {\Recon{bákaną}} \textgreater{}
{\emph{bacan}}.

\section{\texorpdfstring{beech --- OE
\emph{bōc}}{beech --- OE bōc}}\label{beech-oe-bux14dc}

\index[iv]{02oe@\textbf{Old English}!boc@\emph{bōc}}
\index[iv]{03pgmc@\textbf{Proto-Germanic}!boko@*bōkō}

Derivation: \emph{*bōkō} \textgreater{} \emph{bōc} (regular).

\subsection{Derivation trace}\label{derivation-trace-2}

Proto input: \emph{*bōkō}

\begingroup
\setlength{\fboxsep}{6pt}

\noindent\fbox{%
\begin{minipage}{0.97\linewidth}
\small
\begin{tabularx}{\linewidth}{@{}>{\raggedright\arraybackslash}p{0.460\linewidth}>{\centering\arraybackslash}X>{\raggedright\arraybackslash}p{0.460\linewidth}@{}}
\begin{minipage}[t]{\linewidth}
\centering\textbf{Earlier Germanic changes}\par
\vspace{0.35em}
\raggedright
\centering\textbf{}\par
\raggedright
\vspace{0.2em}
\begin{tabular}{@{}>{\raggedright\arraybackslash}p{0.74\linewidth}@{\hspace{0.55em}}>{\raggedright\arraybackslash}p{0.16\linewidth}@{\hspace{0.25em}}}
\multicolumn{2}{@{}l@{}}{*Proto-West Germanic*} \\
\multicolumn{2}{@{}l@{}}{[no change]} \\
\multicolumn{2}{@{}l@{}}{*Northwest Germanic*} \\
\mbox{NWGmc Final Long O Raising} & \emph{*bōku} \\
\end{tabular}
\end{minipage}
&
&
\begin{minipage}[t]{\linewidth}
\centering\textbf{Old English changes}\par
\vspace{0.35em}
\raggedright
\centering\textbf{}\par
\raggedright
\vspace{0.2em}
\begin{tabular}{@{}>{\raggedright\arraybackslash}p{0.74\linewidth}@{\hspace{0.55em}}>{\raggedright\arraybackslash}p{0.16\linewidth}@{\hspace{0.25em}}}
\multicolumn{2}{@{}l@{}}{*Old English*} \\
\mbox{OE High Vowel Apocope} & \emph{*bōk} \\
\end{tabular}
\end{minipage}
\\
\end{tabularx}
\end{minipage}%
} \endgroup

Old English form: \emph{bōc}

\subsection{Reconstruction and comparative
evidence}\label{reconstruction-and-comparative-evidence-2}

Kroonen gives the beech noun as \emph{*bōk(j)ō-} and cites Old English
\emph{boc} `beech', {\emph{bēce}} `beech' among its reflexes
(\citeproc{ref-Kroonen2013}{Kroonen 2013}). The derivational input
\Recon{bōkō} `beech' is the nominative-singular shape of that family,
which is the relevant comparison form here.

\subsection{Old English evidence}\label{old-english-evidence-2}

Kroonen's Old English evidence already separates the paradigm material:
\emph{boc} `beech' as the nominative form and {\emph{bēce}} `beech' as
an oblique form (\citeproc{ref-Kroonen2013}{Kroonen 2013}). The relevant
comparator is therefore \emph{bōc} `beech'; {\emph{bēċe}} `beech'
remains related paradigm evidence rather than the form chosen for this
comparison.

\subsection{Development to Old
English}\label{development-to-old-english-2}

With nominative input \Recon{bōkō} `beech', the development is compact.
Northwest Germanic final long \emph{ō} raises to \emph{u}, and later
high-vowel apocope leaves \emph{bōc} `beech'. The regular comparison is
therefore \emph{*bōkō} \textgreater{} \emph{bōc}.

\section{\texorpdfstring{begin --- OE
\emph{beġinnan}}{begin --- OE beġinnan}}\label{begin-oe-beux121innan}

\index[iv]{02oe@\textbf{Old English}!beginnan@\emph{beġinnan}}
\index[iv]{03pgmc@\textbf{Proto-Germanic}!biginnana@*bigínnaną}

Derivation: \emph{*bigínnaną} \textgreater{} \emph{beġinnan} (regular).

\subsection{Derivation trace}\label{derivation-trace-3}

Proto input: \emph{*bigínnaną}

\begingroup
\setlength{\fboxsep}{6pt}

\noindent\fbox{%
\begin{minipage}{0.97\linewidth}
\footnotesize
\begin{tabularx}{\linewidth}{@{}>{\raggedright\arraybackslash}p{0.320\linewidth}>{\centering\arraybackslash}X>{\raggedright\arraybackslash}p{0.520\linewidth}@{}}
\begin{minipage}[t]{\linewidth}
\centering\textbf{Earlier Germanic changes}\par
\vspace{0.35em}
\raggedright
\centering\textbf{}\par
\raggedright
\vspace{0.2em}
\begin{tabular}{@{}>{\raggedright\arraybackslash}p{0.68\linewidth}@{\hspace{0.55em}}>{\raggedright\arraybackslash}p{0.16\linewidth}@{\hspace{0.25em}}}
\multicolumn{2}{@{}l@{}}{*Proto-West Germanic*} \\
\multicolumn{2}{@{}l@{}}{[no change]} \\
\multicolumn{2}{@{}l@{}}{*Northwest Germanic*} \\
\multicolumn{2}{@{}l@{}}{[no change]} \\
\end{tabular}
\end{minipage}
&
&
\begin{minipage}[t]{\linewidth}
\centering\textbf{Old English changes}\par
\vspace{0.35em}
\raggedright
\centering\textbf{}\par
\raggedright
\vspace{0.2em}
\begin{tabular}{@{}>{\raggedright\arraybackslash}p{0.66\linewidth}@{\hspace{0.55em}}>{\raggedright\arraybackslash}p{0.24\linewidth}@{\hspace{0.25em}}}
\multicolumn{2}{@{}l@{}}{*Old English*} \\
OE Heavy Syllable Nasal Apocope & \emph{*bigínnan} \\
OE Secondary Nasalization & \emph{*bigínnąn} \\
OE Velar Palatalization & \emph{*biʤínnąn} \\
OE Prefix I Reduction & \emph{*bĕʤínnąn} \\
OE Weak Tail Reduction & \emph{*bĕʤínnan} \\
\end{tabular}
\end{minipage}
\\
\end{tabularx}
\end{minipage}%
} \endgroup

Old English form: \emph{beġinnan}

\subsection{Reconstruction and comparative
evidence}\label{reconstruction-and-comparative-evidence-3}

The verb is modeled here as inherited {\Recon{bigínnaną} `begin'}. Ringe
and Taylor state that intervocalic \emph{*g} is palatalized between
front vowels in Old English (\citeproc{ref-RingeTaylor2014}{Ringe and
Taylor 2014}), and Campbell lists {\emph{ginnan}} `begin' among familiar
examples of palatal \emph{g} in this verb family
(\citeproc{ref-Campbell1959}{Campbell 1959, 174}).

\subsection{Old English evidence}\label{old-english-evidence-3}

Bosworth-Toller and Clark Hall lemmatize the verb as {\emph{be-ginnan}}
`begin' / {\emph{beginnan}} `begin'
(\citeproc{ref-BosworthToller1898}{Bosworth and Toller 1898, 84};
\citeproc{ref-ClarkHall1960}{Clark Hall 1960}). Those plain-\emph{g}
dictionary spellings support the same verb that appears here in
normalized form as {\emph{beġinnan}} `begin'.

\subsection{Development note}\label{development-note}

The prefix deserves separate notice. Ringe and Taylor explicitly cite
\emph{bi- \textgreater{} be-} as an Old English unstressed-prefix
development (\citeproc{ref-RingeTaylor2014}{Ringe and Taylor 2014}).

\subsection{Development to Old
English}\label{development-to-old-english-3}

From {\Recon{bigínnaną} `begin'}, heavy-syllable nasal apocope yields
\Recon{bigínnan} `begin'. Intervocalic \emph{*g} between front vowels
then palatalizes to \emph{ġ}, and the unstressed prefix reduces
\emph{bi-} to \emph{be-}, giving {\emph{beġinnan}} `begin'.

\section{\texorpdfstring{bier --- OE
\emph{bǣr}}{bier --- OE bǣr}}\label{bier-oe-bux1e3r}

\index[iv]{02oe@\textbf{Old English}!baer@\emph{bǣr}}
\index[iv]{03pgmc@\textbf{Proto-Germanic}!bero@*bḗrō}

Derivation: \emph{*bḗrō} \textgreater{} \emph{bǣr} (regular).

\subsection{Derivation trace}\label{derivation-trace-4}

Proto input: \emph{*bḗrō}

\begingroup
\setlength{\fboxsep}{6pt}

\noindent\fbox{%
\begin{minipage}{0.97\linewidth}
\small
\begin{tabularx}{\linewidth}{@{}>{\raggedright\arraybackslash}p{0.460\linewidth}>{\centering\arraybackslash}X>{\raggedright\arraybackslash}p{0.460\linewidth}@{}}
\begin{minipage}[t]{\linewidth}
\centering\textbf{Earlier Germanic changes}\par
\vspace{0.35em}
\raggedright
\centering\textbf{}\par
\raggedright
\vspace{0.2em}
\begin{tabular}{@{}>{\raggedright\arraybackslash}p{0.74\linewidth}@{\hspace{0.55em}}>{\raggedright\arraybackslash}p{0.16\linewidth}@{\hspace{0.25em}}}
\multicolumn{2}{@{}l@{}}{*Proto-West Germanic*} \\
\multicolumn{2}{@{}l@{}}{[no change]} \\
\multicolumn{2}{@{}l@{}}{*Northwest Germanic*} \\
\mbox{NWGmc Final Long O Raising} & \emph{*bḗru} \\
\mbox{NWGmc Long E Lowering} & \emph{*bǣru} \\
\end{tabular}
\end{minipage}
&
&
\begin{minipage}[t]{\linewidth}
\centering\textbf{Old English changes}\par
\vspace{0.35em}
\raggedright
\centering\textbf{}\par
\raggedright
\vspace{0.2em}
\begin{tabular}{@{}>{\raggedright\arraybackslash}p{0.74\linewidth}@{\hspace{0.55em}}>{\raggedright\arraybackslash}p{0.16\linewidth}@{\hspace{0.25em}}}
\multicolumn{2}{@{}l@{}}{*Old English*} \\
\mbox{OE High Vowel Apocope} & \emph{*bǣr} \\
\end{tabular}
\end{minipage}
\\
\end{tabularx}
\end{minipage}%
} \endgroup

Old English form: \emph{bǣr}

\subsection{Reconstruction and comparative
evidence}\label{reconstruction-and-comparative-evidence-4}

Kroonen reconstructs the noun as {\emph{*bērō-}} f.~`bier' and cites Old
English {\emph{bar}} `bier', {\emph{bær}} `bier' among the reflexes
(\citeproc{ref-Kroonen2013}{Kroonen 2013, 717}). The derivational input
{\Recon{bḗrō} `bier'} is the same lexeme in the accent notation used
here.

\subsection{Old English evidence}\label{old-english-evidence-4}

Clark Hall and Bosworth-Toller lemmatize the noun as {\emph{bær}}
`bier', and Kroonen also records {\emph{bar}} `bier' beside it
(\citeproc{ref-ClarkHall1960}{Clark Hall 1960};
\citeproc{ref-BosworthToller1898}{Bosworth and Toller 1898, 73};
\citeproc{ref-Kroonen2013}{Kroonen 2013, 717}). The target {\emph{bǣr}}
`bier' is therefore a normalized long-vowel spelling of the same noun.

\subsection{Source note}\label{source-note}

Lexicographic spellings vary between {\emph{bær}} `bier' and
{\emph{bar}} `bier'. The normalized target {\emph{bǣr}} `bier'
represents the same long vowel (\citeproc{ref-ClarkHall1960}{Clark Hall
1960}; \citeproc{ref-BosworthToller1898}{Bosworth and Toller 1898, 73};
\citeproc{ref-Kroonen2013}{Kroonen 2013, 717}).

\subsection{Development to Old
English}\label{development-to-old-english-4}

From {\Recon{bḗrō} `bier'}, Northwest Germanic final long \emph{ō}
raises to \emph{u}, long \emph{ē} lowers to \emph{ǣ}, and high-vowel
apocope yields {\emph{bǣr}} `bier'. The resulting noun matches the
normalized Old English target.

\section{\texorpdfstring{birth --- OE
\emph{byrd}}{birth --- OE byrd}}\label{birth-oe-byrd}

\index[iv]{02oe@\textbf{Old English}!byrd@\emph{byrd}}
\index[iv]{03pgmc@\textbf{Proto-Germanic}!burdiz@*búrdiz}

Derivation: \emph{*búrdiz} \textgreater{} \emph{byrd} (regular).

\subsection{Derivation trace}\label{derivation-trace-5}

Proto input: \emph{*búrdiz}

\begingroup
\setlength{\fboxsep}{6pt}

\noindent\fbox{%
\begin{minipage}{0.97\linewidth}
\small
\begin{tabularx}{\linewidth}{@{}>{\raggedright\arraybackslash}p{0.460\linewidth}>{\centering\arraybackslash}X>{\raggedright\arraybackslash}p{0.460\linewidth}@{}}
\begin{minipage}[t]{\linewidth}
\centering\textbf{Earlier Germanic changes}\par
\vspace{0.35em}
\raggedright
\centering\textbf{}\par
\raggedright
\vspace{0.2em}
\begin{tabular}{@{}>{\raggedright\arraybackslash}p{0.74\linewidth}@{\hspace{0.55em}}>{\raggedright\arraybackslash}p{0.16\linewidth}@{\hspace{0.25em}}}
\multicolumn{2}{@{}l@{}}{*Proto-West Germanic*} \\
\multicolumn{2}{@{}l@{}}{[no change]} \\
\multicolumn{2}{@{}l@{}}{*Northwest Germanic*} \\
\mbox{PGmc Final Z Deletion} & \emph{*búrdi} \\
\end{tabular}
\end{minipage}
&
&
\begin{minipage}[t]{\linewidth}
\centering\textbf{Old English changes}\par
\vspace{0.35em}
\raggedright
\centering\textbf{}\par
\raggedright
\vspace{0.2em}
\begin{tabular}{@{}>{\raggedright\arraybackslash}p{0.74\linewidth}@{\hspace{0.55em}}>{\raggedright\arraybackslash}p{0.16\linewidth}@{\hspace{0.25em}}}
\multicolumn{2}{@{}l@{}}{*Old English*} \\
OE I Umlaut & \emph{*byrdi} \\
\mbox{OE High Vowel Apocope} & \emph{*byrd} \\
\end{tabular}
\end{minipage}
\\
\end{tabularx}
\end{minipage}%
} \endgroup

Old English form: \emph{byrd}

\subsection{Reconstruction and comparative
evidence}\label{reconstruction-and-comparative-evidence-5}

Kroonen cites the noun under stem-level \emph{*burdi-} and gives Old
English \emph{(ge-)byrd} among the reflexes
(\citeproc{ref-Kroonen2013}{Kroonen 2013}). The derivational input
\Recon{búrdiz} `birth' is the nominative-style form that stands behind
that stem label.

\subsection{Old English evidence}\label{old-english-evidence-5}

Clark Hall and Bosworth-Toller both attest simplex {\emph{byrd}} `birth'
as an Old English noun meaning `birth'
(\citeproc{ref-ClarkHall1960}{Clark Hall 1960};
\citeproc{ref-BosworthToller1898}{Bosworth and Toller 1898, 125}). The
prefixed form \emph{gebyrd} `birth' is also well established in the
tradition: Kroonen lists \emph{(ge-)byrd}, Bosworth-Toller has a
separate \emph{ge-byrd} `birth' entry, and Campbell cites \emph{gebyrd}
and \emph{gebyrdu} in his grammatical discussion
(\citeproc{ref-Kroonen2013}{Kroonen 2013};
\citeproc{ref-BosworthToller1898}{Bosworth and Toller 1898, 125};
\citeproc{ref-Campbell1959}{Campbell 1959}).

\subsection{Form note}\label{form-note}

The relevant comparator here is the simplex noun \emph{byrd}. The
prefixed forms remain related attested material within the same lexical
family, and Hogg's discussion of deverbal feminines provides the broader
derivational setting (\citeproc{ref-Hogg1992}{Hogg 1992}).

\subsection{Development to Old
English}\label{development-to-old-english-5}

From \Recon{búrdiz} `birth', loss of final \emph{z} gives \Recon{búrdi}
`birth'. I-umlaut fronts \emph{u} to \emph{y}, and high-vowel apocope
then yields \emph{byrd}. The result is the ordinary simplex Old English
noun.

\section{\texorpdfstring{bone --- OE
\emph{bān}}{bone --- OE bān}}\label{bone-oe-bux101n}

\index[iv]{02oe@\textbf{Old English}!ban@\emph{bān}}
\index[iv]{03pgmc@\textbf{Proto-Germanic}!baina@*báiną}

Derivation: \emph{*báiną} \textgreater{} \emph{bān} (regular).

\subsection{Derivation trace}\label{derivation-trace-6}

Proto input: \emph{*báiną}

\begingroup
\setlength{\fboxsep}{6pt}

\noindent\fbox{%
\begin{minipage}{0.97\linewidth}
\small
\begin{tabularx}{\linewidth}{@{}>{\raggedright\arraybackslash}p{0.440\linewidth}>{\centering\arraybackslash}X>{\raggedright\arraybackslash}p{0.440\linewidth}@{}}
\begin{minipage}[t]{\linewidth}
\centering\textbf{Earlier Germanic changes}\par
\vspace{0.35em}
\raggedright
\centering\textbf{}\par
\raggedright
\vspace{0.2em}
\begin{tabular}{@{}>{\raggedright\arraybackslash}p{0.70\linewidth}@{\hspace{0.55em}}>{\raggedright\arraybackslash}p{0.16\linewidth}@{\hspace{0.25em}}}
\multicolumn{2}{@{}l@{}}{*Proto-West Germanic*} \\
PWGmc Ai Monophthongization & \emph{*bāną} \\
\multicolumn{2}{@{}l@{}}{*Northwest Germanic*} \\
\multicolumn{2}{@{}l@{}}{[no change]} \\
\end{tabular}
\end{minipage}
&
&
\begin{minipage}[t]{\linewidth}
\centering\textbf{Old English changes}\par
\vspace{0.35em}
\raggedright
\centering\textbf{}\par
\raggedright
\vspace{0.2em}
\begin{tabular}{@{}>{\raggedright\arraybackslash}p{0.74\linewidth}@{\hspace{0.55em}}>{\raggedright\arraybackslash}p{0.16\linewidth}@{\hspace{0.25em}}}
\multicolumn{2}{@{}l@{}}{*Old English*} \\
OE Heavy Syllable Nasal Apocope & \emph{*bān} \\
\end{tabular}
\end{minipage}
\\
\end{tabularx}
\end{minipage}%
} \endgroup

Old English form: \emph{bān}

\subsection{Reconstruction and comparative
evidence}\label{reconstruction-and-comparative-evidence-6}

Kroonen cites the noun as {\emph{*baina-}}, and Orel gives the same
lexeme under {\Recon{bainan} `bone'} (\citeproc{ref-Kroonen2013}{Kroonen
2013}; \citeproc{ref-Orel2003}{Orel 2003}). Both are comparative
headword conventions for the same neuter noun whose Old English reflex
is {\emph{bān}} `bone'.

\subsection{Old English evidence}\label{old-english-evidence-6}

Clark Hall and Bosworth-Toller record {\emph{bān}} `bone' as the
ordinary Old English noun (\citeproc{ref-ClarkHall1960}{Clark Hall
1960}; \citeproc{ref-BosworthToller1898}{Bosworth and Toller 1898}).
Bright's glossary also distinguishes citation-form {\emph{bān}} `bone'
from oblique {\emph{bāne}} `bone'
(\citeproc{ref-BrightCassidyRingler1971}{Bright 1971}).

\subsection{Source note}\label{source-note-1}

The comparative headwords {\emph{*baina-}} and {\Recon{bainan} `bone'}
provide lexeme background. The relevant comparison form here is the
nominative-accusative singular {\Recon{báiną} `bone'}.

\subsection{Development to Old
English}\label{development-to-old-english-6}

West Germanic monophthongization turns stressed \emph{*ai} into
\emph{ā}, giving \Recon{bāną} `bone'; heavy-syllable nasal apocope then
yields {\emph{bān}} `bone'. The resulting form matches the attested Old
English citation noun.

\section{\texorpdfstring{both --- OE
\emph{bū}}{both --- OE bū}}\label{both-oe-bux16b}

\index[iv]{02oe@\textbf{Old English}!bu@\emph{bū}}
\index[iv]{03pgmc@\textbf{Proto-Germanic}!bo@*bō}

Derivation: \emph{*bō} \textgreater{} \emph{bū} (regular).

\subsection{Derivation trace}\label{derivation-trace-7}

Proto input: \emph{*bō}

\begingroup
\setlength{\fboxsep}{6pt}

\noindent\fbox{%
\begin{minipage}{0.97\linewidth}
\small
\begin{tabularx}{\linewidth}{@{}>{\raggedright\arraybackslash}p{0.540\linewidth}>{\centering\arraybackslash}X>{\raggedright\arraybackslash}p{0.300\linewidth}@{}}
\begin{minipage}[t]{\linewidth}
\centering\textbf{Earlier Germanic changes}\par
\vspace{0.35em}
\raggedright
\centering\textbf{}\par
\raggedright
\vspace{0.2em}
\begin{tabular}{@{}>{\raggedright\arraybackslash}p{0.74\linewidth}@{\hspace{0.55em}}>{\raggedright\arraybackslash}p{0.16\linewidth}@{\hspace{0.25em}}}
\multicolumn{2}{@{}l@{}}{*Proto-West Germanic*} \\
\multicolumn{2}{@{}l@{}}{[no change]} \\
\multicolumn{2}{@{}l@{}}{*Northwest Germanic*} \\
NWGmc Stressed Monosyllable O Raising & \emph{*bū} \\
\end{tabular}
\end{minipage}
&
&
\begin{minipage}[t]{\linewidth}
\centering\textbf{Old English changes}\par
\vspace{0.35em}
\raggedright
\centering\textbf{}\par
\raggedright
\vspace{0.2em}
\begin{tabular}{@{}>{\raggedright\arraybackslash}p{0.64\linewidth}@{\hspace{0.55em}}>{\raggedright\arraybackslash}p{0.16\linewidth}@{\hspace{0.25em}}}
\multicolumn{2}{@{}l@{}}{*Old English*} \\
\multicolumn{2}{@{}l@{}}{[no change]} \\
\end{tabular}
\end{minipage}
\\
\end{tabularx}
\end{minipage}%
} \endgroup

Old English form: \emph{bū}

\subsection{Reconstruction and comparative
evidence}\label{reconstruction-and-comparative-evidence-7}

Kroonen treats the Germanic numeral under \emph{*ba-} and gives the
inherited paradigm \Recon{bai} `both', \Recon{bans} `both (pl.)',
\Recon{bōz} `both'/ \Recon{bōns} `both', \Recon{bō} `both', with Old
English {\emph{bēġen}} `both', {\emph{bā}} `both', and neuter
{\emph{bū}} `both' (\citeproc{ref-Kroonen2013}{Kroonen 2013, 47}). For
the present entry, the relevant inherited form is the unextended neuter
dual {\emph{*bō}} `both'.

The older explanation of {\emph{bēġen}} `both' derives it from
\emph{*bō-jen-}, and Orel still gives OE {\emph{bezen}} `both'
(\textless{} \Recon{bō-jenō} `both') beside ON {\emph{báðir}} `both',
OFris {\emph{bēthe}} `both', OS {\emph{be-thia}} `both', and OHG
{\emph{bēde}} `both' (\citeproc{ref-Orel2003}{Orel 2003, 65}). Fulk
reports that explanation cautiously and notes Seebold's preference for a
\emph{*bō-þ-} analysis instead (\citeproc{ref-Fulk2018}{Fulk 2018, sect.
10.1}). That debate concerns {\emph{bēġen}} and the extended forms
behind Modern English {\emph{both}} `both', German {\emph{beide}}
`both', and Dutch {\emph{beide}} `both'; it does not displace the
inherited neuter {\emph{*bō}} `both' \textgreater{} {\emph{bū}} `both'
treated here.

\subsection{Old English evidence}\label{old-english-evidence-7}

The Old English dual paradigm is well established. Brunner gives
masculine {\emph{bēġen}} `both', feminine {\emph{bā}} `both', and neuter
{\emph{bū}} `both' beside \emph{bā}, with compounds such as \emph{bā}
\emph{twā} `both (two)', \emph{bū} \emph{tū} `both (two, neut.)', and
\emph{bām} `both' \emph{twām} `two'
(\citeproc{ref-SieversBrunner1965}{Brunner 1965, sect. 324} Anm. 2).
Campbell and Fulk present the same basic pattern: masculine
{\emph{bēġen}} `both', feminine \emph{bā}, neuter \emph{bā}, \emph{bū},
genitive \emph{bēġra} `both', \emph{bēġ(e)a}, and dative {\emph{bǣm}}
`both' (\citeproc{ref-Campbell1959}{Campbell 1959, sect. 683};
\citeproc{ref-Fulk2018}{Fulk 2018, sect. 10.1}).

{\emph{bū}} `both' is therefore an attested neuter dual form, not a
reconstruction. It is the cleanest target for this entry because
{\emph{bēġen}} `both' belongs to the historically more contested
\emph{*bō-jen-} / analogical zone, while \emph{bā} `both' remains a
partner form within the dual paradigm rather than the most
straightforward monosyllabic comparison.

\subsection{Development to Old
English}\label{development-to-old-english-7}

{\emph{*bō}} `both' is a stressed monosyllabic form. Campbell cites
\emph{cū} `cow', \emph{hū} `how', \emph{tū} `two', and {\emph{bū}}
`both' as examples of final accented \emph{ō} \textgreater{} \emph{ū} in
the West Germanic stage leading to Old English
(\citeproc{ref-Campbell1959}{Campbell 1959, sect. 122}). Brunner states
the same development more directly: Auslautendes \emph{ō} erscheint als
û in \emph{bū} \ldots{} cu \ldots{} \emph{hū}, \emph{tū}
(\citeproc{ref-SieversBrunner1965}{Brunner 1965, sect. 69}).

The development is therefore straightforward: {\emph{*bō}} `both'
\textgreater{} {\emph{bū}} `both'.

\subsection{Form comparison}\label{form-comparison}

The comparison below sets the relevant forms side by side. It separates
the inherited OE target from the other forms that belong to the same
broader lexical history.

{\def\LTcaptype{none} % do not increment counter
\begin{longtable}[]{@{}
  >{\raggedright\arraybackslash}p{(\linewidth - 6\tabcolsep) * \real{0.2500}}
  >{\raggedright\arraybackslash}p{(\linewidth - 6\tabcolsep) * \real{0.2500}}
  >{\raggedright\arraybackslash}p{(\linewidth - 6\tabcolsep) * \real{0.2500}}
  >{\raggedright\arraybackslash}p{(\linewidth - 6\tabcolsep) * \real{0.2500}}@{}}
\toprule\noalign{}
\begin{minipage}[b]{\linewidth}\raggedright
Form
\end{minipage} & \begin{minipage}[b]{\linewidth}\raggedright
Source / stage
\end{minipage} & \begin{minipage}[b]{\linewidth}\raggedright
Status
\end{minipage} & \begin{minipage}[b]{\linewidth}\raggedright
Relevance to this entry
\end{minipage} \\
\midrule\noalign{}
\endhead
\bottomrule\noalign{}
\endlastfoot
{\emph{*bō}} \textgreater{} {\emph{bū}} & PGmc neuter dual
\textgreater{} OE neuter dual & selected regular comparison & main line
of the entry \\
{\emph{bēġen}} & OE masculine dual & attested, but historically
contested and at least partly analogical in Kroonen & real OE evidence,
not the Old English form here \\
{\emph{bā}} & OE feminine dual; also neuter variant & attested partner
form & part of the OE paradigm, but not the chosen monosyllabic
comparator \\
{\emph{báðir}}, German {\emph{beide}}, Dutch {\emph{beide}}, Modern
English {\emph{both}} & Norse, continental West Germanic, Modern English
extended forms & related but different formation & useful background,
not the direct continuation of OE \emph{bū} \\
\end{longtable}
}

\section{\texorpdfstring{bow --- OE
\emph{bīeġan}}{bow --- OE bīeġan}}\label{bow-oe-bux12beux121an}

\index[iv]{02oe@\textbf{Old English}!biegan@\emph{bīeġan}}
\index[iv]{03pgmc@\textbf{Proto-Germanic}!baugijana@*báugijaną}

Derivation: \emph{*báugijaną} \textgreater{} \emph{bīeġan} (regular).

\subsection{Derivation trace}\label{derivation-trace-8}

Proto input: \emph{*báugijaną}

\begingroup
\setlength{\fboxsep}{6pt}

\noindent\fbox{%
\begin{minipage}{0.97\linewidth}
\footnotesize
\begin{tabularx}{\linewidth}{@{}>{\raggedright\arraybackslash}p{0.320\linewidth}>{\centering\arraybackslash}X>{\raggedright\arraybackslash}p{0.520\linewidth}@{}}
\begin{minipage}[t]{\linewidth}
\centering\textbf{Earlier Germanic changes}\par
\vspace{0.35em}
\raggedright
\centering\textbf{}\par
\raggedright
\vspace{0.2em}
\begin{tabular}{@{}>{\raggedright\arraybackslash}p{0.68\linewidth}@{\hspace{0.55em}}>{\raggedright\arraybackslash}p{0.16\linewidth}@{\hspace{0.25em}}}
\multicolumn{2}{@{}l@{}}{*Proto-West Germanic*} \\
\multicolumn{2}{@{}l@{}}{[no change]} \\
\multicolumn{2}{@{}l@{}}{*Northwest Germanic*} \\
\multicolumn{2}{@{}l@{}}{[no change]} \\
\end{tabular}
\end{minipage}
&
&
\begin{minipage}[t]{\linewidth}
\centering\textbf{Old English changes}\par
\vspace{0.35em}
\raggedright
\centering\textbf{}\par
\raggedright
\vspace{0.2em}
\begin{tabular}{@{}>{\raggedright\arraybackslash}p{0.62\linewidth}@{\hspace{0.55em}}>{\raggedright\arraybackslash}p{0.28\linewidth}@{\hspace{0.25em}}}
\multicolumn{2}{@{}l@{}}{*Old English*} \\
OE Au Fronting & \emph{*báeugijaną} \\
OE Diphthong Leveling & \emph{*bēagijaną} \\
OE Heavy Syllable Nasal Apocope & \emph{*bēagijan} \\
OE Secondary Nasalization & \emph{*bēagijąn} \\
Sievers Law Syncope & \emph{*bēagjąn} \\
OE Velar Palatalization & \emph{*bēaʤjąn} \\
OE I Umlaut & \emph{*bīeʤjąn} \\
OE Weak Tail Reduction & \emph{*bīeʤjan} \\
OE J Loss After Heavy & \emph{*bīeʤan} \\
\end{tabular}
\end{minipage}
\\
\end{tabularx}
\end{minipage}%
} \endgroup

Old English form: \emph{bīeġan}

\subsection{Reconstruction and comparative
evidence}\label{reconstruction-and-comparative-evidence-8}

Kroonen reconstructs the weak verb as \emph{*baugjan-} `to (make) bend'
and cites Old English {\emph{biegan}} `bow, bend' among its reflexes
(\citeproc{ref-Kroonen2013}{Kroonen 2013}). Ringe and Taylor give the
northwest Germanic preform \Recon{baugijana} `bow', with later
\Recon{béagjan} `bow' behind West Saxon Old English {\emph{biegan}}
`bow, bend' (\citeproc{ref-RingeTaylor2014}{Ringe and Taylor 2014}). The
entry therefore concerns the weak causative member of the bend-family,
alongside the related strong verb and noun.

\subsection{Old English evidence}\label{old-english-evidence-8}

Clark Hall lemmatizes \emph{biegan} `bow, bend' as an Old English weak
verb, and Bosworth-Toller records \emph{bigan} `bow, bend' as an
attested form (\citeproc{ref-ClarkHall1960}{Clark Hall 1960};
\citeproc{ref-BosworthToller1898}{Bosworth and Toller 1898, 102}). The
form \emph{bīeġan} `bow, bend' used here is a normalized spelling of
that attested Old English weak verb.

\subsection{Development to Old
English}\label{development-to-old-english-8}

From \Recon{báugijaną} `bow', the stem reaches pre-Old-English
\Recon{bēagjan} `bow', after which palatalization of \emph{*gj} and
i-umlaut yield West Saxon \emph{biegan} `bow'; Campbell lists
\emph{biegan} `bow' among the regular \emph{ie} outcomes of \emph{*éa}
under i-umlaut (\citeproc{ref-RingeTaylor2014}{Ringe and Taylor 2014};
\citeproc{ref-Campbell1959}{Campbell 1959, 80}). The development is
therefore straightforward: \emph{*báugijaną} \textgreater{}
\emph{bīeġan} `bow'.

\section{\texorpdfstring{breeches --- OE
\emph{brēċ}}{breeches --- OE brēċ}}\label{breeches-oe-brux113ux10b}

\index[iv]{02oe@\textbf{Old English}!brec@\emph{brēċ}}
\index[iv]{03pgmc@\textbf{Proto-Germanic}!brokiz@*brōkiz}

Derivation: \emph{*brōkiz} \textgreater{} \emph{brēċ} (regular).

\subsection{Derivation trace}\label{derivation-trace-9}

Proto input: \emph{*brōkiz}

\begingroup
\setlength{\fboxsep}{6pt}

\noindent\fbox{%
\begin{minipage}{0.97\linewidth}
\small
\begin{tabularx}{\linewidth}{@{}>{\raggedright\arraybackslash}p{0.460\linewidth}>{\centering\arraybackslash}X>{\raggedright\arraybackslash}p{0.460\linewidth}@{}}
\begin{minipage}[t]{\linewidth}
\centering\textbf{Earlier Germanic changes}\par
\vspace{0.35em}
\raggedright
\centering\textbf{}\par
\raggedright
\vspace{0.2em}
\begin{tabular}{@{}>{\raggedright\arraybackslash}p{0.74\linewidth}@{\hspace{0.55em}}>{\raggedright\arraybackslash}p{0.16\linewidth}@{\hspace{0.25em}}}
\multicolumn{2}{@{}l@{}}{*Proto-West Germanic*} \\
\multicolumn{2}{@{}l@{}}{[no change]} \\
\multicolumn{2}{@{}l@{}}{*Northwest Germanic*} \\
\mbox{PGmc Final Z Deletion} & \emph{*brōki} \\
\end{tabular}
\end{minipage}
&
&
\begin{minipage}[t]{\linewidth}
\centering\textbf{Old English changes}\par
\vspace{0.35em}
\raggedright
\centering\textbf{}\par
\raggedright
\vspace{0.2em}
\begin{tabular}{@{}>{\raggedright\arraybackslash}p{0.74\linewidth}@{\hspace{0.55em}}>{\raggedright\arraybackslash}p{0.16\linewidth}@{\hspace{0.25em}}}
\multicolumn{2}{@{}l@{}}{*Old English*} \\
OE Velar Palatalization & \emph{*brōʧi} \\
OE I Umlaut & \emph{*brēʧi} \\
\mbox{OE High Vowel Apocope} & \emph{*brēʧ} \\
\end{tabular}
\end{minipage}
\\
\end{tabularx}
\end{minipage}%
} \endgroup

Old English form: \emph{brēċ}

\subsection{Reconstruction and comparative
evidence}\label{reconstruction-and-comparative-evidence-9}

Kroonen cites the noun under \emph{*brōk-}, with Old English \emph{brōc}
`breeches' and plural `breeches' among its reflexes
(\citeproc{ref-Kroonen2013}{Kroonen 2013}). Ringe and Taylor give the
plural development directly as PNWGmc \Recon{brōkiz} `breeches'
\textgreater{} \Recon{breeci} `breeches' \textgreater{} OE \emph{bréc}
`breeches' (\citeproc{ref-RingeTaylor2014}{Ringe and Taylor 2014}). The
deeper verbal base belongs to the noun's etymological background, while
the derivational input here is the plural noun form \Recon{brōkiz}
`breeches'.

\subsection{Old English evidence}\label{old-english-evidence-9}

Bright notes \emph{brōc} `breeches' with plural \emph{brēc} `breeches',
and Clark Hall gives \emph{brēc} fp. breeches while also listing
\emph{broc} `breeches' as a feminine noun probably represented chiefly
in the plural (\citeproc{ref-BrightCassidyRingler1971}{Bright 1971};
\citeproc{ref-ClarkHall1960}{Clark Hall 1960, 64}). I write \emph{brēċ}
`breeches' for the long vowel and palatal consonant; the attested plural
is \emph{brēc}.

\subsection{Development to Old
English}\label{development-to-old-english-9}

After loss of final \emph{-z}, the stem ends in \emph{-ki}, so the velar
palatalizes and \emph{ō} `breeches' undergoes i-umlaut to \emph{ē}
`breeches'; final high-vowel apocope then yields \emph{brēċ} `breeches'
(\citeproc{ref-RingeTaylor2014}{Ringe and Taylor 2014}). The development
is therefore regular: \Recon{brōkiz} `breeches' \textgreater{}
\emph{brēċ} `breeches'.

\section{\texorpdfstring{calf --- OE
\emph{ċealf}}{calf --- OE ċealf}}\label{calf-oe-ux10bealf}

\index[iv]{02oe@\textbf{Old English}!cealf@\emph{ċealf}}
\index[iv]{03pgmc@\textbf{Proto-Germanic}!kalbaz@*kálbaz}

Derivation: \emph{*kálbaz} \textgreater{} \emph{ċealf} (regular).

\subsection{Derivation trace}\label{derivation-trace-10}

Proto input: \emph{*kálbaz}

\begingroup
\setlength{\fboxsep}{6pt}

\noindent\fbox{%
\begin{minipage}{0.97\linewidth}
\small
\begin{tabularx}{\linewidth}{@{}>{\raggedright\arraybackslash}p{0.440\linewidth}>{\centering\arraybackslash}X>{\raggedright\arraybackslash}p{0.440\linewidth}@{}}
\begin{minipage}[t]{\linewidth}
\centering\textbf{Earlier Germanic changes}\par
\vspace{0.35em}
\raggedright
\centering\textbf{}\par
\raggedright
\vspace{0.2em}
\begin{tabular}{@{}>{\raggedright\arraybackslash}p{0.74\linewidth}@{\hspace{0.55em}}>{\raggedright\arraybackslash}p{0.16\linewidth}@{\hspace{0.25em}}}
\multicolumn{2}{@{}l@{}}{*Proto-West Germanic*} \\
\multicolumn{2}{@{}l@{}}{[no change]} \\
\multicolumn{2}{@{}l@{}}{*Northwest Germanic*} \\
\mbox{PGmc Final Z Deletion} & \emph{*kálba} \\
\end{tabular}
\end{minipage}
&
&
\begin{minipage}[t]{\linewidth}
\centering\textbf{Old English changes}\par
\vspace{0.35em}
\raggedright
\centering\textbf{}\par
\raggedright
\vspace{0.2em}
\begin{tabular}{@{}>{\raggedright\arraybackslash}p{0.70\linewidth}@{\hspace{0.55em}}>{\raggedright\arraybackslash}p{0.16\linewidth}@{\hspace{0.25em}}}
\multicolumn{2}{@{}l@{}}{*Old English*} \\
PWGmc Final Bare A Loss & \emph{*kálb} \\
Anglo Frisian Brightening & \emph{*kælb} \\
OE Breaking & \emph{*kealb} \\
PGmc B Allophony & \emph{*kealβ} \\
OE Velar Palatalization & \emph{*ʧealβ} \\
\end{tabular}
\end{minipage}
\\
\end{tabularx}
\end{minipage}%
} \endgroup

Old English form: \emph{ċealf}

\subsection{Reconstruction and comparative
evidence}\label{reconstruction-and-comparative-evidence-10}

Kroonen treats the noun under \emph{*kalbiz-} and notes an older s-stem
\Recon{kalbaz} `calf' with plural \Recon{kalbizō} `calves', while Orel
cites \Recon{kalbaz} `calf' as the citation form and Ringe and Taylor
derive West Saxon \emph{Cealf} from \emph{*kalbaz}, \emph{*kalbiz-}
(\citeproc{ref-Kroonen2013}{Kroonen 2013}; \citeproc{ref-Orel2003}{Orel
2003, 248}; \citeproc{ref-RingeTaylor2014}{Ringe and Taylor 2014, 220}).
The derivational input here is the singular \Recon{kálbaz} `calf', since
the entry concerns the citation-form noun.

\subsection{Old English evidence}\label{old-english-evidence-10}

Clark Hall gives {\emph{cealf}} `calf' I. (æ, e) nm. (nap.
{\emph{cealfru}} `calves'), and Bosworth-Toller likewise records
\emph{Caelf} / \emph{Cealf} `calf' beside plural forms such as
\emph{calfur} and {\emph{cealfru}} `calves'
(\citeproc{ref-ClarkHall1960}{Clark Hall 1960};
\citeproc{ref-BosworthToller1898}{Bosworth and Toller 1898, 131}).
Campbell and Fulk show the same singular-plus-\emph{-r-} plural pattern
(\citeproc{ref-Campbell1959}{Campbell 1959};
\citeproc{ref-Fulk2018}{Fulk 2018, 193}). I write \emph{ċealf} `calf'
for the palatalized initial; the attested dictionary headword is
{\emph{cealf}} `calf'.

\subsection{Development to Old
English}\label{development-to-old-english-10}

After loss of final \emph{-z} and bare \emph{-a}, Anglo-Frisian
brightening gives \Recon{kælb} `calf', and breaking before \emph{l} plus
consonant yields \Recon{kealb} `calf'. Ringe and Taylor's account of the
lexeme and their rule for initial \emph{k} in front-vocalic environments
support the West Saxon palatalized onset represented here as \emph{ċ-},
so \Recon{kálbaz} `calf' develops regularly to \emph{ċealf} `calf'
(\citeproc{ref-RingeTaylor2014}{Ringe and Taylor 2014, 220}).

\section{\texorpdfstring{corn --- OE
\emph{corn}}{corn --- OE corn}}\label{corn-oe-corn}

\index[iv]{02oe@\textbf{Old English}!corn@\emph{corn}}
\index[iv]{03pgmc@\textbf{Proto-Germanic}!kurna@*kúrną}

Derivation: \emph{*kúrną} \textgreater{} \emph{corn} (regular).

\subsection{Derivation trace}\label{derivation-trace-11}

Proto input: \emph{*kúrną}

\begingroup
\setlength{\fboxsep}{6pt}

\noindent\fbox{%
\begin{minipage}{0.97\linewidth}
\small
\begin{tabularx}{\linewidth}{@{}>{\raggedright\arraybackslash}p{0.460\linewidth}>{\centering\arraybackslash}X>{\raggedright\arraybackslash}p{0.460\linewidth}@{}}
\begin{minipage}[t]{\linewidth}
\centering\textbf{Earlier Germanic changes}\par
\vspace{0.35em}
\raggedright
\centering\textbf{}\par
\raggedright
\vspace{0.2em}
\begin{tabular}{@{}>{\raggedright\arraybackslash}p{0.74\linewidth}@{\hspace{0.55em}}>{\raggedright\arraybackslash}p{0.16\linewidth}@{\hspace{0.25em}}}
\multicolumn{2}{@{}l@{}}{*Proto-West Germanic*} \\
\multicolumn{2}{@{}l@{}}{[no change]} \\
\multicolumn{2}{@{}l@{}}{*Northwest Germanic*} \\
\mbox{NWGmc U Lowering} & \emph{*kórną} \\
\end{tabular}
\end{minipage}
&
&
\begin{minipage}[t]{\linewidth}
\centering\textbf{Old English changes}\par
\vspace{0.35em}
\raggedright
\centering\textbf{}\par
\raggedright
\vspace{0.2em}
\begin{tabular}{@{}>{\raggedright\arraybackslash}p{0.74\linewidth}@{\hspace{0.55em}}>{\raggedright\arraybackslash}p{0.16\linewidth}@{\hspace{0.25em}}}
\multicolumn{2}{@{}l@{}}{*Old English*} \\
OE Heavy Syllable Nasal Apocope & \emph{*kórn} \\
\end{tabular}
\end{minipage}
\\
\end{tabularx}
\end{minipage}%
} \endgroup

Old English form: \emph{corn}

\subsection{Reconstruction and comparative
evidence}\label{reconstruction-and-comparative-evidence-11}

Kroonen cites the noun as \emph{*kurna-}, and Orel gives the citation
form \Recon{kurnan} `corn', both with Old English {\emph{corn}} `corn,
grain' among the reflexes (\citeproc{ref-Kroonen2013}{Kroonen 2013};
\citeproc{ref-Orel2003}{Orel 2003, 264}). The singular form
\Recon{kúrną} `corn' is the nominative-accusative singular appropriate
to the citation noun.

\subsection{Old English evidence}\label{old-english-evidence-11}

Clark Hall gives \emph{corn n.~`corn,' grain}, Bright's glossary lists
\emph{corn, n.} with genitive singular \emph{cornes}, and
Bosworth-Toller treats {\emph{corn}} `corn, grain' as an ordinary noun
headword (\citeproc{ref-ClarkHall1960}{Clark Hall 1960};
\citeproc{ref-BrightCassidyRingler1971}{Bright 1971, 347};
\citeproc{ref-BosworthToller1898}{Bosworth and Toller 1898, 144}). The
target is therefore an attested citation form, while forms such as
\emph{cornes} simply provide paradigm background.

\subsection{Development to Old
English}\label{development-to-old-english-11}

With northwest Germanic lowering, \Recon{kúrną} `corn' becomes
\Recon{kórną} `corn', and later loss of final nasal after a heavy
syllable yields \Recon{kórn} `corn', whence {\emph{corn}} `corn, grain'.
The oblique form \Recon{kurnăn} `corn' belongs to comparative background
rather than to the derivational input of this entry.

\section{\texorpdfstring{deed --- OE
\emph{dǣd}}{deed --- OE dǣd}}\label{deed-oe-dux1e3d}

\index[iv]{02oe@\textbf{Old English}!daed@\emph{dǣd}}
\index[iv]{03pgmc@\textbf{Proto-Germanic}!dediz@*dḗdiz}

Derivation: \emph{*dḗdiz} \textgreater{} \emph{dǣd} (regular).

\subsection{Derivation trace}\label{derivation-trace-12}

Proto input: \emph{*dḗdiz}

\begingroup
\setlength{\fboxsep}{6pt}

\noindent\fbox{%
\begin{minipage}{0.97\linewidth}
\small
\begin{tabularx}{\linewidth}{@{}>{\raggedright\arraybackslash}p{0.460\linewidth}>{\centering\arraybackslash}X>{\raggedright\arraybackslash}p{0.460\linewidth}@{}}
\begin{minipage}[t]{\linewidth}
\centering\textbf{Earlier Germanic changes}\par
\vspace{0.35em}
\raggedright
\centering\textbf{}\par
\raggedright
\vspace{0.2em}
\begin{tabular}{@{}>{\raggedright\arraybackslash}p{0.74\linewidth}@{\hspace{0.55em}}>{\raggedright\arraybackslash}p{0.16\linewidth}@{\hspace{0.25em}}}
\multicolumn{2}{@{}l@{}}{*Proto-West Germanic*} \\
\multicolumn{2}{@{}l@{}}{[no change]} \\
\multicolumn{2}{@{}l@{}}{*Northwest Germanic*} \\
\mbox{PGmc Final Z Deletion} & \emph{*dḗdi} \\
\mbox{NWGmc Long E Lowering} & \emph{*dǣdi} \\
\end{tabular}
\end{minipage}
&
&
\begin{minipage}[t]{\linewidth}
\centering\textbf{Old English changes}\par
\vspace{0.35em}
\raggedright
\centering\textbf{}\par
\raggedright
\vspace{0.2em}
\begin{tabular}{@{}>{\raggedright\arraybackslash}p{0.74\linewidth}@{\hspace{0.55em}}>{\raggedright\arraybackslash}p{0.16\linewidth}@{\hspace{0.25em}}}
\multicolumn{2}{@{}l@{}}{*Old English*} \\
\mbox{OE High Vowel Apocope} & \emph{*dǣd} \\
\end{tabular}
\end{minipage}
\\
\end{tabularx}
\end{minipage}%
} \endgroup

Old English form: \emph{dǣd}

\subsection{Reconstruction and comparative
evidence}\label{reconstruction-and-comparative-evidence-12}

Orel reconstructs the noun as \Recon{dēdiz} `deed', and Ringe and Taylor
derive the same inherited i-stem from Proto-Germanic \Recon{dédiz}
`deed' through northwest Germanic \Recon{dadiz} `deed'
(\citeproc{ref-Orel2003}{Orel 2003};
\citeproc{ref-RingeTaylor2014}{Ringe and Taylor 2014}). The acute accent
in \Recon{dḗdiz} `deed' marks stress on the same reconstructed long
vowel.

\subsection{Old English evidence}\label{old-english-evidence-12}

Campbell states that Primitive Germanic \emph{ē} appears as West Saxon
\emph{ǣ} but in other Old English dialects mostly as \emph{ē}, and
Brunner gives the contrast explicitly as West Saxon \emph{dǣd} `deed'
beside non-West-Saxon \emph{dēd} `deed'
(\citeproc{ref-Campbell1959}{Campbell 1959};
\citeproc{ref-SieversBrunner1965}{Brunner 1965}). Clark Hall likewise
lists \emph{dæd} `deed' and cross-refers Anglian \emph{dēd} to it
(\citeproc{ref-ClarkHall1960}{Clark Hall 1960}). West Saxon \emph{dǣd}
is therefore the relevant Old English form here, with Anglian \emph{dēd}
as a dialectal doublet.

\subsection{Development to Old
English}\label{development-to-old-english-12}

From inherited \Recon{dēdiz} `deed', loss of final \emph{-z} and the
West Saxon lowering of stressed long \emph{ē} yield \emph{dǣd} `deed';
Anglian \emph{dēd} `deed' preserves the non-West-Saxon outcome
(\citeproc{ref-Campbell1959}{Campbell 1959};
\citeproc{ref-SieversBrunner1965}{Brunner 1965}). The development
treated here is therefore the regular West Saxon line.

\section{\texorpdfstring{door --- OE
\emph{dor}}{door --- OE dor}}\label{door-oe-dor}

\index[iv]{02oe@\textbf{Old English}!dor@\emph{dor}}
\index[iv]{03pgmc@\textbf{Proto-Germanic}!dura@*dúrą}
\index[iv]{09ohg@\textbf{Old High German}!tura@tura}
\index[iv]{10ofris@\textbf{Old Frisian}!dore@dore}

Derivation: \emph{*dúrą} \textgreater{} \emph{dor} (regular).

\subsection{Derivation trace}\label{derivation-trace-13}

Proto input: \emph{*dúrą}

\begingroup
\setlength{\fboxsep}{6pt}

\noindent\fbox{%
\begin{minipage}{0.97\linewidth}
\small
\begin{tabularx}{\linewidth}{@{}>{\raggedright\arraybackslash}p{0.460\linewidth}>{\centering\arraybackslash}X>{\raggedright\arraybackslash}p{0.460\linewidth}@{}}
\begin{minipage}[t]{\linewidth}
\centering\textbf{Earlier Germanic changes}\par
\vspace{0.35em}
\raggedright
\centering\textbf{}\par
\raggedright
\vspace{0.2em}
\begin{tabular}{@{}>{\raggedright\arraybackslash}p{0.74\linewidth}@{\hspace{0.55em}}>{\raggedright\arraybackslash}p{0.16\linewidth}@{\hspace{0.25em}}}
\multicolumn{2}{@{}l@{}}{*Proto-West Germanic*} \\
\multicolumn{2}{@{}l@{}}{[no change]} \\
\multicolumn{2}{@{}l@{}}{*Northwest Germanic*} \\
\mbox{NWGmc U Lowering} & \emph{*dórą} \\
\end{tabular}
\end{minipage}
&
&
\begin{minipage}[t]{\linewidth}
\centering\textbf{Old English changes}\par
\vspace{0.35em}
\raggedright
\centering\textbf{}\par
\raggedright
\vspace{0.2em}
\begin{tabular}{@{}>{\raggedright\arraybackslash}p{0.74\linewidth}@{\hspace{0.55em}}>{\raggedright\arraybackslash}p{0.16\linewidth}@{\hspace{0.25em}}}
\multicolumn{2}{@{}l@{}}{*Old English*} \\
OE Heavy Syllable Nasal Apocope & \emph{*dór} \\
\end{tabular}
\end{minipage}
\\
\end{tabularx}
\end{minipage}%
} \endgroup

Old English form: \emph{dor}

\subsection{Reconstruction and comparative
evidence}\label{reconstruction-and-comparative-evidence-13}

Kroonen reconstructs a neuter \emph{*dura-} `gate, (single) door' and
cites Old English {\emph{dor}} `door' among its reflexes. In the same
entry he separates Old English {\emph{duru}} `door', Old Frisian
\emph{dore}, and Old High German {\emph{tura}} `door' as reflexes of
\emph{*durō-} instead (\citeproc{ref-Kroonen2013}{Kroonen 2013}).

\subsection{Old English evidence}\label{old-english-evidence-13}

Clark Hall records {\emph{dor}} `door' as a neuter noun and separately
records feminine {\emph{duru}} `door' with its own inflection
(\citeproc{ref-ClarkHall1960}{Clark Hall 1960}). Ringe and Taylor
likewise treat {\emph{duru}} `door' as an early Old English u-stem,
originally a root noun shifted into that class
(\citeproc{ref-RingeTaylor2014}{Ringe and Taylor 2014}). The Old English
form here is therefore the attested neuter {\emph{dor}} `door', while
{\emph{duru}} `door' remains a parallel Old English reflex from another
stem history.

\subsection{Development to Old
English}\label{development-to-old-english-13}

From \Recon{dúrą} `door', Northwest Germanic u-lowering gives
\Recon{dórą} `door', and heavy-syllable nasal apocope then yields
{\emph{dor}} `door'. The regular development treated in this entry is
therefore \emph{*dúrą} \textgreater{} \emph{dor}; the feminine
{\emph{duru}} `door' belongs to the separate line identified by Kroonen
and Ringe-Taylor (\citeproc{ref-Kroonen2013}{Kroonen 2013};
\citeproc{ref-RingeTaylor2014}{Ringe and Taylor 2014}).

\section{\texorpdfstring{fare --- OE
\emph{faran}}{fare --- OE faran}}\label{fare-oe-faran}

\index[iv]{02oe@\textbf{Old English}!faran@\emph{faran}}
\index[iv]{03pgmc@\textbf{Proto-Germanic}!farana@*fáraną}

Derivation: \emph{*fáraną} \textgreater{} \emph{faran} (regular).

\subsection{Derivation trace}\label{derivation-trace-14}

Proto input: \emph{*fáraną}

\begingroup
\setlength{\fboxsep}{6pt}

\noindent\fbox{%
\begin{minipage}{0.97\linewidth}
\small
\begin{tabularx}{\linewidth}{@{}>{\raggedright\arraybackslash}p{0.320\linewidth}>{\centering\arraybackslash}X>{\raggedright\arraybackslash}p{0.520\linewidth}@{}}
\begin{minipage}[t]{\linewidth}
\centering\textbf{Earlier Germanic changes}\par
\vspace{0.35em}
\raggedright
\centering\textbf{}\par
\raggedright
\vspace{0.2em}
\begin{tabular}{@{}>{\raggedright\arraybackslash}p{0.68\linewidth}@{\hspace{0.55em}}>{\raggedright\arraybackslash}p{0.16\linewidth}@{\hspace{0.25em}}}
\multicolumn{2}{@{}l@{}}{*Proto-West Germanic*} \\
\multicolumn{2}{@{}l@{}}{[no change]} \\
\multicolumn{2}{@{}l@{}}{*Northwest Germanic*} \\
\multicolumn{2}{@{}l@{}}{[no change]} \\
\end{tabular}
\end{minipage}
&
&
\begin{minipage}[t]{\linewidth}
\centering\textbf{Old English changes}\par
\vspace{0.35em}
\raggedright
\centering\textbf{}\par
\raggedright
\vspace{0.2em}
\begin{tabular}{@{}>{\raggedright\arraybackslash}p{0.74\linewidth}@{\hspace{0.55em}}>{\raggedright\arraybackslash}p{0.16\linewidth}@{\hspace{0.25em}}}
\multicolumn{2}{@{}l@{}}{*Old English*} \\
Anglo Frisian Brightening & \emph{*færaną} \\
OE A Restoration & \emph{*faraną} \\
OE Heavy Syllable Nasal Apocope & \emph{*faran} \\
OE Secondary Nasalization & \emph{*farąn} \\
OE Weak Tail Reduction & \emph{*faran} \\
\end{tabular}
\end{minipage}
\\
\end{tabularx}
\end{minipage}%
} \endgroup

Old English form: \emph{faran}

\subsection{Reconstruction and comparative
evidence}\label{reconstruction-and-comparative-evidence-14}

Kroonen gives the inherited strong verb as \emph{*faran-}, and Orel
gives the same lexeme as \Recon{faranan} `fare', both with Old English
\emph{faran} `fare' among the reflexes
(\citeproc{ref-Kroonen2013}{Kroonen 2013}; \citeproc{ref-Orel2003}{Orel
2003, 132}). Campbell also uses \emph{faran} `fare' as a standard
example of Old English A-restoration
(\citeproc{ref-Campbell1959}{Campbell 1959, 61}).

\subsection{Old English evidence}\label{old-english-evidence-14}

Clark Hall lemmatizes the strong verb as \emph{faran} `fare' and
separately records weak \emph{færan} `to frighten'; {\emph{fære}}
`fare', {\emph{færst}} `fare', and {\emph{færð}} `fare' belong to
present-tense forms of \emph{faran} rather than to the infinitive itself
(\citeproc{ref-ClarkHall1960}{Clark Hall 1960}). Bosworth-Toller
preserves the same distinction
(\citeproc{ref-BosworthToller1898}{Bosworth and Toller 1898, 108}). The
Old English form here is therefore the attested citation infinitive
\emph{faran} `fare'.

\subsection{Development to Old
English}\label{development-to-old-english-14}

From \Recon{fáraną} `fare', Anglo-Frisian brightening first gives
\Recon{færaną} `fare', but A-restoration before single \emph{r} returns
\Recon{faraną} `fare'; later apocope and weak-tail reduction yield
\emph{faran} `fare' (\citeproc{ref-Campbell1959}{Campbell 1959, 61}).
Fulk's contrast with participial faren- \textless{} \emph{*faræn-}
\textless{} \emph{*faran-} shows why fronting elsewhere in the paradigm
does not alter the infinitive headword (\citeproc{ref-Fulk2018}{Fulk
2018}).

\section{\texorpdfstring{fell --- OE
\emph{fell}}{fell --- OE fell}}\label{fell-oe-fell}

\index[iv]{02oe@\textbf{Old English}!fell@\emph{fell}}
\index[iv]{03pgmc@\textbf{Proto-Germanic}!fella@*féllą}

Derivation: \emph{*féllą} \textgreater{} \emph{fell} (regular).

\subsection{Derivation trace}\label{derivation-trace-15}

Proto input: \emph{*féllą}

\begingroup
\setlength{\fboxsep}{6pt}

\noindent\fbox{%
\begin{minipage}{0.97\linewidth}
\small
\begin{tabularx}{\linewidth}{@{}>{\raggedright\arraybackslash}p{0.440\linewidth}>{\centering\arraybackslash}X>{\raggedright\arraybackslash}p{0.440\linewidth}@{}}
\begin{minipage}[t]{\linewidth}
\centering\textbf{Earlier Germanic changes}\par
\vspace{0.35em}
\raggedright
\centering\textbf{}\par
\raggedright
\vspace{0.2em}
\begin{tabular}{@{}>{\raggedright\arraybackslash}p{0.68\linewidth}@{\hspace{0.55em}}>{\raggedright\arraybackslash}p{0.16\linewidth}@{\hspace{0.25em}}}
\multicolumn{2}{@{}l@{}}{*Proto-West Germanic*} \\
\multicolumn{2}{@{}l@{}}{[no change]} \\
\multicolumn{2}{@{}l@{}}{*Northwest Germanic*} \\
\multicolumn{2}{@{}l@{}}{[no change]} \\
\end{tabular}
\end{minipage}
&
&
\begin{minipage}[t]{\linewidth}
\centering\textbf{Old English changes}\par
\vspace{0.35em}
\raggedright
\centering\textbf{}\par
\raggedright
\vspace{0.2em}
\begin{tabular}{@{}>{\raggedright\arraybackslash}p{0.74\linewidth}@{\hspace{0.55em}}>{\raggedright\arraybackslash}p{0.16\linewidth}@{\hspace{0.25em}}}
\multicolumn{2}{@{}l@{}}{*Old English*} \\
OE Heavy Syllable Nasal Apocope & \emph{*féll} \\
\end{tabular}
\end{minipage}
\\
\end{tabularx}
\end{minipage}%
} \endgroup

Old English form: \emph{fell}

\subsection{Reconstruction and comparative
evidence}\label{reconstruction-and-comparative-evidence-15}

Kroonen reconstructs the noun as \emph{*fella-} `membrane, skin, hide'
and cites Old English {\emph{fell}} `skin, fell' beside Dutch
{\emph{vel}} `skin, fell' and German \emph{Fell}
(\citeproc{ref-Kroonen2013}{Kroonen 2013}). The derivational input
\Recon{féllą} `fell' is the derivable singular form of that same
inherited noun.

\subsection{Old English evidence}\label{old-english-evidence-15}

Clark Hall records {\emph{fell}} `skin, fell' as the noun `fell, skin,
hide', and Bright's glossary likewise gives {\emph{fell}} `skin, fell'
with inflected forms such as accusative singular \emph{fel} and dative
plural \emph{fellum} (\citeproc{ref-ClarkHall1960}{Clark Hall 1960};
\citeproc{ref-BrightCassidyRingler1971}{Bright 1971}). The target is
therefore the attested noun {\emph{fell}} `skin, fell', not the verb
\emph{fellan} or the preterite \emph{feoll}.

\subsection{Development to Old
English}\label{development-to-old-english-15}

With \Recon{féllą} `fell', no special earlier reshaping is needed:
heavy-syllable nasal apocope yields \Recon{féll} `fell', surfacing as
{\emph{fell}} `skin, fell'. The regular development treated here is
therefore \emph{*féllą} \textgreater{} \emph{fell}.

\section{\texorpdfstring{fern --- OE
\emph{fearn}}{fern --- OE fearn}}\label{fern-oe-fearn}

\index[iv]{02oe@\textbf{Old English}!fearn@\emph{fearn}}
\index[iv]{03pgmc@\textbf{Proto-Germanic}!farnaz@*fárnaz}

Derivation: \emph{*fárnaz} \textgreater{} \emph{fearn} (regular).

\subsection{Derivation trace}\label{derivation-trace-16}

Proto input: \emph{*fárnaz}

\begingroup
\setlength{\fboxsep}{6pt}

\noindent\fbox{%
\begin{minipage}{0.97\linewidth}
\small
\begin{tabularx}{\linewidth}{@{}>{\raggedright\arraybackslash}p{0.440\linewidth}>{\centering\arraybackslash}X>{\raggedright\arraybackslash}p{0.440\linewidth}@{}}
\begin{minipage}[t]{\linewidth}
\centering\textbf{Earlier Germanic changes}\par
\vspace{0.35em}
\raggedright
\centering\textbf{}\par
\raggedright
\vspace{0.2em}
\begin{tabular}{@{}>{\raggedright\arraybackslash}p{0.74\linewidth}@{\hspace{0.55em}}>{\raggedright\arraybackslash}p{0.16\linewidth}@{\hspace{0.25em}}}
\multicolumn{2}{@{}l@{}}{*Proto-West Germanic*} \\
\multicolumn{2}{@{}l@{}}{[no change]} \\
\multicolumn{2}{@{}l@{}}{*Northwest Germanic*} \\
\mbox{PGmc Final Z Deletion} & \emph{*fárna} \\
\end{tabular}
\end{minipage}
&
&
\begin{minipage}[t]{\linewidth}
\centering\textbf{Old English changes}\par
\vspace{0.35em}
\raggedright
\centering\textbf{}\par
\raggedright
\vspace{0.2em}
\begin{tabular}{@{}>{\raggedright\arraybackslash}p{0.70\linewidth}@{\hspace{0.55em}}>{\raggedright\arraybackslash}p{0.16\linewidth}@{\hspace{0.25em}}}
\multicolumn{2}{@{}l@{}}{*Old English*} \\
PWGmc Final Bare A Loss & \emph{*fárn} \\
Anglo Frisian Brightening & \emph{*færn} \\
OE Breaking & \emph{*fearn} \\
\end{tabular}
\end{minipage}
\\
\end{tabularx}
\end{minipage}%
} \endgroup

Old English form: \emph{fearn}

\subsection{Reconstruction and comparative
evidence}\label{reconstruction-and-comparative-evidence-16}

Kroonen cites the noun as masculine \emph{*farna-} and gives Old English
\emph{fearn, fern}, while Orel gives the same lexeme as neuter
\Recon{farnan} `fern' with Old English {\emph{fearn}} `fern'
(\citeproc{ref-Kroonen2013}{Kroonen 2013}; \citeproc{ref-Orel2003}{Orel
2003, 133}). Those are comparative headword conventions rather than
competing Old English outcomes; the modeled input here is the
nominative-style \Recon{fárnaz} `fern'.

\subsection{Old English evidence}\label{old-english-evidence-16}

Clark Hall gives {\emph{fearn}} `fern' as an Old English noun, and
Bosworth-Toller records {\emph{fearn}} `fern' with inflected forms such
as \emph{fearnes}, \emph{fearna}, and \emph{fearne}
(\citeproc{ref-ClarkHall1960}{Clark Hall 1960};
\citeproc{ref-BosworthToller1898}{Bosworth and Toller 1898, 219}).
Kroonen also records {\emph{fern}} `fern', but the local lexical sources
give stronger support to {\emph{fearn}} `fern' as the citation target
(\citeproc{ref-Kroonen2013}{Kroonen 2013}).

\subsection{Development to Old
English}\label{development-to-old-english-16}

From \Recon{fárnaz} `fern', loss of final \emph{-z} and final \emph{-a}
gives \Recon{fárn} `fern'; Anglo-Frisian brightening then yields
\Recon{færn} `fern', and breaking before \emph{r} plus consonant gives
{\emph{fearn}} `fern' (\citeproc{ref-Campbell1959}{Campbell 1959};
\citeproc{ref-RingeTaylor2014}{Ringe and Taylor 2014}). The development
treated here is therefore the regular \emph{rC}-breaking line.

\section{\texorpdfstring{field --- OE
\emph{feld}}{field --- OE feld}}\label{field-oe-feld}

\index[iv]{02oe@\textbf{Old English}!feld@\emph{feld}}
\index[iv]{03pgmc@\textbf{Proto-Germanic}!felthuz@*félθuz}

Derivation: \emph{*félθuz} \textgreater{} \emph{feld} (regular).

\subsection{Derivation trace}\label{derivation-trace-17}

Proto input: \emph{*félθuz}

\begingroup
\setlength{\fboxsep}{6pt}

\noindent\fbox{%
\begin{minipage}{0.97\linewidth}
\small
\begin{tabularx}{\linewidth}{@{}>{\raggedright\arraybackslash}p{0.460\linewidth}>{\centering\arraybackslash}X>{\raggedright\arraybackslash}p{0.460\linewidth}@{}}
\begin{minipage}[t]{\linewidth}
\centering\textbf{Earlier Germanic changes}\par
\vspace{0.35em}
\raggedright
\centering\textbf{}\par
\raggedright
\vspace{0.2em}
\begin{tabular}{@{}>{\raggedright\arraybackslash}p{0.74\linewidth}@{\hspace{0.55em}}>{\raggedright\arraybackslash}p{0.16\linewidth}@{\hspace{0.25em}}}
\multicolumn{2}{@{}l@{}}{*Proto-West Germanic*} \\
PWGmc L Th Voicing & \emph{*félduz} \\
\multicolumn{2}{@{}l@{}}{*Northwest Germanic*} \\
\mbox{PGmc Final Z Deletion} & \emph{*féldu} \\
\end{tabular}
\end{minipage}
&
&
\begin{minipage}[t]{\linewidth}
\centering\textbf{Old English changes}\par
\vspace{0.35em}
\raggedright
\centering\textbf{}\par
\raggedright
\vspace{0.2em}
\begin{tabular}{@{}>{\raggedright\arraybackslash}p{0.74\linewidth}@{\hspace{0.55em}}>{\raggedright\arraybackslash}p{0.16\linewidth}@{\hspace{0.25em}}}
\multicolumn{2}{@{}l@{}}{*Old English*} \\
\mbox{OE High Vowel Apocope} & \emph{*féld} \\
\end{tabular}
\end{minipage}
\\
\end{tabularx}
\end{minipage}%
} \endgroup

Old English form: \emph{feld}

\subsection{Reconstruction and comparative
evidence}\label{reconstruction-and-comparative-evidence-17}

Ringe and Taylor treat Old English {\emph{feld}} `field' as one of the
cases where earlier \emph{*felþu-} \textasciitilde{} \emph{*feldu-} may
reflect either inherited \emph{þ} \textasciitilde{} \emph{*d}
alternation or the regular West Germanic development \emph{*lþ}
\textgreater{} \emph{ld} (\citeproc{ref-RingeTaylor2014}{Ringe and
Taylor 2014, 170}). The broader proto label \Recon{félθuz} `field' can
remain as comparative background, while that narrower historical
ambiguity does not affect the regular classification.

\subsection{Old English evidence}\label{old-english-evidence-17}

Clark Hall records {\emph{feld}} `field' with oblique forms such as
\emph{felda} and \emph{felde}, and Campbell notes early place-name
spellings in \emph{-felth} beside the later standard form
(\citeproc{ref-ClarkHall1960}{Clark Hall 1960, 114};
\citeproc{ref-Campbell1959}{Campbell 1959, 169}). The Old English form
here is therefore the attested citation noun {\emph{feld}} `field', with
the older \emph{-felth} spellings as historical support rather than as
rival targets.

\subsection{Development to Old
English}\label{development-to-old-english-17}

In the modeled pathway, medial \emph{*lþ} becomes \emph{ld}, final
\emph{-z} is lost, and high-vowel apocope then yields {\emph{feld}}
`field'. Whether the voiced dental ultimately reflects inherited
alternation or the regular \emph{*lþ} \textgreater{} \emph{ld}
development, both accounts converge on the same Old English form
(\citeproc{ref-RingeTaylor2014}{Ringe and Taylor 2014, 170}).

\section{\texorpdfstring{fly --- OE
\emph{flēogan}}{fly --- OE flēogan}}\label{fly-oe-flux113ogan}

\index[iv]{02oe@\textbf{Old English}!fleogan@\emph{flēogan}}
\index[iv]{03pgmc@\textbf{Proto-Germanic}!fleugana@*fléuganą}

Derivation: \emph{*fléuganą} \textgreater{} \emph{flēogan} (regular).

\subsection{Derivation trace}\label{derivation-trace-18}

Proto input: \emph{*fléuganą}

\begingroup
\setlength{\fboxsep}{6pt}

\noindent\fbox{%
\begin{minipage}{0.97\linewidth}
\small
\begin{tabularx}{\linewidth}{@{}>{\raggedright\arraybackslash}p{0.320\linewidth}>{\centering\arraybackslash}X>{\raggedright\arraybackslash}p{0.520\linewidth}@{}}
\begin{minipage}[t]{\linewidth}
\centering\textbf{Earlier Germanic changes}\par
\vspace{0.35em}
\raggedright
\centering\textbf{}\par
\raggedright
\vspace{0.2em}
\begin{tabular}{@{}>{\raggedright\arraybackslash}p{0.68\linewidth}@{\hspace{0.55em}}>{\raggedright\arraybackslash}p{0.16\linewidth}@{\hspace{0.25em}}}
\multicolumn{2}{@{}l@{}}{*Proto-West Germanic*} \\
\multicolumn{2}{@{}l@{}}{[no change]} \\
\multicolumn{2}{@{}l@{}}{*Northwest Germanic*} \\
\multicolumn{2}{@{}l@{}}{[no change]} \\
\end{tabular}
\end{minipage}
&
&
\begin{minipage}[t]{\linewidth}
\centering\textbf{Old English changes}\par
\vspace{0.35em}
\raggedright
\centering\textbf{}\par
\raggedright
\vspace{0.2em}
\begin{tabular}{@{}>{\raggedright\arraybackslash}p{0.68\linewidth}@{\hspace{0.55em}}>{\raggedright\arraybackslash}p{0.22\linewidth}@{\hspace{0.25em}}}
\multicolumn{2}{@{}l@{}}{*Old English*} \\
OE Diphthong Leveling & \emph{*flēoganą} \\
OE Heavy Syllable Nasal Apocope & \emph{*flēogan} \\
OE Secondary Nasalization & \emph{*flēogąn} \\
OE Weak Tail Reduction & \emph{*flēogan} \\
\end{tabular}
\end{minipage}
\\
\end{tabularx}
\end{minipage}%
} \endgroup

Old English form: \emph{flēogan}

\subsection{Reconstruction and comparative
evidence}\label{reconstruction-and-comparative-evidence-18}

Ringe and Taylor derive the verb as \emph{*fleugana} `fly'
\textgreater{} OE \emph{fléogan} `fly' and elsewhere contrast West Saxon
\emph{fléogan} `fly' with Anglian \emph{flégan} `fly', alongside related
forms noted in the form note below (\citeproc{ref-RingeTaylor2014}{Ringe
and Taylor 2014}). The derivational input \Recon{fléuganą} `fly'
represents that inherited strong verb in the notation used here.

\subsection{Old English evidence}\label{old-english-evidence-18}

Clark Hall and Bosworth-Toller record \emph{flēogan} `fly' as the
ordinary Old English strong verb, and Bright gives the familiar paradigm
\emph{flēag} `flew', \emph{flugon} `flew', \emph{flogen} `flown' with
present \emph{fleogeð} `flies' (\citeproc{ref-ClarkHall1960}{Clark Hall
1960}; \citeproc{ref-BosworthToller1898}{Bosworth and Toller 1898};
\citeproc{ref-BrightCassidyRingler1971}{Bright 1971}). The target in
this entry is therefore the attested verbal infinitive itself.

\subsection{Form note}\label{form-note-1}

Ringe and Taylor also list related \emph{fléoge} `fly' / \emph{flége}
`fly' and Anglian \emph{flégan} `fly', which belong to the same family
but do not replace the infinitive \emph{flēogan} `fly' treated here
(\citeproc{ref-RingeTaylor2014}{Ringe and Taylor 2014}).

\subsection{Development to Old
English}\label{development-to-old-english-18}

From \Recon{fléuganą} `fly', Old English diphthong leveling gives
\Recon{flēoganą} `fly'; heavy-syllable nasal apocope and weak-tail
reduction then yield \emph{flēogan} `fly'
(\citeproc{ref-RingeTaylor2014}{Ringe and Taylor 2014}). The development
is therefore regular: \emph{*fléuganą} \textgreater{} \emph{flēogan}
`fly'.

\section{\texorpdfstring{forlorn --- OE
\emph{lēosan}}{forlorn --- OE lēosan}}\label{forlorn-oe-lux113osan}

\index[iv]{02oe@\textbf{Old English}!leosan@\emph{lēosan}}
\index[iv]{03pgmc@\textbf{Proto-Germanic}!leusana@*léusaną}

Derivation: \emph{*léusaną} \textgreater{} \emph{lēosan} (regular).

\subsection{Derivation trace}\label{derivation-trace-19}

Proto input: \emph{*léusaną}

\begingroup
\setlength{\fboxsep}{6pt}

\noindent\fbox{%
\begin{minipage}{0.97\linewidth}
\small
\begin{tabularx}{\linewidth}{@{}>{\raggedright\arraybackslash}p{0.320\linewidth}>{\centering\arraybackslash}X>{\raggedright\arraybackslash}p{0.520\linewidth}@{}}
\begin{minipage}[t]{\linewidth}
\centering\textbf{Earlier Germanic changes}\par
\vspace{0.35em}
\raggedright
\centering\textbf{}\par
\raggedright
\vspace{0.2em}
\begin{tabular}{@{}>{\raggedright\arraybackslash}p{0.68\linewidth}@{\hspace{0.55em}}>{\raggedright\arraybackslash}p{0.16\linewidth}@{\hspace{0.25em}}}
\multicolumn{2}{@{}l@{}}{*Proto-West Germanic*} \\
\multicolumn{2}{@{}l@{}}{[no change]} \\
\multicolumn{2}{@{}l@{}}{*Northwest Germanic*} \\
\multicolumn{2}{@{}l@{}}{[no change]} \\
\end{tabular}
\end{minipage}
&
&
\begin{minipage}[t]{\linewidth}
\centering\textbf{Old English changes}\par
\vspace{0.35em}
\raggedright
\centering\textbf{}\par
\raggedright
\vspace{0.2em}
\begin{tabular}{@{}>{\raggedright\arraybackslash}p{0.70\linewidth}@{\hspace{0.55em}}>{\raggedright\arraybackslash}p{0.20\linewidth}@{\hspace{0.25em}}}
\multicolumn{2}{@{}l@{}}{*Old English*} \\
OE Diphthong Leveling & \emph{*lēosaną} \\
OE Heavy Syllable Nasal Apocope & \emph{*lēosan} \\
OE Secondary Nasalization & \emph{*lēosąn} \\
OE Weak Tail Reduction & \emph{*lēosan} \\
\end{tabular}
\end{minipage}
\\
\end{tabularx}
\end{minipage}%
} \endgroup

Old English form: \emph{lēosan}

\subsection{Reconstruction and comparative
evidence}\label{reconstruction-and-comparative-evidence-19}

Kroonen reconstructs the verb under \emph{*leusan-} and cites prefixed
daughters such as Gothic \emph{fra-liusan} `lose' and Old English
\emph{for-lēosan} `lose, forlorn'; Orel likewise gives Old English
\emph{for-leósan} `lose' (\citeproc{ref-Kroonen2013}{Kroonen 2013};
\citeproc{ref-Orel2003}{Orel 2003}). The inherited verbal base is
therefore clear, though the daughter set often appears with the prefix.

\subsection{Old English evidence}\label{old-english-evidence-19}

The direct Old English evidence behind English \emph{forlorn} lies in
the prefixed verb \emph{forlēosan} `lose, forlorn' and especially in the
participle \emph{forloren} `forlorn', recorded by Ringe and Taylor and
in the dictionaries (\citeproc{ref-RingeTaylor2014}{Ringe and Taylor
2014}; \citeproc{ref-ClarkHall1960}{Clark Hall 1960};
\citeproc{ref-BosworthToller1898}{Bosworth and Toller 1898}). The
simplex infinitive \emph{lēosan} `lose' represents the verbal base
itself.

\subsection{Form note}\label{form-note-2}

As a base-form comparison, the simplex infinitive is \emph{lēosan}
`lose', while the English adjective continues the prefixed Old English
family \emph{forlēosan} `lose' / \emph{forloren} `forlorn'
(\citeproc{ref-RingeTaylor2014}{Ringe and Taylor 2014}).

\subsection{Development to Old
English}\label{development-to-old-english-19}

From \Recon{léusaną} `lose', Old English diphthong leveling gives
\Recon{lēosaną} `lose', and later nasal apocope and weak-tail reduction
yield \emph{lēosan} `lose' (\citeproc{ref-RingeTaylor2014}{Ringe and
Taylor 2014}). The prefixed forms follow the same verbal base with added
\emph{for-}.

\section{\texorpdfstring{gang --- OE
\emph{gang}}{gang --- OE gang}}\label{gang-oe-gang}

\index[iv]{02oe@\textbf{Old English}!gang@\emph{gang}}
\index[iv]{03pgmc@\textbf{Proto-Germanic}!gangaz@*gángaz}
\index[iv]{08on@\textbf{Old Norse}!gangr@gangr}

Derivation: \emph{*gángaz} \textgreater{} \emph{gang} (regular).

\subsection{Derivation trace}\label{derivation-trace-20}

Proto input: \emph{*gángaz}

\begingroup
\setlength{\fboxsep}{6pt}

\noindent\fbox{%
\begin{minipage}{0.97\linewidth}
\small
\begin{tabularx}{\linewidth}{@{}>{\raggedright\arraybackslash}p{0.440\linewidth}>{\centering\arraybackslash}X>{\raggedright\arraybackslash}p{0.440\linewidth}@{}}
\begin{minipage}[t]{\linewidth}
\centering\textbf{Earlier Germanic changes}\par
\vspace{0.35em}
\raggedright
\centering\textbf{}\par
\raggedright
\vspace{0.2em}
\begin{tabular}{@{}>{\raggedright\arraybackslash}p{0.74\linewidth}@{\hspace{0.55em}}>{\raggedright\arraybackslash}p{0.16\linewidth}@{\hspace{0.25em}}}
\multicolumn{2}{@{}l@{}}{*Proto-West Germanic*} \\
\multicolumn{2}{@{}l@{}}{[no change]} \\
\multicolumn{2}{@{}l@{}}{*Northwest Germanic*} \\
\mbox{PGmc Final Z Deletion} & \emph{*gánga} \\
\end{tabular}
\end{minipage}
&
&
\begin{minipage}[t]{\linewidth}
\centering\textbf{Old English changes}\par
\vspace{0.35em}
\raggedright
\centering\textbf{}\par
\raggedright
\vspace{0.2em}
\begin{tabular}{@{}>{\raggedright\arraybackslash}p{0.70\linewidth}@{\hspace{0.55em}}>{\raggedright\arraybackslash}p{0.16\linewidth}@{\hspace{0.25em}}}
\multicolumn{2}{@{}l@{}}{*Old English*} \\
PWGmc Final Bare A Loss & \emph{*gáng} \\
\end{tabular}
\end{minipage}
\\
\end{tabularx}
\end{minipage}%
} \endgroup

Old English form: \emph{gang}

\subsection{Reconstruction and comparative
evidence}\label{reconstruction-and-comparative-evidence-20}

Orel reconstructs the noun as \Recon{gangaz} `gang' and cites Old
English {\emph{gang}} `going, way' beside Old Norse \emph{gangr}, Old
Frisian \emph{gang} / \emph{gong}, Old Saxon {\emph{gang}} `going, way',
and Old High German {\emph{gang}} `going, way'
(\citeproc{ref-Orel2003}{Orel 2003}). The derivational input
\Recon{gángaz} `gang' is the same lexeme in the accent notation used
here.

\subsection{Old English evidence}\label{old-english-evidence-20}

Clark Hall and Bosworth-Toller both record {\emph{gang}} `going, way' as
the noun `going, journey, way', and Bright's glossary gives \emph{gong
(gang), m., path, course} (\citeproc{ref-ClarkHall1960}{Clark Hall
1960}; \citeproc{ref-BosworthToller1898}{Bosworth and Toller 1898, 159};
\citeproc{ref-BrightCassidyRingler1971}{Bright 1971, 392}). The target
is therefore the attested noun headword itself.

\subsection{Form note}\label{form-note-3}

This entry concerns the noun {\emph{gang}} `going, way', not the
separate verb \emph{gangan} (\citeproc{ref-ClarkHall1960}{Clark Hall
1960}; \citeproc{ref-BosworthToller1898}{Bosworth and Toller 1898,
159}).

\subsection{Development to Old
English}\label{development-to-old-english-20}

From \Recon{gángaz} `gang', loss of final \emph{-z} gives \Recon{gánga}
`gang', and later loss of final bare \emph{-a} yields {\emph{gang}}
`going, way'. The development is therefore regular: \emph{*gángaz}
\textgreater{} \emph{gang}.

\section{\texorpdfstring{give --- OE
\emph{ġiefan}}{give --- OE ġiefan}}\label{give-oe-ux121iefan}

\index[iv]{02oe@\textbf{Old English}!giefan@\emph{ġiefan}}
\index[iv]{03pgmc@\textbf{Proto-Germanic}!gebana@*gébaną}

Derivation: \emph{*gébaną} \textgreater{} \emph{ġiefan} (regular).

\subsection{Derivation trace}\label{derivation-trace-21}

Proto input: \emph{*gébaną}

\begingroup
\setlength{\fboxsep}{6pt}

\noindent\fbox{%
\begin{minipage}{0.97\linewidth}
\small
\begin{tabularx}{\linewidth}{@{}>{\raggedright\arraybackslash}p{0.320\linewidth}>{\centering\arraybackslash}X>{\raggedright\arraybackslash}p{0.520\linewidth}@{}}
\begin{minipage}[t]{\linewidth}
\centering\textbf{Earlier Germanic changes}\par
\vspace{0.35em}
\raggedright
\centering\textbf{}\par
\raggedright
\vspace{0.2em}
\begin{tabular}{@{}>{\raggedright\arraybackslash}p{0.68\linewidth}@{\hspace{0.55em}}>{\raggedright\arraybackslash}p{0.16\linewidth}@{\hspace{0.25em}}}
\multicolumn{2}{@{}l@{}}{*Proto-West Germanic*} \\
\multicolumn{2}{@{}l@{}}{[no change]} \\
\multicolumn{2}{@{}l@{}}{*Northwest Germanic*} \\
\multicolumn{2}{@{}l@{}}{[no change]} \\
\end{tabular}
\end{minipage}
&
&
\begin{minipage}[t]{\linewidth}
\centering\textbf{Old English changes}\par
\vspace{0.35em}
\raggedright
\centering\textbf{}\par
\raggedright
\vspace{0.2em}
\begin{tabular}{@{}>{\raggedright\arraybackslash}p{0.74\linewidth}@{\hspace{0.55em}}>{\raggedright\arraybackslash}p{0.16\linewidth}@{\hspace{0.25em}}}
\multicolumn{2}{@{}l@{}}{*Old English*} \\
OE Heavy Syllable Nasal Apocope & \emph{*géban} \\
OE Secondary Nasalization & \emph{*gébąn} \\
PGmc B Allophony & \emph{*géβąn} \\
OE Velar Palatalization & \emph{*ʤéβąn} \\
OE Ws Palatal Diphthongization & \emph{*ʤíeβąn} \\
OE Weak Tail Reduction & \emph{*ʤíeβan} \\
\end{tabular}
\end{minipage}
\\
\end{tabularx}
\end{minipage}%
} \endgroup

Old English form: \emph{ġiefan}

\subsection{Reconstruction and comparative
evidence}\label{reconstruction-and-comparative-evidence-21}

Kroonen reconstructs the strong verb as \emph{*geban-} and cites Old
English {\emph{giefan}} `give' among its reflexes
(\citeproc{ref-Kroonen2013}{Kroonen 2013}). Ringe and Taylor contrast
West Saxon \emph{giefan} `give' with Mercian \emph{for-geofan} `give'
and Northumbrian \emph{geafa} `give', showing that the inherited verb
takes different later dialectal shapes
(\citeproc{ref-RingeTaylor2014}{Ringe and Taylor 2014}).

\subsection{Old English evidence}\label{old-english-evidence-21}

Campbell gives \emph{gefan} (W-S \emph{giefan} `give') among examples of
initial palatalization, and Clark Hall records the verb under plain
\emph{giefan} `give' with forms such as \emph{geaf} and \emph{giefen}
(\citeproc{ref-Campbell1959}{Campbell 1959};
\citeproc{ref-ClarkHall1960}{Clark Hall 1960}). I write \emph{ġiefan}
`give' for the palatal initial.

\subsection{Dialect note}\label{dialect-note}

West Saxon \emph{ie} here reflects palatal diphthongization after
initial palatalization; non-West-Saxon forms such as \emph{geafa} `give'
or \emph{for-geofan} `give' continue the same verb without the West
Saxon vocalism (\citeproc{ref-RingeTaylor2014}{Ringe and Taylor 2014}).

\subsection{Development to Old
English}\label{development-to-old-english-21}

From \Recon{gébaną} `give', initial \emph{g} palatalizes before
\emph{e}; West Saxon palatal diphthongization then yields \emph{ie}, and
later tail reduction gives {\emph{giefan}} `give'
(\citeproc{ref-Campbell1959}{Campbell 1959};
\citeproc{ref-RingeTaylor2014}{Ringe and Taylor 2014}). The result is
therefore the regular West Saxon infinitive.

\section{\texorpdfstring{gold --- OE
\emph{gold}}{gold --- OE gold}}\label{gold-oe-gold}

\index[iv]{02oe@\textbf{Old English}!gold@\emph{gold}}
\index[iv]{03pgmc@\textbf{Proto-Germanic}!gultha@*gúlθą}

Derivation: \emph{*gúlθą} \textgreater{} \emph{gold} (regular).

\subsection{Derivation trace}\label{derivation-trace-22}

Proto input: \emph{*gúlθą}

\begingroup
\setlength{\fboxsep}{6pt}

\noindent\fbox{%
\begin{minipage}{0.97\linewidth}
\small
\begin{tabularx}{\linewidth}{@{}>{\raggedright\arraybackslash}p{0.460\linewidth}>{\centering\arraybackslash}X>{\raggedright\arraybackslash}p{0.460\linewidth}@{}}
\begin{minipage}[t]{\linewidth}
\centering\textbf{Earlier Germanic changes}\par
\vspace{0.35em}
\raggedright
\centering\textbf{}\par
\raggedright
\vspace{0.2em}
\begin{tabular}{@{}>{\raggedright\arraybackslash}p{0.74\linewidth}@{\hspace{0.55em}}>{\raggedright\arraybackslash}p{0.16\linewidth}@{\hspace{0.25em}}}
\multicolumn{2}{@{}l@{}}{*Proto-West Germanic*} \\
PWGmc L Th Voicing & \emph{*gúldą} \\
\multicolumn{2}{@{}l@{}}{*Northwest Germanic*} \\
\mbox{NWGmc U Lowering} & \emph{*góldą} \\
\end{tabular}
\end{minipage}
&
&
\begin{minipage}[t]{\linewidth}
\centering\textbf{Old English changes}\par
\vspace{0.35em}
\raggedright
\centering\textbf{}\par
\raggedright
\vspace{0.2em}
\begin{tabular}{@{}>{\raggedright\arraybackslash}p{0.74\linewidth}@{\hspace{0.55em}}>{\raggedright\arraybackslash}p{0.16\linewidth}@{\hspace{0.25em}}}
\multicolumn{2}{@{}l@{}}{*Old English*} \\
OE Heavy Syllable Nasal Apocope & \emph{*góld} \\
\end{tabular}
\end{minipage}
\\
\end{tabularx}
\end{minipage}%
} \endgroup

Old English form: \emph{gold}

\subsection{Reconstruction and comparative
evidence}\label{reconstruction-and-comparative-evidence-22}

Ringe and Taylor cite the noun as \emph{*gulþa-} / \emph{*gulda-}, and
Kroonen gives the same pair (\citeproc{ref-RingeTaylor2014}{Ringe and
Taylor 2014, 42}; \citeproc{ref-Kroonen2013}{Kroonen 2013}). The
derivational input \Recon{gúlθą} `gold' preserves the older consonantal
form while leaving open whether the medial stop reflects inherited
alternation or regular West Germanic development.

\subsection{Old English evidence}\label{old-english-evidence-22}

Bosworth-Toller and Clark Hall both record \emph{gold} as the ordinary
Old English neuter noun (\citeproc{ref-BosworthToller1898}{Bosworth and
Toller 1898, 121}; \citeproc{ref-ClarkHall1960}{Clark Hall 1960, 152}).
The target is therefore the attested citation form itself.

\subsection{Development note}\label{development-note-1}

Ringe and Taylor note that the medial stop can be understood either as
alternation \emph{*gulþa-} / \emph{*gulda-} or as the ordinary West
Germanic change \emph{*lþ} \textgreater{} \emph{ld}; both routes lead to
the same Old English consonantism (\citeproc{ref-RingeTaylor2014}{Ringe
and Taylor 2014, 42}).

\subsection{Development to Old
English}\label{development-to-old-english-22}

From \Recon{gúlθą} `gold', the regular consonant development gives
\Recon{gúldą} `gold'; Northwest Germanic / Old English lowering then
yields \Recon{góldą} `gold', and apocope gives \emph{gold}
(\citeproc{ref-Campbell1959}{Campbell 1959};
\citeproc{ref-RingeTaylor2014}{Ringe and Taylor 2014, 42}).

\section{\texorpdfstring{grave --- OE
\emph{grafan}}{grave --- OE grafan}}\label{grave-oe-grafan}

\index[iv]{02oe@\textbf{Old English}!grafan@\emph{grafan}}
\index[iv]{03pgmc@\textbf{Proto-Germanic}!grabana@*grábaną}

Derivation: \emph{*grábaną} \textgreater{} \emph{grafan} (regular).

\subsection{Derivation trace}\label{derivation-trace-23}

Proto input: \emph{*grábaną}

\begingroup
\setlength{\fboxsep}{6pt}

\noindent\fbox{%
\begin{minipage}{0.97\linewidth}
\small
\begin{tabularx}{\linewidth}{@{}>{\raggedright\arraybackslash}p{0.320\linewidth}>{\centering\arraybackslash}X>{\raggedright\arraybackslash}p{0.520\linewidth}@{}}
\begin{minipage}[t]{\linewidth}
\centering\textbf{Earlier Germanic changes}\par
\vspace{0.35em}
\raggedright
\centering\textbf{}\par
\raggedright
\vspace{0.2em}
\begin{tabular}{@{}>{\raggedright\arraybackslash}p{0.68\linewidth}@{\hspace{0.55em}}>{\raggedright\arraybackslash}p{0.16\linewidth}@{\hspace{0.25em}}}
\multicolumn{2}{@{}l@{}}{*Proto-West Germanic*} \\
\multicolumn{2}{@{}l@{}}{[no change]} \\
\multicolumn{2}{@{}l@{}}{*Northwest Germanic*} \\
\multicolumn{2}{@{}l@{}}{[no change]} \\
\end{tabular}
\end{minipage}
&
&
\begin{minipage}[t]{\linewidth}
\centering\textbf{Old English changes}\par
\vspace{0.35em}
\raggedright
\centering\textbf{}\par
\raggedright
\vspace{0.2em}
\begin{tabular}{@{}>{\raggedright\arraybackslash}p{0.70\linewidth}@{\hspace{0.55em}}>{\raggedright\arraybackslash}p{0.20\linewidth}@{\hspace{0.25em}}}
\multicolumn{2}{@{}l@{}}{*Old English*} \\
Anglo Frisian Brightening & \emph{*græbaną} \\
OE A Restoration & \emph{*grabaną} \\
OE Heavy Syllable Nasal Apocope & \emph{*graban} \\
OE Secondary Nasalization & \emph{*grabąn} \\
PGmc B Allophony & \emph{*graβąn} \\
OE Weak Tail Reduction & \emph{*graβan} \\
\end{tabular}
\end{minipage}
\\
\end{tabularx}
\end{minipage}%
} \endgroup

Old English form: \emph{grafan}

\subsection{Reconstruction and comparative
evidence}\label{reconstruction-and-comparative-evidence-23}

Campbell gives {\emph{grafan}} `dig, grave' among the standard examples
of Old English a-restoration before a single consonant and a following
back vowel, and Ringe and Taylor describe the same development for Class
VI infinitives (\citeproc{ref-Campbell1959}{Campbell 1959, 61};
\citeproc{ref-RingeTaylor2014}{Ringe and Taylor 2014}).

\subsection{Old English evidence}\label{old-english-evidence-23}

Clark Hall records {\emph{grafan}} `dig, grave' as the verb `to dig,
grave' and separately records noun \emph{græf} `grave, trench'
(\citeproc{ref-ClarkHall1960}{Clark Hall 1960}). The target here is the
attested infinitive headword of the verb.

\subsection{Development to Old
English}\label{development-to-old-english-23}

From \Recon{grábaną} `grave', Anglo-Frisian brightening first gives a
fronted stem vowel. A-restoration then returns \emph{a} before single
\emph{b} plus the back-vocalic infinitive ending, and later apocope and
weak-tail reduction yield {\emph{grafan}} `dig, grave'
(\citeproc{ref-Campbell1959}{Campbell 1959, 61};
\citeproc{ref-RingeTaylor2014}{Ringe and Taylor 2014}).

\subsection{Form note}\label{form-note-4}

Noun \emph{græf} `grave' and verbal forms such as \emph{græfð} `graves'
or past participial \emph{græfen} `graven' belong to other lexical or
paradigm positions and do not replace the infinitive \emph{grafan}
`grave' as the target here (\citeproc{ref-ClarkHall1960}{Clark Hall
1960}).

\section{\texorpdfstring{guest --- OE
\emph{ġiest}}{guest --- OE ġiest}}\label{guest-oe-ux121iest}

\index[iv]{02oe@\textbf{Old English}!giest@\emph{ġiest}}
\index[iv]{03pgmc@\textbf{Proto-Germanic}!gastiz@*gástiz}

Derivation: \emph{*gástiz} \textgreater{} \emph{ġiest} (regular).

\subsection{Derivation trace}\label{derivation-trace-24}

Proto input: \emph{*gástiz}

\begingroup
\setlength{\fboxsep}{6pt}

\noindent\fbox{%
\begin{minipage}{0.97\linewidth}
\small
\begin{tabularx}{\linewidth}{@{}>{\raggedright\arraybackslash}p{0.460\linewidth}>{\centering\arraybackslash}X>{\raggedright\arraybackslash}p{0.460\linewidth}@{}}
\begin{minipage}[t]{\linewidth}
\centering\textbf{Earlier Germanic changes}\par
\vspace{0.35em}
\raggedright
\centering\textbf{}\par
\raggedright
\vspace{0.2em}
\begin{tabular}{@{}>{\raggedright\arraybackslash}p{0.74\linewidth}@{\hspace{0.55em}}>{\raggedright\arraybackslash}p{0.16\linewidth}@{\hspace{0.25em}}}
\multicolumn{2}{@{}l@{}}{*Proto-West Germanic*} \\
\multicolumn{2}{@{}l@{}}{[no change]} \\
\multicolumn{2}{@{}l@{}}{*Northwest Germanic*} \\
\mbox{PGmc Final Z Deletion} & \emph{*gásti} \\
\end{tabular}
\end{minipage}
&
&
\begin{minipage}[t]{\linewidth}
\centering\textbf{Old English changes}\par
\vspace{0.35em}
\raggedright
\centering\textbf{}\par
\raggedright
\vspace{0.2em}
\begin{tabular}{@{}>{\raggedright\arraybackslash}p{0.74\linewidth}@{\hspace{0.55em}}>{\raggedright\arraybackslash}p{0.16\linewidth}@{\hspace{0.25em}}}
\multicolumn{2}{@{}l@{}}{*Old English*} \\
Anglo Frisian Brightening & \emph{*gæsti} \\
OE Velar Palatalization & \emph{*ʤæsti} \\
OE I Umlaut & \emph{*ʤesti} \\
OE Ws Palatal Diphthongization & \emph{*ʤiesti} \\
\mbox{OE High Vowel Apocope} & \emph{*ʤiest} \\
\end{tabular}
\end{minipage}
\\
\end{tabularx}
\end{minipage}%
} \endgroup

Old English form: \emph{ġiest}

\subsection{Reconstruction and comparative
evidence}\label{reconstruction-and-comparative-evidence-24}

Campbell and Ringe-Taylor treat the noun as an ordinary i-stem whose
West Saxon development shows palatal diphthongization, while
non-West-Saxon evidence preserves forms of the \emph{gest} `guest' type
(\citeproc{ref-Campbell1959}{Campbell 1959};
\citeproc{ref-RingeTaylor2014}{Ringe and Taylor 2014}).

\subsection{Old English evidence}\label{old-english-evidence-24}

Bosworth-Toller and Clark Hall record the word under forms such as
\emph{gist} `guest', \emph{gest} `guest', \emph{giest} `guest', and
\emph{gyst} `guest' (\citeproc{ref-BosworthToller1898}{Bosworth and
Toller 1898}; \citeproc{ref-ClarkHall1960}{Clark Hall 1960}). The Old
English form here, \emph{ġiest}, `guest' is the normalized West Saxon
form within that attested family.

\subsection{Development to Old
English}\label{development-to-old-english-24}

From \Recon{gástiz} `guest', Anglo-Frisian brightening gives a
\emph{gæst-} stage, and i-mutation affects the front vowel before the
lost high-vocalic ending. In West Saxon the initial palatal environment
then produces \emph{ie}, so the regular outcome is \emph{ġiest} `guest'
(\citeproc{ref-Campbell1959}{Campbell 1959};
\citeproc{ref-RingeTaylor2014}{Ringe and Taylor 2014}).

\subsection{Dialect note}\label{dialect-note-1}

West Saxon \emph{ġiest} `guest' is the Old English form here. Anglian
\emph{gest} `guest' and related spellings remain real Old English
comparators rather than corrections to that choice
(\citeproc{ref-RingeTaylor2014}{Ringe and Taylor 2014};
\citeproc{ref-BosworthToller1898}{Bosworth and Toller 1898}).

\section{\texorpdfstring{hair --- OE
\emph{hǣr}}{hair --- OE hǣr}}\label{hair-oe-hux1e3r}

\index[iv]{02oe@\textbf{Old English}!haer@\emph{hǣr}}
\index[iv]{03pgmc@\textbf{Proto-Germanic}!xera@*xḗrą}

Derivation: \emph{*xḗrą} \textgreater{} \emph{hǣr} (regular).

\subsection{Derivation trace}\label{derivation-trace-25}

Proto input: \emph{*xḗrą}

\begingroup
\setlength{\fboxsep}{6pt}

\noindent\fbox{%
\begin{minipage}{0.97\linewidth}
\small
\begin{tabularx}{\linewidth}{@{}>{\raggedright\arraybackslash}p{0.460\linewidth}>{\centering\arraybackslash}X>{\raggedright\arraybackslash}p{0.460\linewidth}@{}}
\begin{minipage}[t]{\linewidth}
\centering\textbf{Earlier Germanic changes}\par
\vspace{0.35em}
\raggedright
\centering\textbf{}\par
\raggedright
\vspace{0.2em}
\begin{tabular}{@{}>{\raggedright\arraybackslash}p{0.74\linewidth}@{\hspace{0.55em}}>{\raggedright\arraybackslash}p{0.16\linewidth}@{\hspace{0.25em}}}
\multicolumn{2}{@{}l@{}}{*Proto-West Germanic*} \\
\multicolumn{2}{@{}l@{}}{[no change]} \\
\multicolumn{2}{@{}l@{}}{*Northwest Germanic*} \\
\mbox{NWGmc Long E Lowering} & \emph{*xǣrą} \\
\end{tabular}
\end{minipage}
&
&
\begin{minipage}[t]{\linewidth}
\centering\textbf{Old English changes}\par
\vspace{0.35em}
\raggedright
\centering\textbf{}\par
\raggedright
\vspace{0.2em}
\begin{tabular}{@{}>{\raggedright\arraybackslash}p{0.74\linewidth}@{\hspace{0.55em}}>{\raggedright\arraybackslash}p{0.16\linewidth}@{\hspace{0.25em}}}
\multicolumn{2}{@{}l@{}}{*Old English*} \\
OE Velar Fricative Palatalization & \emph{*çǣrą} \\
OE Heavy Syllable Nasal Apocope & \emph{*çǣr} \\
\end{tabular}
\end{minipage}
\\
\end{tabularx}
\end{minipage}%
} \endgroup

Old English form: \emph{hǣr}

\subsection{Reconstruction and comparative
evidence}\label{reconstruction-and-comparative-evidence-25}

Kroonen cites the ordinary Proto-Germanic hair word as \emph{*hēra-}
(\citeproc{ref-Kroonen2013}{Kroonen 2013}). The derivational input
\Recon{xḗrą} `hair' represents that same long-ē stem in the present
derivation.

\subsection{Old English evidence}\label{old-english-evidence-25}

Clark Hall and Bosworth-Toller record \emph{hær} `hair' / \emph{hǣr}
`hair' as the ordinary Old English noun
(\citeproc{ref-ClarkHall1960}{Clark Hall 1960, 158};
\citeproc{ref-BosworthToller1898}{Bosworth and Toller 1898, 510}). The
target is therefore the attested headword itself.

\subsection{Development to Old
English}\label{development-to-old-english-25}

From \Recon{xḗrą} `hair', Northwest Germanic lowering gives a long front
vowel, and later loss of the final nasal leaves the Old English form
\emph{hǣr} `hair'. The development treated here is straightforward and
does not require any special paradigm choice.

\subsection{Form note}\label{form-note-5}

Older references to \Recon{xazwăz} `hair' belong to a different lexeme,
and the separate \emph{haddr} / \emph{heordan} / \emph{hād-} material
does not displace the ordinary simplex \emph{hǣr} `hair' treated here
(\citeproc{ref-Kroonen2013}{Kroonen 2013};
\citeproc{ref-ClarkHall1960}{Clark Hall 1960, 158}).

\section{\texorpdfstring{harvest --- OE
\emph{hierfest}}{harvest --- OE hierfest}}\label{harvest-oe-hierfest}

\index[iv]{02oe@\textbf{Old English}!hierfest@\emph{hierfest}}
\index[iv]{03pgmc@\textbf{Proto-Germanic}!xarbistuz@*xárbistuz}

Derivation: \emph{*xárbistuz} \textgreater{} \emph{hierfest} (regular).

\subsection{Derivation trace}\label{derivation-trace-26}

Proto input: \emph{*xárbistuz}

\begingroup
\setlength{\fboxsep}{6pt}

\noindent\fbox{%
\begin{minipage}{0.97\linewidth}
\footnotesize
\begin{tabularx}{\linewidth}{@{}>{\raggedright\arraybackslash}p{0.460\linewidth}>{\centering\arraybackslash}X>{\raggedright\arraybackslash}p{0.460\linewidth}@{}}
\begin{minipage}[t]{\linewidth}
\centering\textbf{Earlier Germanic changes}\par
\vspace{0.35em}
\raggedright
\centering\textbf{}\par
\raggedright
\vspace{0.2em}
\begin{tabular}{@{}>{\raggedright\arraybackslash}p{0.68\linewidth}@{\hspace{0.55em}}>{\raggedright\arraybackslash}p{0.22\linewidth}@{\hspace{0.25em}}}
\multicolumn{2}{@{}l@{}}{*Proto-West Germanic*} \\
\multicolumn{2}{@{}l@{}}{[no change]} \\
\multicolumn{2}{@{}l@{}}{*Northwest Germanic*} \\
PGmc Final Z Deletion & \emph{*xárbistu} \\
\end{tabular}
\end{minipage}
&
&
\begin{minipage}[t]{\linewidth}
\centering\textbf{Old English changes}\par
\vspace{0.35em}
\raggedright
\centering\textbf{}\par
\raggedright
\vspace{0.2em}
\begin{tabular}{@{}>{\raggedright\arraybackslash}p{0.64\linewidth}@{\hspace{0.55em}}>{\raggedright\arraybackslash}p{0.26\linewidth}@{\hspace{0.25em}}}
\multicolumn{2}{@{}l@{}}{*Old English*} \\
Anglo Frisian Brightening & \emph{*xærbistu} \\
OE Breaking & \emph{*xearbistu} \\
OE Velar Fricative Palatalization & \emph{*çearbistu} \\
PGmc B Allophony & \emph{*çearβistu} \\
OE I Umlaut & \emph{*çierβistu} \\
OE High Vowel Apocope & \emph{*çierβist} \\
OE Med Unstressed I Lowering1 & \emph{*çierβest} \\
\end{tabular}
\end{minipage}
\\
\end{tabularx}
\end{minipage}%
} \endgroup

Old English form: \emph{hierfest}

\subsection{Reconstruction and comparative
evidence}\label{reconstruction-and-comparative-evidence-26}

Bammesberger and Ringe-Taylor treat \emph{*harbist-} as the inherited
base and explain that the regular native West Saxon development would be
of the \emph{hierfest} / \emph{hyrfest} type
(\citeproc{ref-Bammesberger1997}{Bammesberger 1997};
\citeproc{ref-RingeTaylor2014}{Ringe and Taylor 2014}).

\subsection{Old English evidence}\label{old-english-evidence-26}

Bosworth-Toller and Clark Hall record \emph{hærfest} `harvest', with
\emph{herfest} `harvest' as a variant in the lexical tradition
(\citeproc{ref-BosworthToller1898}{Bosworth and Toller 1898};
\citeproc{ref-ClarkHall1960}{Clark Hall 1960}). Those attested forms
remain the main dictionary evidence for the noun.

\subsection{Development to Old
English}\label{development-to-old-english-26}

From \Recon{xárbistuz} `harvest', Anglo-Frisian brightening, breaking,
and i-mutation produce a \emph{hierbist-} stage, and later lowering of
unstressed medial \emph{i} to \emph{e} gives \emph{hierfest} `harvest'.
That is the regular West Saxon development treated here
(\citeproc{ref-RingeTaylor2014}{Ringe and Taylor 2014};
\citeproc{ref-Campbell1959}{Campbell 1959}).

\subsection{Source note}\label{source-note-2}

The Old English form here, \emph{hierfest}, represents the regular
native West Saxon outcome discussed by Bammesberger and Ringe-Taylor.
The attested Old English lexical tradition, however, is chiefly
\emph{hærfest} `harvest' / \emph{herfest} `harvest', commonly treated as
non-West-Saxon or Anglian material in West Saxon transmission
(\citeproc{ref-Bammesberger1997}{Bammesberger 1997};
\citeproc{ref-RingeTaylor2014}{Ringe and Taylor 2014}).

\section{\texorpdfstring{hedge --- OE
\emph{heġġ}}{hedge --- OE heġġ}}\label{hedge-oe-heux121ux121}

\index[iv]{02oe@\textbf{Old English}!hegg@\emph{heġġ}}
\index[iv]{03pgmc@\textbf{Proto-Germanic}!xagjaz@*xágjaz}

Derivation: \emph{*xágjaz} \textgreater{} \emph{heġġ} (regular).

\subsection{Derivation trace}\label{derivation-trace-27}

Proto input: \emph{*xágjaz}

\begingroup
\setlength{\fboxsep}{6pt}

\noindent\fbox{%
\begin{minipage}{0.97\linewidth}
\small
\begin{tabularx}{\linewidth}{@{}>{\raggedright\arraybackslash}p{0.440\linewidth}>{\centering\arraybackslash}X>{\raggedright\arraybackslash}p{0.440\linewidth}@{}}
\begin{minipage}[t]{\linewidth}
\centering\textbf{Earlier Germanic changes}\par
\vspace{0.35em}
\raggedright
\centering\textbf{}\par
\raggedright
\vspace{0.2em}
\begin{tabular}{@{}>{\raggedright\arraybackslash}p{0.70\linewidth}@{\hspace{0.55em}}>{\raggedright\arraybackslash}p{0.20\linewidth}@{\hspace{0.25em}}}
\multicolumn{2}{@{}l@{}}{*Proto-West Germanic*} \\
PWGmc J Gemination & \emph{*xággjaz} \\
\multicolumn{2}{@{}l@{}}{*Northwest Germanic*} \\
\mbox{PGmc Final Z Deletion} & \emph{*xággja} \\
\end{tabular}
\end{minipage}
&
&
\begin{minipage}[t]{\linewidth}
\centering\textbf{Old English changes}\par
\vspace{0.35em}
\raggedright
\centering\textbf{}\par
\raggedright
\vspace{0.2em}
\begin{tabular}{@{}>{\raggedright\arraybackslash}p{0.74\linewidth}@{\hspace{0.55em}}>{\raggedright\arraybackslash}p{0.16\linewidth}@{\hspace{0.25em}}}
\multicolumn{2}{@{}l@{}}{*Old English*} \\
PWGmc Final Bare A Loss & \emph{*xággj} \\
Anglo Frisian Brightening & \emph{*xæggj} \\
OE Velar Fricative Palatalization & \emph{*çæggj} \\
OE Velar Palatalization & \emph{*çæʤʤj} \\
OE I Umlaut & \emph{*çeʤʤj} \\
OE J Loss After Heavy & \emph{*çeʤʤ} \\
\end{tabular}
\end{minipage}
\\
\end{tabularx}
\end{minipage}%
} \endgroup

Old English form: \emph{heġġ}

\subsection{Reconstruction and comparative
evidence}\label{reconstruction-and-comparative-evidence-27}

The derivational input models a palatal \Recon{-gj-*} `hedge' noun whose
Old English development includes gemination, palatalization, and
i-mutation. The current derivation therefore reaches a palatal-geminate
outcome of the \emph{heġġ} `hedge' type
(\citeproc{ref-Campbell1959}{Campbell 1959}).

\subsection{Old English evidence}\label{old-english-evidence-27}

Bosworth-Toller and Clark Hall record the noun under standard spellings
\emph{hecg} `hedge' / \emph{heċġ} `hedge'
(\citeproc{ref-BosworthToller1898}{Bosworth and Toller 1898};
\citeproc{ref-ClarkHall1960}{Clark Hall 1960}). The lexical item itself
is therefore well attested even though the form compared here is
normalized.

\subsection{Development to Old
English}\label{development-to-old-english-27}

From \Recon{xágjaz} `hedge', West Germanic j-gemination first yields a
geminate stop, and later Old English palatalization and loss of final
\emph{j} produce \emph{heġġ} `hedge'. The development is treated as
regular rather than exceptional.

\subsection{Form note}\label{form-note-6}

Standard dictionary spelling is \emph{heċġ} `hedge' or \emph{hecg}
`hedge'. Normalized \emph{heġġ} `hedge' is the Old English form here,
while the ordinary lexicographic forms remain the main Old English
citation evidence (\citeproc{ref-BosworthToller1898}{Bosworth and Toller
1898}; \citeproc{ref-ClarkHall1960}{Clark Hall 1960}).

\section{\texorpdfstring{helm --- OE
\emph{helm}}{helm --- OE helm}}\label{helm-oe-helm}

\index[iv]{02oe@\textbf{Old English}!helm@\emph{helm}}
\index[iv]{03pgmc@\textbf{Proto-Germanic}!xelmaz@*xélmaz}

Derivation: \emph{*xélmaz} \textgreater{} \emph{helm} (regular).

\subsection{Derivation trace}\label{derivation-trace-28}

Proto input: \emph{*xélmaz}

\begingroup
\setlength{\fboxsep}{6pt}

\noindent\fbox{%
\begin{minipage}{0.97\linewidth}
\small
\begin{tabularx}{\linewidth}{@{}>{\raggedright\arraybackslash}p{0.440\linewidth}>{\centering\arraybackslash}X>{\raggedright\arraybackslash}p{0.440\linewidth}@{}}
\begin{minipage}[t]{\linewidth}
\centering\textbf{Earlier Germanic changes}\par
\vspace{0.35em}
\raggedright
\centering\textbf{}\par
\raggedright
\vspace{0.2em}
\begin{tabular}{@{}>{\raggedright\arraybackslash}p{0.74\linewidth}@{\hspace{0.55em}}>{\raggedright\arraybackslash}p{0.16\linewidth}@{\hspace{0.25em}}}
\multicolumn{2}{@{}l@{}}{*Proto-West Germanic*} \\
\multicolumn{2}{@{}l@{}}{[no change]} \\
\multicolumn{2}{@{}l@{}}{*Northwest Germanic*} \\
\mbox{PGmc Final Z Deletion} & \emph{*xélma} \\
\end{tabular}
\end{minipage}
&
&
\begin{minipage}[t]{\linewidth}
\centering\textbf{Old English changes}\par
\vspace{0.35em}
\raggedright
\centering\textbf{}\par
\raggedright
\vspace{0.2em}
\begin{tabular}{@{}>{\raggedright\arraybackslash}p{0.74\linewidth}@{\hspace{0.55em}}>{\raggedright\arraybackslash}p{0.16\linewidth}@{\hspace{0.25em}}}
\multicolumn{2}{@{}l@{}}{*Old English*} \\
PWGmc Final Bare A Loss & \emph{*xélm} \\
OE Velar Fricative Palatalization & \emph{*çélm} \\
\end{tabular}
\end{minipage}
\\
\end{tabularx}
\end{minipage}%
} \endgroup

Old English form: \emph{helm}

\subsection{Reconstruction and comparative
evidence}\label{reconstruction-and-comparative-evidence-28}

Kroonen cites the helmet noun as \emph{*helma-} and separately
distinguishes a different lexeme \emph{*helman-} `rudder'
(\citeproc{ref-Kroonen2013}{Kroonen 2013}). The derivational input
\Recon{xélmaz} `helm' is the nominative-style form used for the helmet
noun itself.

\subsection{Old English evidence}\label{old-english-evidence-28}

Clark Hall and Bosworth-Toller record {\emph{helm}} `helm, helmet' as
the ordinary Old English noun for `helmet', while {\emph{helma}} `helm'
belongs to a separate rudder lexeme (\citeproc{ref-ClarkHall1960}{Clark
Hall 1960}; \citeproc{ref-BosworthToller1898}{Bosworth and Toller 1898,
542}).

\subsection{Development to Old
English}\label{development-to-old-english-28}

From \Recon{xélmaz} `helm', loss of final \emph{z} and later loss of the
short final vowel yield {\emph{helm}} `helm, helmet'. The development is
therefore a straightforward citation-form match.

\subsection{Form note}\label{form-note-7}

Comparative \emph{*helma-} is headword notation for the helmet cognate
set. It should not be confused with Old English {\emph{helma}} `helm',
which is a different noun meaning `helm, rudder'
(\citeproc{ref-Kroonen2013}{Kroonen 2013};
\citeproc{ref-ClarkHall1960}{Clark Hall 1960}).

\section{\texorpdfstring{help --- OE
\emph{helpan}}{help --- OE helpan}}\label{help-oe-helpan}

\index[iv]{02oe@\textbf{Old English}!helpan@\emph{helpan}}
\index[iv]{03pgmc@\textbf{Proto-Germanic}!xelpana@*xélpaną}

Derivation: \emph{*xélpaną} \textgreater{} \emph{helpan} (regular).

\subsection{Derivation trace}\label{derivation-trace-29}

Proto input: \emph{*xélpaną}

\begingroup
\setlength{\fboxsep}{6pt}

\noindent\fbox{%
\begin{minipage}{0.97\linewidth}
\small
\begin{tabularx}{\linewidth}{@{}>{\raggedright\arraybackslash}p{0.320\linewidth}>{\centering\arraybackslash}X>{\raggedright\arraybackslash}p{0.520\linewidth}@{}}
\begin{minipage}[t]{\linewidth}
\centering\textbf{Earlier Germanic changes}\par
\vspace{0.35em}
\raggedright
\centering\textbf{}\par
\raggedright
\vspace{0.2em}
\begin{tabular}{@{}>{\raggedright\arraybackslash}p{0.68\linewidth}@{\hspace{0.55em}}>{\raggedright\arraybackslash}p{0.16\linewidth}@{\hspace{0.25em}}}
\multicolumn{2}{@{}l@{}}{*Proto-West Germanic*} \\
\multicolumn{2}{@{}l@{}}{[no change]} \\
\multicolumn{2}{@{}l@{}}{*Northwest Germanic*} \\
\multicolumn{2}{@{}l@{}}{[no change]} \\
\end{tabular}
\end{minipage}
&
&
\begin{minipage}[t]{\linewidth}
\centering\textbf{Old English changes}\par
\vspace{0.35em}
\raggedright
\centering\textbf{}\par
\raggedright
\vspace{0.2em}
\begin{tabular}{@{}>{\raggedright\arraybackslash}p{0.70\linewidth}@{\hspace{0.55em}}>{\raggedright\arraybackslash}p{0.20\linewidth}@{\hspace{0.25em}}}
\multicolumn{2}{@{}l@{}}{*Old English*} \\
OE Velar Fricative Palatalization & \emph{*çélpaną} \\
OE Heavy Syllable Nasal Apocope & \emph{*çélpan} \\
OE Secondary Nasalization & \emph{*çélpąn} \\
OE Weak Tail Reduction & \emph{*çélpan} \\
\end{tabular}
\end{minipage}
\\
\end{tabularx}
\end{minipage}%
} \endgroup

Old English form: \emph{helpan}

\subsection{Reconstruction and comparative
evidence}\label{reconstruction-and-comparative-evidence-29}

The entry treats the strong verb itself rather than the separate noun
{\emph{help}} `help'. Bright's principal parts \emph{helpan, healp,
hulpon, holpen} show the ordinary Old English strong-verb family
continued by this input (\citeproc{ref-BrightCassidyRingler1971}{Bright
1971}).

\subsection{Old English evidence}\label{old-english-evidence-29}

Clark Hall and Bosworth-Toller record {\emph{helpan}} `help' as the
verbal headword `to help' (\citeproc{ref-ClarkHall1960}{Clark Hall
1960}; \citeproc{ref-BosworthToller1898}{Bosworth and Toller 1898,
542}). The target is therefore the attested infinitive citation form.

\subsection{Development to Old
English}\label{development-to-old-english-29}

From \Recon{xélpaną} `help', no special repair is needed beyond the
ordinary reduction of the infinitive ending. The derivation therefore
reaches {\emph{helpan}} `help' directly.

\subsection{Form note}\label{form-note-8}

Noun {\emph{help}} `help' belongs to a separate lexical line and should
not replace verbal {\emph{helpan}} `help' as the target here
(\citeproc{ref-ClarkHall1960}{Clark Hall 1960};
\citeproc{ref-BosworthToller1898}{Bosworth and Toller 1898, 542}).

\section{\texorpdfstring{hind --- OE
\emph{hind}}{hind --- OE hind}}\label{hind-oe-hind}

\index[iv]{02oe@\textbf{Old English}!hind@\emph{hind}}
\index[iv]{03pgmc@\textbf{Proto-Germanic}!xendjo@*xéndjō}

Derivation: \emph{*xéndjō} \textgreater{} \emph{hind} (regular).

\subsection{Derivation trace}\label{derivation-trace-30}

Proto input: \emph{*xéndjō}

\begingroup
\setlength{\fboxsep}{6pt}

\noindent\fbox{%
\begin{minipage}{0.97\linewidth}
\small
\begin{tabularx}{\linewidth}{@{}>{\raggedright\arraybackslash}p{0.460\linewidth}>{\centering\arraybackslash}X>{\raggedright\arraybackslash}p{0.460\linewidth}@{}}
\begin{minipage}[t]{\linewidth}
\centering\textbf{Earlier Germanic changes}\par
\vspace{0.35em}
\raggedright
\centering\textbf{}\par
\raggedright
\vspace{0.2em}
\begin{tabular}{@{}>{\raggedright\arraybackslash}p{0.74\linewidth}@{\hspace{0.55em}}>{\raggedright\arraybackslash}p{0.16\linewidth}@{\hspace{0.25em}}}
\multicolumn{2}{@{}l@{}}{*Proto-West Germanic*} \\
\multicolumn{2}{@{}l@{}}{[no change]} \\
\multicolumn{2}{@{}l@{}}{*Northwest Germanic*} \\
\mbox{NWGmc Final Long O Raising} & \emph{*xéndju} \\
\end{tabular}
\end{minipage}
&
&
\begin{minipage}[t]{\linewidth}
\centering\textbf{Old English changes}\par
\vspace{0.35em}
\raggedright
\centering\textbf{}\par
\raggedright
\vspace{0.2em}
\begin{tabular}{@{}>{\raggedright\arraybackslash}p{0.74\linewidth}@{\hspace{0.55em}}>{\raggedright\arraybackslash}p{0.16\linewidth}@{\hspace{0.25em}}}
\multicolumn{2}{@{}l@{}}{*Old English*} \\
OE Velar Fricative Palatalization & \emph{*çéndju} \\
OE I Umlaut & \emph{*çindju} \\
\mbox{OE High Vowel Apocope} & \emph{*çindj} \\
OE J Loss After Heavy & \emph{*çind} \\
\end{tabular}
\end{minipage}
\\
\end{tabularx}
\end{minipage}%
} \endgroup

Old English form: \emph{hind}

\subsection{Reconstruction and comparative
evidence}\label{reconstruction-and-comparative-evidence-30}

Kroonen cites the animal name as \emph{*hindō-} f.~`hind'
(\citeproc{ref-Kroonen2013}{Kroonen 2013}). The derivational input
\Recon{xéndjō} `hind' represents that same noun in the present
derivation.

\subsection{Old English evidence}\label{old-english-evidence-30}

Clark Hall and Bosworth-Toller record {\emph{hind}} `hind, doe' as the
noun `hind, female deer' (\citeproc{ref-ClarkHall1960}{Clark Hall 1960};
\citeproc{ref-BosworthToller1898}{Bosworth and Toller 1898, 554}). The
target is therefore the attested lexical item itself.

\subsection{Development to Old
English}\label{development-to-old-english-30}

From \Recon{xéndjō} `hind', i-mutation produces the front-vocalic Old
English stem, and later apocope plus loss of final \emph{j} yield
\emph{hind}. The outcome is therefore a regular citation-form
derivation.

\subsection{Form note}\label{form-note-9}

\emph{hindan} `from behind, behind' is a different Old English lexeme
and does not belong to the noun history of \emph{hind}
(\citeproc{ref-ClarkHall1960}{Clark Hall 1960};
\citeproc{ref-BosworthToller1898}{Bosworth and Toller 1898, 554}).

\section{\texorpdfstring{hold --- OE
\emph{healdan}}{hold --- OE healdan}}\label{hold-oe-healdan}

\index[iv]{02oe@\textbf{Old English}!healdan@\emph{healdan}}
\index[iv]{03pgmc@\textbf{Proto-Germanic}!xaldana@*xáldaną}

Derivation: \emph{*xáldaną} \textgreater{} \emph{healdan} (regular).

\subsection{Derivation trace}\label{derivation-trace-31}

Proto input: \emph{*xáldaną}

\begingroup
\setlength{\fboxsep}{6pt}

\noindent\fbox{%
\begin{minipage}{0.97\linewidth}
\footnotesize
\begin{tabularx}{\linewidth}{@{}>{\raggedright\arraybackslash}p{0.320\linewidth}>{\centering\arraybackslash}X>{\raggedright\arraybackslash}p{0.520\linewidth}@{}}
\begin{minipage}[t]{\linewidth}
\centering\textbf{Earlier Germanic changes}\par
\vspace{0.35em}
\raggedright
\centering\textbf{}\par
\raggedright
\vspace{0.2em}
\begin{tabular}{@{}>{\raggedright\arraybackslash}p{0.68\linewidth}@{\hspace{0.55em}}>{\raggedright\arraybackslash}p{0.16\linewidth}@{\hspace{0.25em}}}
\multicolumn{2}{@{}l@{}}{*Proto-West Germanic*} \\
\multicolumn{2}{@{}l@{}}{[no change]} \\
\multicolumn{2}{@{}l@{}}{*Northwest Germanic*} \\
\multicolumn{2}{@{}l@{}}{[no change]} \\
\end{tabular}
\end{minipage}
&
&
\begin{minipage}[t]{\linewidth}
\centering\textbf{Old English changes}\par
\vspace{0.35em}
\raggedright
\centering\textbf{}\par
\raggedright
\vspace{0.2em}
\begin{tabular}{@{}>{\raggedright\arraybackslash}p{0.66\linewidth}@{\hspace{0.55em}}>{\raggedright\arraybackslash}p{0.24\linewidth}@{\hspace{0.25em}}}
\multicolumn{2}{@{}l@{}}{*Old English*} \\
Anglo Frisian Brightening & \emph{*xældaną} \\
OE Breaking & \emph{*xealdaną} \\
OE Velar Fricative Palatalization & \emph{*çealdaną} \\
OE Heavy Syllable Nasal Apocope & \emph{*çealdan} \\
OE Secondary Nasalization & \emph{*çealdąn} \\
OE Weak Tail Reduction & \emph{*çealdan} \\
\end{tabular}
\end{minipage}
\\
\end{tabularx}
\end{minipage}%
} \endgroup

Old English form: \emph{healdan}

\subsection{Reconstruction and comparative
evidence}\label{reconstruction-and-comparative-evidence-31}

Campbell and Ringe-Taylor treat the verb as a regular \emph{a + lC}
breaking case, with West Saxon {\emph{healdan}} `hold' opposed to
Anglian and Mercian {\emph{haldan}} `hold'
(\citeproc{ref-Campbell1959}{Campbell 1959};
\citeproc{ref-RingeTaylor2014}{Ringe and Taylor 2014}).

\subsection{Old English evidence}\label{old-english-evidence-31}

Bright gives the ordinary strong-verb citation form and principal parts
\emph{healdan, heold, heoldon, healden}
(\citeproc{ref-BrightCassidyRingler1971}{Bright 1971}). The target is
therefore the attested infinitive headword itself.

\subsection{Development to Old
English}\label{development-to-old-english-31}

From \Recon{xáldaną} `hold', Anglo-Frisian brightening first yields a
fronted vowel, and West Saxon breaking then produces \emph{ea} before
\emph{ld}. Later reduction of the infinitive ending gives
{\emph{healdan}} `hold' (\citeproc{ref-Campbell1959}{Campbell 1959};
\citeproc{ref-RingeTaylor2014}{Ringe and Taylor 2014}).

\subsection{Dialect note}\label{dialect-note-2}

West Saxon {\emph{healdan}} `hold' is the Old English form here. Anglian
and Mercian {\emph{haldan}} `hold' are genuine non-West-Saxon doublets
rather than corrections to that choice
(\citeproc{ref-Campbell1959}{Campbell 1959};
\citeproc{ref-RingeTaylor2014}{Ringe and Taylor 2014}).

\section{\texorpdfstring{horn --- OE
\emph{horn}}{horn --- OE horn}}\label{horn-oe-horn}

\index[iv]{02oe@\textbf{Old English}!horn@\emph{horn}}
\index[iv]{03pgmc@\textbf{Proto-Germanic}!xurna@*xúrną}

Derivation: \emph{*xúrną} \textgreater{} \emph{horn} (regular).

\subsection{Derivation trace}\label{derivation-trace-32}

Proto input: \emph{*xúrną}

\begingroup
\setlength{\fboxsep}{6pt}

\noindent\fbox{%
\begin{minipage}{0.97\linewidth}
\small
\begin{tabularx}{\linewidth}{@{}>{\raggedright\arraybackslash}p{0.460\linewidth}>{\centering\arraybackslash}X>{\raggedright\arraybackslash}p{0.460\linewidth}@{}}
\begin{minipage}[t]{\linewidth}
\centering\textbf{Earlier Germanic changes}\par
\vspace{0.35em}
\raggedright
\centering\textbf{}\par
\raggedright
\vspace{0.2em}
\begin{tabular}{@{}>{\raggedright\arraybackslash}p{0.74\linewidth}@{\hspace{0.55em}}>{\raggedright\arraybackslash}p{0.16\linewidth}@{\hspace{0.25em}}}
\multicolumn{2}{@{}l@{}}{*Proto-West Germanic*} \\
\multicolumn{2}{@{}l@{}}{[no change]} \\
\multicolumn{2}{@{}l@{}}{*Northwest Germanic*} \\
\mbox{NWGmc U Lowering} & \emph{*xórną} \\
\end{tabular}
\end{minipage}
&
&
\begin{minipage}[t]{\linewidth}
\centering\textbf{Old English changes}\par
\vspace{0.35em}
\raggedright
\centering\textbf{}\par
\raggedright
\vspace{0.2em}
\begin{tabular}{@{}>{\raggedright\arraybackslash}p{0.74\linewidth}@{\hspace{0.55em}}>{\raggedright\arraybackslash}p{0.16\linewidth}@{\hspace{0.25em}}}
\multicolumn{2}{@{}l@{}}{*Old English*} \\
OE Heavy Syllable Nasal Apocope & \emph{*xórn} \\
\end{tabular}
\end{minipage}
\\
\end{tabularx}
\end{minipage}%
} \endgroup

Old English form: \emph{horn}

\subsection{Reconstruction and comparative
evidence}\label{reconstruction-and-comparative-evidence-32}

Kroonen and Orel cite lemma-style Proto-Germanic headwords of the
\emph{*hurna-} / \Recon{xurnan} `horn' type for this noun
(\citeproc{ref-Kroonen2013}{Kroonen 2013}; \citeproc{ref-Orel2003}{Orel
2003, 234}). The derivational input \Recon{xúrną} `horn' is the
nominative-style form used in the derivation here.

\subsection{Old English evidence}\label{old-english-evidence-32}

Clark Hall, Bosworth-Toller, and Bright all record \emph{horn} as the
ordinary Old English noun (\citeproc{ref-ClarkHall1960}{Clark Hall
1960}; \citeproc{ref-BosworthToller1898}{Bosworth and Toller 1898, 108};
\citeproc{ref-BrightCassidyRingler1971}{Bright 1971}). The target is
therefore the attested citation form.

\subsection{Development to Old
English}\label{development-to-old-english-32}

From \Recon{xúrną} `horn', Northwest Germanic u-lowering gives
\Recon{xórną} `horn', and later loss of the final nasal leaves
\emph{horn}. The development treated here is fully regular.

\subsection{Form note}\label{form-note-10}

The note's oblique \Recon{xurnăn} `horn' belongs to comparative stem
background only. It does not replace the derivational input
\Recon{xúrną} `horn' as the derivational form used here
(\citeproc{ref-Kroonen2013}{Kroonen 2013}; \citeproc{ref-Orel2003}{Orel
2003, 234}).

\section{\texorpdfstring{lead --- OE
\emph{lǣdan}}{lead --- OE lǣdan}}\label{lead-oe-lux1e3dan}

\index[iv]{02oe@\textbf{Old English}!laedan@\emph{lǣdan}}
\index[iv]{03pgmc@\textbf{Proto-Germanic}!laidijana@*láidijaną}

Derivation: \emph{*láidijaną} \textgreater{} \emph{lǣdan} (regular).

\subsection{Derivation trace}\label{derivation-trace-33}

Proto input: \emph{*láidijaną}

\begingroup
\setlength{\fboxsep}{6pt}

\noindent\fbox{%
\begin{minipage}{0.97\linewidth}
\small
\begin{tabularx}{\linewidth}{@{}>{\raggedright\arraybackslash}p{0.440\linewidth}>{\centering\arraybackslash}X>{\raggedright\arraybackslash}p{0.440\linewidth}@{}}
\begin{minipage}[t]{\linewidth}
\centering\textbf{Earlier Germanic changes}\par
\vspace{0.35em}
\raggedright
\centering\textbf{}\par
\raggedright
\vspace{0.2em}
\begin{tabular}{@{}>{\raggedright\arraybackslash}p{0.66\linewidth}@{\hspace{0.55em}}>{\raggedright\arraybackslash}p{0.22\linewidth}@{\hspace{0.25em}}}
\multicolumn{2}{@{}l@{}}{*Proto-West Germanic*} \\
PWGmc Ai Monophthongization & \emph{*lādijaną} \\
\multicolumn{2}{@{}l@{}}{*Northwest Germanic*} \\
\multicolumn{2}{@{}l@{}}{[no change]} \\
\end{tabular}
\end{minipage}
&
&
\begin{minipage}[t]{\linewidth}
\centering\textbf{Old English changes}\par
\vspace{0.35em}
\raggedright
\centering\textbf{}\par
\raggedright
\vspace{0.2em}
\begin{tabular}{@{}>{\raggedright\arraybackslash}p{0.70\linewidth}@{\hspace{0.55em}}>{\raggedright\arraybackslash}p{0.20\linewidth}@{\hspace{0.25em}}}
\multicolumn{2}{@{}l@{}}{*Old English*} \\
OE Heavy Syllable Nasal Apocope & \emph{*lādijan} \\
OE Secondary Nasalization & \emph{*lādijąn} \\
Sievers Law Syncope & \emph{*lādjąn} \\
OE I Umlaut & \emph{*lǣdjąn} \\
OE Weak Tail Reduction & \emph{*lǣdjan} \\
OE J Loss After Heavy & \emph{*lǣdan} \\
\end{tabular}
\end{minipage}
\\
\end{tabularx}
\end{minipage}%
} \endgroup

Old English form: \emph{lǣdan}

\subsection{Reconstruction and comparative
evidence}\label{reconstruction-and-comparative-evidence-33}

Ringe and Taylor derive Old English \emph{lǣdan} `lead' from
Proto-Germanic \Recon{laidijaną} `lead', and Kroonen likewise cites a
weak verb of the \emph{*laidjan-} type for `lead'
(\citeproc{ref-RingeTaylor2014}{Ringe and Taylor 2014};
\citeproc{ref-Kroonen2013}{Kroonen 2013, 363}).

\subsection{Old English evidence}\label{old-english-evidence-33}

Clark Hall and Bosworth-Toller both record \emph{lædan} `lead' /
\emph{lǣdan} `lead' as the ordinary Old English verb `to lead, guide,
conduct' (\citeproc{ref-ClarkHall1960}{Clark Hall 1960};
\citeproc{ref-BosworthToller1898}{Bosworth and Toller 1898}).

\subsection{Development to Old
English}\label{development-to-old-english-33}

From \Recon{láidijaną} `lead', monophthongization of \emph{*ai} first
gives a \emph{*lād-} stage. Later syncope, i-mutation, weak-tail
reduction, and loss of \emph{j} after a heavy stem yield \emph{lǣdan}
`lead', so the development represented here is fully regular
(\citeproc{ref-RingeTaylor2014}{Ringe and Taylor 2014}).

\section{\texorpdfstring{learn --- OE
\emph{liornian}}{learn --- OE liornian}}\label{learn-oe-liornian}

\index[iv]{02oe@\textbf{Old English}!liornian@\emph{liornian}}
\index[iv]{03pgmc@\textbf{Proto-Germanic}!liznojana@*líznōjaną}

Derivation: \emph{*líznōjaną} \textgreater{} \emph{liornian} (regular).

\subsection{Derivation trace}\label{derivation-trace-34}

Proto input: \emph{*líznōjaną}

\begingroup
\setlength{\fboxsep}{6pt}

\noindent\fbox{%
\begin{minipage}{0.97\linewidth}
\footnotesize
\begin{tabularx}{\linewidth}{@{}>{\raggedright\arraybackslash}p{0.320\linewidth}>{\centering\arraybackslash}X>{\raggedright\arraybackslash}p{0.520\linewidth}@{}}
\begin{minipage}[t]{\linewidth}
\centering\textbf{Earlier Germanic changes}\par
\vspace{0.35em}
\raggedright
\centering\textbf{}\par
\raggedright
\vspace{0.2em}
\begin{tabular}{@{}>{\raggedright\arraybackslash}p{0.68\linewidth}@{\hspace{0.55em}}>{\raggedright\arraybackslash}p{0.16\linewidth}@{\hspace{0.25em}}}
\multicolumn{2}{@{}l@{}}{*Proto-West Germanic*} \\
\multicolumn{2}{@{}l@{}}{[no change]} \\
\multicolumn{2}{@{}l@{}}{*Northwest Germanic*} \\
\multicolumn{2}{@{}l@{}}{[no change]} \\
\end{tabular}
\end{minipage}
&
&
\begin{minipage}[t]{\linewidth}
\centering\textbf{Old English changes}\par
\vspace{0.35em}
\raggedright
\centering\textbf{}\par
\raggedright
\vspace{0.2em}
\begin{tabular}{@{}>{\raggedright\arraybackslash}p{0.62\linewidth}@{\hspace{0.55em}}>{\raggedright\arraybackslash}p{0.28\linewidth}@{\hspace{0.25em}}}
\multicolumn{2}{@{}l@{}}{*Old English*} \\
OE Breaking & \emph{*líornōjaną} \\
OE Heavy Syllable Nasal Apocope & \emph{*líornōjan} \\
OE Secondary Nasalization & \emph{*líornōjąn} \\
OE I Umlaut & \emph{*líornējąn} \\
OE Unstressed Long Vowel Shortening & \emph{*líornejąn} \\
OE Weak Tail Reduction & \emph{*líornejan} \\
OE Intervocalic J Vocalization & \emph{*líorneian} \\
OE Unstressed EI Contraction & \emph{*líornian} \\
\end{tabular}
\end{minipage}
\\
\end{tabularx}
\end{minipage}%
} \endgroup

Old English form: \emph{liornian}

\subsection{Reconstruction and comparative
evidence}\label{reconstruction-and-comparative-evidence-34}

Ringe and Taylor place the verb in a class-II weak family of the
\emph{*liznō-} type (\citeproc{ref-RingeTaylor2014}{Ringe and Taylor
2014}), and Kroonen keeps the same comparative base for the Germanic
lexeme (\citeproc{ref-Kroonen2013}{Kroonen 2013}). The derivational
input therefore requires no change of stem class or paradigm cell.

The non-obvious issue is dialectal. Campbell records Northumbrian
\emph{liornian} `learn' beside \emph{leornian} `learn', and he states
that the \emph{eo} of \emph{leornian} `learn' is secondary, while
Northumbrian preserves forms with \emph{io}, where original \emph{eo}
and \emph{io} remain distinct (\citeproc{ref-Campbell1959}{Campbell
1959, sect. 123} n.~2).

\subsection{Old English evidence}\label{old-english-evidence-34}

\Recon{liornian} `learn' is the Northumbrian member of the Old English
family. Dictionary practice more often privileges \emph{leornian}
`learn' as the ordinary headword (\citeproc{ref-ClarkHall1960}{Clark
Hall 1960}; \citeproc{ref-BrightCassidyRingler1971}{Bright 1971}), but
Campbell's dialect evidence shows that \emph{liornian} `learn' is a
genuine Old English form (\citeproc{ref-Campbell1959}{Campbell 1959,
sect. 123} n.~2).

This entry therefore remains compact. The point is to state clearly that
the Old English form here belongs to the Northumbrian side of the OE
evidence rather than to the leveled \emph{leornian} `learn' headword
tradition.

\subsection{Development to Old
English}\label{development-to-old-english-34}

From \Recon{líznōjaną} `learn', the expected Old English developments
include rhotacism of \emph{z}, followed by breaking before \emph{r} plus
consonant, and the ordinary reduction of the weak verbal ending. The
result is \emph{liornian} `learn', preserving the \emph{io} spelling
that Campbell associates with Northumbrian where original \emph{eo} and
\emph{io} remain distinct (\citeproc{ref-Campbell1959}{Campbell 1959,
sect. 123} n.~2).

\emph{Leornian} reflects the later \emph{eo} development that Campbell
treats as secondary {[}Campbell (\citeproc{ref-Campbell1959}{1959}),
§123 n.~2; §296{]}. The selected Northumbrian target is therefore the
regular comparison form for the \emph{i}-grade member of the family.

\subsection{Form comparison}\label{form-comparison-1}

The comparison below sets the relevant forms side by side. It
distinguishes the regular Northumbrian form modeled here from the
better-known West Saxon headword.

{\def\LTcaptype{none} % do not increment counter
\begin{longtable}[]{@{}
  >{\raggedright\arraybackslash}p{(\linewidth - 4\tabcolsep) * \real{0.3333}}
  >{\raggedright\arraybackslash}p{(\linewidth - 4\tabcolsep) * \real{0.3333}}
  >{\raggedright\arraybackslash}p{(\linewidth - 4\tabcolsep) * \real{0.3333}}@{}}
\toprule\noalign{}
\begin{minipage}[b]{\linewidth}\raggedright
Form
\end{minipage} & \begin{minipage}[b]{\linewidth}\raggedright
Status
\end{minipage} & \begin{minipage}[b]{\linewidth}\raggedright
Relevance to this entry
\end{minipage} \\
\midrule\noalign{}
\endhead
\bottomrule\noalign{}
\endlastfoot
\emph{*líznōjaną} \textgreater{} \emph{liornian} & computed regular
output; attested Northumbrian comparison form & selected comparison \\
\emph{leornian} & attested later \emph{eo} form and dictionary headword
& useful control, but not the target of this entry \\
\end{longtable}
}

\section{\texorpdfstring{lid --- OE
\emph{hlid}}{lid --- OE hlid}}\label{lid-oe-hlid}

\index[iv]{02oe@\textbf{Old English}!hlid@\emph{hlid}}
\index[iv]{03pgmc@\textbf{Proto-Germanic}!xlida@*xlídą}
\index[iv]{08on@\textbf{Old Norse}!hlitho@hliþó}

Derivation: \emph{*xlídą} \textgreater{} \emph{hlid} (regular).

\subsection{Derivation trace}\label{derivation-trace-35}

Proto input: \emph{*xlídą}

\begingroup
\setlength{\fboxsep}{6pt}

\noindent\fbox{%
\begin{minipage}{0.97\linewidth}
\small
\begin{tabularx}{\linewidth}{@{}>{\raggedright\arraybackslash}p{0.440\linewidth}>{\centering\arraybackslash}X>{\raggedright\arraybackslash}p{0.440\linewidth}@{}}
\begin{minipage}[t]{\linewidth}
\centering\textbf{Earlier Germanic changes}\par
\vspace{0.35em}
\raggedright
\centering\textbf{}\par
\raggedright
\vspace{0.2em}
\begin{tabular}{@{}>{\raggedright\arraybackslash}p{0.68\linewidth}@{\hspace{0.55em}}>{\raggedright\arraybackslash}p{0.16\linewidth}@{\hspace{0.25em}}}
\multicolumn{2}{@{}l@{}}{*Proto-West Germanic*} \\
\multicolumn{2}{@{}l@{}}{[no change]} \\
\multicolumn{2}{@{}l@{}}{*Northwest Germanic*} \\
\multicolumn{2}{@{}l@{}}{[no change]} \\
\end{tabular}
\end{minipage}
&
&
\begin{minipage}[t]{\linewidth}
\centering\textbf{Old English changes}\par
\vspace{0.35em}
\raggedright
\centering\textbf{}\par
\raggedright
\vspace{0.2em}
\begin{tabular}{@{}>{\raggedright\arraybackslash}p{0.74\linewidth}@{\hspace{0.55em}}>{\raggedright\arraybackslash}p{0.16\linewidth}@{\hspace{0.25em}}}
\multicolumn{2}{@{}l@{}}{*Old English*} \\
OE Heavy Syllable Nasal Apocope & \emph{*xlíd} \\
\end{tabular}
\end{minipage}
\\
\end{tabularx}
\end{minipage}%
} \endgroup

Old English form: \emph{hlid}

\subsection{Reconstruction and comparative
evidence}\label{reconstruction-and-comparative-evidence-35}

Orel cites a neuter lexeme of the \emph{*xliđ-} type with Old English
{\emph{hlid}} `lid, cover', and Lloyd includes OE \emph{hlid} `lid'
beside ON \emph{hliþó} `lid' and OHG \emph{(h)lit} `lid' among forms
that retain \emph{i} (\citeproc{ref-Orel2003}{Orel 2003};
\citeproc{ref-Lloyd1966}{Lloyd 1966}).

\subsection{Old English evidence}\label{old-english-evidence-35}

Clark Hall and Bosworth-Toller record {\emph{hlid}} `lid, cover' as the
noun `lid, cover, door, gate' (\citeproc{ref-ClarkHall1960}{Clark Hall
1960}; \citeproc{ref-BosworthToller1898}{Bosworth and Toller 1898,
563}).

\subsection{Development to Old
English}\label{development-to-old-english-35}

The derivational input already represents the later Germanic
\emph{hliđ-} stage used for the derivation here. From \Recon{xlídą}
`lid', heavy-syllable apocope yields \emph{hlid} `lid', and the form
belongs to the retained-\emph{i} set noted by Lloyd rather than to the
lowered \emph{e} type (\citeproc{ref-Lloyd1966}{Lloyd 1966}).

\subsection{Form note}\label{form-note-11}

An earlier etymological stage \Recon{liþuz} `lid' belongs to comparative
background only. The form represented here is the later \Recon{xlídą}
`lid' \textgreater{} \emph{hlid} `lid' line that matches the attested
Old English noun (\citeproc{ref-Orel2003}{Orel 2003};
\citeproc{ref-Lloyd1966}{Lloyd 1966}).

\section{\texorpdfstring{light --- OE
\emph{līehtan}}{light --- OE līehtan}}\label{light-oe-lux12behtan}

\index[iv]{02oe@\textbf{Old English}!liehtan@\emph{līehtan}}
\index[iv]{03pgmc@\textbf{Proto-Germanic}!leuxtijana@*léuxtijaną}

Derivation: \emph{*léuxtijaną} \textgreater{} \emph{līehtan} (regular).

\subsection{Derivation trace}\label{derivation-trace-36}

Proto input: \emph{*léuxtijaną}

\begingroup
\setlength{\fboxsep}{6pt}

\noindent\fbox{%
\begin{minipage}{0.97\linewidth}
\footnotesize
\begin{tabularx}{\linewidth}{@{}>{\raggedright\arraybackslash}p{0.320\linewidth}>{\centering\arraybackslash}X>{\raggedright\arraybackslash}p{0.520\linewidth}@{}}
\begin{minipage}[t]{\linewidth}
\centering\textbf{Earlier Germanic changes}\par
\vspace{0.35em}
\raggedright
\centering\textbf{}\par
\raggedright
\vspace{0.2em}
\begin{tabular}{@{}>{\raggedright\arraybackslash}p{0.68\linewidth}@{\hspace{0.55em}}>{\raggedright\arraybackslash}p{0.16\linewidth}@{\hspace{0.25em}}}
\multicolumn{2}{@{}l@{}}{*Proto-West Germanic*} \\
\multicolumn{2}{@{}l@{}}{[no change]} \\
\multicolumn{2}{@{}l@{}}{*Northwest Germanic*} \\
\multicolumn{2}{@{}l@{}}{[no change]} \\
\end{tabular}
\end{minipage}
&
&
\begin{minipage}[t]{\linewidth}
\centering\textbf{Old English changes}\par
\vspace{0.35em}
\raggedright
\centering\textbf{}\par
\raggedright
\vspace{0.2em}
\begin{tabular}{@{}>{\raggedright\arraybackslash}p{0.62\linewidth}@{\hspace{0.55em}}>{\raggedright\arraybackslash}p{0.28\linewidth}@{\hspace{0.25em}}}
\multicolumn{2}{@{}l@{}}{*Old English*} \\
OE Diphthong Leveling & \emph{*lēoxtijaną} \\
OE Heavy Syllable Nasal Apocope & \emph{*lēoxtijan} \\
OE Secondary Nasalization & \emph{*lēoxtijąn} \\
Sievers Law Syncope & \emph{*lēoxtjąn} \\
OE I Umlaut & \emph{*līextjąn} \\
OE Weak Tail Reduction & \emph{*līextjan} \\
OE J Loss After Heavy & \emph{*līextan} \\
\end{tabular}
\end{minipage}
\\
\end{tabularx}
\end{minipage}%
} \endgroup

Old English form: \emph{līehtan}

\subsection{Reconstruction and comparative
evidence}\label{reconstruction-and-comparative-evidence-36}

Fulk gives Proto-Germanic \Recon{liuxtijanan} `light' with Old English
\emph{līehtan} `illuminate', and Ringe and Taylor likewise derive West
Saxon {\emph{liehtan}} `light, illuminate' from the same weak-verb
formation (\citeproc{ref-Fulk2018}{Fulk 2018};
\citeproc{ref-RingeTaylor2014}{Ringe and Taylor 2014}).

\subsection{Old English evidence}\label{old-english-evidence-36}

Clark Hall and Bosworth-Toller preserve the verb family under spellings
such as \emph{liehtan} `illuminate', {\emph{lihtan}} `light,
illuminate', and \emph{līhtan} `illuminate', distinct from the related
noun \emph{lēoht} `light' and adjective \emph{leoht} / \emph{liht}
`light' (\citeproc{ref-ClarkHall1960}{Clark Hall 1960};
\citeproc{ref-BosworthToller1898}{Bosworth and Toller 1898}).

\subsection{Development to Old
English}\label{development-to-old-english-36}

From \Recon{léuxtijaną} `light', the regular verbal line preserves
\emph{*xt}, passes through a West Saxon \emph{liehtan} `illuminate'
stage, and is represented here by normalized \emph{līehtan}
`illuminate'. The word treated in this entry is therefore the verb `to
light, illuminate', not the related noun from \Recon{leuxtą} `light'
(\citeproc{ref-Fulk2018}{Fulk 2018};
\citeproc{ref-RingeTaylor2014}{Ringe and Taylor 2014}).

\subsection{Dialect note}\label{dialect-note-3}

Ringe and Taylor and Campbell distinguish West Saxon {\emph{liehtan}}
`light, illuminate' from Anglian {\emph{lihtan}} `light, illuminate',
while later West Saxon also shows {\emph{lyhtan}} `light, illuminate'
(\citeproc{ref-RingeTaylor2014}{Ringe and Taylor 2014};
\citeproc{ref-Campbell1959}{Campbell 1959}).

\section{\texorpdfstring{linden --- OE
\emph{lind}}{linden --- OE lind}}\label{linden-oe-lind}

\index[iv]{02oe@\textbf{Old English}!lind@\emph{lind}}
\index[iv]{03pgmc@\textbf{Proto-Germanic}!lindo@*líndō}

Derivation: \emph{*líndō} \textgreater{} \emph{lind} (regular).

\subsection{Derivation trace}\label{derivation-trace-37}

Proto input: \emph{*líndō}

\begingroup
\setlength{\fboxsep}{6pt}

\noindent\fbox{%
\begin{minipage}{0.97\linewidth}
\small
\begin{tabularx}{\linewidth}{@{}>{\raggedright\arraybackslash}p{0.460\linewidth}>{\centering\arraybackslash}X>{\raggedright\arraybackslash}p{0.460\linewidth}@{}}
\begin{minipage}[t]{\linewidth}
\centering\textbf{Earlier Germanic changes}\par
\vspace{0.35em}
\raggedright
\centering\textbf{}\par
\raggedright
\vspace{0.2em}
\begin{tabular}{@{}>{\raggedright\arraybackslash}p{0.74\linewidth}@{\hspace{0.55em}}>{\raggedright\arraybackslash}p{0.16\linewidth}@{\hspace{0.25em}}}
\multicolumn{2}{@{}l@{}}{*Proto-West Germanic*} \\
\multicolumn{2}{@{}l@{}}{[no change]} \\
\multicolumn{2}{@{}l@{}}{*Northwest Germanic*} \\
\mbox{NWGmc Final Long O Raising} & \emph{*líndu} \\
\end{tabular}
\end{minipage}
&
&
\begin{minipage}[t]{\linewidth}
\centering\textbf{Old English changes}\par
\vspace{0.35em}
\raggedright
\centering\textbf{}\par
\raggedright
\vspace{0.2em}
\begin{tabular}{@{}>{\raggedright\arraybackslash}p{0.74\linewidth}@{\hspace{0.55em}}>{\raggedright\arraybackslash}p{0.16\linewidth}@{\hspace{0.25em}}}
\multicolumn{2}{@{}l@{}}{*Old English*} \\
\mbox{OE High Vowel Apocope} & \emph{*línd} \\
\end{tabular}
\end{minipage}
\\
\end{tabularx}
\end{minipage}%
} \endgroup

Old English form: \emph{lind}

\subsection{Reconstruction and comparative
evidence}\label{reconstruction-and-comparative-evidence-37}

Kroonen cites \emph{*lindō-} `lime tree' and gives Old English
{\emph{lind}} `linden, lime tree' as the relevant reflex
(\citeproc{ref-Kroonen2013}{Kroonen 2013}).

\subsection{Old English evidence}\label{old-english-evidence-37}

Clark Hall and Bosworth-Toller both record {\emph{lind}} `linden, lime
tree' as the noun `lime-tree, linden'
(\citeproc{ref-ClarkHall1960}{Clark Hall 1960};
\citeproc{ref-BosworthToller1898}{Bosworth and Toller 1898, 630}).

\subsection{Development to Old
English}\label{development-to-old-english-37}

From \Recon{líndō} `linden', Northwest Germanic final \emph{*ō} raising
first gives a \Recon{líndu} `linden' stage, and later high-vowel apocope
yields {\emph{lind}} `linden, lime tree'. The development is therefore
straightforward and regular.

\subsection{Form note}\label{form-note-12}

The Old English noun represented here is {\emph{lind}} `linden, lime
tree'. Clark Hall also has a separate adjectival \emph{linden} `made of
linden-wood', but that is not the noun counterpart for this entry
(\citeproc{ref-ClarkHall1960}{Clark Hall 1960}).

\section{\texorpdfstring{milk --- OE
\emph{meoloc}}{milk --- OE meoloc}}\label{milk-oe-meoloc}

\index[iv]{02oe@\textbf{Old English}!meoloc@\emph{meoloc}}
\index[iv]{03pgmc@\textbf{Proto-Germanic}!melukz@*mélukz}

Derivation: \emph{*mélukz} \textgreater{} \emph{meoloc} (regular).

\subsection{Derivation trace}\label{derivation-trace-38}

Proto input: \emph{*mélukz}

\begingroup
\setlength{\fboxsep}{6pt}

\noindent\fbox{%
\begin{minipage}{0.97\linewidth}
\small
\begin{tabularx}{\linewidth}{@{}>{\raggedright\arraybackslash}p{0.440\linewidth}>{\centering\arraybackslash}X>{\raggedright\arraybackslash}p{0.440\linewidth}@{}}
\begin{minipage}[t]{\linewidth}
\centering\textbf{Earlier Germanic changes}\par
\vspace{0.35em}
\raggedright
\centering\textbf{}\par
\raggedright
\vspace{0.2em}
\begin{tabular}{@{}>{\raggedright\arraybackslash}p{0.74\linewidth}@{\hspace{0.55em}}>{\raggedright\arraybackslash}p{0.16\linewidth}@{\hspace{0.25em}}}
\multicolumn{2}{@{}l@{}}{*Proto-West Germanic*} \\
\multicolumn{2}{@{}l@{}}{[no change]} \\
\multicolumn{2}{@{}l@{}}{*Northwest Germanic*} \\
\mbox{PGmc Final Z Deletion} & \emph{*méluk} \\
\end{tabular}
\end{minipage}
&
&
\begin{minipage}[t]{\linewidth}
\centering\textbf{Old English changes}\par
\vspace{0.35em}
\raggedright
\centering\textbf{}\par
\raggedright
\vspace{0.2em}
\begin{tabular}{@{}>{\raggedright\arraybackslash}p{0.74\linewidth}@{\hspace{0.55em}}>{\raggedright\arraybackslash}p{0.16\linewidth}@{\hspace{0.25em}}}
\multicolumn{2}{@{}l@{}}{*Old English*} \\
OE Med Unstressed U Lowering & \emph{*mélok} \\
OE Back Mutation & \emph{*méolok} \\
\end{tabular}
\end{minipage}
\\
\end{tabularx}
\end{minipage}%
} \endgroup

Old English form: \emph{meoloc}

\subsection{Reconstruction and comparative
evidence}\label{reconstruction-and-comparative-evidence-38}

Kroonen and Orel reconstruct the noun with stem notation \emph{*meluk-}
and nominative \Recon{melukz} `milk', and the nominative-style input
used here is \Recon{mélukz} `milk' (\citeproc{ref-Kroonen2013}{Kroonen
2013}; \citeproc{ref-Orel2003}{Orel 2003, 306}).

\subsection{Old English evidence}\label{old-english-evidence-38}

Old English preserves a mixed dossier for this noun. Ringe and Taylor
describe West Saxon \emph{meolc} `milk' \textless{} \emph{meoluc} `milk'
\textless{} \Recon{mélukz} `milk', Campbell likewise discusses
\emph{meoluc} `milk' and {\emph{meoloc}} `milk', and Anglian shows
\emph{milc} `milk' (\citeproc{ref-RingeTaylor2014}{Ringe and Taylor
2014}; \citeproc{ref-Campbell1959}{Campbell 1959}).

\subsection{Development to Old
English}\label{development-to-old-english-38}

The unsyncopated line from \Recon{mélukz} `milk' loses final \emph{*z},
lowers unstressed \emph{u} to \emph{o}, and with back mutation yields
{\emph{meoloc}} `milk'. That fuller unsyncopated outcome is the form
represented here (\citeproc{ref-RingeTaylor2014}{Ringe and Taylor
2014}).

\subsection{Form comparison}\label{form-comparison-2}

Syncopated \emph{meolc} and Anglian \emph{milc} belong to the competing
leveled tradition associated with oblique forms, whereas \emph{meoloc} /
\emph{meoluc} preserves the fuller nominal shape
(\citeproc{ref-Campbell1959}{Campbell 1959};
\citeproc{ref-RingeTaylor2014}{Ringe and Taylor 2014}).

\section{\texorpdfstring{mother --- OE
\emph{mōder}}{mother --- OE mōder}}\label{mother-oe-mux14dder}

\index[iv]{02oe@\textbf{Old English}!moder@\emph{mōder}}
\index[iv]{03pgmc@\textbf{Proto-Germanic}!moder@*mōdēr}

Derivation: \emph{*mōdēr} \textgreater{} \emph{mōder} (regular).

\subsection{Derivation trace}\label{derivation-trace-39}

Proto input: \emph{*mōdēr}

\begingroup
\setlength{\fboxsep}{6pt}

\noindent\fbox{%
\begin{minipage}{0.97\linewidth}
\small
\begin{tabularx}{\linewidth}{@{}>{\raggedright\arraybackslash}p{0.460\linewidth}>{\centering\arraybackslash}X>{\raggedright\arraybackslash}p{0.460\linewidth}@{}}
\begin{minipage}[t]{\linewidth}
\centering\textbf{Earlier Germanic changes}\par
\vspace{0.35em}
\raggedright
\centering\textbf{}\par
\raggedright
\vspace{0.2em}
\begin{tabular}{@{}>{\raggedright\arraybackslash}p{0.74\linewidth}@{\hspace{0.55em}}>{\raggedright\arraybackslash}p{0.16\linewidth}@{\hspace{0.25em}}}
\multicolumn{2}{@{}l@{}}{*Proto-West Germanic*} \\
\multicolumn{2}{@{}l@{}}{[no change]} \\
\multicolumn{2}{@{}l@{}}{*Northwest Germanic*} \\
\mbox{NWGmc Long E Lowering} & \emph{*mōdǣr} \\
\end{tabular}
\end{minipage}
&
&
\begin{minipage}[t]{\linewidth}
\centering\textbf{Old English changes}\par
\vspace{0.35em}
\raggedright
\centering\textbf{}\par
\raggedright
\vspace{0.2em}
\begin{tabular}{@{}>{\raggedright\arraybackslash}p{0.74\linewidth}@{\hspace{0.55em}}>{\raggedright\arraybackslash}p{0.16\linewidth}@{\hspace{0.25em}}}
\multicolumn{2}{@{}l@{}}{*Old English*} \\
OE Unstressed Long Vowel Shortening & \emph{*mōdær} \\
OE Unstressed AE Merger & \emph{*mōder} \\
\end{tabular}
\end{minipage}
\\
\end{tabularx}
\end{minipage}%
} \endgroup

Old English form: \emph{mōder}

\subsection{Reconstruction and comparative
evidence}\label{reconstruction-and-comparative-evidence-39}

Kroonen and Orel cite the Proto-Germanic r-stem kinship noun as
\emph{*mōder-} with citation \Recon{mōdēr} `mother'
(\citeproc{ref-Kroonen2013}{Kroonen 2013}; \citeproc{ref-Orel2003}{Orel
2003}).

\subsection{Old English evidence}\label{old-english-evidence-39}

The transmitted Old English headword tradition is \emph{mōdor} `mother'
/ \emph{modor} `mother', with oblique \emph{mēder} `mother' in the
paradigm. Clark Hall, Campbell, and Ringe and Taylor all preserve that
contrast (\citeproc{ref-ClarkHall1960}{Clark Hall 1960};
\citeproc{ref-Campbell1959}{Campbell 1959};
\citeproc{ref-RingeTaylor2014}{Ringe and Taylor 2014}).

\subsection{Development to Old
English}\label{development-to-old-english-39}

From \Recon{mōdēr} `mother', the regular suffixal development yields
\emph{mōder} `mother'. That regular nominative reflex is the form
represented here, while the more familiar citation form \emph{mōdor}
`mother' reflects later levelling within the r-stem paradigm
(\citeproc{ref-Campbell1959}{Campbell 1959};
\citeproc{ref-RingeTaylor2014}{Ringe and Taylor 2014}).

\subsection{Form comparison}\label{form-comparison-3}

The note therefore concerns inherited vocalism rather than a different
lexeme: \emph{mōder} `mother' is the regularized nominative represented
here, but dictionaries usually print \emph{mōdor} / \emph{modor}
`mother', and the oblique evidence survives in \emph{mēder} `mother'
(\citeproc{ref-ClarkHall1960}{Clark Hall 1960};
\citeproc{ref-RingeTaylor2014}{Ringe and Taylor 2014}).

\section{\texorpdfstring{net --- OE
\emph{nett}}{net --- OE nett}}\label{net-oe-nett}

\index[iv]{02oe@\textbf{Old English}!nett@\emph{nett}}
\index[iv]{03pgmc@\textbf{Proto-Germanic}!natja@*nátją}

Derivation: \emph{*nátją} \textgreater{} \emph{nett} (regular).

\subsection{Derivation trace}\label{derivation-trace-40}

Proto input: \emph{*nátją}

\begingroup
\setlength{\fboxsep}{6pt}

\noindent\fbox{%
\begin{minipage}{0.97\linewidth}
\small
\begin{tabularx}{\linewidth}{@{}>{\raggedright\arraybackslash}p{0.440\linewidth}>{\centering\arraybackslash}X>{\raggedright\arraybackslash}p{0.440\linewidth}@{}}
\begin{minipage}[t]{\linewidth}
\centering\textbf{Earlier Germanic changes}\par
\vspace{0.35em}
\raggedright
\centering\textbf{}\par
\raggedright
\vspace{0.2em}
\begin{tabular}{@{}>{\raggedright\arraybackslash}p{0.68\linewidth}@{\hspace{0.55em}}>{\raggedright\arraybackslash}p{0.16\linewidth}@{\hspace{0.25em}}}
\multicolumn{2}{@{}l@{}}{*Proto-West Germanic*} \\
PWGmc J Gemination & \emph{*náttją} \\
\multicolumn{2}{@{}l@{}}{*Northwest Germanic*} \\
\multicolumn{2}{@{}l@{}}{[no change]} \\
\end{tabular}
\end{minipage}
&
&
\begin{minipage}[t]{\linewidth}
\centering\textbf{Old English changes}\par
\vspace{0.35em}
\raggedright
\centering\textbf{}\par
\raggedright
\vspace{0.2em}
\begin{tabular}{@{}>{\raggedright\arraybackslash}p{0.74\linewidth}@{\hspace{0.55em}}>{\raggedright\arraybackslash}p{0.16\linewidth}@{\hspace{0.25em}}}
\multicolumn{2}{@{}l@{}}{*Old English*} \\
Anglo Frisian Brightening & \emph{*nættją} \\
OE Heavy Syllable Nasal Apocope & \emph{*nættj} \\
OE I Umlaut & \emph{*nettj} \\
OE J Loss After Heavy & \emph{*nett} \\
\end{tabular}
\end{minipage}
\\
\end{tabularx}
\end{minipage}%
} \endgroup

Old English form: \emph{nett}

\subsection{Reconstruction and comparative
evidence}\label{reconstruction-and-comparative-evidence-40}

Orel gives \Recon{natjan} `net' with Old English {\emph{nett}} `net',
and Fulk's account of West Germanic gemination before \emph{j} explains
the geminate outcome after a short vowel (\citeproc{ref-Orel2003}{Orel
2003}; \citeproc{ref-Fulk2018}{Fulk 2018}).

\subsection{Old English evidence}\label{old-english-evidence-40}

Clark Hall and Bosworth-Toller record {\emph{nett}} `net' as the noun,
and Campbell notes that final geminates are often graphically simplified
in Old English spelling (\citeproc{ref-ClarkHall1960}{Clark Hall 1960};
\citeproc{ref-BosworthToller1898}{Bosworth and Toller 1898, 29};
\citeproc{ref-Campbell1959}{Campbell 1959}).

\subsection{Development to Old
English}\label{development-to-old-english-40}

From \Recon{nátją} `net', West Germanic j-gemination first gives
\Recon{náttją} `net'. Later brightening, loss of the weak ending, and
loss of final \emph{j} after a heavy stem yield {\emph{nett}} `net', so
the development represented here is regular
(\citeproc{ref-Fulk2018}{Fulk 2018}).

\subsection{Form note}\label{form-note-13}

Spellings in \emph{net} can therefore be graphic simplifications, but
the lexical target supported by the dictionary evidence is {\emph{nett}}
`net' (\citeproc{ref-Campbell1959}{Campbell 1959};
\citeproc{ref-Orel2003}{Orel 2003}).

\section{\texorpdfstring{nightmare --- OE
\emph{mare}}{nightmare --- OE mare}}\label{nightmare-oe-mare}

\index[iv]{02oe@\textbf{Old English}!mare@\emph{mare}}
\index[iv]{03pgmc@\textbf{Proto-Germanic}!maron@*márōn}

Derivation: \emph{*márōn} \textgreater{} \emph{mare} (regular).

\subsection{Derivation trace}\label{derivation-trace-41}

Proto input: \emph{*márōn}

\begingroup
\setlength{\fboxsep}{6pt}

\noindent\fbox{%
\begin{minipage}{0.97\linewidth}
\small
\begin{tabularx}{\linewidth}{@{}>{\raggedright\arraybackslash}p{0.440\linewidth}>{\centering\arraybackslash}X>{\raggedright\arraybackslash}p{0.440\linewidth}@{}}
\begin{minipage}[t]{\linewidth}
\centering\textbf{Earlier Germanic changes}\par
\vspace{0.35em}
\raggedright
\centering\textbf{}\par
\raggedright
\vspace{0.2em}
\begin{tabular}{@{}>{\raggedright\arraybackslash}p{0.68\linewidth}@{\hspace{0.55em}}>{\raggedright\arraybackslash}p{0.16\linewidth}@{\hspace{0.25em}}}
\multicolumn{2}{@{}l@{}}{*Proto-West Germanic*} \\
\multicolumn{2}{@{}l@{}}{[no change]} \\
\multicolumn{2}{@{}l@{}}{*Northwest Germanic*} \\
NWGmc N Stem N Loss & \emph{*márǭ} \\
\end{tabular}
\end{minipage}
&
&
\begin{minipage}[t]{\linewidth}
\centering\textbf{Old English changes}\par
\vspace{0.35em}
\raggedright
\centering\textbf{}\par
\raggedright
\vspace{0.2em}
\begin{tabular}{@{}>{\raggedright\arraybackslash}p{0.74\linewidth}@{\hspace{0.55em}}>{\raggedright\arraybackslash}p{0.16\linewidth}@{\hspace{0.25em}}}
\multicolumn{2}{@{}l@{}}{*Old English*} \\
Anglo Frisian Brightening & \emph{*mærǭ} \\
OE A Restoration & \emph{*marǭ} \\
OE Unstressed Long Vowel Shortening & \emph{*maræ} \\
OE Unstressed AE Merger & \emph{*mare} \\
\end{tabular}
\end{minipage}
\\
\end{tabularx}
\end{minipage}%
} \endgroup

Old English form: \emph{mare}

\subsection{Reconstruction and comparative
evidence}\label{reconstruction-and-comparative-evidence-41}

Ringe and Taylor treat the lexeme as Proto-Germanic /
Proto-Northwest-Germanic \emph{*marōn-}, with Old English \emph{mare,
maran}, and variant \emph{mere}; Orel preserves the same comparative
lemma though with a different Old English headword tradition
(\citeproc{ref-RingeTaylor2014}{Ringe and Taylor 2014};
\citeproc{ref-Orel2003}{Orel 2003}).

\subsection{Old English evidence}\label{old-english-evidence-41}

Clark Hall records \emph{mare} `nightmare, monster' and also preserves
related variant forms \emph{mera} / \emph{mere}
(\citeproc{ref-ClarkHall1960}{Clark Hall 1960, 213}).

\subsection{Development to Old
English}\label{development-to-old-english-41}

The selected simplex input \Recon{márōn} `nightmare' regularly gives
\emph{mare} `nightmare' after brightening, A-restoration before the
n-stem ending, and later reduction of the final vowel. The word
represented here is the attested simplex noun, not an attested compound
(\citeproc{ref-RingeTaylor2014}{Ringe and Taylor 2014}).

\subsection{Form note}\label{form-note-14}

The concept corresponds to an unattested compound \Recon{nihtmare}
`nightmare', but the Old English lexical evidence is for simplex
\emph{mare} `nightmare', with oblique \emph{maran} and variant
\emph{mere} / \emph{mera} (\citeproc{ref-RingeTaylor2014}{Ringe and
Taylor 2014}; \citeproc{ref-ClarkHall1960}{Clark Hall 1960, 213}).

\section{\texorpdfstring{coat --- OE
\emph{rocc}}{coat --- OE rocc}}\label{coat-oe-rocc}

\index[iv]{02oe@\textbf{Old English}!rocc@\emph{rocc}}
\index[iv]{03pgmc@\textbf{Proto-Germanic}!rukkaz@*rúkkaz}

Derivation: \emph{*rúkkaz} \textgreater{} \emph{rocc} (regular).

\subsection{Derivation trace}\label{derivation-trace-42}

Proto input: \emph{*rúkkaz}

\begingroup
\setlength{\fboxsep}{6pt}

\noindent\fbox{%
\begin{minipage}{0.97\linewidth}
\small
\begin{tabularx}{\linewidth}{@{}>{\raggedright\arraybackslash}p{0.520\linewidth}>{\centering\arraybackslash}X>{\raggedright\arraybackslash}p{0.320\linewidth}@{}}
\begin{minipage}[t]{\linewidth}
\centering\textbf{Earlier Germanic changes}\par
\vspace{0.35em}
\raggedright
\centering\textbf{}\par
\raggedright
\vspace{0.2em}
\begin{tabular}{@{}>{\raggedright\arraybackslash}p{0.74\linewidth}@{\hspace{0.55em}}>{\raggedright\arraybackslash}p{0.16\linewidth}@{\hspace{0.25em}}}
\multicolumn{2}{@{}l@{}}{*Proto-West Germanic*} \\
\multicolumn{2}{@{}l@{}}{[no change]} \\
\multicolumn{2}{@{}l@{}}{*Northwest Germanic*} \\
\mbox{NWGmc U Lowering} & \emph{*rókkaz} \\
\mbox{PGmc Final Z Deletion} & \emph{*rókka} \\
\end{tabular}
\end{minipage}
&
&
\begin{minipage}[t]{\linewidth}
\centering\textbf{Old English changes}\par
\vspace{0.35em}
\raggedright
\centering\textbf{}\par
\raggedright
\vspace{0.2em}
\begin{tabular}{@{}>{\raggedright\arraybackslash}p{0.70\linewidth}@{\hspace{0.55em}}>{\raggedright\arraybackslash}p{0.16\linewidth}@{\hspace{0.25em}}}
\multicolumn{2}{@{}l@{}}{*Old English*} \\
PWGmc Final Bare A Loss & \emph{*rókk} \\
\end{tabular}
\end{minipage}
\\
\end{tabularx}
\end{minipage}%
} \endgroup

Old English form: \emph{rocc}

\subsection{Reconstruction and comparative
evidence}\label{reconstruction-and-comparative-evidence-42}

Orel cites a masculine \Recon{rukkaz} `coat' for the garment word, while
Kroonen gives \emph{*hrukka-}. Both treat this as the garment lexeme and
not as the separate stone word (\citeproc{ref-Orel2003}{Orel 2003};
\citeproc{ref-Kroonen2013}{Kroonen 2013, 290}).

\subsection{Old English evidence}\label{old-english-evidence-42}

Clark Hall and Bosworth-Toller record \emph{rocc} `coat' as an
over-garment or tunic and preserve compounds such as \emph{bisceoprocc}
and \emph{breóstrocc} `breast-coat' (\citeproc{ref-ClarkHall1960}{Clark
Hall 1960}; \citeproc{ref-BosworthToller1898}{Bosworth and Toller
1898}).

\subsection{Development to Old
English}\label{development-to-old-english-42}

With \Recon{rúkkaz} `coat' as the derivational input, Northwest Germanic
u-lowering and later loss of final \emph{-a} yield \emph{rocc} as a
regular outcome.

\subsection{Source note}\label{source-note-3}

This entry concerns the garment noun only. The stone word seen in
\emph{stānrocc} `stone-coat' belongs to a different lexical history
(\citeproc{ref-ClarkHall1960}{Clark Hall 1960};
\citeproc{ref-BosworthToller1898}{Bosworth and Toller 1898}).

\section{\texorpdfstring{sheep --- OE
\emph{sċēap}}{sheep --- OE sċēap}}\label{sheep-oe-sux10bux113ap}

\index[iv]{02oe@\textbf{Old English}!sceap@\emph{sċēap}}
\index[iv]{03pgmc@\textbf{Proto-Germanic}!skepa@*skḗpą}

Derivation: \emph{*skḗpą} \textgreater{} \emph{sċēap} (regular).

\subsection{Derivation trace}\label{derivation-trace-43}

Proto input: \emph{*skḗpą}

\begingroup
\setlength{\fboxsep}{6pt}

\noindent\fbox{%
\begin{minipage}{0.97\linewidth}
\small
\begin{tabularx}{\linewidth}{@{}>{\raggedright\arraybackslash}p{0.460\linewidth}>{\centering\arraybackslash}X>{\raggedright\arraybackslash}p{0.460\linewidth}@{}}
\begin{minipage}[t]{\linewidth}
\centering\textbf{Earlier Germanic changes}\par
\vspace{0.35em}
\raggedright
\centering\textbf{}\par
\raggedright
\vspace{0.2em}
\begin{tabular}{@{}>{\raggedright\arraybackslash}p{0.74\linewidth}@{\hspace{0.55em}}>{\raggedright\arraybackslash}p{0.16\linewidth}@{\hspace{0.25em}}}
\multicolumn{2}{@{}l@{}}{*Proto-West Germanic*} \\
\multicolumn{2}{@{}l@{}}{[no change]} \\
\multicolumn{2}{@{}l@{}}{*Northwest Germanic*} \\
\mbox{NWGmc Long E Lowering} & \emph{*skǣpą} \\
\end{tabular}
\end{minipage}
&
&
\begin{minipage}[t]{\linewidth}
\centering\textbf{Old English changes}\par
\vspace{0.35em}
\raggedright
\centering\textbf{}\par
\raggedright
\vspace{0.2em}
\begin{tabular}{@{}>{\raggedright\arraybackslash}p{0.74\linewidth}@{\hspace{0.55em}}>{\raggedright\arraybackslash}p{0.16\linewidth}@{\hspace{0.25em}}}
\multicolumn{2}{@{}l@{}}{*Old English*} \\
OE Heavy Syllable Nasal Apocope & \emph{*skǣp} \\
OE Sk Palatalization & \emph{*ʃǣp} \\
OE Ws Palatal Diphthongization & \emph{*ʃēap} \\
\end{tabular}
\end{minipage}
\\
\end{tabularx}
\end{minipage}%
} \endgroup

Old English form: \emph{sċēap}

\subsection{Reconstruction and comparative
evidence}\label{reconstruction-and-comparative-evidence-43}

Ringe and Taylor cite a later West Germanic \emph{*skap} `sheep'
\textgreater{} WS \emph{scéap} `sheep', while Orel preserves a
Proto-Germanic noun of the \emph{*skēp-} type for the same lexeme
(\citeproc{ref-RingeTaylor2014}{Ringe and Taylor 2014};
\citeproc{ref-Orel2003}{Orel 2003}).

\subsection{Old English evidence}\label{old-english-evidence-43}

Clark Hall records \emph{scēap} `sheep' with spelling variation, and
Campbell likewise lists West Saxon \emph{scéap} `sheep' among the
palatal-diphthongized forms (\citeproc{ref-ClarkHall1960}{Clark Hall
1960}; \citeproc{ref-Campbell1959}{Campbell 1959}).

\subsection{Development to Old
English}\label{development-to-old-english-43}

From \Recon{skḗpą} `sheep', Northwest Germanic lowering gives
\Recon{skǣpą} `sheep'; after apocope and palatalization the West Saxon
branch diphthongizes to \emph{sċēap} `sheep'. The development
represented here is therefore fully regular.

\subsection{Dialect note}\label{dialect-note-4}

Ringe and Taylor contrast West Saxon \emph{scéap} `sheep' with Mercian
and Kentish \emph{scép} `sheep', and Campbell also notes Northumbrian
\emph{scip} `sheep'. The form represented here is the West Saxon
headword (\citeproc{ref-RingeTaylor2014}{Ringe and Taylor 2014};
\citeproc{ref-Campbell1959}{Campbell 1959}).

\section{\texorpdfstring{shilling --- OE
\emph{sċilling}}{shilling --- OE sċilling}}\label{shilling-oe-sux10billing}

\index[iv]{02oe@\textbf{Old English}!scilling@\emph{sċilling}}
\index[iv]{03pgmc@\textbf{Proto-Germanic}!skillingaz@*skíllingaz}

Derivation: \emph{*skíllingaz} \textgreater{} \emph{sċilling} (regular).

\subsection{Derivation trace}\label{derivation-trace-44}

Proto input: \emph{*skíllingaz}

\begingroup
\setlength{\fboxsep}{6pt}

\noindent\fbox{%
\begin{minipage}{0.97\linewidth}
\small
\begin{tabularx}{\linewidth}{@{}>{\raggedright\arraybackslash}p{0.440\linewidth}>{\centering\arraybackslash}X>{\raggedright\arraybackslash}p{0.440\linewidth}@{}}
\begin{minipage}[t]{\linewidth}
\centering\textbf{Earlier Germanic changes}\par
\vspace{0.35em}
\raggedright
\centering\textbf{}\par
\raggedright
\vspace{0.2em}
\begin{tabular}{@{}>{\raggedright\arraybackslash}p{0.66\linewidth}@{\hspace{0.55em}}>{\raggedright\arraybackslash}p{0.24\linewidth}@{\hspace{0.25em}}}
\multicolumn{2}{@{}l@{}}{*Proto-West Germanic*} \\
\multicolumn{2}{@{}l@{}}{[no change]} \\
\multicolumn{2}{@{}l@{}}{*Northwest Germanic*} \\
PGmc Final Z Deletion & \emph{*skíllinga} \\
\end{tabular}
\end{minipage}
&
&
\begin{minipage}[t]{\linewidth}
\centering\textbf{Old English changes}\par
\vspace{0.35em}
\raggedright
\centering\textbf{}\par
\raggedright
\vspace{0.2em}
\begin{tabular}{@{}>{\raggedright\arraybackslash}p{0.68\linewidth}@{\hspace{0.55em}}>{\raggedright\arraybackslash}p{0.22\linewidth}@{\hspace{0.25em}}}
\multicolumn{2}{@{}l@{}}{*Old English*} \\
PWGmc Final Bare A Loss & \emph{*skílling} \\
OE Sk Palatalization & \emph{*ʃílling} \\
OE Med Unstressed I Lowering1 & \emph{*ʃílleng} \\
OE Med Unstressed I Lowering & \emph{*ʃílling} \\
\end{tabular}
\end{minipage}
\\
\end{tabularx}
\end{minipage}%
} \endgroup

Old English form: \emph{sċilling}

\subsection{Reconstruction and comparative
evidence}\label{reconstruction-and-comparative-evidence-44}

Kroonen treats the cognate set under \emph{*skellinga-}
\textasciitilde{} \emph{*skillinga-} and connects it with
\emph{*skeld-linga-}, while Orel likewise gives the coin word with OE
{\emph{scilling}} `shilling' among the reflexes
(\citeproc{ref-Kroonen2013}{Kroonen 2013}; \citeproc{ref-Orel2003}{Orel
2003}). The derivational input \Recon{skíllingaz} `shilling' is the
nominative-style form used here to represent that inherited
\emph{*-ing-} derivative.

\subsection{Old English evidence}\label{old-english-evidence-44}

Clark Hall records \emph{scilling} `shilling', and Campbell cites it
among nouns with unstressed \emph{i} in derivational \emph{-ing}
(\citeproc{ref-ClarkHall1960}{Clark Hall 1960};
\citeproc{ref-Campbell1959}{Campbell 1959}). The target is the ordinary
OE citation form, normalized as \emph{sċilling} `shilling'.

\subsection{Development to Old
English}\label{development-to-old-english-44}

From \Recon{skíllingaz} `shilling', loss of final \emph{-az} yields
\Recon{skílling} `shilling'. Old English palatalization of initial
\emph{sk} before front vocalism then gives \emph{sċilling} `shilling'.
The \emph{i} of derivational \emph{-ing-} remains, so the regular
outcome is \emph{sċilling} `shilling', not \Recon{sċilleng} `shilling'
(\citeproc{ref-Campbell1959}{Campbell 1959};
\citeproc{ref-Hogg1992}{Hogg 1992}).

\subsection{Form note}\label{form-note-15}

Kroonen's \emph{*skellinga-} \textasciitilde{} \emph{*skillinga-} and
his internal analysis \emph{*skeld-linga-} belong to the etymological
background of the cognate set. The derivational input \Recon{skíllingaz}
`shilling' is the specific form used for the derivation represented here
(\citeproc{ref-Kroonen2013}{Kroonen 2013}).

\section{\texorpdfstring{show --- OE
\emph{sċēawian}}{show --- OE sċēawian}}\label{show-oe-sux10bux113awian}

\index[iv]{02oe@\textbf{Old English}!sceawian@\emph{sċēawian}}
\index[iv]{03pgmc@\textbf{Proto-Germanic}!skawojana@*skáwōjaną}

Derivation: \emph{*skáwōjaną} \textgreater{} \emph{sċēawian} (regular).

\subsection{Derivation trace}\label{derivation-trace-45}

Proto input: \emph{*skáwōjaną}

\begingroup
\setlength{\fboxsep}{6pt}

\noindent\fbox{%
\begin{minipage}{0.97\linewidth}
\footnotesize
\begin{tabularx}{\linewidth}{@{}>{\raggedright\arraybackslash}p{0.320\linewidth}>{\centering\arraybackslash}X>{\raggedright\arraybackslash}p{0.520\linewidth}@{}}
\begin{minipage}[t]{\linewidth}
\centering\textbf{Earlier Germanic changes}\par
\vspace{0.35em}
\raggedright
\centering\textbf{}\par
\raggedright
\vspace{0.2em}
\begin{tabular}{@{}>{\raggedright\arraybackslash}p{0.68\linewidth}@{\hspace{0.55em}}>{\raggedright\arraybackslash}p{0.16\linewidth}@{\hspace{0.25em}}}
\multicolumn{2}{@{}l@{}}{*Proto-West Germanic*} \\
\multicolumn{2}{@{}l@{}}{[no change]} \\
\multicolumn{2}{@{}l@{}}{*Northwest Germanic*} \\
\multicolumn{2}{@{}l@{}}{[no change]} \\
\end{tabular}
\end{minipage}
&
&
\begin{minipage}[t]{\linewidth}
\centering\textbf{Old English changes}\par
\vspace{0.35em}
\raggedright
\centering\textbf{}\par
\raggedright
\vspace{0.2em}
\begin{tabular}{@{}>{\raggedright\arraybackslash}p{0.62\linewidth}@{\hspace{0.55em}}>{\raggedright\arraybackslash}p{0.28\linewidth}@{\hspace{0.25em}}}
\multicolumn{2}{@{}l@{}}{*Old English*} \\
OE Aw Long Diphthong & \emph{*skḗawōjaną} \\
OE Heavy Syllable Nasal Apocope & \emph{*skḗawōjan} \\
OE Secondary Nasalization & \emph{*skḗawōjąn} \\
OE Sk Palatalization & \emph{*ʃḗawōjąn} \\
OE I Umlaut & \emph{*ʃḗawējąn} \\
OE Unstressed Long Vowel Shortening & \emph{*ʃḗawejąn} \\
OE Weak Tail Reduction & \emph{*ʃḗawejan} \\
OE Intervocalic J Vocalization & \emph{*ʃḗaweian} \\
OE Unstressed EI Contraction & \emph{*ʃḗawian} \\
\end{tabular}
\end{minipage}
\\
\end{tabularx}
\end{minipage}%
} \endgroup

Old English form: \emph{sċēawian}

\subsection{Reconstruction and comparative
evidence}\label{reconstruction-and-comparative-evidence-45}

Orel and Kroonen cite a Class II verb of the type {\emph{*skawōjan-}},
with OE {\emph{scēawian}} `show' among the reflexes
(\citeproc{ref-Orel2003}{Orel 2003}; \citeproc{ref-Kroonen2013}{Kroonen
2013, 482}). Brunner likewise records the Old English family as
{\emph{scēawian}} `show', {\emph{scāwian}} `show', confirming that it
belongs to the ordinary show-verb set rather than to a special
finite-cell formation (\citeproc{ref-SieversBrunner1965}{Brunner 1965}).

\subsection{Old English evidence}\label{old-english-evidence-45}

Bright lists {\emph{scēawian}} `show' (W. II.) and also the related form
{\emph{scēawa}} `show' (\citeproc{ref-BrightCassidyRingler1971}{Bright
1971}). The sources therefore use {\emph{scēawian}} `show', while
{\emph{sċēawian}} `show' supplies a normalized spelling.

\subsection{Development to Old
English}\label{development-to-old-english-45}

From {\Recon{skáwōjaną} `show'}, the inherited \emph{aw} sequence before
a following vowel develops into the \emph{ēaw} diphthong, and \emph{*ō}
survives between \emph{*w} and \emph{*j} in the Class II suffix. The
development therefore runs regularly to {\emph{sċēawian}} `show',
without the direct \emph{aw+j} problem seen in other verb types
(\citeproc{ref-Campbell1959}{Campbell 1959};
\citeproc{ref-Orel2003}{Orel 2003}).

\subsection{Form note}\label{form-note-16}

The difference between {\emph{scēawian}} `show' and {\emph{sċēawian}}
`show' is orthographic normalization of initial
\textless{}\emph{sc}\textgreater, not a difference of lexeme or paradigm
cell (\citeproc{ref-Campbell1959}{Campbell 1959};
\citeproc{ref-Hogg1992}{Hogg 1992}).

\section{\texorpdfstring{sleep --- OE
\emph{slǣpan}}{sleep --- OE slǣpan}}\label{sleep-oe-slux1e3pan}

\index[iv]{02oe@\textbf{Old English}!slaepan@\emph{slǣpan}}
\index[iv]{03pgmc@\textbf{Proto-Germanic}!slepana@*slḗpaną}

Derivation: \emph{*slḗpaną} \textgreater{} \emph{slǣpan} (regular).

\subsection{Derivation trace}\label{derivation-trace-46}

Proto input: \emph{*slḗpaną}

\begingroup
\setlength{\fboxsep}{6pt}

\noindent\fbox{%
\begin{minipage}{0.97\linewidth}
\small
\begin{tabularx}{\linewidth}{@{}>{\raggedright\arraybackslash}p{0.460\linewidth}>{\centering\arraybackslash}X>{\raggedright\arraybackslash}p{0.460\linewidth}@{}}
\begin{minipage}[t]{\linewidth}
\centering\textbf{Earlier Germanic changes}\par
\vspace{0.35em}
\raggedright
\centering\textbf{}\par
\raggedright
\vspace{0.2em}
\begin{tabular}{@{}>{\raggedright\arraybackslash}p{0.70\linewidth}@{\hspace{0.55em}}>{\raggedright\arraybackslash}p{0.20\linewidth}@{\hspace{0.25em}}}
\multicolumn{2}{@{}l@{}}{*Proto-West Germanic*} \\
\multicolumn{2}{@{}l@{}}{[no change]} \\
\multicolumn{2}{@{}l@{}}{*Northwest Germanic*} \\
\mbox{NWGmc Long E Lowering} & \emph{*slǣpaną} \\
\end{tabular}
\end{minipage}
&
&
\begin{minipage}[t]{\linewidth}
\centering\textbf{Old English changes}\par
\vspace{0.35em}
\raggedright
\centering\textbf{}\par
\raggedright
\vspace{0.2em}
\begin{tabular}{@{}>{\raggedright\arraybackslash}p{0.74\linewidth}@{\hspace{0.55em}}>{\raggedright\arraybackslash}p{0.16\linewidth}@{\hspace{0.25em}}}
\multicolumn{2}{@{}l@{}}{*Old English*} \\
OE Heavy Syllable Nasal Apocope & \emph{*slǣpan} \\
OE Secondary Nasalization & \emph{*slǣpąn} \\
OE Weak Tail Reduction & \emph{*slǣpan} \\
\end{tabular}
\end{minipage}
\\
\end{tabularx}
\end{minipage}%
} \endgroup

Old English form: \emph{slǣpan}

\subsection{Reconstruction and comparative
evidence}\label{reconstruction-and-comparative-evidence-46}

Kroonen preserves the comparative verb as \emph{*slēpan-}, and Fulk
cites the same family under root \emph{*slēb-}
(\citeproc{ref-Kroonen2013}{Kroonen 2013}; \citeproc{ref-Fulk2018}{Fulk
2018, 120}). The derivational input \Recon{slḗpaną} `sleep' is the
infinitive-style form used here for that inherited sleep-verb.

\subsection{Old English evidence}\label{old-english-evidence-46}

Clark Hall gives {\emph{slæpan}} `sleep' with preterite \emph{slēp}
`slept', \emph{slēap} `slept', and Bright likewise lists {\emph{slæpan}}
`sleep' ({\emph{slāpan}} `sleep'), \emph{slēp} `slept', \emph{slēpon}
`slept', \emph{slēpen} `slept' (\citeproc{ref-ClarkHall1960}{Clark Hall
1960}; \citeproc{ref-BrightCassidyRingler1971}{Bright 1971, 435}). The
target represented here is therefore the normalized infinitive
\emph{slǣpan} `sleep', not the preterite forms and not the separate noun
\emph{slǣp} `sleep'.

\subsection{Development to Old
English}\label{development-to-old-english-46}

From \Recon{slḗpaną} `sleep', Northwest Germanic lowering gives
\Recon{slǣpaną} `sleep'. The later OE tail developments then yield
{\emph{slǣpan}} `sleep' regularly. Brunner and Bülbring show that the OE
tradition also has variant spellings such as West Saxon {\emph{slāpan}}
`sleep' / {\emph{slæpan}} `sleep' and Anglian or Kentish {\emph{slēpan}}
`sleep', but those do not displace the infinitive chosen here
(\citeproc{ref-SieversBrunner1965}{Brunner 1965};
\citeproc{ref-Bulbring1902}{Bülbring 1902}).

\subsection{Form note}\label{form-note-17}

The note concerns lemma type rather than a special derivational problem:
this row represents the verb \emph{slǣpan} `sleep', whereas \emph{slǣp}
`sleep' belongs to noun or lookup background and \emph{slēp} `slept' /
\emph{slēap} `slept' are preterite forms
(\citeproc{ref-ClarkHall1960}{Clark Hall 1960};
\citeproc{ref-BrightCassidyRingler1971}{Bright 1971, 435}).

\section{\texorpdfstring{smear --- OE
\emph{smierwan}}{smear --- OE smierwan}}\label{smear-oe-smierwan}

\index[iv]{02oe@\textbf{Old English}!smierwan@\emph{smierwan}}
\index[iv]{03pgmc@\textbf{Proto-Germanic}!smerwijana@*smérwijaną}

Derivation: \emph{*smérwijaną} \textgreater{} \emph{smierwan} (regular).

\subsection{Derivation trace}\label{derivation-trace-47}

Proto input: \emph{*smérwijaną}

\begingroup
\setlength{\fboxsep}{6pt}

\noindent\fbox{%
\begin{minipage}{0.97\linewidth}
\footnotesize
\begin{tabularx}{\linewidth}{@{}>{\raggedright\arraybackslash}p{0.320\linewidth}>{\centering\arraybackslash}X>{\raggedright\arraybackslash}p{0.520\linewidth}@{}}
\begin{minipage}[t]{\linewidth}
\centering\textbf{Earlier Germanic changes}\par
\vspace{0.35em}
\raggedright
\centering\textbf{}\par
\raggedright
\vspace{0.2em}
\begin{tabular}{@{}>{\raggedright\arraybackslash}p{0.68\linewidth}@{\hspace{0.55em}}>{\raggedright\arraybackslash}p{0.16\linewidth}@{\hspace{0.25em}}}
\multicolumn{2}{@{}l@{}}{*Proto-West Germanic*} \\
\multicolumn{2}{@{}l@{}}{[no change]} \\
\multicolumn{2}{@{}l@{}}{*Northwest Germanic*} \\
\multicolumn{2}{@{}l@{}}{[no change]} \\
\end{tabular}
\end{minipage}
&
&
\begin{minipage}[t]{\linewidth}
\centering\textbf{Old English changes}\par
\vspace{0.35em}
\raggedright
\centering\textbf{}\par
\raggedright
\vspace{0.2em}
\begin{tabular}{@{}>{\raggedright\arraybackslash}p{0.62\linewidth}@{\hspace{0.55em}}>{\raggedright\arraybackslash}p{0.28\linewidth}@{\hspace{0.25em}}}
\multicolumn{2}{@{}l@{}}{*Old English*} \\
OE Breaking & \emph{*sméorwijaną} \\
OE Heavy Syllable Nasal Apocope & \emph{*sméorwijan} \\
OE Secondary Nasalization & \emph{*sméorwijąn} \\
Sievers Law Syncope & \emph{*sméorwjąn} \\
OE I Umlaut & \emph{*smíerwjąn} \\
OE Weak Tail Reduction & \emph{*smíerwjan} \\
OE J Loss After Heavy & \emph{*smíerwan} \\
\end{tabular}
\end{minipage}
\\
\end{tabularx}
\end{minipage}%
} \endgroup

Old English form: \emph{smierwan}

\subsection{Reconstruction and comparative
evidence}\label{reconstruction-and-comparative-evidence-47}

Kroonen gives the comparative headword as \emph{*smerwjan-}
(\citeproc{ref-Kroonen2013}{Kroonen 2013}). Ringe and Taylor instead
cite a later-stage \Recon{smirwijana} `smear', from which they derive
West Saxon {\emph{smierwan}} `smear', Mercian {\emph{smirwan}} `smear',
and Northumbrian \emph{smiriga} (\citeproc{ref-RingeTaylor2014}{Ringe
and Taylor 2014}). The derivational input \Recon{smérwijaną} `smear'
therefore represents the Kroonen-aligned PGmc layer, while the
later-stage dialect split belongs to a different chronological level.

\subsection{Old English evidence}\label{old-english-evidence-47}

The target represented here is the West Saxon citation form
{\emph{smierwan}} `smear'. Campbell's Anglian discussion explains the
contrasting {\emph{smirwan}} `smear', and Clark Hall, Brunner, and
Bright show that the same lexical family later also includes forms such
as \emph{smirian}, \emph{smyrian}, and preterite \emph{smyrode}
(\citeproc{ref-Campbell1959}{Campbell 1959};
\citeproc{ref-ClarkHall1960}{Clark Hall 1960};
\citeproc{ref-SieversBrunner1965}{Brunner 1965};
\citeproc{ref-BrightCassidyRingler1971}{Bright 1971}).

\subsection{Development to Old
English}\label{development-to-old-english-47}

From \Recon{smérwijaną} `smear', breaking before \emph{r + consonant}
yields \emph{eo}, and later i-umlaut produces \emph{ie}. The result is
West Saxon {\emph{smierwan}} `smear'. Anglian {\emph{smirwan}} `smear'
reflects the well-known failure of breaking in this environment, not a
different lexeme (\citeproc{ref-Campbell1959}{Campbell 1959};
\citeproc{ref-RingeTaylor2014}{Ringe and Taylor 2014}).

\subsection{Dialect note}\label{dialect-note-5}

The entry therefore represents the West Saxon member of a broader OE
family: {\emph{smierwan}} `smear' in West Saxon, {\emph{smirwan}}
`smear' in Anglian or Mercian, and related later class-II forms such as
\emph{smirian} or \emph{smyrian} in the same lexical field
(\citeproc{ref-RingeTaylor2014}{Ringe and Taylor 2014};
\citeproc{ref-ClarkHall1960}{Clark Hall 1960}).

\section{\texorpdfstring{span --- OE
\emph{spannan}}{span --- OE spannan}}\label{span-oe-spannan}

\index[iv]{02oe@\textbf{Old English}!spannan@\emph{spannan}}
\index[iv]{03pgmc@\textbf{Proto-Germanic}!spannana@*spánnaną}

Derivation: \emph{*spánnaną} \textgreater{} \emph{spannan} (regular).

\subsection{Derivation trace}\label{derivation-trace-48}

Proto input: \emph{*spánnaną}

\begingroup
\setlength{\fboxsep}{6pt}

\noindent\fbox{%
\begin{minipage}{0.97\linewidth}
\small
\begin{tabularx}{\linewidth}{@{}>{\raggedright\arraybackslash}p{0.320\linewidth}>{\centering\arraybackslash}X>{\raggedright\arraybackslash}p{0.520\linewidth}@{}}
\begin{minipage}[t]{\linewidth}
\centering\textbf{Earlier Germanic changes}\par
\vspace{0.35em}
\raggedright
\centering\textbf{}\par
\raggedright
\vspace{0.2em}
\begin{tabular}{@{}>{\raggedright\arraybackslash}p{0.68\linewidth}@{\hspace{0.55em}}>{\raggedright\arraybackslash}p{0.16\linewidth}@{\hspace{0.25em}}}
\multicolumn{2}{@{}l@{}}{*Proto-West Germanic*} \\
\multicolumn{2}{@{}l@{}}{[no change]} \\
\multicolumn{2}{@{}l@{}}{*Northwest Germanic*} \\
\multicolumn{2}{@{}l@{}}{[no change]} \\
\end{tabular}
\end{minipage}
&
&
\begin{minipage}[t]{\linewidth}
\centering\textbf{Old English changes}\par
\vspace{0.35em}
\raggedright
\centering\textbf{}\par
\raggedright
\vspace{0.2em}
\begin{tabular}{@{}>{\raggedright\arraybackslash}p{0.70\linewidth}@{\hspace{0.55em}}>{\raggedright\arraybackslash}p{0.20\linewidth}@{\hspace{0.25em}}}
\multicolumn{2}{@{}l@{}}{*Old English*} \\
OE Heavy Syllable Nasal Apocope & \emph{*spánnan} \\
OE Secondary Nasalization & \emph{*spánnąn} \\
OE Weak Tail Reduction & \emph{*spánnan} \\
\end{tabular}
\end{minipage}
\\
\end{tabularx}
\end{minipage}%
} \endgroup

Old English form: \emph{spannan}

\subsection{Reconstruction and comparative
evidence}\label{reconstruction-and-comparative-evidence-48}

Kroonen cites the inherited verb as \emph{*spannan-}, with OE
{\emph{spannan}} `span, fasten' among the reflexes
(\citeproc{ref-Kroonen2013}{Kroonen 2013}). The derivational input
\Recon{spánnaną} `span' is the infinitive-style form used here for that
same verbal lexeme.

\subsection{Old English evidence}\label{old-english-evidence-48}

Clark Hall lists noun \emph{spann} and verb {\emph{spannan}} `span,
fasten' as separate headwords, and Brunner likewise records
\emph{sponnan, spannan stv.} (\citeproc{ref-ClarkHall1960}{Clark Hall
1960}; \citeproc{ref-SieversBrunner1965}{Brunner 1965}). The target is
the strong-verb infinitive, not the noun.

\subsection{Development to Old
English}\label{development-to-old-english-48}

From \Recon{spánnaną} `span', the final nasal ending is lost and the
regular OE weak-tail steps surface {\emph{spannan}} `span, fasten'. No
paradigm-cell substitution is needed: the current derivation already
lands on the infinitive directly.

\subsection{Form note}\label{form-note-18}

English \emph{span} can also reach the noun \emph{spann} in local lookup
material. The form represented here is the verb {\emph{spannan}} `span,
fasten', with the noun treated elsewhere
(\citeproc{ref-ClarkHall1960}{Clark Hall 1960}).

\section{\texorpdfstring{spar --- OE
\emph{spearra}}{spar --- OE spearra}}\label{spar-oe-spearra}

\index[iv]{02oe@\textbf{Old English}!spearra@\emph{spearra}}
\index[iv]{03pgmc@\textbf{Proto-Germanic}!sparro@*spárrô}

Derivation: \emph{*spárrô} \textgreater{} \emph{spearra} (regular).

\subsection{Derivation trace}\label{derivation-trace-49}

Proto input: \emph{*spárrô}

\begingroup
\setlength{\fboxsep}{6pt}

\noindent\fbox{%
\begin{minipage}{0.97\linewidth}
\small
\begin{tabularx}{\linewidth}{@{}>{\raggedright\arraybackslash}p{0.320\linewidth}>{\centering\arraybackslash}X>{\raggedright\arraybackslash}p{0.520\linewidth}@{}}
\begin{minipage}[t]{\linewidth}
\centering\textbf{Earlier Germanic changes}\par
\vspace{0.35em}
\raggedright
\centering\textbf{}\par
\raggedright
\vspace{0.2em}
\begin{tabular}{@{}>{\raggedright\arraybackslash}p{0.68\linewidth}@{\hspace{0.55em}}>{\raggedright\arraybackslash}p{0.16\linewidth}@{\hspace{0.25em}}}
\multicolumn{2}{@{}l@{}}{*Proto-West Germanic*} \\
\multicolumn{2}{@{}l@{}}{[no change]} \\
\multicolumn{2}{@{}l@{}}{*Northwest Germanic*} \\
\multicolumn{2}{@{}l@{}}{[no change]} \\
\end{tabular}
\end{minipage}
&
&
\begin{minipage}[t]{\linewidth}
\centering\textbf{Old English changes}\par
\vspace{0.35em}
\raggedright
\centering\textbf{}\par
\raggedright
\vspace{0.2em}
\begin{tabular}{@{}>{\raggedright\arraybackslash}p{0.70\linewidth}@{\hspace{0.55em}}>{\raggedright\arraybackslash}p{0.20\linewidth}@{\hspace{0.25em}}}
\multicolumn{2}{@{}l@{}}{*Old English*} \\
Anglo Frisian Brightening & \emph{*spærrô} \\
OE Breaking & \emph{*spearrô} \\
OE Unstressed Long Vowel Shortening & \emph{*spearra} \\
\end{tabular}
\end{minipage}
\\
\end{tabularx}
\end{minipage}%
} \endgroup

Old English form: \emph{spearra}

\subsection{Reconstruction and comparative
evidence}\label{reconstruction-and-comparative-evidence-49}

Kroonen and Orel place this noun in the beam or rafter set
\emph{*spar(r)an-}, with cognates such as Old Saxon and Old High German
{\emph{sparro}} `spar, rafter' (\citeproc{ref-Kroonen2013}{Kroonen
2013}; \citeproc{ref-Orel2003}{Orel 2003}). The derivational input
\Recon{spárrô} `spar' is the OE-facing nominal form used here for that
same lexeme.

\subsection{Old English evidence}\label{old-english-evidence-49}

The noun is {\emph{spearra}} `spar, rafter'. English gloss overlap also
reaches the unrelated verb \emph{sperran} `to bar', which does not
belong to this row.

\subsection{Development to Old
English}\label{development-to-old-english-49}

From \Recon{spárrô} `spar', Anglo-Frisian brightening gives
\Recon{spærrô} `spar', and OE breaking before geminate \emph{rr} yields
\Recon{spearrô} `spar', later {\emph{spearra}} `spar, rafter'. The
development is therefore regular for a breaking-conditioned noun of this
type (\citeproc{ref-Luick1914}{Luick 1914-\/-1940}).

\subsection{Form note}\label{form-note-19}

This entry concerns the noun {\emph{spearra}} `spar, rafter' only. It
should be kept separate from verb \emph{sperran}, even though the Modern
English glosses overlap (\citeproc{ref-Kroonen2013}{Kroonen 2013};
\citeproc{ref-Orel2003}{Orel 2003}).

\section{\texorpdfstring{still --- OE
\emph{stillan}}{still --- OE stillan}}\label{still-oe-stillan}

\index[iv]{02oe@\textbf{Old English}!stillan@\emph{stillan}}
\index[iv]{03pgmc@\textbf{Proto-Germanic}!stellijana@*stéllijaną}

Derivation: \emph{*stéllijaną} \textgreater{} \emph{stillan} (regular).

\subsection{Derivation trace}\label{derivation-trace-50}

Proto input: \emph{*stéllijaną}

\begingroup
\setlength{\fboxsep}{6pt}

\noindent\fbox{%
\begin{minipage}{0.97\linewidth}
\footnotesize
\begin{tabularx}{\linewidth}{@{}>{\raggedright\arraybackslash}p{0.320\linewidth}>{\centering\arraybackslash}X>{\raggedright\arraybackslash}p{0.520\linewidth}@{}}
\begin{minipage}[t]{\linewidth}
\centering\textbf{Earlier Germanic changes}\par
\vspace{0.35em}
\raggedright
\centering\textbf{}\par
\raggedright
\vspace{0.2em}
\begin{tabular}{@{}>{\raggedright\arraybackslash}p{0.68\linewidth}@{\hspace{0.55em}}>{\raggedright\arraybackslash}p{0.16\linewidth}@{\hspace{0.25em}}}
\multicolumn{2}{@{}l@{}}{*Proto-West Germanic*} \\
\multicolumn{2}{@{}l@{}}{[no change]} \\
\multicolumn{2}{@{}l@{}}{*Northwest Germanic*} \\
\multicolumn{2}{@{}l@{}}{[no change]} \\
\end{tabular}
\end{minipage}
&
&
\begin{minipage}[t]{\linewidth}
\centering\textbf{Old English changes}\par
\vspace{0.35em}
\raggedright
\centering\textbf{}\par
\raggedright
\vspace{0.2em}
\begin{tabular}{@{}>{\raggedright\arraybackslash}p{0.64\linewidth}@{\hspace{0.55em}}>{\raggedright\arraybackslash}p{0.26\linewidth}@{\hspace{0.25em}}}
\multicolumn{2}{@{}l@{}}{*Old English*} \\
OE Heavy Syllable Nasal Apocope & \emph{*stéllijan} \\
OE Secondary Nasalization & \emph{*stéllijąn} \\
Sievers Law Syncope & \emph{*stélljąn} \\
OE I Umlaut & \emph{*stilljąn} \\
OE Weak Tail Reduction & \emph{*stilljan} \\
OE J Loss After Heavy & \emph{*stillan} \\
\end{tabular}
\end{minipage}
\\
\end{tabularx}
\end{minipage}%
} \endgroup

Old English form: \emph{stillan}

\subsection{Reconstruction and comparative
evidence}\label{reconstruction-and-comparative-evidence-50}

The wider West Germanic family includes adjective {\emph{still}} `still,
quiet' and verb \emph{stillen} (\citeproc{ref-KlugeSeebold2011}{Kluge
2011}). The derivational input \Recon{stéllijaną} `still' represents the
verbal j-formation used for the OE row.

\subsection{Old English evidence}\label{old-english-evidence-50}

Clark Hall gives \emph{stillan} as the verb and separately
{\emph{stille}} `still, quiet' as the adjective
(\citeproc{ref-ClarkHall1960}{Clark Hall 1960}). Bosworth-Toller
likewise preserves a substantial prefixed verbal family under
\emph{ge-stillan} `still, stop' and related forms
(\citeproc{ref-BosworthToller1898}{Bosworth and Toller 1898, 724}). The
Old English form here is the verb \emph{stillan}, not the adjective.

\subsection{Development to Old
English}\label{development-to-old-english-50}

As a heavy-stem Class I weak verb, \Recon{stéllijaną} `still' undergoes
the expected syncope and i-umlaut, and later loss of \emph{j} after a
heavy stem yields \emph{stillan}. The development represented here is
regular.

\subsection{Form note}\label{form-note-20}

The note concerns lexical framing rather than sound law: \emph{stillan}
is the verb represented here, while {\emph{stille}} `still, quiet'
belongs to the related adjectival branch of the family
(\citeproc{ref-ClarkHall1960}{Clark Hall 1960};
\citeproc{ref-KlugeSeebold2011}{Kluge 2011}).

\section{\texorpdfstring{summer --- OE
\emph{sumer}}{summer --- OE sumer}}\label{summer-oe-sumer}

\index[iv]{02oe@\textbf{Old English}!sumer@\emph{sumer}}
\index[iv]{03pgmc@\textbf{Proto-Germanic}!sumaraz@*súmaraz}

Derivation: \emph{*súmaraz} \textgreater{} \emph{sumer} (regular).

\subsection{Derivation trace}\label{derivation-trace-51}

Proto input: \emph{*súmaraz}

\begingroup
\setlength{\fboxsep}{6pt}

\noindent\fbox{%
\begin{minipage}{0.97\linewidth}
\small
\begin{tabularx}{\linewidth}{@{}>{\raggedright\arraybackslash}p{0.440\linewidth}>{\centering\arraybackslash}X>{\raggedright\arraybackslash}p{0.440\linewidth}@{}}
\begin{minipage}[t]{\linewidth}
\centering\textbf{Earlier Germanic changes}\par
\vspace{0.35em}
\raggedright
\centering\textbf{}\par
\raggedright
\vspace{0.2em}
\begin{tabular}{@{}>{\raggedright\arraybackslash}p{0.74\linewidth}@{\hspace{0.55em}}>{\raggedright\arraybackslash}p{0.16\linewidth}@{\hspace{0.25em}}}
\multicolumn{2}{@{}l@{}}{*Proto-West Germanic*} \\
\multicolumn{2}{@{}l@{}}{[no change]} \\
\multicolumn{2}{@{}l@{}}{*Northwest Germanic*} \\
\mbox{PGmc Final Z Deletion} & \emph{*súmara} \\
\end{tabular}
\end{minipage}
&
&
\begin{minipage}[t]{\linewidth}
\centering\textbf{Old English changes}\par
\vspace{0.35em}
\raggedright
\centering\textbf{}\par
\raggedright
\vspace{0.2em}
\begin{tabular}{@{}>{\raggedright\arraybackslash}p{0.70\linewidth}@{\hspace{0.55em}}>{\raggedright\arraybackslash}p{0.16\linewidth}@{\hspace{0.25em}}}
\multicolumn{2}{@{}l@{}}{*Old English*} \\
PWGmc Final Bare A Loss & \emph{*súmar} \\
Anglo Frisian Brightening & \emph{*súmær} \\
OE Unstressed AE Merger & \emph{*súmer} \\
\end{tabular}
\end{minipage}
\\
\end{tabularx}
\end{minipage}%
} \endgroup

Old English form: \emph{sumer}

\subsection{Reconstruction and comparative
evidence}\label{reconstruction-and-comparative-evidence-51}

Kroonen gives the lexeme as \emph{*sumara-}, and Ringe and Taylor
likewise use \Recon{sumaraz} `summer', while Orel preserves an alternate
\Recon{sumeraz} `summer' (\citeproc{ref-Kroonen2013}{Kroonen 2013};
\citeproc{ref-RingeTaylor2014}{Ringe and Taylor 2014};
\citeproc{ref-Orel2003}{Orel 2003, 425}). The derivational input
\Recon{súmaraz} `summer' follows the \emph{*a} vocalism that underlies
the regular development represented here.

\subsection{Old English evidence}\label{old-english-evidence-51}

Clark Hall gives {\emph{sumor}} `summer' m., gs. {\emph{sumeres}}
`summer', ds. {\emph{sumera}} `summer' / {\emph{sumere}} `summer', and
Bright likewise lists \emph{sumor (sumer)} with genitive
{\emph{sumeres}} `summer (gen.)' (\citeproc{ref-ClarkHall1960}{Clark
Hall 1960}; \citeproc{ref-BrightCassidyRingler1971}{Bright 1971, 440}).
The tradition therefore preserves both {\emph{sumor}} `summer' and
{\emph{sumer}} `summer', with the oblique forms strongly supporting
second-syllable \emph{e}.

\subsection{Development to Old
English}\label{development-to-old-english-51}

From \Recon{súmaraz} `summer', loss of final \emph{-az} is followed by
fronting and merger in the unstressed second syllable, yielding
{\emph{sumer}} `summer'. The form compared here is the regularized
\emph{e}-form, while the common citation form {\emph{sumor}} `summer'
remains part of the attested OE tradition
(\citeproc{ref-RingeTaylor2014}{Ringe and Taylor 2014}).

\subsection{Form note}\label{form-note-21}

The entry does not deny {\emph{sumor}} `summer'. It represents
{\emph{sumer}} `summer' as the regular outcome chosen here, while
{\emph{sumor}} `summer' remains a common headword spelling and
{\emph{sumeres}} `summer (gen.)' / {\emph{sumere}} `summer (dat.)' show
that the \emph{e}-vocalism was also real in Old English
(\citeproc{ref-ClarkHall1960}{Clark Hall 1960};
\citeproc{ref-BrightCassidyRingler1971}{Bright 1971, 440}).

\section{\texorpdfstring{sunder --- OE
\emph{sundrian}}{sunder --- OE sundrian}}\label{sunder-oe-sundrian}

\index[iv]{02oe@\textbf{Old English}!sundrian@\emph{sundrian}}
\index[iv]{03pgmc@\textbf{Proto-Germanic}!sundrojana@*súndrōjaną}

Derivation: \emph{*súndrōjaną} \textgreater{} \emph{sundrian} (regular).

\subsection{Derivation trace}\label{derivation-trace-52}

Proto input: \emph{*súndrōjaną}

\begingroup
\setlength{\fboxsep}{6pt}

\noindent\fbox{%
\begin{minipage}{0.97\linewidth}
\footnotesize
\begin{tabularx}{\linewidth}{@{}>{\raggedright\arraybackslash}p{0.320\linewidth}>{\centering\arraybackslash}X>{\raggedright\arraybackslash}p{0.520\linewidth}@{}}
\begin{minipage}[t]{\linewidth}
\centering\textbf{Earlier Germanic changes}\par
\vspace{0.35em}
\raggedright
\centering\textbf{}\par
\raggedright
\vspace{0.2em}
\begin{tabular}{@{}>{\raggedright\arraybackslash}p{0.68\linewidth}@{\hspace{0.55em}}>{\raggedright\arraybackslash}p{0.16\linewidth}@{\hspace{0.25em}}}
\multicolumn{2}{@{}l@{}}{*Proto-West Germanic*} \\
\multicolumn{2}{@{}l@{}}{[no change]} \\
\multicolumn{2}{@{}l@{}}{*Northwest Germanic*} \\
\multicolumn{2}{@{}l@{}}{[no change]} \\
\end{tabular}
\end{minipage}
&
&
\begin{minipage}[t]{\linewidth}
\centering\textbf{Old English changes}\par
\vspace{0.35em}
\raggedright
\centering\textbf{}\par
\raggedright
\vspace{0.2em}
\begin{tabular}{@{}>{\raggedright\arraybackslash}p{0.64\linewidth}@{\hspace{0.55em}}>{\raggedright\arraybackslash}p{0.26\linewidth}@{\hspace{0.25em}}}
\multicolumn{2}{@{}l@{}}{*Old English*} \\
OE Heavy Syllable Nasal Apocope & \emph{*súndrōjan} \\
OE Secondary Nasalization & \emph{*súndrōjąn} \\
OE I Umlaut & \emph{*súndrējąn} \\
OE Unstressed Long Vowel Shortening & \emph{*súndrejąn} \\
OE Weak Tail Reduction & \emph{*súndrejan} \\
OE Intervocalic J Vocalization & \emph{*súndreian} \\
OE Unstressed EI Contraction & \emph{*súndrian} \\
\end{tabular}
\end{minipage}
\\
\end{tabularx}
\end{minipage}%
} \endgroup

Old English form: \emph{sundrian}

\subsection{Reconstruction and comparative
evidence}\label{reconstruction-and-comparative-evidence-52}

Orel distinguishes three related formations: adverbial \Recon{sunþraz}
`asunder' \textgreater{} {\emph{sundor}} `apart, asunder', Class I
verbal \Recon{sunþrjanan} `sunder' \textgreater{} {\emph{syndrian}}
`separate, sunder', and Class II verbal \Recon{sunþrōjanan} `sunder'
\textgreater{} \emph{sundrian} `sunder' (\citeproc{ref-Orel2003}{Orel
2003}).

Kluge-Seebold aligns the cognate set with German {\emph{sondern}}
`separate' and OE {\emph{gesundrian}} `separate', so this entry belongs
with the Class II verb, not the adverb
(\citeproc{ref-KlugeSeebold2011}{Kluge 2011}).

\subsection{Old English evidence}\label{old-english-evidence-52}

Clark Hall and Bosworth-Toller list \emph{sundrian} `sunder' and
\emph{syndrian} `sunder' separately from adverbial \emph{sundor}
`asunder'; both also record the prefixed verbal family \emph{ā-sundrian}
`sunder apart' (\citeproc{ref-ClarkHall1960}{Clark Hall 1960, 296};
\citeproc{ref-BosworthToller1898}{Bosworth and Toller 1898}). The target
is therefore the weak verb \emph{sundrian} `sunder'.

\subsection{Development to Old
English}\label{development-to-old-english-52}

From \Recon{súndrōjaną} `sunder', the Class II weak-verb suffix yields
regular OE \emph{-ian}, producing \emph{sundrian} `sunder'. Because this
is the \emph{*-ōjan-} verb and not the Class I \emph{*-jan-} formation,
the form represented here does not belong to the umlauted
{\emph{syndrian}} `separate' branch.

\subsection{Form note}\label{form-note-22}

The earlier confusion was lexical, not phonological: {\emph{sundor}}
`apart' is the separate adverb, and {\emph{syndrian}} `separate' is a
related but different verb. The verb treated here is the Class II verb
\emph{sundrian} (\citeproc{ref-Orel2003}{Orel 2003};
\citeproc{ref-ClarkHall1960}{Clark Hall 1960, 296}).

\section{\texorpdfstring{swallow --- OE
\emph{swealwe}}{swallow --- OE swealwe}}\label{swallow-oe-swealwe}

\index[iv]{02oe@\textbf{Old English}!swealwe@\emph{swealwe}}
\index[iv]{03pgmc@\textbf{Proto-Germanic}!swalwon@*swálwōn}

Derivation: \emph{*swálwōn} \textgreater{} \emph{swealwe} (regular).

\subsection{Derivation trace}\label{derivation-trace-53}

Proto input: \emph{*swálwōn}

\begingroup
\setlength{\fboxsep}{6pt}

\noindent\fbox{%
\begin{minipage}{0.97\linewidth}
\small
\begin{tabularx}{\linewidth}{@{}>{\raggedright\arraybackslash}p{0.440\linewidth}>{\centering\arraybackslash}X>{\raggedright\arraybackslash}p{0.440\linewidth}@{}}
\begin{minipage}[t]{\linewidth}
\centering\textbf{Earlier Germanic changes}\par
\vspace{0.35em}
\raggedright
\centering\textbf{}\par
\raggedright
\vspace{0.2em}
\begin{tabular}{@{}>{\raggedright\arraybackslash}p{0.68\linewidth}@{\hspace{0.55em}}>{\raggedright\arraybackslash}p{0.16\linewidth}@{\hspace{0.25em}}}
\multicolumn{2}{@{}l@{}}{*Proto-West Germanic*} \\
\multicolumn{2}{@{}l@{}}{[no change]} \\
\multicolumn{2}{@{}l@{}}{*Northwest Germanic*} \\
NWGmc N Stem N Loss & \emph{*swálwǭ} \\
\end{tabular}
\end{minipage}
&
&
\begin{minipage}[t]{\linewidth}
\centering\textbf{Old English changes}\par
\vspace{0.35em}
\raggedright
\centering\textbf{}\par
\raggedright
\vspace{0.2em}
\begin{tabular}{@{}>{\raggedright\arraybackslash}p{0.70\linewidth}@{\hspace{0.55em}}>{\raggedright\arraybackslash}p{0.20\linewidth}@{\hspace{0.25em}}}
\multicolumn{2}{@{}l@{}}{*Old English*} \\
Anglo Frisian Brightening & \emph{*swælwǭ} \\
OE Breaking & \emph{*swealwǭ} \\
OE Unstressed Long Vowel Shortening & \emph{*swealwæ} \\
OE Unstressed AE Merger & \emph{*swealwe} \\
\end{tabular}
\end{minipage}
\\
\end{tabularx}
\end{minipage}%
} \endgroup

Old English form: \emph{swealwe}

\subsection{Reconstruction and comparative
evidence}\label{reconstruction-and-comparative-evidence-53}

Kroonen gives the bird name as \emph{*swalwōn-}, and Ringe and Taylor
cite the later West Germanic stage \Recon{swalwa} `swallow', from which
West Saxon {\emph{swealwe}} `swallow' and Mercian {\emph{swalwe}}
`swallow' develop (\citeproc{ref-Kroonen2013}{Kroonen 2013, 535};
\citeproc{ref-RingeTaylor2014}{Ringe and Taylor 2014, 200}). The
selected etymological comparison belongs to the swallow-bird family, not
to the verb \emph{swelgan}.

\subsection{Old English evidence}\label{old-english-evidence-53}

Clark Hall records \emph{swealwe (a, o)} as the noun headword
(\citeproc{ref-ClarkHall1960}{Clark Hall 1960}). Campbell and Brunner
also preserve later or oblique-family forms such as \emph{swaluwe},
\emph{swalewan}, and \emph{swealuwe}, but those belong to wider
variation around the noun rather than to the citation form represented
here (\citeproc{ref-Campbell1959}{Campbell 1959};
\citeproc{ref-SieversBrunner1965}{Brunner 1965}).

\subsection{Development to Old
English}\label{development-to-old-english-53}

From \Recon{swálwōn} `swallow', brightening yields \emph{*swælw-}, and
breaking before \emph{lw} gives \emph{*swealw-}. The later noun ending
develops regularly to \emph{swealwe}. The relevant point is that the
bird name has no inherited \emph{*g}: that consonant belongs to the
separate verb \emph{swelgan} (\citeproc{ref-RingeTaylor2014}{Ringe and
Taylor 2014, 200}; \citeproc{ref-Kroonen2013}{Kroonen 2013, 535}).

\section{\texorpdfstring{swine --- OE
\emph{swīn}}{swine --- OE swīn}}\label{swine-oe-swux12bn}

\index[iv]{02oe@\textbf{Old English}!swin@\emph{swīn}}
\index[iv]{03pgmc@\textbf{Proto-Germanic}!swina@*swī́ną}
\index[iv]{03pgmc@\textbf{Proto-Germanic}!swina@*swḯną}

Derivation: citation reconstruction \emph{*swī́ną}; form followed here
\emph{*swī́ną} \textgreater{} \emph{swīn} (regular).

\subsection{Derivation trace}\label{derivation-trace-54}

Proto input: \emph{*swī́ną}

\begingroup
\setlength{\fboxsep}{6pt}

\noindent\fbox{%
\begin{minipage}{0.97\linewidth}
\small
\begin{tabularx}{\linewidth}{@{}>{\raggedright\arraybackslash}p{0.440\linewidth}>{\centering\arraybackslash}X>{\raggedright\arraybackslash}p{0.440\linewidth}@{}}
\begin{minipage}[t]{\linewidth}
\centering\textbf{Earlier Germanic changes}\par
\vspace{0.35em}
\raggedright
\centering\textbf{}\par
\raggedright
\vspace{0.2em}
\begin{tabular}{@{}>{\raggedright\arraybackslash}p{0.68\linewidth}@{\hspace{0.55em}}>{\raggedright\arraybackslash}p{0.16\linewidth}@{\hspace{0.25em}}}
\multicolumn{2}{@{}l@{}}{*Proto-West Germanic*} \\
\multicolumn{2}{@{}l@{}}{[no change]} \\
\multicolumn{2}{@{}l@{}}{*Northwest Germanic*} \\
\multicolumn{2}{@{}l@{}}{[no change]} \\
\end{tabular}
\end{minipage}
&
&
\begin{minipage}[t]{\linewidth}
\centering\textbf{Old English changes}\par
\vspace{0.35em}
\raggedright
\centering\textbf{}\par
\raggedright
\vspace{0.2em}
\begin{tabular}{@{}>{\raggedright\arraybackslash}p{0.74\linewidth}@{\hspace{0.55em}}>{\raggedright\arraybackslash}p{0.16\linewidth}@{\hspace{0.25em}}}
\multicolumn{2}{@{}l@{}}{*Old English*} \\
OE Heavy Syllable Nasal Apocope & \emph{*swī́n} \\
\end{tabular}
\end{minipage}
\\
\end{tabularx}
\end{minipage}%
} \endgroup

Old English form: \emph{swīn}

\subsection{Reconstruction and comparative
evidence}\label{reconstruction-and-comparative-evidence-54}

Kroonen cites the noun as \emph{*swina-} `pig'
(\citeproc{ref-Kroonen2013}{Kroonen 2013}). The selected comparative
form here is the bare neuter citation cell \Recon{swī́ną} `swine', which
matches the singular noun aimed at in Old English rather than an oblique
stem form.

\subsection{Old English evidence}\label{old-english-evidence-54}

Clark Hall records \emph{swin (y)} as the ordinary noun headword
(\citeproc{ref-ClarkHall1960}{Clark Hall 1960}). The target here is that
singular citation form \emph{swīn} `swine', not a plural glossed in
Modern English as \emph{swine}.

\subsection{Development to Old
English}\label{development-to-old-english-54}

From derivational input \Recon{swī́ną} `swine', loss of the final nasal
vowel yields {\emph{swīn}} `swine, pig'. The outcome is therefore the
regular monosyllabic noun with preserved long root \emph{ī}.

\subsection{Source note}\label{source-note-4}

The derivational input writes stressed long \emph{ī} as \emph{*ī́}, so
comparative \Recon{swī́ną} `swine' and derivational \Recon{swī́ną} `swine'
represent the same lexical form.

\section{\texorpdfstring{think --- OE
\emph{þenċan}}{think --- OE þenċan}}\label{think-oe-uxfeenux10ban}

\index[iv]{02oe@\textbf{Old English}!thencan@\emph{þenċan}}
\index[iv]{03pgmc@\textbf{Proto-Germanic}!thankijana@*θánkijaną}

Derivation: \emph{*θánkijaną} \textgreater{} \emph{þenċan} (regular).

\subsection{Derivation trace}\label{derivation-trace-55}

Proto input: \emph{*θánkijaną}

\begingroup
\setlength{\fboxsep}{6pt}

\noindent\fbox{%
\begin{minipage}{0.97\linewidth}
\footnotesize
\begin{tabularx}{\linewidth}{@{}>{\raggedright\arraybackslash}p{0.320\linewidth}>{\centering\arraybackslash}X>{\raggedright\arraybackslash}p{0.520\linewidth}@{}}
\begin{minipage}[t]{\linewidth}
\centering\textbf{Earlier Germanic changes}\par
\vspace{0.35em}
\raggedright
\centering\textbf{}\par
\raggedright
\vspace{0.2em}
\begin{tabular}{@{}>{\raggedright\arraybackslash}p{0.68\linewidth}@{\hspace{0.55em}}>{\raggedright\arraybackslash}p{0.16\linewidth}@{\hspace{0.25em}}}
\multicolumn{2}{@{}l@{}}{*Proto-West Germanic*} \\
\multicolumn{2}{@{}l@{}}{[no change]} \\
\multicolumn{2}{@{}l@{}}{*Northwest Germanic*} \\
\multicolumn{2}{@{}l@{}}{[no change]} \\
\end{tabular}
\end{minipage}
&
&
\begin{minipage}[t]{\linewidth}
\centering\textbf{Old English changes}\par
\vspace{0.35em}
\raggedright
\centering\textbf{}\par
\raggedright
\vspace{0.2em}
\begin{tabular}{@{}>{\raggedright\arraybackslash}p{0.66\linewidth}@{\hspace{0.55em}}>{\raggedright\arraybackslash}p{0.24\linewidth}@{\hspace{0.25em}}}
\multicolumn{2}{@{}l@{}}{*Old English*} \\
OE Heavy Syllable Nasal Apocope & \emph{*θánkijan} \\
OE Secondary Nasalization & \emph{*θánkijąn} \\
Sievers Law Syncope & \emph{*θánkjąn} \\
OE Velar Palatalization & \emph{*θánʧjąn} \\
OE I Umlaut & \emph{*θenʧjąn} \\
OE Weak Tail Reduction & \emph{*θenʧjan} \\
OE J Loss After Heavy & \emph{*θenʧan} \\
\end{tabular}
\end{minipage}
\\
\end{tabularx}
\end{minipage}%
} \endgroup

Old English form: \emph{þenċan}

\subsection{Reconstruction and comparative
evidence}\label{reconstruction-and-comparative-evidence-55}

Kroonen gives the verb as \emph{*þankjan-} `to think', and Ringe and
Taylor cite fully inflected \Recon{þankijaną} `think' beside OE
\emph{þenċan} `think' (\citeproc{ref-Kroonen2013}{Kroonen 2013};
\citeproc{ref-RingeTaylor2014}{Ringe and Taylor 2014}). The noun
\Recon{þankaz} `think' belongs only to the wider derivational
background.

\subsection{Old English evidence}\label{old-english-evidence-55}

Bosworth-Toller preserves the verb under \emph{þencan} `think' /
\emph{geþencan} `think', and the citation form here is the ordinary
infinitive \emph{þenċan} `think'
(\citeproc{ref-BosworthToller1898}{Bosworth and Toller 1898}).

\subsection{Development to Old
English}\label{development-to-old-english-55}

From \Recon{θánkijaną} `think', palatalization before \emph{*j} and
i-umlaut produce \emph{þenċan} `think'. The infinitive is therefore a
straightforward weak-verb outcome.

\subsection{Lexical note}\label{lexical-note}

Campbell's assibilation discussion uses the same verb \emph{þencan}
`think'; the class-III relic \emph{hycgan} is a different lexeme
(\citeproc{ref-Campbell1959}{Campbell 1959};
\citeproc{ref-Hogg1992}{Hogg 1992}).

\section{\texorpdfstring{thorn --- OE
\emph{þorn}}{thorn --- OE þorn}}\label{thorn-oe-uxfeorn}

\index[iv]{02oe@\textbf{Old English}!thorn@\emph{þorn}}
\index[iv]{03pgmc@\textbf{Proto-Germanic}!thurnaz@*θúrnaz}

Derivation: \emph{*θúrnaz} \textgreater{} \emph{þorn} (regular).

\subsection{Derivation trace}\label{derivation-trace-56}

Proto input: \emph{*θúrnaz}

\begingroup
\setlength{\fboxsep}{6pt}

\noindent\fbox{%
\begin{minipage}{0.97\linewidth}
\small
\begin{tabularx}{\linewidth}{@{}>{\raggedright\arraybackslash}p{0.520\linewidth}>{\centering\arraybackslash}X>{\raggedright\arraybackslash}p{0.320\linewidth}@{}}
\begin{minipage}[t]{\linewidth}
\centering\textbf{Earlier Germanic changes}\par
\vspace{0.35em}
\raggedright
\centering\textbf{}\par
\raggedright
\vspace{0.2em}
\begin{tabular}{@{}>{\raggedright\arraybackslash}p{0.74\linewidth}@{\hspace{0.55em}}>{\raggedright\arraybackslash}p{0.16\linewidth}@{\hspace{0.25em}}}
\multicolumn{2}{@{}l@{}}{*Proto-West Germanic*} \\
\multicolumn{2}{@{}l@{}}{[no change]} \\
\multicolumn{2}{@{}l@{}}{*Northwest Germanic*} \\
\mbox{NWGmc U Lowering} & \emph{*θórnaz} \\
\mbox{PGmc Final Z Deletion} & \emph{*θórna} \\
\end{tabular}
\end{minipage}
&
&
\begin{minipage}[t]{\linewidth}
\centering\textbf{Old English changes}\par
\vspace{0.35em}
\raggedright
\centering\textbf{}\par
\raggedright
\vspace{0.2em}
\begin{tabular}{@{}>{\raggedright\arraybackslash}p{0.70\linewidth}@{\hspace{0.55em}}>{\raggedright\arraybackslash}p{0.16\linewidth}@{\hspace{0.25em}}}
\multicolumn{2}{@{}l@{}}{*Old English*} \\
PWGmc Final Bare A Loss & \emph{*θórn} \\
\end{tabular}
\end{minipage}
\\
\end{tabularx}
\end{minipage}%
} \endgroup

Old English form: \emph{þorn}

\subsection{Reconstruction and comparative
evidence}\label{reconstruction-and-comparative-evidence-56}

Kroonen gives \emph{*þurna-} `thorn, briar', while Orel preserves the
masculine pair \Recon{þurnuz} `thorn' and \Recon{þurnaz} `thorn'
(\citeproc{ref-Kroonen2013}{Kroonen 2013}; \citeproc{ref-Orel2003}{Orel
2003}). The derivational input \Recon{θúrnaz} `thorn' belongs to that
same comparative family.

\subsection{Old English evidence}\label{old-english-evidence-56}

Bright lists \emph{þorn} `thorn' (m.) and Clark Hall likewise treats
\emph{þorn} `thorn' as the ordinary noun headword
(\citeproc{ref-BrightCassidyRingler1971}{Bright 1971};
\citeproc{ref-ClarkHall1960}{Clark Hall 1960}).

\subsection{Development to Old
English}\label{development-to-old-english-56}

The inherited stem shows regular lowering of \emph{u} to \emph{o} before
\emph{r}, and final loss yields \emph{þorn} `thorn'. The noun is
therefore a regular Old English continuation of the Proto-Germanic
thorn-family.

\subsection{Source note}\label{source-note-5}

The comparative sources preserve more than one stem formation, but the
Old English target itself is simply the citation form \emph{þorn}
`thorn'.

\section{\texorpdfstring{tide --- OE
\emph{tīd}}{tide --- OE tīd}}\label{tide-oe-tux12bd}

\index[iv]{02oe@\textbf{Old English}!tid@\emph{tīd}}
\index[iv]{03pgmc@\textbf{Proto-Germanic}!tidiz@*tī́diz}
\index[iv]{03pgmc@\textbf{Proto-Germanic}!tidiz@*tḯdiz}

Derivation: citation reconstruction \emph{*tī́diz}; form followed here
\emph{*tī́diz} \textgreater{} \emph{tīd} (regular).

\subsection{Derivation trace}\label{derivation-trace-57}

Proto input: \emph{*tī́diz}

\begingroup
\setlength{\fboxsep}{6pt}

\noindent\fbox{%
\begin{minipage}{0.97\linewidth}
\small
\begin{tabularx}{\linewidth}{@{}>{\raggedright\arraybackslash}p{0.460\linewidth}>{\centering\arraybackslash}X>{\raggedright\arraybackslash}p{0.460\linewidth}@{}}
\begin{minipage}[t]{\linewidth}
\centering\textbf{Earlier Germanic changes}\par
\vspace{0.35em}
\raggedright
\centering\textbf{}\par
\raggedright
\vspace{0.2em}
\begin{tabular}{@{}>{\raggedright\arraybackslash}p{0.74\linewidth}@{\hspace{0.55em}}>{\raggedright\arraybackslash}p{0.16\linewidth}@{\hspace{0.25em}}}
\multicolumn{2}{@{}l@{}}{*Proto-West Germanic*} \\
\multicolumn{2}{@{}l@{}}{[no change]} \\
\multicolumn{2}{@{}l@{}}{*Northwest Germanic*} \\
\mbox{PGmc Final Z Deletion} & \emph{*tī́di} \\
\end{tabular}
\end{minipage}
&
&
\begin{minipage}[t]{\linewidth}
\centering\textbf{Old English changes}\par
\vspace{0.35em}
\raggedright
\centering\textbf{}\par
\raggedright
\vspace{0.2em}
\begin{tabular}{@{}>{\raggedright\arraybackslash}p{0.74\linewidth}@{\hspace{0.55em}}>{\raggedright\arraybackslash}p{0.16\linewidth}@{\hspace{0.25em}}}
\multicolumn{2}{@{}l@{}}{*Old English*} \\
\mbox{OE High Vowel Apocope} & \emph{*tī́d} \\
\end{tabular}
\end{minipage}
\\
\end{tabularx}
\end{minipage}%
} \endgroup

Old English form: \emph{tīd}

\subsection{Reconstruction and comparative
evidence}\label{reconstruction-and-comparative-evidence-57}

Kroonen gives \emph{*tīdi-} `time', and Orel's \Recon{tīđiz} `tide'
points to the same feminine noun (\citeproc{ref-Kroonen2013}{Kroonen
2013}; \citeproc{ref-Orel2003}{Orel 2003}). The related verb
\emph{tīdan} `happen' is separate.

\subsection{Old English evidence}\label{old-english-evidence-57}

Bright records \emph{tīd} `tide' with singular \emph{tīde} `tide' and
plural \emph{tīda} `tides', and Clark Hall treats \emph{tīd} `tide' as
the ordinary noun `time, period, season'
(\citeproc{ref-BrightCassidyRingler1971}{Bright 1971};
\citeproc{ref-ClarkHall1960}{Clark Hall 1960, 309}).

\subsection{Development to Old
English}\label{development-to-old-english-57}

From \Recon{tī́diz} `tide', final \emph{z} is lost and the high final
vowel drops, leaving \emph{tīd} `tide'. The development is
straightforward for a feminine i-stem.

\subsection{Lexical note}\label{lexical-note-1}

English \emph{tide} can also lead to the separate weak verb \emph{tīdan}
`happen'; the noun is \emph{tīd} `tide'.

\section{\texorpdfstring{token --- OE
\emph{tācn}}{token --- OE tācn}}\label{token-oe-tux101cn}

\index[iv]{02oe@\textbf{Old English}!tacn@\emph{tācn}}
\index[iv]{03pgmc@\textbf{Proto-Germanic}!taikna@*táikną}

Derivation: \emph{*táikną} \textgreater{} \emph{tācn} (regular).

\subsection{Derivation trace}\label{derivation-trace-58}

Proto input: \emph{*táikną}

\begingroup
\setlength{\fboxsep}{6pt}

\noindent\fbox{%
\begin{minipage}{0.97\linewidth}
\small
\begin{tabularx}{\linewidth}{@{}>{\raggedright\arraybackslash}p{0.440\linewidth}>{\centering\arraybackslash}X>{\raggedright\arraybackslash}p{0.440\linewidth}@{}}
\begin{minipage}[t]{\linewidth}
\centering\textbf{Earlier Germanic changes}\par
\vspace{0.35em}
\raggedright
\centering\textbf{}\par
\raggedright
\vspace{0.2em}
\begin{tabular}{@{}>{\raggedright\arraybackslash}p{0.70\linewidth}@{\hspace{0.55em}}>{\raggedright\arraybackslash}p{0.16\linewidth}@{\hspace{0.25em}}}
\multicolumn{2}{@{}l@{}}{*Proto-West Germanic*} \\
PWGmc Ai Monophthongization & \emph{*tākną} \\
\multicolumn{2}{@{}l@{}}{*Northwest Germanic*} \\
\multicolumn{2}{@{}l@{}}{[no change]} \\
\end{tabular}
\end{minipage}
&
&
\begin{minipage}[t]{\linewidth}
\centering\textbf{Old English changes}\par
\vspace{0.35em}
\raggedright
\centering\textbf{}\par
\raggedright
\vspace{0.2em}
\begin{tabular}{@{}>{\raggedright\arraybackslash}p{0.74\linewidth}@{\hspace{0.55em}}>{\raggedright\arraybackslash}p{0.16\linewidth}@{\hspace{0.25em}}}
\multicolumn{2}{@{}l@{}}{*Old English*} \\
OE Heavy Syllable Nasal Apocope & \emph{*tākn} \\
\end{tabular}
\end{minipage}
\\
\end{tabularx}
\end{minipage}%
} \endgroup

Old English form: \emph{tācn}

\subsection{Reconstruction and comparative
evidence}\label{reconstruction-and-comparative-evidence-58}

Kroonen cites \emph{*taikna-} and Orel \Recon{taiknan} `token' for the
noun `sign, token' (\citeproc{ref-Kroonen2013}{Kroonen 2013};
\citeproc{ref-Orel2003}{Orel 2003, 438}). The derivational input
\Recon{táikną} `token' is the simple citation-form noun used for the
derivation.

\subsection{Old English evidence}\label{old-english-evidence-58}

Campbell and Sievers-Brunner preserve both unbroken \emph{tācn} `token'
and broken \emph{tācen} `token', with oblique \emph{tācnes} `token'
remaining unbroken (\citeproc{ref-Campbell1959}{Campbell 1959};
\citeproc{ref-SieversBrunner1965}{Brunner 1965}).

\subsection{Development to Old
English}\label{development-to-old-english-58}

Monophthongization of \emph{ai} yields \emph{ā}, and loss of the final
nasal vowel leaves \emph{tācn} `token'. The unbroken form is therefore a
regular Old English outcome.

\subsection{Form note}\label{form-note-23}

\emph{tācn} `token' is the attested unbroken citation form selected
here. Later West Saxon prose often prefers \emph{tācen} `token', but
that does not displace the older unbroken form
(\citeproc{ref-Campbell1959}{Campbell 1959};
\citeproc{ref-SieversBrunner1965}{Brunner 1965}).

\section{\texorpdfstring{town --- OE
\emph{tūn}}{town --- OE tūn}}\label{town-oe-tux16bn}

\index[iv]{02oe@\textbf{Old English}!tun@\emph{tūn}}
\index[iv]{03pgmc@\textbf{Proto-Germanic}!tuna@*tūną}

Derivation: \emph{*tūną} \textgreater{} \emph{tūn} (regular).

\subsection{Derivation trace}\label{derivation-trace-59}

Proto input: \emph{*tūną}

\begingroup
\setlength{\fboxsep}{6pt}

\noindent\fbox{%
\begin{minipage}{0.97\linewidth}
\small
\begin{tabularx}{\linewidth}{@{}>{\raggedright\arraybackslash}p{0.440\linewidth}>{\centering\arraybackslash}X>{\raggedright\arraybackslash}p{0.440\linewidth}@{}}
\begin{minipage}[t]{\linewidth}
\centering\textbf{Earlier Germanic changes}\par
\vspace{0.35em}
\raggedright
\centering\textbf{}\par
\raggedright
\vspace{0.2em}
\begin{tabular}{@{}>{\raggedright\arraybackslash}p{0.68\linewidth}@{\hspace{0.55em}}>{\raggedright\arraybackslash}p{0.16\linewidth}@{\hspace{0.25em}}}
\multicolumn{2}{@{}l@{}}{*Proto-West Germanic*} \\
\multicolumn{2}{@{}l@{}}{[no change]} \\
\multicolumn{2}{@{}l@{}}{*Northwest Germanic*} \\
\multicolumn{2}{@{}l@{}}{[no change]} \\
\end{tabular}
\end{minipage}
&
&
\begin{minipage}[t]{\linewidth}
\centering\textbf{Old English changes}\par
\vspace{0.35em}
\raggedright
\centering\textbf{}\par
\raggedright
\vspace{0.2em}
\begin{tabular}{@{}>{\raggedright\arraybackslash}p{0.74\linewidth}@{\hspace{0.55em}}>{\raggedright\arraybackslash}p{0.16\linewidth}@{\hspace{0.25em}}}
\multicolumn{2}{@{}l@{}}{*Old English*} \\
OE Heavy Syllable Nasal Apocope & \emph{*tūn} \\
\end{tabular}
\end{minipage}
\\
\end{tabularx}
\end{minipage}%
} \endgroup

Old English form: \emph{tūn}

\subsection{Reconstruction and comparative
evidence}\label{reconstruction-and-comparative-evidence-59}

Kroonen cites \emph{*tūna-} `fenced area', while Orel gives
\Recon{tūnan} `town' and \Recon{tūnaz} `town'
(\citeproc{ref-Kroonen2013}{Kroonen 2013}; \citeproc{ref-Orel2003}{Orel
2003, 452}). The derivational input \Recon{tūną} `town' is the simple
citation-form noun used in the derivation.

\subsection{Old English evidence}\label{old-english-evidence-59}

Clark Hall records \emph{tūn} `town' as the ordinary headword
`enclosure, yard, village, town' (\citeproc{ref-ClarkHall1960}{Clark
Hall 1960}).

\subsection{Development to Old
English}\label{development-to-old-english-59}

The inherited long \emph{ū} is preserved, and loss of the final nasal
vowel yields \emph{tūn} `town' regularly
(\citeproc{ref-SieversBrunner1965}{Brunner 1965}).

\subsection{Source note}\label{source-note-6}

The comparative headwords vary, but the Old English target here is the
direct citation form \emph{tūn} `town', not an oblique \Recon{tūnăn}
`town'.

\section{\texorpdfstring{wade --- OE
\emph{wadan}}{wade --- OE wadan}}\label{wade-oe-wadan}

\index[iv]{02oe@\textbf{Old English}!wadan@\emph{wadan}}
\index[iv]{03pgmc@\textbf{Proto-Germanic}!wadana@*wádaną}

Derivation: \emph{*wádaną} \textgreater{} \emph{wadan} (regular).

\subsection{Derivation trace}\label{derivation-trace-60}

Proto input: \emph{*wádaną}

\begingroup
\setlength{\fboxsep}{6pt}

\noindent\fbox{%
\begin{minipage}{0.97\linewidth}
\small
\begin{tabularx}{\linewidth}{@{}>{\raggedright\arraybackslash}p{0.320\linewidth}>{\centering\arraybackslash}X>{\raggedright\arraybackslash}p{0.520\linewidth}@{}}
\begin{minipage}[t]{\linewidth}
\centering\textbf{Earlier Germanic changes}\par
\vspace{0.35em}
\raggedright
\centering\textbf{}\par
\raggedright
\vspace{0.2em}
\begin{tabular}{@{}>{\raggedright\arraybackslash}p{0.68\linewidth}@{\hspace{0.55em}}>{\raggedright\arraybackslash}p{0.16\linewidth}@{\hspace{0.25em}}}
\multicolumn{2}{@{}l@{}}{*Proto-West Germanic*} \\
\multicolumn{2}{@{}l@{}}{[no change]} \\
\multicolumn{2}{@{}l@{}}{*Northwest Germanic*} \\
\multicolumn{2}{@{}l@{}}{[no change]} \\
\end{tabular}
\end{minipage}
&
&
\begin{minipage}[t]{\linewidth}
\centering\textbf{Old English changes}\par
\vspace{0.35em}
\raggedright
\centering\textbf{}\par
\raggedright
\vspace{0.2em}
\begin{tabular}{@{}>{\raggedright\arraybackslash}p{0.74\linewidth}@{\hspace{0.55em}}>{\raggedright\arraybackslash}p{0.16\linewidth}@{\hspace{0.25em}}}
\multicolumn{2}{@{}l@{}}{*Old English*} \\
Anglo Frisian Brightening & \emph{*wædaną} \\
OE A Restoration & \emph{*wadaną} \\
OE Heavy Syllable Nasal Apocope & \emph{*wadan} \\
OE Secondary Nasalization & \emph{*wadąn} \\
OE Weak Tail Reduction & \emph{*wadan} \\
\end{tabular}
\end{minipage}
\\
\end{tabularx}
\end{minipage}%
} \endgroup

Old English form: \emph{wadan}

\subsection{Reconstruction and comparative
evidence}\label{reconstruction-and-comparative-evidence-60}

Campbell and Ringe and Taylor describe A-restoration before a following
back vowel, and Luick explicitly includes {\emph{wadan}} `wade, go
through' among the standard open-syllable examples
(\citeproc{ref-Campbell1959}{Campbell 1959};
\citeproc{ref-RingeTaylor2014}{Ringe and Taylor 2014};
\citeproc{ref-Luick1914}{Luick 1914-\/-1940, 239}).

\subsection{Old English evidence}\label{old-english-evidence-60}

Clark Hall gives {\emph{wadan}} `wade, go through' as the verb `to go,
move, stride, advance', and Bright lists the same infinitive in the
strong-verb paradigm (\citeproc{ref-ClarkHall1960}{Clark Hall 1960};
\citeproc{ref-BrightCassidyRingler1971}{Bright 1971}).

\subsection{Development to Old
English}\label{development-to-old-english-60}

From \Recon{wádaną} `wade', Anglo-Frisian brightening first gives
\Recon{wædaną} `wade'. A-restoration before the back-vocalic infinitive
ending then returns \emph{a}, and later reduction yields {\emph{wadan}}
`wade, go through' (\citeproc{ref-Campbell1959}{Campbell 1959};
\citeproc{ref-RingeTaylor2014}{Ringe and Taylor 2014}).

\subsection{Development note}\label{development-note-2}

The infinitive belongs to the A-restoration class. The citation form is
therefore \emph{wadan} `wade', not a fronted \emph{wæden} `waded'-type
output.

\section{\texorpdfstring{warp --- OE
\emph{weorpan}}{warp --- OE weorpan}}\label{warp-oe-weorpan}

\index[iv]{02oe@\textbf{Old English}!weorpan@\emph{weorpan}}
\index[iv]{03pgmc@\textbf{Proto-Germanic}!werpana@*wérpaną}

Derivation: \emph{*wérpaną} \textgreater{} \emph{weorpan} (regular).

\subsection{Derivation trace}\label{derivation-trace-61}

Proto input: \emph{*wérpaną}

\begingroup
\setlength{\fboxsep}{6pt}

\noindent\fbox{%
\begin{minipage}{0.97\linewidth}
\small
\begin{tabularx}{\linewidth}{@{}>{\raggedright\arraybackslash}p{0.320\linewidth}>{\centering\arraybackslash}X>{\raggedright\arraybackslash}p{0.520\linewidth}@{}}
\begin{minipage}[t]{\linewidth}
\centering\textbf{Earlier Germanic changes}\par
\vspace{0.35em}
\raggedright
\centering\textbf{}\par
\raggedright
\vspace{0.2em}
\begin{tabular}{@{}>{\raggedright\arraybackslash}p{0.68\linewidth}@{\hspace{0.55em}}>{\raggedright\arraybackslash}p{0.16\linewidth}@{\hspace{0.25em}}}
\multicolumn{2}{@{}l@{}}{*Proto-West Germanic*} \\
\multicolumn{2}{@{}l@{}}{[no change]} \\
\multicolumn{2}{@{}l@{}}{*Northwest Germanic*} \\
\multicolumn{2}{@{}l@{}}{[no change]} \\
\end{tabular}
\end{minipage}
&
&
\begin{minipage}[t]{\linewidth}
\centering\textbf{Old English changes}\par
\vspace{0.35em}
\raggedright
\centering\textbf{}\par
\raggedright
\vspace{0.2em}
\begin{tabular}{@{}>{\raggedright\arraybackslash}p{0.68\linewidth}@{\hspace{0.55em}}>{\raggedright\arraybackslash}p{0.22\linewidth}@{\hspace{0.25em}}}
\multicolumn{2}{@{}l@{}}{*Old English*} \\
OE Breaking & \emph{*wéorpaną} \\
OE Heavy Syllable Nasal Apocope & \emph{*wéorpan} \\
OE Secondary Nasalization & \emph{*wéorpąn} \\
OE Weak Tail Reduction & \emph{*wéorpan} \\
\end{tabular}
\end{minipage}
\\
\end{tabularx}
\end{minipage}%
} \endgroup

Old English form: \emph{weorpan}

\subsection{Reconstruction and comparative
evidence}\label{reconstruction-and-comparative-evidence-61}

Ringe and Taylor distinguish preterite \Recon{warp} `warp' from
infinitive \Recon{werpana} `warp', and the derivational input here is
the verbal form \Recon{wérpaną} `warp'
(\citeproc{ref-RingeTaylor2014}{Ringe and Taylor 2014}).

\subsection{Old English evidence}\label{old-english-evidence-61}

Clark Hall records \emph{weorpan} as the strong verb headword and
separately lists \emph{wearp} as both noun and preterite. Bright gives
the paradigm \emph{weorpan, wearp, wurpon, worpen}
(\citeproc{ref-ClarkHall1960}{Clark Hall 1960};
\citeproc{ref-BrightCassidyRingler1971}{Bright 1971}).

\subsection{Development to Old
English}\label{development-to-old-english-61}

Breaking before \emph{r + C} yields \emph{weor-}, and the infinitive
develops regularly to \emph{weorpan}
(\citeproc{ref-Campbell1959}{Campbell 1959};
\citeproc{ref-Hogg1992}{Hogg 1992}).

\subsection{Lexical note}\label{lexical-note-2}

English \emph{warp} also points to related \emph{wearp} material. The
target is the infinitive \emph{weorpan}.

\section{\texorpdfstring{wash --- OE
\emph{wascan}}{wash --- OE wascan}}\label{wash-oe-wascan}

\index[iv]{02oe@\textbf{Old English}!wascan@\emph{wascan}}
\index[iv]{03pgmc@\textbf{Proto-Germanic}!waskana@*wáskaną}

Derivation: \emph{*wáskaną} \textgreater{} \emph{wascan} (regular).

\subsection{Derivation trace}\label{derivation-trace-62}

Proto input: \emph{*wáskaną}

\begingroup
\setlength{\fboxsep}{6pt}

\noindent\fbox{%
\begin{minipage}{0.97\linewidth}
\small
\begin{tabularx}{\linewidth}{@{}>{\raggedright\arraybackslash}p{0.320\linewidth}>{\centering\arraybackslash}X>{\raggedright\arraybackslash}p{0.520\linewidth}@{}}
\begin{minipage}[t]{\linewidth}
\centering\textbf{Earlier Germanic changes}\par
\vspace{0.35em}
\raggedright
\centering\textbf{}\par
\raggedright
\vspace{0.2em}
\begin{tabular}{@{}>{\raggedright\arraybackslash}p{0.68\linewidth}@{\hspace{0.55em}}>{\raggedright\arraybackslash}p{0.16\linewidth}@{\hspace{0.25em}}}
\multicolumn{2}{@{}l@{}}{*Proto-West Germanic*} \\
\multicolumn{2}{@{}l@{}}{[no change]} \\
\multicolumn{2}{@{}l@{}}{*Northwest Germanic*} \\
\multicolumn{2}{@{}l@{}}{[no change]} \\
\end{tabular}
\end{minipage}
&
&
\begin{minipage}[t]{\linewidth}
\centering\textbf{Old English changes}\par
\vspace{0.35em}
\raggedright
\centering\textbf{}\par
\raggedright
\vspace{0.2em}
\begin{tabular}{@{}>{\raggedright\arraybackslash}p{0.70\linewidth}@{\hspace{0.55em}}>{\raggedright\arraybackslash}p{0.20\linewidth}@{\hspace{0.25em}}}
\multicolumn{2}{@{}l@{}}{*Old English*} \\
Anglo Frisian Brightening & \emph{*wæskaną} \\
OE A Restoration & \emph{*waskaną} \\
OE Heavy Syllable Nasal Apocope & \emph{*waskan} \\
OE Secondary Nasalization & \emph{*waskąn} \\
OE Weak Tail Reduction & \emph{*waskan} \\
\end{tabular}
\end{minipage}
\\
\end{tabularx}
\end{minipage}%
} \endgroup

Old English form: \emph{wascan}

\subsection{Reconstruction and comparative
evidence}\label{reconstruction-and-comparative-evidence-62}

Kroonen cites \emph{*waskan-}, Orel \Recon{waskanan} `wash', and Ringe
and Taylor likewise derive Old English \emph{wascan} from the same verb
family (\citeproc{ref-Kroonen2013}{Kroonen 2013};
\citeproc{ref-Orel2003}{Orel 2003, 489};
\citeproc{ref-RingeTaylor2014}{Ringe and Taylor 2014, 142}).

\subsection{Old English evidence}\label{old-english-evidence-62}

Clark Hall heads the verb as \emph{wascan} `wash', while Sievers-Brunner
also notes the variant \emph{wæscan} `wash'
(\citeproc{ref-ClarkHall1960}{Clark Hall 1960};
\citeproc{ref-SieversBrunner1965}{Brunner 1965}).

\subsection{Development to Old
English}\label{development-to-old-english-62}

From \Recon{wáskaną} `wash', brightening gives \Recon{wæskaną} `wash'.
A-restoration before the \emph{sC} cluster restores \emph{a}, and medial
\emph{sc} remains unpalatalized before the following back vowel,
yielding \emph{wascan} (\citeproc{ref-Campbell1959}{Campbell 1959};
\citeproc{ref-RingeTaylor2014}{Ringe and Taylor 2014, 142}).

\subsection{Form note}\label{form-note-24}

The conservative citation form \emph{wascan} `wash' is selected here.
Spellings such as \emph{wæscan} `wash' or \emph{wasċan} `wash' belong to
variant or normalized background rather than to the target of this
entry.

\section{\texorpdfstring{wax --- OE
\emph{weaxan}}{wax --- OE weaxan}}\label{wax-oe-weaxan}

\index[iv]{02oe@\textbf{Old English}!weaxan@\emph{weaxan}}
\index[iv]{03pgmc@\textbf{Proto-Germanic}!waxsana@*wáxsaną}

Derivation: \emph{*wáxsaną} \textgreater{} \emph{weaxan} (regular).

\subsection{Derivation trace}\label{derivation-trace-63}

Proto input: \emph{*wáxsaną}

\begingroup
\setlength{\fboxsep}{6pt}

\noindent\fbox{%
\begin{minipage}{0.97\linewidth}
\footnotesize
\begin{tabularx}{\linewidth}{@{}>{\raggedright\arraybackslash}p{0.320\linewidth}>{\centering\arraybackslash}X>{\raggedright\arraybackslash}p{0.520\linewidth}@{}}
\begin{minipage}[t]{\linewidth}
\centering\textbf{Earlier Germanic changes}\par
\vspace{0.35em}
\raggedright
\centering\textbf{}\par
\raggedright
\vspace{0.2em}
\begin{tabular}{@{}>{\raggedright\arraybackslash}p{0.68\linewidth}@{\hspace{0.55em}}>{\raggedright\arraybackslash}p{0.16\linewidth}@{\hspace{0.25em}}}
\multicolumn{2}{@{}l@{}}{*Proto-West Germanic*} \\
\multicolumn{2}{@{}l@{}}{[no change]} \\
\multicolumn{2}{@{}l@{}}{*Northwest Germanic*} \\
\multicolumn{2}{@{}l@{}}{[no change]} \\
\end{tabular}
\end{minipage}
&
&
\begin{minipage}[t]{\linewidth}
\centering\textbf{Old English changes}\par
\vspace{0.35em}
\raggedright
\centering\textbf{}\par
\raggedright
\vspace{0.2em}
\begin{tabular}{@{}>{\raggedright\arraybackslash}p{0.66\linewidth}@{\hspace{0.55em}}>{\raggedright\arraybackslash}p{0.24\linewidth}@{\hspace{0.25em}}}
\multicolumn{2}{@{}l@{}}{*Old English*} \\
Anglo Frisian Brightening & \emph{*wæxsaną} \\
OE Breaking & \emph{*weaxsaną} \\
OE Heavy Syllable Nasal Apocope & \emph{*weaxsan} \\
OE Secondary Nasalization & \emph{*weaxsąn} \\
OE Weak Tail Reduction & \emph{*weaxsan} \\
OE Xs Merge & \emph{*weaXSan} \\
\end{tabular}
\end{minipage}
\\
\end{tabularx}
\end{minipage}%
} \endgroup

Old English form: \emph{weaxan}

\subsection{Reconstruction and comparative
evidence}\label{reconstruction-and-comparative-evidence-63}

Kroonen cites the verb as \emph{*wahs(j)an-}, Orel as \Recon{waxsanan}
`wax', and Ringe and Taylor discuss the prehistory of Old English
\emph{weaxan} within the same verbal family
(\citeproc{ref-Kroonen2013}{Kroonen 2013}; \citeproc{ref-Orel2003}{Orel
2003, 478}; \citeproc{ref-RingeTaylor2014}{Ringe and Taylor 2014}).

\subsection{Old English evidence}\label{old-english-evidence-63}

Clark Hall gives \emph{weaxan} `wax' as the verb headword and separately
records bare \emph{wax} `wax' as a preterite form; Bright likewise
treats \emph{weaxan} `wax' as the infinitive
(\citeproc{ref-ClarkHall1960}{Clark Hall 1960};
\citeproc{ref-BrightCassidyRingler1971}{Bright 1971}).

\subsection{Development to Old
English}\label{development-to-old-english-63}

From \Recon{wáxsaną} `wax', brightening and breaking yield \emph{weax-},
and the infinitive develops regularly to \emph{weaxan}. The cluster is
preserved here, since \emph{xs} \textgreater{} \emph{s} belongs to forms
where another consonant follows, such as \emph{wæstm} `growth', not to
the infinitive itself (\citeproc{ref-Campbell1959}{Campbell 1959};
\citeproc{ref-SieversBrunner1965}{Brunner 1965}).

\subsection{Lexical note}\label{lexical-note-3}

The target here is the infinitive \emph{weaxan} `wax'. Noun \emph{weax}
`wax' and preterite \emph{wax} / \emph{wēox} `waxed' belong to different
lexical or paradigm slots.

\section{\texorpdfstring{way --- OE
\emph{weġ}}{way --- OE weġ}}\label{way-oe-weux121}

\index[iv]{02oe@\textbf{Old English}!weg@\emph{weġ}}
\index[iv]{03pgmc@\textbf{Proto-Germanic}!wegaz@*wégaz}

Derivation: \emph{*wégaz} \textgreater{} \emph{weġ} (regular).

\subsection{Derivation trace}\label{derivation-trace-64}

Proto input: \emph{*wégaz}

\begingroup
\setlength{\fboxsep}{6pt}

\noindent\fbox{%
\begin{minipage}{0.97\linewidth}
\small
\begin{tabularx}{\linewidth}{@{}>{\raggedright\arraybackslash}p{0.440\linewidth}>{\centering\arraybackslash}X>{\raggedright\arraybackslash}p{0.440\linewidth}@{}}
\begin{minipage}[t]{\linewidth}
\centering\textbf{Earlier Germanic changes}\par
\vspace{0.35em}
\raggedright
\centering\textbf{}\par
\raggedright
\vspace{0.2em}
\begin{tabular}{@{}>{\raggedright\arraybackslash}p{0.74\linewidth}@{\hspace{0.55em}}>{\raggedright\arraybackslash}p{0.16\linewidth}@{\hspace{0.25em}}}
\multicolumn{2}{@{}l@{}}{*Proto-West Germanic*} \\
\multicolumn{2}{@{}l@{}}{[no change]} \\
\multicolumn{2}{@{}l@{}}{*Northwest Germanic*} \\
\mbox{PGmc Final Z Deletion} & \emph{*wéga} \\
\end{tabular}
\end{minipage}
&
&
\begin{minipage}[t]{\linewidth}
\centering\textbf{Old English changes}\par
\vspace{0.35em}
\raggedright
\centering\textbf{}\par
\raggedright
\vspace{0.2em}
\begin{tabular}{@{}>{\raggedright\arraybackslash}p{0.70\linewidth}@{\hspace{0.55em}}>{\raggedright\arraybackslash}p{0.16\linewidth}@{\hspace{0.25em}}}
\multicolumn{2}{@{}l@{}}{*Old English*} \\
PWGmc Final Bare A Loss & \emph{*wég} \\
OE Velar Palatalization & \emph{*wéʤ} \\
\end{tabular}
\end{minipage}
\\
\end{tabularx}
\end{minipage}%
} \endgroup

Old English form: \emph{weġ}

\subsection{Reconstruction and comparative
evidence}\label{reconstruction-and-comparative-evidence-64}

Kroonen cites the noun as \emph{*wega-} `way, road', while the selected
derivational form here is nominative-singular \Recon{wégaz} `way'
(\citeproc{ref-Kroonen2013}{Kroonen 2013}). Campbell, Hogg, and Ringe
and Taylor use the same word as the standard contrast between singular
palatal \emph{weġ} `way' and inflected \emph{wegas} / \emph{wegum}
(\citeproc{ref-Campbell1959}{Campbell 1959};
\citeproc{ref-Hogg1992}{Hogg 1992}; \citeproc{ref-RingeTaylor2014}{Ringe
and Taylor 2014, 341}).

\subsection{Old English evidence}\label{old-english-evidence-64}

The Old English singular is the ordinary noun \emph{weg} `way', here
normalized as \emph{weġ} `way' to show the palatal final. The
contrasting plural and oblique forms \emph{wegas} `ways', \emph{wegum}
`ways' keep a velar stop before the following back vowel
(\citeproc{ref-Campbell1959}{Campbell 1959};
\citeproc{ref-RingeTaylor2014}{Ringe and Taylor 2014, 341}).

\subsection{Development to Old
English}\label{development-to-old-english-64}

From \Recon{wégaz} `way', final \emph{*z} is lost and the weak tail
apocopates, leaving word-final \emph{*g} after a front vowel. In that
environment Old English palatalization yields \emph{weġ} `way', whereas
\emph{wegas} `ways' remains velar because the following \emph{a} blocks
the same outcome (\citeproc{ref-Campbell1959}{Campbell 1959};
\citeproc{ref-Hogg1992}{Hogg 1992}; \citeproc{ref-RingeTaylor2014}{Ringe
and Taylor 2014, 341}).

\subsection{Form note}\label{form-note-25}

Normalized \emph{weġ} `way' and dictionary \emph{weg} `way' represent
the same noun. No Old English form \emph{wē} `we' is supported in the
checked evidence for `way'.

\section{\texorpdfstring{weapon --- OE
\emph{wǣpn}}{weapon --- OE wǣpn}}\label{weapon-oe-wux1e3pn}

\index[iv]{02oe@\textbf{Old English}!waepn@\emph{wǣpn}}
\index[iv]{03pgmc@\textbf{Proto-Germanic}!wepna@*wḗpną}

Derivation: \emph{*wḗpną} \textgreater{} \emph{wǣpn} (regular).

\subsection{Derivation trace}\label{derivation-trace-65}

Proto input: \emph{*wḗpną}

\begingroup
\setlength{\fboxsep}{6pt}

\noindent\fbox{%
\begin{minipage}{0.97\linewidth}
\small
\begin{tabularx}{\linewidth}{@{}>{\raggedright\arraybackslash}p{0.460\linewidth}>{\centering\arraybackslash}X>{\raggedright\arraybackslash}p{0.460\linewidth}@{}}
\begin{minipage}[t]{\linewidth}
\centering\textbf{Earlier Germanic changes}\par
\vspace{0.35em}
\raggedright
\centering\textbf{}\par
\raggedright
\vspace{0.2em}
\begin{tabular}{@{}>{\raggedright\arraybackslash}p{0.74\linewidth}@{\hspace{0.55em}}>{\raggedright\arraybackslash}p{0.16\linewidth}@{\hspace{0.25em}}}
\multicolumn{2}{@{}l@{}}{*Proto-West Germanic*} \\
\multicolumn{2}{@{}l@{}}{[no change]} \\
\multicolumn{2}{@{}l@{}}{*Northwest Germanic*} \\
\mbox{NWGmc Long E Lowering} & \emph{*wǣpną} \\
\end{tabular}
\end{minipage}
&
&
\begin{minipage}[t]{\linewidth}
\centering\textbf{Old English changes}\par
\vspace{0.35em}
\raggedright
\centering\textbf{}\par
\raggedright
\vspace{0.2em}
\begin{tabular}{@{}>{\raggedright\arraybackslash}p{0.74\linewidth}@{\hspace{0.55em}}>{\raggedright\arraybackslash}p{0.16\linewidth}@{\hspace{0.25em}}}
\multicolumn{2}{@{}l@{}}{*Old English*} \\
OE Heavy Syllable Nasal Apocope & \emph{*wǣpn} \\
\end{tabular}
\end{minipage}
\\
\end{tabularx}
\end{minipage}%
} \endgroup

Old English form: \emph{wǣpn}

\subsection{Reconstruction and comparative
evidence}\label{reconstruction-and-comparative-evidence-65}

Kroonen reconstructs a double-stem noun \emph{*wēbna-} \textasciitilde{}
\emph{*wēpna-} and cites OE \emph{wæpn} `weapon' among its reflexes
(\citeproc{ref-Kroonen2013}{Kroonen 2013, 617}). The derivational input
\Recon{wḗpną} `weapon' represents the unbroken citation-form noun rather
than the later broken simplex.

\subsection{Old English evidence}\label{old-english-evidence-65}

Campbell's cluster-noun discussion preserves unbroken \emph{wépn}
`weapon' beside broken \emph{wépen} `weapon'-type forms
(\citeproc{ref-Campbell1959}{Campbell 1959, 150};
\citeproc{ref-Campbell1959}{1959, 226--27}). Bright contrasts broken
nominative \emph{wǣpen} `weapon' / \emph{wapen} `weapon' with unbroken
oblique \emph{wǣpnes} `weapon', while Clark Hall lemmatizes the noun
under \emph{wapen} `weapon' and also preserves unbroken forms in
compounds and related spellings
(\citeproc{ref-BrightCassidyRingler1971}{Bright 1971, 29};
\citeproc{ref-ClarkHall1960}{Clark Hall 1960, 355}).

\subsection{Development to Old
English}\label{development-to-old-english-65}

Northwest Germanic lowering gives \emph{wǣpn} `weapon', and loss of the
final nasal vowel leaves the unbroken cluster word-finally. The Old
English form here is the attested unbroken form \emph{wǣpn} `weapon'.

\subsection{Form note}\label{form-note-26}

The ordinary late West Saxon simplex headword is \emph{wǣpen} `weapon',
but Campbell's noun-class discussion also preserves unbroken \emph{wépn}
`weapon' beside broken \emph{wépen} `weapon'-type forms
(\citeproc{ref-Campbell1959}{Campbell 1959, 150};
\citeproc{ref-Campbell1959}{1959, 226--27}). \emph{wǣpnes} `weapon'
remains the regular unbroken oblique comparator
(\citeproc{ref-BrightCassidyRingler1971}{Bright 1971, 29};
\citeproc{ref-ClarkHall1960}{Clark Hall 1960, 355}).

\section{\texorpdfstring{will --- OE
\emph{willa}}{will --- OE willa}}\label{will-oe-willa}

\index[iv]{02oe@\textbf{Old English}!willa@\emph{willa}}
\index[iv]{03pgmc@\textbf{Proto-Germanic}!weljo@*wéljô}

Derivation: \emph{*wéljô} \textgreater{} \emph{willa} (regular).

\subsection{Derivation trace}\label{derivation-trace-66}

Proto input: \emph{*wéljô}

\begingroup
\setlength{\fboxsep}{6pt}

\noindent\fbox{%
\begin{minipage}{0.97\linewidth}
\small
\begin{tabularx}{\linewidth}{@{}>{\raggedright\arraybackslash}p{0.440\linewidth}>{\centering\arraybackslash}X>{\raggedright\arraybackslash}p{0.440\linewidth}@{}}
\begin{minipage}[t]{\linewidth}
\centering\textbf{Earlier Germanic changes}\par
\vspace{0.35em}
\raggedright
\centering\textbf{}\par
\raggedright
\vspace{0.2em}
\begin{tabular}{@{}>{\raggedright\arraybackslash}p{0.68\linewidth}@{\hspace{0.55em}}>{\raggedright\arraybackslash}p{0.16\linewidth}@{\hspace{0.25em}}}
\multicolumn{2}{@{}l@{}}{*Proto-West Germanic*} \\
PWGmc J Gemination & \emph{*wélljô} \\
\multicolumn{2}{@{}l@{}}{*Northwest Germanic*} \\
\multicolumn{2}{@{}l@{}}{[no change]} \\
\end{tabular}
\end{minipage}
&
&
\begin{minipage}[t]{\linewidth}
\centering\textbf{Old English changes}\par
\vspace{0.35em}
\raggedright
\centering\textbf{}\par
\raggedright
\vspace{0.2em}
\begin{tabular}{@{}>{\raggedright\arraybackslash}p{0.74\linewidth}@{\hspace{0.55em}}>{\raggedright\arraybackslash}p{0.16\linewidth}@{\hspace{0.25em}}}
\multicolumn{2}{@{}l@{}}{*Old English*} \\
OE I Umlaut & \emph{*willjô} \\
OE Unstressed Long Vowel Shortening & \emph{*willja} \\
OE J Loss After Heavy & \emph{*willa} \\
\end{tabular}
\end{minipage}
\\
\end{tabularx}
\end{minipage}%
} \endgroup

Old English form: \emph{willa}

\subsection{Reconstruction and comparative
evidence}\label{reconstruction-and-comparative-evidence-66}

Kroonen separates noun \Recon{weljan} `will, wish' from verb
\Recon{weljan} `to want', while Orel and Kluge represent the noun as
\Recon{weljōn} `will' (\citeproc{ref-Kroonen2013}{Kroonen 2013};
\citeproc{ref-Orel2003}{Orel 2003};
\citeproc{ref-KlugeSeebold2011}{Kluge 2011}). The selected derivational
form \Recon{wéljô} `will' is the noun-side input used for this row.

\subsection{Old English evidence}\label{old-english-evidence-66}

Clark Hall lemmatizes noun \emph{willa m.} separately from verb
{\emph{willan}} `will, want' (\citeproc{ref-ClarkHall1960}{Clark Hall
1960, 368}). The Old English form here is the noun citation form, not
the related verb.

\subsection{Development to Old
English}\label{development-to-old-english-66}

From \Recon{wéljô} `will', j-gemination yields a heavy stem, i-umlaut
gives \emph{will-}, and later shortening plus j-loss produce
{\emph{willa}} `will, desire'. The noun is therefore a regular weak
masculine outcome.

\subsection{Lexical note}\label{lexical-note-4}

The target here is the noun \emph{willa} `will, wish'. Related verb
{\emph{willan}} `will, want' belongs to a separate lexeme and should not
be substituted for the noun row (\citeproc{ref-Kroonen2013}{Kroonen
2013}; \citeproc{ref-ClarkHall1960}{Clark Hall 1960, 368}).

\section{\texorpdfstring{wind --- OE
\emph{windan}}{wind --- OE windan}}\label{wind-oe-windan}

\index[iv]{02oe@\textbf{Old English}!windan@\emph{windan}}
\index[iv]{03pgmc@\textbf{Proto-Germanic}!windana@*wíndaną}

Derivation: \emph{*wíndaną} \textgreater{} \emph{windan} (regular).

\subsection{Derivation trace}\label{derivation-trace-67}

Proto input: \emph{*wíndaną}

\begingroup
\setlength{\fboxsep}{6pt}

\noindent\fbox{%
\begin{minipage}{0.97\linewidth}
\small
\begin{tabularx}{\linewidth}{@{}>{\raggedright\arraybackslash}p{0.320\linewidth}>{\centering\arraybackslash}X>{\raggedright\arraybackslash}p{0.520\linewidth}@{}}
\begin{minipage}[t]{\linewidth}
\centering\textbf{Earlier Germanic changes}\par
\vspace{0.35em}
\raggedright
\centering\textbf{}\par
\raggedright
\vspace{0.2em}
\begin{tabular}{@{}>{\raggedright\arraybackslash}p{0.68\linewidth}@{\hspace{0.55em}}>{\raggedright\arraybackslash}p{0.16\linewidth}@{\hspace{0.25em}}}
\multicolumn{2}{@{}l@{}}{*Proto-West Germanic*} \\
\multicolumn{2}{@{}l@{}}{[no change]} \\
\multicolumn{2}{@{}l@{}}{*Northwest Germanic*} \\
\multicolumn{2}{@{}l@{}}{[no change]} \\
\end{tabular}
\end{minipage}
&
&
\begin{minipage}[t]{\linewidth}
\centering\textbf{Old English changes}\par
\vspace{0.35em}
\raggedright
\centering\textbf{}\par
\raggedright
\vspace{0.2em}
\begin{tabular}{@{}>{\raggedright\arraybackslash}p{0.74\linewidth}@{\hspace{0.55em}}>{\raggedright\arraybackslash}p{0.16\linewidth}@{\hspace{0.25em}}}
\multicolumn{2}{@{}l@{}}{*Old English*} \\
OE Heavy Syllable Nasal Apocope & \emph{*wíndan} \\
OE Secondary Nasalization & \emph{*wíndąn} \\
OE Weak Tail Reduction & \emph{*wíndan} \\
\end{tabular}
\end{minipage}
\\
\end{tabularx}
\end{minipage}%
} \endgroup

Old English form: \emph{windan}

\subsection{Reconstruction and comparative
evidence}\label{reconstruction-and-comparative-evidence-67}

Kroonen distinguishes noun \emph{*winda-} from verb \emph{*windan-}, and
the present row belongs to the verb (\citeproc{ref-Kroonen2013}{Kroonen
2013}). Fulk and Ringe and Taylor derive the dental directly from PIE
\emph{*wendh-}, not from a Verner alternant
(\citeproc{ref-Fulk2018}{Fulk 2018};
\citeproc{ref-RingeTaylor2014}{Ringe and Taylor 2014}).

\subsection{Old English evidence}\label{old-english-evidence-67}

Clark Hall and Bosworth-Toller record {\emph{windan}} `wind, coil' as
the verb headword (\citeproc{ref-ClarkHall1960}{Clark Hall 1960};
\citeproc{ref-BosworthToller1898}{Bosworth and Toller 1898, 101}). The
Old English form here is the ordinary infinitive of the strong verb.

\subsection{Development to Old
English}\label{development-to-old-english-67}

The derivational input \Recon{wíndaną} `wind' yields the regular
infinitive \emph{windan} by ordinary heavy-syllable apocope and
weak-tail reduction. The form is therefore a straightforward strong-verb
outcome.

\subsection{Lexical note}\label{lexical-note-5}

English \emph{wind} also names the noun. This row represents the
class-III verb, not the noun (\citeproc{ref-Kroonen2013}{Kroonen 2013};
\citeproc{ref-ClarkHall1960}{Clark Hall 1960}).

\section{\texorpdfstring{wold --- OE
\emph{weald}}{wold --- OE weald}}\label{wold-oe-weald}

\index[iv]{02oe@\textbf{Old English}!weald@\emph{weald}}
\index[iv]{03pgmc@\textbf{Proto-Germanic}!walthuz@*wálθuz}

Derivation: \emph{*wálθuz} \textgreater{} \emph{weald} (regular).

\subsection{Derivation trace}\label{derivation-trace-68}

Proto input: \emph{*wálθuz}

\begingroup
\setlength{\fboxsep}{6pt}

\noindent\fbox{%
\begin{minipage}{0.97\linewidth}
\small
\begin{tabularx}{\linewidth}{@{}>{\raggedright\arraybackslash}p{0.460\linewidth}>{\centering\arraybackslash}X>{\raggedright\arraybackslash}p{0.460\linewidth}@{}}
\begin{minipage}[t]{\linewidth}
\centering\textbf{Earlier Germanic changes}\par
\vspace{0.35em}
\raggedright
\centering\textbf{}\par
\raggedright
\vspace{0.2em}
\begin{tabular}{@{}>{\raggedright\arraybackslash}p{0.74\linewidth}@{\hspace{0.55em}}>{\raggedright\arraybackslash}p{0.16\linewidth}@{\hspace{0.25em}}}
\multicolumn{2}{@{}l@{}}{*Proto-West Germanic*} \\
PWGmc L Th Voicing & \emph{*wálduz} \\
\multicolumn{2}{@{}l@{}}{*Northwest Germanic*} \\
\mbox{PGmc Final Z Deletion} & \emph{*wáldu} \\
\end{tabular}
\end{minipage}
&
&
\begin{minipage}[t]{\linewidth}
\centering\textbf{Old English changes}\par
\vspace{0.35em}
\raggedright
\centering\textbf{}\par
\raggedright
\vspace{0.2em}
\begin{tabular}{@{}>{\raggedright\arraybackslash}p{0.74\linewidth}@{\hspace{0.55em}}>{\raggedright\arraybackslash}p{0.16\linewidth}@{\hspace{0.25em}}}
\multicolumn{2}{@{}l@{}}{*Old English*} \\
Anglo Frisian Brightening & \emph{*wældu} \\
OE Breaking & \emph{*wealdu} \\
\mbox{OE High Vowel Apocope} & \emph{*weald} \\
\end{tabular}
\end{minipage}
\\
\end{tabularx}
\end{minipage}%
} \endgroup

Old English form: \emph{weald}

\subsection{Reconstruction and comparative
evidence}\label{reconstruction-and-comparative-evidence-68}

Kroonen reconstructs the noun as \emph{*walþu-} and gives OE
{\emph{weald}} `forest, woodland' beside other West Germanic
{\emph{wald}} `forest' forms (\citeproc{ref-Kroonen2013}{Kroonen 2013}).
The derivational input \Recon{wálθuz} `wold' is the nominative singular
used for the derivation.

\subsection{Old English evidence}\label{old-english-evidence-68}

Clark Hall makes {\emph{weald}} `forest, woodland' the main noun
headword and cross-refers {\emph{wald}} `forest' and {\emph{wold}}
`plain, open country' to it (\citeproc{ref-ClarkHall1960}{Clark Hall
1960}). The Anglian-looking {\emph{wald}} `forest' therefore remains
variant background rather than the main target.

\subsection{Development to Old
English}\label{development-to-old-english-68}

\emph{*lþ} voices to \emph{ld}, Anglo-Frisian brightening yields
\emph{wæld-}, and breaking before the cluster gives \emph{weald-};
apocope then yields {\emph{weald}} `forest, woodland'. The noun is
therefore a regular breaking outcome.

\subsection{Dialect note}\label{dialect-note-6}

{\emph{wald}} `forest' survives as an Anglian-type variant in the same
family. The Old English form here is normalized {\emph{weald}} `forest,
woodland', not the variant form (\citeproc{ref-ClarkHall1960}{Clark Hall
1960}; \citeproc{ref-RingeTaylor2014}{Ringe and Taylor 2014}).

\section{\texorpdfstring{yarn --- OE
\emph{ġearn}}{yarn --- OE ġearn}}\label{yarn-oe-ux121earn}

\index[iv]{02oe@\textbf{Old English}!gearn@\emph{ġearn}}
\index[iv]{03pgmc@\textbf{Proto-Germanic}!garna@*gárną}

Derivation: \emph{*gárną} \textgreater{} \emph{ġearn} (regular).

\subsection{Derivation trace}\label{derivation-trace-69}

Proto input: \emph{*gárną}

\begingroup
\setlength{\fboxsep}{6pt}

\noindent\fbox{%
\begin{minipage}{0.97\linewidth}
\small
\begin{tabularx}{\linewidth}{@{}>{\raggedright\arraybackslash}p{0.320\linewidth}>{\centering\arraybackslash}X>{\raggedright\arraybackslash}p{0.520\linewidth}@{}}
\begin{minipage}[t]{\linewidth}
\centering\textbf{Earlier Germanic changes}\par
\vspace{0.35em}
\raggedright
\centering\textbf{}\par
\raggedright
\vspace{0.2em}
\begin{tabular}{@{}>{\raggedright\arraybackslash}p{0.68\linewidth}@{\hspace{0.55em}}>{\raggedright\arraybackslash}p{0.16\linewidth}@{\hspace{0.25em}}}
\multicolumn{2}{@{}l@{}}{*Proto-West Germanic*} \\
\multicolumn{2}{@{}l@{}}{[no change]} \\
\multicolumn{2}{@{}l@{}}{*Northwest Germanic*} \\
\multicolumn{2}{@{}l@{}}{[no change]} \\
\end{tabular}
\end{minipage}
&
&
\begin{minipage}[t]{\linewidth}
\centering\textbf{Old English changes}\par
\vspace{0.35em}
\raggedright
\centering\textbf{}\par
\raggedright
\vspace{0.2em}
\begin{tabular}{@{}>{\raggedright\arraybackslash}p{0.74\linewidth}@{\hspace{0.55em}}>{\raggedright\arraybackslash}p{0.16\linewidth}@{\hspace{0.25em}}}
\multicolumn{2}{@{}l@{}}{*Old English*} \\
Anglo Frisian Brightening & \emph{*gærną} \\
OE Breaking & \emph{*gearną} \\
OE Heavy Syllable Nasal Apocope & \emph{*gearn} \\
OE Velar Palatalization & \emph{*ʤearn} \\
\end{tabular}
\end{minipage}
\\
\end{tabularx}
\end{minipage}%
} \endgroup

Old English form: \emph{ġearn}

\subsection{Reconstruction and comparative
evidence}\label{reconstruction-and-comparative-evidence-69}

Kroonen cites the noun as \emph{*garna-}, and Ringe and Taylor give the
early chain \Recon{garna} `yarn' \textgreater{} \Recon{geern} `yarn'
\textgreater{} \Recon{gearn} `yarn' \textgreater{} OE {\emph{gearn}}
`yarn' (\citeproc{ref-Kroonen2013}{Kroonen 2013};
\citeproc{ref-RingeTaylor2014}{Ringe and Taylor 2014, 220}). The
derivational input \Recon{gárną} `yarn' is the nominal citation form
used here, while oblique \Recon{garnăn} `yarn' belongs only to
comparative background.

\subsection{Old English evidence}\label{old-english-evidence-69}

Clark Hall records \emph{gearn (e) n.} `yarn, spun wool', and
Bosworth-Toller glosses {\emph{gearn}} `yarn' with Latin \emph{filatum}
`spun thread' (\citeproc{ref-ClarkHall1960}{Clark Hall 1960};
\citeproc{ref-BosworthToller1898}{Bosworth and Toller 1898}).

\subsection{Development to Old
English}\label{development-to-old-english-69}

From \Recon{gárną} `yarn', brightening and breaking before \emph{rn}
yield \emph{gearn} `yarn'; palatalization of initial \emph{g} before the
resulting front-vocalic sequence gives normalized \emph{ġearn} `yarn'.
The derivation is regular.

\subsection{Form note}\label{form-note-27}

Dictionary \emph{gearn} `yarn' and normalized \emph{ġearn} `yarn' refer
to the same noun. The comparative stem \emph{*garna-} and oblique
\Recon{garnăn} `yarn' do not replace the derivational input
\Recon{gárną} `yarn'.

\clearpage

\chapter{Attested variants and comparison
forms}\label{attested-variants-and-comparison-forms}

Here the Old English comparator belongs to a documented set of variants.
The variation is part of the evidence and not an inconvenience to be
normalized away.

\section{\texorpdfstring{cud --- OE
\emph{cwedu}}{cud --- OE cwedu}}\label{cud-oe-cwedu}

\index[iv]{02oe@\textbf{Old English}!cwedu@\emph{cwedu}}
\index[iv]{03pgmc@\textbf{Proto-Germanic}!kweduz@*kwéðuz}
\index[iv]{03pgmc@\textbf{Proto-Germanic}!kwithuz@*kwíθuz}

Derivation: citation reconstruction \emph{*kwíθuz}; form followed here
\emph{*kwéðuz} \textgreater{} \emph{cwedu} (attested variant).

\subsection{Derivation trace}\label{derivation-trace-70}

Proto input: \emph{*kwéðuz}

\begingroup
\setlength{\fboxsep}{6pt}

\noindent\fbox{%
\begin{minipage}{0.97\linewidth}
\small
\begin{tabularx}{\linewidth}{@{}>{\raggedright\arraybackslash}p{0.540\linewidth}>{\centering\arraybackslash}X>{\raggedright\arraybackslash}p{0.300\linewidth}@{}}
\begin{minipage}[t]{\linewidth}
\centering\textbf{Earlier Germanic changes}\par
\vspace{0.35em}
\raggedright
\centering\textbf{}\par
\raggedright
\vspace{0.2em}
\begin{tabular}{@{}>{\raggedright\arraybackslash}p{0.74\linewidth}@{\hspace{0.55em}}>{\raggedright\arraybackslash}p{0.16\linewidth}@{\hspace{0.25em}}}
\multicolumn{2}{@{}l@{}}{*Proto-West Germanic*} \\
PWGmc Dental Hardening & \emph{*kwéduz} \\
\multicolumn{2}{@{}l@{}}{*Northwest Germanic*} \\
\mbox{PGmc Final Z Deletion} & \emph{*kwédu} \\
\end{tabular}
\end{minipage}
&
&
\begin{minipage}[t]{\linewidth}
\centering\textbf{Old English changes}\par
\vspace{0.35em}
\raggedright
\centering\textbf{}\par
\raggedright
\vspace{0.2em}
\begin{tabular}{@{}>{\raggedright\arraybackslash}p{0.64\linewidth}@{\hspace{0.55em}}>{\raggedright\arraybackslash}p{0.16\linewidth}@{\hspace{0.25em}}}
\multicolumn{2}{@{}l@{}}{*Old English*} \\
\multicolumn{2}{@{}l@{}}{[no change]} \\
\end{tabular}
\end{minipage}
\\
\end{tabularx}
\end{minipage}%
} \endgroup

Old English form: \emph{cwedu}

\subsection{Reconstruction and comparative
evidence}\label{reconstruction-and-comparative-evidence-70}

Kroonen reconstructs the resin word as \Recon{kwedu-2} `cud' and gives
Old English variants {\emph{cwidu}} `cud', {\emph{cweodu}} `cud', and
{\emph{c(w)udu}} `cud' (\citeproc{ref-Kroonen2013}{Kroonen 2013, 355}).
Orel likewise lists {\emph{cwidu}} under the cognate set
(\citeproc{ref-Orel2003}{Orel 2003, 266}). The derivational input
{\Recon{kwéðuz} `cud'} therefore represents the older e-grade,
voiced-dental form behind the chosen variant {\emph{cwedu}} `cud'.

\subsection{Old English evidence}\label{old-english-evidence-70}

The Old English word survives in a wider variant set than one dictionary
headword suggests. Ringe and Taylor discuss {\emph{cwidu}} `cud'
\textgreater{} {\emph{cwudu}} `cud' \textgreater{} {\emph{cudu}} `cud'
and also note late West Saxon {\emph{cweodu}} `cud'; Clark Hall gives
{\emph{cwudu}} `cud', {\emph{cweodu}} `cud', and {\emph{cudu}} `cud'
(\citeproc{ref-RingeTaylor2014}{Ringe and Taylor 2014, 338};
\citeproc{ref-ClarkHall1960}{Clark Hall 1960, 84}). Attested
\emph{cwedu} `cud' is treated here as the conservative variant within
that set.

\subsection{Development to Old
English}\label{development-to-old-english-70}

From {\Recon{kwéðuz} `cud'}, the West Germanic voiced dental hardens in
the expected way and the regular Old English development yields
{\emph{cwedu}} `cud'. The other Old English spellings belong to the same
lexical family, but reflect later leveling, back-umlaut, or further
reduction rather than a need to replace the selected input.

\subsection{Variant comparison}\label{variant-comparison}

{\def\LTcaptype{none} % do not increment counter
\begin{longtable}[]{@{}
  >{\raggedright\arraybackslash}p{(\linewidth - 4\tabcolsep) * \real{0.3333}}
  >{\raggedright\arraybackslash}p{(\linewidth - 4\tabcolsep) * \real{0.3333}}
  >{\raggedright\arraybackslash}p{(\linewidth - 4\tabcolsep) * \real{0.3333}}@{}}
\toprule\noalign{}
\begin{minipage}[b]{\linewidth}\raggedright
Variant type
\end{minipage} & \begin{minipage}[b]{\linewidth}\raggedright
Old English form
\end{minipage} & \begin{minipage}[b]{\linewidth}\raggedright
Comment
\end{minipage} \\
\midrule\noalign{}
\endhead
\bottomrule\noalign{}
\endlastfoot
conservative target & {\emph{cwedu}} & selected attested variant
represented here \\
leveled i-grade form & {\emph{cwidu}} & common lexical variant in the
same family \\
back-umlauted forms & {\emph{cweodu}}, {\emph{cwudu}} & later
developments within the same OE tradition \\
reduced form & {\emph{cudu}} & further reduced member of the same
variant set \\
\end{longtable}
}

\section{\texorpdfstring{ten --- OE
\emph{tēon}}{ten --- OE tēon}}\label{ten-oe-tux113on}

\index[iv]{02oe@\textbf{Old English}!teon@\emph{tēon}}
\index[iv]{03pgmc@\textbf{Proto-Germanic}!texun@*téxun}

Derivation: \emph{*téxun} \textgreater{} \emph{tēon} (attested variant).

\subsection{Derivation trace}\label{derivation-trace-71}

Proto input: \emph{*téxun}

\begingroup
\setlength{\fboxsep}{6pt}

\noindent\fbox{%
\begin{minipage}{0.97\linewidth}
\small
\begin{tabularx}{\linewidth}{@{}>{\raggedright\arraybackslash}p{0.320\linewidth}>{\centering\arraybackslash}X>{\raggedright\arraybackslash}p{0.520\linewidth}@{}}
\begin{minipage}[t]{\linewidth}
\centering\textbf{Earlier Germanic changes}\par
\vspace{0.35em}
\raggedright
\centering\textbf{}\par
\raggedright
\vspace{0.2em}
\begin{tabular}{@{}>{\raggedright\arraybackslash}p{0.68\linewidth}@{\hspace{0.55em}}>{\raggedright\arraybackslash}p{0.16\linewidth}@{\hspace{0.25em}}}
\multicolumn{2}{@{}l@{}}{*Proto-West Germanic*} \\
\multicolumn{2}{@{}l@{}}{[no change]} \\
\multicolumn{2}{@{}l@{}}{*Northwest Germanic*} \\
\multicolumn{2}{@{}l@{}}{[no change]} \\
\end{tabular}
\end{minipage}
&
&
\begin{minipage}[t]{\linewidth}
\centering\textbf{Old English changes}\par
\vspace{0.35em}
\raggedright
\centering\textbf{}\par
\raggedright
\vspace{0.2em}
\begin{tabular}{@{}>{\raggedright\arraybackslash}p{0.74\linewidth}@{\hspace{0.55em}}>{\raggedright\arraybackslash}p{0.16\linewidth}@{\hspace{0.25em}}}
\multicolumn{2}{@{}l@{}}{*Old English*} \\
OE Med Unstressed U Lowering & \emph{*téxon} \\
OE Breaking & \emph{*téoxon} \\
OE H Loss & \emph{*téoon} \\
OE Contraction & \emph{*tḗon} \\
\end{tabular}
\end{minipage}
\\
\end{tabularx}
\end{minipage}%
} \endgroup

Old English form: \emph{tēon}

\subsection{Reconstruction and comparative
evidence}\label{reconstruction-and-comparative-evidence-71}

\Recon{tēon} `ten' is the bare cardinal form. Fulk states that Old
English \emph{tien} `ten' shows umlaut from the inflected forms, whereas
the uninflected form without umlaut is reflected in \emph{tēon} `ten'
(\citeproc{ref-Fulk2018}{Fulk 2018, sect. 10.2}). Brunner gives the same
contrast more broadly: \Recon{tēon} `ten' develops from \Recon{tëhun}
`ten', while West Saxon \emph{tien} `ten', \emph{tȳn} `ten' belong to a
different, umlauted branch of the numeral history
(\citeproc{ref-SieversBrunner1965}{Brunner 1965, sects 129.2, 129} Anm.
6, 234).

The comparison form \emph{tēon} `ten' therefore represents the bare
cardinal's un-umlauted line, not the later umlauted simplex tradition.

\subsection{Old English evidence}\label{old-english-evidence-71}

The attested simplex forms are varied. Campbell gives \emph{tien} `ten',
north-western West Saxon \emph{tēn} `ten', and late Northumbrian
\emph{tēo} `ten', \emph{tēa} `ten' (\citeproc{ref-Campbell1959}{Campbell
1959, sect. 682}). Brunner likewise lists West Saxon \emph{tien} `ten',
\emph{tȳn} `ten' beside \emph{tēn} `ten', \emph{tēo} `ten', \emph{tēa}
`ten' in other dialects (\citeproc{ref-SieversBrunner1965}{Brunner 1965,
sect. 325}).

Exact simplex \emph{tēon} `ten' is weaker as a directly cited headword
than those spellings. The un-umlauted stem is, however, explicit in
\emph{tēoða} `tenth' and \emph{tēontig} `ten'
(\citeproc{ref-SieversBrunner1965}{Brunner 1965, sect. 129.2};
\citeproc{ref-Fulk2018}{Fulk 2018, sect. 10.2}). The comparison form
\emph{tēon} `ten' is therefore a normalized spelling of that un-umlauted
base.

\subsection{Development to Old
English}\label{development-to-old-english-71}

From \textbf{\Recon{téxun} `ten'}, lowering of medial unstressed
\emph{u} gives \textbf{\Recon{téxon} `ten'}, breaking gives
\textbf{\Recon{téoxon} `ten'}, loss of intervocalic \emph{h}/\emph{x}
yields \textbf{\Recon{téoon} `ten'}, and contraction produces
\textbf{\Recon{tḗon} `ten'}, written \emph{tēon} `ten'. This is the
regular bare-cardinal path.

The umlauted forms \emph{tien} / \emph{tīen} `ten' belong to a different
branch, created when the numeral was levelled from inflected forms with
a front-vocalic trigger (\citeproc{ref-Fulk2018}{Fulk 2018, sect. 10.2};
\citeproc{ref-SieversBrunner1965}{Brunner 1965, sect. 129} Anm. 6).

\subsection{Variant comparison}\label{variant-comparison-1}

The comparison below sets the relevant forms side by side. It
distinguishes the normalized un-umlauted base from the attested simplex
variants.

{\def\LTcaptype{none} % do not increment counter
\begin{longtable}[]{@{}
  >{\raggedright\arraybackslash}p{(\linewidth - 4\tabcolsep) * \real{0.3333}}
  >{\raggedright\arraybackslash}p{(\linewidth - 4\tabcolsep) * \real{0.3333}}
  >{\raggedright\arraybackslash}p{(\linewidth - 4\tabcolsep) * \real{0.3333}}@{}}
\toprule\noalign{}
\begin{minipage}[b]{\linewidth}\raggedright
Form or branch
\end{minipage} & \begin{minipage}[b]{\linewidth}\raggedright
Status
\end{minipage} & \begin{minipage}[b]{\linewidth}\raggedright
Relevance to this entry
\end{minipage} \\
\midrule\noalign{}
\endhead
\bottomrule\noalign{}
\endlastfoot
\emph{tēon} & normalized un-umlauted comparison form; trace-supported &
Old English form here \\
\emph{tien} / \emph{tīen} & attested West Saxon umlauted simplex forms &
genuine OE variants, but not the bare-cardinal line modeled here \\
\emph{tēn} / \emph{tēo} / \emph{tēa} & attested un-umlauted simplex
variants in other dialects & support the same branch as the comparison
form \\
\end{longtable}
}

\section{\texorpdfstring{three --- OE
\emph{þrīe}}{three --- OE þrīe}}\label{three-oe-uxferux12be}

\index[iv]{02oe@\textbf{Old English}!thrie@\emph{þrīe}}
\index[iv]{03pgmc@\textbf{Proto-Germanic}!threjez@*θréjez}

Derivation: \emph{*θréjez} \textgreater{} \emph{þrīe} (attested
variant).

\subsection{Derivation trace}\label{derivation-trace-72}

Proto input: \emph{*θréjez}

\begingroup
\setlength{\fboxsep}{6pt}

\noindent\fbox{%
\begin{minipage}{0.97\linewidth}
\small
\begin{tabularx}{\linewidth}{@{}>{\raggedright\arraybackslash}p{0.440\linewidth}>{\centering\arraybackslash}X>{\raggedright\arraybackslash}p{0.440\linewidth}@{}}
\begin{minipage}[t]{\linewidth}
\centering\textbf{Earlier Germanic changes}\par
\vspace{0.35em}
\raggedright
\centering\textbf{}\par
\raggedright
\vspace{0.2em}
\begin{tabular}{@{}>{\raggedright\arraybackslash}p{0.74\linewidth}@{\hspace{0.55em}}>{\raggedright\arraybackslash}p{0.16\linewidth}@{\hspace{0.25em}}}
\multicolumn{2}{@{}l@{}}{*Proto-West Germanic*} \\
\multicolumn{2}{@{}l@{}}{[no change]} \\
\multicolumn{2}{@{}l@{}}{*Northwest Germanic*} \\
\mbox{PGmc Final Z Deletion} & \emph{*θréje} \\
\end{tabular}
\end{minipage}
&
&
\begin{minipage}[t]{\linewidth}
\centering\textbf{Old English changes}\par
\vspace{0.35em}
\raggedright
\centering\textbf{}\par
\raggedright
\vspace{0.2em}
\begin{tabular}{@{}>{\raggedright\arraybackslash}p{0.74\linewidth}@{\hspace{0.55em}}>{\raggedright\arraybackslash}p{0.16\linewidth}@{\hspace{0.25em}}}
\multicolumn{2}{@{}l@{}}{*Old English*} \\
OE I Umlaut & \emph{*θrije} \\
OE Intervocalic J Vocalization & \emph{*θriie} \\
OE Contraction & \emph{*θrīe} \\
\end{tabular}
\end{minipage}
\\
\end{tabularx}
\end{minipage}%
} \endgroup

Old English form: \emph{þrīe}

\subsection{Reconstruction and comparative
evidence}\label{reconstruction-and-comparative-evidence-72}

Kroonen cites the numeral under a broader stem-style reconstruction
rather than under one Old English-ready paradigm cell
(\citeproc{ref-Kroonen2013}{Kroonen 2013, 586}). The input
\textbf{{\Recon{θréjez} `three'}} is therefore best understood as the
inherited masculine nominative-accusative singular.

The Old English numeral has no uniform citation form across the
paradigm. The masculine singular line must be kept apart from
feminine-neuter \emph{þrēo} `three' and from later reduced spellings of
the masculine form.

\subsection{Old English evidence}\label{old-english-evidence-72}

Campbell gives masculine nominative-accusative {\emph{þrīe}} `three',
feminine and neuter nominative-accusative {\emph{þrēo}} `three',
genitive {\emph{þrēora}} `three', and dative {\emph{þrim}} `three',
adding that late West Saxon has \emph{þry} `three' / \emph{þri} `three'
for masculine nominative-accusative beside {\emph{þrīe}} `three'
(\citeproc{ref-Campbell1959}{Campbell 1959, sect. 683}). Fulk presents
the same masculine \emph{þrīe} `three' beside the wider numeral paradigm
(\citeproc{ref-Fulk2018}{Fulk 2018, sect. 10.1}).

The target is therefore an attested Old English paradigm cell.
\emph{þrī} `three' / \emph{þry} `three' belongs to later reduction or
headword-style citation, whereas {\emph{þrīe}} `three' is the
conservative masculine nominative-accusative form.

\subsection{Development to Old
English}\label{development-to-old-english-72}

From \Recon{θréjez} `three', loss of final \emph{-z} leaves
\Recon{θréje} `three' as an intermediate form. The following
\Recon{θrije} `three' fronts the stem vowel, then vocalizes between
vowels, and contraction yields \Recon{þrīe} `three'. The compact trace
records the same sequence as θrije \textgreater{} θriie \textgreater{}
\emph{þrīe}.

\subsection{Variant comparison}\label{variant-comparison-2}

The comparison below sets the relevant forms side by side. It separates
the selected masculine cell from the later reduced form and from the
rest of the numeral paradigm.

{\def\LTcaptype{none} % do not increment counter
\begin{longtable}[]{@{}
  >{\raggedright\arraybackslash}p{(\linewidth - 4\tabcolsep) * \real{0.3333}}
  >{\raggedright\arraybackslash}p{(\linewidth - 4\tabcolsep) * \real{0.3333}}
  >{\raggedright\arraybackslash}p{(\linewidth - 4\tabcolsep) * \real{0.3333}}@{}}
\toprule\noalign{}
\begin{minipage}[b]{\linewidth}\raggedright
Form
\end{minipage} & \begin{minipage}[b]{\linewidth}\raggedright
Status
\end{minipage} & \begin{minipage}[b]{\linewidth}\raggedright
Relevance to this entry
\end{minipage} \\
\midrule\noalign{}
\endhead
\bottomrule\noalign{}
\endlastfoot
{\emph{þrīe}} & attested masculine nom./acc.; regular output & Old
English form here \\
\emph{þrī} / \emph{þry} & later reduced masculine variant & genuine OE
variant, but not the conservative comparison form \\
{\emph{þrēo}} & attested feminine-neuter nom./acc. & same numeral,
different paradigm cell \\
{\emph{þrēora}}, {\emph{þrim}} & attested genitive and dative forms &
confirm the wider paradigm, not the cell compared here \\
\end{longtable}
}

\section{\texorpdfstring{wasp --- OE
\emph{wæfs}}{wasp --- OE wæfs}}\label{wasp-oe-wuxe6fs}

\index[iv]{02oe@\textbf{Old English}!waefs@\emph{wæfs}}
\index[iv]{03pgmc@\textbf{Proto-Germanic}!wabsaz@*wábsaz}

Derivation: \emph{*wábsaz} \textgreater{} \emph{wæfs} (attested
variant).

\subsection{Derivation trace}\label{derivation-trace-73}

Proto input: \emph{*wábsaz}

\begingroup
\setlength{\fboxsep}{6pt}

\noindent\fbox{%
\begin{minipage}{0.97\linewidth}
\small
\begin{tabularx}{\linewidth}{@{}>{\raggedright\arraybackslash}p{0.440\linewidth}>{\centering\arraybackslash}X>{\raggedright\arraybackslash}p{0.440\linewidth}@{}}
\begin{minipage}[t]{\linewidth}
\centering\textbf{Earlier Germanic changes}\par
\vspace{0.35em}
\raggedright
\centering\textbf{}\par
\raggedright
\vspace{0.2em}
\begin{tabular}{@{}>{\raggedright\arraybackslash}p{0.74\linewidth}@{\hspace{0.55em}}>{\raggedright\arraybackslash}p{0.16\linewidth}@{\hspace{0.25em}}}
\multicolumn{2}{@{}l@{}}{*Proto-West Germanic*} \\
\multicolumn{2}{@{}l@{}}{[no change]} \\
\multicolumn{2}{@{}l@{}}{*Northwest Germanic*} \\
\mbox{PGmc Final Z Deletion} & \emph{*wábsa} \\
\end{tabular}
\end{minipage}
&
&
\begin{minipage}[t]{\linewidth}
\centering\textbf{Old English changes}\par
\vspace{0.35em}
\raggedright
\centering\textbf{}\par
\raggedright
\vspace{0.2em}
\begin{tabular}{@{}>{\raggedright\arraybackslash}p{0.70\linewidth}@{\hspace{0.55em}}>{\raggedright\arraybackslash}p{0.16\linewidth}@{\hspace{0.25em}}}
\multicolumn{2}{@{}l@{}}{*Old English*} \\
PWGmc Final Bare A Loss & \emph{*wábs} \\
Anglo Frisian Brightening & \emph{*wæbs} \\
PGmc B Allophony & \emph{*wæβs} \\
\end{tabular}
\end{minipage}
\\
\end{tabularx}
\end{minipage}%
} \endgroup

Old English form: \emph{wæfs}

\subsection{Reconstruction and comparative
evidence}\label{reconstruction-and-comparative-evidence-73}

The Proto-Germanic form {\Recon{wábsaz} `wasp'} reaches Old English
without any special change of stem or paradigm cell. The question in
this entry is instead which attested Old English member of the variant
set should serve as the comparison form.

Fulk presents the Old English forms together as {\emph{wæfs}} `wasp'
with variants {\emph{wæsp}} `wasp' and {\emph{wæps}} `wasp'
(\citeproc{ref-Fulk2018}{Fulk 2018, sect. 6.5}). Bülbring and Brunner
then make the chronology more explicit by deriving later \emph{wæps} and
late West Saxon {\emph{wasp}} `wasp' from earlier \emph{waefs} /
\emph{wæfs} through restricted metatheses
(\citeproc{ref-Bulbring1902}{Bülbring 1902, sect. 484} Anm. 3;
\citeproc{ref-SieversBrunner1965}{Brunner 1965, sects 193, 204}).

\subsection{Old English evidence}\label{old-english-evidence-73}

The earliest directly cited Old English form is {\emph{wæfs}} `wasp',
written {\emph{waefs}} `wasp' in the Épinal-Corpus material discussed by
Bülbring and Brunner (\citeproc{ref-Bulbring1902}{Bülbring 1902, sect.
484} Anm. 3; \citeproc{ref-SieversBrunner1965}{Brunner 1965, sect.
193}). Later Old English also shows {\emph{wæps}} `wasp' and
{\emph{wæsp}} `wasp' / {\emph{wasp}} `wasp', and dictionary practice
often favors {\emph{wæps}} or later spellings as headwords
(\citeproc{ref-ClarkHall1960}{Clark Hall 1960, 341}).

This entry therefore distinguishes chronological priority from headword
habit. {\emph{wæfs}} `wasp' is not a convenient reconstruction: it is an
attested Old English form and also the one that matches the regular
development most closely.

\subsection{Development to Old
English}\label{development-to-old-english-73}

From {\Recon{wábsaz} `wasp'}, the regular Old English path passes
through loss of final \emph{z}, Anglo-Frisian fronting, and the
allophonic development of \emph{b} to a fricative before \emph{s},
yielding {\emph{wæfs}} `wasp'.

The later forms {\emph{wæps}} `wasp' and {\emph{wæsp}} `wasp' /
{\emph{wasp}} `wasp' belong to subsequent, lexically restricted
metatheses. They are genuine Old English forms, but they are later
within the variant history.

\subsection{Variant comparison}\label{variant-comparison-3}

The comparison below sets the relevant forms side by side. It separates
the earliest attested and regular form from the later metathesized
doublets.

{\def\LTcaptype{none} % do not increment counter
\begin{longtable}[]{@{}
  >{\raggedright\arraybackslash}p{(\linewidth - 4\tabcolsep) * \real{0.3333}}
  >{\raggedright\arraybackslash}p{(\linewidth - 4\tabcolsep) * \real{0.3333}}
  >{\raggedright\arraybackslash}p{(\linewidth - 4\tabcolsep) * \real{0.3333}}@{}}
\toprule\noalign{}
\begin{minipage}[b]{\linewidth}\raggedright
Form
\end{minipage} & \begin{minipage}[b]{\linewidth}\raggedright
Status
\end{minipage} & \begin{minipage}[b]{\linewidth}\raggedright
Relevance to this entry
\end{minipage} \\
\midrule\noalign{}
\endhead
\bottomrule\noalign{}
\endlastfoot
{\emph{wæfs}} & earliest attested OE form; regular output & Old English
form here \\
{\emph{wæps}} & later attested metathesized variant & genuine OE
doublet, but secondary \\
{\emph{wæsp}} / {\emph{wasp}} & later West Saxon metathesized variant &
genuine OE doublet, but not the form compared here \\
\end{longtable}
}

\clearpage

\chapter{Early analogy and pre-Old-English input
selection}\label{early-analogy-and-pre-old-english-input-selection}

Here an analogical change already separates the transducer input from
the lexeme's citation reconstruction before the specifically Old English
changes apply. Each entry must therefore justify the remodeled input
independently of the successful output.

\section{\texorpdfstring{bottom --- OE
\emph{botm}}{bottom --- OE botm}}\label{bottom-oe-botm}

\index[iv]{02oe@\textbf{Old English}!botm@\emph{botm}}
\index[iv]{03pgmc@\textbf{Proto-Germanic}!budmaz@*búdmaz}
\index[iv]{03pgmc@\textbf{Proto-Germanic}!buttmaz@*búttmaz}

Derivation: citation reconstruction \emph{*búdmaz}; form followed here
\emph{*búttmaz} \textgreater{} \emph{botm} (early analogy).

\subsection{Derivation trace}\label{derivation-trace-74}

Proto input: \emph{*búttmaz}

\begingroup
\setlength{\fboxsep}{6pt}

\noindent\fbox{%
\begin{minipage}{0.97\linewidth}
\small
\begin{tabularx}{\linewidth}{@{}>{\raggedright\arraybackslash}p{0.440\linewidth}>{\centering\arraybackslash}X>{\raggedright\arraybackslash}p{0.440\linewidth}@{}}
\begin{minipage}[t]{\linewidth}
\centering\textbf{Earlier Germanic changes}\par
\vspace{0.35em}
\raggedright
\centering\textbf{}\par
\raggedright
\vspace{0.2em}
\begin{tabular}{@{}>{\raggedright\arraybackslash}p{0.70\linewidth}@{\hspace{0.55em}}>{\raggedright\arraybackslash}p{0.20\linewidth}@{\hspace{0.25em}}}
\multicolumn{2}{@{}l@{}}{*Proto-West Germanic*} \\
\multicolumn{2}{@{}l@{}}{[no change]} \\
\multicolumn{2}{@{}l@{}}{*Northwest Germanic*} \\
\mbox{NWGmc U Lowering} & \emph{*bóttmaz} \\
\mbox{PGmc Final Z Deletion} & \emph{*bóttma} \\
\end{tabular}
\end{minipage}
&
&
\begin{minipage}[t]{\linewidth}
\centering\textbf{Old English changes}\par
\vspace{0.35em}
\raggedright
\centering\textbf{}\par
\raggedright
\vspace{0.2em}
\begin{tabular}{@{}>{\raggedright\arraybackslash}p{0.74\linewidth}@{\hspace{0.55em}}>{\raggedright\arraybackslash}p{0.16\linewidth}@{\hspace{0.25em}}}
\multicolumn{2}{@{}l@{}}{*Old English*} \\
PWGmc Final Bare A Loss & \emph{*bóttm} \\
OE Preconsonantal Degemination & \emph{*bótm} \\
\end{tabular}
\end{minipage}
\\
\end{tabularx}
\end{minipage}%
} \endgroup

Old English form: \emph{botm}

\subsection{Reconstruction and comparative
evidence}\label{reconstruction-and-comparative-evidence-74}

Kroonen reconstructs the word as a stem complex \Recon{budmō} `bottom'
with genitive \Recon{buttaz} `bottom', summarized as \emph{*budman-}
\textasciitilde{} \emph{*buttman-}, and gives Old English {\emph{botm}}
`bottom' as the reflex (\citeproc{ref-Kroonen2013}{Kroonen 2013, 120}).
The comparative label \Recon{búdmaz} `bottom' names the lexeme-level
stem complex, while the derivational input \Recon{búttmaz} `bottom'
represents the pre-Old-English form with oblique \emph{*butt-}
generalized into the nominative formation.

Orel likewise preserves both sides of the comparison under
\Recon{budmaz} `bottom' and \Recon{butmaz} `bottom'
(\citeproc{ref-Orel2003}{Orel 2003, 100}). The derivational input is
thus a historical stem choice, not an arbitrary respelling.

\subsection{Old English evidence}\label{old-english-evidence-74}

The Old English noun itself is secure. Clark Hall gives {\emph{botm}}
`bottom' (\citeproc{ref-ClarkHall1960}{Clark Hall 1960, 63}).
Bosworth-Toller cross-references {\emph{bodan}} `bottom, floor' to
{\emph{botm}} `bottom', showing the wider reflex family without
weakening the attested lemma (\citeproc{ref-BosworthToller1898}{Bosworth
and Toller 1898, 112}).

\subsection{Development to Old
English}\label{development-to-old-english-74}

Once the oblique \emph{*butt-} stem has been generalized, the
derivational input \Recon{búttmaz} `bottom' develops regularly to
{\emph{botm}} `bottom'. The analogical step is therefore early: it
belongs to pre-Old-English stem formation rather than to a later choice
among Old English paradigm cells.

\section{\texorpdfstring{brand --- OE
\emph{brandes}}{brand --- OE brandes}}\label{brand-oe-brandes}

\index[iv]{02oe@\textbf{Old English}!brandes@\emph{brandes}}
\index[iv]{03pgmc@\textbf{Proto-Germanic}!brandas@*brándas}
\index[iv]{03pgmc@\textbf{Proto-Germanic}!brandaz@*brándaz}

Derivation: citation reconstruction \emph{*brándaz}; form followed here
\emph{*brándas} \textgreater{} \emph{brandes} (early analogy).

\subsection{Derivation trace}\label{derivation-trace-75}

Proto input: \emph{*brándas}

\begingroup
\setlength{\fboxsep}{6pt}

\noindent\fbox{%
\begin{minipage}{0.97\linewidth}
\small
\begin{tabularx}{\linewidth}{@{}>{\raggedright\arraybackslash}p{0.440\linewidth}>{\centering\arraybackslash}X>{\raggedright\arraybackslash}p{0.440\linewidth}@{}}
\begin{minipage}[t]{\linewidth}
\centering\textbf{Earlier Germanic changes}\par
\vspace{0.35em}
\raggedright
\centering\textbf{}\par
\raggedright
\vspace{0.2em}
\begin{tabular}{@{}>{\raggedright\arraybackslash}p{0.68\linewidth}@{\hspace{0.55em}}>{\raggedright\arraybackslash}p{0.16\linewidth}@{\hspace{0.25em}}}
\multicolumn{2}{@{}l@{}}{*Proto-West Germanic*} \\
\multicolumn{2}{@{}l@{}}{[no change]} \\
\multicolumn{2}{@{}l@{}}{*Northwest Germanic*} \\
\multicolumn{2}{@{}l@{}}{[no change]} \\
\end{tabular}
\end{minipage}
&
&
\begin{minipage}[t]{\linewidth}
\centering\textbf{Old English changes}\par
\vspace{0.35em}
\raggedright
\centering\textbf{}\par
\raggedright
\vspace{0.2em}
\begin{tabular}{@{}>{\raggedright\arraybackslash}p{0.70\linewidth}@{\hspace{0.55em}}>{\raggedright\arraybackslash}p{0.20\linewidth}@{\hspace{0.25em}}}
\multicolumn{2}{@{}l@{}}{*Old English*} \\
Anglo Frisian Brightening & \emph{*brándæs} \\
OE Unstressed AE Merger & \emph{*brándes} \\
\end{tabular}
\end{minipage}
\\
\end{tabularx}
\end{minipage}%
} \endgroup

Old English form: \emph{brandes}

\subsection{Reconstruction and comparative
evidence}\label{reconstruction-and-comparative-evidence-75}

The inherited noun is the masculine a-stem {\Recon{brándaz} `brand'},
continued by Old English {\emph{brand}} `brand' and its continental
cognates (\citeproc{ref-Orel2003}{Orel 2003, 53}). The selected input
{\Recon{brándas} `brand'} is not a different lexeme but the genitive
singular of that same a-stem noun.

Both forms belong to the same root and stem class but occupy different
inherited inflectional cells. The derivational input is the oblique
cell.

\subsection{Old English evidence}\label{old-english-evidence-75}

Old English dictionaries lemmatize the noun as {\emph{brand}} `brand'
(\citeproc{ref-ClarkHall1960}{Clark Hall 1960, 49};
\citeproc{ref-BosworthToller1898}{Bosworth and Toller 1898, 116}).
Bosworth-Toller also records inflectional forms such as {\emph{brandas}}
`brand', {\emph{branda}} `brand', and {\emph{brandum}} `brand' under the
same entry (\citeproc{ref-BosworthToller1898}{Bosworth and Toller 1898,
116}).

The specific comparison form in this entry, {\emph{brandes}} `brand', is
the expected genitive singular of that a-stem noun. It is therefore an
inferred Old English paradigm form rather than the ordinary dictionary
headword.

\subsection{Development to Old
English}\label{development-to-old-english-75}

From {\Recon{brándas} `brand'}, the regular Old English development
passes through the usual unstressed-vowel weakening of the inflectional
ending, yielding {\emph{brandes}} `brand'. Nothing in the stem itself
requires a special repair. The root consonants and the stressed vowel
are the same as in the citation lemma {\emph{brand}} `brand'.

The analytical weight of the entry lies in the ending. By choosing the
oblique singular rather than the nominative citation form, the entry
presents the same lexeme in a different inherited cell.

\subsection{Form comparison}\label{form-comparison-4}

The comparison below sets the relevant forms side by side. It separates
the citation lemma from the selected oblique singular.

{\def\LTcaptype{none} % do not increment counter
\begin{longtable}[]{@{}
  >{\raggedright\arraybackslash}p{(\linewidth - 8\tabcolsep) * \real{0.2000}}
  >{\raggedright\arraybackslash}p{(\linewidth - 8\tabcolsep) * \real{0.2000}}
  >{\raggedright\arraybackslash}p{(\linewidth - 8\tabcolsep) * \real{0.2000}}
  >{\raggedright\arraybackslash}p{(\linewidth - 8\tabcolsep) * \real{0.2000}}
  >{\raggedright\arraybackslash}p{(\linewidth - 8\tabcolsep) * \real{0.2000}}@{}}
\toprule\noalign{}
\begin{minipage}[b]{\linewidth}\raggedright
Form / interpretation
\end{minipage} & \begin{minipage}[b]{\linewidth}\raggedright
Candidate input
\end{minipage} & \begin{minipage}[b]{\linewidth}\raggedright
Expected or documented OE outcome
\end{minipage} & \begin{minipage}[b]{\linewidth}\raggedright
OE comparison form
\end{minipage} & \begin{minipage}[b]{\linewidth}\raggedright
Result
\end{minipage} \\
\midrule\noalign{}
\endhead
\bottomrule\noalign{}
\endlastfoot
citation nominative singular & {\emph{*brándaz}} & expected OE lemma
{\emph{brand}} & {\emph{brand}} & regular headword-level outcome \\
genitive singular & {\emph{*brándas}} & regular output: {\emph{brandes}}
& {\emph{brandes}} & exact match for the oblique cell \\
\end{longtable}
}

The noun itself is straightforwardly inherited. The main point of the
entry is that {\emph{brandes}} `brand (gen.)' belongs to the same
regular a-stem paradigm as {\emph{brand}} `brand', even though the
citation lemma remains the nominative singular.

\section{\texorpdfstring{breast --- OE
\emph{brēost}}{breast --- OE brēost}}\label{breast-oe-brux113ost}

\index[iv]{02oe@\textbf{Old English}!breost@\emph{brēost}}
\index[iv]{03pgmc@\textbf{Proto-Germanic}!breusta@*bréustą}
\index[iv]{03pgmc@\textbf{Proto-Germanic}!brustz@*brústz}

Derivation: citation reconstruction \emph{*brústz}; form followed here
\emph{*bréustą} \textgreater{} \emph{brēost} (early analogy).

\subsection{Derivation trace}\label{derivation-trace-76}

Proto input: \emph{*bréustą}

\begingroup
\setlength{\fboxsep}{6pt}

\noindent\fbox{%
\begin{minipage}{0.97\linewidth}
\small
\begin{tabularx}{\linewidth}{@{}>{\raggedright\arraybackslash}p{0.440\linewidth}>{\centering\arraybackslash}X>{\raggedright\arraybackslash}p{0.440\linewidth}@{}}
\begin{minipage}[t]{\linewidth}
\centering\textbf{Earlier Germanic changes}\par
\vspace{0.35em}
\raggedright
\centering\textbf{}\par
\raggedright
\vspace{0.2em}
\begin{tabular}{@{}>{\raggedright\arraybackslash}p{0.68\linewidth}@{\hspace{0.55em}}>{\raggedright\arraybackslash}p{0.16\linewidth}@{\hspace{0.25em}}}
\multicolumn{2}{@{}l@{}}{*Proto-West Germanic*} \\
\multicolumn{2}{@{}l@{}}{[no change]} \\
\multicolumn{2}{@{}l@{}}{*Northwest Germanic*} \\
\multicolumn{2}{@{}l@{}}{[no change]} \\
\end{tabular}
\end{minipage}
&
&
\begin{minipage}[t]{\linewidth}
\centering\textbf{Old English changes}\par
\vspace{0.35em}
\raggedright
\centering\textbf{}\par
\raggedright
\vspace{0.2em}
\begin{tabular}{@{}>{\raggedright\arraybackslash}p{0.70\linewidth}@{\hspace{0.55em}}>{\raggedright\arraybackslash}p{0.20\linewidth}@{\hspace{0.25em}}}
\multicolumn{2}{@{}l@{}}{*Old English*} \\
OE Diphthong Leveling & \emph{*brēostą} \\
OE Heavy Syllable Nasal Apocope & \emph{*brēost} \\
\end{tabular}
\end{minipage}
\\
\end{tabularx}
\end{minipage}%
} \endgroup

Old English form: \emph{brēost}

\subsection{Reconstruction and comparative
evidence}\label{reconstruction-and-comparative-evidence-76}

The word family shows two related but distinct Proto-Germanic
formations. The root noun {\emph{*brust-}} lies behind forms such as
Gothic {\emph{brusts}} `breast', whereas Old English {\emph{brēost}}
`breast' belongs to a thematic formation {\emph{*breusta-}}, alongside
Old Norse {\emph{brjóst}} `breast' and Old Saxon {\emph{briost}}
`breast' (\citeproc{ref-Kroonen2013}{Kroonen 2013, 114};
\citeproc{ref-Orel2003}{Orel 2003, 95};
\citeproc{ref-RingeTaylor2014}{Ringe and Taylor 2014, 43}).

The derivational input {\Recon{bréustą} `breast'} therefore differs from
the citation label {\Recon{brústz} `breast'} because Old English
reflects the thematic branch rather than the root noun. The
morphological choice comes before the Old English sound changes
themselves.

\subsection{Old English evidence}\label{old-english-evidence-76}

Clark Hall records the noun as {\emph{brēost}} `breast' /
{\emph{breóst}} `breast' (\citeproc{ref-ClarkHall1960}{Clark Hall 1960,
65}). The form is an established Old English lexeme, not a reconstructed
target assembled from comparative evidence alone.

What requires explanation is not the Old English attestation but the
relation between that attested noun and the broader Germanic word
family. The relevant comparison form is therefore the thematic Old
English noun {\emph{brēost}} `breast'.

\subsection{Development to Old
English}\label{development-to-old-english-76}

From {\Recon{bréustą} `breast'}, the regular Old English development
gives {\emph{brēost}} `breast', with the expected \emph{eu}
\textgreater{} \emph{ēo} vowel history
(\citeproc{ref-Campbell1959}{Campbell 1959, sect. 115}). No special
repair is needed once the correct thematic formation is chosen.

The earlier mismatch arose only if the word was forced into the
root-noun line. The Old English noun itself continues the thematic
branch cleanly and directly.

\subsection{Formation comparison}\label{formation-comparison}

The comparison below sets the relevant forms side by side. It separates
the broader root-noun family label from the thematic formation actually
continued in Old English.

{\def\LTcaptype{none} % do not increment counter
\begin{longtable}[]{@{}
  >{\raggedright\arraybackslash}p{(\linewidth - 8\tabcolsep) * \real{0.2000}}
  >{\raggedright\arraybackslash}p{(\linewidth - 8\tabcolsep) * \real{0.2000}}
  >{\raggedright\arraybackslash}p{(\linewidth - 8\tabcolsep) * \real{0.2000}}
  >{\raggedright\arraybackslash}p{(\linewidth - 8\tabcolsep) * \real{0.2000}}
  >{\raggedright\arraybackslash}p{(\linewidth - 8\tabcolsep) * \real{0.2000}}@{}}
\toprule\noalign{}
\begin{minipage}[b]{\linewidth}\raggedright
Formation
\end{minipage} & \begin{minipage}[b]{\linewidth}\raggedright
Candidate input
\end{minipage} & \begin{minipage}[b]{\linewidth}\raggedright
Expected or documented OE outcome
\end{minipage} & \begin{minipage}[b]{\linewidth}\raggedright
OE comparison form
\end{minipage} & \begin{minipage}[b]{\linewidth}\raggedright
Result
\end{minipage} \\
\midrule\noalign{}
\endhead
\bottomrule\noalign{}
\endlastfoot
broader root-noun family & {\emph{*brústz}} & root-noun type outcomes
outside OE & non-OE comparanda & useful family label, but not the direct
source of \emph{brēost} \\
selected thematic formation & {\emph{*bréustą}} & regular output:
\emph{brēost} & {\emph{brēost}} & exact match between formation and
attested OE noun \\
\end{longtable}
}

The relevant point is the formation split. {\emph{brēost}} `breast' is
the regular Old English outcome of the thematic {\emph{*breusta-}}
branch, not of the root noun {\emph{*brust-}}.

\section{\texorpdfstring{craft --- OE
\emph{cræft}}{craft --- OE cræft}}\label{craft-oe-cruxe6ft}

\index[iv]{02oe@\textbf{Old English}!craeft@\emph{cræft}}
\index[iv]{03pgmc@\textbf{Proto-Germanic}!kraftaz@*kráftaz}
\index[iv]{03pgmc@\textbf{Proto-Germanic}!kraftiz@*kráftiz}

Derivation: citation reconstruction \emph{*kráftiz}; form followed here
\emph{*kráftaz} \textgreater{} \emph{cræft} (early analogy).

\subsection{Derivation trace}\label{derivation-trace-77}

Proto input: \emph{*kráftaz}

\begingroup
\setlength{\fboxsep}{6pt}

\noindent\fbox{%
\begin{minipage}{0.97\linewidth}
\small
\begin{tabularx}{\linewidth}{@{}>{\raggedright\arraybackslash}p{0.440\linewidth}>{\centering\arraybackslash}X>{\raggedright\arraybackslash}p{0.440\linewidth}@{}}
\begin{minipage}[t]{\linewidth}
\centering\textbf{Earlier Germanic changes}\par
\vspace{0.35em}
\raggedright
\centering\textbf{}\par
\raggedright
\vspace{0.2em}
\begin{tabular}{@{}>{\raggedright\arraybackslash}p{0.74\linewidth}@{\hspace{0.55em}}>{\raggedright\arraybackslash}p{0.16\linewidth}@{\hspace{0.25em}}}
\multicolumn{2}{@{}l@{}}{*Proto-West Germanic*} \\
\multicolumn{2}{@{}l@{}}{[no change]} \\
\multicolumn{2}{@{}l@{}}{*Northwest Germanic*} \\
\mbox{PGmc Final Z Deletion} & \emph{*kráfta} \\
\end{tabular}
\end{minipage}
&
&
\begin{minipage}[t]{\linewidth}
\centering\textbf{Old English changes}\par
\vspace{0.35em}
\raggedright
\centering\textbf{}\par
\raggedright
\vspace{0.2em}
\begin{tabular}{@{}>{\raggedright\arraybackslash}p{0.70\linewidth}@{\hspace{0.55em}}>{\raggedright\arraybackslash}p{0.16\linewidth}@{\hspace{0.25em}}}
\multicolumn{2}{@{}l@{}}{*Old English*} \\
PWGmc Final Bare A Loss & \emph{*kráft} \\
Anglo Frisian Brightening & \emph{*kræft} \\
\end{tabular}
\end{minipage}
\\
\end{tabularx}
\end{minipage}%
} \endgroup

Old English form: \emph{cræft}

\subsection{Reconstruction and comparative
evidence}\label{reconstruction-and-comparative-evidence-77}

Comparative sources disagree about the older stem class. Kroonen gives a
u-stem \emph{*kraftu-}, while Orel prints \Recon{kraftiz} `craft' and
\Recon{kraftuz} `craft' (\citeproc{ref-Kroonen2013}{Kroonen 2013, 340};
\citeproc{ref-Orel2003}{Orel 2003, 259}). The comparative label
\Recon{kráftiz} `craft' remains in view as a lexeme-level shorthand,
while \Recon{kráftaz} `craft' is the pre-Old-English form used for the
Old English derivation.

\subsection{Old English evidence}\label{old-english-evidence-77}

The Old English noun itself is secure. Clark Hall and Bosworth-Toller
both give \emph{cræft} `craft' as the headword
(\citeproc{ref-ClarkHall1960}{Clark Hall 1960, 19};
\citeproc{ref-BosworthToller1898}{Bosworth and Toller 1898, 145}).

\subsection{Development to Old
English}\label{development-to-old-english-77}

The comparison is between possible pre-Old-English inputs. The i-stem
comparator \Recon{kráftiz} `craft' gives \emph{creft}, while the u-stem
comparator \Recon{kráftuz} `craft' gives \emph{craft}. The a-stem-shaped
input \Recon{kráftaz} `craft' yields \emph{cræft} `craft' and is
therefore the form used for the Old English derivation. This does not
require treating the comparative dictionaries as identical; it shows the
narrower point that the Old English derivation needs a pre-Old-English
form without the i-umlaut trigger of \Recon{-iz} `craft' and without the
back-vowel outcome associated with the u-stem comparator.

\subsection{Form comparison}\label{form-comparison-5}

{\def\LTcaptype{none} % do not increment counter
\begin{longtable}[]{@{}lll@{}}
\toprule\noalign{}
Candidate input & OE output & Result \\
\midrule\noalign{}
\endhead
\bottomrule\noalign{}
\endlastfoot
*kráftiz & \emph{creft} & non-match; i-stem comparator \\
*kráftuz & \emph{craft} & non-match; u-stem comparator \\
*kráftaz & \emph{cræft} & exact match; selected pre-OE input \\
\end{longtable}
}

\section{\texorpdfstring{dill --- OE
\emph{dile}}{dill --- OE dile}}\label{dill-oe-dile}

\index[iv]{02oe@\textbf{Old English}!dile@\emph{dile}}
\index[iv]{03pgmc@\textbf{Proto-Germanic}!deliz@*déliz}
\index[iv]{03pgmc@\textbf{Proto-Germanic}!deljaz@*déljaz}

Derivation: citation reconstruction \emph{*déljaz}; form followed here
\emph{*déliz} \textgreater{} \emph{dile} (early analogy).

\subsection{Derivation trace}\label{derivation-trace-78}

Proto input: \emph{*déliz}

\begingroup
\setlength{\fboxsep}{6pt}

\noindent\fbox{%
\begin{minipage}{0.97\linewidth}
\small
\begin{tabularx}{\linewidth}{@{}>{\raggedright\arraybackslash}p{0.440\linewidth}>{\centering\arraybackslash}X>{\raggedright\arraybackslash}p{0.440\linewidth}@{}}
\begin{minipage}[t]{\linewidth}
\centering\textbf{Earlier Germanic changes}\par
\vspace{0.35em}
\raggedright
\centering\textbf{}\par
\raggedright
\vspace{0.2em}
\begin{tabular}{@{}>{\raggedright\arraybackslash}p{0.74\linewidth}@{\hspace{0.55em}}>{\raggedright\arraybackslash}p{0.16\linewidth}@{\hspace{0.25em}}}
\multicolumn{2}{@{}l@{}}{*Proto-West Germanic*} \\
\multicolumn{2}{@{}l@{}}{[no change]} \\
\multicolumn{2}{@{}l@{}}{*Northwest Germanic*} \\
\mbox{PGmc Final Z Deletion} & \emph{*déli} \\
\end{tabular}
\end{minipage}
&
&
\begin{minipage}[t]{\linewidth}
\centering\textbf{Old English changes}\par
\vspace{0.35em}
\raggedright
\centering\textbf{}\par
\raggedright
\vspace{0.2em}
\begin{tabular}{@{}>{\raggedright\arraybackslash}p{0.74\linewidth}@{\hspace{0.55em}}>{\raggedright\arraybackslash}p{0.16\linewidth}@{\hspace{0.25em}}}
\multicolumn{2}{@{}l@{}}{*Old English*} \\
OE I Umlaut & \emph{*dili} \\
OE Med Unstressed I Lowering1 & \emph{*dile} \\
\end{tabular}
\end{minipage}
\\
\end{tabularx}
\end{minipage}%
} \endgroup

Old English form: \emph{dile}

\subsection{Reconstruction and comparative
evidence}\label{reconstruction-and-comparative-evidence-78}

Comparative evidence preserves both an i-stem and a ja-stem formation,
with Old English {\emph{dile}} `dill' on one side and continental forms
such as Old Saxon {\emph{dilli}} `dill' and Old High German
{\emph{tilli}} `dill' on the other (\citeproc{ref-Fulk2018}{Fulk 2018,
170}). The derivational input {\Recon{déliz} `dill'} therefore
represents the i-stem side of the paradigm, whereas the citation label
{\Recon{déljaz} `dill'} is a broader comparative headword.

The stem class determines the Old English consonant shape. A ja-stem
with \emph{*-lj-} would be expected to produce gemination, but the Old
English noun shows a single \emph{l}. Fulk's discussion of ja-stems
transferred to the i-stems provides the relevant morphological
background for the OE side (\citeproc{ref-Fulk2018}{Fulk 2018, 170}).

\subsection{Old English evidence}\label{old-english-evidence-78}

Old English dictionaries record the plant name as {\emph{dile}} `dill',
alongside the variant {\emph{dili}} `dill'
(\citeproc{ref-BosworthToller1898}{Bosworth and Toller 1898, 164};
\citeproc{ref-ClarkHall1960}{Clark Hall 1960, 95}). The form discussed
here is therefore an attested Old English noun with single \emph{l}.

The Old English evidence is the relevant point. Whatever broader
comparative headword is chosen for the family, the inherited form
reflected in OE is the i-stem type {\emph{dile}} `dill', not a geminated
{\emph{dill}} `dill' outcome.

\subsection{Development to Old
English}\label{development-to-old-english-78}

From {\Recon{déliz} `dill'}, regular loss of final \emph{z} and the
later lowering of unstressed \emph{i} yield {\emph{dile}} `dill'. The
stem itself remains ungeminated throughout that path.

The contrast is morphological rather than phonological. If the word were
forced through a ja-stem \emph{*-lj-} pathway, the expected result would
show \emph{ll}. The attested Old English noun instead matches the i-stem
development.

\subsection{Formation comparison}\label{formation-comparison-1}

The comparison below sets the relevant forms side by side. It separates
the broader comparative headword from the stem class actually reflected
in Old English.

{\def\LTcaptype{none} % do not increment counter
\begin{longtable}[]{@{}
  >{\raggedright\arraybackslash}p{(\linewidth - 8\tabcolsep) * \real{0.2000}}
  >{\raggedright\arraybackslash}p{(\linewidth - 8\tabcolsep) * \real{0.2000}}
  >{\raggedright\arraybackslash}p{(\linewidth - 8\tabcolsep) * \real{0.2000}}
  >{\raggedright\arraybackslash}p{(\linewidth - 8\tabcolsep) * \real{0.2000}}
  >{\raggedright\arraybackslash}p{(\linewidth - 8\tabcolsep) * \real{0.2000}}@{}}
\toprule\noalign{}
\begin{minipage}[b]{\linewidth}\raggedright
Formation
\end{minipage} & \begin{minipage}[b]{\linewidth}\raggedright
Candidate input
\end{minipage} & \begin{minipage}[b]{\linewidth}\raggedright
Expected or documented OE outcome
\end{minipage} & \begin{minipage}[b]{\linewidth}\raggedright
OE comparison form
\end{minipage} & \begin{minipage}[b]{\linewidth}\raggedright
Result
\end{minipage} \\
\midrule\noalign{}
\endhead
\bottomrule\noalign{}
\endlastfoot
comparative ja-stem label & {\emph{*déljaz}} & ja-stem type outcome with
gemination & {\emph{dill}}-type comparison & useful comparative label,
but not the OE form \\
selected i-stem formation & {\emph{*déliz}} & regular output:
{\emph{dile}} & {\emph{dile}} & exact match between formation and
attested OE noun \\
\end{longtable}
}

The single \emph{l} is the decisive diagnostic. It identifies
{\emph{dile}} `dill' with the i-stem formation rather than with the
continental ja-stem branch.

\section{\texorpdfstring{fast --- OE
\emph{festan}}{fast --- OE festan}}\label{fast-oe-festan}

\index[iv]{02oe@\textbf{Old English}!festan@\emph{festan}}
\index[iv]{03pgmc@\textbf{Proto-Germanic}!fastena@*fastēną}
\index[iv]{03pgmc@\textbf{Proto-Germanic}!fastijana@*fástijaną}

Derivation: citation reconstruction \emph{*fastēną}; form followed here
\emph{*fástijaną} \textgreater{} \emph{festan} (early analogy).

\subsection{Derivation trace}\label{derivation-trace-79}

Proto input: \emph{*fástijaną}

\begingroup
\setlength{\fboxsep}{6pt}

\noindent\fbox{%
\begin{minipage}{0.97\linewidth}
\footnotesize
\begin{tabularx}{\linewidth}{@{}>{\raggedright\arraybackslash}p{0.320\linewidth}>{\centering\arraybackslash}X>{\raggedright\arraybackslash}p{0.520\linewidth}@{}}
\begin{minipage}[t]{\linewidth}
\centering\textbf{Earlier Germanic changes}\par
\vspace{0.35em}
\raggedright
\centering\textbf{}\par
\raggedright
\vspace{0.2em}
\begin{tabular}{@{}>{\raggedright\arraybackslash}p{0.68\linewidth}@{\hspace{0.55em}}>{\raggedright\arraybackslash}p{0.16\linewidth}@{\hspace{0.25em}}}
\multicolumn{2}{@{}l@{}}{*Proto-West Germanic*} \\
\multicolumn{2}{@{}l@{}}{[no change]} \\
\multicolumn{2}{@{}l@{}}{*Northwest Germanic*} \\
\multicolumn{2}{@{}l@{}}{[no change]} \\
\end{tabular}
\end{minipage}
&
&
\begin{minipage}[t]{\linewidth}
\centering\textbf{Old English changes}\par
\vspace{0.35em}
\raggedright
\centering\textbf{}\par
\raggedright
\vspace{0.2em}
\begin{tabular}{@{}>{\raggedright\arraybackslash}p{0.64\linewidth}@{\hspace{0.55em}}>{\raggedright\arraybackslash}p{0.26\linewidth}@{\hspace{0.25em}}}
\multicolumn{2}{@{}l@{}}{*Old English*} \\
Anglo Frisian Brightening & \emph{*fæstijaną} \\
OE Heavy Syllable Nasal Apocope & \emph{*fæstijan} \\
OE Secondary Nasalization & \emph{*fæstijąn} \\
Sievers Law Syncope & \emph{*fæstjąn} \\
OE I Umlaut & \emph{*festjąn} \\
OE Weak Tail Reduction & \emph{*festjan} \\
OE J Loss After Heavy & \emph{*festan} \\
\end{tabular}
\end{minipage}
\\
\end{tabularx}
\end{minipage}%
} \endgroup

Old English form: \emph{festan}

\subsection{Reconstruction and comparative
evidence}\label{reconstruction-and-comparative-evidence-79}

Kroonen places the verb within the wider {\emph{*fastu-}} adjective
family and its derived {\emph{*fasten-}} verbal line, the comparative
background behind Old English `to fast'
(\citeproc{ref-Kroonen2013}{Kroonen 2013, 171}). Ringe and Taylor,
however, distinguish the Old English verb more closely: they treat OE
`to fast' as originally a class-I weak verb that later acquired the
stative meaning through lexical confusion
(\citeproc{ref-RingeTaylor2014}{Ringe and Taylor 2014, 110}).

The derivational input {\Recon{fástijaną} `fast'} therefore represents
the inherited class-I formation reflected in Old English, whereas the
citation label {\Recon{fastēną} `fast'} belongs to the broader
comparative presentation of the lexeme.

\subsection{Old English evidence}\label{old-english-evidence-79}

Old English dictionaries record forms such as {\emph{festan}} `fast',
alongside related {\emph{fæstan}} `fast' / {\emph{fǣstan}} `fast'
spellings and meanings (\citeproc{ref-BosworthToller1898}{Bosworth and
Toller 1898, 213}). The form selected here is {\emph{festan}}, which
fits the regular class-I phonological development.

The \emph{æ}-forms remain relevant, but they do not control the entry.
In the present analysis they belong to a later analogical reshaping
under the adjective {\emph{fæst}} `fast', whereas {\emph{festan}} `fast'
is the regular inherited class-I comparison form.

\subsection{Development to Old
English}\label{development-to-old-english-79}

From {\Recon{fástijaną} `fast'}, Anglo-Frisian brightening and
subsequent i-umlaut produce the fronted vowel seen in {\emph{festan}}
`fast'. The later weak-tail reductions and loss of \emph{j} after a
heavy syllable complete the regular Old English outcome.

What makes the entry non-regular is not the phonology of {\emph{festan}}
`fast' itself, but the choice of formation. Old English continues the
class-I verb, even though the comparative headword is often given under
the parallel {\emph{*fastēn-}} family.

\subsection{Class comparison}\label{class-comparison}

The comparison below sets the relevant forms side by side. It
distinguishes the comparative class-III headword from the class-I
formation actually reflected in Old English.

{\def\LTcaptype{none} % do not increment counter
\begin{longtable}[]{@{}
  >{\raggedright\arraybackslash}p{(\linewidth - 8\tabcolsep) * \real{0.2000}}
  >{\raggedright\arraybackslash}p{(\linewidth - 8\tabcolsep) * \real{0.2000}}
  >{\raggedright\arraybackslash}p{(\linewidth - 8\tabcolsep) * \real{0.2000}}
  >{\raggedright\arraybackslash}p{(\linewidth - 8\tabcolsep) * \real{0.2000}}
  >{\raggedright\arraybackslash}p{(\linewidth - 8\tabcolsep) * \real{0.2000}}@{}}
\toprule\noalign{}
\begin{minipage}[b]{\linewidth}\raggedright
Formation / class
\end{minipage} & \begin{minipage}[b]{\linewidth}\raggedright
Candidate input
\end{minipage} & \begin{minipage}[b]{\linewidth}\raggedright
Expected or documented OE outcome
\end{minipage} & \begin{minipage}[b]{\linewidth}\raggedright
OE comparison form
\end{minipage} & \begin{minipage}[b]{\linewidth}\raggedright
Result
\end{minipage} \\
\midrule\noalign{}
\endhead
\bottomrule\noalign{}
\endlastfoot
comparative class-III headword & {\emph{*fastēną}} & class-III type
outcome, not {\emph{festan}} & wider family context & useful family
label, but not the direct source of the target \\
selected class-I weak verb & {\emph{*fástijaną}} & regular output:
{\emph{festan}} & {\emph{festan}} & exact match between formation and
attested OE verb \\
later analogical reshaping & adjective-driven {\emph{fæst}} influence &
{\emph{fæstan}} / {\emph{fǣstan}} type spellings & fæstan-type evidence
& genuine later OE reshaping, but secondary to the Old English form
here \\
\end{longtable}
}

The relevant point is the class split. {\emph{festan}} `fast' is the
regular Old English outcome of the class-I formation, while the
better-known \emph{æ}-forms belong to a later analogical layer.

\section{\texorpdfstring{flask --- OE
\emph{flasce}}{flask --- OE flasce}}\label{flask-oe-flasce}

\index[iv]{02oe@\textbf{Old English}!flasce@\emph{flasce}}
\index[iv]{03pgmc@\textbf{Proto-Germanic}!flasko@*flaskō}
\index[iv]{03pgmc@\textbf{Proto-Germanic}!flaskon@*fláskōn}

Derivation: citation reconstruction \emph{*flaskō}; form followed here
\emph{*fláskōn} \textgreater{} \emph{flasce} (early analogy).

\subsection{Derivation trace}\label{derivation-trace-80}

Proto input: \emph{*fláskōn}

\begingroup
\setlength{\fboxsep}{6pt}

\noindent\fbox{%
\begin{minipage}{0.97\linewidth}
\small
\begin{tabularx}{\linewidth}{@{}>{\raggedright\arraybackslash}p{0.440\linewidth}>{\centering\arraybackslash}X>{\raggedright\arraybackslash}p{0.440\linewidth}@{}}
\begin{minipage}[t]{\linewidth}
\centering\textbf{Earlier Germanic changes}\par
\vspace{0.35em}
\raggedright
\centering\textbf{}\par
\raggedright
\vspace{0.2em}
\begin{tabular}{@{}>{\raggedright\arraybackslash}p{0.68\linewidth}@{\hspace{0.55em}}>{\raggedright\arraybackslash}p{0.16\linewidth}@{\hspace{0.25em}}}
\multicolumn{2}{@{}l@{}}{*Proto-West Germanic*} \\
\multicolumn{2}{@{}l@{}}{[no change]} \\
\multicolumn{2}{@{}l@{}}{*Northwest Germanic*} \\
NWGmc N Stem N Loss & \emph{*fláskǭ} \\
\end{tabular}
\end{minipage}
&
&
\begin{minipage}[t]{\linewidth}
\centering\textbf{Old English changes}\par
\vspace{0.35em}
\raggedright
\centering\textbf{}\par
\raggedright
\vspace{0.2em}
\begin{tabular}{@{}>{\raggedright\arraybackslash}p{0.74\linewidth}@{\hspace{0.55em}}>{\raggedright\arraybackslash}p{0.16\linewidth}@{\hspace{0.25em}}}
\multicolumn{2}{@{}l@{}}{*Old English*} \\
Anglo Frisian Brightening & \emph{*flæskǭ} \\
OE A Restoration & \emph{*flaskǭ} \\
OE Unstressed Long Vowel Shortening & \emph{*flaskæ} \\
OE Unstressed AE Merger & \emph{*flaske} \\
\end{tabular}
\end{minipage}
\\
\end{tabularx}
\end{minipage}%
} \endgroup

Old English form: \emph{flasce}

\subsection{Reconstruction and comparative
evidence}\label{reconstruction-and-comparative-evidence-80}

The wider Germanic family is often cited under a form such as
{\Recon{flaskō} `flask'}, but the evidence relevant for Old English
points instead to a weak feminine formation {\Recon{fláskōn} `flask'} /
{\Recon{flaskǭ} `flask'} (\citeproc{ref-Orel2003}{Orel 2003, 104}). That
distinction is crucial for the suffixal history of the noun.

The derivational input therefore differs from the citation label in stem
class. Old English {\emph{flasce}} `flask' belongs with the weak
feminine line, and the plural or oblique forms {\emph{flascan}} `flask'
support that analysis (\citeproc{ref-RingeTaylor2014}{Ringe and Taylor
2014, 192}).

\subsection{Old English evidence}\label{old-english-evidence-80}

Old English dictionaries record the noun as {\emph{flasce}} `flask',
with inflectional support from forms such as {\emph{flascan}} `flask'; a
later West Saxon {\emph{flaxe}} `flask' is also noted as a secondary
variant (\citeproc{ref-BosworthToller1898}{Bosworth and Toller 1898,
235}; \citeproc{ref-ClarkHall1960}{Clark Hall 1960, 121}).

The relevant comparison form is therefore the weak feminine noun
{\emph{flasce}} `flask'. The plural and oblique evidence helps explain
why the vowel and ending are preserved as they are in the singular.

\subsection{Development to Old
English}\label{development-to-old-english-80}

From {\Recon{fláskōn} `flask'}, the weak feminine passes through the
expected loss of \emph{n} and the later Old English development of the
unstressed ending, reaching {\emph{flasce}} `flask'. Campbell cites
restored \emph{a} in exactly this environment, including {\emph{flasce}}
`flask' after inflected {\emph{flascan}} `flask'
(\citeproc{ref-Campbell1959}{Campbell 1959, sect. 158}). Once the weak
feminine formation is chosen, the noun follows a regular path to its Old
English shape.

The decisive issue is morphological rather than phonological. A simple
strong feminine citation form does not capture the OE weak noun as
cleanly as the selected \Recon{fláskōn} `flask' does.

\subsection{Formation comparison}\label{formation-comparison-2}

The comparison below sets the relevant forms side by side. It separates
the broader comparative headword from the weak feminine formation
actually reflected in Old English.

{\def\LTcaptype{none} % do not increment counter
\begin{longtable}[]{@{}
  >{\raggedright\arraybackslash}p{(\linewidth - 8\tabcolsep) * \real{0.2000}}
  >{\raggedright\arraybackslash}p{(\linewidth - 8\tabcolsep) * \real{0.2000}}
  >{\raggedright\arraybackslash}p{(\linewidth - 8\tabcolsep) * \real{0.2000}}
  >{\raggedright\arraybackslash}p{(\linewidth - 8\tabcolsep) * \real{0.2000}}
  >{\raggedright\arraybackslash}p{(\linewidth - 8\tabcolsep) * \real{0.2000}}@{}}
\toprule\noalign{}
\begin{minipage}[b]{\linewidth}\raggedright
Formation
\end{minipage} & \begin{minipage}[b]{\linewidth}\raggedright
Candidate input
\end{minipage} & \begin{minipage}[b]{\linewidth}\raggedright
Expected or documented OE outcome
\end{minipage} & \begin{minipage}[b]{\linewidth}\raggedright
OE comparison form
\end{minipage} & \begin{minipage}[b]{\linewidth}\raggedright
Result
\end{minipage} \\
\midrule\noalign{}
\endhead
\bottomrule\noalign{}
\endlastfoot
broader comparative headword & {\emph{*flaskō}} & broader family label &
wider family context & useful lexeme label, but not the cleanest
OE-facing derivation \\
selected weak feminine formation & {\emph{*fláskōn}} & regular output:
{\emph{flasce}} & {\emph{flasce}} & exact match between formation and
attested OE noun \\
\end{longtable}
}

The weak feminine suffix is the relevant point. It aligns the inherited
form with attested {\emph{flasce}} `flask' and its supporting paradigm
forms.

\section{\texorpdfstring{follow --- OE
\emph{fylġan}}{follow --- OE fylġan}}\label{follow-oe-fylux121an}

\index[iv]{02oe@\textbf{Old English}!fylgan@\emph{fylġan}}
\index[iv]{03pgmc@\textbf{Proto-Germanic}!fulgena@*fulgēną}
\index[iv]{03pgmc@\textbf{Proto-Germanic}!fulgijana@*fúlgijaną}

Derivation: citation reconstruction \emph{*fulgēną}; form followed here
\emph{*fúlgijaną} \textgreater{} \emph{fylġan} (early analogy).

\subsection{Derivation trace}\label{derivation-trace-81}

Proto input: \emph{*fúlgijaną}

\begingroup
\setlength{\fboxsep}{6pt}

\noindent\fbox{%
\begin{minipage}{0.97\linewidth}
\footnotesize
\begin{tabularx}{\linewidth}{@{}>{\raggedright\arraybackslash}p{0.320\linewidth}>{\centering\arraybackslash}X>{\raggedright\arraybackslash}p{0.520\linewidth}@{}}
\begin{minipage}[t]{\linewidth}
\centering\textbf{Earlier Germanic changes}\par
\vspace{0.35em}
\raggedright
\centering\textbf{}\par
\raggedright
\vspace{0.2em}
\begin{tabular}{@{}>{\raggedright\arraybackslash}p{0.68\linewidth}@{\hspace{0.55em}}>{\raggedright\arraybackslash}p{0.16\linewidth}@{\hspace{0.25em}}}
\multicolumn{2}{@{}l@{}}{*Proto-West Germanic*} \\
\multicolumn{2}{@{}l@{}}{[no change]} \\
\multicolumn{2}{@{}l@{}}{*Northwest Germanic*} \\
\multicolumn{2}{@{}l@{}}{[no change]} \\
\end{tabular}
\end{minipage}
&
&
\begin{minipage}[t]{\linewidth}
\centering\textbf{Old English changes}\par
\vspace{0.35em}
\raggedright
\centering\textbf{}\par
\raggedright
\vspace{0.2em}
\begin{tabular}{@{}>{\raggedright\arraybackslash}p{0.66\linewidth}@{\hspace{0.55em}}>{\raggedright\arraybackslash}p{0.24\linewidth}@{\hspace{0.25em}}}
\multicolumn{2}{@{}l@{}}{*Old English*} \\
OE Heavy Syllable Nasal Apocope & \emph{*fúlgijan} \\
OE Secondary Nasalization & \emph{*fúlgijąn} \\
Sievers Law Syncope & \emph{*fúlgjąn} \\
OE Velar Palatalization & \emph{*fúlʤjąn} \\
OE I Umlaut & \emph{*fylʤjąn} \\
OE Weak Tail Reduction & \emph{*fylʤjan} \\
OE J Loss After Heavy & \emph{*fylʤan} \\
\end{tabular}
\end{minipage}
\\
\end{tabularx}
\end{minipage}%
} \endgroup

Old English form: \emph{fylġan}

\subsection{Reconstruction and comparative
evidence}\label{reconstruction-and-comparative-evidence-81}

Kroonen reconstructs the verb as \emph{*fulgen-} and gives Old English
{\emph{fylgan}} `follow', {\emph{folgian}} `follow', adding that Old
Norse {\emph{fylgja}} `follow' and Old English {\emph{fylg(e)an}}
continue a formation \emph{*fulgjan-}
(\citeproc{ref-Kroonen2013}{Kroonen 2013, 159}). The comparative
headword and the class-I formation are therefore related but not
identical.

Ringe and Taylor distinguish PNWGmc \Recon{fulgija-} `follow'
\textasciitilde{} \Recon{fulgai-} `follow' \textgreater{} OE
{\emph{fylgan}} `follow' \textasciitilde{} {\emph{folgian}} `follow' and
describe it as a dual formation that probably reflects an older
alternation between j-present and e-stative
(\citeproc{ref-RingeTaylor2014}{Ringe and Taylor 2014, 293--94}). This
is a stem-class choice, not a spelling choice. The derivational input
{\Recon{fúlgijaną} `follow'} belongs to the class-I \emph{*fulgija-} /
\emph{*fulgjan-} branch; the citation form {\Recon{fulgēną} `follow'}
belongs to the parallel class-II history behind {\emph{folgian}}
`follow'.

\subsection{Old English evidence}\label{old-english-evidence-81}

The Old English evidence preserves both formations. Clark Hall lists
{\emph{fylgan}} `follow' with variant spellings {\emph{fylgian}}
`follow' and {\emph{fyligan}} `follow'
(\citeproc{ref-ClarkHall1960}{Clark Hall 1960, 125}). Bosworth-Toller
likewise has a separate {\emph{fylgean}} `follow' entry
(\citeproc{ref-BosworthToller1898}{Bosworth and Toller 1898, 275}).

Bright notes traces of the older conjugation in {\emph{fylg(e)an}}
`follow' (\citeproc{ref-BrightCassidyRingler1971}{Bright 1971, 77}) and
lists {\emph{folgian}} `follow' ({\emph{fylgean}} `follow') in the
glossary (\citeproc{ref-BrightCassidyRingler1971}{Bright 1971, 364}).
The relevant comparison form in this entry is therefore the class-I verb
{\emph{fylgan}} `follow' / {\emph{fylgean}}, here normalized as
{\emph{fylġan}} `follow'. The spelling with ġ represents the palatalized
velar before a front-vocalic environment.

\subsection{Development to Old
English}\label{development-to-old-english-81}

{\Recon{fúlgijaną} `follow'} is a class-I weak-verb formation. In the
class-I branch the \emph{*j} blocks NWGmc lowering of \emph{u} to
\emph{o}, since Ringe and Taylor formulate that lowering for
environments in which no \emph{*j} intervened
(\citeproc{ref-RingeTaylor2014}{Ringe and Taylor 2014, 96}). The same
front-vocalic environment then triggers i-umlaut, so \emph{u} becomes
\emph{y} (\citeproc{ref-RingeTaylor2014}{Ringe and Taylor 2014, sect.
6.6.2}).

The subsequent Old English developments are palatalization of the velar,
weak-tail reduction, and loss of \emph{j} after a heavy syllable,
yielding {\emph{fylġan}} `follow'. This is the regular outcome of the
class-I formation. The class-II form {\emph{folgian}} `follow' belongs
to the parallel \emph{*-ē-} / \emph{*-ai-} branch and is not the form
modeled here.

\subsection{Class comparison}\label{class-comparison-1}

A class comparison identifies which inherited formation corresponds to
the established Old English form under discussion. The comparison below
is manual; no full automatic class probe is presented here.

{\def\LTcaptype{none} % do not increment counter
\begin{longtable}[]{@{}
  >{\raggedright\arraybackslash}p{(\linewidth - 8\tabcolsep) * \real{0.2000}}
  >{\raggedright\arraybackslash}p{(\linewidth - 8\tabcolsep) * \real{0.2000}}
  >{\raggedright\arraybackslash}p{(\linewidth - 8\tabcolsep) * \real{0.2000}}
  >{\raggedright\arraybackslash}p{(\linewidth - 8\tabcolsep) * \real{0.2000}}
  >{\raggedright\arraybackslash}p{(\linewidth - 8\tabcolsep) * \real{0.2000}}@{}}
\toprule\noalign{}
\begin{minipage}[b]{\linewidth}\raggedright
Formation / class
\end{minipage} & \begin{minipage}[b]{\linewidth}\raggedright
Candidate input
\end{minipage} & \begin{minipage}[b]{\linewidth}\raggedright
Expected or documented OE outcome
\end{minipage} & \begin{minipage}[b]{\linewidth}\raggedright
OE comparison form
\end{minipage} & \begin{minipage}[b]{\linewidth}\raggedright
Result
\end{minipage} \\
\midrule\noalign{}
\endhead
\bottomrule\noalign{}
\endlastfoot
citation class-II formation & {\emph{*fulgēną}} & probe output:
{\emph{folgon}} & {\emph{folgian}} `follow' & mismatch: the regular
output is not the remodeled infinitive \emph{folgian} \\
parallel class-II branch & PNWGmc \emph{*fulgai-} & Ringe-Taylor: OE
{\emph{folgian}} `follow' & {\emph{folgian}} `follow' & documents the
separate class-II branch, but not the target of this entry \\
selected class-I formation & {\emph{*fúlgijaną}} & regular output:
{\emph{fylġan}} `follow' & {\emph{fylġan}} `follow' / {\emph{fylgan}}
`follow' & exact match between input, output, and class \\
\end{longtable}
}

The relevant point is the class split. {\emph{fylġan}} `follow' is the
regular Old English outcome of the class-I \emph{*fulgija-} /
\emph{*fulgjan-} formation, whereas {\emph{folgian}} `follow' belongs to
the parallel class-II branch.

\section{\texorpdfstring{gall --- OE
\emph{ġealla}}{gall --- OE ġealla}}\label{gall-oe-ux121ealla}

\index[iv]{02oe@\textbf{Old English}!gealla@\emph{ġealla}}
\index[iv]{03pgmc@\textbf{Proto-Germanic}!galla@*gállą}
\index[iv]{03pgmc@\textbf{Proto-Germanic}!gallo@*gállô}

Derivation: citation reconstruction \emph{*gállą}; form followed here
\emph{*gállô} \textgreater{} \emph{ġealla} (early analogy).

\subsection{Derivation trace}\label{derivation-trace-82}

Proto input: \emph{*gállô}

\begingroup
\setlength{\fboxsep}{6pt}

\noindent\fbox{%
\begin{minipage}{0.97\linewidth}
\small
\begin{tabularx}{\linewidth}{@{}>{\raggedright\arraybackslash}p{0.320\linewidth}>{\centering\arraybackslash}X>{\raggedright\arraybackslash}p{0.520\linewidth}@{}}
\begin{minipage}[t]{\linewidth}
\centering\textbf{Earlier Germanic changes}\par
\vspace{0.35em}
\raggedright
\centering\textbf{}\par
\raggedright
\vspace{0.2em}
\begin{tabular}{@{}>{\raggedright\arraybackslash}p{0.68\linewidth}@{\hspace{0.55em}}>{\raggedright\arraybackslash}p{0.16\linewidth}@{\hspace{0.25em}}}
\multicolumn{2}{@{}l@{}}{*Proto-West Germanic*} \\
\multicolumn{2}{@{}l@{}}{[no change]} \\
\multicolumn{2}{@{}l@{}}{*Northwest Germanic*} \\
\multicolumn{2}{@{}l@{}}{[no change]} \\
\end{tabular}
\end{minipage}
&
&
\begin{minipage}[t]{\linewidth}
\centering\textbf{Old English changes}\par
\vspace{0.35em}
\raggedright
\centering\textbf{}\par
\raggedright
\vspace{0.2em}
\begin{tabular}{@{}>{\raggedright\arraybackslash}p{0.74\linewidth}@{\hspace{0.55em}}>{\raggedright\arraybackslash}p{0.16\linewidth}@{\hspace{0.25em}}}
\multicolumn{2}{@{}l@{}}{*Old English*} \\
Anglo Frisian Brightening & \emph{*gællô} \\
OE Breaking & \emph{*geallô} \\
OE Velar Palatalization & \emph{*ʤeallô} \\
OE Unstressed Long Vowel Shortening & \emph{*ʤealla} \\
\end{tabular}
\end{minipage}
\\
\end{tabularx}
\end{minipage}%
} \endgroup

Old English form: \emph{ġealla}

\subsection{Reconstruction and comparative
evidence}\label{reconstruction-and-comparative-evidence-82}

The wider cognate family can be presented under a form such as
{\Recon{gállą} `gall'}, but the Old English noun itself belongs with a
weak noun \emph{*gallōn-} `gall', cited here as {\Recon{gállô} `gall'}
(\citeproc{ref-Kroonen2013}{Kroonen 2013, 165}). The derivational input
therefore differs from the broader comparative headword in stem class.

The stem class determines the Old English shape. The weak masculine
pathway preserves the ending needed for \emph{ġealla} `gall', whereas a
simple strong-noun headword does not align as closely with the attested
OE noun.

\subsection{Old English evidence}\label{old-english-evidence-82}

Old English dictionaries record the noun as {\emph{gealla}} `gall', and
Bright also gives the dative {\emph{geallan}} `gall', confirming a
weak-noun paradigm (\citeproc{ref-BosworthToller1898}{Bosworth and
Toller 1898, 297}; \citeproc{ref-ClarkHall1960}{Clark Hall 1960, 145};
\citeproc{ref-BrightCassidyRingler1971}{Bright 1971, 372}). The
normalized spelling {\emph{ġealla}} `gall' uses ġ for the palatal
consonant.

Campbell also notes dialectal variation, contrasting West Saxon or
Kentish {\emph{gealla}} `gall' with Anglian {\emph{galla}} `gall'
(\citeproc{ref-Campbell1959}{Campbell 1959, sect. 486}). The target of
this entry is the West Saxon type {\emph{ġealla}} `gall'.

\subsection{Development to Old
English}\label{development-to-old-english-82}

From {\Recon{gállô} `gall'}, the weak noun develops through the expected
Old English history of the suffix and the regular breaking environment
before \emph{ll}, yielding {\emph{ġealla}} `gall'
(\citeproc{ref-Campbell1959}{Campbell 1959, sect. 486}). Once the weak
masculine input is chosen, the noun follows a regular path to its
attested Old English form.

The decisive issue is therefore morphological. Old English reflects the
weak noun, while the broader family label belongs to a different way of
presenting the cognate set.

\subsection{Stem comparison}\label{stem-comparison}

The comparison below sets the relevant forms side by side. It separates
the broader comparative headword from the weak noun formation actually
reflected in Old English.

{\def\LTcaptype{none} % do not increment counter
\begin{longtable}[]{@{}
  >{\raggedright\arraybackslash}p{(\linewidth - 8\tabcolsep) * \real{0.2000}}
  >{\raggedright\arraybackslash}p{(\linewidth - 8\tabcolsep) * \real{0.2000}}
  >{\raggedright\arraybackslash}p{(\linewidth - 8\tabcolsep) * \real{0.2000}}
  >{\raggedright\arraybackslash}p{(\linewidth - 8\tabcolsep) * \real{0.2000}}
  >{\raggedright\arraybackslash}p{(\linewidth - 8\tabcolsep) * \real{0.2000}}@{}}
\toprule\noalign{}
\begin{minipage}[b]{\linewidth}\raggedright
Formation
\end{minipage} & \begin{minipage}[b]{\linewidth}\raggedright
Candidate input
\end{minipage} & \begin{minipage}[b]{\linewidth}\raggedright
Expected or documented OE outcome
\end{minipage} & \begin{minipage}[b]{\linewidth}\raggedright
OE comparison form
\end{minipage} & \begin{minipage}[b]{\linewidth}\raggedright
Result
\end{minipage} \\
\midrule\noalign{}
\endhead
\bottomrule\noalign{}
\endlastfoot
broader family label & {\emph{*gállą}} & broader cognate-set headword &
wider family context & useful lexeme label, but not the direct source of
\emph{ġealla} \\
selected weak noun & {\emph{*gállô}} & regular output: {\emph{ġealla}}
`gall' & {\emph{ġealla}} `gall' & exact match between formation and
attested OE noun \\
dialectal Anglian continuation & weak noun branch & Anglian
{\emph{galla}} `gall' type & {\emph{galla}} `gall' & genuine OE variant,
but not the West Saxon form used here \\
\end{longtable}
}

The weak-noun stem class is the relevant point. It gives a direct route
to attested \emph{ġealla} `gall', while the broader comparative label
serves only as a family heading.

\section{\texorpdfstring{knight --- OE
\emph{cniht}}{knight --- OE cniht}}\label{knight-oe-cniht}

\index[iv]{02oe@\textbf{Old English}!cniht@\emph{cniht}}
\index[iv]{03pgmc@\textbf{Proto-Germanic}!knextaz@*knéxtaz}
\index[iv]{03pgmc@\textbf{Proto-Germanic}!knixtaz@*kníxtaz}

Derivation: citation reconstruction \emph{*kníxtaz}; form followed here
\emph{*knéxtaz} \textgreater{} \emph{cniht} (early analogy).

\subsection{Derivation trace}\label{derivation-trace-83}

Proto input: \emph{*knéxtaz}

\begingroup
\setlength{\fboxsep}{6pt}

\noindent\fbox{%
\begin{minipage}{0.97\linewidth}
\small
\begin{tabularx}{\linewidth}{@{}>{\raggedright\arraybackslash}p{0.440\linewidth}>{\centering\arraybackslash}X>{\raggedright\arraybackslash}p{0.440\linewidth}@{}}
\begin{minipage}[t]{\linewidth}
\centering\textbf{Earlier Germanic changes}\par
\vspace{0.35em}
\raggedright
\centering\textbf{}\par
\raggedright
\vspace{0.2em}
\begin{tabular}{@{}>{\raggedright\arraybackslash}p{0.74\linewidth}@{\hspace{0.55em}}>{\raggedright\arraybackslash}p{0.16\linewidth}@{\hspace{0.25em}}}
\multicolumn{2}{@{}l@{}}{*Proto-West Germanic*} \\
\multicolumn{2}{@{}l@{}}{[no change]} \\
\multicolumn{2}{@{}l@{}}{*Northwest Germanic*} \\
\mbox{PGmc Final Z Deletion} & \emph{*knéxta} \\
\end{tabular}
\end{minipage}
&
&
\begin{minipage}[t]{\linewidth}
\centering\textbf{Old English changes}\par
\vspace{0.35em}
\raggedright
\centering\textbf{}\par
\raggedright
\vspace{0.2em}
\begin{tabular}{@{}>{\raggedright\arraybackslash}p{0.70\linewidth}@{\hspace{0.55em}}>{\raggedright\arraybackslash}p{0.16\linewidth}@{\hspace{0.25em}}}
\multicolumn{2}{@{}l@{}}{*Old English*} \\
PWGmc Final Bare A Loss & \emph{*knéxt} \\
OE Breaking & \emph{*knéoxt} \\
OE Ws Palatal Umlaut & \emph{*knixt} \\
\end{tabular}
\end{minipage}
\\
\end{tabularx}
\end{minipage}%
} \endgroup

Old English form: \emph{cniht}

\subsection{Reconstruction and comparative
evidence}\label{reconstruction-and-comparative-evidence-83}

The comparative sources align on an \emph{e}-grade reconstruction for
this noun. Ringe and Taylor cite \Recon{kneht} `knight', and Orel gives
{\Recon{knextaz} `knight'} (\citeproc{ref-RingeTaylor2014}{Ringe and
Taylor 2014, 142}; \citeproc{ref-Orel2003}{Orel 2003, 256}).
Kluge-Seebold likewise points to \emph{*knehta-}
(\citeproc{ref-KlugeSeebold2011}{Kluge 2011, 506}). The derivational
input {\Recon{knéxtaz} `knight'} follows that comparative evidence.

A competing citation reconstruction {\Recon{kníxtaz} `knight'} remains
possible as a label for the word family, but it is not the
reconstruction followed here. The Old English development discussed
below is based on {\Recon{knéxtaz} `knight'}.

\subsection{Old English evidence}\label{old-english-evidence-83}

Old English dictionaries record the noun as {\emph{cniht}} `knight'
(\citeproc{ref-ClarkHall1960}{Clark Hall 1960, 63};
\citeproc{ref-BosworthToller1898}{Bosworth and Toller 1898, 71}).
Campbell cites plural {\emph{cneohtas}} `knights' among the broken
forms, showing the same vowel environment from another point in the
paradigm (\citeproc{ref-Campbell1959}{Campbell 1959, sect. 146}).

The target is therefore an ordinary attested Old English noun. No
reconstructed OE comparator is needed here.

\subsection{Development to Old
English}\label{development-to-old-english-83}

From {\Recon{knéxtaz} `knight'}, the relevant Old English changes
include breaking before the velar cluster and then the later reduction
that yields {\emph{cniht}} `knight'. Campbell later notes the early
West-Saxon alternation {\emph{cniht}} `knight' beside plural
{\emph{cneohtas}} `knights' (\citeproc{ref-Campbell1959}{Campbell 1959,
sect. 305}). Sievers-Brunner gives the same contrast as \emph{cniht
\ldots{} cneohtas} (\citeproc{ref-SieversBrunner1965}{Brunner 1965,
sect. 122}). With that corrected input, the derivation is
straightforward.

\subsection{Stem comparison}\label{stem-comparison-1}

The comparison below sets the relevant forms side by side. It separates
the handbook-supported \emph{e}-grade input from a competing citation
reconstruction.

{\def\LTcaptype{none} % do not increment counter
\begin{longtable}[]{@{}
  >{\raggedright\arraybackslash}p{(\linewidth - 8\tabcolsep) * \real{0.2000}}
  >{\raggedright\arraybackslash}p{(\linewidth - 8\tabcolsep) * \real{0.2000}}
  >{\raggedright\arraybackslash}p{(\linewidth - 8\tabcolsep) * \real{0.2000}}
  >{\raggedright\arraybackslash}p{(\linewidth - 8\tabcolsep) * \real{0.2000}}
  >{\raggedright\arraybackslash}p{(\linewidth - 8\tabcolsep) * \real{0.2000}}@{}}
\toprule\noalign{}
\begin{minipage}[b]{\linewidth}\raggedright
Formation / label
\end{minipage} & \begin{minipage}[b]{\linewidth}\raggedright
Candidate input
\end{minipage} & \begin{minipage}[b]{\linewidth}\raggedright
Expected or documented OE outcome
\end{minipage} & \begin{minipage}[b]{\linewidth}\raggedright
OE comparison form
\end{minipage} & \begin{minipage}[b]{\linewidth}\raggedright
Result
\end{minipage} \\
\midrule\noalign{}
\endhead
\bottomrule\noalign{}
\endlastfoot
competing citation reconstruction & {\emph{*kníxtaz}} & not the
reconstruction followed here & broader citation tradition & useful as a
competing label, but not the source-based choice used for the OE
derivation \\
handbook-supported reconstruction & {\emph{*knéxtaz}} & regular output:
{\emph{cniht}} `knight' & {\emph{cniht}} `knight' & exact match between
comparative reconstruction and attested OE noun \\
related plural evidence & same stem family & plural {\emph{cneohtas}}
`knights' type background & {\emph{cneohtas}} `knights' & supports the
vowel environment, but not the Old English form here cell \\
\end{longtable}
}

\section{\texorpdfstring{lade --- OE
\emph{hladan}}{lade --- OE hladan}}\label{lade-oe-hladan}

\index[iv]{02oe@\textbf{Old English}!hladan@\emph{hladan}}
\index[iv]{03pgmc@\textbf{Proto-Germanic}!lathojana@*laθōjaną}
\index[iv]{03pgmc@\textbf{Proto-Germanic}!xladana@*xláðaną}

Derivation: citation reconstruction \emph{*laθōjaną}; form followed here
\emph{*xláðaną} \textgreater{} \emph{hladan} (early analogy).

\subsection{Derivation trace}\label{derivation-trace-84}

Proto input: \emph{*xláðaną}

\begingroup
\setlength{\fboxsep}{6pt}

\noindent\fbox{%
\begin{minipage}{0.97\linewidth}
\small
\begin{tabularx}{\linewidth}{@{}>{\raggedright\arraybackslash}p{0.440\linewidth}>{\centering\arraybackslash}X>{\raggedright\arraybackslash}p{0.440\linewidth}@{}}
\begin{minipage}[t]{\linewidth}
\centering\textbf{Earlier Germanic changes}\par
\vspace{0.35em}
\raggedright
\centering\textbf{}\par
\raggedright
\vspace{0.2em}
\begin{tabular}{@{}>{\raggedright\arraybackslash}p{0.70\linewidth}@{\hspace{0.55em}}>{\raggedright\arraybackslash}p{0.20\linewidth}@{\hspace{0.25em}}}
\multicolumn{2}{@{}l@{}}{*Proto-West Germanic*} \\
PWGmc Dental Hardening & \emph{*xládaną} \\
\multicolumn{2}{@{}l@{}}{*Northwest Germanic*} \\
\multicolumn{2}{@{}l@{}}{[no change]} \\
\end{tabular}
\end{minipage}
&
&
\begin{minipage}[t]{\linewidth}
\centering\textbf{Old English changes}\par
\vspace{0.35em}
\raggedright
\centering\textbf{}\par
\raggedright
\vspace{0.2em}
\begin{tabular}{@{}>{\raggedright\arraybackslash}p{0.70\linewidth}@{\hspace{0.55em}}>{\raggedright\arraybackslash}p{0.20\linewidth}@{\hspace{0.25em}}}
\multicolumn{2}{@{}l@{}}{*Old English*} \\
Anglo Frisian Brightening & \emph{*xlædaną} \\
OE A Restoration & \emph{*xladaną} \\
OE Heavy Syllable Nasal Apocope & \emph{*xladan} \\
OE Secondary Nasalization & \emph{*xladąn} \\
OE Weak Tail Reduction & \emph{*xladan} \\
\end{tabular}
\end{minipage}
\\
\end{tabularx}
\end{minipage}%
} \endgroup

Old English form: \emph{hladan}

\subsection{Reconstruction and comparative
evidence}\label{reconstruction-and-comparative-evidence-84}

Ringe and Taylor cite the strong verb {\emph{hladan}} `lade, load'
directly (\citeproc{ref-RingeTaylor2014}{Ringe and Taylor 2014, 248}).
The citation label \Recon{laθōjaną} `lade' is used here only as a
broader family heading, not as the direct source of the OE derivation.

The derivational input therefore marks an early stem choice. The entry
follows the strong Verner-grade form that reaches Old English
\emph{hladan} directly.

\subsection{Old English evidence}\label{old-english-evidence-84}

Bosworth-Toller records the verb as \emph{hladan} and preserves the
expected strong-verb paradigm material around it
(\citeproc{ref-BosworthToller1898}{Bosworth and Toller 1898, 559}).
Clark Hall likewise records \emph{hladan}
(\citeproc{ref-ClarkHall1960}{Clark Hall 1960, 159}). The target is an
attested infinitive rather than a reconstructed paradigm cell.

For this entry the relevant comparison form is the infinitive
\emph{hladan} itself. The question is how that attested strong verb
relates to the broader comparative family.

\subsection{Development to Old
English}\label{development-to-old-english-84}

From \Recon{xláðaną} `lade', the verb passes through the expected early
voiced stop stage, Anglo-Frisian brightening, and the later
A-restoration that returns the root vowel to \emph{a} before the full
infinitival ending. Campbell treats verbs such as \emph{hladan} as
showing restored \emph{a} in the present system
(\citeproc{ref-Campbell1959}{Campbell 1959, sect. 744}). Ringe and
Taylor likewise derive \emph{hladan} from the voiced strong-verb line
(\citeproc{ref-RingeTaylor2014}{Ringe and Taylor 2014, 248}). The
resulting Old English infinitive is \emph{hladan}.

\subsection{Class comparison}\label{class-comparison-2}

The comparison below sets the relevant forms side by side. It separates
the wider weak-verb family label from the strong verb actually reflected
in Old English.

{\def\LTcaptype{none} % do not increment counter
\begin{longtable}[]{@{}
  >{\raggedright\arraybackslash}p{(\linewidth - 8\tabcolsep) * \real{0.2000}}
  >{\raggedright\arraybackslash}p{(\linewidth - 8\tabcolsep) * \real{0.2000}}
  >{\raggedright\arraybackslash}p{(\linewidth - 8\tabcolsep) * \real{0.2000}}
  >{\raggedright\arraybackslash}p{(\linewidth - 8\tabcolsep) * \real{0.2000}}
  >{\raggedright\arraybackslash}p{(\linewidth - 8\tabcolsep) * \real{0.2000}}@{}}
\toprule\noalign{}
\begin{minipage}[b]{\linewidth}\raggedright
Formation / class
\end{minipage} & \begin{minipage}[b]{\linewidth}\raggedright
Candidate input
\end{minipage} & \begin{minipage}[b]{\linewidth}\raggedright
Expected or documented OE outcome
\end{minipage} & \begin{minipage}[b]{\linewidth}\raggedright
OE comparison form
\end{minipage} & \begin{minipage}[b]{\linewidth}\raggedright
Result
\end{minipage} \\
\midrule\noalign{}
\endhead
\bottomrule\noalign{}
\endlastfoot
comparative weak-verb label & *laθōjaną & wider family background &
broader family context & useful family label, but not the direct source
of \emph{hladan} \\
selected strong Verner-grade input & *xláðaną & regular output:
\emph{hladan} & \emph{hladan} & exact match between formation and
attested OE infinitive \\
\end{longtable}
}

\section{\texorpdfstring{lap --- OE
\emph{lappa}}{lap --- OE lappa}}\label{lap-oe-lappa}

\index[iv]{02oe@\textbf{Old English}!lappa@\emph{lappa}}
\index[iv]{03pgmc@\textbf{Proto-Germanic}!labbaz@*lábbaz}
\index[iv]{03pgmc@\textbf{Proto-Germanic}!lappo@*láppô}

Derivation: citation reconstruction \emph{*lábbaz}; form followed here
\emph{*láppô} \textgreater{} \emph{lappa} (early analogy).

\subsection{Derivation trace}\label{derivation-trace-85}

Proto input: \emph{*láppô}

\begingroup
\setlength{\fboxsep}{6pt}

\noindent\fbox{%
\begin{minipage}{0.97\linewidth}
\small
\begin{tabularx}{\linewidth}{@{}>{\raggedright\arraybackslash}p{0.320\linewidth}>{\centering\arraybackslash}X>{\raggedright\arraybackslash}p{0.520\linewidth}@{}}
\begin{minipage}[t]{\linewidth}
\centering\textbf{Earlier Germanic changes}\par
\vspace{0.35em}
\raggedright
\centering\textbf{}\par
\raggedright
\vspace{0.2em}
\begin{tabular}{@{}>{\raggedright\arraybackslash}p{0.68\linewidth}@{\hspace{0.55em}}>{\raggedright\arraybackslash}p{0.16\linewidth}@{\hspace{0.25em}}}
\multicolumn{2}{@{}l@{}}{*Proto-West Germanic*} \\
\multicolumn{2}{@{}l@{}}{[no change]} \\
\multicolumn{2}{@{}l@{}}{*Northwest Germanic*} \\
\multicolumn{2}{@{}l@{}}{[no change]} \\
\end{tabular}
\end{minipage}
&
&
\begin{minipage}[t]{\linewidth}
\centering\textbf{Old English changes}\par
\vspace{0.35em}
\raggedright
\centering\textbf{}\par
\raggedright
\vspace{0.2em}
\begin{tabular}{@{}>{\raggedright\arraybackslash}p{0.74\linewidth}@{\hspace{0.55em}}>{\raggedright\arraybackslash}p{0.16\linewidth}@{\hspace{0.25em}}}
\multicolumn{2}{@{}l@{}}{*Old English*} \\
Anglo Frisian Brightening & \emph{*læppô} \\
OE A Restoration & \emph{*lappô} \\
OE Unstressed Long Vowel Shortening & \emph{*lappa} \\
\end{tabular}
\end{minipage}
\\
\end{tabularx}
\end{minipage}%
} \endgroup

Old English form: \emph{lappa}

\subsection{Reconstruction and comparative
evidence}\label{reconstruction-and-comparative-evidence-85}

The comparative sources point to a weak noun with \emph{pp}: Orel gives
\Recon{lappōn} `lap', and the Old English dictionary tradition preserves
variant \emph{læppa} `lap' (\citeproc{ref-Orel2003}{Orel 2003, 236};
\citeproc{ref-ClarkHall1960}{Clark Hall 1960, 180}). The derivational
input \Recon{láppô} `lap' follows that evidence.

A competing comparative label \Recon{lábbaz} `lap' has also circulated
for the word family, but the cited handbooks do not make it the direct
source of the Old English weak noun. The form relevant to the OE
development is the weak masculine input \Recon{láppô} `lap'.

\subsection{Old English evidence}\label{old-english-evidence-85}

Campbell cites \emph{lappa} `lap' as a case of restored \emph{a}
(\citeproc{ref-Campbell1959}{Campbell 1959, sect. 158}). Sievers-Brunner
records \emph{lappa} `lap' beside variant \emph{læppa} `lap'
(\citeproc{ref-SieversBrunner1965}{Brunner 1965, sect. 10}). The
dictionary tradition also preserves \emph{læppa} `lap'
(\citeproc{ref-ClarkHall1960}{Clark Hall 1960, 180};
\citeproc{ref-BosworthToller1898}{Bosworth and Toller 1898, 613}).

The target of this entry is the restored singular \emph{lappa} `lap'.
The variant \emph{læppa} `lap' and the oblique or plural \emph{leappan}
remain part of the Old English record and help frame the noun's vowel
history.

\subsection{Development to Old
English}\label{development-to-old-english-85}

Campbell explicitly lists \emph{lappa} `lap' among the forms with
restored \emph{a} (\citeproc{ref-Campbell1959}{Campbell 1959, sect.
158}). Sievers-Brunner records \emph{lappa} `lap' beside variant
\emph{læppa} `lap' at the same Old English stage
(\citeproc{ref-SieversBrunner1965}{Brunner 1965, sect. 10}). With the
weak masculine input chosen, the selected \emph{lappa} `lap' outcome is
therefore the regular Old English comparison.

\subsection{Stem comparison}\label{stem-comparison-2}

The comparison below sets the relevant forms side by side. It separates
the weak masculine formation from a competing voiced comparative label.

{\def\LTcaptype{none} % do not increment counter
\begin{longtable}[]{@{}
  >{\raggedright\arraybackslash}p{(\linewidth - 8\tabcolsep) * \real{0.2000}}
  >{\raggedright\arraybackslash}p{(\linewidth - 8\tabcolsep) * \real{0.2000}}
  >{\raggedright\arraybackslash}p{(\linewidth - 8\tabcolsep) * \real{0.2000}}
  >{\raggedright\arraybackslash}p{(\linewidth - 8\tabcolsep) * \real{0.2000}}
  >{\raggedright\arraybackslash}p{(\linewidth - 8\tabcolsep) * \real{0.2000}}@{}}
\toprule\noalign{}
\begin{minipage}[b]{\linewidth}\raggedright
Formation / label
\end{minipage} & \begin{minipage}[b]{\linewidth}\raggedright
Candidate input
\end{minipage} & \begin{minipage}[b]{\linewidth}\raggedright
Expected or documented OE outcome
\end{minipage} & \begin{minipage}[b]{\linewidth}\raggedright
OE comparison form
\end{minipage} & \begin{minipage}[b]{\linewidth}\raggedright
Result
\end{minipage} \\
\midrule\noalign{}
\endhead
\bottomrule\noalign{}
\endlastfoot
competing voiced comparative label & *lábbaz & not the form followed for
the OE weak-noun derivation & broader comparative background & useful as
a competing label, but not the source-based choice used here \\
selected weak masculine noun & *láppô & regular output: \emph{lappa} &
\emph{lappa} & exact match between formation and attested OE noun \\
attested OE variant line & same noun family & \emph{læppa},
\emph{leappan} & \emph{læppa} / \emph{leappan} & useful control forms
within the same OE tradition \\
\end{longtable}
}

\section{\texorpdfstring{laugh --- OE
\emph{hliehhan}}{laugh --- OE hliehhan}}\label{laugh-oe-hliehhan}

\index[iv]{02oe@\textbf{Old English}!hliehhan@\emph{hliehhan}}
\index[iv]{03pgmc@\textbf{Proto-Germanic}!lakana@*lákaną}
\index[iv]{03pgmc@\textbf{Proto-Germanic}!xlaxjana@*xláxjaną}

Derivation: citation reconstruction \emph{*lákaną}; form followed here
\emph{*xláxjaną} \textgreater{} \emph{hliehhan} (early analogy).

\subsection{Derivation trace}\label{derivation-trace-86}

Proto input: \emph{*xláxjaną}

\begingroup
\setlength{\fboxsep}{6pt}

\noindent\fbox{%
\begin{minipage}{0.97\linewidth}
\footnotesize
\begin{tabularx}{\linewidth}{@{}>{\raggedright\arraybackslash}p{0.470\linewidth}>{\centering\arraybackslash}X>{\raggedright\arraybackslash}p{0.470\linewidth}@{}}
\begin{minipage}[t]{\linewidth}
\centering\textbf{Earlier Germanic changes}\par
\vspace{0.35em}
\raggedright
\centering\textbf{}\par
\raggedright
\vspace{0.2em}
\begin{tabular}{@{}>{\raggedright\arraybackslash}p{0.64\linewidth}@{\hspace{0.55em}}>{\raggedright\arraybackslash}p{0.24\linewidth}@{\hspace{0.25em}}}
\multicolumn{2}{@{}l@{}}{*Proto-West Germanic*} \\
PWGmc J Gemination & \emph{*xláxxjaną} \\
\multicolumn{2}{@{}l@{}}{*Northwest Germanic*} \\
\multicolumn{2}{@{}l@{}}{[no change]} \\
\end{tabular}
\end{minipage}
&
&
\begin{minipage}[t]{\linewidth}
\centering\textbf{Old English changes}\par
\vspace{0.35em}
\raggedright
\centering\textbf{}\par
\raggedright
\vspace{0.2em}
\begin{tabular}{@{}>{\raggedright\arraybackslash}p{0.62\linewidth}@{\hspace{0.55em}}>{\raggedright\arraybackslash}p{0.28\linewidth}@{\hspace{0.25em}}}
\multicolumn{2}{@{}l@{}}{*Old English*} \\
Anglo Frisian Brightening & \emph{*xlæxxjaną} \\
OE Breaking & \emph{*xleaxxjaną} \\
OE Velar Fricative Palatalization & \emph{*xleaxçjaną} \\
OE Heavy Syllable Nasal Apocope & \emph{*xleaxçjan} \\
OE Secondary Nasalization & \emph{*xleaxçjąn} \\
OE I Umlaut & \emph{*xliexçjąn} \\
OE Weak Tail Reduction & \emph{*xliexçjan} \\
OE J Loss After Heavy & \emph{*xliexçan} \\
\end{tabular}
\end{minipage}
\\
\end{tabularx}
\end{minipage}%
} \endgroup

Old English form: \emph{hliehhan}

\subsection{Reconstruction and comparative
evidence}\label{reconstruction-and-comparative-evidence-86}

The wider Germanic family includes a non-j branch represented here by
the citation label \Recon{lákaną} `laugh', while the derivational input
\Recon{xláxjaną} `laugh' reflects the j-present line behind Old English
{\emph{hliehhan}} `laugh'.

This branch supplies the geminate fricative and the vowel development
characteristic of the Old English verb. The comparative family label and
the OE-facing input are therefore related but not identical.

\subsection{Old English evidence}\label{old-english-evidence-86}

Bosworth-Toller records {\emph{hlihhan}} `laugh' as the verb `to laugh'
(\citeproc{ref-BosworthToller1898}{Bosworth and Toller 1898, 551}).
Clark Hall cross-references \emph{hlæhan} `laugh', \emph{hlehhan}
`laugh', and \emph{hlihhan} `laugh' to {\emph{hliehhan}} `laugh'
(\citeproc{ref-ClarkHall1960}{Clark Hall 1960, 160--61}). Bright's
glossary likewise gives \emph{hlihhan (hliehhan, hlyhhan)}
(\citeproc{ref-BrightCassidyRingler1971}{Bright 1971, 315}). The target
of this entry is the West Saxon {\emph{hliehhan}} `laugh'.

The variants belong to the same background, but the attested lemma
{\emph{hliehhan}} `laugh' supplies the evidence followed here.

\subsection{Development to Old
English}\label{development-to-old-english-86}

From \Recon{xláxjaną} `laugh', West Germanic j-gemination yields the
doubled consonant (\citeproc{ref-Campbell1959}{Campbell 1959, sect.
407}). Ringe and Taylor derive Old English {\emph{hliehhan}} `laugh'
from the j-present branch via breaking before the palatalized geminate
(\citeproc{ref-RingeTaylor2014}{Ringe and Taylor 2014, 240}). They
separately compare the related noun \emph{hleahtor} as the outcome of
\Recon{hlahtraz} `laugh' (\citeproc{ref-RingeTaylor2014}{Ringe and
Taylor 2014, 328}).

\subsection{Branch comparison}\label{branch-comparison}

The comparison below sets the relevant forms side by side. It separates
the wider non-j family label from the j-present branch actually
reflected in Old English.

{\def\LTcaptype{none} % do not increment counter
\begin{longtable}[]{@{}
  >{\raggedright\arraybackslash}p{(\linewidth - 8\tabcolsep) * \real{0.2000}}
  >{\raggedright\arraybackslash}p{(\linewidth - 8\tabcolsep) * \real{0.2000}}
  >{\raggedright\arraybackslash}p{(\linewidth - 8\tabcolsep) * \real{0.2000}}
  >{\raggedright\arraybackslash}p{(\linewidth - 8\tabcolsep) * \real{0.2000}}
  >{\raggedright\arraybackslash}p{(\linewidth - 8\tabcolsep) * \real{0.2000}}@{}}
\toprule\noalign{}
\begin{minipage}[b]{\linewidth}\raggedright
Formation / branch
\end{minipage} & \begin{minipage}[b]{\linewidth}\raggedright
Candidate input
\end{minipage} & \begin{minipage}[b]{\linewidth}\raggedright
Expected or documented OE outcome
\end{minipage} & \begin{minipage}[b]{\linewidth}\raggedright
OE comparison form
\end{minipage} & \begin{minipage}[b]{\linewidth}\raggedright
Result
\end{minipage} \\
\midrule\noalign{}
\endhead
\bottomrule\noalign{}
\endlastfoot
wider non-j family & *lákaną & comparative background outside the
selected OE line & wider family context & useful family label, but not
the direct source of \emph{hliehhan} \\
selected j-present branch & *xláxjaną & regular output: \emph{hliehhan}
& \emph{hliehhan} & exact match between branch and attested OE lemma \\
attested OE variants & same OE verb line & \emph{hlæhhan},
\emph{hlehhan} & \emph{hlæhhan} / \emph{hlehhan} & genuine variant
evidence, but secondary to the form compared here \\
\end{longtable}
}

\section{\texorpdfstring{loam --- OE
\emph{lām}}{loam --- OE lām}}\label{loam-oe-lux101m}

\index[iv]{02oe@\textbf{Old English}!lam@\emph{lām}}
\index[iv]{03pgmc@\textbf{Proto-Germanic}!laima@*láimą}
\index[iv]{03pgmc@\textbf{Proto-Germanic}!laimon@*laimōn}

Derivation: citation reconstruction \emph{*laimōn}; form followed here
\emph{*láimą} \textgreater{} \emph{lām} (early analogy).

\subsection{Derivation trace}\label{derivation-trace-87}

Proto input: \emph{*láimą}

\begingroup
\setlength{\fboxsep}{6pt}

\noindent\fbox{%
\begin{minipage}{0.97\linewidth}
\small
\begin{tabularx}{\linewidth}{@{}>{\raggedright\arraybackslash}p{0.440\linewidth}>{\centering\arraybackslash}X>{\raggedright\arraybackslash}p{0.440\linewidth}@{}}
\begin{minipage}[t]{\linewidth}
\centering\textbf{Earlier Germanic changes}\par
\vspace{0.35em}
\raggedright
\centering\textbf{}\par
\raggedright
\vspace{0.2em}
\begin{tabular}{@{}>{\raggedright\arraybackslash}p{0.70\linewidth}@{\hspace{0.55em}}>{\raggedright\arraybackslash}p{0.16\linewidth}@{\hspace{0.25em}}}
\multicolumn{2}{@{}l@{}}{*Proto-West Germanic*} \\
PWGmc Ai Monophthongization & \emph{*lāmą} \\
\multicolumn{2}{@{}l@{}}{*Northwest Germanic*} \\
\multicolumn{2}{@{}l@{}}{[no change]} \\
\end{tabular}
\end{minipage}
&
&
\begin{minipage}[t]{\linewidth}
\centering\textbf{Old English changes}\par
\vspace{0.35em}
\raggedright
\centering\textbf{}\par
\raggedright
\vspace{0.2em}
\begin{tabular}{@{}>{\raggedright\arraybackslash}p{0.74\linewidth}@{\hspace{0.55em}}>{\raggedright\arraybackslash}p{0.16\linewidth}@{\hspace{0.25em}}}
\multicolumn{2}{@{}l@{}}{*Old English*} \\
OE Heavy Syllable Nasal Apocope & \emph{*lām} \\
\end{tabular}
\end{minipage}
\\
\end{tabularx}
\end{minipage}%
} \endgroup

Old English form: \emph{lām}

\subsection{Reconstruction and comparative
evidence}\label{reconstruction-and-comparative-evidence-87}

The inherited comparative noun is given as \Recon{laimōn} `loam' or
\emph{*laiman-}, and both Orel and Kroonen identify Old English
\emph{lām} `loam' as a neuter reflex of that family
(\citeproc{ref-Orel2003}{Orel 2003, 272};
\citeproc{ref-Kroonen2013}{Kroonen 2013, 363}). The form followed here,
\Recon{láimą} `loam', differs from the comparative headword because it
represents the stem class that matches the Old English noun most
directly.

This is therefore a class shift within the history of the English branch
rather than a dispute about the OE target itself.

\subsection{Old English evidence}\label{old-english-evidence-87}

Old English dictionaries record the noun as \emph{lām} `loam', a neuter
word for `loam, clay, mud' (\citeproc{ref-BosworthToller1898}{Bosworth
and Toller 1898, 604}; \citeproc{ref-ClarkHall1960}{Clark Hall 1960,
196}). The target is an attested citation form rather than a
reconstructed comparator.

The relevant question is not whether \emph{lām} `loam' is Old English,
but which inherited formation best accounts for that attested neuter
noun.

\subsection{Development to Old
English}\label{development-to-old-english-87}

From \Recon{láimą} `loam', regular monophthongization of \emph{ai} and
the later loss of the final nasal syllable yield \emph{lām} `loam'. With
that OE-facing input, the phonological development is straightforward.

\subsection{Class comparison}\label{class-comparison-3}

The comparison below sets the inherited comparative n-stem label beside
the OE-facing stem class used to derive the attested noun.

{\def\LTcaptype{none} % do not increment counter
\begin{longtable}[]{@{}
  >{\raggedright\arraybackslash}p{(\linewidth - 8\tabcolsep) * \real{0.2000}}
  >{\raggedright\arraybackslash}p{(\linewidth - 8\tabcolsep) * \real{0.2000}}
  >{\raggedright\arraybackslash}p{(\linewidth - 8\tabcolsep) * \real{0.2000}}
  >{\raggedright\arraybackslash}p{(\linewidth - 8\tabcolsep) * \real{0.2000}}
  >{\raggedright\arraybackslash}p{(\linewidth - 8\tabcolsep) * \real{0.2000}}@{}}
\toprule\noalign{}
\begin{minipage}[b]{\linewidth}\raggedright
Formation / class
\end{minipage} & \begin{minipage}[b]{\linewidth}\raggedright
Candidate input
\end{minipage} & \begin{minipage}[b]{\linewidth}\raggedright
Expected or documented OE outcome
\end{minipage} & \begin{minipage}[b]{\linewidth}\raggedright
OE comparison form
\end{minipage} & \begin{minipage}[b]{\linewidth}\raggedright
Result
\end{minipage} \\
\midrule\noalign{}
\endhead
\bottomrule\noalign{}
\endlastfoot
inherited comparative noun & *laimōn & comparative family background &
wider family context & useful headword, but not the direct OE-facing
input \\
OE-facing stem class followed here & *láimą & regular output: \emph{lām}
& \emph{lām} & exact match between the form followed here and the
attested OE noun \\
\end{longtable}
}

\section{\texorpdfstring{lung --- OE
\emph{lungen}}{lung --- OE lungen}}\label{lung-oe-lungen}

\index[iv]{02oe@\textbf{Old English}!lungen@\emph{lungen}}
\index[iv]{03pgmc@\textbf{Proto-Germanic}!lunganjo@*lúnganjō}
\index[iv]{03pgmc@\textbf{Proto-Germanic}!lungo@*lungō}

Derivation: citation reconstruction \emph{*lungō}; form followed here
\emph{*lúnganjō} \textgreater{} \emph{lungen} (early analogy).

\subsection{Derivation trace}\label{derivation-trace-88}

Proto input: \emph{*lúnganjō}

\begingroup
\setlength{\fboxsep}{6pt}

\noindent\fbox{%
\begin{minipage}{0.97\linewidth}
\small
\begin{tabularx}{\linewidth}{@{}>{\raggedright\arraybackslash}p{0.470\linewidth}>{\centering\arraybackslash}X>{\raggedright\arraybackslash}p{0.470\linewidth}@{}}
\begin{minipage}[t]{\linewidth}
\centering\textbf{Earlier Germanic changes}\par
\vspace{0.35em}
\raggedright
\centering\textbf{}\par
\raggedright
\vspace{0.2em}
\begin{tabular}{@{}>{\raggedright\arraybackslash}p{0.64\linewidth}@{\hspace{0.55em}}>{\raggedright\arraybackslash}p{0.26\linewidth}@{\hspace{0.25em}}}
\multicolumn{2}{@{}l@{}}{*Proto-West Germanic*} \\
PWGmc J Gemination & \emph{*lúngannjō} \\
\multicolumn{2}{@{}l@{}}{*Northwest Germanic*} \\
NWGmc Final Long O Raising & \emph{*lúngannju} \\
\end{tabular}
\end{minipage}
&
&
\begin{minipage}[t]{\linewidth}
\centering\textbf{Old English changes}\par
\vspace{0.35em}
\raggedright
\centering\textbf{}\par
\raggedright
\vspace{0.2em}
\begin{tabular}{@{}>{\raggedright\arraybackslash}p{0.64\linewidth}@{\hspace{0.55em}}>{\raggedright\arraybackslash}p{0.26\linewidth}@{\hspace{0.25em}}}
\multicolumn{2}{@{}l@{}}{*Old English*} \\
OE I Umlaut & \emph{*lúngennju} \\
OE High Vowel Apocope & \emph{*lúngennj} \\
OE J Loss After Heavy & \emph{*lúngenn} \\
OE Final Geminate Simplification & \emph{*lúngen} \\
\end{tabular}
\end{minipage}
\\
\end{tabularx}
\end{minipage}%
} \endgroup

Old English form: \emph{lungen}

\subsection{Reconstruction and comparative
evidence}\label{reconstruction-and-comparative-evidence-88}

Kroonen treats the basic noun as \emph{*lungōn-} and also cites an
OE-facing derivative \emph{*lungunjō-}, continued by Old English
{\emph{lungen}} `lung' and close West Germanic cognates
(\citeproc{ref-Kroonen2013}{Kroonen 2013, 384}). The derivational input
\Recon{lúnganjō} `lung' models that derived feminine formation rather
than the base noun. The notation differs slightly from Kroonen's
\emph{*lungunjō-}, but both point to the same derived feminine line.

The difference between the citation label and the derivational input is
therefore derivational. Old English {\emph{lungen}} `lung' is not a
direct reflex of the bare base noun \Recon{lungō} `lung'; it belongs to
an expanded feminine formation.

\subsection{Old English evidence}\label{old-english-evidence-88}

Old English dictionaries record the noun as {\emph{lungen}} `lung', with
inflected forms such as \emph{lungenne} and \emph{lungena}
(\citeproc{ref-BosworthToller1898}{Bosworth and Toller 1898, 634}).
Clark Hall also preserves a small family of compounds such as
\emph{lungenādl} `lung-disease', \emph{lungensealf}, and
\emph{lungenwyrt} (\citeproc{ref-ClarkHall1960}{Clark Hall 1960, 191}).

The target is an attested Old English lexeme with its own paradigm, not
a rescued inflectional cell.

\subsection{Development to Old
English}\label{development-to-old-english-88}

From the selected derived input, the expected derivational consonant and
vowel adjustments lead to {\emph{lungen}} `lung'. Once the expanded
feminine formation is chosen, the Old English outcome is regular.

\subsection{Formation comparison}\label{formation-comparison-3}

The comparison below sets the relevant forms side by side. It separates
the base noun from the derived feminine formation reflected in Old
English.

{\def\LTcaptype{none} % do not increment counter
\begin{longtable}[]{@{}
  >{\raggedright\arraybackslash}p{(\linewidth - 8\tabcolsep) * \real{0.2000}}
  >{\raggedright\arraybackslash}p{(\linewidth - 8\tabcolsep) * \real{0.2000}}
  >{\raggedright\arraybackslash}p{(\linewidth - 8\tabcolsep) * \real{0.2000}}
  >{\raggedright\arraybackslash}p{(\linewidth - 8\tabcolsep) * \real{0.2000}}
  >{\raggedright\arraybackslash}p{(\linewidth - 8\tabcolsep) * \real{0.2000}}@{}}
\toprule\noalign{}
\begin{minipage}[b]{\linewidth}\raggedright
Formation
\end{minipage} & \begin{minipage}[b]{\linewidth}\raggedright
Candidate input
\end{minipage} & \begin{minipage}[b]{\linewidth}\raggedright
Expected or documented OE outcome
\end{minipage} & \begin{minipage}[b]{\linewidth}\raggedright
OE comparison form
\end{minipage} & \begin{minipage}[b]{\linewidth}\raggedright
Result
\end{minipage} \\
\midrule\noalign{}
\endhead
\bottomrule\noalign{}
\endlastfoot
base noun & *lungō & base-noun outcome without the OE derivative suffix
& broader family context & useful headword, but not the direct source of
\emph{lungen} \\
derived OE-facing formation & *lúnganjō & regular output: \emph{lungen}
& \emph{lungen} & exact match between the derived formation and the
attested OE noun \\
Kroonen's cited derivative & \emph{*lungunjō-} & comparative support for
the same OE-facing formation & \emph{lungen} and cognate set & supports
the derived feminine formation, with notation differing from the
normalized input form used here \\
\end{longtable}
}

\section{\texorpdfstring{navel --- OE
\emph{nafola}}{navel --- OE nafola}}\label{navel-oe-nafola}

\index[iv]{02oe@\textbf{Old English}!nafola@\emph{nafola}}
\index[iv]{03pgmc@\textbf{Proto-Germanic}!nablo@*nablô}
\index[iv]{03pgmc@\textbf{Proto-Germanic}!nabulo@*nábulô}

Derivation: citation reconstruction \emph{*nablô}; form followed here
\emph{*nábulô} \textgreater{} \emph{nafola} (early analogy).

\subsection{Derivation trace}\label{derivation-trace-89}

Proto input: \emph{*nábulô}

\begingroup
\setlength{\fboxsep}{6pt}

\noindent\fbox{%
\begin{minipage}{0.97\linewidth}
\small
\begin{tabularx}{\linewidth}{@{}>{\raggedright\arraybackslash}p{0.320\linewidth}>{\centering\arraybackslash}X>{\raggedright\arraybackslash}p{0.520\linewidth}@{}}
\begin{minipage}[t]{\linewidth}
\centering\textbf{Earlier Germanic changes}\par
\vspace{0.35em}
\raggedright
\centering\textbf{}\par
\raggedright
\vspace{0.2em}
\begin{tabular}{@{}>{\raggedright\arraybackslash}p{0.68\linewidth}@{\hspace{0.55em}}>{\raggedright\arraybackslash}p{0.16\linewidth}@{\hspace{0.25em}}}
\multicolumn{2}{@{}l@{}}{*Proto-West Germanic*} \\
\multicolumn{2}{@{}l@{}}{[no change]} \\
\multicolumn{2}{@{}l@{}}{*Northwest Germanic*} \\
\multicolumn{2}{@{}l@{}}{[no change]} \\
\end{tabular}
\end{minipage}
&
&
\begin{minipage}[t]{\linewidth}
\centering\textbf{Old English changes}\par
\vspace{0.35em}
\raggedright
\centering\textbf{}\par
\raggedright
\vspace{0.2em}
\begin{tabular}{@{}>{\raggedright\arraybackslash}p{0.74\linewidth}@{\hspace{0.55em}}>{\raggedright\arraybackslash}p{0.16\linewidth}@{\hspace{0.25em}}}
\multicolumn{2}{@{}l@{}}{*Old English*} \\
OE Med Unstressed U Lowering & \emph{*nábolô} \\
Anglo Frisian Brightening & \emph{*næbolô} \\
OE A Restoration & \emph{*nabolô} \\
PGmc B Allophony & \emph{*naβolô} \\
OE Unstressed Long Vowel Shortening & \emph{*naβola} \\
\end{tabular}
\end{minipage}
\\
\end{tabularx}
\end{minipage}%
} \endgroup

Old English form: \emph{nafola}

\subsection{Reconstruction and comparative
evidence}\label{reconstruction-and-comparative-evidence-89}

Kroonen instead gives a nasal-suffix navel formation with Old English
{\emph{nafela}} `navel' among its reflexes
(\citeproc{ref-Kroonen2013}{Kroonen 2013, 420}), while Ringe and Taylor
give the derivational pathway \Recon{nabulō} `navel' \textgreater{}
\Recon{næbula} `navel' \textgreater{} \emph{nafola} `navel'
(\citeproc{ref-RingeTaylor2014}{Ringe and Taylor 2014, 270}). The
difference is one of stage and notation rather than of lexeme identity:
the derivational input {\Recon{nábulô} `navel'} is the pre-syncope form
needed for the Old English development.

For the Old English comparison, the crucial point is simply that the
pre-OE form still contains a medial vowel.

\subsection{Old English evidence}\label{old-english-evidence-89}

Ringe and Taylor note the early West Saxon shift {\emph{nafola}} `navel'
\textgreater{} {\emph{nafela}} `navel'
(\citeproc{ref-RingeTaylor2014}{Ringe and Taylor 2014, 336}). Campbell
likewise records \emph{nafela} `navel' beside Corpus {\emph{nabula}}
`navel' (\citeproc{ref-Campbell1959}{Campbell 1959, sect. 159}). The
target of this entry is the nominative singular {\emph{nafola}} `navel',
the form that matches the selected derivational pathway most directly.

{\emph{nafela}} `navel' is the better-known later West Saxon spelling,
while {\emph{nabula}} `navel' preserves a less reduced medial vowel.
These forms belong to the same lexical history, but this entry is
centered on {\emph{nafola}} `navel'.

\subsection{Development to Old
English}\label{development-to-old-english-89}

Ringe and Taylor give the pre-OE line \Recon{nabulō} `navel'
\textgreater{} \Recon{næbula} `navel' \textgreater{} OE \emph{nafola}
`navel' (\citeproc{ref-RingeTaylor2014}{Ringe and Taylor 2014, 270}).
The trace represents that same development with stress-marked notation
and explicit intermediate weakening. Intervocalic \emph{b} then surfaces
as \emph{f}, and final weak-tail shortening gives \emph{nafola} `navel'.

The medial vowel is still present when A-restoration applies. That is
why the selected pre-syncope input differs from the syncopated
comparative headword.

\subsection{Stage comparison}\label{stage-comparison}

The comparison below sets the relevant forms side by side. It separates
the comparative citation form from the pre-syncope input and from the
later OE spellings.

{\def\LTcaptype{none} % do not increment counter
\begin{longtable}[]{@{}
  >{\raggedright\arraybackslash}p{(\linewidth - 8\tabcolsep) * \real{0.2000}}
  >{\raggedright\arraybackslash}p{(\linewidth - 8\tabcolsep) * \real{0.2000}}
  >{\raggedright\arraybackslash}p{(\linewidth - 8\tabcolsep) * \real{0.2000}}
  >{\raggedright\arraybackslash}p{(\linewidth - 8\tabcolsep) * \real{0.2000}}
  >{\raggedright\arraybackslash}p{(\linewidth - 8\tabcolsep) * \real{0.2000}}@{}}
\toprule\noalign{}
\begin{minipage}[b]{\linewidth}\raggedright
Formation / stage
\end{minipage} & \begin{minipage}[b]{\linewidth}\raggedright
Candidate input
\end{minipage} & \begin{minipage}[b]{\linewidth}\raggedright
Expected or documented OE outcome
\end{minipage} & \begin{minipage}[b]{\linewidth}\raggedright
OE comparison form
\end{minipage} & \begin{minipage}[b]{\linewidth}\raggedright
Result
\end{minipage} \\
\midrule\noalign{}
\endhead
\bottomrule\noalign{}
\endlastfoot
syncopated comparative headword & {\emph{*nablô}} & reduced
\emph{næfla}-type outcome rather than {\emph{nafola}} `navel' & not the
Old English form here & useful citation form, but too reduced for the
pathway modeled here \\
selected pre-syncope input & {\emph{*nábulô}} & regular output:
{\emph{nafola}} `navel' & {\emph{nafola}} `navel' & exact match between
derivational input and target \\
later OE reduction stages & same lexical history & attested
{\emph{nafela}} `navel'; Corpus {\emph{nabula}} `navel' &
{\emph{nafela}} `navel' / {\emph{nabula}} `navel' & related OE
spellings, but not the chosen comparator \\
\end{longtable}
}

\section{\texorpdfstring{neck --- OE
\emph{hnecca}}{neck --- OE hnecca}}\label{neck-oe-hnecca}

\index[iv]{02oe@\textbf{Old English}!hnecca@\emph{hnecca}}
\index[iv]{03pgmc@\textbf{Proto-Germanic}!xnakkaz@*xnákkaz}
\index[iv]{03pgmc@\textbf{Proto-Germanic}!xnekko@*xnékkô}

Derivation: citation reconstruction \emph{*xnákkaz}; form followed here
\emph{*xnékkô} \textgreater{} \emph{hnecca} (early analogy).

\subsection{Derivation trace}\label{derivation-trace-90}

Proto input: \emph{*xnékkô}

\begingroup
\setlength{\fboxsep}{6pt}

\noindent\fbox{%
\begin{minipage}{0.97\linewidth}
\small
\begin{tabularx}{\linewidth}{@{}>{\raggedright\arraybackslash}p{0.440\linewidth}>{\centering\arraybackslash}X>{\raggedright\arraybackslash}p{0.440\linewidth}@{}}
\begin{minipage}[t]{\linewidth}
\centering\textbf{Earlier Germanic changes}\par
\vspace{0.35em}
\raggedright
\centering\textbf{}\par
\raggedright
\vspace{0.2em}
\begin{tabular}{@{}>{\raggedright\arraybackslash}p{0.68\linewidth}@{\hspace{0.55em}}>{\raggedright\arraybackslash}p{0.16\linewidth}@{\hspace{0.25em}}}
\multicolumn{2}{@{}l@{}}{*Proto-West Germanic*} \\
\multicolumn{2}{@{}l@{}}{[no change]} \\
\multicolumn{2}{@{}l@{}}{*Northwest Germanic*} \\
\multicolumn{2}{@{}l@{}}{[no change]} \\
\end{tabular}
\end{minipage}
&
&
\begin{minipage}[t]{\linewidth}
\centering\textbf{Old English changes}\par
\vspace{0.35em}
\raggedright
\centering\textbf{}\par
\raggedright
\vspace{0.2em}
\begin{tabular}{@{}>{\raggedright\arraybackslash}p{0.74\linewidth}@{\hspace{0.55em}}>{\raggedright\arraybackslash}p{0.16\linewidth}@{\hspace{0.25em}}}
\multicolumn{2}{@{}l@{}}{*Old English*} \\
OE Unstressed Long Vowel Shortening & \emph{*xnékka} \\
\end{tabular}
\end{minipage}
\\
\end{tabularx}
\end{minipage}%
} \endgroup

Old English form: \emph{hnecca}

\subsection{Reconstruction and comparative
evidence}\label{reconstruction-and-comparative-evidence-90}

The noun belongs to an ablauting n-stem family. Kroonen reconstructs a
paradigm with nominative \Recon{hnekkō} `neck', genitive \Recon{hnukkaz}
`neck', and accusative plural \Recon{hnakkuns} `neck', and he places Old
English {\emph{hnecca}} `neck, nape' among the e-grade descendants
(\citeproc{ref-Kroonen2011}{Kroonen 2011, 167}). Kluge-Seebold likewise
identifies \emph{ae. hnecca} as an ablaut partner of the a-grade
\emph{Nacken} family (\citeproc{ref-KlugeSeebold2011}{Kluge 2011, 347}).

A competing comparative label {\Recon{xnákkaz} `neck'} belongs to the
wider family, and Orel also gives an a-grade headword line
(\citeproc{ref-Orel2003}{Orel 2003, 218}). The derivational input
{\Recon{xnékkô} `neck'}, however, is the form that matches the Old
English branch.

\subsection{Old English evidence}\label{old-english-evidence-90}

Clark Hall records the weak masculine noun {\emph{hnecca}} `neck, nape'
(\citeproc{ref-ClarkHall1960}{Clark Hall 1960, 162}). Bosworth-Toller
likewise records {\emph{hnecca}} `neck, nape'
(\citeproc{ref-BosworthToller1898}{Bosworth and Toller 1898, 567}). The
target is therefore an attested citation form, not an oblique cell or a
reconstructed lemma.

The phonological question is upstream of the Old English evidence. The
attested noun already shows that the branch continued an e-grade form
rather than the a-grade seen in much of the continental material.

\subsection{Development to Old
English}\label{development-to-old-english-90}

From {\Recon{xnékkô} `neck'}, the derivation is straightforward. The
trace shortens the final long vowel to \Recon{xnékka} `neck', and Old
English orthography gives {\emph{hnecca}} `neck'.

The derivation depends on the earlier selection of the e-grade weak-noun
form continued by Old English.

\subsection{Stem comparison}\label{stem-comparison-3}

The comparison below sets the relevant forms side by side. It separates
the wider a-grade family from the selected e-grade Old English branch.

{\def\LTcaptype{none} % do not increment counter
\begin{longtable}[]{@{}
  >{\raggedright\arraybackslash}p{(\linewidth - 8\tabcolsep) * \real{0.2000}}
  >{\raggedright\arraybackslash}p{(\linewidth - 8\tabcolsep) * \real{0.2000}}
  >{\raggedright\arraybackslash}p{(\linewidth - 8\tabcolsep) * \real{0.2000}}
  >{\raggedright\arraybackslash}p{(\linewidth - 8\tabcolsep) * \real{0.2000}}
  >{\raggedright\arraybackslash}p{(\linewidth - 8\tabcolsep) * \real{0.2000}}@{}}
\toprule\noalign{}
\begin{minipage}[b]{\linewidth}\raggedright
Formation / label
\end{minipage} & \begin{minipage}[b]{\linewidth}\raggedright
Candidate input
\end{minipage} & \begin{minipage}[b]{\linewidth}\raggedright
Expected or documented OE outcome
\end{minipage} & \begin{minipage}[b]{\linewidth}\raggedright
OE comparison form
\end{minipage} & \begin{minipage}[b]{\linewidth}\raggedright
Result
\end{minipage} \\
\midrule\noalign{}
\endhead
\bottomrule\noalign{}
\endlastfoot
competing comparative label & {\emph{*xnákkaz}} & broader a-grade family
rather than the selected OE source & continental {\emph{Nacken}} line &
useful family label, but not the input followed for the Old English
derivation \\
weak noun with a-grade & {\emph{*xnakkô}} & expected {\emph{hnacca}}
type outcome & {\emph{hnacca}} & fixes the class, but not the vowel
grade \\
selected e-grade nominative & {\emph{*xnékkô}} & regular output:
{\emph{hnecca}} & {\emph{hnecca}} & exact match between derivational
input and attested OE noun \\
oblique paradigm background & {\emph{*hnukkaz}}, {\emph{*hnakkuns}} &
ON/OHG/German a-grade continuation & {\emph{hnakki}} / {\emph{Nacken}} &
shows the wider ablaut family, but not the chosen OE branch \\
\end{longtable}
}

\section{\texorpdfstring{needle --- OE
\emph{nǣdl}}{needle --- OE nǣdl}}\label{needle-oe-nux1e3dl}

\index[iv]{02oe@\textbf{Old English}!naedl@\emph{nǣdl}}
\index[iv]{03pgmc@\textbf{Proto-Germanic}!nedlo@*nḗðlō}
\index[iv]{03pgmc@\textbf{Proto-Germanic}!nethlo@*nḗθlō}

Derivation: citation reconstruction \emph{*nḗθlō}; form followed here
\emph{*nḗðlō} \textgreater{} \emph{nǣdl} (early analogy).

\subsection{Derivation trace}\label{derivation-trace-91}

Proto input: \emph{*nḗðlō}

\begingroup
\setlength{\fboxsep}{6pt}

\noindent\fbox{%
\begin{minipage}{0.97\linewidth}
\small
\begin{tabularx}{\linewidth}{@{}>{\raggedright\arraybackslash}p{0.460\linewidth}>{\centering\arraybackslash}X>{\raggedright\arraybackslash}p{0.460\linewidth}@{}}
\begin{minipage}[t]{\linewidth}
\centering\textbf{Earlier Germanic changes}\par
\vspace{0.35em}
\raggedright
\centering\textbf{}\par
\raggedright
\vspace{0.2em}
\begin{tabular}{@{}>{\raggedright\arraybackslash}p{0.74\linewidth}@{\hspace{0.55em}}>{\raggedright\arraybackslash}p{0.16\linewidth}@{\hspace{0.25em}}}
\multicolumn{2}{@{}l@{}}{*Proto-West Germanic*} \\
PWGmc Dental Hardening & \emph{*nḗdlō} \\
\multicolumn{2}{@{}l@{}}{*Northwest Germanic*} \\
\mbox{NWGmc Final Long O Raising} & \emph{*nḗdlu} \\
\mbox{NWGmc Long E Lowering} & \emph{*nǣdlu} \\
\end{tabular}
\end{minipage}
&
&
\begin{minipage}[t]{\linewidth}
\centering\textbf{Old English changes}\par
\vspace{0.35em}
\raggedright
\centering\textbf{}\par
\raggedright
\vspace{0.2em}
\begin{tabular}{@{}>{\raggedright\arraybackslash}p{0.74\linewidth}@{\hspace{0.55em}}>{\raggedright\arraybackslash}p{0.16\linewidth}@{\hspace{0.25em}}}
\multicolumn{2}{@{}l@{}}{*Old English*} \\
\mbox{OE High Vowel Apocope} & \emph{*nǣdl} \\
\end{tabular}
\end{minipage}
\\
\end{tabularx}
\end{minipage}%
} \endgroup

Old English form: \emph{nǣdl}

\subsection{Reconstruction and comparative
evidence}\label{reconstruction-and-comparative-evidence-91}

Ringe and Taylor treat the word as a voiced/voiceless alternant, citing
\Recon{nēþlō} `needle' \textasciitilde{} \Recon{nēdlō} `needle'
\textgreater{} OE \emph{nédl} `needle'
(\citeproc{ref-RingeTaylor2014}{Ringe and Taylor 2014, 329}). The form
followed here, {\Recon{nḗðlō} `needle'}, is the voiced Verner-grade form
used for the Old English comparison, while the citation form
{\Recon{nḗθlō} `needle'} remains the broader lexeme label.

The development discussed here follows the Ringe-Taylor alternant
framework.

\subsection{Old English evidence}\label{old-english-evidence-91}

Clark Hall records the attested citation form \emph{nǣdl} `needle'
(\citeproc{ref-ClarkHall1960}{Clark Hall 1960, 210}). Campbell lists
\emph{nédl} `needle' among the expected unbroken forms after \emph{t}
and \emph{d} (\citeproc{ref-Campbell1959}{Campbell 1959, sect. 367}).
Hogg also includes \emph{nidi} / \emph{nǣdl} `needle' in the same
broader cluster history (\citeproc{ref-Hogg1992}{Hogg 1992, 95}).

The target is therefore an attested citation form. No oblique-cell
substitution is involved in this entry.

\subsection{Development to Old
English}\label{development-to-old-english-91}

Ringe and Taylor give the historical line \Recon{nēþlō} `needle'
\textasciitilde{} \Recon{nēdlō} `needle' \textgreater{} OE \emph{nédl}
`needle' (\citeproc{ref-RingeTaylor2014}{Ringe and Taylor 2014, 329}).
Campbell likewise lists \emph{nédl} `needle' among the expected unbroken
forms after \emph{t} and \emph{d} (\citeproc{ref-Campbell1959}{Campbell
1959, sect. 367}). The trace expresses the same pathway with the voiced
alternant followed here for the Old English comparison.

The essential choice lies in which Proto-Germanic alternant is taken as
the starting point. Once the voiced form is chosen, the rest of the
pathway is regular.

\subsection{Alternant comparison}\label{alternant-comparison}

The comparison below sets the relevant forms side by side. It separates
the broader citation headword from the voiced alternant used for Old
English.

{\def\LTcaptype{none} % do not increment counter
\begin{longtable}[]{@{}
  >{\raggedright\arraybackslash}p{(\linewidth - 8\tabcolsep) * \real{0.2000}}
  >{\raggedright\arraybackslash}p{(\linewidth - 8\tabcolsep) * \real{0.2000}}
  >{\raggedright\arraybackslash}p{(\linewidth - 8\tabcolsep) * \real{0.2000}}
  >{\raggedright\arraybackslash}p{(\linewidth - 8\tabcolsep) * \real{0.2000}}
  >{\raggedright\arraybackslash}p{(\linewidth - 8\tabcolsep) * \real{0.2000}}@{}}
\toprule\noalign{}
\begin{minipage}[b]{\linewidth}\raggedright
Formation / stage
\end{minipage} & \begin{minipage}[b]{\linewidth}\raggedright
Candidate input
\end{minipage} & \begin{minipage}[b]{\linewidth}\raggedright
Expected or documented OE outcome
\end{minipage} & \begin{minipage}[b]{\linewidth}\raggedright
OE comparison form
\end{minipage} & \begin{minipage}[b]{\linewidth}\raggedright
Result
\end{minipage} \\
\midrule\noalign{}
\endhead
\bottomrule\noalign{}
\endlastfoot
comparative voiceless headword & {\emph{*nḗθlō}} & broader word-family
label rather than the OE-facing alternant & \emph{*nēþlō} line & useful
citation form, but not the derivational input for the Old English
comparison \\
voiced Verner alternant followed here & {\emph{*nḗðlō}} & regular
output: {\emph{nǣdl}} & {\emph{nǣdl}} & exact match between the form
followed here and the attested OE noun \\
later hardening stage & *nḗdlō & intermediate pre-OE stage in the same
derivation & \emph{nǣdl} & genuine stage in the pathway, but not the
Proto-Germanic form followed here \\
\end{longtable}
}

\section{\texorpdfstring{nose --- OE
\emph{nosu}}{nose --- OE nosu}}\label{nose-oe-nosu}

\index[iv]{02oe@\textbf{Old English}!nosu@\emph{nosu}}
\index[iv]{03pgmc@\textbf{Proto-Germanic}!naso@*nasō}
\index[iv]{03pgmc@\textbf{Proto-Germanic}!nuso@*núsō}

Derivation: citation reconstruction \emph{*nasō}; form followed here
\emph{*núsō} \textgreater{} \emph{nosu} (early analogy).

\subsection{Derivation trace}\label{derivation-trace-92}

Proto input: \emph{*núsō}

\begingroup
\setlength{\fboxsep}{6pt}

\noindent\fbox{%
\begin{minipage}{0.97\linewidth}
\small
\begin{tabularx}{\linewidth}{@{}>{\raggedright\arraybackslash}p{0.540\linewidth}>{\centering\arraybackslash}X>{\raggedright\arraybackslash}p{0.300\linewidth}@{}}
\begin{minipage}[t]{\linewidth}
\centering\textbf{Earlier Germanic changes}\par
\vspace{0.35em}
\raggedright
\centering\textbf{}\par
\raggedright
\vspace{0.2em}
\begin{tabular}{@{}>{\raggedright\arraybackslash}p{0.74\linewidth}@{\hspace{0.55em}}>{\raggedright\arraybackslash}p{0.16\linewidth}@{\hspace{0.25em}}}
\multicolumn{2}{@{}l@{}}{*Proto-West Germanic*} \\
\multicolumn{2}{@{}l@{}}{[no change]} \\
\multicolumn{2}{@{}l@{}}{*Northwest Germanic*} \\
\mbox{NWGmc U Lowering} & \emph{*nósō} \\
\mbox{NWGmc Final Long O Raising} & \emph{*nósu} \\
\end{tabular}
\end{minipage}
&
&
\begin{minipage}[t]{\linewidth}
\centering\textbf{Old English changes}\par
\vspace{0.35em}
\raggedright
\centering\textbf{}\par
\raggedright
\vspace{0.2em}
\begin{tabular}{@{}>{\raggedright\arraybackslash}p{0.64\linewidth}@{\hspace{0.55em}}>{\raggedright\arraybackslash}p{0.16\linewidth}@{\hspace{0.25em}}}
\multicolumn{2}{@{}l@{}}{*Old English*} \\
\multicolumn{2}{@{}l@{}}{[no change]} \\
\end{tabular}
\end{minipage}
\\
\end{tabularx}
\end{minipage}%
} \endgroup

Old English form: \emph{nosu}

\subsection{Reconstruction and comparative
evidence}\label{reconstruction-and-comparative-evidence-92}

Kroonen reconstructs a Germanic ablaut pair \emph{*nasō-}
\textasciitilde{} \emph{*nusō-} and adds that the root \emph{*nus-} is
likely to have arisen as a secondary zero grade after a remodeling of
the older paradigm (\citeproc{ref-Kroonen2013}{Kroonen 2013, 423}).
Campbell is more specific for Old English, citing \emph{nosu} `nose'
\textless{} \Recon{nusō} `nose' (\citeproc{ref-Campbell1959}{Campbell
1959, 44}).

The citation reconstruction \Recon{nasō} `nose' is therefore best
treated as the full-grade comparative headword, while the derivational
input \Recon{núsō} `nose' represents the remodeled zero-grade line
continued by the Old English form discussed here. Orel's \Recon{nasō}
`nose' alongside OE \emph{nasu} `nose' preserves the competing
full-grade notation and shows that the two lines should not be collapsed
without comment (\citeproc{ref-Orel2003}{Orel 2003, 320}).

\subsection{Old English evidence}\label{old-english-evidence-92}

\emph{Nosu} `nose' is an attested Old English noun. Ringe and Taylor
list it among the few surviving early Old English feminine u-stems
(\citeproc{ref-RingeTaylor2014}{Ringe and Taylor 2014, 385}). Clark Hall
likewise gives \emph{nosu f.} `nose', with genitive-dative singular
\emph{nosa}, and cross-refers {\emph{nasu}} `nose' to \emph{nosu} `nose'
(\citeproc{ref-ClarkHall1960}{Clark Hall 1960, 810}).

The selected OE target is therefore attested \emph{nosu} `nose', not a
reconstructed placeholder. The lexicographical record also gives
{\emph{nasu}} `nose' for the full-grade side of the tradition.

\subsection{Development to Old
English}\label{development-to-old-english-92}

From \Recon{núsō} `nose', the regular path is the one documented by the
current trace: \Recon{núsō} `nose' \textgreater{} \Recon{nósō} `nose'
\textgreater{} \Recon{nósu} `nose' \textgreater{} \emph{nosu} `nose'.
The early special step lies in the choice of the zero-grade input, not
in any late Old English repair.

With that input chosen, the OE development is straightforward. The
full-grade line behind \Recon{nasō} `nose' instead points toward
{\emph{nasu}} `nose', not to the form treated here.

\subsection{Stem comparison}\label{stem-comparison-4}

The comparison below sets the relevant forms side by side. It separates
the full-grade comparative line from the remodeled zero-grade input that
yields the Old English form.

{\def\LTcaptype{none} % do not increment counter
\begin{longtable}[]{@{}
  >{\raggedright\arraybackslash}p{(\linewidth - 8\tabcolsep) * \real{0.2000}}
  >{\raggedright\arraybackslash}p{(\linewidth - 8\tabcolsep) * \real{0.2000}}
  >{\raggedright\arraybackslash}p{(\linewidth - 8\tabcolsep) * \real{0.2000}}
  >{\raggedright\arraybackslash}p{(\linewidth - 8\tabcolsep) * \real{0.2000}}
  >{\raggedright\arraybackslash}p{(\linewidth - 8\tabcolsep) * \real{0.2000}}@{}}
\toprule\noalign{}
\begin{minipage}[b]{\linewidth}\raggedright
Formation / label
\end{minipage} & \begin{minipage}[b]{\linewidth}\raggedright
Candidate input
\end{minipage} & \begin{minipage}[b]{\linewidth}\raggedright
Expected or documented OE outcome
\end{minipage} & \begin{minipage}[b]{\linewidth}\raggedright
OE comparison form
\end{minipage} & \begin{minipage}[b]{\linewidth}\raggedright
Result
\end{minipage} \\
\midrule\noalign{}
\endhead
\bottomrule\noalign{}
\endlastfoot
full-grade comparative line & *nasō & expected full-grade continuation
\emph{nasu} & \emph{nasu} & useful comparative background, but not the
Old English-facing input \\
remodeled zero-grade line & *núsō & regular output: \emph{nosu} &
\emph{nosu} & exact match between derivational input and attested OE
noun \\
\end{longtable}
}

\section{\texorpdfstring{sap --- OE
\emph{sæp}}{sap --- OE sæp}}\label{sap-oe-suxe6p}

\index[iv]{02oe@\textbf{Old English}!saep@\emph{sæp}}
\index[iv]{03pgmc@\textbf{Proto-Germanic}!sapa@*sápą}
\index[iv]{03pgmc@\textbf{Proto-Germanic}!sapon@*sapōn}

Derivation: citation reconstruction \emph{*sapōn}; form followed here
\emph{*sápą} \textgreater{} \emph{sæp} (early analogy).

\subsection{Derivation trace}\label{derivation-trace-93}

Proto input: \emph{*sápą}

\begingroup
\setlength{\fboxsep}{6pt}

\noindent\fbox{%
\begin{minipage}{0.97\linewidth}
\small
\begin{tabularx}{\linewidth}{@{}>{\raggedright\arraybackslash}p{0.440\linewidth}>{\centering\arraybackslash}X>{\raggedright\arraybackslash}p{0.440\linewidth}@{}}
\begin{minipage}[t]{\linewidth}
\centering\textbf{Earlier Germanic changes}\par
\vspace{0.35em}
\raggedright
\centering\textbf{}\par
\raggedright
\vspace{0.2em}
\begin{tabular}{@{}>{\raggedright\arraybackslash}p{0.68\linewidth}@{\hspace{0.55em}}>{\raggedright\arraybackslash}p{0.16\linewidth}@{\hspace{0.25em}}}
\multicolumn{2}{@{}l@{}}{*Proto-West Germanic*} \\
\multicolumn{2}{@{}l@{}}{[no change]} \\
\multicolumn{2}{@{}l@{}}{*Northwest Germanic*} \\
\multicolumn{2}{@{}l@{}}{[no change]} \\
\end{tabular}
\end{minipage}
&
&
\begin{minipage}[t]{\linewidth}
\centering\textbf{Old English changes}\par
\vspace{0.35em}
\raggedright
\centering\textbf{}\par
\raggedright
\vspace{0.2em}
\begin{tabular}{@{}>{\raggedright\arraybackslash}p{0.74\linewidth}@{\hspace{0.55em}}>{\raggedright\arraybackslash}p{0.16\linewidth}@{\hspace{0.25em}}}
\multicolumn{2}{@{}l@{}}{*Old English*} \\
Anglo Frisian Brightening & \emph{*sæpą} \\
OE Heavy Syllable Nasal Apocope & \emph{*sæp} \\
\end{tabular}
\end{minipage}
\\
\end{tabularx}
\end{minipage}%
} \endgroup

Old English form: \emph{sæp}

\subsection{Reconstruction and comparative
evidence}\label{reconstruction-and-comparative-evidence-93}

The comparative sources do not give one uniform inherited stem. Kroonen
preserves the word family as \Recon{saf} `sap' / \emph{*ppan-}, with Old
English \emph{sæp} `sap' (\citeproc{ref-Kroonen2013}{Kroonen 2013,
420}). Orel preserves the comparative notation \Recon{sapōn} `sap' and
\Recon{sapan} `sap' (\citeproc{ref-Orel2003}{Orel 2003, 319}).

The derivational input {\Recon{sápą} `sap'} therefore does not replace
those comparative labels. It identifies the OE-facing stem shape that
yields the attested noun treated here.

\subsection{Old English evidence}\label{old-english-evidence-93}

Clark Hall records {\emph{sæp}} `sap' (e) n.
(\citeproc{ref-ClarkHall1960}{Clark Hall 1960, 247}). The target is
therefore an attested neuter Old English noun. Orel's plain {\emph{sap}}
`sap, juice' notation belongs to comparative normalization, not to the
spelling adopted here for the Old English form
(\citeproc{ref-Orel2003}{Orel 2003, 319}).

\subsection{Development to Old
English}\label{development-to-old-english-93}

From {\Recon{sápą} `sap'}, Anglo-Frisian brightening yields \emph{sæ}
`sap', and heavy-syllable nasal apocope then produces {\emph{sæp}}
`sap'. That is the regular path documented by the current trace.

The competing comparative lines do not give the same result. The
inherited n-stem notation {\Recon{sapōn} `sap'} yields {\emph{sape}}
`sap', while an i-stem continuation from the {\emph{*sapi-}} line leads
to {\emph{sep}} `sap' / {\emph{sepe}} `sap' rather than to {\emph{sæp}}
`sap'. The special step in this entry is therefore the early stem
choice, not a late OE paradigm-cell selection.

\subsection{Stem comparison}\label{stem-comparison-5}

The comparison below sets the relevant forms side by side. It separates
the competing comparative stem lines from the Old English-facing input.

{\def\LTcaptype{none} % do not increment counter
\begin{longtable}[]{@{}
  >{\raggedright\arraybackslash}p{(\linewidth - 8\tabcolsep) * \real{0.2000}}
  >{\raggedright\arraybackslash}p{(\linewidth - 8\tabcolsep) * \real{0.2000}}
  >{\raggedright\arraybackslash}p{(\linewidth - 8\tabcolsep) * \real{0.2000}}
  >{\raggedright\arraybackslash}p{(\linewidth - 8\tabcolsep) * \real{0.2000}}
  >{\raggedright\arraybackslash}p{(\linewidth - 8\tabcolsep) * \real{0.2000}}@{}}
\toprule\noalign{}
\begin{minipage}[b]{\linewidth}\raggedright
Formation / label
\end{minipage} & \begin{minipage}[b]{\linewidth}\raggedright
Candidate input
\end{minipage} & \begin{minipage}[b]{\linewidth}\raggedright
Expected or documented OE outcome
\end{minipage} & \begin{minipage}[b]{\linewidth}\raggedright
OE comparison form
\end{minipage} & \begin{minipage}[b]{\linewidth}\raggedright
Result
\end{minipage} \\
\midrule\noalign{}
\endhead
\bottomrule\noalign{}
\endlastfoot
comparative n-stem line & {\emph{*sapōn}} & local comparator output:
{\emph{sape}} `sap' & {\emph{sape}} `sap' & useful comparative
background, but not the source of attested \emph{sæp} \\
inferred i-stem comparator from \emph{*sapi-} & {\emph{*sapiz}} & local
comparator output: {\emph{sepe}} `sap' & {\emph{sepe}} `sap' & confirms
that an i-triggering stem does not reach the target \\
selected a-stem input & {\emph{*sápą}} & regular output: {\emph{sæp}}
`sap' & {\emph{sæp}} `sap' & exact match between derivational input and
attested OE noun \\
\end{longtable}
}

\section{\texorpdfstring{sea --- OE
\emph{sǣ}}{sea --- OE sǣ}}\label{sea-oe-sux1e3}

\index[iv]{02oe@\textbf{Old English}!sae@\emph{sǣ}}
\index[iv]{03pgmc@\textbf{Proto-Germanic}!sai@*sái}
\index[iv]{03pgmc@\textbf{Proto-Germanic}!saiwiz@*sáiwiz}

Derivation: citation reconstruction \emph{*sái}; form followed here
\emph{*sáiwiz} \textgreater{} \emph{sǣ} (early analogy).

\subsection{Derivation trace}\label{derivation-trace-94}

Proto input: \emph{*sáiwiz}

\begingroup
\setlength{\fboxsep}{6pt}

\noindent\fbox{%
\begin{minipage}{0.97\linewidth}
\small
\begin{tabularx}{\linewidth}{@{}>{\raggedright\arraybackslash}p{0.460\linewidth}>{\centering\arraybackslash}X>{\raggedright\arraybackslash}p{0.460\linewidth}@{}}
\begin{minipage}[t]{\linewidth}
\centering\textbf{Earlier Germanic changes}\par
\vspace{0.35em}
\raggedright
\centering\textbf{}\par
\raggedright
\vspace{0.2em}
\begin{tabular}{@{}>{\raggedright\arraybackslash}p{0.74\linewidth}@{\hspace{0.55em}}>{\raggedright\arraybackslash}p{0.16\linewidth}@{\hspace{0.25em}}}
\multicolumn{2}{@{}l@{}}{*Proto-West Germanic*} \\
PWGmc Ai Monophthongization & \emph{*sāwiz} \\
\multicolumn{2}{@{}l@{}}{*Northwest Germanic*} \\
\mbox{PGmc Final Z Deletion} & \emph{*sāwi} \\
\end{tabular}
\end{minipage}
&
&
\begin{minipage}[t]{\linewidth}
\centering\textbf{Old English changes}\par
\vspace{0.35em}
\raggedright
\centering\textbf{}\par
\raggedright
\vspace{0.2em}
\begin{tabular}{@{}>{\raggedright\arraybackslash}p{0.74\linewidth}@{\hspace{0.55em}}>{\raggedright\arraybackslash}p{0.16\linewidth}@{\hspace{0.25em}}}
\multicolumn{2}{@{}l@{}}{*Old English*} \\
OE W Loss Before I & \emph{*sāi} \\
OE I Umlaut & \emph{*sǣi} \\
\mbox{OE High Vowel Apocope} & \emph{*sǣ} \\
\end{tabular}
\end{minipage}
\\
\end{tabularx}
\end{minipage}%
} \endgroup

Old English form: \emph{sǣ}

\subsection{Reconstruction and comparative
evidence}\label{reconstruction-and-comparative-evidence-94}

Kroonen gives the noun in stem notation as \emph{*saiwi-}, an i-stem
whose English reflex is cited as OE {\emph{sæ}} `sea'
(\citeproc{ref-Kroonen2013}{Kroonen 2013, 423}). Ringe and Taylor write
the fuller form {\Recon{saiwiz} `sea'} and derive it through
\Recon{sawi} `sea' \textgreater{} \Recon{sei} `sea' \textgreater{} OE
\emph{sǣ} `sea' (\citeproc{ref-RingeTaylor2014}{Ringe and Taylor 2014,
sect. 6.7.1}). The comparative headword is therefore shorter than the
form required for the English history: {\Recon{sái} `sea'} names the
lexeme, but {\Recon{sáiwiz} `sea'} preserves the medial \emph{*w} and
the final high vowel that control the later development.

\subsection{Old English evidence}\label{old-english-evidence-94}

The Old English noun is the ordinary word for `sea'. Kroonen cites it as
{\emph{sæ}} `sea'; the normalized form here is {\emph{sǣ}} `sea'
(\citeproc{ref-Kroonen2013}{Kroonen 2013, 423}). Campbell likewise
treats \emph{sea} as continuing the same \Recon{saiui} `sea'
\textgreater{} \Recon{sǣi} `sea' history, with loss of \emph{u}/\emph{w}
before \emph{i} (\citeproc{ref-Campbell1959}{Campbell 1959, sect. 406}).

\subsection{Development to Old
English}\label{development-to-old-english-94}

Once the fuller i-stem input is chosen, the development is regular.
After Proto-West-Germanic monophthongization \Recon{sáiwiz} `sea'
\textgreater{} \Recon{sāwiz} `sea' and final \emph{*-z} loss
\Recon{sāwiz} `sea' \textgreater{} \Recon{sāwi} `sea', the non-initial
\emph{*w} disappears before unstressed \emph{*i}, and the following high
vowel fronts the root vowel before final apocope. The documented chain
is \Recon{sáiwiz} `sea' \textgreater{} \Recon{sāwiz} `sea'
\textgreater{} \Recon{sāwi} `sea' \textgreater{} \Recon{sāi} `sea'
\textgreater{} \Recon{sǣi} `sea' \textgreater{} \emph{sǣ} `sea'
(\citeproc{ref-RingeTaylor2014}{Ringe and Taylor 2014, sect. 6.7.1};
\citeproc{ref-Campbell1959}{Campbell 1959, sect. 406}).

\subsection{Stem and stage comparison}\label{stem-and-stage-comparison}

The comparison below sets the relevant forms side by side. It separates
the abbreviated comparative headword from the fuller i-stem input that
yields the Old English form.

{\def\LTcaptype{none} % do not increment counter
\begin{longtable}[]{@{}
  >{\raggedright\arraybackslash}p{(\linewidth - 8\tabcolsep) * \real{0.2000}}
  >{\raggedright\arraybackslash}p{(\linewidth - 8\tabcolsep) * \real{0.2000}}
  >{\raggedright\arraybackslash}p{(\linewidth - 8\tabcolsep) * \real{0.2000}}
  >{\raggedright\arraybackslash}p{(\linewidth - 8\tabcolsep) * \real{0.2000}}
  >{\raggedright\arraybackslash}p{(\linewidth - 8\tabcolsep) * \real{0.2000}}@{}}
\toprule\noalign{}
\begin{minipage}[b]{\linewidth}\raggedright
Formation / label
\end{minipage} & \begin{minipage}[b]{\linewidth}\raggedright
Candidate input
\end{minipage} & \begin{minipage}[b]{\linewidth}\raggedright
OE output or comparison
\end{minipage} & \begin{minipage}[b]{\linewidth}\raggedright
OE comparison form
\end{minipage} & \begin{minipage}[b]{\linewidth}\raggedright
Result
\end{minipage} \\
\midrule\noalign{}
\endhead
\bottomrule\noalign{}
\endlastfoot
abbreviated comparative headword & {\emph{*sái}} & too short to preserve
the \emph{*w} \ldots{} \emph{*i} environment needed for the documented
chronology & {\emph{sǣ}} & useful comparative label, but not the Old
English-facing input \\
selected i-stem input & {\emph{*sáiwiz}} & documented regular output:
{\emph{sǣ}} & {\emph{sǣ}} & exact match between derivational input and
Old English target \\
\end{longtable}
}

\section{\texorpdfstring{sieve --- OE
\emph{sife}}{sieve --- OE sife}}\label{sieve-oe-sife}

\index[iv]{02oe@\textbf{Old English}!sife@\emph{sife}}
\index[iv]{03pgmc@\textbf{Proto-Germanic}!sibaz@*síbaz}
\index[iv]{03pgmc@\textbf{Proto-Germanic}!sibi@*síbi}

Derivation: citation reconstruction \emph{*síbaz}; form followed here
\emph{*síbi} \textgreater{} \emph{sife} (early analogy).

\subsection{Derivation trace}\label{derivation-trace-95}

Proto input: \emph{*síbi}

\begingroup
\setlength{\fboxsep}{6pt}

\noindent\fbox{%
\begin{minipage}{0.97\linewidth}
\small
\begin{tabularx}{\linewidth}{@{}>{\raggedright\arraybackslash}p{0.440\linewidth}>{\centering\arraybackslash}X>{\raggedright\arraybackslash}p{0.440\linewidth}@{}}
\begin{minipage}[t]{\linewidth}
\centering\textbf{Earlier Germanic changes}\par
\vspace{0.35em}
\raggedright
\centering\textbf{}\par
\raggedright
\vspace{0.2em}
\begin{tabular}{@{}>{\raggedright\arraybackslash}p{0.68\linewidth}@{\hspace{0.55em}}>{\raggedright\arraybackslash}p{0.16\linewidth}@{\hspace{0.25em}}}
\multicolumn{2}{@{}l@{}}{*Proto-West Germanic*} \\
\multicolumn{2}{@{}l@{}}{[no change]} \\
\multicolumn{2}{@{}l@{}}{*Northwest Germanic*} \\
\multicolumn{2}{@{}l@{}}{[no change]} \\
\end{tabular}
\end{minipage}
&
&
\begin{minipage}[t]{\linewidth}
\centering\textbf{Old English changes}\par
\vspace{0.35em}
\raggedright
\centering\textbf{}\par
\raggedright
\vspace{0.2em}
\begin{tabular}{@{}>{\raggedright\arraybackslash}p{0.74\linewidth}@{\hspace{0.55em}}>{\raggedright\arraybackslash}p{0.16\linewidth}@{\hspace{0.25em}}}
\multicolumn{2}{@{}l@{}}{*Old English*} \\
PGmc B Allophony & \emph{*síβi} \\
OE Med Unstressed I Lowering1 & \emph{*síβe} \\
\end{tabular}
\end{minipage}
\\
\end{tabularx}
\end{minipage}%
} \endgroup

Old English form: \emph{sife}

\subsection{Reconstruction and comparative
evidence}\label{reconstruction-and-comparative-evidence-95}

Kluge-Seebold gives wg. \emph{*sibi-} n.~\ldots{} ae. \emph{sife}, and
Campbell groups {\emph{sife}} `sieve' with short neuter i-stems such as
\emph{spere} (\citeproc{ref-KlugeSeebold2011}{Kluge 2011, 847};
\citeproc{ref-Campbell1959}{Campbell 1959, sect. 609}). The older
morphological background is the s-stem \Recon{sib-iz} `sieve', but the
derivational input is the normalized i-stem form {\Recon{síbi} `sieve'}.

Kroonen's nearby \emph{*sebjō-} entry belongs to the separate kinship
lexeme that yields Old English {\emph{sibb}} `sieve', not to the sieve
word. Orel's {\Recon{sibaz} `sieve'} \ldots{} OE {\emph{sife}} `sieve'
preserves a broader handbook notation, but that a-stem shape does not
fit the Old English form treated here (\citeproc{ref-Orel2003}{Orel
2003, 328}).

\subsection{Old English evidence}\label{old-english-evidence-95}

Clark Hall gives {\emph{sibi}} `sieve' (GL) \ldots{} = {\emph{sife}}
`sieve' and also {\emph{sife}} `sieve' n.~`sieve'
(\citeproc{ref-ClarkHall1960}{Clark Hall 1960, 263}). Campbell likewise
cites Corpus Glossary \emph{sibi} `sieve' and treats \emph{sife} `sieve'
as a short neuter i-stem (\citeproc{ref-Campbell1959}{Campbell 1959,
sects 444, 609}). The normalized Old English target is therefore
\emph{sife} `sieve', while \emph{sibi} `sieve' is an attested earlier
spelling rather than a separate lexeme.

\subsection{Development to Old
English}\label{development-to-old-english-95}

From {\Recon{síbi} `sieve'}, the regular derivation gives \Recon{síβi}
`sieve' \textgreater{} \Recon{síβe} `sieve' \textgreater{} {\emph{sife}}
`sieve'. Medial \emph{b} is realized as a spirant and later written
\emph{f}, while the final unstressed \emph{i} lowers to \emph{e}. The
older s-stem background \Recon{sib-iz} `sieve' explains the morphology,
but the derivational input {\Recon{síbi}} is the immediate
pre-Old-English form.

\subsection{Stem comparison}\label{stem-comparison-6}

The comparison below sets the relevant forms side by side. It
distinguishes the accepted i-stem line from its rejected competitors.

{\def\LTcaptype{none} % do not increment counter
\begin{longtable}[]{@{}
  >{\raggedright\arraybackslash}p{(\linewidth - 8\tabcolsep) * \real{0.2000}}
  >{\raggedright\arraybackslash}p{(\linewidth - 8\tabcolsep) * \real{0.2000}}
  >{\raggedright\arraybackslash}p{(\linewidth - 8\tabcolsep) * \real{0.2000}}
  >{\raggedright\arraybackslash}p{(\linewidth - 8\tabcolsep) * \real{0.2000}}
  >{\raggedright\arraybackslash}p{(\linewidth - 8\tabcolsep) * \real{0.2000}}@{}}
\toprule\noalign{}
\begin{minipage}[b]{\linewidth}\raggedright
Formation / label
\end{minipage} & \begin{minipage}[b]{\linewidth}\raggedright
Candidate input
\end{minipage} & \begin{minipage}[b]{\linewidth}\raggedright
Expected or documented OE outcome
\end{minipage} & \begin{minipage}[b]{\linewidth}\raggedright
OE comparison form
\end{minipage} & \begin{minipage}[b]{\linewidth}\raggedright
Result
\end{minipage} \\
\midrule\noalign{}
\endhead
\bottomrule\noalign{}
\endlastfoot
ja-stem kinship line & Kroonen \emph{*sebjō-} / comparator
{\emph{*sibja}} & OE {\emph{sibb}} `sieve' & {\emph{sibb}} `sieve' &
separate lexeme, not the target treated here \\
a-stem handbook line & {\emph{*síbaz}} & expected {\emph{sif}} &
{\emph{sif}} & wrong ending for the attested noun \\
selected i-stem line from older \emph{*sib-iz} & {\emph{*síbi}} &
documented regular output: {\emph{sife}} `sieve' & {\emph{sife}}
`sieve'; early spelling {\emph{sibi}} `sieve' & exact match between
derivational input and Old English evidence \\
\end{longtable}
}

\section{\texorpdfstring{spare --- OE
\emph{sparian}}{spare --- OE sparian}}\label{spare-oe-sparian}

\index[iv]{02oe@\textbf{Old English}!sparian@\emph{sparian}}
\index[iv]{03pgmc@\textbf{Proto-Germanic}!sparena@*sparēną}
\index[iv]{03pgmc@\textbf{Proto-Germanic}!sparojana@*spárōjaną}

Derivation: citation reconstruction \emph{*sparēną}; form followed here
\emph{*spárōjaną} \textgreater{} \emph{sparian} (early analogy).

\subsection{Derivation trace}\label{derivation-trace-96}

Proto input: \emph{*spárōjaną}

\begingroup
\setlength{\fboxsep}{6pt}

\noindent\fbox{%
\begin{minipage}{0.97\linewidth}
\footnotesize
\begin{tabularx}{\linewidth}{@{}>{\raggedright\arraybackslash}p{0.320\linewidth}>{\centering\arraybackslash}X>{\raggedright\arraybackslash}p{0.520\linewidth}@{}}
\begin{minipage}[t]{\linewidth}
\centering\textbf{Earlier Germanic changes}\par
\vspace{0.35em}
\raggedright
\centering\textbf{}\par
\raggedright
\vspace{0.2em}
\begin{tabular}{@{}>{\raggedright\arraybackslash}p{0.68\linewidth}@{\hspace{0.55em}}>{\raggedright\arraybackslash}p{0.16\linewidth}@{\hspace{0.25em}}}
\multicolumn{2}{@{}l@{}}{*Proto-West Germanic*} \\
\multicolumn{2}{@{}l@{}}{[no change]} \\
\multicolumn{2}{@{}l@{}}{*Northwest Germanic*} \\
\multicolumn{2}{@{}l@{}}{[no change]} \\
\end{tabular}
\end{minipage}
&
&
\begin{minipage}[t]{\linewidth}
\centering\textbf{Old English changes}\par
\vspace{0.35em}
\raggedright
\centering\textbf{}\par
\raggedright
\vspace{0.2em}
\begin{tabular}{@{}>{\raggedright\arraybackslash}p{0.64\linewidth}@{\hspace{0.55em}}>{\raggedright\arraybackslash}p{0.26\linewidth}@{\hspace{0.25em}}}
\multicolumn{2}{@{}l@{}}{*Old English*} \\
Anglo Frisian Brightening & \emph{*spærōjaną} \\
OE A Restoration & \emph{*sparōjaną} \\
OE Heavy Syllable Nasal Apocope & \emph{*sparōjan} \\
OE Secondary Nasalization & \emph{*sparōjąn} \\
OE I Umlaut & \emph{*sparējąn} \\
OE Unstressed Long Vowel Shortening & \emph{*sparejąn} \\
OE Weak Tail Reduction & \emph{*sparejan} \\
OE Intervocalic J Vocalization & \emph{*spareian} \\
OE Unstressed EI Contraction & \emph{*sparian} \\
\end{tabular}
\end{minipage}
\\
\end{tabularx}
\end{minipage}%
} \endgroup

Old English form: \emph{sparian}

\subsection{Reconstruction and comparative
evidence}\label{reconstruction-and-comparative-evidence-96}

Kroonen keeps the inherited verb under class-III {\emph{*sparēn-}}
(\citeproc{ref-Kroonen2013}{Kroonen 2013, 465}). Orel similarly
preserves {\Recon{sparēnan} `spare'} (\citeproc{ref-Orel2003}{Orel 2003,
362}). Ringe and Taylor, however, reconstruct {\Recon{sparai-}}
\textasciitilde{} {\Recon{sparja-}} for the English branch and derive
the citation verb from a class-II line
(\citeproc{ref-RingeTaylor2014}{Ringe and Taylor 2014, 162, 191}). The
derivational input {\Recon{spárōjaną} `spare'} therefore represents the
refashioned class-II formation behind Old English {\emph{sparian}}
`spare', while the citation reconstruction {\Recon{sparēną} `spare'}
remains the inherited comparative headword.

\subsection{Old English evidence}\label{old-english-evidence-96}

Campbell says that {\emph{sparian}} `spare' does not show the ordinary
class-III characteristics, but the Ritual forms, normalized here as
{\emph{spæria}} `spare', {\emph{spær}} `spare', and {\emph{spærede}}
`spare', together with Vespasian Psalter {\emph{spearad}} `spare', point
to primitive Old English forms both with and without back vowels
(\citeproc{ref-Campbell1959}{Campbell 1959, sect. 764}). Brunner
likewise records Northumbrian {\emph{spæria}} `spare', {\emph{spærede}}
`spare' beside common Old English {\emph{sparian}} `spare' and Vespasian
Psalter {\emph{spearad}} `spare'
(\citeproc{ref-SieversBrunner1965}{Brunner 1965, sect. 364} Anm. 11).
The citation form treated here is {\emph{sparian}} `spare'; the Anglian
forms are relics of the older formation, not alternative headwords of
equal status.

\subsection{Development to Old
English}\label{development-to-old-english-96}

Once the class-II formation {\Recon{spárōjaną} `spare'} is chosen, the
remaining development is regular. The regular derivation shows
brightening, restoration of \emph{a} before the back vocalism of the
suffix, later i-mutation within the weak ending, weak-tail reduction,
and contraction to {\emph{sparian}} `spare'. By contrast, Brunner's rule
against further apocope of final \emph{-e} explains why Ritual
{\emph{spær}} `spare' cannot be the regular continuation of inherited
{\Recon{spárē} `spare'} (\citeproc{ref-SieversBrunner1965}{Brunner 1965,
sect. 150}).

\subsection{Formation comparison}\label{formation-comparison-4}

The comparison below sets the relevant forms side by side. It contrasts
the inherited class-III formation with the refashioned class-II one that
yields the citation verb.

{\def\LTcaptype{none} % do not increment counter
\begin{longtable}[]{@{}
  >{\raggedright\arraybackslash}p{(\linewidth - 8\tabcolsep) * \real{0.2000}}
  >{\raggedright\arraybackslash}p{(\linewidth - 8\tabcolsep) * \real{0.2000}}
  >{\raggedright\arraybackslash}p{(\linewidth - 8\tabcolsep) * \real{0.2000}}
  >{\raggedright\arraybackslash}p{(\linewidth - 8\tabcolsep) * \real{0.2000}}
  >{\raggedright\arraybackslash}p{(\linewidth - 8\tabcolsep) * \real{0.2000}}@{}}
\toprule\noalign{}
\begin{minipage}[b]{\linewidth}\raggedright
Formation / comparison
\end{minipage} & \begin{minipage}[b]{\linewidth}\raggedright
Candidate input
\end{minipage} & \begin{minipage}[b]{\linewidth}\raggedright
Expected or documented OE outcome
\end{minipage} & \begin{minipage}[b]{\linewidth}\raggedright
OE comparison form
\end{minipage} & \begin{minipage}[b]{\linewidth}\raggedright
Result
\end{minipage} \\
\midrule\noalign{}
\endhead
\bottomrule\noalign{}
\endlastfoot
inherited class-III infinitive & {\emph{*spárēną}} & paradigm comparison
/ probe output: {\emph{sparen}} & {\emph{sparian}} `spare' & wrong class
and wrong ending for the citation verb \\
inherited class-III imperative singular & \Recon{spárē} & paradigm
comparison / probe output: {\emph{spære}} & Ritual {\emph{spær}} `spare'
& loss of final \emph{-e} is not regular, so the relic form cannot
control the entry \\
inherited class-III finite present & \Recon{spárēθi} & paradigm
comparison / probe output: {\emph{spæreþ}} & {\emph{spearad}} `spare' &
attested form is mixed, not a direct continuation of the inherited
cell \\
selected class-II formation & {\emph{*spárōjaną}} & documented regular
output: {\emph{sparian}} `spare' & {\emph{sparian}} `spare' & exact
match between derivational input and Old English citation form \\
\end{longtable}
}

\section{\texorpdfstring{staff --- OE
\emph{stæf}}{staff --- OE stæf}}\label{staff-oe-stuxe6f}

\index[iv]{02oe@\textbf{Old English}!staef@\emph{stæf}}
\index[iv]{03pgmc@\textbf{Proto-Germanic}!stabaz@*stábaz}
\index[iv]{03pgmc@\textbf{Proto-Germanic}!stabiz@*stábiz}

Derivation: citation reconstruction \emph{*stábiz}; form followed here
\emph{*stábaz} \textgreater{} \emph{stæf} (early analogy).

\subsection{Derivation trace}\label{derivation-trace-97}

Proto input: \emph{*stábaz}

\begingroup
\setlength{\fboxsep}{6pt}

\noindent\fbox{%
\begin{minipage}{0.97\linewidth}
\small
\begin{tabularx}{\linewidth}{@{}>{\raggedright\arraybackslash}p{0.440\linewidth}>{\centering\arraybackslash}X>{\raggedright\arraybackslash}p{0.440\linewidth}@{}}
\begin{minipage}[t]{\linewidth}
\centering\textbf{Earlier Germanic changes}\par
\vspace{0.35em}
\raggedright
\centering\textbf{}\par
\raggedright
\vspace{0.2em}
\begin{tabular}{@{}>{\raggedright\arraybackslash}p{0.74\linewidth}@{\hspace{0.55em}}>{\raggedright\arraybackslash}p{0.16\linewidth}@{\hspace{0.25em}}}
\multicolumn{2}{@{}l@{}}{*Proto-West Germanic*} \\
\multicolumn{2}{@{}l@{}}{[no change]} \\
\multicolumn{2}{@{}l@{}}{*Northwest Germanic*} \\
\mbox{PGmc Final Z Deletion} & \emph{*stába} \\
\end{tabular}
\end{minipage}
&
&
\begin{minipage}[t]{\linewidth}
\centering\textbf{Old English changes}\par
\vspace{0.35em}
\raggedright
\centering\textbf{}\par
\raggedright
\vspace{0.2em}
\begin{tabular}{@{}>{\raggedright\arraybackslash}p{0.70\linewidth}@{\hspace{0.55em}}>{\raggedright\arraybackslash}p{0.16\linewidth}@{\hspace{0.25em}}}
\multicolumn{2}{@{}l@{}}{*Old English*} \\
PWGmc Final Bare A Loss & \emph{*stáb} \\
Anglo Frisian Brightening & \emph{*stæb} \\
PGmc B Allophony & \emph{*stæβ} \\
\end{tabular}
\end{minipage}
\\
\end{tabularx}
\end{minipage}%
} \endgroup

Old English form: \emph{stæf}

\subsection{Reconstruction and comparative
evidence}\label{reconstruction-and-comparative-evidence-97}

The comparative dictionaries do not give one uniform stem class. Kroonen
reconstructs an a-stem \emph{*staba-}
(\citeproc{ref-Kroonen2013}{Kroonen 2013, 471}). Orel writes
{\Recon{stábiz} `staff'} \textasciitilde{} {\Recon{stábaz} `staff'}
(\citeproc{ref-Orel2003}{Orel 2003, 368}). A direct i-stem input in
\Recon{-iz} `staff' would predict i-mutation in Old English, whereas the
attested noun keeps \emph{æ}.

\subsection{Old English evidence}\label{old-english-evidence-97}

The Old English noun itself is the ordinary citation form {\emph{stæf}}
`staff'. Luick lists {\emph{stæf}} `staff' among closed monosyllables
with \emph{æ} (\citeproc{ref-Luick1914}{Luick 1914-\/-1940, 176}). Ringe
and Taylor pair singular {\emph{stæf}} `staff' with plural
{\emph{stafas}} `staffs' (\citeproc{ref-RingeTaylor2014}{Ringe and
Taylor 2014, 193}). The normalized form here is therefore \emph{stæf};
later English \emph{staff} with \emph{a} belongs to a later stage of the
word's history.

\subsection{Development to Old
English}\label{development-to-old-english-97}

With the selected a-stem input, the development is regular. Final
\emph{*-z} disappears, bare final \emph{-a} is lost, Anglo-Frisian
brightening gives \emph{æ} in the closed monosyllable, and medial
\emph{b} surfaces as a fricative written \emph{f}. The documented chain
is \Recon{stábaz} `staff' \textgreater{} \Recon{stába} `staff'
\textgreater{} \Recon{stáb} `staff' \textgreater{} \emph{stæb} `staff'
\textgreater{} \emph{stæf} `staff'. A direct continuation of
\Recon{stábiz} `staff', by contrast, would produce i-mutated
{\emph{stefe}} `stave, staff' rather than the attested singular.

\subsection{Formation comparison}\label{formation-comparison-5}

The comparison below sets the relevant forms side by side. It separates
the rejected i-stem line from the selected a-stem input.

{\def\LTcaptype{none} % do not increment counter
\begin{longtable}[]{@{}
  >{\raggedright\arraybackslash}p{(\linewidth - 8\tabcolsep) * \real{0.2000}}
  >{\raggedright\arraybackslash}p{(\linewidth - 8\tabcolsep) * \real{0.2000}}
  >{\raggedright\arraybackslash}p{(\linewidth - 8\tabcolsep) * \real{0.2000}}
  >{\raggedright\arraybackslash}p{(\linewidth - 8\tabcolsep) * \real{0.2000}}
  >{\raggedright\arraybackslash}p{(\linewidth - 8\tabcolsep) * \real{0.2000}}@{}}
\toprule\noalign{}
\begin{minipage}[b]{\linewidth}\raggedright
Formation / label
\end{minipage} & \begin{minipage}[b]{\linewidth}\raggedright
Candidate input
\end{minipage} & \begin{minipage}[b]{\linewidth}\raggedright
OE output or comparison
\end{minipage} & \begin{minipage}[b]{\linewidth}\raggedright
OE comparison form
\end{minipage} & \begin{minipage}[b]{\linewidth}\raggedright
Result
\end{minipage} \\
\midrule\noalign{}
\endhead
\bottomrule\noalign{}
\endlastfoot
comparative i-stem line & {\emph{*stábiz}} & expected {\emph{stefe}}
after i-mutation & {\emph{stæf}} & wrong vowel for the attested
singular \\
mixed comparative notation & Orel \emph{*stabiz} \textasciitilde{}
\emph{*stabaz}; Kluge \emph{*stabi-}/a- & source-level stem-class
uncertainty & \emph{stæf} & useful comparative background, but not a
single OE-facing input \\
selected a-stem input & {\emph{*stábaz}} & documented regular output:
{\emph{stæf}} & {\emph{stæf}} & exact match between derivational input
and Old English target \\
\end{longtable}
}

\section{\texorpdfstring{swan --- OE
\emph{swanes}}{swan --- OE swanes}}\label{swan-oe-swanes}

\index[iv]{02oe@\textbf{Old English}!swanes@\emph{swanes}}
\index[iv]{03pgmc@\textbf{Proto-Germanic}!swanas@*swánas}
\index[iv]{03pgmc@\textbf{Proto-Germanic}!swanaz@*swánaz}

Derivation: citation reconstruction \emph{*swánaz}; form followed here
\emph{*swánas} \textgreater{} \emph{swanes} (early analogy).

\subsection{Derivation trace}\label{derivation-trace-98}

Proto input: \emph{*swánas}

\begingroup
\setlength{\fboxsep}{6pt}

\noindent\fbox{%
\begin{minipage}{0.97\linewidth}
\small
\begin{tabularx}{\linewidth}{@{}>{\raggedright\arraybackslash}p{0.440\linewidth}>{\centering\arraybackslash}X>{\raggedright\arraybackslash}p{0.440\linewidth}@{}}
\begin{minipage}[t]{\linewidth}
\centering\textbf{Earlier Germanic changes}\par
\vspace{0.35em}
\raggedright
\centering\textbf{}\par
\raggedright
\vspace{0.2em}
\begin{tabular}{@{}>{\raggedright\arraybackslash}p{0.68\linewidth}@{\hspace{0.55em}}>{\raggedright\arraybackslash}p{0.16\linewidth}@{\hspace{0.25em}}}
\multicolumn{2}{@{}l@{}}{*Proto-West Germanic*} \\
\multicolumn{2}{@{}l@{}}{[no change]} \\
\multicolumn{2}{@{}l@{}}{*Northwest Germanic*} \\
\multicolumn{2}{@{}l@{}}{[no change]} \\
\end{tabular}
\end{minipage}
&
&
\begin{minipage}[t]{\linewidth}
\centering\textbf{Old English changes}\par
\vspace{0.35em}
\raggedright
\centering\textbf{}\par
\raggedright
\vspace{0.2em}
\begin{tabular}{@{}>{\raggedright\arraybackslash}p{0.70\linewidth}@{\hspace{0.55em}}>{\raggedright\arraybackslash}p{0.16\linewidth}@{\hspace{0.25em}}}
\multicolumn{2}{@{}l@{}}{*Old English*} \\
Anglo Frisian Brightening & \emph{*swánæs} \\
OE Unstressed AE Merger & \emph{*swánes} \\
\end{tabular}
\end{minipage}
\\
\end{tabularx}
\end{minipage}%
} \endgroup

Old English form: \emph{swanes}

\subsection{Reconstruction and comparative
evidence}\label{reconstruction-and-comparative-evidence-98}

The Germanic noun is ordinarily cited as the masculine a-stem
{\Recon{swánaz} `swan'} (\citeproc{ref-Orel2003}{Orel 2003, 367}). The
derivational input \Recon{swánas} `swan' is not a competing lexeme
reconstruction. It is the genitive singular {\Recon{swánas}} of the same
paradigm.

The question here is therefore one of paradigm cell rather than stem
history. The citation form remains {\Recon{swánaz} `swan'}
\textgreater{} {\emph{swan}} `swan'; the comparison form is the genitive
singular {\Recon{swánas} `swan'} \textgreater{} {\emph{swanes}} `swan'.

\subsection{Old English evidence}\label{old-english-evidence-98}

Bright's glossary records the ordinary noun as {\emph{swan}} `swan', m.
and also gives the exact inflected form {\emph{swanes}} `swan', citing
the phrase \emph{swanes feðre}
(\citeproc{ref-BrightCassidyRingler1971}{Bright 1971, 441}).

The target is therefore an attested Old English genitive singular, not a
reconstruction or the ordinary citation lemma.

\subsection{Development to Old
English}\label{development-to-old-english-98}

From \Recon{swánas} `swan', Anglo-Frisian brightening gives
\Recon{swánæs} `swan', and subsequent merger of unstressed \emph{æ} with
\emph{e} yields \emph{swanes} `swan'. The comparison is straightforward
once the genitive singular is chosen as the relevant cell.

\subsection{Paradigm-cell comparison}\label{paradigm-cell-comparison}

The comparison below sets the relevant forms side by side. It separates
the ordinary citation form from the selected inflected cell.

{\def\LTcaptype{none} % do not increment counter
\begin{longtable}[]{@{}
  >{\raggedright\arraybackslash}p{(\linewidth - 8\tabcolsep) * \real{0.2000}}
  >{\raggedright\arraybackslash}p{(\linewidth - 8\tabcolsep) * \real{0.2000}}
  >{\raggedright\arraybackslash}p{(\linewidth - 8\tabcolsep) * \real{0.2000}}
  >{\raggedright\arraybackslash}p{(\linewidth - 8\tabcolsep) * \real{0.2000}}
  >{\raggedright\arraybackslash}p{(\linewidth - 8\tabcolsep) * \real{0.2000}}@{}}
\toprule\noalign{}
\begin{minipage}[b]{\linewidth}\raggedright
PGmc cell / interpretation
\end{minipage} & \begin{minipage}[b]{\linewidth}\raggedright
Candidate input
\end{minipage} & \begin{minipage}[b]{\linewidth}\raggedright
OE output or comparison
\end{minipage} & \begin{minipage}[b]{\linewidth}\raggedright
OE comparison form
\end{minipage} & \begin{minipage}[b]{\linewidth}\raggedright
Result
\end{minipage} \\
\midrule\noalign{}
\endhead
\bottomrule\noalign{}
\endlastfoot
citation nominative singular & {\emph{*swánaz}} & OE headword
{\emph{swan}} `swan' & {\emph{swan}} `swan' & ordinary lexeme line \\
genitive singular & {\emph{*swánas}} & regular output: {\emph{swanes}}
`swan' & {\emph{swanes}} `swan' & attested cell \\
\end{longtable}
}

\section{\texorpdfstring{thousand --- OE
\emph{þūsend}}{thousand --- OE þūsend}}\label{thousand-oe-uxfeux16bsend}

\index[iv]{02oe@\textbf{Old English}!aerende@\emph{ærende}}
\index[iv]{02oe@\textbf{Old English}!thusend@\emph{þūsend}}
\index[iv]{03pgmc@\textbf{Proto-Germanic}!thusendi@*θūs-èndi}
\index[iv]{03pgmc@\textbf{Proto-Germanic}!thusendi@*θūsèndi}
\index[iv]{03pgmc@\textbf{Proto-Germanic}!thusundi@*θūs-undī}

Derivation: citation reconstruction \emph{*θūs-undī}; form followed here
\emph{*θūs-èndi} \textgreater{} \emph{þūsend} (early analogy).

\subsection{Derivation trace}\label{derivation-trace-99}

Proto input: \emph{*θūsèndi}

\begingroup
\setlength{\fboxsep}{6pt}

\noindent\fbox{%
\begin{minipage}{0.97\linewidth}
\small
\begin{tabularx}{\linewidth}{@{}>{\raggedright\arraybackslash}p{0.440\linewidth}>{\centering\arraybackslash}X>{\raggedright\arraybackslash}p{0.440\linewidth}@{}}
\begin{minipage}[t]{\linewidth}
\centering\textbf{Earlier Germanic changes}\par
\vspace{0.35em}
\raggedright
\centering\textbf{}\par
\raggedright
\vspace{0.2em}
\begin{tabular}{@{}>{\raggedright\arraybackslash}p{0.68\linewidth}@{\hspace{0.55em}}>{\raggedright\arraybackslash}p{0.16\linewidth}@{\hspace{0.25em}}}
\multicolumn{2}{@{}l@{}}{*Proto-West Germanic*} \\
PWGmc Early I Apocope & \emph{*θūsènd} \\
\multicolumn{2}{@{}l@{}}{*Northwest Germanic*} \\
\multicolumn{2}{@{}l@{}}{[no change]} \\
\end{tabular}
\end{minipage}
&
&
\begin{minipage}[t]{\linewidth}
\centering\textbf{Old English changes}\par
\vspace{0.35em}
\raggedright
\centering\textbf{}\par
\raggedright
\vspace{0.2em}
\begin{tabular}{@{}>{\raggedright\arraybackslash}p{0.70\linewidth}@{\hspace{0.55em}}>{\raggedright\arraybackslash}p{0.16\linewidth}@{\hspace{0.25em}}}
\multicolumn{2}{@{}l@{}}{*Old English*} \\
OE Strip Secondary Stress & \emph{*θūsend} \\
\end{tabular}
\end{minipage}
\\
\end{tabularx}
\end{minipage}%
} \endgroup

Old English form: \emph{þūsend}

\subsection{Reconstruction and comparative
evidence}\label{reconstruction-and-comparative-evidence-99}

Kroonen reconstructs the Germanic numeral as {\emph{*þūsundī-}} and
cites Old English {\emph{þūsend}} `thousand' among its continuations
(\citeproc{ref-Kroonen2013}{Kroonen 2013, 554}). The derivational input
\textbf{{\Recon{θūs-èndi} `thousand'}} is not the same claim. It is an
OE-oriented transponent with the second-member vowel already resolved to
\emph{e} and the final high vowel already shortened for apocope.

The chronology must explain why Old English shows {\emph{þūsend}}
`thousand' while related languages such as Old Saxon and Old High German
keep \emph{u} in the second syllable?
(\citeproc{ref-Kroonen2013}{Kroonen 2013, 554}).

\subsection{Old English evidence}\label{old-english-evidence-99}

Old English \emph{þūsend} `thousand' is an ordinary citation form, not a
selected oblique or paradigm cell. Campbell treats it as a neuter noun
with normal case forms (\citeproc{ref-Campbell1959}{Campbell 1959, sect.
689}). The problem lies in the internal history of the word, not in its
lexical status.

\subsection{Development to Old
English}\label{development-to-old-english-99}

If the old final \emph{-ī} had remained long enough to trigger ordinary
double umlaut, Campbell's rule would point toward a form of the
{\emph{*þȳsend}} `thousand' type rather than attested {\emph{þūsend}}
`thousand' (\citeproc{ref-Campbell1959}{Campbell 1959, sect. 203}).
Preserved root \emph{ū} therefore argues that the umlaut-triggering
vowel was lost or neutralized before the ordinary OE umlaut could
develop.

That early loss, however, does not by itself explain the medial
\emph{e}. Luick compares the word with {\emph{ærende}} `errand/message'
and later groups \emph{thousand} with forms reshaped on that pattern
(\citeproc{ref-Luick1914}{Luick 1914-\/-1940, sects 198, 492}). Viredaz
is more cautious, arguing that Old English \emph{e} in this weak
position may simply write schwa and so need not prove a unique
\emph{ærende} `message'-type analogy
(\citeproc{ref-GermanicSlavicBaltic2025}{Viredaz 2025, sect. 2.1.4}).

The selected transponent \textbf{{\Recon{θūs-èndi} `thousand'}} captures
the OE-side state from which the regular derivation reaches
{\emph{þūsend}} `thousand'.

\subsection{Stage comparison}\label{stage-comparison-1}

The comparison below sets the relevant forms side by side. It separates
the secure chronology from the more interpretive account of the
second-syllable vowel.

{\def\LTcaptype{none} % do not increment counter
\begin{longtable}[]{@{}
  >{\raggedright\arraybackslash}p{(\linewidth - 6\tabcolsep) * \real{0.2500}}
  >{\raggedright\arraybackslash}p{(\linewidth - 6\tabcolsep) * \real{0.2500}}
  >{\raggedright\arraybackslash}p{(\linewidth - 6\tabcolsep) * \real{0.2500}}
  >{\raggedright\arraybackslash}p{(\linewidth - 6\tabcolsep) * \real{0.2500}}@{}}
\toprule\noalign{}
\begin{minipage}[b]{\linewidth}\raggedright
Stage / interpretation
\end{minipage} & \begin{minipage}[b]{\linewidth}\raggedright
Candidate form
\end{minipage} & \begin{minipage}[b]{\linewidth}\raggedright
OE relation
\end{minipage} & \begin{minipage}[b]{\linewidth}\raggedright
Result
\end{minipage} \\
\midrule\noalign{}
\endhead
\bottomrule\noalign{}
\endlastfoot
surviving \emph{-ī} with ordinary double umlaut & {\emph{*þūsundī-}}
treated as still umlaut-active in OE & would point toward
{\emph{*þȳsend}} & excluded by preserved \emph{ū} \\
early loss of the trigger without further reshaping & \emph{*þūsund-}
type & explains \emph{ū}, but not why OE alone has medial \emph{e} &
incomplete account \\
selected OE-oriented transponent & {\emph{*θūs-èndi}} & regular output:
{\emph{þūsend}} `thousand' & selected modeling input \\
\end{longtable}
}

\section{\texorpdfstring{timber --- OE
\emph{timber}}{timber --- OE timber}}\label{timber-oe-timber}

\index[iv]{02oe@\textbf{Old English}!timber@\emph{timber}}
\index[iv]{03pgmc@\textbf{Proto-Germanic}!timbra@*tímbrą}
\index[iv]{03pgmc@\textbf{Proto-Germanic}!timra@*tímrą}

Derivation: citation reconstruction \emph{*tímrą}; form followed here
\emph{*tímbrą} \textgreater{} \emph{timber} (early analogy).

\subsection{Derivation trace}\label{derivation-trace-100}

Proto input: \emph{*tímbrą}

\begingroup
\setlength{\fboxsep}{6pt}

\noindent\fbox{%
\begin{minipage}{0.97\linewidth}
\small
\begin{tabularx}{\linewidth}{@{}>{\raggedright\arraybackslash}p{0.440\linewidth}>{\centering\arraybackslash}X>{\raggedright\arraybackslash}p{0.440\linewidth}@{}}
\begin{minipage}[t]{\linewidth}
\centering\textbf{Earlier Germanic changes}\par
\vspace{0.35em}
\raggedright
\centering\textbf{}\par
\raggedright
\vspace{0.2em}
\begin{tabular}{@{}>{\raggedright\arraybackslash}p{0.68\linewidth}@{\hspace{0.55em}}>{\raggedright\arraybackslash}p{0.16\linewidth}@{\hspace{0.25em}}}
\multicolumn{2}{@{}l@{}}{*Proto-West Germanic*} \\
\multicolumn{2}{@{}l@{}}{[no change]} \\
\multicolumn{2}{@{}l@{}}{*Northwest Germanic*} \\
\multicolumn{2}{@{}l@{}}{[no change]} \\
\end{tabular}
\end{minipage}
&
&
\begin{minipage}[t]{\linewidth}
\centering\textbf{Old English changes}\par
\vspace{0.35em}
\raggedright
\centering\textbf{}\par
\raggedright
\vspace{0.2em}
\begin{tabular}{@{}>{\raggedright\arraybackslash}p{0.74\linewidth}@{\hspace{0.55em}}>{\raggedright\arraybackslash}p{0.16\linewidth}@{\hspace{0.25em}}}
\multicolumn{2}{@{}l@{}}{*Old English*} \\
OE Heavy Syllable Nasal Apocope & \emph{*tímbr} \\
OE Epenthetic Vowel & \emph{*tímber} \\
\end{tabular}
\end{minipage}
\\
\end{tabularx}
\end{minipage}%
} \endgroup

Old English form: \emph{timber}

\subsection{Reconstruction and comparative
evidence}\label{reconstruction-and-comparative-evidence-100}

Kroonen reconstructs the noun as {\emph{*timbra-}} and cites Old English
{\emph{timber}} `timber' among its continuations
(\citeproc{ref-Kroonen2013}{Kroonen 2013, 517}). Ringe and Taylor
instead state the history from \Recon{timra} `timber' through West
Germanic {\Recon{timbr} `timber'} to Old English {\emph{timber}}
`timber' (\citeproc{ref-RingeTaylor2014}{Ringe and Taylor 2014, 327}).

The difference is therefore not over the Old English noun itself. It
concerns whether medial \emph{b} belongs in the comparative citation
form or appears in an early pre-Old-English stage of the cluster.

\subsection{Old English evidence}\label{old-english-evidence-100}

Clark Hall lemmatizes the noun as {\emph{timber}} `timber' and also
records {\emph{timbor}} `timber' as a variant spelling
(\citeproc{ref-ClarkHall1960}{Clark Hall 1960, 294}). The Old English
form is thus an ordinary citation noun, not a selected oblique cell or a
reconstructed convenience form.

\subsection{Development to Old
English}\label{development-to-old-english-100}

With the consonantal frame \Recon{timbr-} `timber' in place, the rest of
the development is straightforward. Loss of final \emph{-ą} leaves
\Recon{tímbr} `timber', and epenthetic \emph{e} in the final cluster
yields {\emph{timber}} `timber'. Ringe and Taylor's \emph{timra}
\textgreater{} \emph{timbr} \textgreater{} OE \emph{timber} and the
handbook treatment of this epenthetic vowel point to the same Old
English result (\citeproc{ref-RingeTaylor2014}{Ringe and Taylor 2014,
327}; \citeproc{ref-Campbell1959}{Campbell 1959, sects 463--464}).

\subsection{Formation comparison}\label{formation-comparison-6}

The comparison below sets the relevant forms side by side. It separates
the comparative headword from the OE-facing consonantal input.

{\def\LTcaptype{none} % do not increment counter
\begin{longtable}[]{@{}
  >{\raggedright\arraybackslash}p{(\linewidth - 6\tabcolsep) * \real{0.2500}}
  >{\raggedright\arraybackslash}p{(\linewidth - 6\tabcolsep) * \real{0.2500}}
  >{\raggedright\arraybackslash}p{(\linewidth - 6\tabcolsep) * \real{0.2500}}
  >{\raggedright\arraybackslash}p{(\linewidth - 6\tabcolsep) * \real{0.2500}}@{}}
\toprule\noalign{}
\begin{minipage}[b]{\linewidth}\raggedright
Formation or notation
\end{minipage} & \begin{minipage}[b]{\linewidth}\raggedright
Candidate form
\end{minipage} & \begin{minipage}[b]{\linewidth}\raggedright
OE relation
\end{minipage} & \begin{minipage}[b]{\linewidth}\raggedright
Result
\end{minipage} \\
\midrule\noalign{}
\endhead
\bottomrule\noalign{}
\endlastfoot
Kroonen's comparative citation & {\emph{*timbra-}} & already matches the
consonantal frame of OE {\emph{timber}} `timber' & closest comparative
support for the derivational input \\
Ringe-Taylor citation line & {\emph{*timra}} \textgreater{}
{\Recon{timbr} `timber'} & reaches the same OE noun through early
cluster expansion & compatible comparative background \\
modeled input & {\emph{*tímbrą}} & regular output: {\emph{timber}}
`timber' & Old English-facing input \\
\end{longtable}
}

\section{\texorpdfstring{wake --- OE
\emph{wacan}}{wake --- OE wacan}}\label{wake-oe-wacan}

\index[iv]{02oe@\textbf{Old English}!wacan@\emph{wacan}}
\index[iv]{03pgmc@\textbf{Proto-Germanic}!wakana@*wákaną}
\index[iv]{03pgmc@\textbf{Proto-Germanic}!wakena@*wakēną}

Derivation: citation reconstruction \emph{*wakēną}; form followed here
\emph{*wákaną} \textgreater{} \emph{wacan} (early analogy).

\subsection{Derivation trace}\label{derivation-trace-101}

Proto input: \emph{*wákaną}

\begingroup
\setlength{\fboxsep}{6pt}

\noindent\fbox{%
\begin{minipage}{0.97\linewidth}
\small
\begin{tabularx}{\linewidth}{@{}>{\raggedright\arraybackslash}p{0.320\linewidth}>{\centering\arraybackslash}X>{\raggedright\arraybackslash}p{0.520\linewidth}@{}}
\begin{minipage}[t]{\linewidth}
\centering\textbf{Earlier Germanic changes}\par
\vspace{0.35em}
\raggedright
\centering\textbf{}\par
\raggedright
\vspace{0.2em}
\begin{tabular}{@{}>{\raggedright\arraybackslash}p{0.68\linewidth}@{\hspace{0.55em}}>{\raggedright\arraybackslash}p{0.16\linewidth}@{\hspace{0.25em}}}
\multicolumn{2}{@{}l@{}}{*Proto-West Germanic*} \\
\multicolumn{2}{@{}l@{}}{[no change]} \\
\multicolumn{2}{@{}l@{}}{*Northwest Germanic*} \\
\multicolumn{2}{@{}l@{}}{[no change]} \\
\end{tabular}
\end{minipage}
&
&
\begin{minipage}[t]{\linewidth}
\centering\textbf{Old English changes}\par
\vspace{0.35em}
\raggedright
\centering\textbf{}\par
\raggedright
\vspace{0.2em}
\begin{tabular}{@{}>{\raggedright\arraybackslash}p{0.74\linewidth}@{\hspace{0.55em}}>{\raggedright\arraybackslash}p{0.16\linewidth}@{\hspace{0.25em}}}
\multicolumn{2}{@{}l@{}}{*Old English*} \\
Anglo Frisian Brightening & \emph{*wækaną} \\
OE A Restoration & \emph{*wakaną} \\
OE Heavy Syllable Nasal Apocope & \emph{*wakan} \\
OE Secondary Nasalization & \emph{*wakąn} \\
OE Weak Tail Reduction & \emph{*wakan} \\
\end{tabular}
\end{minipage}
\\
\end{tabularx}
\end{minipage}%
} \endgroup

Old English form: \emph{wacan}

\subsection{Reconstruction and comparative
evidence}\label{reconstruction-and-comparative-evidence-101}

Kroonen gives the strong verb as \emph{*wakan-} with Old English
{\emph{wacan}} `wake' (\citeproc{ref-Kroonen2013}{Kroonen 2013, 568}).
Ringe and Taylor separately derive Old English {\emph{wacian}} `wake'
from weak \emph{*wakai-} \textasciitilde{} \emph{*wakja-}
(\citeproc{ref-RingeTaylor2014}{Ringe and Taylor 2014, sect. 3.3.2}).

The difference is therefore lexical and class-based, not graphic. Strong
{\emph{wacan}} `wake' `wake up, arise' and weak {\emph{wacian}} `wake'
`be awake, watch' belong to related but distinct histories.

\subsection{Old English evidence}\label{old-english-evidence-101}

Clark Hall lists {\emph{wacan}} `wake' and {\emph{wacian}} `wake' as
separate headwords (\citeproc{ref-ClarkHall1960}{Clark Hall 1960, 338}).
Bosworth-Toller cautions under {\emph{wacan}} `wake': the simplex
infinitive itself does not occur, its place seeming to be taken by
{\emph{wæcnan}} `wake' (\citeproc{ref-BosworthToller1898}{Bosworth and
Toller 1898, 226}).

The target {\emph{wacan}} `wake' is therefore best understood as a
normalized strong headword for the verb family, not as a directly quoted
simplex infinitive. It still remains the correct Old English comparison
form for the strong branch.

\subsection{Development to Old
English}\label{development-to-old-english-101}

With strong \textbf{{\Recon{wákaną} `wake'}}, Anglo-Frisian brightening
first gives a form of the \textbf{\Recon{wækaną} `wake'} type.
A-restoration then returns \emph{a}, and the ordinary tail reductions
yield {\emph{wacan}} `wake'. The weak verb {\emph{wacian}} `wake'
belongs to a different prehistory and is not the expected outcome of
this input.

\subsection{Class comparison}\label{class-comparison-4}

The comparison below sets the relevant forms side by side. It separates
the strong and weak verb lines.

{\def\LTcaptype{none} % do not increment counter
\begin{longtable}[]{@{}
  >{\raggedright\arraybackslash}p{(\linewidth - 6\tabcolsep) * \real{0.2500}}
  >{\raggedright\arraybackslash}p{(\linewidth - 6\tabcolsep) * \real{0.2500}}
  >{\raggedright\arraybackslash}p{(\linewidth - 6\tabcolsep) * \real{0.2500}}
  >{\raggedright\arraybackslash}p{(\linewidth - 6\tabcolsep) * \real{0.2500}}@{}}
\toprule\noalign{}
\begin{minipage}[b]{\linewidth}\raggedright
Formation / class
\end{minipage} & \begin{minipage}[b]{\linewidth}\raggedright
Candidate input
\end{minipage} & \begin{minipage}[b]{\linewidth}\raggedright
OE outcome or comparison
\end{minipage} & \begin{minipage}[b]{\linewidth}\raggedright
Result
\end{minipage} \\
\midrule\noalign{}
\endhead
\bottomrule\noalign{}
\endlastfoot
weak class-III / class-II branch & {\emph{*wakēną}}, \emph{*wakai-}
\textasciitilde{} \emph{*wakja-} & OE {\emph{wacian}} `wake' and related
weak forms & related lexeme, but not the target of this entry \\
strong class-VI branch & {\emph{*wákaną}} & regular output:
{\emph{wacan}} `wake' & Old English-facing input \\
strong normalized headword & {\emph{wacan}} `wake' & dictionary
comparison form beside attested strong-family forms & correct Old
English comparator, though not a directly quoted simplex infinitive \\
\end{longtable}
}

\section{\texorpdfstring{water --- OE
\emph{wæter}}{water --- OE wæter}}\label{water-oe-wuxe6ter}

\index[iv]{02oe@\textbf{Old English}!waeter@\emph{wæter}}
\index[iv]{03pgmc@\textbf{Proto-Germanic}!watna@*wátną}
\index[iv]{03pgmc@\textbf{Proto-Germanic}!wator@*wátōr}

Derivation: citation reconstruction \emph{*wátną}; form followed here
\emph{*wátōr} \textgreater{} \emph{wæter} (early analogy).

\subsection{Derivation trace}\label{derivation-trace-102}

Proto input: \emph{*wátōr}

\begingroup
\setlength{\fboxsep}{6pt}

\noindent\fbox{%
\begin{minipage}{0.97\linewidth}
\small
\begin{tabularx}{\linewidth}{@{}>{\raggedright\arraybackslash}p{0.440\linewidth}>{\centering\arraybackslash}X>{\raggedright\arraybackslash}p{0.440\linewidth}@{}}
\begin{minipage}[t]{\linewidth}
\centering\textbf{Earlier Germanic changes}\par
\vspace{0.35em}
\raggedright
\centering\textbf{}\par
\raggedright
\vspace{0.2em}
\begin{tabular}{@{}>{\raggedright\arraybackslash}p{0.70\linewidth}@{\hspace{0.55em}}>{\raggedright\arraybackslash}p{0.16\linewidth}@{\hspace{0.25em}}}
\multicolumn{2}{@{}l@{}}{*Proto-West Germanic*} \\
PWGmc Final Or Lowering & \emph{*wátar} \\
\multicolumn{2}{@{}l@{}}{*Northwest Germanic*} \\
\multicolumn{2}{@{}l@{}}{[no change]} \\
\end{tabular}
\end{minipage}
&
&
\begin{minipage}[t]{\linewidth}
\centering\textbf{Old English changes}\par
\vspace{0.35em}
\raggedright
\centering\textbf{}\par
\raggedright
\vspace{0.2em}
\begin{tabular}{@{}>{\raggedright\arraybackslash}p{0.70\linewidth}@{\hspace{0.55em}}>{\raggedright\arraybackslash}p{0.16\linewidth}@{\hspace{0.25em}}}
\multicolumn{2}{@{}l@{}}{*Old English*} \\
Anglo Frisian Brightening & \emph{*wætær} \\
OE Unstressed AE Merger & \emph{*wæter} \\
\end{tabular}
\end{minipage}
\\
\end{tabularx}
\end{minipage}%
} \endgroup

Old English form: \emph{wæter}

\subsection{Reconstruction and comparative
evidence}\label{reconstruction-and-comparative-evidence-102}

Kroonen reconstructs a heteroclitic noun \emph{*watar-}
\textasciitilde{} \emph{*watan-} and states that the Proto-Germanic
material points to \Recon{watōr} `water' and \Recon{watenaz} `water'
(\citeproc{ref-Kroonen2013}{Kroonen 2013, 616}). Ringe and Taylor
likewise start from singular \Recon{wator} `water' before the Old
English branch (\citeproc{ref-RingeTaylor2014}{Ringe and Taylor 2014,
sect. 3.1.4}).

The generalized comparative label is therefore broader than the singular
form that actually corresponds to Old English {\emph{wæter}} `water'.
The relevant comparator is the inherited nominative-accusative singular
{\Recon{wátōr} `water'}.

\subsection{Old English evidence}\label{old-english-evidence-102}

Bright gives the noun as {\emph{wæter}} `water' with the regular
paradigm {\emph{wæteres}} `water (gen.)', {\emph{wætere}} `water
(dat.)', \emph{wæter(u)}, {\emph{wætera}} `water (gen.pl.)',
{\emph{wæterum}} `water (dat.pl.)'
(\citeproc{ref-BrightCassidyRingler1971}{Bright 1971, 29}). Ringe and
Taylor add the dialectal contrast between West Saxon \emph{weeter}
`water' and Mercian \emph{weter} `water'
(\citeproc{ref-RingeTaylor2014}{Ringe and Taylor 2014, sect. 6.5.2}).

The target is therefore an attested Old English citation form within a
normal paradigm. The complication lies on the comparative side of the
lexeme, not in Old English attestation.

\subsection{Development to Old
English}\label{development-to-old-english-102}

From \Recon{wátōr} `water', pre-final \emph{ō} becomes \emph{a},
yielding \Recon{watar} `water' (\citeproc{ref-RingeTaylor2014}{Ringe and
Taylor 2014, sect. 3.1.4}). Anglo-Frisian brightening then gives
\Recon{wætær} `water', and merger of unstressed \emph{æ}/\emph{e} yields
{\emph{wæter}} `water'.

\subsection{Stage comparison}\label{stage-comparison-2}

The comparison below sets the relevant forms side by side. It separates
the generalized lexeme label from the singular input that matches the
Old English citation form.

{\def\LTcaptype{none} % do not increment counter
\begin{longtable}[]{@{}
  >{\raggedright\arraybackslash}p{(\linewidth - 6\tabcolsep) * \real{0.2500}}
  >{\raggedright\arraybackslash}p{(\linewidth - 6\tabcolsep) * \real{0.2500}}
  >{\raggedright\arraybackslash}p{(\linewidth - 6\tabcolsep) * \real{0.2500}}
  >{\raggedright\arraybackslash}p{(\linewidth - 6\tabcolsep) * \real{0.2500}}@{}}
\toprule\noalign{}
\begin{minipage}[b]{\linewidth}\raggedright
Stage or notation
\end{minipage} & \begin{minipage}[b]{\linewidth}\raggedright
Candidate form
\end{minipage} & \begin{minipage}[b]{\linewidth}\raggedright
OE relation
\end{minipage} & \begin{minipage}[b]{\linewidth}\raggedright
Result
\end{minipage} \\
\midrule\noalign{}
\endhead
\bottomrule\noalign{}
\endlastfoot
generalized comparative label & {\emph{*wátną}} & broader lexeme
shorthand, not the singular that corresponds directly to {\emph{wæter}}
`water' & useful background only \\
heteroclitic stem notation & \emph{*watar-} \textasciitilde{}
\emph{*watan-} & source-faithful comparative reconstruction & explains
why a singular comparator is needed \\
inherited singular input & {\emph{*wátōr}} & regular output:
{\emph{wæter}} `water' & Old English-facing input \\
\end{longtable}
}

\section{\texorpdfstring{whale --- OE
\emph{hwæl}}{whale --- OE hwæl}}\label{whale-oe-hwuxe6l}

\index[iv]{02oe@\textbf{Old English}!hwael@\emph{hwæl}}
\index[iv]{03pgmc@\textbf{Proto-Germanic}!walaz@*wálaz}
\index[iv]{03pgmc@\textbf{Proto-Germanic}!xwalaz@*xwálaz}

Derivation: citation reconstruction \emph{*wálaz}; form followed here
\emph{*xwálaz} \textgreater{} \emph{hwæl} (early analogy).

\subsection{Derivation trace}\label{derivation-trace-103}

Proto input: \emph{*xwálaz}

\begingroup
\setlength{\fboxsep}{6pt}

\noindent\fbox{%
\begin{minipage}{0.97\linewidth}
\small
\begin{tabularx}{\linewidth}{@{}>{\raggedright\arraybackslash}p{0.440\linewidth}>{\centering\arraybackslash}X>{\raggedright\arraybackslash}p{0.440\linewidth}@{}}
\begin{minipage}[t]{\linewidth}
\centering\textbf{Earlier Germanic changes}\par
\vspace{0.35em}
\raggedright
\centering\textbf{}\par
\raggedright
\vspace{0.2em}
\begin{tabular}{@{}>{\raggedright\arraybackslash}p{0.74\linewidth}@{\hspace{0.55em}}>{\raggedright\arraybackslash}p{0.16\linewidth}@{\hspace{0.25em}}}
\multicolumn{2}{@{}l@{}}{*Proto-West Germanic*} \\
\multicolumn{2}{@{}l@{}}{[no change]} \\
\multicolumn{2}{@{}l@{}}{*Northwest Germanic*} \\
\mbox{PGmc Final Z Deletion} & \emph{*xwála} \\
\end{tabular}
\end{minipage}
&
&
\begin{minipage}[t]{\linewidth}
\centering\textbf{Old English changes}\par
\vspace{0.35em}
\raggedright
\centering\textbf{}\par
\raggedright
\vspace{0.2em}
\begin{tabular}{@{}>{\raggedright\arraybackslash}p{0.70\linewidth}@{\hspace{0.55em}}>{\raggedright\arraybackslash}p{0.16\linewidth}@{\hspace{0.25em}}}
\multicolumn{2}{@{}l@{}}{*Old English*} \\
PWGmc Final Bare A Loss & \emph{*xwál} \\
Anglo Frisian Brightening & \emph{*xwæl} \\
\end{tabular}
\end{minipage}
\\
\end{tabularx}
\end{minipage}%
} \endgroup

Old English form: \emph{hwæl}

\subsection{Reconstruction and comparative
evidence}\label{reconstruction-and-comparative-evidence-103}

The comparative sources are not uniform. Orel gives {\Recon{xwalaz}
`whale'} and notes some mixed {\Recon{xwaliz} `whale'} evidence
(\citeproc{ref-Orel2003}{Orel 2003, 197}). Kroonen instead cites
{\emph{*hwali-}} (\citeproc{ref-Kroonen2013}{Kroonen 2013, 262}).

Both notations agree on inherited initial \emph{hw-/xw-}, but they
differ in stem label. The a-stem-like input followed here is closer to
Orel's notation than to Kroonen's citation form.

\subsection{Old English evidence}\label{old-english-evidence-103}

Clark Hall lemmatizes the noun as {\emph{hwal}} `whale', and
Bosworth-Toller preserves plural {\emph{hwalas}} `whales'
(\citeproc{ref-ClarkHall1960}{Clark Hall 1960, 170};
\citeproc{ref-BosworthToller1898}{Bosworth and Toller 1898, 326}). The
comparison form is normalized here as {\emph{hwæl}} `whale' for the
singular citation form with Anglo-Frisian fronting.

The plural {\emph{hwalas}} `whales' supplies control evidence. It shows
the same lexeme with \emph{a} in an open syllable beside singular
{\emph{hwæl}} `whale' in a closed monosyllable.

\subsection{Development to Old
English}\label{development-to-old-english-103}

From \Recon{xwálaz} `whale', final \emph{-z} disappears and bare final
\emph{-a} is lost. Anglo-Frisian fronting then yields \emph{æ} in the
closed monosyllable, and Old English orthography writes {\emph{hwæl}}
`whale'.

\subsection{Formation comparison}\label{formation-comparison-7}

The comparison below sets the relevant forms side by side. It separates
the competing comparative notations from the normalized Old English
singular.

{\def\LTcaptype{none} % do not increment counter
\begin{longtable}[]{@{}
  >{\raggedright\arraybackslash}p{(\linewidth - 6\tabcolsep) * \real{0.2500}}
  >{\raggedright\arraybackslash}p{(\linewidth - 6\tabcolsep) * \real{0.2500}}
  >{\raggedright\arraybackslash}p{(\linewidth - 6\tabcolsep) * \real{0.2500}}
  >{\raggedright\arraybackslash}p{(\linewidth - 6\tabcolsep) * \real{0.2500}}@{}}
\toprule\noalign{}
\begin{minipage}[b]{\linewidth}\raggedright
Comparative line
\end{minipage} & \begin{minipage}[b]{\linewidth}\raggedright
Candidate form
\end{minipage} & \begin{minipage}[b]{\linewidth}\raggedright
OE relation
\end{minipage} & \begin{minipage}[b]{\linewidth}\raggedright
Result
\end{minipage} \\
\midrule\noalign{}
\endhead
\bottomrule\noalign{}
\endlastfoot
Orel's citation & {\emph{*xwalaz}} & same stem notation as the modeled
singular line & closest comparative support for the derivational
input \\
Kroonen's citation & {\emph{*hwali-}} & same initial cluster, different
stem label & important comparative rival, but not the notation followed
here \\
modeled input & {\emph{*xwálaz}} & regular output: {\emph{hwæl}} `whale'
& Old English-facing input \\
plural control & {\emph{hwalas}} `whales' & attested open-syllable
plural beside singular {\emph{hwæl}} `whale' & confirms that the lexeme
also preserves an \emph{a}-vocalism branch \\
\end{longtable}
}

\section{\texorpdfstring{whine --- OE
\emph{hwīnan}}{whine --- OE hwīnan}}\label{whine-oe-hwux12bnan}

\index[iv]{02oe@\textbf{Old English}!hwinan@\emph{hwīnan}}
\index[iv]{03pgmc@\textbf{Proto-Germanic}!wainojana@*wainōjaną}
\index[iv]{03pgmc@\textbf{Proto-Germanic}!xwinana@*xwḯnaną}

Derivation: citation reconstruction \emph{*wainōjaną}; form followed
here \emph{*xwī́naną} \textgreater{} \emph{hwīnan} (early analogy).

\subsection{Derivation trace}\label{derivation-trace-104}

Proto input: \emph{*xwī́naną}

\begingroup
\setlength{\fboxsep}{6pt}

\noindent\fbox{%
\begin{minipage}{0.97\linewidth}
\small
\begin{tabularx}{\linewidth}{@{}>{\raggedright\arraybackslash}p{0.320\linewidth}>{\centering\arraybackslash}X>{\raggedright\arraybackslash}p{0.520\linewidth}@{}}
\begin{minipage}[t]{\linewidth}
\centering\textbf{Earlier Germanic changes}\par
\vspace{0.35em}
\raggedright
\centering\textbf{}\par
\raggedright
\vspace{0.2em}
\begin{tabular}{@{}>{\raggedright\arraybackslash}p{0.68\linewidth}@{\hspace{0.55em}}>{\raggedright\arraybackslash}p{0.16\linewidth}@{\hspace{0.25em}}}
\multicolumn{2}{@{}l@{}}{*Proto-West Germanic*} \\
\multicolumn{2}{@{}l@{}}{[no change]} \\
\multicolumn{2}{@{}l@{}}{*Northwest Germanic*} \\
\multicolumn{2}{@{}l@{}}{[no change]} \\
\end{tabular}
\end{minipage}
&
&
\begin{minipage}[t]{\linewidth}
\centering\textbf{Old English changes}\par
\vspace{0.35em}
\raggedright
\centering\textbf{}\par
\raggedright
\vspace{0.2em}
\begin{tabular}{@{}>{\raggedright\arraybackslash}p{0.74\linewidth}@{\hspace{0.55em}}>{\raggedright\arraybackslash}p{0.16\linewidth}@{\hspace{0.25em}}}
\multicolumn{2}{@{}l@{}}{*Old English*} \\
OE Heavy Syllable Nasal Apocope & \emph{*xwī́nan} \\
OE Secondary Nasalization & \emph{*xwī́nąn} \\
OE Weak Tail Reduction & \emph{*xwī́nan} \\
\end{tabular}
\end{minipage}
\\
\end{tabularx}
\end{minipage}%
} \endgroup

Old English form: \emph{hwīnan}

\subsection{Reconstruction and comparative
evidence}\label{reconstruction-and-comparative-evidence-104}

The citation reconstruction preserved in the header belongs to the
lament-family verb seen in German {\emph{weinen}} `whine' and Old
English {\emph{wānian}} `wane'. Kroonen instead separates Old English
{\emph{hwīnan}} `whine' under \emph{*hwinan-}
(\citeproc{ref-Kroonen2013}{Kroonen 2013, 267}). Orel likewise
distinguishes strong {\Recon{xwinanan} `whine'} from weak
{\Recon{wainōjanan} `whine'} (\citeproc{ref-Orel2003}{Orel 2003, 201}).
Ringe and Taylor make the same split at the Northwest Germanic level,
linking Old Norse {\emph{hvina}} `whine' and Old English {\emph{hwinan}}
`whine' to the same strong verb (\citeproc{ref-RingeTaylor2014}{Ringe
and Taylor 2014, 130}).

The two families also differ phonologically and morphologically. The
lament family has initial \emph{w-}, diphthongal \emph{ai}, and weak-II
morphology, whereas the verb behind Old English \emph{hwīnan} `whine'
has initial \emph{hw-/xw-}, long \emph{ī}, and strong-verb inflection.
The derivational input {\Recon{xwī́naną} `whine'} therefore represents a
competing comparative identification rather than a hidden cell of
\Recon{wainōjaną} `whine'.

\subsection{Old English evidence}\label{old-english-evidence-104}

Clark Hall records {\emph{hwinan}} `whine' with the gloss `to hiss,
whizz, whistle' (\citeproc{ref-ClarkHall1960}{Clark Hall 1960, 171}).
Seebold keeps the verb among the strong verbs and notes that only a
present-tense attestation is directly preserved
(\citeproc{ref-Seebold1970}{Seebold 1970, 280}).

The Old English form is normalized here as {\emph{hwīnan}} `whine'. That
normalization adds the usual vowel length marking to the dictionary
spelling {\emph{hwinan}} `whine'; it does not turn an unattested verb
into a reconstructed one.

\subsection{Development to Old
English}\label{development-to-old-english-104}

Once the strong-verb input {\Recon{xwī́naną} `whine'} is selected, the
path to Old English is straightforward. The compact trace shows
heavy-syllable nasal apocope, secondary nasalization, and weak-tail
reduction, after which the form surfaces as {\emph{hwīnan}} `whine'.

No special paradigm maneuver is needed for this verb. The comparison is
between two different Germanic verb families: the Old English form
belongs with the strong verb \emph{*hwīnan-}, not with the weak lament
verb.

\subsection{Verb-family comparison}\label{verb-family-comparison}

The comparison below sets the relevant forms side by side. It separates
the competing comparative labels that stand behind the inherited Old
English forms.

{\def\LTcaptype{none} % do not increment counter
\begin{longtable}[]{@{}
  >{\raggedright\arraybackslash}p{(\linewidth - 8\tabcolsep) * \real{0.2000}}
  >{\raggedright\arraybackslash}p{(\linewidth - 8\tabcolsep) * \real{0.2000}}
  >{\raggedright\arraybackslash}p{(\linewidth - 8\tabcolsep) * \real{0.2000}}
  >{\raggedright\arraybackslash}p{(\linewidth - 8\tabcolsep) * \real{0.2000}}
  >{\raggedright\arraybackslash}p{(\linewidth - 8\tabcolsep) * \real{0.2000}}@{}}
\toprule\noalign{}
\begin{minipage}[b]{\linewidth}\raggedright
Verb family / interpretation
\end{minipage} & \begin{minipage}[b]{\linewidth}\raggedright
Candidate input
\end{minipage} & \begin{minipage}[b]{\linewidth}\raggedright
Old English outcome or comparison
\end{minipage} & \begin{minipage}[b]{\linewidth}\raggedright
OE comparison form
\end{minipage} & \begin{minipage}[b]{\linewidth}\raggedright
Result
\end{minipage} \\
\midrule\noalign{}
\endhead
\bottomrule\noalign{}
\endlastfoot
lament-family weak verb & {\emph{*wainōjaną}} & comparative continuation
in OE {\emph{wānian}} `wane' & {\emph{wānian}} `wane' & competing
citation reconstruction, but not the source of \emph{hwīnan} \\
selected strong verb & {\emph{*xwī́naną}} & regular output:
{\emph{hwīnan}} `whine' & {\emph{hwīnan}} `whine' & exact match between
derivational input and OE verb \\
comparative North Germanic cognate & Northwest Germanic strong verb
behind ON {\emph{hvina}} `whine' / OE {\emph{hwinan}} `whine' & ON
{\emph{hvina}} `whine' / OE {\emph{hwinan}} `whine' & {\emph{hwīnan}}
`whine' & supports the strong-verb identification \\
\end{longtable}
}

\section{\texorpdfstring{withy --- OE
\emph{wīþiġ}}{withy --- OE wīþiġ}}\label{withy-oe-wux12buxfeiux121}

\index[iv]{02oe@\textbf{Old English}!withig@\emph{wīþiġ}}
\index[iv]{03pgmc@\textbf{Proto-Germanic}!waithiz@*wáiθiz}
\index[iv]{03pgmc@\textbf{Proto-Germanic}!withaga@*wḯθagą}

Derivation: citation reconstruction \emph{*wáiθiz}; form followed here
\emph{*wī́θagą} \textgreater{} \emph{wīþiġ} (early analogy).

\subsection{Derivation trace}\label{derivation-trace-105}

Proto input: \emph{*wī́θagą}

\begingroup
\setlength{\fboxsep}{6pt}

\noindent\fbox{%
\begin{minipage}{0.97\linewidth}
\small
\begin{tabularx}{\linewidth}{@{}>{\raggedright\arraybackslash}p{0.320\linewidth}>{\centering\arraybackslash}X>{\raggedright\arraybackslash}p{0.520\linewidth}@{}}
\begin{minipage}[t]{\linewidth}
\centering\textbf{Earlier Germanic changes}\par
\vspace{0.35em}
\raggedright
\centering\textbf{}\par
\raggedright
\vspace{0.2em}
\begin{tabular}{@{}>{\raggedright\arraybackslash}p{0.68\linewidth}@{\hspace{0.55em}}>{\raggedright\arraybackslash}p{0.16\linewidth}@{\hspace{0.25em}}}
\multicolumn{2}{@{}l@{}}{*Proto-West Germanic*} \\
\multicolumn{2}{@{}l@{}}{[no change]} \\
\multicolumn{2}{@{}l@{}}{*Northwest Germanic*} \\
\multicolumn{2}{@{}l@{}}{[no change]} \\
\end{tabular}
\end{minipage}
&
&
\begin{minipage}[t]{\linewidth}
\centering\textbf{Old English changes}\par
\vspace{0.35em}
\raggedright
\centering\textbf{}\par
\raggedright
\vspace{0.2em}
\begin{tabular}{@{}>{\raggedright\arraybackslash}p{0.74\linewidth}@{\hspace{0.55em}}>{\raggedright\arraybackslash}p{0.16\linewidth}@{\hspace{0.25em}}}
\multicolumn{2}{@{}l@{}}{*Old English*} \\
Anglo Frisian Brightening & \emph{*wī́θægą} \\
OE Heavy Syllable Nasal Apocope & \emph{*wī́θæg} \\
OE Velar Palatalization & \emph{*wī́θæʤ} \\
OE Unstressed AE Merger & \emph{*wī́θeʤ} \\
OE Late Unstressed Ag Suffix & \emph{*wī́θiʤ} \\
\end{tabular}
\end{minipage}
\\
\end{tabularx}
\end{minipage}%
} \endgroup

Old English form: \emph{wīþiġ}

\subsection{Reconstruction and comparative
evidence}\label{reconstruction-and-comparative-evidence-105}

The comparative evidence groups the word with Germanic forms of the
\Recon{wīþja} `withy' (or \emph{*wīþjō-} variant) or \emph{*wiþ-} type
(\citeproc{ref-Orel2003}{Orel 2003, 503}). These forms establish the
cognate set but do not by themselves explain the Old English suffix of
\emph{wīþiġ} `withy'.

For Old English, the relevant point is the suffix history. Campbell's
account of OE \emph{-ig}, including forms such as \emph{hunig}, supports
an analysis in which the \emph{-iġ} of \emph{wīþiġ} `withy' continues a
derivational \emph{*-ag-} sequence rather than a heavy ja-stem
\emph{*-ij-} (\citeproc{ref-Campbell1959}{Campbell 1959, sects 275,
376}). The derivational input {\Recon{wī́θagą} `withy'} is thus a
formation choice rather than a mere respelling of the comparative
headword.

\subsection{Old English evidence}\label{old-english-evidence-105}

Clark Hall records the noun as {\emph{wiðig}} `withy'
(\citeproc{ref-ClarkHall1960}{Clark Hall 1960, 358}). The form used
here, {\emph{wīþiġ}} `withy', is a normalized Old English spelling with
macrons and palatal ġ marked explicitly.

The relevant comparison form is therefore not a reconstructed dictionary
convenience but an established Old English noun. What requires
explanation is why the selected Proto-Germanic input is {\Recon{wī́θagą}
`withy'} rather than a comparative headword of the \Recon{wīþja} `withy'
type.

\subsection{Development to Old
English}\label{development-to-old-english-105}

From {\Recon{wī́θagą} `withy'}, Anglo-Frisian brightening gives a fronted
vowel in the suffixal syllable, and, on the Campbell analysis adopted
here, the later Old English development of \emph{*-ag-} yields
\emph{-iġ} (\citeproc{ref-Campbell1959}{Campbell 1959, sects 275, 376}).
Palatalization supplies the final \emph{ġ}, and the full development
reaches {\emph{wīþiġ}} `withy'.

This derivation is regular for the form compared here. The central claim
of the entry is therefore morphological: Old English \emph{wīþiġ}
`withy' belongs with an \emph{*-ag-} derivative, whereas the comparative
\Recon{wīþja} `withy' label belongs to a different way of presenting the
cognate family.

\subsection{Formation comparison}\label{formation-comparison-8}

The comparison below sets the relevant formations side by side. It
distinguishes the comparative headword from the Old English-facing
formation that actually yields the attested noun.

{\def\LTcaptype{none} % do not increment counter
\begin{longtable}[]{@{}
  >{\raggedright\arraybackslash}p{(\linewidth - 8\tabcolsep) * \real{0.2000}}
  >{\raggedright\arraybackslash}p{(\linewidth - 8\tabcolsep) * \real{0.2000}}
  >{\raggedright\arraybackslash}p{(\linewidth - 8\tabcolsep) * \real{0.2000}}
  >{\raggedright\arraybackslash}p{(\linewidth - 8\tabcolsep) * \real{0.2000}}
  >{\raggedright\arraybackslash}p{(\linewidth - 8\tabcolsep) * \real{0.2000}}@{}}
\toprule\noalign{}
\begin{minipage}[b]{\linewidth}\raggedright
Formation
\end{minipage} & \begin{minipage}[b]{\linewidth}\raggedright
Candidate input
\end{minipage} & \begin{minipage}[b]{\linewidth}\raggedright
Expected or documented OE outcome
\end{minipage} & \begin{minipage}[b]{\linewidth}\raggedright
OE comparison form
\end{minipage} & \begin{minipage}[b]{\linewidth}\raggedright
Result
\end{minipage} \\
\midrule\noalign{}
\endhead
\bottomrule\noalign{}
\endlastfoot
comparative family label & {\emph{*wáiθiz}} & broader cognate-set
headword & OE family context & useful lexeme label, but not the direct
source of \emph{wīþiġ} \\
heavy ja-stem analysis & \Recon{wīþja} `withy' type &
Campbell/Adamczyk-style heavy ja-stem \emph{-e} / zero outcome &
{\emph{wīþiġ}} `withy' & does not account cleanly for the OE suffix \\
\emph{*-ag-} derivative followed here & {\emph{*wī́θagą}} & regular
output: {\emph{wīþiġ}} `withy' & {\emph{wīþiġ}} `withy' & exact match
between formation and target \\
\end{longtable}
}

\section{\texorpdfstring{world --- OE
\emph{weorold}}{world --- OE weorold}}\label{world-oe-weorold}

\index[iv]{02oe@\textbf{Old English}!weorold@\emph{weorold}}
\index[iv]{03pgmc@\textbf{Proto-Germanic}!wiraaldiz@*wíra-àldiz}
\index[iv]{03pgmc@\textbf{Proto-Germanic}!wiraldu@*wír-àldu}
\index[iv]{03pgmc@\textbf{Proto-Germanic}!wiraldu@*wíràldu}

Derivation: citation reconstruction \emph{*wíra-àldiz}; form followed
here \emph{*wír-àldu} \textgreater{} \emph{weorold} (early analogy).

\subsection{Derivation trace}\label{derivation-trace-106}

Proto input: \emph{*wíràldu}

\begingroup
\setlength{\fboxsep}{6pt}

\noindent\fbox{%
\begin{minipage}{0.97\linewidth}
\small
\begin{tabularx}{\linewidth}{@{}>{\raggedright\arraybackslash}p{0.440\linewidth}>{\centering\arraybackslash}X>{\raggedright\arraybackslash}p{0.440\linewidth}@{}}
\begin{minipage}[t]{\linewidth}
\centering\textbf{Earlier Germanic changes}\par
\vspace{0.35em}
\raggedright
\centering\textbf{}\par
\raggedright
\vspace{0.2em}
\begin{tabular}{@{}>{\raggedright\arraybackslash}p{0.68\linewidth}@{\hspace{0.55em}}>{\raggedright\arraybackslash}p{0.20\linewidth}@{\hspace{0.25em}}}
\multicolumn{2}{@{}l@{}}{*Proto-West Germanic*} \\
\multicolumn{2}{@{}l@{}}{[no change]} \\
\multicolumn{2}{@{}l@{}}{*Northwest Germanic*} \\
NWGmc I Lowering & \emph{*wéràldu} \\
\end{tabular}
\end{minipage}
&
&
\begin{minipage}[t]{\linewidth}
\centering\textbf{Old English changes}\par
\vspace{0.35em}
\raggedright
\centering\textbf{}\par
\raggedright
\vspace{0.2em}
\begin{tabular}{@{}>{\raggedright\arraybackslash}p{0.68\linewidth}@{\hspace{0.55em}}>{\raggedright\arraybackslash}p{0.22\linewidth}@{\hspace{0.25em}}}
\multicolumn{2}{@{}l@{}}{*Old English*} \\
OE Inter Stress Raising & \emph{*wéruldu} \\
OE Med Unstressed U Lowering & \emph{*wéroldu} \\
OE Back Mutation & \emph{*wéoroldu} \\
OE High Vowel Apocope & \emph{*wéorold} \\
\end{tabular}
\end{minipage}
\\
\end{tabularx}
\end{minipage}%
} \endgroup

Old English form: \emph{weorold}

\subsection{Reconstruction and comparative
evidence}\label{reconstruction-and-comparative-evidence-106}

The word is the old compound `age of men'. Orel and the \emph{*wira-}
tradition reconstruct the older \emph{i}-vocalism, while Ringe and
Taylor discuss the lowered form \Recon{weraldiz} `world' and its
pre-Old-English chain \Recon{weraldu} `world' \textgreater{}
\Recon{weruld} `world' (\citeproc{ref-Orel2003}{Orel 2003, 501};
\citeproc{ref-RingeTaylor2014}{Ringe and Taylor 2014, 341}).
Kluge-Seebold likewise gives the compound \Recon{wira-aldō} `world' and
explicitly includes Old English \emph{weorold} `world'
(\citeproc{ref-KlugeSeebold2011}{Kluge 2011, 981}).

The derivational input {\Recon{wír-àldu} `world'} therefore combines the
older \emph{*wir-} vowel of the comparative headword with the early
shift of the compound into the ō-stems that Ringe and Taylor note for
this lexeme (\citeproc{ref-RingeTaylor2014}{Ringe and Taylor 2014,
341}). The early analogical step lies in that stem-class reassignment;
the later phonological developments can then run regularly.

\subsection{Old English evidence}\label{old-english-evidence-106}

Old English does not preserve a single isolated form. Ringe and Taylor
give West Saxon {\emph{weorold}} `world' \textasciitilde{}
{\emph{worold}} `world', Mercian {\emph{weoruld}} `world', Northumbrian
{\emph{woruld}} `world', and Kentish {\emph{wiarald}} `world'
(\citeproc{ref-RingeTaylor2014}{Ringe and Taylor 2014, 341}).
Sievers-Brunner and Bright present the same wider set, including the
syncopated {\emph{world}} `world' and later rounded {\emph{wurold}}
`world' (\citeproc{ref-SieversBrunner1965}{Brunner 1965, sect. 113};
\citeproc{ref-BrightCassidyRingler1971}{Bright 1971, 465}).

The Old English form used here is the West Saxon form {\emph{weorold}}
`world'. It is an attested Old English form within that broader variant
cluster, not the only form the lexeme ever shows.

\subsection{Development to Old
English}\label{development-to-old-english-106}

From the derivational input \Recon{wír-àldu} `world', Northwest Germanic
\emph{i}-lowering gives \Recon{wér-àldu} `world'. Inter-stress raising
then changes the medial \emph{a} to \emph{u}, producing \Recon{wér-uldu}
`world'. In the Old English branch that unstressed \emph{u} lowers to
\emph{o}, and back mutation yields \Recon{wéor-oldu} `world'; final
high-vowel apocope then gives \emph{weorold} `world'.

This sequence matches the comparative background in Ringe and Taylor's
\Recon{weraldiz} `world' \textgreater{} \Recon{weraldu} `world'
\textgreater{} \Recon{weruld} `world' chain while preserving the
\emph{*wir-} notation of the comparative label
(\citeproc{ref-RingeTaylor2014}{Ringe and Taylor 2014, 341}). The
modeled Old English form therefore stands at the meeting point of an
early stem-class reshaping and later regular sound change.

\subsection{Stage comparison}\label{stage-comparison-3}

The comparison below sets the relevant forms side by side. It separates
the comparative headword from the OE-facing stage chosen for the
derivation.

{\def\LTcaptype{none} % do not increment counter
\begin{longtable}[]{@{}
  >{\raggedright\arraybackslash}p{(\linewidth - 6\tabcolsep) * \real{0.2500}}
  >{\raggedright\arraybackslash}p{(\linewidth - 6\tabcolsep) * \real{0.2500}}
  >{\raggedright\arraybackslash}p{(\linewidth - 6\tabcolsep) * \real{0.2500}}
  >{\raggedright\arraybackslash}p{(\linewidth - 6\tabcolsep) * \real{0.2500}}@{}}
\toprule\noalign{}
\begin{minipage}[b]{\linewidth}\raggedright
Stage / interpretation
\end{minipage} & \begin{minipage}[b]{\linewidth}\raggedright
Candidate form
\end{minipage} & \begin{minipage}[b]{\linewidth}\raggedright
Old English outcome or comparison
\end{minipage} & \begin{minipage}[b]{\linewidth}\raggedright
Relevance to this entry
\end{minipage} \\
\midrule\noalign{}
\endhead
\bottomrule\noalign{}
\endlastfoot
comparative compound with older first-element vowel &
{\emph{*wíra-àldiz}} & citation reconstruction / lexeme label &
preserves the older \emph{*wir-} tradition of the compound \\
literature-stage lowered compound after early stem-class shift &
{\Recon{weraldiz}} \textgreater{} {\Recon{weraldu}} \textgreater{}
{\Recon{weruld}} & Ringe-Taylor background chain to OE {\emph{weorold}}
`world' \textasciitilde{} {\emph{worold}} `world' & explains the older
comparative literature cited for the word \\
Old English-facing input & {\emph{*wír-àldu}} & regular output:
{\emph{weorold}} `world' & exact match for the West Saxon form used
here \\
broader OE variant cluster & --- & {\emph{worold}} `world',
{\emph{weoruld}} `world', {\emph{woruld}} `world', {\emph{wiarald}}
`world', {\emph{world}} `world' & real attested comparanda that remain
outside that West Saxon line \\
\end{longtable}
}

\section{\texorpdfstring{youth --- OE
\emph{ġeoguþ}}{youth --- OE ġeoguþ}}\label{youth-oe-ux121eoguuxfe}

\index[iv]{02oe@\textbf{Old English}!geoguth@\emph{ġeoguþ}}
\index[iv]{03pgmc@\textbf{Proto-Germanic}!jugunth@*júgunθ}
\index[iv]{03pgmc@\textbf{Proto-Germanic}!jugunthiz@*júgunθiz}

Derivation: citation reconstruction \emph{*júgunθiz}; form followed here
\emph{*júgunθ} \textgreater{} \emph{ġeoguþ} (early analogy).

\subsection{Derivation trace}\label{derivation-trace-107}

Proto input: \emph{*júgunθ}

\begingroup
\setlength{\fboxsep}{6pt}

\noindent\fbox{%
\begin{minipage}{0.97\linewidth}
\small
\begin{tabularx}{\linewidth}{@{}>{\raggedright\arraybackslash}p{0.440\linewidth}>{\centering\arraybackslash}X>{\raggedright\arraybackslash}p{0.440\linewidth}@{}}
\begin{minipage}[t]{\linewidth}
\centering\textbf{Earlier Germanic changes}\par
\vspace{0.35em}
\raggedright
\centering\textbf{}\par
\raggedright
\vspace{0.2em}
\begin{tabular}{@{}>{\raggedright\arraybackslash}p{0.70\linewidth}@{\hspace{0.55em}}>{\raggedright\arraybackslash}p{0.20\linewidth}@{\hspace{0.25em}}}
\multicolumn{2}{@{}l@{}}{*Proto-West Germanic*} \\
\multicolumn{2}{@{}l@{}}{[no change]} \\
\multicolumn{2}{@{}l@{}}{*Northwest Germanic*} \\
OE Ws Palatal Glide & \emph{*jéugunθ} \\
NWGmc Nasal Spirant Lengthening & \emph{*jéugūnθ} \\
NWGmc Nasal Spirant Loss & \emph{*jéugūθ} \\
\end{tabular}
\end{minipage}
&
&
\begin{minipage}[t]{\linewidth}
\centering\textbf{Old English changes}\par
\vspace{0.35em}
\raggedright
\centering\textbf{}\par
\raggedright
\vspace{0.2em}
\begin{tabular}{@{}>{\raggedright\arraybackslash}p{0.74\linewidth}@{\hspace{0.55em}}>{\raggedright\arraybackslash}p{0.16\linewidth}@{\hspace{0.25em}}}
\multicolumn{2}{@{}l@{}}{*Old English*} \\
OE Diphthong Leveling & \emph{*jéogūθ} \\
OE Unstressed Long Vowel Shortening & \emph{*jéoguθ} \\
\end{tabular}
\end{minipage}
\\
\end{tabularx}
\end{minipage}%
} \endgroup

Old English form: \emph{ġeoguþ}

\subsection{Reconstruction and comparative
evidence}\label{reconstruction-and-comparative-evidence-107}

The wider etymological tradition reconstructs an earlier form of the
word as {\emph{*ju(w)unþi-}}. The comparative label {\Recon{júgunθiz}
`youth'} already stands at a later Germanic stage with \emph{g}, and the
derivational input {\Recon{júgunθ} `youth'} is later again: it
represents the form after final \emph{-i} has been lost.

Ringe and Taylor explicitly give the sequence \Recon{jugunþi} `youth'
\textgreater{} \Recon{juguþ} `youth' \textgreater{} OE geoguþ
\textasciitilde{} iuguþ (\citeproc{ref-RingeTaylor2014}{Ringe and Taylor
2014, 141}). The derivational input therefore differs from the broader
comparative headword because the Old English development must begin
after early loss of final \emph{-i}.

\subsection{Old English evidence}\label{old-english-evidence-107}

The Old English noun is attested with varying spellings. Ringe and
Taylor cite {\emph{geoguþ}} `youth' \textasciitilde{} {\emph{iuguþ}}
`youth' (\citeproc{ref-RingeTaylor2014}{Ringe and Taylor 2014, 141}).
The form is normalized here as {\emph{ġeoguþ}} `youth': the initial
palatal is written with \emph{ġ}, and the attested spelling variation is
treated as orthographic rather than lexical.

Nothing in the source stack suggests that a different paradigm cell
should be chosen. The relevant Old English comparison form is the noun
{\emph{ġeoguþ}} `youth' itself.

\subsection{Development to Old
English}\label{development-to-old-english-107}

The decisive early step is the loss of final \emph{-i} before the Old
English umlaut stage. If that high vowel remained, the word would
develop an over-umlauted \emph{y}-type vowel instead of the attested
form (\citeproc{ref-RingeTaylor2014}{Ringe and Taylor 2014, 141}).

From the derivational input {\Recon{júgunθ} `youth'}, the later
development is regular: palatal fronting yields \Recon{jéugunθ} `youth';
nasal-spirant lengthening and loss give \Recon{jéogūθ} `youth'
(\citeproc{ref-Fulk2018}{Fulk 2018, 109}); unstressed long-vowel
shortening then produces \Recon{jéoguθ} `youth', which surfaces as
{\emph{ġeoguþ}} `youth'. Campbell preserves \emph{u} after accented
\emph{u} in forms such as {\emph{duguþ}} `troop' and {\emph{munuc}}
`monk' (\citeproc{ref-Campbell1959}{Campbell 1959, sect. 374}). Brunner
likewise cites {\emph{iuzuð}} `youth' and {\emph{munuc}} `monk' under
the same development (\citeproc{ref-SieversBrunner1965}{Brunner 1965,
sect. 150.3}).

\subsection{Stage comparison}\label{stage-comparison-4}

The comparison below sets the relevant forms side by side. It separates
the broader comparative headword from the later stages relevant to the
Old English noun.

{\def\LTcaptype{none} % do not increment counter
\begin{longtable}[]{@{}
  >{\raggedright\arraybackslash}p{(\linewidth - 6\tabcolsep) * \real{0.2500}}
  >{\raggedright\arraybackslash}p{(\linewidth - 6\tabcolsep) * \real{0.2500}}
  >{\raggedright\arraybackslash}p{(\linewidth - 6\tabcolsep) * \real{0.2500}}
  >{\raggedright\arraybackslash}p{(\linewidth - 6\tabcolsep) * \real{0.2500}}@{}}
\toprule\noalign{}
\begin{minipage}[b]{\linewidth}\raggedright
Stage / interpretation
\end{minipage} & \begin{minipage}[b]{\linewidth}\raggedright
Candidate form
\end{minipage} & \begin{minipage}[b]{\linewidth}\raggedright
Old English outcome or comparison
\end{minipage} & \begin{minipage}[b]{\linewidth}\raggedright
Relevance to this entry
\end{minipage} \\
\midrule\noalign{}
\endhead
\bottomrule\noalign{}
\endlastfoot
earlier etymological headword & {\emph{*ju(w)unþi-}} & comparative
family background & older comparative reconstruction of the lexeme \\
later g-bearing comparative label & {\emph{*júgunθiz}} & citation
reconstruction / lexeme label & preserves the later Germanic stage
behind the selected entry \\
Old English-facing input & {\emph{*júgunθ}} & regular output:
{\emph{ġeoguþ}} & exact match for the Old English form used here \\
full \emph{-i} stage retained too long & {\emph{*jugunþi}} & expected
over-umlauted \emph{y}-type result & negative control showing why early
\emph{-i} loss must precede the OE umlaut stage \\
\end{longtable}
}

\clearpage

\chapter{Late analogy and paradigm-cell
selection}\label{late-analogy-and-paradigm-cell-selection}

These Old English forms continue a particular paradigm cell or a later
analogical remodeling, not the unaltered citation form. The phonology
may be regular once the proper morphological history has been
identified.

\section{\texorpdfstring{ban --- OE
\emph{bannes}}{ban --- OE bannes}}\label{ban-oe-bannes}

\index[iv]{02oe@\textbf{Old English}!bannes@\emph{bannes}}
\index[iv]{03pgmc@\textbf{Proto-Germanic}!banna@*bánną}
\index[iv]{03pgmc@\textbf{Proto-Germanic}!bannas@*bánnas}

Derivation: citation reconstruction \emph{*bánną}; form followed here
\emph{*bánnas} \textgreater{} \emph{bannes} (late analogy).

\subsection{Derivation trace}\label{derivation-trace-108}

Proto input: \emph{*bánnas}

\begingroup
\setlength{\fboxsep}{6pt}

\noindent\fbox{%
\begin{minipage}{0.97\linewidth}
\small
\begin{tabularx}{\linewidth}{@{}>{\raggedright\arraybackslash}p{0.440\linewidth}>{\centering\arraybackslash}X>{\raggedright\arraybackslash}p{0.440\linewidth}@{}}
\begin{minipage}[t]{\linewidth}
\centering\textbf{Earlier Germanic changes}\par
\vspace{0.35em}
\raggedright
\centering\textbf{}\par
\raggedright
\vspace{0.2em}
\begin{tabular}{@{}>{\raggedright\arraybackslash}p{0.68\linewidth}@{\hspace{0.55em}}>{\raggedright\arraybackslash}p{0.16\linewidth}@{\hspace{0.25em}}}
\multicolumn{2}{@{}l@{}}{*Proto-West Germanic*} \\
\multicolumn{2}{@{}l@{}}{[no change]} \\
\multicolumn{2}{@{}l@{}}{*Northwest Germanic*} \\
\multicolumn{2}{@{}l@{}}{[no change]} \\
\end{tabular}
\end{minipage}
&
&
\begin{minipage}[t]{\linewidth}
\centering\textbf{Old English changes}\par
\vspace{0.35em}
\raggedright
\centering\textbf{}\par
\raggedright
\vspace{0.2em}
\begin{tabular}{@{}>{\raggedright\arraybackslash}p{0.70\linewidth}@{\hspace{0.55em}}>{\raggedright\arraybackslash}p{0.16\linewidth}@{\hspace{0.25em}}}
\multicolumn{2}{@{}l@{}}{*Old English*} \\
Anglo Frisian Brightening & \emph{*bánnæs} \\
OE Unstressed AE Merger & \emph{*bánnes} \\
\end{tabular}
\end{minipage}
\\
\end{tabularx}
\end{minipage}%
} \endgroup

Old English form: \emph{bannes}

\subsection{Reconstruction and comparative
evidence}\label{reconstruction-and-comparative-evidence-108}

Orel cites a bann-noun under {\Recon{bannan} `ban'}, while Seebold
distinguishes bann-stems of both masculine and neuter type and gives Old
English {\emph{gebann}} `ban' as the noun reflex
(\citeproc{ref-Orel2003}{Orel 2003, 35};
\citeproc{ref-Seebold1970}{Seebold 1970, 89}). The citation
reconstruction {\Recon{bánną} `ban'} names the lexeme, but the
comparison here turns on the genitive singular {\Recon{bánnas} `ban'}.

The analysis therefore depends on medial, not final, gemination.

\subsection{Old English evidence}\label{old-english-evidence-108}

Old English lexicographic evidence securely supports the noun itself.
Bosworth-Toller records the noun under nominative {\emph{ge-bann}}
`ban', with oblique usage such as {\emph{gebanne}} `ban'
(\citeproc{ref-BosworthToller1898}{Bosworth and Toller 1898, 303}). The
exact unprefixed genitive {\emph{bannes}} `ban' is less directly cited
in the dictionaries, so it is best treated here as the regular genitive
form used for comparison rather than as a dictionary headword.

\subsection{Development to Old
English}\label{development-to-old-english-108}

From {\Recon{bánnas} `ban'}, the geminate remains medial before the case
ending and the unstressed vowel develops regularly to give
{\emph{bannes}} `ban'. The paradigm comparison therefore sets the
genitive against nominative {\emph{ban}} `ban', the ordinary nominative
form of the same noun, rather than against a directly cited genitive
headword.

\subsection{Paradigm comparison}\label{paradigm-comparison}

The comparison below sets the relevant forms side by side. It shows why
the genitive singular is the conservative cell used for the entry.

{\def\LTcaptype{none} % do not increment counter
\begin{longtable}[]{@{}
  >{\raggedright\arraybackslash}p{(\linewidth - 8\tabcolsep) * \real{0.2000}}
  >{\raggedright\arraybackslash}p{(\linewidth - 8\tabcolsep) * \real{0.2000}}
  >{\raggedright\arraybackslash}p{(\linewidth - 8\tabcolsep) * \real{0.2000}}
  >{\raggedright\arraybackslash}p{(\linewidth - 8\tabcolsep) * \real{0.2000}}
  >{\raggedright\arraybackslash}p{(\linewidth - 8\tabcolsep) * \real{0.2000}}@{}}
\toprule\noalign{}
\begin{minipage}[b]{\linewidth}\raggedright
PGmc cell / interpretation
\end{minipage} & \begin{minipage}[b]{\linewidth}\raggedright
Candidate input
\end{minipage} & \begin{minipage}[b]{\linewidth}\raggedright
OE output or comparison
\end{minipage} & \begin{minipage}[b]{\linewidth}\raggedright
OE comparison form
\end{minipage} & \begin{minipage}[b]{\linewidth}\raggedright
Result
\end{minipage} \\
\midrule\noalign{}
\endhead
\bottomrule\noalign{}
\endlastfoot
citation nominative singular & {\emph{*bánną}} & regular output:
{\emph{ban}} & ban & regular nominative outcome, but not the Old English
form here \\
genitive singular & {\emph{*bánnas}} & regular output: \emph{bannes} &
{\emph{bannes}} & direct match for the conservative genitive \\
\end{longtable}
}

\section{\texorpdfstring{berry --- OE
\emph{berġes}}{berry --- OE berġes}}\label{berry-oe-berux121es}

\index[iv]{02oe@\textbf{Old English}!berges@\emph{berġes}}
\index[iv]{03pgmc@\textbf{Proto-Germanic}!bazja@*bázją}
\index[iv]{03pgmc@\textbf{Proto-Germanic}!bazjas@*bázjas}

Derivation: citation reconstruction \emph{*bázją}; form followed here
\emph{*bázjas} \textgreater{} \emph{berġes} (late analogy).

\subsection{Derivation trace}\label{derivation-trace-109}

Proto input: \emph{*bázjas}

\begingroup
\setlength{\fboxsep}{6pt}

\noindent\fbox{%
\begin{minipage}{0.97\linewidth}
\small
\begin{tabularx}{\linewidth}{@{}>{\raggedright\arraybackslash}p{0.320\linewidth}>{\centering\arraybackslash}X>{\raggedright\arraybackslash}p{0.520\linewidth}@{}}
\begin{minipage}[t]{\linewidth}
\centering\textbf{Earlier Germanic changes}\par
\vspace{0.35em}
\raggedright
\centering\textbf{}\par
\raggedright
\vspace{0.2em}
\begin{tabular}{@{}>{\raggedright\arraybackslash}p{0.68\linewidth}@{\hspace{0.55em}}>{\raggedright\arraybackslash}p{0.16\linewidth}@{\hspace{0.25em}}}
\multicolumn{2}{@{}l@{}}{*Proto-West Germanic*} \\
\multicolumn{2}{@{}l@{}}{[no change]} \\
\multicolumn{2}{@{}l@{}}{*Northwest Germanic*} \\
\multicolumn{2}{@{}l@{}}{[no change]} \\
\end{tabular}
\end{minipage}
&
&
\begin{minipage}[t]{\linewidth}
\centering\textbf{Old English changes}\par
\vspace{0.35em}
\raggedright
\centering\textbf{}\par
\raggedright
\vspace{0.2em}
\begin{tabular}{@{}>{\raggedright\arraybackslash}p{0.70\linewidth}@{\hspace{0.55em}}>{\raggedright\arraybackslash}p{0.16\linewidth}@{\hspace{0.25em}}}
\multicolumn{2}{@{}l@{}}{*Old English*} \\
Anglo Frisian Brightening & \emph{*bærjæs} \\
OE I Umlaut & \emph{*berjæs} \\
OE Unstressed AE Merger & \emph{*berjes} \\
\end{tabular}
\end{minipage}
\\
\end{tabularx}
\end{minipage}%
} \endgroup

Old English form: \emph{berġes}

\subsection{Reconstruction and comparative
evidence}\label{reconstruction-and-comparative-evidence-109}

Kroonen reconstructs the berry noun as {\emph{*basja-}}
\textasciitilde{} {\emph{*bazja-}} (\citeproc{ref-Kroonen2013}{Kroonen
2013, 54}). The derivational input {\Recon{bázjas} `berry'} is therefore
not a rival lexeme headword, but a specific genitive singular cell drawn
from that paradigm.

The relevant point is that \emph{*rj} did not geminate in Proto-West
Germanic. Ringe and Taylor's {\emph{here}} `army', {\emph{herges}}
`army's' comparison shows the same \emph{rj} environment in an Old
English paradigm without any hidden gemination repair
(\citeproc{ref-RingeTaylor2014}{Ringe and Taylor 2014, 181}).

\subsection{Old English evidence}\label{old-english-evidence-109}

Campbell cites feminine {\emph{berige}} `berry' and notes that
\emph{-j-} is retained after \emph{r} in this type
(\citeproc{ref-Campbell1959}{Campbell 1959, 250}). The reviewed evidence
therefore supports the citation form more directly than the exact
genitive {\emph{berġes}} `berry', which is best read here as the regular
genitive comparison form rather than as a dictionary headword.

\subsection{Development to Old
English}\label{development-to-old-english-109}

Citation {\Recon{bázją} `berry'} gives {\emph{bere}} `berry (variant)',
not the Old English form here. The genitive singular {\Recon{bázjas}
`berry'}, however, gives {\emph{berġes}} `berry', with medial
\emph{-rġ-} preserved in the same way that Ringe and Taylor cite
{\emph{herges}} `army's' beside {\emph{here}} `army'
(\citeproc{ref-RingeTaylor2014}{Ringe and Taylor 2014, 181}). This
points to paradigm choice rather than to an extra phonological rule.

\subsection{Paradigm comparison}\label{paradigm-comparison-1}

The comparison below sets the relevant forms side by side. It shows the
contrast between the citation form and the genitive singular cell.

{\def\LTcaptype{none} % do not increment counter
\begin{longtable}[]{@{}
  >{\raggedright\arraybackslash}p{(\linewidth - 8\tabcolsep) * \real{0.2000}}
  >{\raggedright\arraybackslash}p{(\linewidth - 8\tabcolsep) * \real{0.2000}}
  >{\raggedright\arraybackslash}p{(\linewidth - 8\tabcolsep) * \real{0.2000}}
  >{\raggedright\arraybackslash}p{(\linewidth - 8\tabcolsep) * \real{0.2000}}
  >{\raggedright\arraybackslash}p{(\linewidth - 8\tabcolsep) * \real{0.2000}}@{}}
\toprule\noalign{}
\begin{minipage}[b]{\linewidth}\raggedright
PGmc cell / interpretation
\end{minipage} & \begin{minipage}[b]{\linewidth}\raggedright
Candidate input
\end{minipage} & \begin{minipage}[b]{\linewidth}\raggedright
OE output or comparison
\end{minipage} & \begin{minipage}[b]{\linewidth}\raggedright
OE comparison form
\end{minipage} & \begin{minipage}[b]{\linewidth}\raggedright
Result
\end{minipage} \\
\midrule\noalign{}
\endhead
\bottomrule\noalign{}
\endlastfoot
citation nominative singular & {\emph{*bázją}} & regular output:
{\emph{bere}} & {\emph{berige}} / {\emph{berġe}} & useful citation-form
background, but not the Old English form here \\
genitive singular & {\emph{*bázjas}} & regular output: \emph{berġes} &
{\emph{berġes}} & exact match for the conservative cell \\
\end{longtable}
}

\section{\texorpdfstring{bow --- OE
\emph{bēag}}{bow --- OE bēag}}\label{bow-oe-bux113ag}

\index[iv]{02oe@\textbf{Old English}!beag@\emph{bēag}}
\index[iv]{03pgmc@\textbf{Proto-Germanic}!baug@*báug}
\index[iv]{03pgmc@\textbf{Proto-Germanic}!beugana@*béuganą}

Derivation: citation reconstruction \emph{*béuganą}; form followed here
\emph{*báug} \textgreater{} \emph{bēag} (late analogy).

\subsection{Derivation trace}\label{derivation-trace-110}

Proto input: \emph{*báug}

\begingroup
\setlength{\fboxsep}{6pt}

\noindent\fbox{%
\begin{minipage}{0.97\linewidth}
\small
\begin{tabularx}{\linewidth}{@{}>{\raggedright\arraybackslash}p{0.440\linewidth}>{\centering\arraybackslash}X>{\raggedright\arraybackslash}p{0.440\linewidth}@{}}
\begin{minipage}[t]{\linewidth}
\centering\textbf{Earlier Germanic changes}\par
\vspace{0.35em}
\raggedright
\centering\textbf{}\par
\raggedright
\vspace{0.2em}
\begin{tabular}{@{}>{\raggedright\arraybackslash}p{0.68\linewidth}@{\hspace{0.55em}}>{\raggedright\arraybackslash}p{0.16\linewidth}@{\hspace{0.25em}}}
\multicolumn{2}{@{}l@{}}{*Proto-West Germanic*} \\
\multicolumn{2}{@{}l@{}}{[no change]} \\
\multicolumn{2}{@{}l@{}}{*Northwest Germanic*} \\
\multicolumn{2}{@{}l@{}}{[no change]} \\
\end{tabular}
\end{minipage}
&
&
\begin{minipage}[t]{\linewidth}
\centering\textbf{Old English changes}\par
\vspace{0.35em}
\raggedright
\centering\textbf{}\par
\raggedright
\vspace{0.2em}
\begin{tabular}{@{}>{\raggedright\arraybackslash}p{0.68\linewidth}@{\hspace{0.55em}}>{\raggedright\arraybackslash}p{0.16\linewidth}@{\hspace{0.25em}}}
\multicolumn{2}{@{}l@{}}{*Old English*} \\
OE Au Fronting & \emph{*báeug} \\
OE Diphthong Leveling & \emph{*bēag} \\
\end{tabular}
\end{minipage}
\\
\end{tabularx}
\end{minipage}%
} \endgroup

Old English form: \emph{bēag}

\subsection{Reconstruction and comparative
evidence}\label{reconstruction-and-comparative-evidence-110}

The inherited verb belongs to the class-II strong-verb family
{\Recon{béuganą} `bow'} (\citeproc{ref-RingeTaylor2014}{Ringe and Taylor
2014, 55}). Within that paradigm, however, the infinitive and the
singular preterite continue different ablaut grades. The derivational
input {\Recon{báug} `bow'} is the singular preterite cell, whereas the
citation form {\Recon{béuganą}} is the infinitive.

Campbell's account of Old English class-II strong verbs treats the
singular preterite \emph{au} \textgreater{} \emph{ēa} development as
regular in this environment (\citeproc{ref-Campbell1959}{Campbell 1959,
53}). That is the phonological path relevant for {\emph{bēag}} `bowed',
whereas the analogical \emph{ū} of the present stem belongs to the
separate history behind {\emph{būgan}} `to bow'
(\citeproc{ref-RingeTaylor2014}{Ringe and Taylor 2014, 55}).

\subsection{Old English evidence}\label{old-english-evidence-110}

Bosworth-Toller and Clark Hall both record {\emph{bēag}} `bowed' as a
preterite form of {\emph{būgan}} `to bow'
(\citeproc{ref-BosworthToller1898}{Bosworth and Toller 1898, 122};
\citeproc{ref-ClarkHall1960}{Clark Hall 1960, 45}). The form discussed
here is therefore an attested Old English verbal form, not a
reconstructed substitute for the infinitive.

The ordinary dictionary headword remains {\emph{būgan}} `to bow', but
the relevant comparison form for this entry is the singular preterite
{\emph{bēag}} `bow, ring' `bowed'. That is the paradigm cell in which
the inherited \emph{*au} grade is preserved most directly.

\subsection{Development to Old
English}\label{development-to-old-english-110}

From {\Recon{báug} `bow'}, Anglo-Frisian fronting and the later leveling
of the diphthong produce {\emph{bēag}} `bow, ring'
(\citeproc{ref-Campbell1959}{Campbell 1959, 53}). No special analogical
repair is needed for that cell. The form is the regular Old English
outcome of the singular-preterite grade.

The analogical element in the wider lexeme belongs instead to the
present stem seen in {\emph{būgan}} `bow, bend'. The derivational input
differs from the citation form because the regular inherited pathway
survives more transparently in the preterite than in the infinitive.

\subsection{Paradigm comparison}\label{paradigm-comparison-2}

The comparison below sets the relevant forms side by side. It
distinguishes the regular singular preterite from the more familiar
infinitival citation form.

{\def\LTcaptype{none} % do not increment counter
\begin{longtable}[]{@{}
  >{\raggedright\arraybackslash}p{(\linewidth - 8\tabcolsep) * \real{0.2000}}
  >{\raggedright\arraybackslash}p{(\linewidth - 8\tabcolsep) * \real{0.2000}}
  >{\raggedright\arraybackslash}p{(\linewidth - 8\tabcolsep) * \real{0.2000}}
  >{\raggedright\arraybackslash}p{(\linewidth - 8\tabcolsep) * \real{0.2000}}
  >{\raggedright\arraybackslash}p{(\linewidth - 8\tabcolsep) * \real{0.2000}}@{}}
\toprule\noalign{}
\begin{minipage}[b]{\linewidth}\raggedright
PGmc cell / interpretation
\end{minipage} & \begin{minipage}[b]{\linewidth}\raggedright
Candidate input
\end{minipage} & \begin{minipage}[b]{\linewidth}\raggedright
Expected or documented OE outcome
\end{minipage} & \begin{minipage}[b]{\linewidth}\raggedright
OE comparison form
\end{minipage} & \begin{minipage}[b]{\linewidth}\raggedright
Result
\end{minipage} \\
\midrule\noalign{}
\endhead
\bottomrule\noalign{}
\endlastfoot
citation infinitive & {\emph{*béuganą}} & inherited present-stem history
behind {\emph{būgan}} & {\emph{būgan}} & establishes the lexeme, but not
the Old English form here \\
singular preterite & {\emph{*báug}} & regular output: {\emph{bēag}} &
{\emph{bēag}} & exact match between input, output, and attested cell \\
past participial branch & participial \emph{*bugan-} type & later
participial outcomes & bogen-type evidence & relevant to the paradigm,
but not the clearest match for this entry \\
\end{longtable}
}

The singular preterite is the relevant comparison form. It gives a
direct lautgesetzlich path to attested {\emph{bēag}} `bow, ring', while
the citation form {\emph{būgan}} `bow, bend' belongs to a paradigm whose
present stem has already undergone later leveling.

\section{\texorpdfstring{cow --- OE
\emph{cȳ}}{cow --- OE cȳ}}\label{cow-oe-cux233}

\index[iv]{02oe@\textbf{Old English}!cy@\emph{cȳ}}
\index[iv]{03pgmc@\textbf{Proto-Germanic}!koz@*kōz}
\index[iv]{03pgmc@\textbf{Proto-Germanic}!kui@*kūi}

Derivation: citation reconstruction \emph{*kōz}; form followed here
\emph{*kūi} \textgreater{} \emph{cȳ} (late analogy).

\subsection{Derivation trace}\label{derivation-trace-111}

Proto input: \emph{*kūi}

\begingroup
\setlength{\fboxsep}{6pt}

\noindent\fbox{%
\begin{minipage}{0.97\linewidth}
\small
\begin{tabularx}{\linewidth}{@{}>{\raggedright\arraybackslash}p{0.440\linewidth}>{\centering\arraybackslash}X>{\raggedright\arraybackslash}p{0.440\linewidth}@{}}
\begin{minipage}[t]{\linewidth}
\centering\textbf{Earlier Germanic changes}\par
\vspace{0.35em}
\raggedright
\centering\textbf{}\par
\raggedright
\vspace{0.2em}
\begin{tabular}{@{}>{\raggedright\arraybackslash}p{0.68\linewidth}@{\hspace{0.55em}}>{\raggedright\arraybackslash}p{0.16\linewidth}@{\hspace{0.25em}}}
\multicolumn{2}{@{}l@{}}{*Proto-West Germanic*} \\
\multicolumn{2}{@{}l@{}}{[no change]} \\
\multicolumn{2}{@{}l@{}}{*Northwest Germanic*} \\
\multicolumn{2}{@{}l@{}}{[no change]} \\
\end{tabular}
\end{minipage}
&
&
\begin{minipage}[t]{\linewidth}
\centering\textbf{Old English changes}\par
\vspace{0.35em}
\raggedright
\centering\textbf{}\par
\raggedright
\vspace{0.2em}
\begin{tabular}{@{}>{\raggedright\arraybackslash}p{0.74\linewidth}@{\hspace{0.55em}}>{\raggedright\arraybackslash}p{0.16\linewidth}@{\hspace{0.25em}}}
\multicolumn{2}{@{}l@{}}{*Old English*} \\
OE I Umlaut & \emph{*kȳi} \\
\mbox{OE High Vowel Apocope} & \emph{*kȳ} \\
\end{tabular}
\end{minipage}
\\
\end{tabularx}
\end{minipage}%
} \endgroup

Old English form: \emph{cȳ}

\subsection{Reconstruction and comparative
evidence}\label{reconstruction-and-comparative-evidence-111}

Kroonen reconstructs a root noun with the inherited alternation
\emph{*kō-} \textasciitilde{} \emph{*ku-}, explicitly nom.
\emph{\Recon{kōz}} `cow', obl. \emph{*kū-}
(\citeproc{ref-Kroonen2013}{Kroonen 2013, 299}). The citation form
therefore belongs to the nominative singular, whereas the derivational
input \Recon{kūi} `cow' belongs to the oblique stem.

Ringe and Taylor also posit a later PNWGmc nominative \Recon{kūaz} `cow'
\textgreater{} \Recon{kūz} `cow' behind Old English {\emph{cū}} `cow'
(\citeproc{ref-RingeTaylor2014}{Ringe and Taylor 2014, sect. 3.1.3}).
That nominative history accounts for the headword, whereas the form cȳ
`cow' depends on the oblique \emph{*kū-} stem and a following \emph{*i}.

\subsection{Old English evidence}\label{old-english-evidence-111}

Clark Hall lemmatizes the noun under {\emph{cū}} `cow' and records
gen.sg. {\emph{cū}} `cow' (or {\emph{cū(e)}}), {\emph{cȳ}} `cow',
{\emph{cūs}} `cow (gen.)', dat.sg. {\emph{cȳ}} `cow', nom.-acc.pl.
{\emph{cȳ}} `cow', and dat.pl. \emph{cūm} `cow'
(\citeproc{ref-ClarkHall1960}{Clark Hall 1960, 67}). Ringe and Taylor
likewise state that the root-noun {\emph{cū}} `cow' exhibits dat.sg.,
nom.-acc.pl. {\emph{cȳ}} `cow' \textless{} \Recon{cūi} `cow', dat.pl.
\emph{cūm} \textless{} \emph{*cūm}, and apparently gen.sg. {\emph{cā}}
`cow (gen.)' \textless{} \Recon{cūiz} `cow (gen.)'
(\citeproc{ref-RingeTaylor2014}{Ringe and Taylor 2014, sect. 6.6.1}).

The dictionary headword is therefore {\emph{cū}} `cow', but {\emph{cȳ}}
`cow' is an established Old English paradigm form rather than a
convenient reconstruction. It serves as dative singular and also as
nominative-accusative plural. The genitive singular is less stable, with
{\emph{cā}} `cow (gen.sg.)', {\emph{cȳ}} `cow', \emph{cū(e)}, and
{\emph{cūs}} `cow (gen.sg.)' all appearing in the local source record.

\subsection{Development to Old
English}\label{development-to-old-english-111}

\Recon{kūi} `cow' is the dative singular of the oblique \emph{*kū-}
stem. The following \emph{*i} triggers i-umlaut, so \emph{ū} becomes
\emph{ȳ}, and loss of the final high vowel leaves {\emph{cȳ}} `cow'. The
development is \Recon{kūi} `cow' \textgreater{} \Recon{kȳi} `cow'
\textgreater{} \Recon{kȳ} `cow' \textgreater{} cȳ.

This is the regular oblique-cell path recognized by Ringe and Taylor's
paradigm statement (\citeproc{ref-RingeTaylor2014}{Ringe and Taylor
2014, sect. 6.6.1}). It also explains why {\emph{cȳ}} `cow' is the
cleanest comparison form for the present entry, even though the ordinary
headword is {\emph{cū}} `cow'.

\subsection{Paradigm comparison}\label{paradigm-comparison-3}

A paradigm comparison identifies the Proto-Germanic inflectional cell
that corresponds to an established Old English paradigm form. The
comparison below sets the relevant forms side by side.

{\def\LTcaptype{none} % do not increment counter
\begin{longtable}[]{@{}
  >{\raggedright\arraybackslash}p{(\linewidth - 8\tabcolsep) * \real{0.2000}}
  >{\raggedright\arraybackslash}p{(\linewidth - 8\tabcolsep) * \real{0.2000}}
  >{\raggedright\arraybackslash}p{(\linewidth - 8\tabcolsep) * \real{0.2000}}
  >{\raggedright\arraybackslash}p{(\linewidth - 8\tabcolsep) * \real{0.2000}}
  >{\raggedright\arraybackslash}p{(\linewidth - 8\tabcolsep) * \real{0.2000}}@{}}
\toprule\noalign{}
\begin{minipage}[b]{\linewidth}\raggedright
PGmc cell / interpretation
\end{minipage} & \begin{minipage}[b]{\linewidth}\raggedright
Candidate input
\end{minipage} & \begin{minipage}[b]{\linewidth}\raggedright
OE output or comparison
\end{minipage} & \begin{minipage}[b]{\linewidth}\raggedright
OE comparison form
\end{minipage} & \begin{minipage}[b]{\linewidth}\raggedright
Result
\end{minipage} \\
\midrule\noalign{}
\endhead
\bottomrule\noalign{}
\endlastfoot
citation nominative singular & *kōz & OE headword {\emph{cū}} `cow'
belongs to the nominative history of the lexeme & \emph{cū} & useful
background, but not the chosen comparison for {\emph{cȳ}} `cow' \\
later generalized nominative & PNWGmc \emph{kūaz \textgreater{} }kūz &
inferred nominative {\emph{cū}} `cow' & \emph{cū} & explains the leveled
headword, not the oblique target \\
dative singular oblique & *kūi & regular output: {\emph{cȳ}} `cow' &
\emph{cȳ} & exact match between input, output, and paradigm cell \\
genitive singular oblique & *kūiz & Ringe-Taylor: apparently {\emph{cā}}
`cow (gen.sg.)'; Hall: \emph{cū(e)}, \emph{cȳ}, \emph{cūs} & gen.sg.
variable & too unstable to control the entry \\
\end{longtable}
}

The dative singular is the relevant comparison form. It gives a regular
path to attested {\emph{cȳ}} `cow', while the broader Old English
paradigm shows how far the oblique \emph{*kū-} grade spread beyond that
one cell.

\section{\texorpdfstring{find --- OE
\emph{fundene}}{find --- OE fundene}}\label{find-oe-fundene}

\index[iv]{02oe@\textbf{Old English}!fundene@\emph{fundene}}
\index[iv]{03pgmc@\textbf{Proto-Germanic}!finthana@*fínθaną}
\index[iv]{03pgmc@\textbf{Proto-Germanic}!fundano@*fúnðanǭ}

Derivation: citation reconstruction \emph{*fínθaną}; form followed here
\emph{*fúnðanǭ} \textgreater{} \emph{fundene} (late analogy).

\subsection{Derivation trace}\label{derivation-trace-112}

Proto input: \emph{*fúnðanǭ}

\begingroup
\setlength{\fboxsep}{6pt}

\noindent\fbox{%
\begin{minipage}{0.97\linewidth}
\small
\begin{tabularx}{\linewidth}{@{}>{\raggedright\arraybackslash}p{0.440\linewidth}>{\centering\arraybackslash}X>{\raggedright\arraybackslash}p{0.440\linewidth}@{}}
\begin{minipage}[t]{\linewidth}
\centering\textbf{Earlier Germanic changes}\par
\vspace{0.35em}
\raggedright
\centering\textbf{}\par
\raggedright
\vspace{0.2em}
\begin{tabular}{@{}>{\raggedright\arraybackslash}p{0.70\linewidth}@{\hspace{0.55em}}>{\raggedright\arraybackslash}p{0.20\linewidth}@{\hspace{0.25em}}}
\multicolumn{2}{@{}l@{}}{*Proto-West Germanic*} \\
PWGmc Dental Hardening & \emph{*fúndanǭ} \\
\multicolumn{2}{@{}l@{}}{*Northwest Germanic*} \\
\multicolumn{2}{@{}l@{}}{[no change]} \\
\end{tabular}
\end{minipage}
&
&
\begin{minipage}[t]{\linewidth}
\centering\textbf{Old English changes}\par
\vspace{0.35em}
\raggedright
\centering\textbf{}\par
\raggedright
\vspace{0.2em}
\begin{tabular}{@{}>{\raggedright\arraybackslash}p{0.70\linewidth}@{\hspace{0.55em}}>{\raggedright\arraybackslash}p{0.20\linewidth}@{\hspace{0.25em}}}
\multicolumn{2}{@{}l@{}}{*Old English*} \\
OE Unstressed Fronting Early & \emph{*fúndænǭ} \\
OE Unstressed Long Vowel Shortening & \emph{*fúndænæ} \\
OE Unstressed AE Merger & \emph{*fúndene} \\
\end{tabular}
\end{minipage}
\\
\end{tabularx}
\end{minipage}%
} \endgroup

Old English form: \emph{fundene}

\subsection{Reconstruction and comparative
evidence}\label{reconstruction-and-comparative-evidence-112}

The inherited verb is the strong verb {\Recon{fínθaną} `find'},
continued by Old English {\emph{findan}} `find'
(\citeproc{ref-RingeTaylor2014}{Ringe and Taylor 2014, 344}). The form
followed here, {\Recon{fúnðanǭ} `find'}, belongs to the past-participial
paradigm rather than to the infinitive. It represents an oblique
singular form of the participle.

The familiar dictionary form {\emph{funden}} `found' is not the form
compared here. The derivational input instead models an attested
participial form directly, rather than treating the infinitive or the
ordinary dictionary headword as primary. It therefore reaches
{\emph{fundene}} `found (p.p.)' in the form where the trace and the
attested evidence match directly.

\subsection{Old English evidence}\label{old-english-evidence-112}

Bosworth-Toller records {\emph{fundene}} `found (p.p.)' `found' under
{\emph{findan}} `find', citing the form in \emph{Beón þá herigeata swa
fundene} (\citeproc{ref-BosworthToller1898}{Bosworth and Toller 1898,
219}). Clark Hall likewise preserves the participial stem in forms such
as {\emph{funden}} `found' and \emph{tō-fundennes} `finding'
(\citeproc{ref-ClarkHall1960}{Clark Hall 1960, 124}).

The ordinary dictionary headword for the participle is {\emph{funden}}
`found', but the relevant comparison form for this entry is the attested
oblique {\emph{fundene}} `found (p.p.)' `found'. It is an Old English
form in its own right, not a merely convenient probe.

\subsection{Development to Old
English}\label{development-to-old-english-112}

From {\Recon{fúnðanǭ} `find'}, the participial oblique develops through
regular loss and weakening of the final ending, yielding
{\emph{fundene}} `found (p.p.)'. In that cell both the consonantism and
the medial vowel history remain regular.

The broader participial paradigm fixes the interpretation. The more
familiar nominative {\emph{funden}} `found (p.p.)' is the ordinary
dictionary form, whereas the oblique form {\emph{fundene}} `found
(p.p.)' is the attested form compared here.

\subsection{Paradigm comparison}\label{paradigm-comparison-4}

{\def\LTcaptype{none} % do not increment counter
\begin{longtable}[]{@{}
  >{\raggedright\arraybackslash}p{(\linewidth - 8\tabcolsep) * \real{0.2000}}
  >{\raggedright\arraybackslash}p{(\linewidth - 8\tabcolsep) * \real{0.2000}}
  >{\raggedright\arraybackslash}p{(\linewidth - 8\tabcolsep) * \real{0.2000}}
  >{\raggedright\arraybackslash}p{(\linewidth - 8\tabcolsep) * \real{0.2000}}
  >{\raggedright\arraybackslash}p{(\linewidth - 8\tabcolsep) * \real{0.2000}}@{}}
\toprule\noalign{}
\begin{minipage}[b]{\linewidth}\raggedright
PGmc cell / interpretation
\end{minipage} & \begin{minipage}[b]{\linewidth}\raggedright
Candidate input
\end{minipage} & \begin{minipage}[b]{\linewidth}\raggedright
Expected or documented OE outcome
\end{minipage} & \begin{minipage}[b]{\linewidth}\raggedright
OE comparison form
\end{minipage} & \begin{minipage}[b]{\linewidth}\raggedright
Result
\end{minipage} \\
\midrule\noalign{}
\endhead
\bottomrule\noalign{}
\endlastfoot
citation infinitive & {\emph{*fínθaną}} & inherited verb {\emph{findan}}
& {\emph{findan}} & establishes the lexeme, but not the form compared
here \\
nominative participial line & {\emph{*fúnðanaz}} & ordinary dictionary
{\emph{funden}} type & {\emph{funden}} & important paradigm background,
but not the form compared here \\
oblique participle compared here & {\emph{*fúnðanǭ}} & regular output:
{\emph{fundene}} & {\emph{fundene}} & exact match between input, output,
and attested cell \\
\end{longtable}
}

The oblique participle is the relevant comparison form. It matches the
derivational input and Old English form directly, while the nominative
participial headword remains a different presentation cell within the
same paradigm.

\section{\texorpdfstring{fright --- OE
\emph{fyrhte}}{fright --- OE fyrhte}}\label{fright-oe-fyrhte}

\index[iv]{02oe@\textbf{Old English}!fyrhte@\emph{fyrhte}}
\index[iv]{03pgmc@\textbf{Proto-Germanic}!furxtin@*furxtīn}
\index[iv]{03pgmc@\textbf{Proto-Germanic}!furxtinaz@*fúrxtīnaz}

Derivation: citation reconstruction \emph{*furxtīn}; form followed here
\emph{*fúrxtīnaz} \textgreater{} \emph{fyrhte} (late analogy).

\subsection{Derivation trace}\label{derivation-trace-113}

Proto input: \emph{*fúrxtīnaz}

\begingroup
\setlength{\fboxsep}{6pt}

\noindent\fbox{%
\begin{minipage}{0.97\linewidth}
\small
\begin{tabularx}{\linewidth}{@{}>{\raggedright\arraybackslash}p{0.460\linewidth}>{\centering\arraybackslash}X>{\raggedright\arraybackslash}p{0.460\linewidth}@{}}
\begin{minipage}[t]{\linewidth}
\centering\textbf{Earlier Germanic changes}\par
\vspace{0.35em}
\raggedright
\centering\textbf{}\par
\raggedright
\vspace{0.2em}
\begin{tabular}{@{}>{\raggedright\arraybackslash}p{0.68\linewidth}@{\hspace{0.55em}}>{\raggedright\arraybackslash}p{0.22\linewidth}@{\hspace{0.25em}}}
\multicolumn{2}{@{}l@{}}{*Proto-West Germanic*} \\
\multicolumn{2}{@{}l@{}}{[no change]} \\
\multicolumn{2}{@{}l@{}}{*Northwest Germanic*} \\
PGmc Final Z Deletion & \emph{*fúrxtīna} \\
\end{tabular}
\end{minipage}
&
&
\begin{minipage}[t]{\linewidth}
\centering\textbf{Old English changes}\par
\vspace{0.35em}
\raggedright
\centering\textbf{}\par
\raggedright
\vspace{0.2em}
\begin{tabular}{@{}>{\raggedright\arraybackslash}p{0.70\linewidth}@{\hspace{0.55em}}>{\raggedright\arraybackslash}p{0.20\linewidth}@{\hspace{0.25em}}}
\multicolumn{2}{@{}l@{}}{*Old English*} \\
PWGmc Final Bare A Loss & \emph{*fúrxtīn} \\
OE I Umlaut & \emph{*fyrxtīn} \\
NWGmc In Stem N Loss & \emph{*fyrxtī} \\
OE Unstressed Long Vowel Shortening & \emph{*fyrxti} \\
OE Med Unstressed I Lowering1 & \emph{*fyrxte} \\
\end{tabular}
\end{minipage}
\\
\end{tabularx}
\end{minipage}%
} \endgroup

Old English form: \emph{fyrhte}

\subsection{Reconstruction and comparative
evidence}\label{reconstruction-and-comparative-evidence-113}

The noun belongs to the inherited in-stem abstract {\Recon{furxtīn}
`fright'}, the same family as Gothic {\emph{faurhtei}} `fear, fright'
(\citeproc{ref-Orel2003}{Orel 2003, 120}). The derivational input
{\Recon{fúrxtīnaz} `fright'} is not a different lexeme but an oblique
singular cell within that in-stem paradigm.

Ringe and Taylor treat the later nominative forms with \emph{-u} or
\emph{-o} as analogically remodeled
(\citeproc{ref-RingeTaylor2014}{Ringe and Taylor 2014, 395--96}). I
therefore use the oblique in-stem form, which continues the older
formation rather than the remodeled nominative.

\subsection{Old English evidence}\label{old-english-evidence-113}

Bosworth-Toller records {\emph{fyrhte}} `fright' with textual
attestation, and it also records nominative forms such as
{\emph{fyrhtu}} `fright' and {\emph{fyrhto}} `fright'
(\citeproc{ref-BosworthToller1898}{Bosworth and Toller 1898, 160}).
Clark Hall lists adjective and verb material separately under
\emph{fyrht} / \emph{fyrhtan} (\citeproc{ref-ClarkHall1960}{Clark Hall
1960, 141}).

The relevant comparison form is therefore the attested oblique
{\emph{fyrhte}} `fright'. The nominative lemma forms remain part of the
Old English evidence, but the Old English form here of this entry is the
oblique cell.

\subsection{Development to Old
English}\label{development-to-old-english-113}

From {\Recon{fúrxtīnaz} `fright'}, the oblique in-stem develops through
the loss and weakening of the final ending, yielding {\emph{fyrhte}}
`fright'. The form compared here therefore follows the ordinary Old
English reduction of the abstract ending in this paradigm.

The later nominative forms with \emph{-u} or \emph{-o} belong to a
subsequent analogical reshaping of the paradigm. The Old English form
here is earlier in that sense: it is the attested OE cell in which the
inherited in-stem development remains most transparent.

\subsection{Paradigm comparison}\label{paradigm-comparison-5}

The comparison below sets the relevant forms side by side. It separates
the attested oblique form from the later remodeled nominative line.

{\def\LTcaptype{none} % do not increment counter
\begin{longtable}[]{@{}
  >{\raggedright\arraybackslash}p{(\linewidth - 8\tabcolsep) * \real{0.2000}}
  >{\raggedright\arraybackslash}p{(\linewidth - 8\tabcolsep) * \real{0.2000}}
  >{\raggedright\arraybackslash}p{(\linewidth - 8\tabcolsep) * \real{0.2000}}
  >{\raggedright\arraybackslash}p{(\linewidth - 8\tabcolsep) * \real{0.2000}}
  >{\raggedright\arraybackslash}p{(\linewidth - 8\tabcolsep) * \real{0.2000}}@{}}
\toprule\noalign{}
\begin{minipage}[b]{\linewidth}\raggedright
PGmc cell / interpretation
\end{minipage} & \begin{minipage}[b]{\linewidth}\raggedright
Candidate input
\end{minipage} & \begin{minipage}[b]{\linewidth}\raggedright
Expected or documented OE outcome
\end{minipage} & \begin{minipage}[b]{\linewidth}\raggedright
OE comparison form
\end{minipage} & \begin{minipage}[b]{\linewidth}\raggedright
Result
\end{minipage} \\
\midrule\noalign{}
\endhead
\bottomrule\noalign{}
\endlastfoot
citation in-stem headword & {\emph{*furxtīn}} & broader noun-class label
& wider family context & useful lexeme label, but not the cell compared
here \\
remodeled nominative line & nominative in-stem forms & {\emph{fyrhtu}}
`fright' / {\emph{fyrhto}} `fright' type lemma forms & {\emph{fyrhtu}}
`fright' / {\emph{fyrhto}} `fright' & genuine OE evidence, but later
remodeled \\
selected oblique singular & {\emph{*fúrxtīnaz}} & regular output:
{\emph{fyrhte}} `fright' & {\emph{fyrhte}} `fright' & exact match
between input, output, and attested cell \\
\end{longtable}
}

The oblique in-stem form is the relevant comparison form. It yields
attested {\emph{fyrhte}} `fright' directly, while the more familiar
nominative forms belong to a later analogical layer.

\section{\texorpdfstring{hammer --- OE
\emph{hameres}}{hammer --- OE hameres}}\label{hammer-oe-hameres}

\index[iv]{02oe@\textbf{Old English}!hameres@\emph{hameres}}
\index[iv]{03pgmc@\textbf{Proto-Germanic}!xamaras@*xámaras}
\index[iv]{03pgmc@\textbf{Proto-Germanic}!xamaraz@*xámaraz}

Derivation: citation reconstruction \emph{*xámaraz}; form followed here
\emph{*xámaras} \textgreater{} \emph{hameres} (late analogy).

\subsection{Derivation trace}\label{derivation-trace-114}

Proto input: \emph{*xámaras}

\begingroup
\setlength{\fboxsep}{6pt}

\noindent\fbox{%
\begin{minipage}{0.97\linewidth}
\small
\begin{tabularx}{\linewidth}{@{}>{\raggedright\arraybackslash}p{0.440\linewidth}>{\centering\arraybackslash}X>{\raggedright\arraybackslash}p{0.440\linewidth}@{}}
\begin{minipage}[t]{\linewidth}
\centering\textbf{Earlier Germanic changes}\par
\vspace{0.35em}
\raggedright
\centering\textbf{}\par
\raggedright
\vspace{0.2em}
\begin{tabular}{@{}>{\raggedright\arraybackslash}p{0.68\linewidth}@{\hspace{0.55em}}>{\raggedright\arraybackslash}p{0.16\linewidth}@{\hspace{0.25em}}}
\multicolumn{2}{@{}l@{}}{*Proto-West Germanic*} \\
\multicolumn{2}{@{}l@{}}{[no change]} \\
\multicolumn{2}{@{}l@{}}{*Northwest Germanic*} \\
\multicolumn{2}{@{}l@{}}{[no change]} \\
\end{tabular}
\end{minipage}
&
&
\begin{minipage}[t]{\linewidth}
\centering\textbf{Old English changes}\par
\vspace{0.35em}
\raggedright
\centering\textbf{}\par
\raggedright
\vspace{0.2em}
\begin{tabular}{@{}>{\raggedright\arraybackslash}p{0.70\linewidth}@{\hspace{0.55em}}>{\raggedright\arraybackslash}p{0.20\linewidth}@{\hspace{0.25em}}}
\multicolumn{2}{@{}l@{}}{*Old English*} \\
Anglo Frisian Brightening & \emph{*xámæræs} \\
OE Unstressed AE Merger & \emph{*xámeres} \\
\end{tabular}
\end{minipage}
\\
\end{tabularx}
\end{minipage}%
} \endgroup

Old English form: \emph{hameres}

\subsection{Reconstruction and comparative
evidence}\label{reconstruction-and-comparative-evidence-114}

The inherited noun is the masculine a-stem {\Recon{xámaraz} `hammer'},
reflected in Old English citation forms such as {\emph{hamor}} `hammer'
and {\emph{hamer}} `hammer' (\citeproc{ref-Kroonen2013}{Kroonen 2013,
206}; \citeproc{ref-Orel2003}{Orel 2003, 197};
\citeproc{ref-ClarkHall1960}{Clark Hall 1960, 160}). The derivational
input {\Recon{xámaras} `hammer'} is the genitive singular of that same
noun rather than a different lexeme.

The citation tradition is already mixed in its unstressed vowel, while
the genitive singular gives a closer comparison form. This is a cell
choice within one paradigm, not a change of stem class.

\subsection{Old English evidence}\label{old-english-evidence-114}

Bosworth-Toller directly records {\emph{hameres}} `hammer' `hammer' in
an Old English genitival phrase
(\citeproc{ref-BosworthToller1898}{Bosworth and Toller 1898, 78}). Clark
Hall preserves the simplex headword as {\emph{hamer}} `hammer' /
{\emph{hamor}} `hammer' (\citeproc{ref-ClarkHall1960}{Clark Hall 1960,
160}).

Sievers-Brunner gives a paradigm line {\emph{hamor}} `hammer' ---
{\emph{hamores}} `hammer', which shows that the oblique tradition itself
was not entirely uniform (\citeproc{ref-SieversBrunner1965}{Brunner
1965, sect. 245}). The relevant comparison form here is the attested
genitive singular {\emph{hameres}} `hammer'.

\subsection{Development to Old
English}\label{development-to-old-english-114}

From {\Recon{xámaras} `hammer'}, Anglo-Frisian brightening and the
subsequent merger of unstressed \emph{æ} with \emph{e} yield
{\emph{hameres}} `hammer'. The derivation of that oblique form is
straightforward once the genitive singular cell is selected.

The noun as a whole retains a mixed citation tradition in {\emph{hamor}}
`hammer' and {\emph{hamer}} `hammer', and the selected oblique cell
avoids making that variation carry the argument of the entry.

\subsection{Paradigm comparison}\label{paradigm-comparison-6}

The comparison below sets the relevant forms side by side. It separates
the attested genitive singular from the less stable citation tradition.

{\def\LTcaptype{none} % do not increment counter
\begin{longtable}[]{@{}
  >{\raggedright\arraybackslash}p{(\linewidth - 8\tabcolsep) * \real{0.2000}}
  >{\raggedright\arraybackslash}p{(\linewidth - 8\tabcolsep) * \real{0.2000}}
  >{\raggedright\arraybackslash}p{(\linewidth - 8\tabcolsep) * \real{0.2000}}
  >{\raggedright\arraybackslash}p{(\linewidth - 8\tabcolsep) * \real{0.2000}}
  >{\raggedright\arraybackslash}p{(\linewidth - 8\tabcolsep) * \real{0.2000}}@{}}
\toprule\noalign{}
\begin{minipage}[b]{\linewidth}\raggedright
PGmc cell / interpretation
\end{minipage} & \begin{minipage}[b]{\linewidth}\raggedright
Candidate input
\end{minipage} & \begin{minipage}[b]{\linewidth}\raggedright
Expected or documented OE outcome
\end{minipage} & \begin{minipage}[b]{\linewidth}\raggedright
OE comparison form
\end{minipage} & \begin{minipage}[b]{\linewidth}\raggedright
Result
\end{minipage} \\
\midrule\noalign{}
\endhead
\bottomrule\noalign{}
\endlastfoot
citation nominative singular & {\emph{*xámaraz}} & regular citation form
{\emph{hamer}} / {\emph{hamor}} & {\emph{hamor}} / {\emph{hamer}} & good
lexical background, but not the Old English form here \\
genitive singular & {\emph{*xámaras}} & regular output: {\emph{hameres}}
& {\emph{hameres}} & exact match between input, output, and attested
cell \\
later oblique tradition & oblique a-stem forms & {\emph{hamores}} type
evidence & {\emph{hamores}} & attested background variant, but not the
chosen comparison form \\
\end{longtable}
}

\section{\texorpdfstring{have --- OE
\emph{hæfeþ}}{have --- OE hæfeþ}}\label{have-oe-huxe6feuxfe}

\index[iv]{02oe@\textbf{Old English}!haefeth@\emph{hæfeþ}}
\index[iv]{03pgmc@\textbf{Proto-Germanic}!xabena@*xabēną}
\index[iv]{03pgmc@\textbf{Proto-Germanic}!xabethi@*xábēθi}

Derivation: citation reconstruction \emph{*xabēną}; form followed here
\emph{*xábēθi} \textgreater{} \emph{hæfeþ} (late analogy).

\subsection{Derivation trace}\label{derivation-trace-115}

Proto input: \emph{*xábēθi}

\begingroup
\setlength{\fboxsep}{6pt}

\noindent\fbox{%
\begin{minipage}{0.97\linewidth}
\small
\begin{tabularx}{\linewidth}{@{}>{\raggedright\arraybackslash}p{0.460\linewidth}>{\centering\arraybackslash}X>{\raggedright\arraybackslash}p{0.460\linewidth}@{}}
\begin{minipage}[t]{\linewidth}
\centering\textbf{Earlier Germanic changes}\par
\vspace{0.35em}
\raggedright
\centering\textbf{}\par
\raggedright
\vspace{0.2em}
\begin{tabular}{@{}>{\raggedright\arraybackslash}p{0.74\linewidth}@{\hspace{0.55em}}>{\raggedright\arraybackslash}p{0.16\linewidth}@{\hspace{0.25em}}}
\multicolumn{2}{@{}l@{}}{*Proto-West Germanic*} \\
PWGmc Early I Apocope & \emph{*xábēθ} \\
\multicolumn{2}{@{}l@{}}{*Northwest Germanic*} \\
\mbox{NWGmc Long E Lowering} & \emph{*xábǣθ} \\
\end{tabular}
\end{minipage}
&
&
\begin{minipage}[t]{\linewidth}
\centering\textbf{Old English changes}\par
\vspace{0.35em}
\raggedright
\centering\textbf{}\par
\raggedright
\vspace{0.2em}
\begin{tabular}{@{}>{\raggedright\arraybackslash}p{0.74\linewidth}@{\hspace{0.55em}}>{\raggedright\arraybackslash}p{0.16\linewidth}@{\hspace{0.25em}}}
\multicolumn{2}{@{}l@{}}{*Old English*} \\
Anglo Frisian Brightening & \emph{*xæbǣθ} \\
OE Velar Fricative Palatalization & \emph{*çæbǣθ} \\
PGmc B Allophony & \emph{*çæβǣθ} \\
OE Unstressed Long Vowel Shortening & \emph{*çæβæθ} \\
OE Unstressed AE Merger & \emph{*çæβeθ} \\
\end{tabular}
\end{minipage}
\\
\end{tabularx}
\end{minipage}%
} \endgroup

Old English form: \emph{hæfeþ}

\subsection{Reconstruction and comparative
evidence}\label{reconstruction-and-comparative-evidence-115}

The verb belongs to the inherited class-III weak paradigm usually cited
under {\Recon{xabēną} `have'} and Old English {\emph{habban}} `have'
(\citeproc{ref-Kroonen2013}{Kroonen 2013, 237};
\citeproc{ref-RingeTaylor2014}{Ringe and Taylor 2014, 93}). Within that
paradigm, however, the infinitive and the singular present indicative do
not continue the same stem. Ringe and Taylor distinguish a \emph{-ja-}
stem in the infinitive from a non-geminating -ai- / \emph{-ē-} stem in
the 2sg and 3sg present forms (\citeproc{ref-RingeTaylor2014}{Ringe and
Taylor 2014, 93}).

The derivational input {\Recon{xábēθi} `have'} is therefore the 3sg
present cell rather than a rephrasing of the infinitive. For the present
analysis, that finite cell is the cleaner comparator for the inherited
non-geminating stem.

\subsection{Old English evidence}\label{old-english-evidence-115}

The ordinary Old English headword is {\emph{habban}} `have'
(\citeproc{ref-ClarkHall1960}{Clark Hall 1960, 157}). Campbell's Anglian
paradigm includes unsyncopated 3sg forms of the {\emph{hæfed}} `has'
type, and the present paradigm therefore shows forms of the \emph{hæf-}
type that support the normalized target {\emph{hæfeþ}} `has'
(\citeproc{ref-Campbell1959}{Campbell 1959, sect. 762}).

The target form is therefore a normalized finite cell rather than the
ordinary dictionary lemma. It represents the inherited non-geminating
present stem more directly than {\emph{habban}} `have' does.

\subsection{Development to Old
English}\label{development-to-old-english-115}

From {\Recon{xábēθi} `have'}, the finite form yields {\emph{hæfeþ}}
`has' regularly. Ringe and Taylor discuss this non-geminating present
stem under {\emph{habban}} `have' (\citeproc{ref-RingeTaylor2014}{Ringe
and Taylor 2014, 364}). Campbell's Anglian paradigms include
unsyncopated 3sg forms of the {\emph{hæfeþ}} `has' / {\emph{hæfed}}
`has' type (\citeproc{ref-Campbell1959}{Campbell 1959, sect. 762}).

The wider lexeme is less straightforward only because the infinitive
{\emph{habban}} `have' shows later leveling. The 3sg present cell is
therefore the closer comparison form.

\subsection{Paradigm comparison}\label{paradigm-comparison-7}

The comparison below sets the relevant forms side by side. It separates
the analogically leveled citation form from the regular 3sg present
line.

{\def\LTcaptype{none} % do not increment counter
\begin{longtable}[]{@{}
  >{\raggedright\arraybackslash}p{(\linewidth - 8\tabcolsep) * \real{0.2000}}
  >{\raggedright\arraybackslash}p{(\linewidth - 8\tabcolsep) * \real{0.2000}}
  >{\raggedright\arraybackslash}p{(\linewidth - 8\tabcolsep) * \real{0.2000}}
  >{\raggedright\arraybackslash}p{(\linewidth - 8\tabcolsep) * \real{0.2000}}
  >{\raggedright\arraybackslash}p{(\linewidth - 8\tabcolsep) * \real{0.2000}}@{}}
\toprule\noalign{}
\begin{minipage}[b]{\linewidth}\raggedright
PGmc cell / interpretation
\end{minipage} & \begin{minipage}[b]{\linewidth}\raggedright
Candidate input
\end{minipage} & \begin{minipage}[b]{\linewidth}\raggedright
Expected or documented OE outcome
\end{minipage} & \begin{minipage}[b]{\linewidth}\raggedright
OE comparison form
\end{minipage} & \begin{minipage}[b]{\linewidth}\raggedright
Result
\end{minipage} \\
\midrule\noalign{}
\endhead
\bottomrule\noalign{}
\endlastfoot
citation infinitive & \emph{-ja-} stem of {\emph{*xabēną}} & citation
form {\emph{habban}} & {\emph{habban}} & important headword, but shaped
by later leveling \\
3sg present & {\emph{*xábēθi}} & regular output: {\emph{hæfeþ}} &
{\emph{hæfeþ}} & exact match between input, output, and finite form
compared here \\
syncopated finite tradition & same present stem & {\emph{hæfþ}} type
evidence & {\emph{hæfþ}} & genuine later OE finite form, but not the
normalized target used here \\
\end{longtable}
}

\section{\texorpdfstring{heaven --- OE
\emph{heofon}}{heaven --- OE heofon}}\label{heaven-oe-heofon}

\index[iv]{02oe@\textbf{Old English}!heofon@\emph{heofon}}
\index[iv]{03pgmc@\textbf{Proto-Germanic}!xemenaz@*xémenaz}
\index[iv]{03pgmc@\textbf{Proto-Germanic}!xemonu@*xémonų}

Derivation: citation reconstruction \emph{*xémenaz}; form followed here
\emph{*xémonų} \textgreater{} \emph{heofon} (late analogy).

\subsection{Derivation trace}\label{derivation-trace-116}

Proto input: \emph{*xémonų}

\begingroup
\setlength{\fboxsep}{6pt}

\noindent\fbox{%
\begin{minipage}{0.97\linewidth}
\small
\begin{tabularx}{\linewidth}{@{}>{\raggedright\arraybackslash}p{0.440\linewidth}>{\centering\arraybackslash}X>{\raggedright\arraybackslash}p{0.440\linewidth}@{}}
\begin{minipage}[t]{\linewidth}
\centering\textbf{Earlier Germanic changes}\par
\vspace{0.35em}
\raggedright
\centering\textbf{}\par
\raggedright
\vspace{0.2em}
\begin{tabular}{@{}>{\raggedright\arraybackslash}p{0.70\linewidth}@{\hspace{0.55em}}>{\raggedright\arraybackslash}p{0.16\linewidth}@{\hspace{0.25em}}}
\multicolumn{2}{@{}l@{}}{*Proto-West Germanic*} \\
\multicolumn{2}{@{}l@{}}{[no change]} \\
\multicolumn{2}{@{}l@{}}{*Northwest Germanic*} \\
NWGmc Unstressed O Raising & \emph{*xémunų} \\
NWGmc Mn Dissimilation & \emph{*xéβunų} \\
\end{tabular}
\end{minipage}
&
&
\begin{minipage}[t]{\linewidth}
\centering\textbf{Old English changes}\par
\vspace{0.35em}
\raggedright
\centering\textbf{}\par
\raggedright
\vspace{0.2em}
\begin{tabular}{@{}>{\raggedright\arraybackslash}p{0.70\linewidth}@{\hspace{0.55em}}>{\raggedright\arraybackslash}p{0.20\linewidth}@{\hspace{0.25em}}}
\multicolumn{2}{@{}l@{}}{*Old English*} \\
OE Med Unstressed U Lowering & \emph{*xéβonų} \\
OE Velar Fricative Palatalization & \emph{*çéβonų} \\
OE Back Mutation & \emph{*çéoβonų} \\
\mbox{OE High Vowel Apocope} & \emph{*çéoβon} \\
\end{tabular}
\end{minipage}
\\
\end{tabularx}
\end{minipage}%
} \endgroup

Old English form: \emph{heofon}

\subsection{Reconstruction and comparative
evidence}\label{reconstruction-and-comparative-evidence-116}

The inherited noun belongs to the mn-stem family cited by Kroonen as
{\Recon{hemina-*} `heaven'} \textasciitilde{} {\Recon{hemna-*} `heaven'}
(\citeproc{ref-Kroonen2013}{Kroonen 2013, 220}). The derivational input
{\Recon{xémonų} `heaven'} is an oblique singular form within that
paradigm rather than the lexeme-level citation form {\Recon{xémenaz}
`heaven'}.

The back-vocalic oblique stem accounts for the West Saxon target. Ringe
and Taylor give northern WGmc {\Recon{hebun} `heaven'} \textgreater{}
West Saxon and Northumbrian {\emph{heofon}} `heaven', Mercian
{\emph{heofen}} `heaven' (\citeproc{ref-RingeTaylor2014}{Ringe and
Taylor 2014, 324}). Campbell likewise gives {\emph{heofon}} `heaven'
beside {\emph{hefen}} `heaven' in the same West-Saxon \emph{u}-umlaut
environment (\citeproc{ref-Campbell1959}{Campbell 1959, sect. 210.1}).

\subsection{Old English evidence}\label{old-english-evidence-116}

Old English dictionaries record the standard West Saxon noun as
{\emph{heofon}} `heaven', alongside Anglian or Mercian {\emph{hefen}}
`heaven' material (\citeproc{ref-ClarkHall1960}{Clark Hall 1960, 188};
\citeproc{ref-BosworthToller1898}{Bosworth and Toller 1898, 43}).
Campbell also cites an earlier stage {\emph{hefzen}} `heaven (Mercian)'
in the history of the word (\citeproc{ref-Campbell1959}{Campbell 1959,
sect. 381}).

The target of this entry is the West Saxon citation form {\emph{heofon}}
`heaven'. Its vowel history points toward the oblique stem rather than
the front-vocalic nominative line.

\subsection{Development to Old
English}\label{development-to-old-english-116}

From {\Recon{xémonų} `heaven'}, the West Saxon line passes through the
oblique-stem type reflected in northern WGmc {\Recon{hebun} `heaven'}
(\citeproc{ref-RingeTaylor2014}{Ringe and Taylor 2014, 324}). Campbell's
{\emph{heofon}} `heaven' beside {\emph{hefen}} `heaven' and earlier
{\emph{hefzen}} `heaven (Mercian)' show the later West-Saxon
back-mutation and suffix reshaping behind {\emph{heofon}} `heaven'
(\citeproc{ref-Campbell1959}{Campbell 1959, sect. 210.1};
\citeproc{ref-Campbell1959}{1959, sect. 381}).

The front-vocalic nominative line explains the dialectal {\emph{hefen}}
`heaven (Mercian)' type. West Saxon {\emph{heofon}} `heaven' reflects
the oblique stem that was generalized into the nominative position.

\subsection{Paradigm comparison}\label{paradigm-comparison-8}

The comparison below sets the relevant forms side by side. It
distinguishes the front-vocalic nominative line from the oblique stem
selected for West Saxon {\emph{heofon}} `heaven'.

{\def\LTcaptype{none} % do not increment counter
\begin{longtable}[]{@{}
  >{\raggedright\arraybackslash}p{(\linewidth - 8\tabcolsep) * \real{0.2000}}
  >{\raggedright\arraybackslash}p{(\linewidth - 8\tabcolsep) * \real{0.2000}}
  >{\raggedright\arraybackslash}p{(\linewidth - 8\tabcolsep) * \real{0.2000}}
  >{\raggedright\arraybackslash}p{(\linewidth - 8\tabcolsep) * \real{0.2000}}
  >{\raggedright\arraybackslash}p{(\linewidth - 8\tabcolsep) * \real{0.2000}}@{}}
\toprule\noalign{}
\begin{minipage}[b]{\linewidth}\raggedright
PGmc cell / interpretation
\end{minipage} & \begin{minipage}[b]{\linewidth}\raggedright
Candidate input
\end{minipage} & \begin{minipage}[b]{\linewidth}\raggedright
Expected or documented OE outcome
\end{minipage} & \begin{minipage}[b]{\linewidth}\raggedright
OE comparison form
\end{minipage} & \begin{minipage}[b]{\linewidth}\raggedright
Result
\end{minipage} \\
\midrule\noalign{}
\endhead
\bottomrule\noalign{}
\endlastfoot
citation nominative singular & {\emph{*xémenaz}} & front-vocalic
{\emph{hefen}} type outcome & {\emph{hefen}} / {\emph{heofen}} & useful
control, but not the West Saxon form used here \\
selected oblique singular & {\emph{*xémonų}} & regular output:
{\emph{heofon}} & {\emph{heofon}} & exact match between input, output,
and target \\
older pre-OE stage & inherited oblique line & earlier {\emph{hefzen}}
stage & {\emph{hefzen}} & historical background for the same West Saxon
development \\
\end{longtable}
}

\section{\texorpdfstring{live --- OE
\emph{lifeþ}}{live --- OE lifeþ}}\label{live-oe-lifeuxfe}

\index[iv]{02oe@\textbf{Old English}!lifeth@\emph{lifeþ}}
\index[iv]{03pgmc@\textbf{Proto-Germanic}!libena@*libēną}
\index[iv]{03pgmc@\textbf{Proto-Germanic}!libethi@*líbēθi}

Derivation: citation reconstruction \emph{*libēną}; form followed here
\emph{*líbēθi} \textgreater{} \emph{lifeþ} (late analogy).

\subsection{Derivation trace}\label{derivation-trace-117}

Proto input: \emph{*líbēθi}

\begingroup
\setlength{\fboxsep}{6pt}

\noindent\fbox{%
\begin{minipage}{0.97\linewidth}
\small
\begin{tabularx}{\linewidth}{@{}>{\raggedright\arraybackslash}p{0.460\linewidth}>{\centering\arraybackslash}X>{\raggedright\arraybackslash}p{0.460\linewidth}@{}}
\begin{minipage}[t]{\linewidth}
\centering\textbf{Earlier Germanic changes}\par
\vspace{0.35em}
\raggedright
\centering\textbf{}\par
\raggedright
\vspace{0.2em}
\begin{tabular}{@{}>{\raggedright\arraybackslash}p{0.74\linewidth}@{\hspace{0.55em}}>{\raggedright\arraybackslash}p{0.16\linewidth}@{\hspace{0.25em}}}
\multicolumn{2}{@{}l@{}}{*Proto-West Germanic*} \\
PWGmc Early I Apocope & \emph{*líbēθ} \\
\multicolumn{2}{@{}l@{}}{*Northwest Germanic*} \\
\mbox{NWGmc Long E Lowering} & \emph{*líbǣθ} \\
\end{tabular}
\end{minipage}
&
&
\begin{minipage}[t]{\linewidth}
\centering\textbf{Old English changes}\par
\vspace{0.35em}
\raggedright
\centering\textbf{}\par
\raggedright
\vspace{0.2em}
\begin{tabular}{@{}>{\raggedright\arraybackslash}p{0.74\linewidth}@{\hspace{0.55em}}>{\raggedright\arraybackslash}p{0.16\linewidth}@{\hspace{0.25em}}}
\multicolumn{2}{@{}l@{}}{*Old English*} \\
PGmc B Allophony & \emph{*líβǣθ} \\
OE Unstressed Long Vowel Shortening & \emph{*líβæθ} \\
OE Unstressed AE Merger & \emph{*líβeθ} \\
\end{tabular}
\end{minipage}
\\
\end{tabularx}
\end{minipage}%
} \endgroup

Old English form: \emph{lifeþ}

\subsection{Reconstruction and comparative
evidence}\label{reconstruction-and-comparative-evidence-117}

The inherited verb belongs to the class-III weak family cited by Kroonen
under \emph{*libēn-}, reflected in Old English {\emph{libban}} `live'
(\citeproc{ref-Kroonen2013}{Kroonen 2013, 336}). Ringe and Taylor show
that the paradigm also contained a separate 3sg present stem, continued
in late Northumbrian {\emph{lifed}} `lives', which they treat as an
archaism (\citeproc{ref-RingeTaylor2014}{Ringe and Taylor 2014, 364}).

The derivational input {\Recon{líbēθi} `live'} therefore represents a
finite present cell rather than the citation infinitive. The later lemma
tradition also includes remodeled forms such as {\emph{lifian}} `live'.

\subsection{Old English evidence}\label{old-english-evidence-117}

The ordinary lemma tradition centers on {\emph{libban}} `live' and, in
later remodeling, {\emph{lifian}} `live'. For this entry, however, the
relevant comparison form is the archaic 3sg present attested as
{\emph{lifed}} `lives', here normalized as {\emph{lifeþ}} `lives'
(\citeproc{ref-RingeTaylor2014}{Ringe and Taylor 2014, 364};
\citeproc{ref-Campbell1959}{Campbell 1959, sect. 762}).

The target is thus a normalized finite form, not the ordinary dictionary
lemma. It preserves the older present-stem history more clearly than the
remodeled lemma tradition does.

\subsection{Development to Old
English}\label{development-to-old-english-117}

From {\Recon{líbēθi} `live'}, regular reduction of the final syllable
and later weakening of the unstressed vowel yield {\emph{lifeþ}}
`lives'. The attested spelling {\emph{lifed}} `lives' belongs to the
same finite form in late Northumbrian orthography
(\citeproc{ref-Campbell1959}{Campbell 1959, sect. 762};
\citeproc{ref-RingeTaylor2014}{Ringe and Taylor 2014, 364}).

\subsection{Paradigm comparison}\label{paradigm-comparison-9}

The comparison below sets the relevant forms side by side. It separates
the archaic finite cell from the ordinary infinitival and later
remodeled lemma lines.

{\def\LTcaptype{none} % do not increment counter
\begin{longtable}[]{@{}
  >{\raggedright\arraybackslash}p{(\linewidth - 8\tabcolsep) * \real{0.2000}}
  >{\raggedright\arraybackslash}p{(\linewidth - 8\tabcolsep) * \real{0.2000}}
  >{\raggedright\arraybackslash}p{(\linewidth - 8\tabcolsep) * \real{0.2000}}
  >{\raggedright\arraybackslash}p{(\linewidth - 8\tabcolsep) * \real{0.2000}}
  >{\raggedright\arraybackslash}p{(\linewidth - 8\tabcolsep) * \real{0.2000}}@{}}
\toprule\noalign{}
\begin{minipage}[b]{\linewidth}\raggedright
PGmc cell / interpretation
\end{minipage} & \begin{minipage}[b]{\linewidth}\raggedright
Candidate input
\end{minipage} & \begin{minipage}[b]{\linewidth}\raggedright
Expected or documented OE outcome
\end{minipage} & \begin{minipage}[b]{\linewidth}\raggedright
OE comparison form
\end{minipage} & \begin{minipage}[b]{\linewidth}\raggedright
Result
\end{minipage} \\
\midrule\noalign{}
\endhead
\bottomrule\noalign{}
\endlastfoot
citation infinitive line & {\emph{*libēną}} & OE {\emph{libban}}
headword tradition & {\emph{libban}} & establishes the lexeme, but not
the Old English form here \\
3sg present & {\emph{*líbēθi}} & regular output: {\emph{lifeþ}};
attested {\emph{lifed}} & {\emph{lifed}}, normalized here as
{\emph{lifeþ}} & selected archaic finite cell \\
later remodeled present tradition & later class-II-type forms &
{\emph{lifian}} and related finite remodeling & {\emph{lifian}} &
genuine OE development, but secondary to the cell compared here \\
\end{longtable}
}

\section{\texorpdfstring{man --- OE
\emph{mannes}}{man --- OE mannes}}\label{man-oe-mannes}

\index[iv]{02oe@\textbf{Old English}!mannes@\emph{mannes}}
\index[iv]{03pgmc@\textbf{Proto-Germanic}!mannas@*mánnas}
\index[iv]{03pgmc@\textbf{Proto-Germanic}!mannaz@*mánnaz}

Derivation: citation reconstruction \emph{*mánnaz}; form followed here
\emph{*mánnas} \textgreater{} \emph{mannes} (late analogy).

\subsection{Derivation trace}\label{derivation-trace-118}

Proto input: \emph{*mánnas}

\begingroup
\setlength{\fboxsep}{6pt}

\noindent\fbox{%
\begin{minipage}{0.97\linewidth}
\small
\begin{tabularx}{\linewidth}{@{}>{\raggedright\arraybackslash}p{0.440\linewidth}>{\centering\arraybackslash}X>{\raggedright\arraybackslash}p{0.440\linewidth}@{}}
\begin{minipage}[t]{\linewidth}
\centering\textbf{Earlier Germanic changes}\par
\vspace{0.35em}
\raggedright
\centering\textbf{}\par
\raggedright
\vspace{0.2em}
\begin{tabular}{@{}>{\raggedright\arraybackslash}p{0.68\linewidth}@{\hspace{0.55em}}>{\raggedright\arraybackslash}p{0.16\linewidth}@{\hspace{0.25em}}}
\multicolumn{2}{@{}l@{}}{*Proto-West Germanic*} \\
\multicolumn{2}{@{}l@{}}{[no change]} \\
\multicolumn{2}{@{}l@{}}{*Northwest Germanic*} \\
\multicolumn{2}{@{}l@{}}{[no change]} \\
\end{tabular}
\end{minipage}
&
&
\begin{minipage}[t]{\linewidth}
\centering\textbf{Old English changes}\par
\vspace{0.35em}
\raggedright
\centering\textbf{}\par
\raggedright
\vspace{0.2em}
\begin{tabular}{@{}>{\raggedright\arraybackslash}p{0.70\linewidth}@{\hspace{0.55em}}>{\raggedright\arraybackslash}p{0.16\linewidth}@{\hspace{0.25em}}}
\multicolumn{2}{@{}l@{}}{*Old English*} \\
Anglo Frisian Brightening & \emph{*mánnæs} \\
OE Unstressed AE Merger & \emph{*mánnes} \\
\end{tabular}
\end{minipage}
\\
\end{tabularx}
\end{minipage}%
} \endgroup

Old English form: \emph{mannes}

\subsection{Reconstruction and comparative
evidence}\label{reconstruction-and-comparative-evidence-118}

The lexeme-level reconstruction is not uniform. Kroonen cites
\emph{*mannan-}, and Orel has \Recon{mannz} `man'
(\citeproc{ref-Kroonen2013}{Kroonen 2013, 354};
\citeproc{ref-Orel2003}{Orel 2003, 299}). The derivational input
{\Recon{mánnas} `man'} belongs to a different level: it is the
genitive-singular cell chosen for the Old English comparison.

The target is not the ordinary citation form. Its medial geminate
precedes the genitive ending.

\subsection{Old English evidence}\label{old-english-evidence-118}

Campbell gives the paradigm {\emph{mann}} `man', {\emph{man}} `man' /
{\emph{mannes}} `man's / {\emph{menn}} 'men'
(\citeproc{ref-Campbell1959}{Campbell 1959, sect. 621}). Sievers-Brunner
likewise cites {\emph{man}} `man' {\emph{mannes}} `man's
(\citeproc{ref-SieversBrunner1965}{Brunner 1965, sect. 226}). He also
explains that word-final simplification underlies forms such as
{\emph{man}} 'man' beside inflected {\emph{monnes}} 'man's
(\citeproc{ref-SieversBrunner1965}{Brunner 1965, sect. 231}). Clark Hall
keeps the dictionary headword under {\emph{mann}}
(\citeproc{ref-ClarkHall1960}{Clark Hall 1960, 197}).

The relevant comparison form is therefore the attested genitive singular
{\emph{mannes}} `man's, not the citation lemma {\emph{mann}} 'man'.

\subsection{Development to Old
English}\label{development-to-old-english-118}

Campbell's paradigm \emph{mann}, \emph{man} / \emph{mannes} /
{\emph{menn}} `men' confirms the selected genitive singular
\emph{mannes} (\citeproc{ref-Campbell1959}{Campbell 1959, sect. 621}).
In the present analysis, \Recon{mánnas} `man' develops through
Anglo-Frisian brightening and later unstressed merger to \emph{mannes}.
In this cell the geminate remains medial before \emph{-es}. The citation
form behaves differently because word-final gemination was simplified in
Old English (\citeproc{ref-SieversBrunner1965}{Brunner 1965, sect.
231}).

\subsection{Paradigm comparison}\label{paradigm-comparison-10}

The comparison below sets the relevant forms side by side. It separates
the citation-form line from the genitive singular.

{\def\LTcaptype{none} % do not increment counter
\begin{longtable}[]{@{}
  >{\raggedright\arraybackslash}p{(\linewidth - 8\tabcolsep) * \real{0.2000}}
  >{\raggedright\arraybackslash}p{(\linewidth - 8\tabcolsep) * \real{0.2000}}
  >{\raggedright\arraybackslash}p{(\linewidth - 8\tabcolsep) * \real{0.2000}}
  >{\raggedright\arraybackslash}p{(\linewidth - 8\tabcolsep) * \real{0.2000}}
  >{\raggedright\arraybackslash}p{(\linewidth - 8\tabcolsep) * \real{0.2000}}@{}}
\toprule\noalign{}
\begin{minipage}[b]{\linewidth}\raggedright
PGmc cell / interpretation
\end{minipage} & \begin{minipage}[b]{\linewidth}\raggedright
Candidate input
\end{minipage} & \begin{minipage}[b]{\linewidth}\raggedright
Expected or documented OE outcome
\end{minipage} & \begin{minipage}[b]{\linewidth}\raggedright
OE comparison form
\end{minipage} & \begin{minipage}[b]{\linewidth}\raggedright
Result
\end{minipage} \\
\midrule\noalign{}
\endhead
\bottomrule\noalign{}
\endlastfoot
citation nominative singular & {\emph{*mannăz}} & expected citation-form
outcome {\emph{man}} & {\emph{mann}} / {\emph{monn}} & establishes the
lexeme, but not the Old English form here \\
accusative singular & {\emph{*manną}} & expected {\emph{man}} &
{\emph{man}} & same word-final simplification as the nominative \\
dative singular & {\emph{*mannăi}} & expected {\emph{manne}} &
{\emph{manne}} & preserves medial \emph{nn}, but not the chosen cell \\
genitive singular & {\emph{*mánnas}} & regular output: {\emph{mannes}} &
{\emph{mannes}} & exact match between input, output, and attested
comparator \\
\end{longtable}
}

\section{\texorpdfstring{meed --- OE
\emph{meorde}}{meed --- OE meorde}}\label{meed-oe-meorde}

\index[iv]{02oe@\textbf{Old English}!meorde@\emph{meorde}}
\index[iv]{03pgmc@\textbf{Proto-Germanic}!mizdai@*mízdai}
\index[iv]{03pgmc@\textbf{Proto-Germanic}!mizdo@*mizdō}

Derivation: citation reconstruction \emph{*mizdō}; form followed here
\emph{*mízdai} \textgreater{} \emph{meorde} (late analogy).

\subsection{Derivation trace}\label{derivation-trace-119}

Proto input: \emph{*mízdai}

\begingroup
\setlength{\fboxsep}{6pt}

\noindent\fbox{%
\begin{minipage}{0.97\linewidth}
\small
\begin{tabularx}{\linewidth}{@{}>{\raggedright\arraybackslash}p{0.440\linewidth}>{\centering\arraybackslash}X>{\raggedright\arraybackslash}p{0.440\linewidth}@{}}
\begin{minipage}[t]{\linewidth}
\centering\textbf{Earlier Germanic changes}\par
\vspace{0.35em}
\raggedright
\centering\textbf{}\par
\raggedright
\vspace{0.2em}
\begin{tabular}{@{}>{\raggedright\arraybackslash}p{0.70\linewidth}@{\hspace{0.55em}}>{\raggedright\arraybackslash}p{0.16\linewidth}@{\hspace{0.25em}}}
\multicolumn{2}{@{}l@{}}{*Proto-West Germanic*} \\
PWGmc Ai Monophthongization & \emph{*mírdē} \\
\multicolumn{2}{@{}l@{}}{*Northwest Germanic*} \\
NWGmc I Lowering & \emph{*mérdē} \\
\end{tabular}
\end{minipage}
&
&
\begin{minipage}[t]{\linewidth}
\centering\textbf{Old English changes}\par
\vspace{0.35em}
\raggedright
\centering\textbf{}\par
\raggedright
\vspace{0.2em}
\begin{tabular}{@{}>{\raggedright\arraybackslash}p{0.74\linewidth}@{\hspace{0.55em}}>{\raggedright\arraybackslash}p{0.16\linewidth}@{\hspace{0.25em}}}
\multicolumn{2}{@{}l@{}}{*Old English*} \\
OE Breaking & \emph{*méordē} \\
OE Unstressed Long Vowel Shortening & \emph{*méorde} \\
\end{tabular}
\end{minipage}
\\
\end{tabularx}
\end{minipage}%
} \endgroup

Old English form: \emph{meorde}

\subsection{Reconstruction and comparative
evidence}\label{reconstruction-and-comparative-evidence-119}

The lexeme-level reconstruction is {\Recon{mizdō} `meed'}, but the
derivational input {\Recon{mízdai} `meed'} is a dative-singular cell
rather than the citation form. The Old English evidence for the
{\emph{meord}} `meed, reward' side is oblique.

The wider history of competing {\emph{mēd}} `meed, reward' remains
disputed. Kroonen and Fulk explain it through some form of \emph{z}-loss
and compensatory lengthening (\citeproc{ref-Kroonen2013}{Kroonen 2013,
410}; \citeproc{ref-Fulk2018}{Fulk 2018, 69}), while Orel keeps a
doublet analysis (\citeproc{ref-Orel2003}{Orel 2003, 311}). The
comparison here concerns the attested oblique line \emph{meorde}.

\subsection{Old English evidence}\label{old-english-evidence-119}

The directly attested forms are obliques: {\emph{meorde}} `meed' as a
dative singular and {\emph{meorda}} `meed' as a genitive plural
(\citeproc{ref-BrightCassidyRingler1971}{Bright 1971, 328};
\citeproc{ref-BosworthToller1898}{Bosworth and Toller 1898, 647}).
Lexicographers reconstruct a bare nominative {\emph{meord}} `meed' from
those obliques, while West Saxon prose more commonly shows the competing
doublet \emph{mēd} `meed' (\citeproc{ref-ClarkHall1960}{Clark Hall 1960,
214}; \citeproc{ref-BosworthToller1898}{Bosworth and Toller 1898, 647}).

The target of this entry is therefore the attested oblique
{\emph{meorde}} `meed', not the reconstructed lemma {\emph{meord}}
`meed' and not the better-known West Saxon citation form \emph{mēd}
`meed'.

\subsection{Development to Old
English}\label{development-to-old-english-119}

Ringe and Taylor give the broader noun history as \Recon{mizdo} `reward'
\textgreater{} \Recon{mizdu} `reward' \textgreater{} OE meord
\textasciitilde{} méd\_(\citeproc{ref-RingeTaylor2014}{Ringe and Taylor
2014, 99}). The fuller oblique-cell path modeled here spells out the
intermediate rhotacism, monophthongization, lowering, breaking, and
unstressed-shortening steps needed for the selected dative-singular
comparison. This entry therefore follows the attested oblique line
within the broader development. It does not depend on a full decision
about the history of the competing {\emph{\_mēd`}} `meed, reward'
tradition.

\subsection{Paradigm comparison}\label{paradigm-comparison-11}

The comparison below sets the relevant forms side by side. It
distinguishes the attested oblique target from the broader lemma
history.

{\def\LTcaptype{none} % do not increment counter
\begin{longtable}[]{@{}
  >{\raggedright\arraybackslash}p{(\linewidth - 8\tabcolsep) * \real{0.2000}}
  >{\raggedright\arraybackslash}p{(\linewidth - 8\tabcolsep) * \real{0.2000}}
  >{\raggedright\arraybackslash}p{(\linewidth - 8\tabcolsep) * \real{0.2000}}
  >{\raggedright\arraybackslash}p{(\linewidth - 8\tabcolsep) * \real{0.2000}}
  >{\raggedright\arraybackslash}p{(\linewidth - 8\tabcolsep) * \real{0.2000}}@{}}
\toprule\noalign{}
\begin{minipage}[b]{\linewidth}\raggedright
PGmc cell / interpretation
\end{minipage} & \begin{minipage}[b]{\linewidth}\raggedright
Candidate input
\end{minipage} & \begin{minipage}[b]{\linewidth}\raggedright
Expected or documented OE outcome
\end{minipage} & \begin{minipage}[b]{\linewidth}\raggedright
OE comparison form
\end{minipage} & \begin{minipage}[b]{\linewidth}\raggedright
Result
\end{minipage} \\
\midrule\noalign{}
\endhead
\bottomrule\noalign{}
\endlastfoot
citation nominative singular & {\emph{*mizdō}} & inferred lemma outcome
{\emph{meord}} `meed' & {\emph{meord}} `meed' & useful background, but
the bare lemma is reconstructed rather than directly attested \\
selected dative singular & {\emph{*mízdai}} & regular output:
{\emph{meorde}} `meed' & {\emph{meorde}} `meed' & exact match between
derivational input and attested target \\
genitive singular & {\emph{*mizdōz}} & regular output: {\emph{meorde}}
`meed' & {\emph{meorde}} `meed' & converges on the same attested string,
but the dat.sg. has the clearest direct support \\
genitive plural control & plural oblique line & attested {\emph{meorda}}
`meed' & {\emph{meorda}} `meed' & confirms the broader oblique
tradition, but not the chosen singular target \\
\end{longtable}
}

\section{\texorpdfstring{night --- OE
\emph{niht}}{night --- OE niht}}\label{night-oe-niht}

\index[iv]{02oe@\textbf{Old English}!niht@\emph{niht}}
\index[iv]{03pgmc@\textbf{Proto-Germanic}!naxti@*náxti}
\index[iv]{03pgmc@\textbf{Proto-Germanic}!naxtz@*náxtz}

Derivation: citation reconstruction \emph{*náxtz}; form followed here
\emph{*náxti} \textgreater{} \emph{niht} (late analogy).

\subsection{Derivation trace}\label{derivation-trace-120}

Proto input: \emph{*náxti}

\begingroup
\setlength{\fboxsep}{6pt}

\noindent\fbox{%
\begin{minipage}{0.97\linewidth}
\small
\begin{tabularx}{\linewidth}{@{}>{\raggedright\arraybackslash}p{0.320\linewidth}>{\centering\arraybackslash}X>{\raggedright\arraybackslash}p{0.520\linewidth}@{}}
\begin{minipage}[t]{\linewidth}
\centering\textbf{Earlier Germanic changes}\par
\vspace{0.35em}
\raggedright
\centering\textbf{}\par
\raggedright
\vspace{0.2em}
\begin{tabular}{@{}>{\raggedright\arraybackslash}p{0.68\linewidth}@{\hspace{0.55em}}>{\raggedright\arraybackslash}p{0.16\linewidth}@{\hspace{0.25em}}}
\multicolumn{2}{@{}l@{}}{*Proto-West Germanic*} \\
\multicolumn{2}{@{}l@{}}{[no change]} \\
\multicolumn{2}{@{}l@{}}{*Northwest Germanic*} \\
\multicolumn{2}{@{}l@{}}{[no change]} \\
\end{tabular}
\end{minipage}
&
&
\begin{minipage}[t]{\linewidth}
\centering\textbf{Old English changes}\par
\vspace{0.35em}
\raggedright
\centering\textbf{}\par
\raggedright
\vspace{0.2em}
\begin{tabular}{@{}>{\raggedright\arraybackslash}p{0.74\linewidth}@{\hspace{0.55em}}>{\raggedright\arraybackslash}p{0.16\linewidth}@{\hspace{0.25em}}}
\multicolumn{2}{@{}l@{}}{*Old English*} \\
Anglo Frisian Brightening & \emph{*næxti} \\
OE Breaking & \emph{*neaxti} \\
OE I Umlaut & \emph{*niexti} \\
OE Ws Palatal Umlaut & \emph{*nixti} \\
\mbox{OE High Vowel Apocope} & \emph{*nixt} \\
\end{tabular}
\end{minipage}
\\
\end{tabularx}
\end{minipage}%
} \endgroup

Old English form: \emph{niht}

\subsection{Reconstruction and comparative
evidence}\label{reconstruction-and-comparative-evidence-120}

Ringe and Taylor cite gen.sg. \Recon{nahtiz} `night', dat.sg.
\Recon{nahti} `night', and nom.pl. \Recon{nahtiz} `night' for the
high-vowel side of the paradigm, and derive West Saxon {\emph{niht}}
`night' from that side (\citeproc{ref-RingeTaylor2014}{Ringe and Taylor
2014, 240}). The citation reconstruction {\Recon{náxtz} `night'}
therefore belongs to the nominative-like headword, while the
derivational input {\Recon{náxti} `night'} represents the
dative-singular cell.

The word later became the model for endingless datives. Ringe and Taylor
explicitly explain forms such as {\emph{dæg}} `day' by analogy with
dat.sg. niht \textless{} \Recon{nahti} `night'
(\citeproc{ref-RingeTaylor2014}{Ringe and Taylor 2014, 380}).

\subsection{Old English evidence}\label{old-english-evidence-120}

Clark Hall lemmatizes {\emph{niht}} `night' and cross-references forms
such as {\emph{neaht}} `night', {\emph{neht}} `night', and
{\emph{nieht}} `night' (\citeproc{ref-ClarkHall1960}{Clark Hall 1960,
215}). Campbell likewise preserves the fluctuation between
{\emph{neaht}} `night' and {\emph{niht}} `night', giving genitive
\emph{nihte, nihtes}, dative \emph{niht, nihte}, nominative plural
\emph{niht} `night', and the contrasting plural-side forms represented
by {\emph{neahtas}} `night' (\citeproc{ref-Campbell1959}{Campbell 1959,
sect. 628.3}).

The comparison form used here is therefore an attested Old English
\emph{niht} `night', not a reconstructed substitute. The broader lexical
record still preserves the non-umlauted side of the paradigm in
{\emph{neaht}} `night'-type forms.

\subsection{Development to Old
English}\label{development-to-old-english-120}

Ringe and Taylor derive West Saxon \emph{niht} `night' from
\Recon{nahti} `night' via \Recon{nehti} `night' and \Recon{neahti}
`night' (\citeproc{ref-RingeTaylor2014}{Ringe and Taylor 2014, 240}).
Campbell and Brunner preserve the contrasting non-umlauted
{\emph{neaht}} `night'-type forms elsewhere in the paradigm
(\citeproc{ref-Campbell1959}{Campbell 1959, sect. 628.3};
\citeproc{ref-SieversBrunner1965}{Brunner 1965, sect. 284}).

The modeled path is therefore \Recon{náxti} `night' \textgreater{}
\Recon{næxti} `night' \textgreater{} \Recon{neaxti} `night'
\textgreater{} \Recon{niexti} `night' \textgreater{} \Recon{nixti}
`night' \textgreater{} niht.

\subsection{Paradigm comparison}\label{paradigm-comparison-12}

The comparison below sets the relevant forms side by side. It identifies
the inherited cell that matches the attested Old English form.

{\def\LTcaptype{none} % do not increment counter
\begin{longtable}[]{@{}
  >{\raggedright\arraybackslash}p{(\linewidth - 8\tabcolsep) * \real{0.2000}}
  >{\raggedright\arraybackslash}p{(\linewidth - 8\tabcolsep) * \real{0.2000}}
  >{\raggedright\arraybackslash}p{(\linewidth - 8\tabcolsep) * \real{0.2000}}
  >{\raggedright\arraybackslash}p{(\linewidth - 8\tabcolsep) * \real{0.2000}}
  >{\raggedright\arraybackslash}p{(\linewidth - 8\tabcolsep) * \real{0.2000}}@{}}
\toprule\noalign{}
\begin{minipage}[b]{\linewidth}\raggedright
PGmc cell / interpretation
\end{minipage} & \begin{minipage}[b]{\linewidth}\raggedright
Candidate input
\end{minipage} & \begin{minipage}[b]{\linewidth}\raggedright
OE output or comparison
\end{minipage} & \begin{minipage}[b]{\linewidth}\raggedright
OE comparison form
\end{minipage} & \begin{minipage}[b]{\linewidth}\raggedright
Result
\end{minipage} \\
\midrule\noalign{}
\endhead
\bottomrule\noalign{}
\endlastfoot
citation nominative singular & {\emph{*náxtz}} & expected non-umlauted
outcome {\emph{neaht}} & {\emph{neaht}} & useful background, but not the
comparison used for \emph{niht} \\
selected dative singular & {\emph{*náxti}} & regular output:
{\emph{niht}} `night' & {\emph{niht}} `night' & exact match between
input, output, and paradigm cell \\
\end{longtable}
}

\section{\texorpdfstring{rest --- OE
\emph{ræste}}{rest --- OE ræste}}\label{rest-oe-ruxe6ste}

\index[iv]{02oe@\textbf{Old English}!raeste@\emph{ræste}}
\index[iv]{03pgmc@\textbf{Proto-Germanic}!rasto@*rastō}
\index[iv]{03pgmc@\textbf{Proto-Germanic}!rastoz@*rástōz}

Derivation: citation reconstruction \emph{*rastō}; form followed here
\emph{*rástōz} \textgreater{} \emph{ræste} (late analogy).

\subsection{Derivation trace}\label{derivation-trace-121}

Proto input: \emph{*rástōz}

\begingroup
\setlength{\fboxsep}{6pt}

\noindent\fbox{%
\begin{minipage}{0.97\linewidth}
\small
\begin{tabularx}{\linewidth}{@{}>{\raggedright\arraybackslash}p{0.460\linewidth}>{\centering\arraybackslash}X>{\raggedright\arraybackslash}p{0.460\linewidth}@{}}
\begin{minipage}[t]{\linewidth}
\centering\textbf{Earlier Germanic changes}\par
\vspace{0.35em}
\raggedright
\centering\textbf{}\par
\raggedright
\vspace{0.2em}
\begin{tabular}{@{}>{\raggedright\arraybackslash}p{0.74\linewidth}@{\hspace{0.55em}}>{\raggedright\arraybackslash}p{0.16\linewidth}@{\hspace{0.25em}}}
\multicolumn{2}{@{}l@{}}{*Proto-West Germanic*} \\
\multicolumn{2}{@{}l@{}}{[no change]} \\
\multicolumn{2}{@{}l@{}}{*Northwest Germanic*} \\
\mbox{PGmc Final Z Deletion} & \emph{*rástō} \\
\end{tabular}
\end{minipage}
&
&
\begin{minipage}[t]{\linewidth}
\centering\textbf{Old English changes}\par
\vspace{0.35em}
\raggedright
\centering\textbf{}\par
\raggedright
\vspace{0.2em}
\begin{tabular}{@{}>{\raggedright\arraybackslash}p{0.74\linewidth}@{\hspace{0.55em}}>{\raggedright\arraybackslash}p{0.16\linewidth}@{\hspace{0.25em}}}
\multicolumn{2}{@{}l@{}}{*Old English*} \\
PWGmc Surviving Bimoric O Unrounding & \emph{*rástā} \\
Anglo Frisian Brightening & \emph{*ræstǣ} \\
OE Unstressed Long Vowel Shortening & \emph{*ræstæ} \\
OE Unstressed AE Merger & \emph{*ræste} \\
\end{tabular}
\end{minipage}
\\
\end{tabularx}
\end{minipage}%
} \endgroup

Old English form: \emph{ræste}

\subsection{Reconstruction and comparative
evidence}\label{reconstruction-and-comparative-evidence-121}

Kroonen treats the noun as a feminine ō-stem {\emph{*rastō-}}, continued
by Old English {\emph{ræst}} `rest' (\citeproc{ref-Kroonen2013}{Kroonen
2013, 445}). The derivational input {\Recon{rástōz} `rest'} therefore
does not replace the lexeme-level headword. It identifies one oblique
singular cell on the side of the paradigm that yields \emph{ræste}
`rest'.

The source tradition used here labels that cell specifically as genitive
singular, but the broader local synthesis of the ō-stem paradigm shows
that the oblique singulars converge on the same front-vocalic
\emph{ræste} `rest' side, in contrast to a nominative singular that
would remain \emph{rast} `rest'.

\subsection{Old English evidence}\label{old-english-evidence-121}

The ordinary Old English citation form is {\emph{ræst}} `rest'
(\citeproc{ref-Kroonen2013}{Kroonen 2013, 445}). Clark Hall likewise
gives {\emph{ræst}} `rest' (\citeproc{ref-ClarkHall1960}{Clark Hall
1960, 239}). Bosworth-Toller also preserves oblique uses of
{\emph{ræste}} `rest', including prepositional examples such as \emph{on
ræste} `at rest' and \emph{tó ræste} `to rest'
(\citeproc{ref-BosworthToller1898}{Bosworth and Toller 1898, 121}).

The comparison form used here is therefore an attested oblique
{\emph{ræste}} `rest', not a reconstructed surrogate. The dictionary
headword {\emph{ræst}} `rest' remains an equally real part of the Old
English record.

\subsection{Development to Old
English}\label{development-to-old-english-121}

Once final \emph{*z} is lost, the derivational input moves through the
front-vocalic oblique side of the paradigm rather than the back-vocalic
nominative side. In the modeled derivation, the surviving final long
vowel is first exposed, unrounded, fronted, shortened, and then reduced
to the final \emph{-e} of \emph{ræste} `rest'.

The key point is the paradigm split. Nominative \Recon{rastō} `rest'
yields a regular \emph{rast} `rest', whereas the selected oblique input
yields \emph{ræste} `rest'. The later citation form \emph{ræst} `rest'
is best explained as leveling from that oblique \emph{ræst-} stem.

\subsection{Paradigm comparison}\label{paradigm-comparison-13}

The comparison below sets the relevant forms side by side. It
distinguishes the nominative citation form from the oblique singular
chosen here.

{\def\LTcaptype{none} % do not increment counter
\begin{longtable}[]{@{}
  >{\raggedright\arraybackslash}p{(\linewidth - 8\tabcolsep) * \real{0.2000}}
  >{\raggedright\arraybackslash}p{(\linewidth - 8\tabcolsep) * \real{0.2000}}
  >{\raggedright\arraybackslash}p{(\linewidth - 8\tabcolsep) * \real{0.2000}}
  >{\raggedright\arraybackslash}p{(\linewidth - 8\tabcolsep) * \real{0.2000}}
  >{\raggedright\arraybackslash}p{(\linewidth - 8\tabcolsep) * \real{0.2000}}@{}}
\toprule\noalign{}
\begin{minipage}[b]{\linewidth}\raggedright
PGmc cell / interpretation
\end{minipage} & \begin{minipage}[b]{\linewidth}\raggedright
Candidate input
\end{minipage} & \begin{minipage}[b]{\linewidth}\raggedright
OE output or comparison
\end{minipage} & \begin{minipage}[b]{\linewidth}\raggedright
OE comparison form
\end{minipage} & \begin{minipage}[b]{\linewidth}\raggedright
Result
\end{minipage} \\
\midrule\noalign{}
\endhead
\bottomrule\noalign{}
\endlastfoot
citation nominative singular & {\emph{*rastō}} & expected regular
outcome \emph{rast} & {\emph{ræst}} `rest' & useful background, but not
the cell that matches attested oblique \emph{ræste} \\
selected oblique singular & {\emph{*rástōz}} & regular output:
{\emph{ræste}} `rest' & {\emph{ræste}} `rest' & exact match between
derivational input and attested OE oblique form \\
\end{longtable}
}

\section{\texorpdfstring{shoulder --- OE
\emph{sċuldrum}}{shoulder --- OE sċuldrum}}\label{shoulder-oe-sux10buldrum}

\index[iv]{02oe@\textbf{Old English}!sculdrum@\emph{sċuldrum}}
\index[iv]{03pgmc@\textbf{Proto-Germanic}!skuldramiz@*skúldramiz}
\index[iv]{03pgmc@\textbf{Proto-Germanic}!skuldro@*skuldrō}

Derivation: citation reconstruction \emph{*skuldrō}; form followed here
\emph{*skúldramiz} \textgreater{} \emph{sċuldrum} (late analogy).

\subsection{Derivation trace}\label{derivation-trace-122}

Proto input: \emph{*skúldramiz}

\begingroup
\setlength{\fboxsep}{6pt}

\noindent\fbox{%
\begin{minipage}{0.97\linewidth}
\small
\begin{tabularx}{\linewidth}{@{}>{\raggedright\arraybackslash}p{0.520\linewidth}>{\centering\arraybackslash}X>{\raggedright\arraybackslash}p{0.320\linewidth}@{}}
\begin{minipage}[t]{\linewidth}
\centering\textbf{Earlier Germanic changes}\par
\vspace{0.35em}
\raggedright
\centering\textbf{}\par
\raggedright
\vspace{0.2em}
\begin{tabular}{@{}>{\raggedright\arraybackslash}p{0.62\linewidth}@{\hspace{0.55em}}>{\raggedright\arraybackslash}p{0.28\linewidth}@{\hspace{0.25em}}}
\multicolumn{2}{@{}l@{}}{*Proto-West Germanic*} \\
NWGmc A To U Before M & \emph{*skúldrumiz} \\
PWGmc Early I Apocope & \emph{*skúldrumz} \\
\multicolumn{2}{@{}l@{}}{*Northwest Germanic*} \\
PGmc Final Z Deletion & \emph{*skúldrum} \\
\end{tabular}
\end{minipage}
&
&
\begin{minipage}[t]{\linewidth}
\centering\textbf{Old English changes}\par
\vspace{0.35em}
\raggedright
\centering\textbf{}\par
\raggedright
\vspace{0.2em}
\begin{tabular}{@{}>{\raggedright\arraybackslash}p{0.68\linewidth}@{\hspace{0.55em}}>{\raggedright\arraybackslash}p{0.20\linewidth}@{\hspace{0.25em}}}
\multicolumn{2}{@{}l@{}}{*Old English*} \\
OE Sk Palatalization & \emph{*ʃúldrum} \\
\end{tabular}
\end{minipage}
\\
\end{tabularx}
\end{minipage}%
} \endgroup

Old English form: \emph{sċuldrum}

\subsection{Reconstruction and comparative
evidence}\label{reconstruction-and-comparative-evidence-122}

The handbooks do not agree on the reconstruction of the Germanic word.
Orel gives \emph{*skuldr(j)ō} as a feminine ō-/jō-stem and explicitly
notes that Old English {\emph{sculdor}} `shoulder' is masculine beside
OFrisian {\emph{skulder}} `shoulder', Middle Low German {schulder}
`shoulder', and Old High German {\emph{scultra}} `shoulder',
{\emph{scultirra}} `shoulder' (\citeproc{ref-Orel2003}{Orel 2003, 345}).
Kroonen reconstructs {\emph{*skuldra-}}, a masculine a-stem, and derives
the Old High German feminine forms from {\emph{*skuldrjōn-}}
(\citeproc{ref-Kroonen2013}{Kroonen 2013, 478}). Ringe and Taylor cite
\textbf{PWGmc {\Recon{skuldru} `shoulder'}} for the Old English branch
(\citeproc{ref-RingeTaylor2014}{Ringe and Taylor 2014, 142}).

These forms imply different stem classes and different expectations for
the Old English inflection. The question is which inflectional cell best
aligns with the Old English evidence.

A dative/instrumental plural form \textbf{{\Recon{skúldramiz}
`shoulder'}} aligns with the inherited plural ending that later yields
Old English \emph{-um}, and it corresponds directly to the attested
dative plural discussed below.

\subsection{Old English evidence}\label{old-english-evidence-122}

The ordinary Old English headword is {\emph{sculdor}} `shoulder'. Clark
Hall lemmatizes {\emph{sculdor}} `shoulder' as the normal dictionary
form (\citeproc{ref-ClarkHall1960}{Clark Hall 1960, 257}).
Bosworth-Toller also preserves the dative plural {\emph{sculdrum}}
`shoulder' (\citeproc{ref-BosworthToller1898}{Bosworth and Toller 1898,
85}).

Bosworth-Toller's Supplement records a weak-feminine {\emph{sculdra}}
`shoulder', an (\citeproc{ref-BosworthToller1898}{Bosworth and Toller
1898, 699}), so {\emph{sculdra}} `shoulder' belongs to the Old English
record beside the stronger masculine paradigm headed by {\emph{sculdor}}
`shoulder'. Brunner and Luick also record later spellings such as
{\emph{sceoldor}} `shoulder' and the i-mutated dative plural
{\emph{scyldrum}} `shoulder', which reflect secondary phonological and
analogical reshaping within Old English
(\citeproc{ref-SieversBrunner1965}{Brunner 1965, sect. 92.2.a};
\citeproc{ref-Luick1914}{Luick 1914-\/-1940, 230}).

The singular and plural evidence point to different parts of the
paradigm. The relevant comparison form here is the attested dative
plural {\emph{sċuldrum}} `shoulder (dat.pl.)'. The spelling with
\emph{sċ-} is a normalized representation of the same Old English
initial cluster.

\subsection{Development to Old
English}\label{development-to-old-english-122}

Proto-Germanic \textbf{\Recon{skúldramiz} `shoulder'} can be interpreted
as a dative/instrumental plural form. In this environment the post-tonic
\emph{a} before \emph{m} is raised to \emph{u}, giving a form of the
\textbf{\Recon{skúldrumiz} `shoulder'} type. Unstressed \emph{u} is
regularly preserved before \emph{m}, especially in the dative plural
ending \emph{-um}: Campbell states this explicitly, and Hogg formulates
the same condition for the dative plural inflexion
(\citeproc{ref-Campbell1959}{Campbell 1959, sect. 373};
\citeproc{ref-Hogg1992}{Hogg 1992, sect. 3.3.1.3}). Brunner points in
the same direction by excluding \emph{m} from the environments in which
medial \emph{o} became general in West Saxon
(\citeproc{ref-SieversBrunner1965}{Brunner 1965, sect. 44} Anm. 7).

Subsequent reduction of the ending removes the final \emph{*i} and
\emph{*z}, so that the inflectional ending appears in Old English as
\emph{-um}. The initial cluster is written here as \emph{sċ-}, and the
development is \Recon{skúldramiz} `shoulder' \textgreater{}
\Recon{skúldrumiz} `shoulder' \textgreater{} \Recon{skúldrum} `shoulder'
\textgreater{} {\emph{sċuldrum}} `shoulder'.

\subsection{Paradigm comparison}\label{paradigm-comparison-14}

A paradigm comparison identifies the Proto-Germanic inflectional cell
that corresponds to an established Old English paradigm form. The
comparison below sets the relevant forms side by side.

{\def\LTcaptype{none} % do not increment counter
\begin{longtable}[]{@{}
  >{\raggedright\arraybackslash}p{(\linewidth - 8\tabcolsep) * \real{0.2000}}
  >{\raggedright\arraybackslash}p{(\linewidth - 8\tabcolsep) * \real{0.2000}}
  >{\raggedright\arraybackslash}p{(\linewidth - 8\tabcolsep) * \real{0.2000}}
  >{\raggedright\arraybackslash}p{(\linewidth - 8\tabcolsep) * \real{0.2000}}
  >{\raggedright\arraybackslash}p{(\linewidth - 8\tabcolsep) * \real{0.2000}}@{}}
\toprule\noalign{}
\begin{minipage}[b]{\linewidth}\raggedright
PGmc cell / interpretation
\end{minipage} & \begin{minipage}[b]{\linewidth}\raggedright
Candidate input
\end{minipage} & \begin{minipage}[b]{\linewidth}\raggedright
OE output or comparison
\end{minipage} & \begin{minipage}[b]{\linewidth}\raggedright
OE comparison form
\end{minipage} & \begin{minipage}[b]{\linewidth}\raggedright
Result
\end{minipage} \\
\midrule\noalign{}
\endhead
\bottomrule\noalign{}
\endlastfoot
singular-oriented citation input & {\emph{*skúldrō}} & probe output:
{\emph{sċoldor}} & {\emph{sculdor}} `shoulder' & fails: the singular
output has root \emph{o}, not the attested \emph{u} \\
serious plural-based singular alternative & {\emph{*skúldru}} & probe
output: {\emph{sċuldor}} & {\emph{sculdor}} `shoulder' & close formally,
but it compares a plural-stage input with a singular form \\
dat./inst.pl. input & {\Recon{skúldramiz} `shoulder'} & regular output:
{\emph{sċuldrum}} `shoulder (dat.pl.)' & {\emph{sculdrum}} `shoulder' &
matches both the output and the dative plural comparison form \\
later weak-feminine singular & --- & OE {\emph{sculdra}} `shoulder' &
{\emph{sculdra}} `shoulder' & secondary doublet, useful as a control
rather than the inherited target \\
\end{longtable}
}

The dative plural line is decisive because it matches both the output
and the paradigm cell of Old English {\emph{sculdrum}} `shoulder'.
Singular-oriented candidates either lower the root vowel or compare
unlike cells.

\section{\texorpdfstring{shove --- OE
\emph{sċēaf}}{shove --- OE sċēaf}}\label{shove-oe-sux10bux113af}

\index[iv]{02oe@\textbf{Old English}!sceaf@\emph{sċēaf}}
\index[iv]{03pgmc@\textbf{Proto-Germanic}!skaub@*skáub}
\index[iv]{03pgmc@\textbf{Proto-Germanic}!skeubana@*skéubaną}

Derivation: citation reconstruction \emph{*skéubaną}; form followed here
\emph{*skáub} \textgreater{} \emph{sċēaf} (late analogy).

\subsection{Derivation trace}\label{derivation-trace-123}

Proto input: \emph{*skáub}

\begingroup
\setlength{\fboxsep}{6pt}

\noindent\fbox{%
\begin{minipage}{0.97\linewidth}
\small
\begin{tabularx}{\linewidth}{@{}>{\raggedright\arraybackslash}p{0.320\linewidth}>{\centering\arraybackslash}X>{\raggedright\arraybackslash}p{0.520\linewidth}@{}}
\begin{minipage}[t]{\linewidth}
\centering\textbf{Earlier Germanic changes}\par
\vspace{0.35em}
\raggedright
\centering\textbf{}\par
\raggedright
\vspace{0.2em}
\begin{tabular}{@{}>{\raggedright\arraybackslash}p{0.68\linewidth}@{\hspace{0.55em}}>{\raggedright\arraybackslash}p{0.16\linewidth}@{\hspace{0.25em}}}
\multicolumn{2}{@{}l@{}}{*Proto-West Germanic*} \\
\multicolumn{2}{@{}l@{}}{[no change]} \\
\multicolumn{2}{@{}l@{}}{*Northwest Germanic*} \\
\multicolumn{2}{@{}l@{}}{[no change]} \\
\end{tabular}
\end{minipage}
&
&
\begin{minipage}[t]{\linewidth}
\centering\textbf{Old English changes}\par
\vspace{0.35em}
\raggedright
\centering\textbf{}\par
\raggedright
\vspace{0.2em}
\begin{tabular}{@{}>{\raggedright\arraybackslash}p{0.68\linewidth}@{\hspace{0.55em}}>{\raggedright\arraybackslash}p{0.16\linewidth}@{\hspace{0.25em}}}
\multicolumn{2}{@{}l@{}}{*Old English*} \\
OE Au Fronting & \emph{*skáeub} \\
OE Diphthong Leveling & \emph{*skēab} \\
PGmc B Allophony & \emph{*skēaβ} \\
OE Sk Palatalization & \emph{*ʃēaβ} \\
\end{tabular}
\end{minipage}
\\
\end{tabularx}
\end{minipage}%
} \endgroup

Old English form: \emph{sċēaf}

\subsection{Reconstruction and comparative
evidence}\label{reconstruction-and-comparative-evidence-123}

Kroonen reconstructs the strong verb as \emph{*skeuban-}
\textasciitilde{} \emph{*skūban-} and cites Old English present forms
{\emph{scēofan}} `shove', {\emph{scūfan}} `shove'
(\citeproc{ref-Kroonen2013}{Kroonen 2013, 444}). Those present-system
forms belong to the same verb family, but the comparison here uses the
singular preterite {\Recon{skáub} `shove'}, not the infinitive.

\subsection{Old English evidence}\label{old-english-evidence-123}

The ordinary dictionary verb is {\emph{scūfan}} `shove'/{\emph{scēofan}}
`shove', but the preterite itself is well attested. Bright gives the
principal parts {\emph{scufan}} `shove', {\emph{sceaf}} `shove',
{\emph{scufon}} `shove', {\emph{scofen}} `shove'
(\citeproc{ref-BrightCassidyRingler1971}{Bright 1971, 347}). Sweet gives
the same paradigm (\citeproc{ref-Sweet1953}{Sweet 1953, 29}). The
normalized form here is {\emph{sċēaf}} `sheaf', regularizing the
attested spellings {\emph{sceaf}} `shove' and prefixed {\emph{āsceaf}}
`shove'.

\subsection{Development to Old
English}\label{development-to-old-english-123}

From \Recon{skáub} `shove', the development is straightforward.
\emph{*au} fronts and levels to \emph{ēa}, final \emph{*b} becomes a
fricative and is written \emph{f}, and initial \emph{*sk-} undergoes the
usual Old English palatalized spelling in this environment. The
derivation therefore gives \Recon{skáub} `shove' \textgreater{}
\Recon{skáeub} `shove' \textgreater{} \Recon{skēab} `shove'
\textgreater{} \Recon{skēaβ} `shove' \textgreater{} {\emph{sċēaf}}
`sheaf'.

\subsection{Paradigm comparison}\label{paradigm-comparison-15}

A paradigm comparison is needed here because the ordinary citation verb
and {\emph{sċēaf}} `sheaf' belong to different cells of the same strong
paradigm. The comparison below is manual.

{\def\LTcaptype{none} % do not increment counter
\begin{longtable}[]{@{}
  >{\raggedright\arraybackslash}p{(\linewidth - 8\tabcolsep) * \real{0.2000}}
  >{\raggedright\arraybackslash}p{(\linewidth - 8\tabcolsep) * \real{0.2000}}
  >{\raggedright\arraybackslash}p{(\linewidth - 8\tabcolsep) * \real{0.2000}}
  >{\raggedright\arraybackslash}p{(\linewidth - 8\tabcolsep) * \real{0.2000}}
  >{\raggedright\arraybackslash}p{(\linewidth - 8\tabcolsep) * \real{0.2000}}@{}}
\toprule\noalign{}
\begin{minipage}[b]{\linewidth}\raggedright
PGmc cell / interpretation
\end{minipage} & \begin{minipage}[b]{\linewidth}\raggedright
Candidate input
\end{minipage} & \begin{minipage}[b]{\linewidth}\raggedright
OE output or comparison
\end{minipage} & \begin{minipage}[b]{\linewidth}\raggedright
OE comparison form
\end{minipage} & \begin{minipage}[b]{\linewidth}\raggedright
Result
\end{minipage} \\
\midrule\noalign{}
\endhead
\bottomrule\noalign{}
\endlastfoot
citation infinitive & {\emph{*skéubaną}} & inherited infinitive line
{\emph{sċēofan}}; present system also leveled {\emph{scūfan}} `shove' &
{\emph{scēofan}} `shove' / {\emph{scūfan}} `shove' & necessary
background, but not the comparison used for {\emph{sċēaf}} `sheaf' \\
1/3 sg. preterite & {\emph{*skáub}} & documented regular output:
{\emph{sċēaf}} `sheaf' & {\emph{sċēaf}} `sheaf' & direct match for the
singular preterite \\
preterite plural & {\emph{*skúbun}} & later leveled plural
{\emph{scufon}} `shove' beside expected {\emph{sċufun}} under the
corrected cascade & {\emph{scufon}} `shove' & poorer comparison for the
singular-preterite target \\
past participle & {\emph{*skúbanaz}} & attested participial line
{\emph{scofen}} `shove' & {\emph{scofen}} `shove' & valid alternative
cell, but not the form compared here \\
\end{longtable}
}

\section{\texorpdfstring{span --- OE
\emph{spanne}}{span --- OE spanne}}\label{span-oe-spanne}

\index[iv]{02oe@\textbf{Old English}!spanne@\emph{spanne}}
\index[iv]{03pgmc@\textbf{Proto-Germanic}!spannai@*spánnai}
\index[iv]{03pgmc@\textbf{Proto-Germanic}!spanno@*spannō}

Derivation: citation reconstruction \emph{*spannō}; form followed here
\emph{*spánnai} \textgreater{} \emph{spanne} (late analogy).

\subsection{Derivation trace}\label{derivation-trace-124}

Proto input: \emph{*spánnai}

\begingroup
\setlength{\fboxsep}{6pt}

\noindent\fbox{%
\begin{minipage}{0.97\linewidth}
\small
\begin{tabularx}{\linewidth}{@{}>{\raggedright\arraybackslash}p{0.440\linewidth}>{\centering\arraybackslash}X>{\raggedright\arraybackslash}p{0.440\linewidth}@{}}
\begin{minipage}[t]{\linewidth}
\centering\textbf{Earlier Germanic changes}\par
\vspace{0.35em}
\raggedright
\centering\textbf{}\par
\raggedright
\vspace{0.2em}
\begin{tabular}{@{}>{\raggedright\arraybackslash}p{0.70\linewidth}@{\hspace{0.55em}}>{\raggedright\arraybackslash}p{0.16\linewidth}@{\hspace{0.25em}}}
\multicolumn{2}{@{}l@{}}{*Proto-West Germanic*} \\
PWGmc Ai Monophthongization & \emph{*spánnē} \\
\multicolumn{2}{@{}l@{}}{*Northwest Germanic*} \\
\multicolumn{2}{@{}l@{}}{[no change]} \\
\end{tabular}
\end{minipage}
&
&
\begin{minipage}[t]{\linewidth}
\centering\textbf{Old English changes}\par
\vspace{0.35em}
\raggedright
\centering\textbf{}\par
\raggedright
\vspace{0.2em}
\begin{tabular}{@{}>{\raggedright\arraybackslash}p{0.74\linewidth}@{\hspace{0.55em}}>{\raggedright\arraybackslash}p{0.16\linewidth}@{\hspace{0.25em}}}
\multicolumn{2}{@{}l@{}}{*Old English*} \\
OE Unstressed Long Vowel Shortening & \emph{*spánne} \\
\end{tabular}
\end{minipage}
\\
\end{tabularx}
\end{minipage}%
} \endgroup

Old English form: \emph{spanne}

\subsection{Reconstruction and comparative
evidence}\label{reconstruction-and-comparative-evidence-124}

Seebold gives Old English {\emph{spann}} `span' under this noun family
(\citeproc{ref-Seebold1970}{Seebold 1970, 450}). The form followed here,
{\Recon{spánnai} `span'}, is therefore not a rival headword, but the
specific dative singular form compared on the model of the feminine
ō-stem paradigm (\citeproc{ref-SieversBrunner1965}{Brunner 1965, sect.
252}; \citeproc{ref-SieversBrunner1965}{1965, sect. 255.2}).

\subsection{Old English evidence}\label{old-english-evidence-124}

The reviewed lexicographic evidence more directly supports the citation
noun {\emph{spann}} `span' than the exact form {\emph{spanne}} `span'.
Clark Hall gives {\emph{spann}} `span'
(\citeproc{ref-ClarkHall1960}{Clark Hall 1960, 286}), and
{\emph{spanne}} `span' is accordingly treated as the regular dative
singular form compared here rather than as a dictionary headword.

\subsection{Development to Old
English}\label{development-to-old-english-124}

Citation {\Recon{spannō} `span'} yields {\emph{span}} `span'. The
oblique cell {\Recon{spánnai} `span'} therefore supplies the
conservative comparison form: it preserves the medial geminate and
yields {\emph{spanne}} `span', while citation {\Recon{spannō}} gives the
nominative background form.

\subsection{Paradigm comparison}\label{paradigm-comparison-16}

The comparison below sets the citation form beside the dative singular
form compared here.

{\def\LTcaptype{none} % do not increment counter
\begin{longtable}[]{@{}
  >{\raggedright\arraybackslash}p{(\linewidth - 8\tabcolsep) * \real{0.2000}}
  >{\raggedright\arraybackslash}p{(\linewidth - 8\tabcolsep) * \real{0.2000}}
  >{\raggedright\arraybackslash}p{(\linewidth - 8\tabcolsep) * \real{0.2000}}
  >{\raggedright\arraybackslash}p{(\linewidth - 8\tabcolsep) * \real{0.2000}}
  >{\raggedright\arraybackslash}p{(\linewidth - 8\tabcolsep) * \real{0.2000}}@{}}
\toprule\noalign{}
\begin{minipage}[b]{\linewidth}\raggedright
PGmc cell / interpretation
\end{minipage} & \begin{minipage}[b]{\linewidth}\raggedright
Candidate input
\end{minipage} & \begin{minipage}[b]{\linewidth}\raggedright
OE output or comparison
\end{minipage} & \begin{minipage}[b]{\linewidth}\raggedright
OE comparison form
\end{minipage} & \begin{minipage}[b]{\linewidth}\raggedright
Result
\end{minipage} \\
\midrule\noalign{}
\endhead
\bottomrule\noalign{}
\endlastfoot
citation nominative singular & {\emph{*spannō}} & regular output:
{\emph{span}} `span' & {\emph{spann}} `span' & useful citation-form
background, but not the form compared here \\
dative singular compared here & {\emph{*spánnai}} & regular output:
{\emph{spanne}} `span' & {\emph{spanne}} `span' & exact match for that
conservative form \\
\end{longtable}
}

\section{\texorpdfstring{thistle --- OE
\emph{þistles}}{thistle --- OE þistles}}\label{thistle-oe-uxfeistles}

\index[iv]{02oe@\textbf{Old English}!thistles@\emph{þistles}}
\index[iv]{03pgmc@\textbf{Proto-Germanic}!thestilaz@*θéstilaz}
\index[iv]{03pgmc@\textbf{Proto-Germanic}!thistilas@*θístilas}

Derivation: citation reconstruction \emph{*θéstilaz}; form followed here
\emph{*θístilas} \textgreater{} \emph{þistles} (late analogy).

\subsection{Derivation trace}\label{derivation-trace-125}

Proto input: \emph{*θístilas}

\begingroup
\setlength{\fboxsep}{6pt}

\noindent\fbox{%
\begin{minipage}{0.97\linewidth}
\small
\begin{tabularx}{\linewidth}{@{}>{\raggedright\arraybackslash}p{0.320\linewidth}>{\centering\arraybackslash}X>{\raggedright\arraybackslash}p{0.520\linewidth}@{}}
\begin{minipage}[t]{\linewidth}
\centering\textbf{Earlier Germanic changes}\par
\vspace{0.35em}
\raggedright
\centering\textbf{}\par
\raggedright
\vspace{0.2em}
\begin{tabular}{@{}>{\raggedright\arraybackslash}p{0.68\linewidth}@{\hspace{0.55em}}>{\raggedright\arraybackslash}p{0.16\linewidth}@{\hspace{0.25em}}}
\multicolumn{2}{@{}l@{}}{*Proto-West Germanic*} \\
\multicolumn{2}{@{}l@{}}{[no change]} \\
\multicolumn{2}{@{}l@{}}{*Northwest Germanic*} \\
\multicolumn{2}{@{}l@{}}{[no change]} \\
\end{tabular}
\end{minipage}
&
&
\begin{minipage}[t]{\linewidth}
\centering\textbf{Old English changes}\par
\vspace{0.35em}
\raggedright
\centering\textbf{}\par
\raggedright
\vspace{0.2em}
\begin{tabular}{@{}>{\raggedright\arraybackslash}p{0.66\linewidth}@{\hspace{0.55em}}>{\raggedright\arraybackslash}p{0.22\linewidth}@{\hspace{0.25em}}}
\multicolumn{2}{@{}l@{}}{*Old English*} \\
Anglo Frisian Brightening & \emph{*θístilæs} \\
OE L Adjacent Syncope & \emph{*θístlæs} \\
OE Unstressed AE Merger & \emph{*θístles} \\
\end{tabular}
\end{minipage}
\\
\end{tabularx}
\end{minipage}%
} \endgroup

Old English form: \emph{þistles}

\subsection{Reconstruction and comparative
evidence}\label{reconstruction-and-comparative-evidence-125}

Orel prints \emph{*þe(x)stilaz} `thistle' for the lexeme
(\citeproc{ref-Orel2003}{Orel 2003, 458}). The comparative label
{\Recon{θéstilaz} `thistle'} therefore remains in view as the
lexeme-level headword, while the derivational input {\Recon{θístilas}
`thistle'} is a specific genitive singular cell.

\subsection{Old English evidence}\label{old-english-evidence-125}

The ordinary simplex headword tradition is broken {\emph{þistel}}
`thistle' / {\emph{ðistel}} `thistle'. Clark Hall gives {\emph{ðistel}}
as the noun headword (\citeproc{ref-ClarkHall1960}{Clark Hall 1960,
326}). The Old English form here here is the genitive singular
{\emph{þistles}} `thistle', which preserves the same stem in an oblique
form where the cluster is medial.

\subsection{Development to Old
English}\label{development-to-old-english-125}

Campbell's discussion of cluster nouns shows the contrast clearly.
Simplex forms often develop a parasite vowel in word-final obstruent +
sonorant clusters, while comparable medial clusters remain unbroken; his
examples include \emph{hrefn}, {\emph{tacn}} `token, sign',
{\emph{wépn}} `weapon', and {\emph{botm}} `bottom' beside forms with
parasitic vowels elsewhere in the same lexical class
(\citeproc{ref-Campbell1959}{Campbell 1959, 151}). The genitive singular
{\Recon{θístilas} `thistle'} therefore supplies the conservative
comparison form: the cluster is medial and the regular development
yields {\emph{þistles}} `thistle', while the simplex nominative belongs
to the broken headword tradition {\emph{þistel}} `thistle'.

\subsection{Paradigm comparison}\label{paradigm-comparison-17}

The comparison below sets the relevant forms side by side. It shows the
contrast between the citation form and the genitive singular cell.

{\def\LTcaptype{none} % do not increment counter
\begin{longtable}[]{@{}
  >{\raggedright\arraybackslash}p{(\linewidth - 8\tabcolsep) * \real{0.2000}}
  >{\raggedright\arraybackslash}p{(\linewidth - 8\tabcolsep) * \real{0.2000}}
  >{\raggedright\arraybackslash}p{(\linewidth - 8\tabcolsep) * \real{0.2000}}
  >{\raggedright\arraybackslash}p{(\linewidth - 8\tabcolsep) * \real{0.2000}}
  >{\raggedright\arraybackslash}p{(\linewidth - 8\tabcolsep) * \real{0.2000}}@{}}
\toprule\noalign{}
\begin{minipage}[b]{\linewidth}\raggedright
PGmc cell / interpretation
\end{minipage} & \begin{minipage}[b]{\linewidth}\raggedright
Candidate input
\end{minipage} & \begin{minipage}[b]{\linewidth}\raggedright
OE output or comparison
\end{minipage} & \begin{minipage}[b]{\linewidth}\raggedright
OE comparison form
\end{minipage} & \begin{minipage}[b]{\linewidth}\raggedright
Result
\end{minipage} \\
\midrule\noalign{}
\endhead
\bottomrule\noalign{}
\endlastfoot
citation nominative singular & {\emph{*θéstilaz}} & computed output:
{\emph{þistl}} & {\emph{þistel}} `thistle' & useful citation-form
background, but not the Old English form here \\
genitive singular & {\emph{*θístilas}} & computed output:
{\emph{þistles}} `thistle' & {\emph{þistles}} `thistle' & exact match
for the conservative cell \\
\end{longtable}
}

\section{\texorpdfstring{make (iptv.2sg) --- OE
\emph{maca}}{make (iptv.2sg) --- OE maca}}\label{make-iptv.2sg-oe-maca}

\index[iv]{02oe@\textbf{Old English}!maca@\emph{maca}}
\index[iv]{03pgmc@\textbf{Proto-Germanic}!mako@*mákô}
\index[iv]{03pgmc@\textbf{Proto-Germanic}!makona@*makōną}

Derivation: citation reconstruction \emph{*makōną}; form followed here
\emph{*mákô} \textgreater{} \emph{maca} (late analogy).

\subsection{Derivation trace}\label{derivation-trace-126}

Proto input: \emph{*mákô}

\begingroup
\setlength{\fboxsep}{6pt}

\noindent\fbox{%
\begin{minipage}{0.97\linewidth}
\small
\begin{tabularx}{\linewidth}{@{}>{\raggedright\arraybackslash}p{0.320\linewidth}>{\centering\arraybackslash}X>{\raggedright\arraybackslash}p{0.520\linewidth}@{}}
\begin{minipage}[t]{\linewidth}
\centering\textbf{Earlier Germanic changes}\par
\vspace{0.35em}
\raggedright
\centering\textbf{}\par
\raggedright
\vspace{0.2em}
\begin{tabular}{@{}>{\raggedright\arraybackslash}p{0.68\linewidth}@{\hspace{0.55em}}>{\raggedright\arraybackslash}p{0.16\linewidth}@{\hspace{0.25em}}}
\multicolumn{2}{@{}l@{}}{*Proto-West Germanic*} \\
\multicolumn{2}{@{}l@{}}{[no change]} \\
\multicolumn{2}{@{}l@{}}{*Northwest Germanic*} \\
\multicolumn{2}{@{}l@{}}{[no change]} \\
\end{tabular}
\end{minipage}
&
&
\begin{minipage}[t]{\linewidth}
\centering\textbf{Old English changes}\par
\vspace{0.35em}
\raggedright
\centering\textbf{}\par
\raggedright
\vspace{0.2em}
\begin{tabular}{@{}>{\raggedright\arraybackslash}p{0.74\linewidth}@{\hspace{0.55em}}>{\raggedright\arraybackslash}p{0.16\linewidth}@{\hspace{0.25em}}}
\multicolumn{2}{@{}l@{}}{*Old English*} \\
Anglo Frisian Brightening & \emph{*mækô} \\
OE A Restoration & \emph{*makô} \\
OE Unstressed Long Vowel Shortening & \emph{*maka} \\
\end{tabular}
\end{minipage}
\\
\end{tabularx}
\end{minipage}%
} \endgroup

Old English form: \emph{maca}

\subsection{Reconstruction and comparative
evidence}\label{reconstruction-and-comparative-evidence-126}

The make-family belongs to the Old English class-II weak verbs. Campbell
cites \emph{lapian,} {\emph{macian}} `make' among verbs with restored
\emph{a} (\citeproc{ref-Campbell1959}{Campbell 1959, sect. 159}). Ringe
and Taylor place the Germanic verb in the same class, comparing West
Germanic continuants such as Old Frisian {\emph{makia}} `make', Old
Saxon {\emph{makon}} `make', and Old High German {\emph{mahhon}} `make'
(\citeproc{ref-RingeTaylor2014}{Ringe and Taylor 2014, 191}).

The derivational input {\Recon{mákô} `make (iptv.2sg)'} is not the
citation form of the lexeme but a finite paradigm cell. Ringe and Taylor
give the class-II weak imperative singular as -a \textless{} \emph{*-ō},
which makes this cell the relevant comparison point for the Old English
form treated here (\citeproc{ref-RingeTaylor2014}{Ringe and Taylor 2014,
314}).

\subsection{Old English evidence}\label{old-english-evidence-126}

The dictionary headword is {\emph{macian}} `make'
(\citeproc{ref-ClarkHall1960}{Clark Hall 1960, 193}). The form compared
here in this entry is therefore not the lemma but the imperative
singular {\emph{maca}} `make', chosen as a paradigm form beside the
headword {\emph{macian}} `make' and the related finite form
{\emph{macaþ}} `makes'.

The lexical history is still that of \emph{macian} `make', but the
finite cell isolates the regular outcome of trimoric \emph{*ō} more
cleanly than the citation form does.

\subsection{Development to Old
English}\label{development-to-old-english-126}

From \Recon{mákô} `make (iptv.2sg)', Anglo-Frisian brightening first
gives \Recon{mækô} `make (iptv.2sg)'. Campbell cites \emph{macian}
`make' among the class-II verbs with restored \emph{a}
(\citeproc{ref-Campbell1959}{Campbell 1959, sect. 159}). Ringe and
Taylor's class-II weak imperative singular -a \textless{} \emph{*-ō}
then supports the later finite ending
(\citeproc{ref-RingeTaylor2014}{Ringe and Taylor 2014, 314}).

The same development explains why earlier fronted forms of the
\emph{mæċa} `make' type do not control the entry. Once trimoric
\emph{*ô} is treated as a back-vocalic trigger for restoration, the
imperative singular falls into line with the broader \emph{macian}
`make' family.

\subsection{Paradigm comparison}\label{paradigm-comparison-18}

A paradigm comparison identifies which finite cell of the make-family
matches the Old English form chosen here. The comparison below sets the
relevant forms side by side.

{\def\LTcaptype{none} % do not increment counter
\begin{longtable}[]{@{}
  >{\raggedright\arraybackslash}p{(\linewidth - 8\tabcolsep) * \real{0.2000}}
  >{\raggedright\arraybackslash}p{(\linewidth - 8\tabcolsep) * \real{0.2000}}
  >{\raggedright\arraybackslash}p{(\linewidth - 8\tabcolsep) * \real{0.2000}}
  >{\raggedright\arraybackslash}p{(\linewidth - 8\tabcolsep) * \real{0.2000}}
  >{\raggedright\arraybackslash}p{(\linewidth - 8\tabcolsep) * \real{0.2000}}@{}}
\toprule\noalign{}
\begin{minipage}[b]{\linewidth}\raggedright
PGmc cell / interpretation
\end{minipage} & \begin{minipage}[b]{\linewidth}\raggedright
Candidate input
\end{minipage} & \begin{minipage}[b]{\linewidth}\raggedright
Old English outcome or comparison
\end{minipage} & \begin{minipage}[b]{\linewidth}\raggedright
OE comparison form
\end{minipage} & \begin{minipage}[b]{\linewidth}\raggedright
Result
\end{minipage} \\
\midrule\noalign{}
\endhead
\bottomrule\noalign{}
\endlastfoot
lexeme-level infinitive & {\emph{*mákōjaną}} & comparative continuation
{\emph{macian}} `make' & {\emph{macian}} `make' & ordinary headword of
the verb, but not the finite form compared here \\
imperative singular & {\emph{*mákô}} & regular output: {\emph{maca}}
`make' & {\emph{maca}} `make' & exact match between input, output, and
selected paradigm form \\
present third singular companion & {\emph{*mákōθi}} & comparative
companion {\emph{macaþ}} & {\emph{macaþ}} & useful family control, but
not the target of this entry \\
\end{longtable}
}

\section{\texorpdfstring{make (3sg) --- OE
\emph{macaþ}}{make (3sg) --- OE macaþ}}\label{make-3sg-oe-macauxfe}

\index[iv]{02oe@\textbf{Old English}!macath@\emph{macaþ}}
\index[iv]{03pgmc@\textbf{Proto-Germanic}!makona@*makōną}
\index[iv]{03pgmc@\textbf{Proto-Germanic}!makothi@*mákōθi}

Derivation: citation reconstruction \emph{*makōną}; form followed here
\emph{*mákōθi} \textgreater{} \emph{macaþ} (late analogy).

\subsection{Derivation trace}\label{derivation-trace-127}

Proto input: \emph{*mákōθi}

\begingroup
\setlength{\fboxsep}{6pt}

\noindent\fbox{%
\begin{minipage}{0.97\linewidth}
\small
\begin{tabularx}{\linewidth}{@{}>{\raggedright\arraybackslash}p{0.440\linewidth}>{\centering\arraybackslash}X>{\raggedright\arraybackslash}p{0.440\linewidth}@{}}
\begin{minipage}[t]{\linewidth}
\centering\textbf{Earlier Germanic changes}\par
\vspace{0.35em}
\raggedright
\centering\textbf{}\par
\raggedright
\vspace{0.2em}
\begin{tabular}{@{}>{\raggedright\arraybackslash}p{0.68\linewidth}@{\hspace{0.55em}}>{\raggedright\arraybackslash}p{0.16\linewidth}@{\hspace{0.25em}}}
\multicolumn{2}{@{}l@{}}{*Proto-West Germanic*} \\
PWGmc Early I Apocope & \emph{*mákōθ} \\
\multicolumn{2}{@{}l@{}}{*Northwest Germanic*} \\
\multicolumn{2}{@{}l@{}}{[no change]} \\
\end{tabular}
\end{minipage}
&
&
\begin{minipage}[t]{\linewidth}
\centering\textbf{Old English changes}\par
\vspace{0.35em}
\raggedright
\centering\textbf{}\par
\raggedright
\vspace{0.2em}
\begin{tabular}{@{}>{\raggedright\arraybackslash}p{0.70\linewidth}@{\hspace{0.55em}}>{\raggedright\arraybackslash}p{0.16\linewidth}@{\hspace{0.25em}}}
\multicolumn{2}{@{}l@{}}{*Old English*} \\
Anglo Frisian Brightening & \emph{*mækōθ} \\
OE A Restoration & \emph{*makōθ} \\
OE Late O Shortening & \emph{*makaθ} \\
\end{tabular}
\end{minipage}
\\
\end{tabularx}
\end{minipage}%
} \endgroup

Old English form: \emph{macaþ}

\subsection{Reconstruction and comparative
evidence}\label{reconstruction-and-comparative-evidence-127}

Kroonen derives the Old English verb from {\emph{*makōjan-}} on the
make-family base \emph{*maka-} (\citeproc{ref-Kroonen2013}{Kroonen 2013,
350}). Ringe and Taylor likewise derive Old English {\emph{macian}}
`make' from PWGmc \Recon{makon} `make' through \Recon{mekojan} `make'
(\citeproc{ref-RingeTaylor2014}{Ringe and Taylor 2014, 191}).

The derivational input {\Recon{mákōθi} `makes'} is therefore a finite
3sg cell of the same family, not the citation form of the verb.

\subsection{Old English evidence}\label{old-english-evidence-127}

Clark Hall lemmatizes the verb as {\emph{macian}} `make'
(\citeproc{ref-ClarkHall1960}{Clark Hall 1960, 193}). The relevant
comparison form here is normalized present-third-singular {\emph{macaþ}}
`makes', set beside the dictionary headword and the related imperative
singular {\emph{maca}} `make'.

Campbell's class-II paradigm gives ordinary 3sg ending \emph{-aþ}, while
his dialect survey allows secondary \emph{-e-} spellings in some
traditions (\citeproc{ref-Campbell1959}{Campbell 1959, sect. 356.4};
\citeproc{ref-Campbell1959}{1959, sect. 757}). {\emph{Macaþ}} `makes' is
thus the regular comparison form for the non-\emph{j} 3sg cell.

\subsection{Development to Old
English}\label{development-to-old-english-127}

After early loss of final \emph{-i}, \Recon{mákōθi} `makes' →
\Recon{mákōθ} `makes' after loss of final -i. Anglo-Frisian brightening
gives \Recon{mækōθ} `makes', but Campbell lists {\emph{macian}} `make'
among class-II verbs with restored \emph{a}, so the stem returns to
\emph{mak-} before the ending is reduced
(\citeproc{ref-Campbell1959}{Campbell 1959, sect. 159}).

The ending then follows ordinary class-II 3sg development: \Recon{makōθ}
`makes' \textgreater{} \Recon{makaθ} `makes' \textgreater{} macaþ\_.
Campbell's\_lufas, -aþ (\textless{} -ōsi, -ōþi)\_and Ringe and Taylor's
discussion of stable class-II\_a\_in finite non-\_j` cells support this
sequence (\citeproc{ref-Campbell1959}{Campbell 1959, sect. 356.4};
\citeproc{ref-RingeTaylor2014}{Ringe and Taylor 2014, 80}).

\subsection{Paradigm comparison}\label{paradigm-comparison-19}

The comparison below sets the relevant forms side by side. It
distinguishes the selected 3sg cell from the make-family lemma and from
the companion imperative form.

{\def\LTcaptype{none} % do not increment counter
\begin{longtable}[]{@{}
  >{\raggedright\arraybackslash}p{(\linewidth - 8\tabcolsep) * \real{0.2000}}
  >{\raggedright\arraybackslash}p{(\linewidth - 8\tabcolsep) * \real{0.2000}}
  >{\raggedright\arraybackslash}p{(\linewidth - 8\tabcolsep) * \real{0.2000}}
  >{\raggedright\arraybackslash}p{(\linewidth - 8\tabcolsep) * \real{0.2000}}
  >{\raggedright\arraybackslash}p{(\linewidth - 8\tabcolsep) * \real{0.2000}}@{}}
\toprule\noalign{}
\begin{minipage}[b]{\linewidth}\raggedright
PGmc cell / interpretation
\end{minipage} & \begin{minipage}[b]{\linewidth}\raggedright
Candidate input
\end{minipage} & \begin{minipage}[b]{\linewidth}\raggedright
OE outcome or comparison
\end{minipage} & \begin{minipage}[b]{\linewidth}\raggedright
OE comparison form
\end{minipage} & \begin{minipage}[b]{\linewidth}\raggedright
Result
\end{minipage} \\
\midrule\noalign{}
\endhead
\bottomrule\noalign{}
\endlastfoot
lexeme-level infinitive & {\emph{*makōjaną}} & dictionary headword
{\emph{macian}} `make' & {\emph{macian}} `make' & family background, but
not the cell compared here \\
3sg present & {\emph{*mákōθi}} & regular output {\emph{macaþ}} `makes' &
{\emph{macaþ}} `makes' & exact match \\
imperative singular companion & {\emph{*mákô}} & related finite form
{\emph{maca}} `make' & {\emph{maca}} `make' & useful control, but not
the target \\
\end{longtable}
}

\section{\texorpdfstring{bore (iptv.2sg) --- OE
\emph{bora}}{bore (iptv.2sg) --- OE bora}}\label{bore-iptv.2sg-oe-bora}

\index[iv]{02oe@\textbf{Old English}!bora@\emph{bora}}
\index[iv]{03pgmc@\textbf{Proto-Germanic}!buro@*búrô}
\index[iv]{03pgmc@\textbf{Proto-Germanic}!burona@*burōną}

Derivation: citation reconstruction \emph{*burōną}; form followed here
\emph{*búrô} \textgreater{} \emph{bora} (late analogy).

\subsection{Derivation trace}\label{derivation-trace-128}

Proto input: \emph{*búrô}

\begingroup
\setlength{\fboxsep}{6pt}

\noindent\fbox{%
\begin{minipage}{0.97\linewidth}
\small
\begin{tabularx}{\linewidth}{@{}>{\raggedright\arraybackslash}p{0.460\linewidth}>{\centering\arraybackslash}X>{\raggedright\arraybackslash}p{0.460\linewidth}@{}}
\begin{minipage}[t]{\linewidth}
\centering\textbf{Earlier Germanic changes}\par
\vspace{0.35em}
\raggedright
\centering\textbf{}\par
\raggedright
\vspace{0.2em}
\begin{tabular}{@{}>{\raggedright\arraybackslash}p{0.74\linewidth}@{\hspace{0.55em}}>{\raggedright\arraybackslash}p{0.16\linewidth}@{\hspace{0.25em}}}
\multicolumn{2}{@{}l@{}}{*Proto-West Germanic*} \\
\multicolumn{2}{@{}l@{}}{[no change]} \\
\multicolumn{2}{@{}l@{}}{*Northwest Germanic*} \\
\mbox{NWGmc U Lowering} & \emph{*bórô} \\
\end{tabular}
\end{minipage}
&
&
\begin{minipage}[t]{\linewidth}
\centering\textbf{Old English changes}\par
\vspace{0.35em}
\raggedright
\centering\textbf{}\par
\raggedright
\vspace{0.2em}
\begin{tabular}{@{}>{\raggedright\arraybackslash}p{0.74\linewidth}@{\hspace{0.55em}}>{\raggedright\arraybackslash}p{0.16\linewidth}@{\hspace{0.25em}}}
\multicolumn{2}{@{}l@{}}{*Old English*} \\
OE Unstressed Long Vowel Shortening & \emph{*bóra} \\
\end{tabular}
\end{minipage}
\\
\end{tabularx}
\end{minipage}%
} \endgroup

Old English form: \emph{bora}

\subsection{Reconstruction and comparative
evidence}\label{reconstruction-and-comparative-evidence-128}

Kroonen reconstructs the bore-family verb as \emph{*burojan-} and cites
Old English {\emph{borian}} `bore' among its continuants
(\citeproc{ref-Kroonen2013}{Kroonen 2013, 85}). Ringe and Taylor give
the class-II weak imperative singular ending as -a \textless{}
\emph{*-ō} (\citeproc{ref-RingeTaylor2014}{Ringe and Taylor 2014, 314}).

The derivational input \Recon{búrô} `bore (iptv.2sg)' is therefore an
imperative cell of the same family, not the citation form of the verb.

\subsection{Old English evidence}\label{old-english-evidence-128}

Clark Hall lemmatizes the verb as {\emph{borian}} `bore'
(\citeproc{ref-ClarkHall1960}{Clark Hall 1960, 48}). The comparison form
here is the normalized imperative singular \emph{bora} `bore
(iptv.2sg)', used beside the headword and the related 3sg form
\emph{boraþ} `bores'.

The imperative is thus a paradigm form rather than a replacement for the
dictionary lemma. It is the most direct Old English comparator for the
non-\emph{j} finite cell represented by \Recon{búrô} `bore (iptv.2sg)'.

\subsection{Development to Old
English}\label{development-to-old-english-128}

Northwest Germanic lowering first gives \Recon{bórô} `bore (iptv.2sg)'
from \Recon{búrô} `bore (iptv.2sg)', and late shortening of the
unstressed long vowel then yields \Recon{bóra} `bore (iptv.2sg)', whence
\emph{bora} `bore'.

Ringe and Taylor's class-II imperative singular -a \textless{}
\emph{*-ō} points to exactly this type of outcome
(\citeproc{ref-RingeTaylor2014}{Ringe and Taylor 2014, 314}). The form
compared here therefore isolates the regular finite-cell development
more cleanly than the remodelled infinitive does.

\subsection{Paradigm comparison}\label{paradigm-comparison-20}

The comparison below sets the relevant forms side by side. It
distinguishes the selected imperative cell from the bore-family lemma
and from the companion 3sg form.

{\def\LTcaptype{none} % do not increment counter
\begin{longtable}[]{@{}
  >{\raggedright\arraybackslash}p{(\linewidth - 8\tabcolsep) * \real{0.2000}}
  >{\raggedright\arraybackslash}p{(\linewidth - 8\tabcolsep) * \real{0.2000}}
  >{\raggedright\arraybackslash}p{(\linewidth - 8\tabcolsep) * \real{0.2000}}
  >{\raggedright\arraybackslash}p{(\linewidth - 8\tabcolsep) * \real{0.2000}}
  >{\raggedright\arraybackslash}p{(\linewidth - 8\tabcolsep) * \real{0.2000}}@{}}
\toprule\noalign{}
\begin{minipage}[b]{\linewidth}\raggedright
PGmc cell / interpretation
\end{minipage} & \begin{minipage}[b]{\linewidth}\raggedright
Candidate input
\end{minipage} & \begin{minipage}[b]{\linewidth}\raggedright
OE outcome or comparison
\end{minipage} & \begin{minipage}[b]{\linewidth}\raggedright
OE comparison form
\end{minipage} & \begin{minipage}[b]{\linewidth}\raggedright
Result
\end{minipage} \\
\midrule\noalign{}
\endhead
\bottomrule\noalign{}
\endlastfoot
lexeme-level infinitive & {\emph{*burōjaną}} & dictionary headword
{\emph{borian}} & {\emph{borian}} & family background, but not the cell
compared here \\
imperative singular & {\emph{*búrô}} & regular output {\emph{bora}}
`bore' & {\emph{bora}} `bore' & exact match \\
3sg present companion & {\emph{*búrōθi}} & related finite form
{\emph{boraþ}} `bores' & {\emph{boraþ}} `bores' & useful control, but
not the target \\
\end{longtable}
}

\section{\texorpdfstring{bore (3sg) --- OE
\emph{boraþ}}{bore (3sg) --- OE boraþ}}\label{bore-3sg-oe-borauxfe}

\index[iv]{02oe@\textbf{Old English}!borath@\emph{boraþ}}
\index[iv]{03pgmc@\textbf{Proto-Germanic}!burona@*burōną}
\index[iv]{03pgmc@\textbf{Proto-Germanic}!burothi@*búrōθi}

Derivation: citation reconstruction \emph{*burōną}; form followed here
\emph{*búrōθi} \textgreater{} \emph{boraþ} (late analogy).

\subsection{Derivation trace}\label{derivation-trace-129}

Proto input: \emph{*búrōθi}

\begingroup
\setlength{\fboxsep}{6pt}

\noindent\fbox{%
\begin{minipage}{0.97\linewidth}
\small
\begin{tabularx}{\linewidth}{@{}>{\raggedright\arraybackslash}p{0.520\linewidth}>{\centering\arraybackslash}X>{\raggedright\arraybackslash}p{0.320\linewidth}@{}}
\begin{minipage}[t]{\linewidth}
\centering\textbf{Earlier Germanic changes}\par
\vspace{0.35em}
\raggedright
\centering\textbf{}\par
\raggedright
\vspace{0.2em}
\begin{tabular}{@{}>{\raggedright\arraybackslash}p{0.74\linewidth}@{\hspace{0.55em}}>{\raggedright\arraybackslash}p{0.16\linewidth}@{\hspace{0.25em}}}
\multicolumn{2}{@{}l@{}}{*Proto-West Germanic*} \\
PWGmc Early I Apocope & \emph{*búrōθ} \\
\multicolumn{2}{@{}l@{}}{*Northwest Germanic*} \\
\mbox{NWGmc U Lowering} & \emph{*bórōθ} \\
\end{tabular}
\end{minipage}
&
&
\begin{minipage}[t]{\linewidth}
\centering\textbf{Old English changes}\par
\vspace{0.35em}
\raggedright
\centering\textbf{}\par
\raggedright
\vspace{0.2em}
\begin{tabular}{@{}>{\raggedright\arraybackslash}p{0.68\linewidth}@{\hspace{0.55em}}>{\raggedright\arraybackslash}p{0.16\linewidth}@{\hspace{0.25em}}}
\multicolumn{2}{@{}l@{}}{*Old English*} \\
OE Late O Shortening & \emph{*bóraθ} \\
\end{tabular}
\end{minipage}
\\
\end{tabularx}
\end{minipage}%
} \endgroup

Old English form: \emph{boraþ}

\subsection{Reconstruction and comparative
evidence}\label{reconstruction-and-comparative-evidence-129}

Kroonen reconstructs the bore-family verb as \emph{*burojan-} and cites
Old English {\emph{borian}} `bore' among its reflexes
(\citeproc{ref-Kroonen2013}{Kroonen 2013, 85}). The form compared here
isolates the finite 3sg cell \Recon{búrōθi} `bores' rather than the
infinitive.

Campbell's class-II pattern lufas, \emph{-aþ} (\textless{} \emph{-ōsi},
\emph{-ōþi)} and Ringe and Taylor's account of stable class-II \emph{a}
in finite forms make this 3sg cell the relevant comparison form for the
ending (\citeproc{ref-Campbell1959}{Campbell 1959, sect. 356.4};
\citeproc{ref-RingeTaylor2014}{Ringe and Taylor 2014, 80}).

\subsection{Old English evidence}\label{old-english-evidence-129}

Clark Hall lemmatizes the verb as {\emph{borian}} `bore'; the imperative
singular \emph{bora} and present-third-singular {\emph{boraþ}} `bores'
are the relevant comparison forms (\citeproc{ref-ClarkHall1960}{Clark
Hall 1960, 48}).

Campbell's dialect survey allows secondary \emph{-e-} and \emph{-o-}
spellings in 2sg and 3sg class-II forms, but the basic ending remains
\emph{-aþ} (\citeproc{ref-Campbell1959}{Campbell 1959, sect. 757}). The
regular comparison form for this non-\emph{j} 3sg cell is {\emph{boraþ}}
`bores'.

\subsection{Development to Old
English}\label{development-to-old-english-129}

From \Recon{búrōθi} `bores', loss of final \emph{-i} first gives
\Recon{búrōθ} `bores'. Northwest Germanic lowering then produces
\Recon{bórōθ} `bores', and late shortening of unstressed \emph{ō} yields
\Recon{bóraθ} `bores'. {\emph{Boraþ}} `bores' is therefore the regular
comparison form for this non-\emph{j} 3sg cell.

Campbell's class-II ending evidence and Ringe and Taylor's discussion of
stable \emph{a} in finite non-\emph{j} cells support this sequence
(\citeproc{ref-Campbell1959}{Campbell 1959, sect. 356.4};
\citeproc{ref-RingeTaylor2014}{Ringe and Taylor 2014, 80}).

\subsection{Paradigm comparison}\label{paradigm-comparison-21}

The comparison below sets the relevant forms side by side. It
distinguishes the selected 3sg cell from the bore-family lemma and from
the companion imperative form.

{\def\LTcaptype{none} % do not increment counter
\begin{longtable}[]{@{}
  >{\raggedright\arraybackslash}p{(\linewidth - 8\tabcolsep) * \real{0.2000}}
  >{\raggedright\arraybackslash}p{(\linewidth - 8\tabcolsep) * \real{0.2000}}
  >{\raggedright\arraybackslash}p{(\linewidth - 8\tabcolsep) * \real{0.2000}}
  >{\raggedright\arraybackslash}p{(\linewidth - 8\tabcolsep) * \real{0.2000}}
  >{\raggedright\arraybackslash}p{(\linewidth - 8\tabcolsep) * \real{0.2000}}@{}}
\toprule\noalign{}
\begin{minipage}[b]{\linewidth}\raggedright
PGmc cell / interpretation
\end{minipage} & \begin{minipage}[b]{\linewidth}\raggedright
Candidate input
\end{minipage} & \begin{minipage}[b]{\linewidth}\raggedright
OE outcome or comparison
\end{minipage} & \begin{minipage}[b]{\linewidth}\raggedright
OE comparison form
\end{minipage} & \begin{minipage}[b]{\linewidth}\raggedright
Result
\end{minipage} \\
\midrule\noalign{}
\endhead
\bottomrule\noalign{}
\endlastfoot
lexeme-level infinitive & {\emph{*burōjaną}} & dictionary headword
{\emph{borian}} & {\emph{borian}} & family background, but not the cell
compared here \\
3sg present & {\emph{*búrōθi}} & regular output {\emph{boraþ}} `bores' &
{\emph{boraþ}} `bores' & exact match \\
imperative singular companion & {\emph{*búrô}} & related finite form
{\emph{bora}} `bore' & {\emph{bora}} `bore' & useful control, but not
the target \\
\end{longtable}
}

\section{\texorpdfstring{learn (iptv.2sg) --- OE
\emph{liorna}}{learn (iptv.2sg) --- OE liorna}}\label{learn-iptv.2sg-oe-liorna}

\index[iv]{02oe@\textbf{Old English}!liorna@\emph{liorna}}
\index[iv]{03pgmc@\textbf{Proto-Germanic}!lizno@*líznô}
\index[iv]{03pgmc@\textbf{Proto-Germanic}!liznojana@*liznōjaną}

Derivation: citation reconstruction \emph{*liznōjaną}; form followed
here \emph{*líznô} \textgreater{} \emph{liorna} (late analogy).

\subsection{Derivation trace}\label{derivation-trace-130}

Proto input: \emph{*líznô}

\begingroup
\setlength{\fboxsep}{6pt}

\noindent\fbox{%
\begin{minipage}{0.97\linewidth}
\small
\begin{tabularx}{\linewidth}{@{}>{\raggedright\arraybackslash}p{0.440\linewidth}>{\centering\arraybackslash}X>{\raggedright\arraybackslash}p{0.440\linewidth}@{}}
\begin{minipage}[t]{\linewidth}
\centering\textbf{Earlier Germanic changes}\par
\vspace{0.35em}
\raggedright
\centering\textbf{}\par
\raggedright
\vspace{0.2em}
\begin{tabular}{@{}>{\raggedright\arraybackslash}p{0.68\linewidth}@{\hspace{0.55em}}>{\raggedright\arraybackslash}p{0.16\linewidth}@{\hspace{0.25em}}}
\multicolumn{2}{@{}l@{}}{*Proto-West Germanic*} \\
\multicolumn{2}{@{}l@{}}{[no change]} \\
\multicolumn{2}{@{}l@{}}{*Northwest Germanic*} \\
\multicolumn{2}{@{}l@{}}{[no change]} \\
\end{tabular}
\end{minipage}
&
&
\begin{minipage}[t]{\linewidth}
\centering\textbf{Old English changes}\par
\vspace{0.35em}
\raggedright
\centering\textbf{}\par
\raggedright
\vspace{0.2em}
\begin{tabular}{@{}>{\raggedright\arraybackslash}p{0.74\linewidth}@{\hspace{0.55em}}>{\raggedright\arraybackslash}p{0.16\linewidth}@{\hspace{0.25em}}}
\multicolumn{2}{@{}l@{}}{*Old English*} \\
OE Breaking & \emph{*líornô} \\
OE Unstressed Long Vowel Shortening & \emph{*líorna} \\
\end{tabular}
\end{minipage}
\\
\end{tabularx}
\end{minipage}%
} \endgroup

Old English form: \emph{liorna}

\subsection{Reconstruction and comparative
evidence}\label{reconstruction-and-comparative-evidence-130}

Ringe and Taylor give Old English {\emph{liornian}} `learn' and
{\emph{leornian}} `learn' from a learn-family base of the \emph{*lizn-}
type (\citeproc{ref-RingeTaylor2014}{Ringe and Taylor 2014, 38}), and
Kroonen likewise keeps the weak verb as \emph{*liznōn-}
(\citeproc{ref-Kroonen2013}{Kroonen 2013, 380}). Fulk cites the same Old
English family from \emph{*liznō-} (\citeproc{ref-Fulk2018}{Fulk 2018,
127}).

The derivational input {\Recon{líznô} `learn (iptv.2sg)'} is a finite
imperative cell of that family, not the citation form of the verb.

\subsection{Old English evidence}\label{old-english-evidence-130}

Clark Hall gives the ordinary headword as \emph{leornian}
(\citeproc{ref-ClarkHall1960}{Clark Hall 1960, 186}). Brunner explicitly
records \emph{leornian, nordh. auch liorna}, and Campbell notes
Northumbrian forms with \emph{io} beside \emph{leornian}
(\citeproc{ref-SieversBrunner1965}{Brunner 1965, sect. 417} Anm. 10;
\citeproc{ref-Campbell1959}{Campbell 1959, sect. 123} n.~2).

{\emph{Liorna}} can therefore be treated as an attested Northumbrian
finite form, while {\emph{leornian}} `learn' remains the better-known
dictionary headword.

\subsection{Development to Old
English}\label{development-to-old-english-130}

The form compared here develops regularly as \Recon{líznô} `learn'
\textgreater{} \Recon{lírnô} `learn' by rhotacism, then \Recon{líornô}
`learn (iptv.2sg)' by breaking before \emph{rn}, and finally
\Recon{líorna} `learn' \textgreater{} \emph{liorna} by late shortening
of the unstressed long vowel.

Campbell's Northumbrian \emph{io} evidence and Ringe and Taylor's
explicit statement that no form of \emph{liornian} stood in an
i-umlauting environment support this stem shape
(\citeproc{ref-Campbell1959}{Campbell 1959, sect. 123} n.~2;
\citeproc{ref-RingeTaylor2014}{Ringe and Taylor 2014, 247}). The
West-Saxon-looking \emph{eo} forms belong to a different dialectal
presentation of the same family.

\subsection{Paradigm comparison}\label{paradigm-comparison-22}

The comparison below sets the relevant forms side by side. It
distinguishes the selected imperative cell from the learn-family
infinitive and from the companion 3sg form.

{\def\LTcaptype{none} % do not increment counter
\begin{longtable}[]{@{}
  >{\raggedright\arraybackslash}p{(\linewidth - 8\tabcolsep) * \real{0.2000}}
  >{\raggedright\arraybackslash}p{(\linewidth - 8\tabcolsep) * \real{0.2000}}
  >{\raggedright\arraybackslash}p{(\linewidth - 8\tabcolsep) * \real{0.2000}}
  >{\raggedright\arraybackslash}p{(\linewidth - 8\tabcolsep) * \real{0.2000}}
  >{\raggedright\arraybackslash}p{(\linewidth - 8\tabcolsep) * \real{0.2000}}@{}}
\toprule\noalign{}
\begin{minipage}[b]{\linewidth}\raggedright
PGmc cell / interpretation
\end{minipage} & \begin{minipage}[b]{\linewidth}\raggedright
Candidate input
\end{minipage} & \begin{minipage}[b]{\linewidth}\raggedright
OE outcome or comparison
\end{minipage} & \begin{minipage}[b]{\linewidth}\raggedright
OE comparison form
\end{minipage} & \begin{minipage}[b]{\linewidth}\raggedright
Result
\end{minipage} \\
\midrule\noalign{}
\endhead
\bottomrule\noalign{}
\endlastfoot
lexeme-level infinitive & {\emph{*liznōjaną}} & Northumbrian
{\emph{liornian}} `learn'; dictionary headword often {\emph{leornian}}
`learn' & {\emph{liornian}} `learn' / {\emph{leornian}} `learn' & family
background, but not the cell compared here \\
imperative singular & {\emph{*líznô}} & regular output and Brunner's
Northumbrian {\emph{liorna}} & {\emph{liorna}} & exact match \\
3sg present companion & {\emph{*líznōθi}} & related finite form
{\emph{liornaþ}} & {\emph{liornaþ}} & useful control, but not the
target \\
\end{longtable}
}

\section{\texorpdfstring{learn (3sg) --- OE
\emph{liornaþ}}{learn (3sg) --- OE liornaþ}}\label{learn-3sg-oe-liornauxfe}

\index[iv]{02oe@\textbf{Old English}!liornath@\emph{liornaþ}}
\index[iv]{03pgmc@\textbf{Proto-Germanic}!liznojana@*liznōjaną}
\index[iv]{03pgmc@\textbf{Proto-Germanic}!liznothi@*líznōθi}

Derivation: citation reconstruction \emph{*liznōjaną}; form followed
here \emph{*líznōθi} \textgreater{} \emph{liornaþ} (late analogy).

\subsection{Derivation trace}\label{derivation-trace-131}

Proto input: \emph{*líznōθi}

\begingroup
\setlength{\fboxsep}{6pt}

\noindent\fbox{%
\begin{minipage}{0.97\linewidth}
\small
\begin{tabularx}{\linewidth}{@{}>{\raggedright\arraybackslash}p{0.440\linewidth}>{\centering\arraybackslash}X>{\raggedright\arraybackslash}p{0.440\linewidth}@{}}
\begin{minipage}[t]{\linewidth}
\centering\textbf{Earlier Germanic changes}\par
\vspace{0.35em}
\raggedright
\centering\textbf{}\par
\raggedright
\vspace{0.2em}
\begin{tabular}{@{}>{\raggedright\arraybackslash}p{0.68\linewidth}@{\hspace{0.55em}}>{\raggedright\arraybackslash}p{0.16\linewidth}@{\hspace{0.25em}}}
\multicolumn{2}{@{}l@{}}{*Proto-West Germanic*} \\
PWGmc Early I Apocope & \emph{*lírnōθ} \\
\multicolumn{2}{@{}l@{}}{*Northwest Germanic*} \\
\multicolumn{2}{@{}l@{}}{[no change]} \\
\end{tabular}
\end{minipage}
&
&
\begin{minipage}[t]{\linewidth}
\centering\textbf{Old English changes}\par
\vspace{0.35em}
\raggedright
\centering\textbf{}\par
\raggedright
\vspace{0.2em}
\begin{tabular}{@{}>{\raggedright\arraybackslash}p{0.68\linewidth}@{\hspace{0.55em}}>{\raggedright\arraybackslash}p{0.20\linewidth}@{\hspace{0.25em}}}
\multicolumn{2}{@{}l@{}}{*Old English*} \\
OE Breaking & \emph{*líornōθ} \\
OE Late O Shortening & \emph{*líornaθ} \\
\end{tabular}
\end{minipage}
\\
\end{tabularx}
\end{minipage}%
} \endgroup

Old English form: \emph{liornaþ}

\subsection{Reconstruction and comparative
evidence}\label{reconstruction-and-comparative-evidence-131}

Ringe and Taylor give Old English {\emph{liornian}} `learn' and
{\emph{leornian}} `learn' from a learn-family base of the \emph{*lizn-}
type (\citeproc{ref-RingeTaylor2014}{Ringe and Taylor 2014, 38}), and
Kroonen likewise keeps the weak verb as \emph{*liznōn-}
(\citeproc{ref-Kroonen2013}{Kroonen 2013, 380}). The derivational input
{\Recon{líznōθi} `learns'} is the finite 3sg cell of that family, not
the citation form of the verb.

For the ending, Campbell's class-II \emph{-aþ} (\textless{} \emph{-ōþi)}
and Ringe and Taylor's discussion of stable class-II \emph{a} in finite
forms make the non-\emph{j} 3sg cell the relevant comparison point
(\citeproc{ref-Campbell1959}{Campbell 1959, sect. 356.4};
\citeproc{ref-RingeTaylor2014}{Ringe and Taylor 2014, 80}).

\subsection{Old English evidence}\label{old-english-evidence-131}

Clark Hall gives the ordinary headword as {\emph{leornian}} `learn'
(\citeproc{ref-ClarkHall1960}{Clark Hall 1960, 186}). Brunner records
Northumbrian finite forms in \emph{liorn-}, including {\emph{liorna}}
`learn' and {\emph{liornes}} `learning', beside the West-Saxon-oriented
{\emph{leornian}} tradition (\citeproc{ref-SieversBrunner1965}{Brunner
1965, sect. 417} Anm. 10). Campbell likewise notes Northumbrian forms
with \emph{io} beside \emph{leornian}
(\citeproc{ref-Campbell1959}{Campbell 1959, sect. 123} n.~2).

The relevant comparison form here is the normalized 3sg {\emph{liornaþ}}
`learns'. The directly cited Old English evidence supports the finite
stem \emph{liorn-}; the \emph{-aþ} ending follows the regular class-II
3sg pattern.

\subsection{Development to Old
English}\label{development-to-old-english-131}

The form compared here develops as \Recon{líznōθi} `learns'
\textgreater{} \Recon{lírnōθi} `learns' by rhotacism, then
\Recon{lírnōθ} `learns' after early apocope of final \emph{-i}, then
\Recon{líornōθ} `learns' by breaking before \emph{rn}, and finally
\Recon{líornaθ} `learns' \textgreater{} {\emph{liornaþ}} `learns' by
late shortening of the unstressed long vowel.

Campbell's Northumbrian \emph{io} evidence and Ringe and Taylor's
statement that no form of \emph{liornian} stood in an i-umlauting
environment support the stem, while Campbell's class-II ending evidence
supports the final \emph{-aþ} (\citeproc{ref-Campbell1959}{Campbell
1959, sect. 123} n.~2; \citeproc{ref-Campbell1959}{Campbell 1959, sect.
356.4}; \citeproc{ref-RingeTaylor2014}{Ringe and Taylor 2014, 247}).

\subsection{Paradigm comparison}\label{paradigm-comparison-23}

The comparison below sets the relevant forms side by side. It
distinguishes the selected 3sg cell from the learn-family infinitive and
from the companion imperative form.

{\def\LTcaptype{none} % do not increment counter
\begin{longtable}[]{@{}
  >{\raggedright\arraybackslash}p{(\linewidth - 8\tabcolsep) * \real{0.2000}}
  >{\raggedright\arraybackslash}p{(\linewidth - 8\tabcolsep) * \real{0.2000}}
  >{\raggedright\arraybackslash}p{(\linewidth - 8\tabcolsep) * \real{0.2000}}
  >{\raggedright\arraybackslash}p{(\linewidth - 8\tabcolsep) * \real{0.2000}}
  >{\raggedright\arraybackslash}p{(\linewidth - 8\tabcolsep) * \real{0.2000}}@{}}
\toprule\noalign{}
\begin{minipage}[b]{\linewidth}\raggedright
PGmc cell / interpretation
\end{minipage} & \begin{minipage}[b]{\linewidth}\raggedright
Candidate input
\end{minipage} & \begin{minipage}[b]{\linewidth}\raggedright
OE outcome or comparison
\end{minipage} & \begin{minipage}[b]{\linewidth}\raggedright
OE comparison form
\end{minipage} & \begin{minipage}[b]{\linewidth}\raggedright
Result
\end{minipage} \\
\midrule\noalign{}
\endhead
\bottomrule\noalign{}
\endlastfoot
lexeme-level infinitive & {\emph{*liznōjaną}} & Northumbrian
{\emph{liornian}} `learn'; dictionary headword often {\emph{leornian}}
`learn' & {\emph{liornian}} `learn' / {\emph{leornian}} `learn' & family
background, but not the cell compared here \\
3sg present & {\emph{*líznōθi}} & regular output {\emph{liornaþ}}
`learns' & {\emph{liornaþ}} `learns' & exact match \\
imperative singular companion & {\emph{*líznô}} & regular output and
Brunner's Northumbrian {\emph{liorna}} & {\emph{liorna}} & useful
control, but not the target \\
\end{longtable}
}

\section{\texorpdfstring{lick (iptv.2sg) --- OE
\emph{licca}}{lick (iptv.2sg) --- OE licca}}\label{lick-iptv.2sg-oe-licca}

\index[iv]{02oe@\textbf{Old English}!licca@\emph{licca}}
\index[iv]{03pgmc@\textbf{Proto-Germanic}!likko@*líkkô}
\index[iv]{03pgmc@\textbf{Proto-Germanic}!likkona@*likkōną}

Derivation: citation reconstruction \emph{*likkōną}; form followed here
\emph{*líkkô} \textgreater{} \emph{licca} (late analogy).

\subsection{Derivation trace}\label{derivation-trace-132}

Proto input: \emph{*líkkô}

\begingroup
\setlength{\fboxsep}{6pt}

\noindent\fbox{%
\begin{minipage}{0.97\linewidth}
\small
\begin{tabularx}{\linewidth}{@{}>{\raggedright\arraybackslash}p{0.440\linewidth}>{\centering\arraybackslash}X>{\raggedright\arraybackslash}p{0.440\linewidth}@{}}
\begin{minipage}[t]{\linewidth}
\centering\textbf{Earlier Germanic changes}\par
\vspace{0.35em}
\raggedright
\centering\textbf{}\par
\raggedright
\vspace{0.2em}
\begin{tabular}{@{}>{\raggedright\arraybackslash}p{0.68\linewidth}@{\hspace{0.55em}}>{\raggedright\arraybackslash}p{0.16\linewidth}@{\hspace{0.25em}}}
\multicolumn{2}{@{}l@{}}{*Proto-West Germanic*} \\
\multicolumn{2}{@{}l@{}}{[no change]} \\
\multicolumn{2}{@{}l@{}}{*Northwest Germanic*} \\
\multicolumn{2}{@{}l@{}}{[no change]} \\
\end{tabular}
\end{minipage}
&
&
\begin{minipage}[t]{\linewidth}
\centering\textbf{Old English changes}\par
\vspace{0.35em}
\raggedright
\centering\textbf{}\par
\raggedright
\vspace{0.2em}
\begin{tabular}{@{}>{\raggedright\arraybackslash}p{0.74\linewidth}@{\hspace{0.55em}}>{\raggedright\arraybackslash}p{0.16\linewidth}@{\hspace{0.25em}}}
\multicolumn{2}{@{}l@{}}{*Old English*} \\
OE Unstressed Long Vowel Shortening & \emph{*líkka} \\
\end{tabular}
\end{minipage}
\\
\end{tabularx}
\end{minipage}%
} \endgroup

Old English form: \emph{licca}

\subsection{Reconstruction and comparative
evidence}\label{reconstruction-and-comparative-evidence-132}

Ringe and Taylor give PWGmc \Recon{ekkōn} `lick (3sg.)' continuing as
Old English {\emph{liccian}} `lick', Old Saxon {\emph{likkon}} `lick',
and Old High German {\emph{lecchon}} `lick'
(\citeproc{ref-RingeTaylor2014}{Ringe and Taylor 2014, 50}). Orel gives
the fuller weak-verb reconstruction {\Recon{likkōjanan} `lick'} with the
same Old English continuation (\citeproc{ref-Orel2003}{Orel 2003, 285}).

Campbell's weak class-II discussion gives present forms such as lufas,
\emph{-aþ} (\textless{} \emph{-ōsi}, \emph{-ōþi)}
(\citeproc{ref-Campbell1959}{Campbell 1959, sect. 356.4}). Ringe and
Taylor likewise note that class-II weak present 2sg. \emph{-as(t)} and
3sg. \emph{-aþ} have stable \emph{a}
(\citeproc{ref-RingeTaylor2014}{Ringe and Taylor 2014, 80}). The form
treated here is therefore not that remodeled infinitive but a finite
cell in bare trimoric \emph{*-ō}.

\subsection{Old English evidence}\label{old-english-evidence-132}

Bosworth-Toller lemmatizes the verb as \emph{liccian} `lick'
(\citeproc{ref-BosworthToller1898}{Bosworth and Toller 1898, 614}).
Campbell cites \emph{liccian} `lick' among Old English forms with
preserved geminate \emph{cc} (\citeproc{ref-Campbell1959}{Campbell 1959,
sect. 398.1}). Brunner likewise cites \emph{liccian} `lick'
(\citeproc{ref-SieversBrunner1965}{Brunner 1965, sect. 45} Anm. 3). The
Old English evidence therefore establishes the verbal headword and its
consonantal frame securely.

The Old English form here in this entry is the imperative singular
{\emph{licca}} `lick'. It is a paradigm form chosen beside the headword
{\emph{liccian}} `lick' and the related present {\emph{liccaþ}} `licks',
not a separately lemmatized citation word.

\subsection{Development to Old
English}\label{development-to-old-english-132}

With the stem \emph{licc-} established, the remaining development is
brief. Campbell's class-II present endings lufas, \emph{-aþ}
(\textless{} \emph{-ōsi}, \emph{-ōþi)} support late \emph{-a} in this
finite cell (\citeproc{ref-Campbell1959}{Campbell 1959, sect. 356.4}).
Ringe and Taylor likewise note stable \emph{a} in the class-II 2sg and
3sg (\citeproc{ref-RingeTaylor2014}{Ringe and Taylor 2014, 80}). The
same stem consonantism that appears in \emph{liccian} `lick' is
preserved here, giving \emph{cc} throughout the finite form.

\subsection{Paradigm comparison}\label{paradigm-comparison-24}

The comparison below sets the relevant forms side by side.

{\def\LTcaptype{none} % do not increment counter
\begin{longtable}[]{@{}
  >{\raggedright\arraybackslash}p{(\linewidth - 8\tabcolsep) * \real{0.2000}}
  >{\raggedright\arraybackslash}p{(\linewidth - 8\tabcolsep) * \real{0.2000}}
  >{\raggedright\arraybackslash}p{(\linewidth - 8\tabcolsep) * \real{0.2000}}
  >{\raggedright\arraybackslash}p{(\linewidth - 8\tabcolsep) * \real{0.2000}}
  >{\raggedright\arraybackslash}p{(\linewidth - 8\tabcolsep) * \real{0.2000}}@{}}
\toprule\noalign{}
\begin{minipage}[b]{\linewidth}\raggedright
PGmc cell / interpretation
\end{minipage} & \begin{minipage}[b]{\linewidth}\raggedright
Candidate input
\end{minipage} & \begin{minipage}[b]{\linewidth}\raggedright
Old English outcome or comparison
\end{minipage} & \begin{minipage}[b]{\linewidth}\raggedright
OE comparison form
\end{minipage} & \begin{minipage}[b]{\linewidth}\raggedright
Result
\end{minipage} \\
\midrule\noalign{}
\endhead
\bottomrule\noalign{}
\endlastfoot
lexeme-level infinitive & {\emph{*líkkōjaną}} & regular output
{\emph{liccian}} `lick' & {\emph{liccian}} `lick' & ordinary dictionary
headword of the verb, but not the finite form compared here \\
imperative singular & {\emph{*líkkô}} & regular output {\emph{licca}}
`lick' & {\emph{licca}} `lick' & exact match between the derivational
input and the Old English form here \\
present third singular companion & {\emph{*líkkōθi}} & regular output
{\emph{liccaþ}} & {\emph{liccaþ}} & useful family control, but not the
target of this entry \\
\end{longtable}
}

\section{\texorpdfstring{lick (3sg) --- OE
\emph{liccaþ}}{lick (3sg) --- OE liccaþ}}\label{lick-3sg-oe-liccauxfe}

\index[iv]{02oe@\textbf{Old English}!liccath@\emph{liccaþ}}
\index[iv]{03pgmc@\textbf{Proto-Germanic}!likkona@*likkōną}
\index[iv]{03pgmc@\textbf{Proto-Germanic}!likkothi@*líkkōθi}

Derivation: citation reconstruction \emph{*likkōną}; form followed here
\emph{*líkkōθi} \textgreater{} \emph{liccaþ} (late analogy).

\subsection{Derivation trace}\label{derivation-trace-133}

Proto input: \emph{*líkkōθi}

\begingroup
\setlength{\fboxsep}{6pt}

\noindent\fbox{%
\begin{minipage}{0.97\linewidth}
\small
\begin{tabularx}{\linewidth}{@{}>{\raggedright\arraybackslash}p{0.440\linewidth}>{\centering\arraybackslash}X>{\raggedright\arraybackslash}p{0.440\linewidth}@{}}
\begin{minipage}[t]{\linewidth}
\centering\textbf{Earlier Germanic changes}\par
\vspace{0.35em}
\raggedright
\centering\textbf{}\par
\raggedright
\vspace{0.2em}
\begin{tabular}{@{}>{\raggedright\arraybackslash}p{0.68\linewidth}@{\hspace{0.55em}}>{\raggedright\arraybackslash}p{0.16\linewidth}@{\hspace{0.25em}}}
\multicolumn{2}{@{}l@{}}{*Proto-West Germanic*} \\
PWGmc Early I Apocope & \emph{*líkkōθ} \\
\multicolumn{2}{@{}l@{}}{*Northwest Germanic*} \\
\multicolumn{2}{@{}l@{}}{[no change]} \\
\end{tabular}
\end{minipage}
&
&
\begin{minipage}[t]{\linewidth}
\centering\textbf{Old English changes}\par
\vspace{0.35em}
\raggedright
\centering\textbf{}\par
\raggedright
\vspace{0.2em}
\begin{tabular}{@{}>{\raggedright\arraybackslash}p{0.68\linewidth}@{\hspace{0.55em}}>{\raggedright\arraybackslash}p{0.16\linewidth}@{\hspace{0.25em}}}
\multicolumn{2}{@{}l@{}}{*Old English*} \\
OE Late O Shortening & \emph{*líkkaθ} \\
\end{tabular}
\end{minipage}
\\
\end{tabularx}
\end{minipage}%
} \endgroup

Old English form: \emph{liccaþ}

\subsection{Reconstruction and comparative
evidence}\label{reconstruction-and-comparative-evidence-133}

Ringe and Taylor give PWGmc \Recon{ekkōn} `lick (3sg.)' continuing as
Old English {\emph{liccian}} `lick', Old Saxon {\emph{likkon}} `lick',
and Old High German {\emph{lecchon}} `lick'
(\citeproc{ref-RingeTaylor2014}{Ringe and Taylor 2014, 50}). Orel gives
the fuller weak-verb reconstruction {\Recon{likkōjanan} `lick'} with the
same Old English continuation (\citeproc{ref-Orel2003}{Orel 2003, 285}).

The form compared here in this entry is the non-\emph{j} present third
singular \Recon{líkkōθi} `licks', not the remodeled infinitive. Campbell
states the class-II present endings as lufas, \emph{-aþ} (\textless{}
\emph{-ōsi}, \emph{-ōþi)} (\citeproc{ref-Campbell1959}{Campbell 1959,
sect. 356.4}). Ringe and Taylor likewise note that class-II weak present
2sg. \emph{-as(t)} and 3sg. \emph{-aþ} have stable \emph{a}
(\citeproc{ref-RingeTaylor2014}{Ringe and Taylor 2014, 80}).

\subsection{Old English evidence}\label{old-english-evidence-133}

Bosworth-Toller lemmatizes the verb as \emph{liccian} `lick'
(\citeproc{ref-BosworthToller1898}{Bosworth and Toller 1898, 614}). The
same consonantal frame appears in Campbell's and Brunner's grammatical
citations of \emph{liccian} `lick' (\citeproc{ref-Campbell1959}{Campbell
1959, sect. 398.1}; \citeproc{ref-SieversBrunner1965}{Brunner 1965,
sect. 45} Anm. 3). The Old English headword is therefore clear even
though the entry here is not about the citation form.

The form treated here is the present third singular {\emph{liccaþ}}
`licks'. It is a selected paradigm form beside the lemma
{\emph{liccian}} `lick' and the related imperative {\emph{licca}}
`lick', not a separately lemmatized headword.

\subsection{Development to Old
English}\label{development-to-old-english-133}

{\Recon{líkkōθi} `licks'} first loses final \emph{-i}, giving
\Recon{líkkōθ} `licks'. Campbell's class-II present endings lufas,
\emph{-aþ} (\textless{} \emph{-ōsi}, \emph{-ōþi)} support the regular
3sg outcome \emph{-aþ} (\citeproc{ref-Campbell1959}{Campbell 1959, sect.
356.4}). Ringe and Taylor likewise note stable \emph{a} in the class-II
2sg and 3sg (\citeproc{ref-RingeTaylor2014}{Ringe and Taylor 2014, 80}).
Because this ending never contains \emph{-j-}, the form does not pass
through an i-umlauted \emph{-eþ} stage.

\subsection{Paradigm comparison}\label{paradigm-comparison-25}

The comparison below sets the relevant forms side by side.

{\def\LTcaptype{none} % do not increment counter
\begin{longtable}[]{@{}
  >{\raggedright\arraybackslash}p{(\linewidth - 8\tabcolsep) * \real{0.2000}}
  >{\raggedright\arraybackslash}p{(\linewidth - 8\tabcolsep) * \real{0.2000}}
  >{\raggedright\arraybackslash}p{(\linewidth - 8\tabcolsep) * \real{0.2000}}
  >{\raggedright\arraybackslash}p{(\linewidth - 8\tabcolsep) * \real{0.2000}}
  >{\raggedright\arraybackslash}p{(\linewidth - 8\tabcolsep) * \real{0.2000}}@{}}
\toprule\noalign{}
\begin{minipage}[b]{\linewidth}\raggedright
PGmc cell / interpretation
\end{minipage} & \begin{minipage}[b]{\linewidth}\raggedright
Candidate input
\end{minipage} & \begin{minipage}[b]{\linewidth}\raggedright
Old English outcome or comparison
\end{minipage} & \begin{minipage}[b]{\linewidth}\raggedright
OE comparison form
\end{minipage} & \begin{minipage}[b]{\linewidth}\raggedright
Result
\end{minipage} \\
\midrule\noalign{}
\endhead
\bottomrule\noalign{}
\endlastfoot
lexeme-level infinitive & {\emph{*líkkōjaną}} & regular output
{\emph{liccian}} `lick' & {\emph{liccian}} `lick' & ordinary dictionary
headword of the verb, but not the finite form compared here \\
imperative singular companion & {\emph{*líkkô}} & regular output
{\emph{licca}} `lick' & {\emph{licca}} `lick' & useful family control,
but not the target of this entry \\
present third singular & {\emph{*líkkōθi}} & regular output
{\emph{liccaþ}} `licks' & {\emph{liccaþ}} `licks' & exact match between
the derivational input and the Old English form here \\
\end{longtable}
}

\section{\texorpdfstring{show (iptv.2sg) --- OE
\emph{sċēawa}}{show (iptv.2sg) --- OE sċēawa}}\label{show-iptv.2sg-oe-sux10bux113awa}

\index[iv]{02oe@\textbf{Old English}!sceawa@\emph{sċēawa}}
\index[iv]{03pgmc@\textbf{Proto-Germanic}!skawo@*skáwô}
\index[iv]{03pgmc@\textbf{Proto-Germanic}!skawona@*skawōną}
\index[iv]{09ohg@\textbf{Old High German}!scouwon@scouwōn}
\index[iv]{10ofris@\textbf{Old Frisian}!skawia@skawia}
\index[iv]{12os@\textbf{Old Saxon}!skawon@skawōn}

Derivation: citation reconstruction \emph{*skawōną}; form followed here
\emph{*skáwô} \textgreater{} \emph{sċēawa} (late analogy).

\subsection{Derivation trace}\label{derivation-trace-134}

Proto input: \emph{*skáwô}

\begingroup
\setlength{\fboxsep}{6pt}

\noindent\fbox{%
\begin{minipage}{0.97\linewidth}
\small
\begin{tabularx}{\linewidth}{@{}>{\raggedright\arraybackslash}p{0.320\linewidth}>{\centering\arraybackslash}X>{\raggedright\arraybackslash}p{0.520\linewidth}@{}}
\begin{minipage}[t]{\linewidth}
\centering\textbf{Earlier Germanic changes}\par
\vspace{0.35em}
\raggedright
\centering\textbf{}\par
\raggedright
\vspace{0.2em}
\begin{tabular}{@{}>{\raggedright\arraybackslash}p{0.68\linewidth}@{\hspace{0.55em}}>{\raggedright\arraybackslash}p{0.16\linewidth}@{\hspace{0.25em}}}
\multicolumn{2}{@{}l@{}}{*Proto-West Germanic*} \\
\multicolumn{2}{@{}l@{}}{[no change]} \\
\multicolumn{2}{@{}l@{}}{*Northwest Germanic*} \\
\multicolumn{2}{@{}l@{}}{[no change]} \\
\end{tabular}
\end{minipage}
&
&
\begin{minipage}[t]{\linewidth}
\centering\textbf{Old English changes}\par
\vspace{0.35em}
\raggedright
\centering\textbf{}\par
\raggedright
\vspace{0.2em}
\begin{tabular}{@{}>{\raggedright\arraybackslash}p{0.74\linewidth}@{\hspace{0.55em}}>{\raggedright\arraybackslash}p{0.16\linewidth}@{\hspace{0.25em}}}
\multicolumn{2}{@{}l@{}}{*Old English*} \\
OE Aw Long Diphthong & \emph{*skḗawô} \\
OE Sk Palatalization & \emph{*ʃḗawô} \\
OE Unstressed Long Vowel Shortening & \emph{*ʃḗawa} \\
\end{tabular}
\end{minipage}
\\
\end{tabularx}
\end{minipage}%
} \endgroup

Old English form: \emph{sċēawa}

\subsection{Reconstruction and comparative
evidence}\label{reconstruction-and-comparative-evidence-134}

Orel reconstructs the verb as \Recon{skawōjanan} `show (iptv.2sg)' and
cites Old English \emph{sceáwian} `show' beside Old Frisian
\emph{skawia} `show', Old Saxon \emph{skawōn} `show', and Old High
German \emph{scouwōn} `show' (\citeproc{ref-Orel2003}{Orel 2003, 337}).
The derivational input in this entry is not that infinitive but the
imperative singular \Recon{skáwô} `show (iptv.2sg)', a finite class-II
cell with imperative -a \textless{} \emph{*-ō}
(\citeproc{ref-RingeTaylor2014}{Ringe and Taylor 2014, 314}).

The imperative singular provides the direct comparison with the Old
English form. The lexical history still belongs to the \emph{sceáwian}
`show' verb, but the cell compared here isolates the finite \emph{-a}
outcome more clearly than the citation form does.

\subsection{Old English evidence}\label{old-english-evidence-134}

Bright lists \emph{scēawian} `show' and explicitly gives the imperative
singular \emph{scēawa} `show' under that headword
(\citeproc{ref-BrightCassidyRingler1971}{Bright 1971, 346}). The form
treated here is therefore an attested finite paradigm form, not a
reconstructed convenience form.

The spelling used in this entry is normalized \emph{sċēawa} `show',
while Bright's glossary gives source spelling \emph{scēawa} `show'. The
ordinary Old English headword remains \emph{scēawian} `show';
\emph{sċēawa} `show' is the imperative singular chosen beside it.

\subsection{Development to Old
English}\label{development-to-old-english-134}

Campbell lists \emph{scéawian} `show' under the West Germanic
\Recon{auw} `show (iptv.2sg)' developments
(\citeproc{ref-Campbell1959}{Campbell 1959, sect. 120}). Ringe and
Taylor's class-II weak imperative singular -a \textless{} \emph{*-ō}
supports the late finite ending that yields \emph{sċēawa} `show'
(\citeproc{ref-RingeTaylor2014}{Ringe and Taylor 2014, 314}). The result
is therefore the expected finite singular form of the \emph{scēawian}
`show' family rather than an analogical replacement of the headword.

\subsection{Paradigm comparison}\label{paradigm-comparison-26}

The comparison below sets the relevant forms side by side.

{\def\LTcaptype{none} % do not increment counter
\begin{longtable}[]{@{}
  >{\raggedright\arraybackslash}p{(\linewidth - 8\tabcolsep) * \real{0.2000}}
  >{\raggedright\arraybackslash}p{(\linewidth - 8\tabcolsep) * \real{0.2000}}
  >{\raggedright\arraybackslash}p{(\linewidth - 8\tabcolsep) * \real{0.2000}}
  >{\raggedright\arraybackslash}p{(\linewidth - 8\tabcolsep) * \real{0.2000}}
  >{\raggedright\arraybackslash}p{(\linewidth - 8\tabcolsep) * \real{0.2000}}@{}}
\toprule\noalign{}
\begin{minipage}[b]{\linewidth}\raggedright
PGmc cell / interpretation
\end{minipage} & \begin{minipage}[b]{\linewidth}\raggedright
Candidate input
\end{minipage} & \begin{minipage}[b]{\linewidth}\raggedright
Old English outcome or comparison
\end{minipage} & \begin{minipage}[b]{\linewidth}\raggedright
OE comparison form
\end{minipage} & \begin{minipage}[b]{\linewidth}\raggedright
Result
\end{minipage} \\
\midrule\noalign{}
\endhead
\bottomrule\noalign{}
\endlastfoot
lexeme-level infinitive & {\emph{*skáwōjaną}} & regular output
{\emph{sċēawian}} & {\emph{scēawian}} `show' & ordinary dictionary
headword of the verb, but not the finite form compared here \\
imperative singular & {\emph{*skáwô}} & regular output {\emph{sċēawa}}
`show' & {\emph{scēawa}} `show' / normalized {\emph{sċēawa}} `show' &
exact match between the derivational input and the Old English form
here \\
present third singular companion & {\emph{*skáwōθi}} & regular output
{\emph{sċēawaþ}} `show' & {\emph{sċēawaþ}} `show' & useful family
control, but not the target of this entry \\
\end{longtable}
}

\section{\texorpdfstring{show (3sg) --- OE
\emph{sċēawaþ}}{show (3sg) --- OE sċēawaþ}}\label{show-3sg-oe-sux10bux113awauxfe}

\index[iv]{02oe@\textbf{Old English}!sceawath@\emph{sċēawaþ}}
\index[iv]{03pgmc@\textbf{Proto-Germanic}!skawona@*skawōną}
\index[iv]{03pgmc@\textbf{Proto-Germanic}!skawothi@*skáwōθi}

Derivation: citation reconstruction \emph{*skawōną}; form followed here
\emph{*skáwōθi} \textgreater{} \emph{sċēawaþ} (late analogy).

\subsection{Derivation trace}\label{derivation-trace-135}

Proto input: \emph{*skáwōθi}

\begingroup
\setlength{\fboxsep}{6pt}

\noindent\fbox{%
\begin{minipage}{0.97\linewidth}
\small
\begin{tabularx}{\linewidth}{@{}>{\raggedright\arraybackslash}p{0.440\linewidth}>{\centering\arraybackslash}X>{\raggedright\arraybackslash}p{0.440\linewidth}@{}}
\begin{minipage}[t]{\linewidth}
\centering\textbf{Earlier Germanic changes}\par
\vspace{0.35em}
\raggedright
\centering\textbf{}\par
\raggedright
\vspace{0.2em}
\begin{tabular}{@{}>{\raggedright\arraybackslash}p{0.68\linewidth}@{\hspace{0.55em}}>{\raggedright\arraybackslash}p{0.16\linewidth}@{\hspace{0.25em}}}
\multicolumn{2}{@{}l@{}}{*Proto-West Germanic*} \\
PWGmc Early I Apocope & \emph{*skáwōθ} \\
\multicolumn{2}{@{}l@{}}{*Northwest Germanic*} \\
\multicolumn{2}{@{}l@{}}{[no change]} \\
\end{tabular}
\end{minipage}
&
&
\begin{minipage}[t]{\linewidth}
\centering\textbf{Old English changes}\par
\vspace{0.35em}
\raggedright
\centering\textbf{}\par
\raggedright
\vspace{0.2em}
\begin{tabular}{@{}>{\raggedright\arraybackslash}p{0.68\linewidth}@{\hspace{0.55em}}>{\raggedright\arraybackslash}p{0.20\linewidth}@{\hspace{0.25em}}}
\multicolumn{2}{@{}l@{}}{*Old English*} \\
OE Aw Long Diphthong & \emph{*skḗawōθ} \\
OE Sk Palatalization & \emph{*ʃḗawōθ} \\
OE Late O Shortening & \emph{*ʃḗawaθ} \\
\end{tabular}
\end{minipage}
\\
\end{tabularx}
\end{minipage}%
} \endgroup

Old English form: \emph{sċēawaþ}

\subsection{Reconstruction and comparative
evidence}\label{reconstruction-and-comparative-evidence-135}

Orel reconstructs the verb as {\Recon{skawōjanan} `show (3sg)'} and
cites Old English {\emph{sceáwian}} `show' beside Old Frisian
{\emph{skawia}} `show', Old Saxon {\emph{skawōn}} `show', and Old High
German {\emph{scouwōn}} `show' (\citeproc{ref-Orel2003}{Orel 2003,
337}). The derivational input in this entry is the present third
singular {\Recon{skáwōθi} `show (3sg)'}, a finite class-II cell with
stable \emph{a} in the 3sg ending (\citeproc{ref-RingeTaylor2014}{Ringe
and Taylor 2014, 80}).

Campbell states the class-II present endings as lufas, \emph{-aþ}
(\textless{} \emph{-ōsi}, \emph{-ōþi)}
(\citeproc{ref-Campbell1959}{Campbell 1959, sect. 356.4}). Ringe and
Taylor likewise note that class-II weak present 2sg. \emph{-as(t)} and
3sg. \emph{-aþ} have stable \emph{a}
(\citeproc{ref-RingeTaylor2014}{Ringe and Taylor 2014, 80}). The
relevant comparison is therefore the 3sg cell itself, not an i-umlauted
alternative.

\subsection{Old English evidence}\label{old-english-evidence-135}

Bright lists the simplex headword {\emph{scēawian}} `show' and the
imperative {\emph{scēawa}} `show', and under \emph{geond-scēawian}
`survey' also records a third singular \emph{sceawað} `show'
(\citeproc{ref-BrightCassidyRingler1971}{Bright 1971, 383}). The
evidence thus establishes the \emph{scēaw-} / \emph{sceawað} finite-cell
pattern directly.

The form written here as {\emph{sċēawaþ}} `shows' is the normalized
simplex comparison form for that weak class-II pattern. It is therefore
not a dictionary headword but a finite comparison form aligned with the
attested \emph{scēaw-} evidence and the directly cited \emph{sceawað}
`show' ending pattern.

\subsection{Development to Old
English}\label{development-to-old-english-135}

Campbell lists {\emph{scéawian}} `show' under the same West Germanic
\Recon{auw} `show (3sg)' development
(\citeproc{ref-Campbell1959}{Campbell 1959, sect. 120}).
{\Recon{skáwōθi} `show (3sg)'} therefore belongs to the \emph{scēaw-}
family before the class-II 3sg ending is applied. Campbell's chronology
and Ringe and Taylor's stable-\emph{a} discussion show that the class-II
3sg ending gives \emph{-aþ}, not \emph{-eþ}
(\citeproc{ref-Campbell1959}{Campbell 1959, sect. 356.4};
\citeproc{ref-RingeTaylor2014}{Ringe and Taylor 2014, 80}). Because the
ending never contains \emph{-j-}, no i-umlaut applies.

\subsection{Paradigm comparison}\label{paradigm-comparison-27}

The comparison below sets the relevant forms side by side.

{\def\LTcaptype{none} % do not increment counter
\begin{longtable}[]{@{}
  >{\raggedright\arraybackslash}p{(\linewidth - 8\tabcolsep) * \real{0.2000}}
  >{\raggedright\arraybackslash}p{(\linewidth - 8\tabcolsep) * \real{0.2000}}
  >{\raggedright\arraybackslash}p{(\linewidth - 8\tabcolsep) * \real{0.2000}}
  >{\raggedright\arraybackslash}p{(\linewidth - 8\tabcolsep) * \real{0.2000}}
  >{\raggedright\arraybackslash}p{(\linewidth - 8\tabcolsep) * \real{0.2000}}@{}}
\toprule\noalign{}
\begin{minipage}[b]{\linewidth}\raggedright
PGmc cell / interpretation
\end{minipage} & \begin{minipage}[b]{\linewidth}\raggedright
Candidate input
\end{minipage} & \begin{minipage}[b]{\linewidth}\raggedright
Old English outcome or comparison
\end{minipage} & \begin{minipage}[b]{\linewidth}\raggedright
OE comparison form
\end{minipage} & \begin{minipage}[b]{\linewidth}\raggedright
Result
\end{minipage} \\
\midrule\noalign{}
\endhead
\bottomrule\noalign{}
\endlastfoot
lexeme-level infinitive & {\emph{*skáwōjaną}} & regular output
{\emph{sċēawian}} & {\emph{scēawian}} `show' & ordinary dictionary
headword of the verb, but not the finite form compared here \\
imperative singular companion & {\emph{*skáwô}} & regular output
{\emph{sċēawa}} `show' & {\emph{scēawa}} `show' & useful family control,
but not the target of this entry \\
present third singular & {\emph{*skáwōθi}} & regular output
{\emph{sċēawaþ}} `shows' & normalized {\emph{sċēawaþ}} `shows';
source-side pattern \emph{sceawað} & exact match for the finite form
compared here \\
\end{longtable}
}

\clearpage

\chapter{Reconstructed Old English
comparators}\label{reconstructed-old-english-comparators}

Direct attestation does not supply the required comparator in these
entries. The target is an explicitly reconstructed Old English form and
carries the corresponding evidential burden.

\section{\texorpdfstring{knob --- OE
\emph{*cnobba}}{knob --- OE *cnobba}}\label{knob-oe-cnobba}

\index[iv]{02oe@\textbf{Old English}!cnobba@\emph{*cnobba}}
\index[iv]{03pgmc@\textbf{Proto-Germanic}!knubbo@*knúbbô}
\index[iv]{03pgmc@\textbf{Proto-Germanic}!knuppaz@*knúppaz}

Derivation: citation reconstruction \emph{*knúppaz}; form followed here
\emph{*knúbbô} \textgreater{} \emph{*cnobba} (reconstructed Old English
comparator).

\subsection{Derivation trace}\label{derivation-trace-136}

Proto input: \emph{*knúbbô}

\begingroup
\setlength{\fboxsep}{6pt}

\noindent\fbox{%
\begin{minipage}{0.97\linewidth}
\small
\begin{tabularx}{\linewidth}{@{}>{\raggedright\arraybackslash}p{0.460\linewidth}>{\centering\arraybackslash}X>{\raggedright\arraybackslash}p{0.460\linewidth}@{}}
\begin{minipage}[t]{\linewidth}
\centering\textbf{Earlier Germanic changes}\par
\vspace{0.35em}
\raggedright
\centering\textbf{}\par
\raggedright
\vspace{0.2em}
\begin{tabular}{@{}>{\raggedright\arraybackslash}p{0.74\linewidth}@{\hspace{0.55em}}>{\raggedright\arraybackslash}p{0.16\linewidth}@{\hspace{0.25em}}}
\multicolumn{2}{@{}l@{}}{*Proto-West Germanic*} \\
\multicolumn{2}{@{}l@{}}{[no change]} \\
\multicolumn{2}{@{}l@{}}{*Northwest Germanic*} \\
\mbox{NWGmc U Lowering} & \emph{*knóbbô} \\
\end{tabular}
\end{minipage}
&
&
\begin{minipage}[t]{\linewidth}
\centering\textbf{Old English changes}\par
\vspace{0.35em}
\raggedright
\centering\textbf{}\par
\raggedright
\vspace{0.2em}
\begin{tabular}{@{}>{\raggedright\arraybackslash}p{0.74\linewidth}@{\hspace{0.55em}}>{\raggedright\arraybackslash}p{0.16\linewidth}@{\hspace{0.25em}}}
\multicolumn{2}{@{}l@{}}{*Old English*} \\
OE Unstressed Long Vowel Shortening & \emph{*knóbba} \\
\end{tabular}
\end{minipage}
\\
\end{tabularx}
\end{minipage}%
} \endgroup

Old English form: \emph{*cnobba}

\subsection{Reconstruction and comparative
evidence}\label{reconstruction-and-comparative-evidence-136}

The wider knob-family is not uniform. Kroonen's discussion of the
Germanic n-stems points to related voiced and voiceless branches within
this group (\citeproc{ref-Kroonen2011}{Kroonen 2011, 297}). The citation
reconstruction \Recon{knúppaz} `knob' represents the broader cognate-set
headword, while the form followed here, \Recon{knúbbô} `knob',
represents the voiced weak-noun branch treated here.

The Old English record is uneven. The better attested OE material
belongs to the voiceless branch, but the present entry represents the
reconstructed OE form that would continue the voiced branch behind later
English knob.

\subsection{Old English evidence}\label{old-english-evidence-136}

Clark Hall preserves Old English evidence of the {\emph{cnoppa}} `knob'
type (\citeproc{ref-ClarkHall1960}{Clark Hall 1960, 79}). Those forms
are genuine Old English evidence, but they belong to the voiceless
branch of the family.

The target \textbf{\Recon{cnobba} `knob'} is a \textbf{reconstructed Old
English form}, not a directly attested one. I use it for the voiced
branch because attested {\emph{cnoppa}} `knob' continues the voiceless
branch and therefore represents a different prehistory.

\subsection{Development to Old
English}\label{development-to-old-english-136}

From the weak-noun form followed here, \Recon{knúbbô} `knob', the
regular Old English outcome is \Recon{cnobba} `knob', with
Proto-Germanic \emph{kn-} represented in Old English as \emph{cn-} and
with the expected weak-noun ending.

The entry therefore does not claim that \Recon{cnobba} `knob' is
attested. Its claim is different: if the voiced weak-noun branch is the
one to be represented, then \Recon{cnobba} is the regular Old English
form corresponding to that branch.

\subsection{Reconstruction status}\label{reconstruction-status}

{\def\LTcaptype{none} % do not increment counter
\begin{longtable}[]{@{}
  >{\raggedright\arraybackslash}p{(\linewidth - 4\tabcolsep) * \real{0.3333}}
  >{\raggedright\arraybackslash}p{(\linewidth - 4\tabcolsep) * \real{0.3333}}
  >{\raggedright\arraybackslash}p{(\linewidth - 4\tabcolsep) * \real{0.3333}}@{}}
\toprule\noalign{}
\begin{minipage}[b]{\linewidth}\raggedright
Form
\end{minipage} & \begin{minipage}[b]{\linewidth}\raggedright
Status
\end{minipage} & \begin{minipage}[b]{\linewidth}\raggedright
Relevance to this entry
\end{minipage} \\
\midrule\noalign{}
\endhead
\bottomrule\noalign{}
\endlastfoot
{\emph{\emph{*knúbbô} \textgreater{} \emph{*cnobba}}} & reconstructed OE
form; regular derivation & reconstructed Old English form compared
here \\
{\emph{cnopp}} / {\emph{cnoppa}} & attested OE branch & important
control form, but belongs to the voiceless branch \\
{\emph{cnæp}} & attested OE form from another family & not part of the
present lexeme line \\
\end{longtable}
}

This remains the most review-sensitive item here, because the choice
between reconstructed \Recon{cnobba} `knob' and attested {\emph{cnoppa}}
`knob' is still a comparator-policy question rather than a settled point
of OE attestation.

\section{\texorpdfstring{reek --- OE
\emph{*rēac}}{reek --- OE *rēac}}\label{reek-oe-rux113ac}

\index[iv]{02oe@\textbf{Old English}!reac@\emph{*rēac}}
\index[iv]{03pgmc@\textbf{Proto-Germanic}!raukaz@*ráukaz}
\index[iv]{03pgmc@\textbf{Proto-Germanic}!raukiz@*ráukiz}

Derivation: citation reconstruction \emph{*ráukiz}; form followed here
\emph{*ráukaz} \textgreater{} \emph{*rēac} (reconstructed Old English
comparator).

\subsection{Derivation trace}\label{derivation-trace-137}

Proto input: \emph{*ráukaz}

\begingroup
\setlength{\fboxsep}{6pt}

\noindent\fbox{%
\begin{minipage}{0.97\linewidth}
\small
\begin{tabularx}{\linewidth}{@{}>{\raggedright\arraybackslash}p{0.440\linewidth}>{\centering\arraybackslash}X>{\raggedright\arraybackslash}p{0.440\linewidth}@{}}
\begin{minipage}[t]{\linewidth}
\centering\textbf{Earlier Germanic changes}\par
\vspace{0.35em}
\raggedright
\centering\textbf{}\par
\raggedright
\vspace{0.2em}
\begin{tabular}{@{}>{\raggedright\arraybackslash}p{0.74\linewidth}@{\hspace{0.55em}}>{\raggedright\arraybackslash}p{0.16\linewidth}@{\hspace{0.25em}}}
\multicolumn{2}{@{}l@{}}{*Proto-West Germanic*} \\
\multicolumn{2}{@{}l@{}}{[no change]} \\
\multicolumn{2}{@{}l@{}}{*Northwest Germanic*} \\
\mbox{PGmc Final Z Deletion} & \emph{*ráuka} \\
\end{tabular}
\end{minipage}
&
&
\begin{minipage}[t]{\linewidth}
\centering\textbf{Old English changes}\par
\vspace{0.35em}
\raggedright
\centering\textbf{}\par
\raggedright
\vspace{0.2em}
\begin{tabular}{@{}>{\raggedright\arraybackslash}p{0.70\linewidth}@{\hspace{0.55em}}>{\raggedright\arraybackslash}p{0.16\linewidth}@{\hspace{0.25em}}}
\multicolumn{2}{@{}l@{}}{*Old English*} \\
OE Au Fronting & \emph{*ráeuka} \\
OE Diphthong Leveling & \emph{*rēaka} \\
PWGmc Final Bare A Loss & \emph{*rēak} \\
\end{tabular}
\end{minipage}
\\
\end{tabularx}
\end{minipage}%
} \endgroup

Old English form: \emph{*rēac}

\subsection{Reconstruction and comparative
evidence}\label{reconstruction-and-comparative-evidence-137}

The wider noun family is represented by {\Recon{ráukiz} `reek'} /
{\emph{*rauki-}}, with Old English {\emph{rēc}} `reek' as the attested
noun reflex in the comparative dictionaries
(\citeproc{ref-Kroonen2013}{Kroonen 2013, 446};
\citeproc{ref-Orel2003}{Orel 2003, 338}). The derivational input
{\Recon{ráukaz} `reek'} is therefore not the lexeme-level headword, but
the form used here for the Old English derivation.

\subsection{Old English evidence}\label{old-english-evidence-137}

The attested noun is {\emph{rēc}} `reek', not {\emph{\Recon{rēac}
`reek'}}. Clark Hall records {\emph{rēc}} `reek' as the noun and also
preserves related forms such as {\emph{rēcels}} `incense'; Kroonen
likewise gives OE {\emph{rēc}} `reek' under the noun family
(\citeproc{ref-ClarkHall1960}{Clark Hall 1960, 255};
\citeproc{ref-Kroonen2013}{Kroonen 2013, 446}). Clark Hall and Seebold
also record verbal {\emph{rēac}} `reek' as the preterite of
{\emph{rēocan}} `reek', but that verbal form is separate from the noun
treated here (\citeproc{ref-ClarkHall1960}{Clark Hall 1960, 254};
\citeproc{ref-Seebold1970}{Seebold 1970, 380}).

The Old English form here {\emph{\Recon{rēac} `reek'}} is therefore a
reconstructed West Saxon noun form, not a directly attested manuscript
headword.

\subsection{Development to Old
English}\label{development-to-old-english-137}

From {\Recon{ráukaz} `reek'}, the regular West Saxon development gives
{\emph{\Recon{rēac} `reek'}}. The attested noun {\emph{rēc}} `reek'
belongs to the same lexical family, but reflects a later smoothed
surface form rather than the regular noun target represented here.

\subsection{Form note}\label{form-note-28}

The distinction here is between an attested noun headword {\emph{rēc}}
`reek' and a reconstructed regular West Saxon target {\emph{\Recon{rēac}
`reek'}}. The latter is treated as the modelling target, while the
former remains philological background.

\section{\texorpdfstring{strew --- OE
\emph{*strīeġan}}{strew --- OE *strīeġan}}\label{strew-oe-strux12beux121an}

\index[iv]{02oe@\textbf{Old English}!striegan@\emph{*strīeġan}}
\index[iv]{03pgmc@\textbf{Proto-Germanic}!strawjana@*stráwjaną}

Derivation: \emph{*stráwjaną} \textgreater{} \emph{*strīeġan}
(reconstructed Old English comparator).

\subsection{Derivation trace}\label{derivation-trace-138}

Proto input: \emph{*stráwjaną}

\begingroup
\setlength{\fboxsep}{6pt}

\noindent\fbox{%
\begin{minipage}{0.97\linewidth}
\footnotesize
\begin{tabularx}{\linewidth}{@{}>{\raggedright\arraybackslash}p{0.320\linewidth}>{\centering\arraybackslash}X>{\raggedright\arraybackslash}p{0.520\linewidth}@{}}
\begin{minipage}[t]{\linewidth}
\centering\textbf{Earlier Germanic changes}\par
\vspace{0.35em}
\raggedright
\centering\textbf{}\par
\raggedright
\vspace{0.2em}
\begin{tabular}{@{}>{\raggedright\arraybackslash}p{0.68\linewidth}@{\hspace{0.55em}}>{\raggedright\arraybackslash}p{0.16\linewidth}@{\hspace{0.25em}}}
\multicolumn{2}{@{}l@{}}{*Proto-West Germanic*} \\
\multicolumn{2}{@{}l@{}}{[no change]} \\
\multicolumn{2}{@{}l@{}}{*Northwest Germanic*} \\
\multicolumn{2}{@{}l@{}}{[no change]} \\
\end{tabular}
\end{minipage}
&
&
\begin{minipage}[t]{\linewidth}
\centering\textbf{Old English changes}\par
\vspace{0.35em}
\raggedright
\centering\textbf{}\par
\raggedright
\vspace{0.2em}
\begin{tabular}{@{}>{\raggedright\arraybackslash}p{0.62\linewidth}@{\hspace{0.55em}}>{\raggedright\arraybackslash}p{0.28\linewidth}@{\hspace{0.25em}}}
\multicolumn{2}{@{}l@{}}{*Old English*} \\
OE Awj Glide Formation & \emph{*stráujaną} \\
OE Au Fronting & \emph{*stráeujaną} \\
OE Diphthong Leveling & \emph{*strēajaną} \\
OE Heavy Syllable Nasal Apocope & \emph{*strēajan} \\
OE Secondary Nasalization & \emph{*strēająn} \\
OE I Umlaut & \emph{*strīejąn} \\
OE Weak Tail Reduction & \emph{*strīejan} \\
OE J Strengthening After Front Diphthong & \emph{*strīeʒan} \\
\end{tabular}
\end{minipage}
\\
\end{tabularx}
\end{minipage}%
} \endgroup

Old English form: \emph{*strīeġan}

\subsection{Reconstruction and comparative
evidence}\label{reconstruction-and-comparative-evidence-138}

Kroonen cites the inherited weak verb as {\emph{*straujan-}} and gives
Old English {\emph{streowian}} `strew' as its dictionary continuation
(\citeproc{ref-Kroonen2013}{Kroonen 2013, 483}). Ringe and Taylor
separate the inherited class-I branch from the remodelled class-II
branch: Anglian {\emph{strēgan}} `strew' continues the inherited verb,
while West Saxon {\emph{streowian}} `strew' is remodelled
(\citeproc{ref-RingeTaylor2014}{Ringe and Taylor 2014, sect. 6.1}
n.~27).

Luick groups \Recon{strauwjan} `strew' with the same set as
\emph{*hauwja-} and \Recon{kauwjan} `chew', yielding Anglian
{\emph{strēzan}} `strew' beside West Saxon forms of the {\emph{hīez}}
`hew', {\emph{ciezan}} `choose' type (\citeproc{ref-Luick1914}{Luick
1914-\/-1940, sect. 98}). Fulk likewise allows an early West Saxon
\Recon{striegan} `strew' directly from Proto-Germanic \Recon{straujana}
`strew' (\citeproc{ref-Fulk2018}{Fulk 2018, sect. 4.10} n.~1).

\subsection{Old English evidence}\label{old-english-evidence-138}

The attested inherited Old English form is {\emph{strēgan}} `strew' in
Anglian. The attested West Saxon citation forms are {\emph{strewian}}
`strew', {\emph{streowian}} `strew', and {\emph{strēawian}} `strew',
which belong to the remodelled class-II branch
(\citeproc{ref-RingeTaylor2014}{Ringe and Taylor 2014, sect. 6.1} n.~27;
\citeproc{ref-Campbell1959}{Campbell 1959, sect. 753.7}).

The target {\emph{\Recon{strīeġan} `strew'}} is therefore a
reconstructed Old English form, not an attested manuscript lemma. It
represents the inferred West Saxon reflex of the inherited class-I
branch.

\subsection{Development to Old
English}\label{development-to-old-english-138}

From \Recon{stráwjaną} `strew', the inherited West Saxon line passes
through \Recon{straujaną} `strew', then a fronted \emph{*strēajan-}
stage, then \Recon{strīejan} `strew' by i-umlaut, with
retained/strengthened glide after the front diphthong to yield
reconstructed \Recon{strīeġan} `strew'.

This differs from Anglian {\emph{strēgan}} `strew', where smoothing
removes the diphthongal sequence, and from West Saxon {\emph{strewian}}
`strew' / {\emph{streowian}} `strew' / {\emph{strēawian}} `strew', where
the verb has already been remodelled into class II
(\citeproc{ref-Fulk2018}{Fulk 2018, sect. 4.10} n.~1;
\citeproc{ref-Campbell1959}{Campbell 1959, sect. 753.7}).

\subsection{Reconstruction status}\label{reconstruction-status-1}

{\def\LTcaptype{none} % do not increment counter
\begin{longtable}[]{@{}
  >{\raggedright\arraybackslash}p{(\linewidth - 4\tabcolsep) * \real{0.3333}}
  >{\raggedright\arraybackslash}p{(\linewidth - 4\tabcolsep) * \real{0.3333}}
  >{\raggedright\arraybackslash}p{(\linewidth - 4\tabcolsep) * \real{0.3333}}@{}}
\toprule\noalign{}
\begin{minipage}[b]{\linewidth}\raggedright
Form or branch
\end{minipage} & \begin{minipage}[b]{\linewidth}\raggedright
Status
\end{minipage} & \begin{minipage}[b]{\linewidth}\raggedright
Relevance to this entry
\end{minipage} \\
\midrule\noalign{}
\endhead
\bottomrule\noalign{}
\endlastfoot
{\emph{strēgan}} `strew' & attested Anglian inherited class-I form &
proves that the inherited verb survived into Old English \\
{\emph{\Recon{strīeġan} `strew'}} & reconstructed West Saxon inherited
class-I form; trace-supported & Old English form here \\
{\emph{strewian}} / {\emph{streowian}} `strew' / {\emph{strēawian}} &
attested remodelled West Saxon class-II forms & genuine OE evidence, but
not the inherited branch modeled here \\
\end{longtable}
}

\clearpage

\chapter{Known but unmodelled
developments}\label{known-but-unmodelled-developments}

The historical development is sufficiently supported by comparative and
philological evidence, but the current transducer cascade does not yet
implement the phonological or morphological process required to derive
the target. Naming the gap explains the status without pretending that
no historical account exists.

\section{\texorpdfstring{stem --- OE
\emph{stefn}}{stem --- OE stefn}}\label{stem-oe-stefn}

\index[iv]{02oe@\textbf{Old English}!stefn@\emph{stefn}}
\index[iv]{03pgmc@\textbf{Proto-Germanic}!stamnaz@*stámnaz}

Derivation: form followed here \emph{*stámnaz}; Old English form
\emph{stefn}; no regular trace was confidently matched for this entry.

\subsection{Reconstruction and comparative
evidence}\label{reconstruction-and-comparative-evidence-139}

Orel reconstructs the Proto-Germanic source as \Recon{stámnaz} `stem,
trunk' (also variant \Recon{stamniz} `stem, trunk'), citing Old Norse
{\emph{stafn}} `stem of a ship', Old Frisian {\emph{stevene}} `stem of a
ship' (fem.), and Old Saxon {\emph{stamn}} `stem' as its main
continuants (\citeproc{ref-Orel2003}{Orel 2003, 371}).

This word is etymologically unrelated to the Old English homonym
\emph{stefn} `voice, sound' / \emph{stemn} `voice, sound', which
descends from a distinct Proto-Germanic \Recon{stebnō} `voice, sound' /
\Recon{stemnō} `voice, sound' etymon. The two OE lexemes are
distinguished in Clark Hall, Bosworth-Toller, Brunner, and Luick; they
collide on the surface forms \emph{stefn} and \emph{stemn} through
historically independent developments.

Kroonen relates the `stem' family to the \emph{stam(m)} `stem, trunk'
group attested in Old High German, Dutch, and German, without
reconstructing a separate `prow' proto-form
(\citeproc{ref-Kroonen2013}{Kroonen 2013, 479--80}).

\subsection{Old English evidence}\label{old-english-evidence-139}

Clark Hall records {\emph{stefn}} `stem, trunk, root, prow, foundation',
the weak n-stem {\emph{stefna}} `prow or stern of a ship', and the
related {\emph{stofn}} `trunk, stem, branch, shoot'
(\citeproc{ref-ClarkHall1960}{Clark Hall 1960, 276, 341}).

Luick §211 explicitly distinguishes OE \emph{stefn} `stem, trunk' /
\emph{stemn} `stem, trunk' --- noting cognates with Old Saxon
{\emph{stamn}} `stem' and Middle English {\emph{stam}} `stem' --- from
the separate OE voice/sound word (\emph{stefn} `voice, sound' /
\emph{stemn} `voice, sound') treated in §75
(\citeproc{ref-Luick1914}{Luick 1914-\/-1940, sect. 211}). Brunner §205
lists \emph{stefn, stemn Stamm} among words showing the
\emph{fn}/\emph{mn} alternation
(\citeproc{ref-SieversBrunner1965}{Brunner 1965, sect. 205}).

The primary comparison form used here is {\emph{stefn}} `stem, trunk,
root, prow', the form attested most directly in the semantic range
relevant to English \emph{stem}. The n-stem {\emph{stefna}} `prow/stern'
is a closely related but more narrowly attested nautical specialization.

\subsection{Development to Old
English}\label{development-to-old-english-139}

The derivation of OE {\emph{stefn}} `stem, trunk' from \Recon{stámnaz}
`stem, trunk' involves a cluster change \emph{mn → fn} that is attested
comparatively: Old Norse {\emph{stafn}} `stem of a ship' and Old Saxon
{\emph{stamn}} `stem' preserve the \emph{fn}/\emph{mn} variants expected
from this family. The precise phonological dating and domain of this
change --- whether it belongs to a North-West Germanic stage, an early
West Germanic stage, or is reconstructed separately in each branch ---
is not definitively resolved by the cited scholarship. The handbooks
treat it as a pre-Old-English development (Luick §211; Brunner §205),
but the term ``North-West Germanic coda dissimilation'' should be taken
as descriptive rather than technically established. This coda
environment (\emph{{[}C{]}mn → {[}C{]}fn}) is distinct from the
cross-syllable m → \emph{β} change (as in \Recon{xémonų} `heaven' →
{\emph{heofon}} `heaven') handled by the \emph{NWGmcMnDissimilation}
rule; it is currently not modeled in the FST.

The correct OE-facing transponent has not yet been established for the
FST pipeline. The classification \emph{known\_unmodelled} reflects that
the derivation involves a historically real development not yet
represented in the FST cascade. The derivation is attested
comparatively; the FST work remains to be done.

\subsection{Homonym note}\label{homonym-note}

Old English \emph{stefn}/\emph{stemn} is a surface homonym for two
etymologically unrelated words:

\begin{enumerate}
\def\labelenumi{\arabic{enumi}.}
\tightlist
\item
  \emph{stefn}/\emph{stemn} `voice, sound': from PGmc \emph{*stebnō}
  (Ringe-Taylor, Orel); via \emph{bn → fn} (b-allophony), giving
  \emph{stefn} `voice, sound'. This is a \textbf{different row} from
  2216.
\item
  \emph{stefn}/\emph{stefna}/\emph{stofn}/\emph{stemn} `stem, trunk,
  prow': from PGmc \Recon{stámnaz} `stem/trunk' (Orel); via the pre-OE
  cluster change \emph{mn → fn}. \textbf{This is row 2216.}
\end{enumerate}

The form followed here, \emph{*stébnō}, used in earlier versions of this
entry was the wrong homonym's transponent and must not be used here.

\subsection{Source comparison}\label{source-comparison}

{\def\LTcaptype{none} % do not increment counter
\begin{longtable}[]{@{}
  >{\raggedright\arraybackslash}p{(\linewidth - 6\tabcolsep) * \real{0.2500}}
  >{\raggedright\arraybackslash}p{(\linewidth - 6\tabcolsep) * \real{0.2500}}
  >{\raggedright\arraybackslash}p{(\linewidth - 6\tabcolsep) * \real{0.2500}}
  >{\raggedright\arraybackslash}p{(\linewidth - 6\tabcolsep) * \real{0.2500}}@{}}
\toprule\noalign{}
\begin{minipage}[b]{\linewidth}\raggedright
Form or label
\end{minipage} & \begin{minipage}[b]{\linewidth}\raggedright
Status
\end{minipage} & \begin{minipage}[b]{\linewidth}\raggedright
OE relation
\end{minipage} & \begin{minipage}[b]{\linewidth}\raggedright
Result
\end{minipage} \\
\midrule\noalign{}
\endhead
\bottomrule\noalign{}
\endlastfoot
\Recon{stámnaz} `stem, trunk' & Orel's PGmc citation for stem/trunk/prow
family & controls this derivation & selected comparative citation \\
{\emph{stefn}} III `stem, trunk, root, prow' & primary OE target (strong
masc.) & stem/trunk sense per Clark Hall & target form \\
{\emph{stefna}} `prow/stern' & OE n-stem specialization & nautical sense
per Clark Hall & comparison form \\
{\emph{stofn}} `trunk' & OE \emph{o}-grade variant or earlier stage &
trunk/stem sense & comparison form \\
{\emph{stemn}} `trunk' & late West Saxon \emph{fn → mn} doublet &
secondary form & comparison form \\
\emph{*stébnō} (voice word) & \textbf{wrong homonym} --- belongs to
voice/sound dossier & no relation to stem/trunk sense & \textbf{must not
be used as row 2216's derivational input} \\
\end{longtable}
}

\section{\texorpdfstring{fire --- OE
\emph{fȳre}}{fire --- OE fȳre}}\label{fire-oe-fux233re}

\index[iv]{02oe@\textbf{Old English}!fyre@\emph{fȳre}}
\index[iv]{03pgmc@\textbf{Proto-Germanic}!furi@*fūri}

Derivation: \emph{*fūri} yields regular \emph{fȳr}; the Old English form
here is \emph{fȳre} (known but unmodelled development).

\subsection{Derivation trace}\label{derivation-trace-139}

Proto input: \emph{*fūri}

\begingroup
\setlength{\fboxsep}{6pt}

\noindent\fbox{%
\begin{minipage}{0.97\linewidth}
\small
\begin{tabularx}{\linewidth}{@{}>{\raggedright\arraybackslash}p{0.440\linewidth}>{\centering\arraybackslash}X>{\raggedright\arraybackslash}p{0.440\linewidth}@{}}
\begin{minipage}[t]{\linewidth}
\centering\textbf{Earlier Germanic changes}\par
\vspace{0.35em}
\raggedright
\centering\textbf{}\par
\raggedright
\vspace{0.2em}
\begin{tabular}{@{}>{\raggedright\arraybackslash}p{0.68\linewidth}@{\hspace{0.55em}}>{\raggedright\arraybackslash}p{0.16\linewidth}@{\hspace{0.25em}}}
\multicolumn{2}{@{}l@{}}{*Proto-West Germanic*} \\
\multicolumn{2}{@{}l@{}}{[no change]} \\
\multicolumn{2}{@{}l@{}}{*Northwest Germanic*} \\
\multicolumn{2}{@{}l@{}}{[no change]} \\
\end{tabular}
\end{minipage}
&
&
\begin{minipage}[t]{\linewidth}
\centering\textbf{Old English changes}\par
\vspace{0.35em}
\raggedright
\centering\textbf{}\par
\raggedright
\vspace{0.2em}
\begin{tabular}{@{}>{\raggedright\arraybackslash}p{0.74\linewidth}@{\hspace{0.55em}}>{\raggedright\arraybackslash}p{0.16\linewidth}@{\hspace{0.25em}}}
\multicolumn{2}{@{}l@{}}{*Old English*} \\
OE I Umlaut & \emph{*fȳri} \\
\mbox{OE High Vowel Apocope} & \emph{*fȳr} \\
\end{tabular}
\end{minipage}
\\
\end{tabularx}
\end{minipage}%
} \endgroup

Regular outcome: \emph{fȳr}

Old English form: \emph{fȳre}

\subsection{Reconstruction and comparative
evidence}\label{reconstruction-and-comparative-evidence-140}

Kroonen places the lexeme in a heteroclitic family \Recon{fōr} `fire'
\textasciitilde{} \emph{*fun-} and explains the front-mutated West
Germanic forms from an oblique form of the \emph{*fu(w)eri} type
(\citeproc{ref-Kroonen2013}{Kroonen 2013, 151}). The derivational input
{\Recon{fūri} `fire'} therefore does not function as an arbitrary
substitute for the headword: it represents the specific inherited cell
that supplies the \emph{i} needed for i-umlaut.

The Old English target combines a regular inherited form {\emph{fȳr}}
`fire' with an attested analogical surface form {\emph{fȳre}} `fire'.

\subsection{Old English evidence}\label{old-english-evidence-140}

Bosworth-Toller records {\emph{fyr}} `fire' as the noun `fire' and also
preserves oblique {\emph{fyre}} `fire' in the Old English record
(\citeproc{ref-BosworthToller1898}{Bosworth and Toller 1898, 288}). The
first is the regular inherited outcome of the phonological development
from the selected input; the second shows the later restoration of a
final \emph{-e} within the paradigm.

The entry therefore concerns the relation between a regular inherited
oblique input and an attested Old English surface form that has
undergone later morphological remodeling.

\subsection{Development to Old
English}\label{development-to-old-english-140}

From {\Recon{fūri} `fire'}, i-umlaut changes \emph{ū} to \emph{ȳ}
(\citeproc{ref-Hogg1992}{Hogg 1992, sect. 3.3.3.1}). Subsequent loss of
the final high vowel after a heavy syllable yields {\emph{fȳr}} `fire'
(\citeproc{ref-Campbell1959}{Campbell 1959, sect. 345}). The inherited
phonology is complete at that point.

{\emph{fȳre}} `fire' is later than that inherited output. Its final
\emph{-e} belongs to analogical restoration in the Old English paradigm
rather than to the original Proto-Germanic ending. The form therefore
remains a known but unmodelled remodelling: the deterministic phonology
is regular, but the attested surface form includes later morphological
rebuilding.

\subsection{Paradigm comparison}\label{paradigm-comparison-28}

The comparison below sets the relevant forms side by side. It
distinguishes the lexeme-level etymological background from the
inherited cell that actually produces the front-mutated form and from
the later analogical surface result.

{\def\LTcaptype{none} % do not increment counter
\begin{longtable}[]{@{}
  >{\raggedright\arraybackslash}p{(\linewidth - 8\tabcolsep) * \real{0.2000}}
  >{\raggedright\arraybackslash}p{(\linewidth - 8\tabcolsep) * \real{0.2000}}
  >{\raggedright\arraybackslash}p{(\linewidth - 8\tabcolsep) * \real{0.2000}}
  >{\raggedright\arraybackslash}p{(\linewidth - 8\tabcolsep) * \real{0.2000}}
  >{\raggedright\arraybackslash}p{(\linewidth - 8\tabcolsep) * \real{0.2000}}@{}}
\toprule\noalign{}
\begin{minipage}[b]{\linewidth}\raggedright
PGmc cell / interpretation
\end{minipage} & \begin{minipage}[b]{\linewidth}\raggedright
Candidate input
\end{minipage} & \begin{minipage}[b]{\linewidth}\raggedright
OE output or comparison
\end{minipage} & \begin{minipage}[b]{\linewidth}\raggedright
OE comparison form
\end{minipage} & \begin{minipage}[b]{\linewidth}\raggedright
Result
\end{minipage} \\
\midrule\noalign{}
\endhead
\bottomrule\noalign{}
\endlastfoot
lexeme-level heteroclitic headword & \Recon{fōr} `fire'
\textasciitilde{} \emph{*fun-} & comparative background only & fire
family & explains the wider lexeme, but not the selected oblique
input \\
inherited oblique cell & {\emph{*fūri}} & regular output: {\emph{fȳr}}
`fire' & {\emph{fȳr}} `fire' & regular inherited output from the
derivational input \\
later analogical surface form & --- & attested {\emph{fȳre}} `fire' with
restored \emph{-e} & {\emph{fȳre}} `fire' & genuine OE target, but not
the direct phonological output \\
\end{longtable}
}

\section{\texorpdfstring{tap --- OE
\emph{tæppa}}{tap --- OE tæppa}}\label{tap-oe-tuxe6ppa}

\index[iv]{02oe@\textbf{Old English}!taeppa@\emph{tæppa}}
\index[iv]{03pgmc@\textbf{Proto-Germanic}!tappo@*táppô}

Derivation: \emph{*táppô} yields regular \emph{tappa}; the Old English
form here is \emph{tæppa} (known but unmodelled development).

\subsection{Derivation trace}\label{derivation-trace-140}

Proto input: \emph{*táppô}

\begingroup
\setlength{\fboxsep}{6pt}

\noindent\fbox{%
\begin{minipage}{0.97\linewidth}
\small
\begin{tabularx}{\linewidth}{@{}>{\raggedright\arraybackslash}p{0.320\linewidth}>{\centering\arraybackslash}X>{\raggedright\arraybackslash}p{0.520\linewidth}@{}}
\begin{minipage}[t]{\linewidth}
\centering\textbf{Earlier Germanic changes}\par
\vspace{0.35em}
\raggedright
\centering\textbf{}\par
\raggedright
\vspace{0.2em}
\begin{tabular}{@{}>{\raggedright\arraybackslash}p{0.68\linewidth}@{\hspace{0.55em}}>{\raggedright\arraybackslash}p{0.16\linewidth}@{\hspace{0.25em}}}
\multicolumn{2}{@{}l@{}}{*Proto-West Germanic*} \\
\multicolumn{2}{@{}l@{}}{[no change]} \\
\multicolumn{2}{@{}l@{}}{*Northwest Germanic*} \\
\multicolumn{2}{@{}l@{}}{[no change]} \\
\end{tabular}
\end{minipage}
&
&
\begin{minipage}[t]{\linewidth}
\centering\textbf{Old English changes}\par
\vspace{0.35em}
\raggedright
\centering\textbf{}\par
\raggedright
\vspace{0.2em}
\begin{tabular}{@{}>{\raggedright\arraybackslash}p{0.74\linewidth}@{\hspace{0.55em}}>{\raggedright\arraybackslash}p{0.16\linewidth}@{\hspace{0.25em}}}
\multicolumn{2}{@{}l@{}}{*Old English*} \\
Anglo Frisian Brightening & \emph{*tæppô} \\
OE A Restoration & \emph{*tappô} \\
OE Unstressed Long Vowel Shortening & \emph{*tappa} \\
\end{tabular}
\end{minipage}
\\
\end{tabularx}
\end{minipage}%
} \endgroup

Regular outcome: \emph{tappa}

Old English form: \emph{tæppa}

\subsection{Reconstruction and comparative
evidence}\label{reconstruction-and-comparative-evidence-141}

Orel gives the noun under \Recon{tappòn} `tap' and already connects it
with Old English {\emph{tæppa}} `tap' (\citeproc{ref-Orel2003}{Orel
2003, 402}). The derivational input {\Recon{táppô} `tap'} is therefore
the inherited noun itself; the entry does not depend on a different
lexeme-level proto or a different inherited noun cell.

\subsection{Old English evidence}\label{old-english-evidence-141}

The Old English noun family is well attested. Orel gives {\emph{tæppa}}
`tap', and Clark Hall records {\emph{tæppa}} `tap' together with
derivatives \emph{tæppere} `tapster' and \emph{tæppestre} `tapster'
(\citeproc{ref-Orel2003}{Orel 2003, 402};
\citeproc{ref-ClarkHall1960}{Clark Hall 1960, 305}). The target is
therefore a real Old English noun form, not a reconstructed convenience
spelling.

\subsection{Development to Old
English}\label{development-to-old-english-141}

From {\Recon{táppô} `tap'}, the regular inherited noun path gives
{\emph{tappa}} `tap'. The attested target {\emph{tæppa}} `tap' therefore
stands outside that regular phonological development.

The mismatch is historically intelligible, but it is not solved here by
a new inherited input. A related j-verb pathway would give
{\emph{teppan}} `tap', not the noun target {\emph{tæppa}} `tap'. The
entry accordingly remains a known but unmodelled case.

\subsection{Form comparison}\label{form-comparison-6}

{\def\LTcaptype{none} % do not increment counter
\begin{longtable}[]{@{}
  >{\raggedright\arraybackslash}p{(\linewidth - 6\tabcolsep) * \real{0.2500}}
  >{\raggedright\arraybackslash}p{(\linewidth - 6\tabcolsep) * \real{0.2500}}
  >{\raggedright\arraybackslash}p{(\linewidth - 6\tabcolsep) * \real{0.2500}}
  >{\raggedright\arraybackslash}p{(\linewidth - 6\tabcolsep) * \real{0.2500}}@{}}
\toprule\noalign{}
\begin{minipage}[b]{\linewidth}\raggedright
Form type
\end{minipage} & \begin{minipage}[b]{\linewidth}\raggedright
Input or form
\end{minipage} & \begin{minipage}[b]{\linewidth}\raggedright
OE output or comparison
\end{minipage} & \begin{minipage}[b]{\linewidth}\raggedright
Result
\end{minipage} \\
\midrule\noalign{}
\endhead
\bottomrule\noalign{}
\endlastfoot
regular inherited noun path & {\emph{*táppô}} & regular output:
{\emph{tappa}} `tap' & regular output, but not the target \\
attested OE target & --- & {\emph{tæppa}} `tap' & genuine target form,
but analogically remodelled in the present classification \\
related j-verb background & {\emph{*táppjaną}} & {\emph{teppan}} `tap' &
related formation, but not the noun target \\
\end{longtable}
}

\clearpage

\chapter{Unexplained or deliberately unmodelled
exceptions}\label{unexplained-or-deliberately-unmodelled-exceptions}

No sufficiently supported account yet reconciles the regular output with
the Old English form. An ad hoc sound law would conceal rather than
explain the mismatch.

\section{\texorpdfstring{buck --- OE
\emph{bucc}}{buck --- OE bucc}}\label{buck-oe-bucc}

\index[iv]{02oe@\textbf{Old English}!bucc@\emph{bucc}}
\index[iv]{03pgmc@\textbf{Proto-Germanic}!bukkaz@*búkkaz}

Derivation: \emph{*búkkaz} yields regular \emph{bocc}; the Old English
form here is \emph{bucc} (unexplained exception).

\subsection{Derivation trace}\label{derivation-trace-141}

Proto input: \emph{*búkkaz}

\begingroup
\setlength{\fboxsep}{6pt}

\noindent\fbox{%
\begin{minipage}{0.97\linewidth}
\small
\begin{tabularx}{\linewidth}{@{}>{\raggedright\arraybackslash}p{0.520\linewidth}>{\centering\arraybackslash}X>{\raggedright\arraybackslash}p{0.320\linewidth}@{}}
\begin{minipage}[t]{\linewidth}
\centering\textbf{Earlier Germanic changes}\par
\vspace{0.35em}
\raggedright
\centering\textbf{}\par
\raggedright
\vspace{0.2em}
\begin{tabular}{@{}>{\raggedright\arraybackslash}p{0.74\linewidth}@{\hspace{0.55em}}>{\raggedright\arraybackslash}p{0.16\linewidth}@{\hspace{0.25em}}}
\multicolumn{2}{@{}l@{}}{*Proto-West Germanic*} \\
\multicolumn{2}{@{}l@{}}{[no change]} \\
\multicolumn{2}{@{}l@{}}{*Northwest Germanic*} \\
\mbox{NWGmc U Lowering} & \emph{*bókkaz} \\
\mbox{PGmc Final Z Deletion} & \emph{*bókka} \\
\end{tabular}
\end{minipage}
&
&
\begin{minipage}[t]{\linewidth}
\centering\textbf{Old English changes}\par
\vspace{0.35em}
\raggedright
\centering\textbf{}\par
\raggedright
\vspace{0.2em}
\begin{tabular}{@{}>{\raggedright\arraybackslash}p{0.70\linewidth}@{\hspace{0.55em}}>{\raggedright\arraybackslash}p{0.16\linewidth}@{\hspace{0.25em}}}
\multicolumn{2}{@{}l@{}}{*Old English*} \\
PWGmc Final Bare A Loss & \emph{*bókk} \\
\end{tabular}
\end{minipage}
\\
\end{tabularx}
\end{minipage}%
} \endgroup

Regular outcome: \emph{bocc}

Old English form: \emph{bucc}

\subsection{Reconstruction and comparative
evidence}\label{reconstruction-and-comparative-evidence-142}

Kroonen and Orel both reconstruct the word with a geminate stop,
{\Recon{bukkaz} `buck'} (\citeproc{ref-Kroonen2013}{Kroonen 2013, 121};
\citeproc{ref-Orel2003}{Orel 2003, 61}). Orel also preserves parallel
n-stem material behind Old English {\emph{bucca}} `buck'
(\citeproc{ref-Orel2003}{Orel 2003, 62}). The derivational input
therefore remains identical with the lexeme label: no alternative
inherited cell accounts for the form.

\subsection{Old English evidence}\label{old-english-evidence-142}

Old English preserves a mixed lexical picture. Campbell cites
{\emph{bucca}} `buck' in the exception set for this phonological
environment (\citeproc{ref-Campbell1959}{Campbell 1959, sect. 115}).
Clark Hall and Bosworth-Toller show that Old English has both
{\emph{bucca}} `buck' and {\emph{bucc}} `buck'
(\citeproc{ref-ClarkHall1960}{Clark Hall 1960, 53};
\citeproc{ref-BosworthToller1898}{Bosworth and Toller 1898, 122}). The
a-stem citation form {\emph{bucc}} `buck' is the target treated here;
{\emph{bucca}} `buck' supplies genuine philological background from the
same lexical family.

\subsection{Development to Old
English}\label{development-to-old-english-142}

From {\Recon{búkkaz} `buck'}, the regular inherited path gives
{\emph{bocc}} `buck'. That is the form expected under the ordinary
lowering pattern in this environment. {\emph{bucc}} `buck' therefore
remains outside the deterministic phonology.

No accepted inherited cell repairs the mismatch. A high-vowel
alternative would introduce i-umlaut and produce a {\emph{byċċ}}
`buck'-type form rather than the target. {\emph{bucc}} `buck' is
therefore best treated as a documented exception, not as a regular
paradigm-cell survival.

\subsection{Form comparison}\label{form-comparison-7}

{\def\LTcaptype{none} % do not increment counter
\begin{longtable}[]{@{}
  >{\raggedright\arraybackslash}p{(\linewidth - 6\tabcolsep) * \real{0.2500}}
  >{\raggedright\arraybackslash}p{(\linewidth - 6\tabcolsep) * \real{0.2500}}
  >{\raggedright\arraybackslash}p{(\linewidth - 6\tabcolsep) * \real{0.2500}}
  >{\raggedright\arraybackslash}p{(\linewidth - 6\tabcolsep) * \real{0.2500}}@{}}
\toprule\noalign{}
\begin{minipage}[b]{\linewidth}\raggedright
Form type
\end{minipage} & \begin{minipage}[b]{\linewidth}\raggedright
Input or form
\end{minipage} & \begin{minipage}[b]{\linewidth}\raggedright
OE output or comparison
\end{minipage} & \begin{minipage}[b]{\linewidth}\raggedright
Result
\end{minipage} \\
\midrule\noalign{}
\endhead
\bottomrule\noalign{}
\endlastfoot
regular inherited noun path & {\emph{*búkkaz}} & regular output:
{\emph{bocc}} `buck' & regular output, but not the target \\
attested OE target & --- & {\emph{bucc}} `buck' & genuine target form,
but unexplained in the present classification \\
parallel OE lexical background & --- & {\emph{bucca}} `buck' & related
n-stem form, not the present target \\
\end{longtable}
}

\section{\texorpdfstring{fowl --- OE
\emph{fugol}}{fowl --- OE fugol}}\label{fowl-oe-fugol}

\index[iv]{02oe@\textbf{Old English}!fugol@\emph{fugol}}
\index[iv]{03pgmc@\textbf{Proto-Germanic}!fuglaz@*fúglaz}
\index[iv]{08on@\textbf{Old Norse}!fugl@fugl}
\index[iv]{09ohg@\textbf{Old High German}!fogal@fogal}

Derivation: \emph{*fúglaz} yields regular \emph{fogol}; the Old English
form here is \emph{fugol} (unexplained exception).

\subsection{Derivation trace}\label{derivation-trace-142}

Proto input: \emph{*fúglaz}

\begingroup
\setlength{\fboxsep}{6pt}

\noindent\fbox{%
\begin{minipage}{0.97\linewidth}
\small
\begin{tabularx}{\linewidth}{@{}>{\raggedright\arraybackslash}p{0.440\linewidth}>{\centering\arraybackslash}X>{\raggedright\arraybackslash}p{0.440\linewidth}@{}}
\begin{minipage}[t]{\linewidth}
\centering\textbf{Earlier Germanic changes}\par
\vspace{0.35em}
\raggedright
\centering\textbf{}\par
\raggedright
\vspace{0.2em}
\begin{tabular}{@{}>{\raggedright\arraybackslash}p{0.74\linewidth}@{\hspace{0.55em}}>{\raggedright\arraybackslash}p{0.16\linewidth}@{\hspace{0.25em}}}
\multicolumn{2}{@{}l@{}}{*Proto-West Germanic*} \\
\multicolumn{2}{@{}l@{}}{[no change]} \\
\multicolumn{2}{@{}l@{}}{*Northwest Germanic*} \\
\mbox{NWGmc U Lowering} & \emph{*fóglaz} \\
\mbox{PGmc Final Z Deletion} & \emph{*fógla} \\
\end{tabular}
\end{minipage}
&
&
\begin{minipage}[t]{\linewidth}
\centering\textbf{Old English changes}\par
\vspace{0.35em}
\raggedright
\centering\textbf{}\par
\raggedright
\vspace{0.2em}
\begin{tabular}{@{}>{\raggedright\arraybackslash}p{0.70\linewidth}@{\hspace{0.55em}}>{\raggedright\arraybackslash}p{0.16\linewidth}@{\hspace{0.25em}}}
\multicolumn{2}{@{}l@{}}{*Old English*} \\
PWGmc Final Bare A Loss & \emph{*fógl} \\
OE Epenthetic Vowel & \emph{*fógol} \\
\end{tabular}
\end{minipage}
\\
\end{tabularx}
\end{minipage}%
} \endgroup

Regular outcome: \emph{fogol}

Old English form: \emph{fugol}

\subsection{Reconstruction and comparative
evidence}\label{reconstruction-and-comparative-evidence-143}

The noun is the ordinary Germanic a-stem \Recon{fúglaz} `fowl',
continued by forms such as Old Norse \emph{fugl} and Old High German
{\emph{fogal}} `bird, fowl' (\citeproc{ref-Kroonen2013}{Kroonen 2013,
197}; \citeproc{ref-Orel2003}{Orel 2003, 155}). There is no stem-class
or paradigm-cell dispute behind this entry. The comparative headword and
the derivational input are the same.

The difficulty lies instead in the stressed root vowel. Under the
regular West Germanic and Old English development, that \emph{u} should
lower before the following non-high vowel, yielding an \emph{o}-vocalism
(\citeproc{ref-RingeTaylor2014}{Ringe and Taylor 2014, 42--43};
\citeproc{ref-Campbell1959}{Campbell 1959, 43}).

\subsection{Old English evidence}\label{old-english-evidence-143}

Old English dictionaries record the noun as \emph{fugol} `fowl', with
variant spelling \emph{fugel} `fowl'
(\citeproc{ref-BosworthToller1898}{Bosworth and Toller 1898, 282};
\citeproc{ref-ClarkHall1960}{Clark Hall 1960, 138}). The target is
therefore an attested ordinary Old English noun, not a reconstructed or
selectively chosen paradigm form.

The attested word already contains the crucial problem. Its medial
\emph{-o-} is the ordinary parasite vowel of Old English cluster
phonology, but the root \emph{fu-} retains \emph{u} where the regular
history predicts \emph{fo-} (\citeproc{ref-Campbell1959}{Campbell 1959,
150}).

\subsection{Development to Old
English}\label{development-to-old-english-143}

From \Recon{fúglaz} `fowl', the regular cascade yields \emph{fogol}
`fowl': the root vowel lowers before the following non-high vowel
(\citeproc{ref-RingeTaylor2014}{Ringe and Taylor 2014, 42--43};
\citeproc{ref-Campbell1959}{Campbell 1959, 43}), final \emph{z} is lost,
and the cluster is resolved by the usual medial vowel
(\citeproc{ref-RingeTaylor2014}{Ringe and Taylor 2014, 345};
\citeproc{ref-Campbell1959}{Campbell 1959, 150}). That is the expected
inherited outcome.

The attested Old English noun is \emph{fugol} `fowl', not \emph{fogol}
`fowl'. Luick and later handbooks treat this preservation of \emph{u} as
a small inherited residue, not as a categorical sound law
(\citeproc{ref-Luick1914}{Luick 1914-\/-1940, 148};
\citeproc{ref-RingeTaylor2014}{Ringe and Taylor 2014, 47}). The item
therefore remains a genuine lexical exception rather than the output of
a recoverable regular mechanism.

\subsection{Expected and attested
forms}\label{expected-and-attested-forms}

The comparison below sets the relevant forms side by side. It
distinguishes the regular prediction from the attested Old English noun.

{\def\LTcaptype{none} % do not increment counter
\begin{longtable}[]{@{}
  >{\raggedright\arraybackslash}p{(\linewidth - 4\tabcolsep) * \real{0.3333}}
  >{\raggedright\arraybackslash}p{(\linewidth - 4\tabcolsep) * \real{0.3333}}
  >{\raggedright\arraybackslash}p{(\linewidth - 4\tabcolsep) * \real{0.3333}}@{}}
\toprule\noalign{}
\begin{minipage}[b]{\linewidth}\raggedright
Form
\end{minipage} & \begin{minipage}[b]{\linewidth}\raggedright
Status
\end{minipage} & \begin{minipage}[b]{\linewidth}\raggedright
Relevance to this entry
\end{minipage} \\
\midrule\noalign{}
\endhead
\bottomrule\noalign{}
\endlastfoot
\emph{fogol} & computed regular output from \emph{*fúglaz} & establishes
the expected inherited outcome \\
\emph{fugol} & attested Old English form & Old English form here;
preserves unexplained root \emph{u} \\
\emph{fugel} & attested variant spelling & secondary spelling variant of
the attested noun \\
\end{longtable}
}

The unresolved point lies only in the root vowel. The medial \emph{-o-}
is regular, but no accepted lautgesetzlich pathway has been found from
\Recon{fúglaz} `fowl' to attested \emph{fugol} `fowl'.

\section{\texorpdfstring{rust --- OE
\emph{rust}}{rust --- OE rust}}\label{rust-oe-rust}

\index[iv]{02oe@\textbf{Old English}!rust@\emph{rust}}
\index[iv]{03pgmc@\textbf{Proto-Germanic}!rusto@*rústō}

Derivation: \emph{*rústō} yields regular \emph{rost}; the Old English
form here is \emph{rust} (unexplained exception).

\subsection{Derivation trace}\label{derivation-trace-143}

Proto input: \emph{*rústō}

\begingroup
\setlength{\fboxsep}{6pt}

\noindent\fbox{%
\begin{minipage}{0.97\linewidth}
\small
\begin{tabularx}{\linewidth}{@{}>{\raggedright\arraybackslash}p{0.460\linewidth}>{\centering\arraybackslash}X>{\raggedright\arraybackslash}p{0.460\linewidth}@{}}
\begin{minipage}[t]{\linewidth}
\centering\textbf{Earlier Germanic changes}\par
\vspace{0.35em}
\raggedright
\centering\textbf{}\par
\raggedright
\vspace{0.2em}
\begin{tabular}{@{}>{\raggedright\arraybackslash}p{0.74\linewidth}@{\hspace{0.55em}}>{\raggedright\arraybackslash}p{0.16\linewidth}@{\hspace{0.25em}}}
\multicolumn{2}{@{}l@{}}{*Proto-West Germanic*} \\
\multicolumn{2}{@{}l@{}}{[no change]} \\
\multicolumn{2}{@{}l@{}}{*Northwest Germanic*} \\
\mbox{NWGmc U Lowering} & \emph{*róstō} \\
\mbox{NWGmc Final Long O Raising} & \emph{*róstu} \\
\end{tabular}
\end{minipage}
&
&
\begin{minipage}[t]{\linewidth}
\centering\textbf{Old English changes}\par
\vspace{0.35em}
\raggedright
\centering\textbf{}\par
\raggedright
\vspace{0.2em}
\begin{tabular}{@{}>{\raggedright\arraybackslash}p{0.74\linewidth}@{\hspace{0.55em}}>{\raggedright\arraybackslash}p{0.16\linewidth}@{\hspace{0.25em}}}
\multicolumn{2}{@{}l@{}}{*Old English*} \\
\mbox{OE High Vowel Apocope} & \emph{*róst} \\
\end{tabular}
\end{minipage}
\\
\end{tabularx}
\end{minipage}%
} \endgroup

Regular outcome: \emph{rost}

Old English form: \emph{rust}

\subsection{Reconstruction and comparative
evidence}\label{reconstruction-and-comparative-evidence-144}

The comparative dictionaries do not support a single citation
reconstruction uniformly. Orel cites \Recon{rustaz} `rust' (sb.m./f.)
with Old English {\emph{rust}} `rust' and Old Saxon and Old High German
\emph{rost} `rust' (\citeproc{ref-Orel2003}{Orel 2003, 308}). The form
\Recon{rústō} `rust' therefore stands here as a competing citation
reconstruction rather than as the best-supported inherited headword.

That disagreement does not remove the central problem. Whether one
starts from \Recon{rústō} `rust' or from source-supported \Recon{rustaz}
`rust', the regular citation-form history points toward \emph{rost}
`rust', not toward the attested Old English noun.

\subsection{Old English evidence}\label{old-english-evidence-144}

The Old English noun is attested, not reconstructed. Clark Hall gives
{\emph{rūst}} `rust' (m.) (\citeproc{ref-ClarkHall1960}{Clark Hall 1960,
245}), and Bosworth-Toller records {\emph{rúst}} `rust' (also
\emph{rust}) (\citeproc{ref-BosworthToller1898}{Bosworth and Toller
1898, 677}). The form is normalized here as {\emph{rust}} `rust' from
that attested record.

Those dictionary entries identify a masculine noun, which aligns better
with Orel's \Recon{rustaz} `rust' than with the competing \Recon{rústō}
`rust' preserved in the header.

\subsection{Development to Old
English}\label{development-to-old-english-144}

Under Campbell's regular lowering of stressed \emph{u} before a
following mid or low vowel, the citation-form input gives \emph{rost}
`rust', not {\emph{rust}} `rust' (\citeproc{ref-Campbell1959}{Campbell
1959, sect. 115}). The same lowering would also affect comparative
citation-form reconstructions such as \Recon{rustaz} `rust'.

A high-vowel comparator such as instrumental-type \Recon{rústu} `rust'
would yield {\emph{rust}} `rust' regularly, but that does not explain
the attested citation form of the noun. No accepted regular pathway from
the citation form to attested {\emph{rust}} `rust' has been established.

\subsection{Expected and attested
forms}\label{expected-and-attested-forms-1}

The comparison below sets the regular inherited outcomes beside the
attested Old English noun.

{\def\LTcaptype{none} % do not increment counter
\begin{longtable}[]{@{}
  >{\raggedright\arraybackslash}p{(\linewidth - 4\tabcolsep) * \real{0.3333}}
  >{\raggedright\arraybackslash}p{(\linewidth - 4\tabcolsep) * \real{0.3333}}
  >{\raggedright\arraybackslash}p{(\linewidth - 4\tabcolsep) * \real{0.3333}}@{}}
\toprule\noalign{}
\begin{minipage}[b]{\linewidth}\raggedright
Form / interpretation
\end{minipage} & \begin{minipage}[b]{\linewidth}\raggedright
Status
\end{minipage} & \begin{minipage}[b]{\linewidth}\raggedright
Relevance to this entry
\end{minipage} \\
\midrule\noalign{}
\endhead
\bottomrule\noalign{}
\endlastfoot
\emph{*rústō} \textgreater{} \emph{rost} & computed regular output from
one competing citation reconstruction & shows that this citation
reconstruction does not reach attested \emph{rust} \\
\emph{*rustaz} \textgreater{} \emph{rost} & expected regular output from
the source-supported citation reconstruction & shows that correcting the
stem class does not solve the vowel problem \\
\emph{*rústu} \textgreater{} \emph{rust} & regular high-vowel comparator
& useful negative control, but not a defensible citation-form
solution \\
\emph{rust} & attested Old English noun, normalized from \emph{rūst} /
\emph{rúst} / \emph{rust} & attested Old English form; the citation-form
development remains unexplained \\
\end{longtable}
}

\section{\texorpdfstring{wolf --- OE
\emph{wulf}}{wolf --- OE wulf}}\label{wolf-oe-wulf}

\index[iv]{02oe@\textbf{Old English}!wulf@\emph{wulf}}
\index[iv]{03pgmc@\textbf{Proto-Germanic}!wulfaz@*wúlfaz}

Derivation: \emph{*wúlfaz} yields regular \emph{wolf}; the Old English
form here is \emph{wulf} (unexplained exception).

\subsection{Derivation trace}\label{derivation-trace-144}

Proto input: \emph{*wúlfaz}

\begingroup
\setlength{\fboxsep}{6pt}

\noindent\fbox{%
\begin{minipage}{0.97\linewidth}
\small
\begin{tabularx}{\linewidth}{@{}>{\raggedright\arraybackslash}p{0.520\linewidth}>{\centering\arraybackslash}X>{\raggedright\arraybackslash}p{0.320\linewidth}@{}}
\begin{minipage}[t]{\linewidth}
\centering\textbf{Earlier Germanic changes}\par
\vspace{0.35em}
\raggedright
\centering\textbf{}\par
\raggedright
\vspace{0.2em}
\begin{tabular}{@{}>{\raggedright\arraybackslash}p{0.74\linewidth}@{\hspace{0.55em}}>{\raggedright\arraybackslash}p{0.16\linewidth}@{\hspace{0.25em}}}
\multicolumn{2}{@{}l@{}}{*Proto-West Germanic*} \\
\multicolumn{2}{@{}l@{}}{[no change]} \\
\multicolumn{2}{@{}l@{}}{*Northwest Germanic*} \\
\mbox{NWGmc U Lowering} & \emph{*wólfaz} \\
\mbox{PGmc Final Z Deletion} & \emph{*wólfa} \\
\end{tabular}
\end{minipage}
&
&
\begin{minipage}[t]{\linewidth}
\centering\textbf{Old English changes}\par
\vspace{0.35em}
\raggedright
\centering\textbf{}\par
\raggedright
\vspace{0.2em}
\begin{tabular}{@{}>{\raggedright\arraybackslash}p{0.70\linewidth}@{\hspace{0.55em}}>{\raggedright\arraybackslash}p{0.16\linewidth}@{\hspace{0.25em}}}
\multicolumn{2}{@{}l@{}}{*Old English*} \\
PWGmc Final Bare A Loss & \emph{*wólf} \\
\end{tabular}
\end{minipage}
\\
\end{tabularx}
\end{minipage}%
} \endgroup

Regular outcome: \emph{wolf}

Old English form: \emph{wulf}

\subsection{Reconstruction and comparative
evidence}\label{reconstruction-and-comparative-evidence-145}

The inherited noun is an a-stem: Kroonen gives \emph{*wulfa-}, and Ringe
and Taylor list PGmc \Recon{wulfaz} `wolf' among the inherited words
that preserve \emph{u} in Old English beside Old High German \emph{wolf}
`wolf' (\citeproc{ref-Kroonen2013}{Kroonen 2013, 638};
\citeproc{ref-RingeTaylor2014}{Ringe and Taylor 2014, 47}). Campbell
accordingly names Old English \emph{wulf} as an exception to the regular
lowering of stressed \emph{u} before a following non-high vowel
(\citeproc{ref-Campbell1959}{Campbell 1959, sect. 115}).

The older literature often notices that the exceptional words cluster
near labials. Bülbring lists \emph{full}, \emph{wulle}, and \emph{wulf}
together, but he also says that the ordinary rule still gives \emph{o}
in comparable forms such as \emph{folc} and \emph{bolt}
(\citeproc{ref-Bulbring1902}{Bülbring 1902, sect. 116}). Luick rejects a
categorical labial blocker on the same grounds and prefers a lexical or
analogical account instead (\citeproc{ref-Luick1914}{Luick 1914-\/-1940,
148}).

\subsection{Old English evidence}\label{old-english-evidence-145}

Campbell treats \emph{wulf} `wolf' as part of the exceptional \emph{u}
set (\citeproc{ref-Campbell1959}{Campbell 1959, sect. 115}).
Sievers-Brunner notes that oblique {\emph{wulfe}} `wolf (dat.)'
continues {\emph{wulfi}} `wolf (dat.)' or older {\emph{wulfai}} `wolf
(dat.)' (\citeproc{ref-SieversBrunner1965}{Brunner 1965, sect. 160}).

The surviving oblique forms do not supply a regular route back to bare
\emph{wulf} `wolf'. They belong to the same lexeme, but they do not
remove the explanatory problem presented by the citation form.

\subsection{Development to Old
English}\label{development-to-old-english-145}

Under the ordinary Northwest Germanic lowering of stressed \emph{u}
before a following non-high vowel, the citation-form input would point
toward \emph{o}-vocalism (\citeproc{ref-RingeTaylor2014}{Ringe and
Taylor 2014, 42--44}). The compact trace shows exactly that path:
\Recon{wúlfaz} `wolf' \textgreater{} \Recon{wólfaz} `wolf'
\textgreater{} \Recon{wólfa} `wolf' \textgreater{} wolf.

A high-vowel oblique input would behave differently. There the following
high vowel would block the lowering of \emph{u}, but the same
environment would also trigger i-umlaut, so the regular control result
would be {\emph{wylf}} `wolf' or \emph{wylfe}, not bare \emph{wulf}. The
attested noun therefore remains unexplained at the citation-form level.

\subsection{Expected and attested
forms}\label{expected-and-attested-forms-2}

The comparison below sets the regular inherited outcomes beside the
attested Old English noun.

{\def\LTcaptype{none} % do not increment counter
\begin{longtable}[]{@{}
  >{\raggedright\arraybackslash}p{(\linewidth - 4\tabcolsep) * \real{0.3333}}
  >{\raggedright\arraybackslash}p{(\linewidth - 4\tabcolsep) * \real{0.3333}}
  >{\raggedright\arraybackslash}p{(\linewidth - 4\tabcolsep) * \real{0.3333}}@{}}
\toprule\noalign{}
\begin{minipage}[b]{\linewidth}\raggedright
Form / interpretation
\end{minipage} & \begin{minipage}[b]{\linewidth}\raggedright
Status
\end{minipage} & \begin{minipage}[b]{\linewidth}\raggedright
Relevance to this entry
\end{minipage} \\
\midrule\noalign{}
\endhead
\bottomrule\noalign{}
\endlastfoot
\emph{*wúlfaz} \textgreater{} \emph{wolf} & computed regular output from
the citation form & shows the regular development expected from the
inherited a-stem \\
\emph{OHG wolf} & comparative regular cognate & confirms that the
\emph{o}-vocalism is the ordinary outcome \\
\emph{*wúlfi} / \emph{*wúlfis} \textgreater{} \emph{wylf} / \emph{wylfe}
& expected high-vowel control forms & shows why oblique high-vowel cells
do not solve the noun's vowel history \\
\emph{wulf} & attested Old English noun & attested Old English form; the
preservation of \emph{u} remains unexplained \\
\end{longtable}
}

\section{\texorpdfstring{wool --- OE
\emph{wull}}{wool --- OE wull}}\label{wool-oe-wull}

\index[iv]{02oe@\textbf{Old English}!wull@\emph{wull}}
\index[iv]{03pgmc@\textbf{Proto-Germanic}!wullo@*wúllō}

Derivation: \emph{*wúllō} yields regular \emph{woll}; the Old English
form here is \emph{wull} (unexplained exception).

\subsection{Derivation trace}\label{derivation-trace-145}

Proto input: \emph{*wúllō}

\begingroup
\setlength{\fboxsep}{6pt}

\noindent\fbox{%
\begin{minipage}{0.97\linewidth}
\small
\begin{tabularx}{\linewidth}{@{}>{\raggedright\arraybackslash}p{0.460\linewidth}>{\centering\arraybackslash}X>{\raggedright\arraybackslash}p{0.460\linewidth}@{}}
\begin{minipage}[t]{\linewidth}
\centering\textbf{Earlier Germanic changes}\par
\vspace{0.35em}
\raggedright
\centering\textbf{}\par
\raggedright
\vspace{0.2em}
\begin{tabular}{@{}>{\raggedright\arraybackslash}p{0.74\linewidth}@{\hspace{0.55em}}>{\raggedright\arraybackslash}p{0.16\linewidth}@{\hspace{0.25em}}}
\multicolumn{2}{@{}l@{}}{*Proto-West Germanic*} \\
\multicolumn{2}{@{}l@{}}{[no change]} \\
\multicolumn{2}{@{}l@{}}{*Northwest Germanic*} \\
\mbox{NWGmc U Lowering} & \emph{*wóllō} \\
\mbox{NWGmc Final Long O Raising} & \emph{*wóllu} \\
\end{tabular}
\end{minipage}
&
&
\begin{minipage}[t]{\linewidth}
\centering\textbf{Old English changes}\par
\vspace{0.35em}
\raggedright
\centering\textbf{}\par
\raggedright
\vspace{0.2em}
\begin{tabular}{@{}>{\raggedright\arraybackslash}p{0.74\linewidth}@{\hspace{0.55em}}>{\raggedright\arraybackslash}p{0.16\linewidth}@{\hspace{0.25em}}}
\multicolumn{2}{@{}l@{}}{*Old English*} \\
\mbox{OE High Vowel Apocope} & \emph{*wóll} \\
\end{tabular}
\end{minipage}
\\
\end{tabularx}
\end{minipage}%
} \endgroup

Regular outcome: \emph{woll}

Old English form: \emph{wull}

\subsection{Reconstruction and comparative
evidence}\label{reconstruction-and-comparative-evidence-146}

The inherited form is a feminine ō-stem \Recon{wúllō} `wool'. In the
ordinary phonological history of West Germanic, stressed \emph{u} lowers
before a following non-high vowel, so the regular Old English outcome is
an \emph{o}-form. Campbell's discussion of the parallel adjective
\emph{full}, with OHG {foll} `full' as the regular comparator, shows
that the handbooks treat this as a genuine exception cluster rather than
as a place where the rule itself is doubtful
(\citeproc{ref-Campbell1959}{Campbell 1959, sect. 115}).

Bülbring likewise lists {\emph{wulle}} `wool' among the traditional
\emph{u}-preserving exceptions (\citeproc{ref-Bulbring1902}{Bülbring
1902, sect. 116}). The comparative evidence therefore establishes two
things at once: the regular result should be \emph{woll}, and Old
English still has a lexical exception of the {\emph{wull}} `wool' /
{\emph{wulle}} `wool' type.

\subsection{Old English evidence}\label{old-english-evidence-146}

The Old English target is given here as {\emph{wull}} `wool', a
normalized lexeme form. Handbook discussion often cites {\emph{wulle}}
`wool', the feminine weak form of the noun
(\citeproc{ref-Bulbring1902}{Bülbring 1902, sect. 116}). Both point to
the same lexical item and to the same exceptional preservation of root
\emph{u}.

The OE evidence therefore does not remove the problem. It confirms that
the language has a \emph{u}-form where the regular phonology would have
produced \emph{o}.

\subsection{Development to Old
English}\label{development-to-old-english-146}

From \Recon{wúllō} `wool', the regular sequence is lowering of stressed
\emph{u} before a non-high vowel, followed by the ordinary later
reductions of the ending. The regular outcome is therefore \emph{woll}.

That regular derivation is not the attested Old English form. Luick
rejects a simple phonological rule that would protect this word alone,
and Ringe and Taylor state the larger problem plainly: we do not really
know why \emph{*u} failed to lower in forms of this sort
(\citeproc{ref-Luick1914}{Luick 1914-\/-1940, 148};
\citeproc{ref-RingeTaylor2014}{Ringe and Taylor 2014, sect. 2.3.1}).

\subsection{What remains unexplained}\label{what-remains-unexplained}

The comparison below sets the regular result beside the attested lexical
exception.

{\def\LTcaptype{none} % do not increment counter
\begin{longtable}[]{@{}
  >{\raggedright\arraybackslash}p{(\linewidth - 4\tabcolsep) * \real{0.3333}}
  >{\raggedright\arraybackslash}p{(\linewidth - 4\tabcolsep) * \real{0.3333}}
  >{\raggedright\arraybackslash}p{(\linewidth - 4\tabcolsep) * \real{0.3333}}@{}}
\toprule\noalign{}
\begin{minipage}[b]{\linewidth}\raggedright
Form
\end{minipage} & \begin{minipage}[b]{\linewidth}\raggedright
Status
\end{minipage} & \begin{minipage}[b]{\linewidth}\raggedright
Relevance to this entry
\end{minipage} \\
\midrule\noalign{}
\endhead
\bottomrule\noalign{}
\endlastfoot
\emph{*wúllō} \textgreater{} \emph{woll} & regular output & shows what
the deterministic sound laws produce \\
\emph{wull} / \emph{wulle} & attested OE exception & attested Old
English form to be recorded \\
high-vowel escape from another paradigm cell & unsupported for this noun
& rejected because the feminine ō-stem paradigm supplies no suitable
escape cell \\
\end{longtable}
}

\clearpage

\section{References}\label{references}

\backmatter

\chapter{References}\label{references-1}

\protect\phantomsection\label{refs}
\begin{CSLReferences}{1}{1}
\bibitem[\citeproctext]{ref-Adamczyk2001}
Adamczyk, Elżbieta. 2001. {`Old {English} Reflexes of {Sievers'} Law'}.
\emph{Studia Anglica Posnaniensia: An International Review of English
Studies} 36: 61--72. \url{http://hdl.handle.net/10593/18335}.

\bibitem[\citeproctext]{ref-Bammesberger1997}
Bammesberger, Alfred. 1997. {`Die Vorform von Altenglisch {h{æ}rfest}'}.
\emph{Anglia: Zeitschrift f{ü}r Englische Philologie} 115 (2): 223--30.
\url{https://doi.org/10.1515/angl.1997.115.2.223}.

\bibitem[\citeproctext]{ref-BosworthToller1898}
Bosworth, Joseph, and T. Northcote Toller. 1898. \emph{An {Anglo-Saxon}
Dictionary}. Clarendon Press.

\bibitem[\citeproctext]{ref-BrightCassidyRingler1971}
Bright, James W. 1971. \emph{Bright's {Old English} Grammar and Reader}.
3rd edn. Edited by Frederic G. Cassidy and Richard N. Ringler. Holt,
Rinehart; Winston.

\bibitem[\citeproctext]{ref-SieversBrunner1965}
Brunner, Karl. 1965. \emph{{Altenglische Grammatik nach der
angels{ä}chsischen Grammatik von Eduard Sievers}}. 3rd edn. Max
Niemeyer.

\bibitem[\citeproctext]{ref-Bulbring1902}
Bülbring, Karl D. 1902. \emph{Altenglisches Elementarbuch. 1:
Lautlehre}. Carl Winter.

\bibitem[\citeproctext]{ref-Campbell1959}
Campbell, Alistair. 1959. \emph{Old {English} Grammar}. Clarendon Press.

\bibitem[\citeproctext]{ref-ClarkHall1960}
Clark Hall, J. R. 1960. \emph{A Concise {Anglo-Saxon} Dictionary}. 4th
edn. Cambridge University Press.

\bibitem[\citeproctext]{ref-Crist2001}
Crist, Sean J. 2001. {`Conspiracy in Historical Phonology'}. Ph.\{D\}.
dissertation, University of Pennsylvania.

\bibitem[\citeproctext]{ref-Crist2002}
Crist, Sean J. 2002. \emph{An Analysis of *{z} Loss in {West Germanic}}.
Paper presented at the Linguistic Society of America Annual Meeting.

\bibitem[\citeproctext]{ref-Fulk2018}
Fulk, R. D. 2018. \emph{A Comparative Grammar of the Early {Germanic}
Languages}. Vol. 3. Studies in Germanic Linguistics. John Benjamins.

\bibitem[\citeproctext]{ref-Hogg1992}
Hogg, Richard M. 1992. \emph{A Grammar of Old English. Volume 1:
Phonology}. Blackwell.

\bibitem[\citeproctext]{ref-Hulden2009}
Hulden, Mans. 2009. {`Foma: A Finite-State Compiler and Library'}.
\emph{Proceedings of the Demonstrations Session at EACL 2009} (Athens),
29--32.

\bibitem[\citeproctext]{ref-KaplanKay1994}
Kaplan, Ronald M., and Martin Kay. 1994. {`Regular Models of
Phonological Rule Systems'}. \emph{Computational Linguistics} 20 (3):
331--78.

\bibitem[\citeproctext]{ref-KlugeSeebold2011}
Kluge, Friedrich. 2011. \emph{{Etymologisches W{ö}rterbuch der deutschen
Sprache}}. 25th edn. Edited by Elmar Seebold. De Gruyter.

\bibitem[\citeproctext]{ref-Kroonen2011}
Kroonen, Guus. 2011. \emph{The {Proto-Germanic} n-Stems: A Study in
Diachronic Morphophonology}. Vol. 18. Leiden Studies in Indo-European.
Rodopi.

\bibitem[\citeproctext]{ref-Kroonen2013}
Kroonen, Guus. 2013. \emph{Etymological Dictionary of {Proto-Germanic}}.
Vol. 11. Leiden Indo-European Etymological Dictionary Series. Brill.

\bibitem[\citeproctext]{ref-Lloyd1966}
Lloyd, Albert L. 1966. {`Is There an {a-}Umlaut of {i} in {Germanic}?'}
\emph{Language} 42 (4): 738--45. \url{https://doi.org/10.2307/411829}.

\bibitem[\citeproctext]{ref-Luick1914}
Luick, Karl. 1914-\/-1940. \emph{{Historische Grammatik der englischen
Sprache}}. Tauchnitz.

\bibitem[\citeproctext]{ref-Orel2003}
Orel, Vladimir. 2003. \emph{A Handbook of {Germanic} Etymology}. Brill.

\bibitem[\citeproctext]{ref-OsthoffBrugmann1881}
Osthoff, Hermann, and Karl Brugmann. 1878-\/-1910. \emph{{Morphologische
Untersuchungen auf dem Gebiete der indogermanischen Sprachen}}. 6 vols.
S. Hirzel.

\bibitem[\citeproctext]{ref-RingeTaylor2014}
Ringe, Don, and Ann Taylor. 2014. \emph{The Development of {Old
English}}. Vol. 2. A Linguistic History of English. Oxford University
Press.

\bibitem[\citeproctext]{ref-Seebold1970}
Seebold, Elmar. 1970. \emph{{Vergleichendes und etymologisches
W{ö}rterbuch der germanischen starken Verben}}. Mouton.

\bibitem[\citeproctext]{ref-SimsWilliams2018}
Sims-Williams, Patrick. 2018. {`Mechanising Historical Phonology'}.
\emph{Transactions of the Philological Society} 116 (3): 555--73.
\url{https://doi.org/10.1111/1467-968X.12138}.

\bibitem[\citeproctext]{ref-Sweet1953}
Sweet, Henry. 1953. \emph{An {Anglo-Saxon} Primer}. 9th edn. Clarendon
Press.

\bibitem[\citeproctext]{ref-GermanicSlavicBaltic2025}
Viredaz, Rémy. 2025. \emph{{Germanic}, {Slavic} and {Baltic}
{`Thousand'} Once More}. Unpublished manuscript (unfinished), Geneva, 20
October 2025.
\url{https://www.academia.edu/144462167/2025_Germanic_Slavic_and_Baltic_thousand_once_more_full_text_unfinished_}.

\end{CSLReferences}

\part*{Index verborum}

\addcontentsline{toc}{part}{Index verborum}

\printindex[iv]

\backmatter
\end{document}
